\documentclass[14pt,a4paper,oneside]{extbook} % Dùng extbook để nhận cỡ chữ 14pt chuẩn
\usepackage[utf8]{vietnam}
\usepackage{amsmath,amssymb,mathrsfs,fancyhdr,enumerate,multirow,makecell,venndiagram}
\usepackage{tkz-euclide,tikz-3dplot,tikz,tkz-tab}
\usetikzlibrary{arrows,calc,intersections,angles,quotes,backgrounds,shapes.geometric,patterns,through}
\usepackage{yhmath}
\usepackage{pgfplots}
\pgfplotsset{compat=1.9}
\usepgfplotslibrary{fillbetween}
\usepackage[tikz]{bclogo}
\usepackage[top=1.8cm, bottom=1.8cm, left=1.8cm, right=1.8cm, headheight=16pt]{geometry}
\usepackage[hidelinks,unicode]{hyperref}
\usepackage{qrcode}

\usepackage[many]{tcolorbox}

% Gói ex_test (Chọn [dethi] để in đề hoặc [loigiai] để in cả lời giải)
%\usepackage[loigiai]{ex_test}
\usepackage[dethi]{ex_test}

\def\colorEX{\color{blue}}

\usepackage{fontawesome}
\usepackage{framed}
\usepackage{polynom}

\colorlet{tcbcol@back}{tcbcolback}
\colorlet{tcbcol@frame}{tcbcolframe}
\newtcolorbox[auto counter]{danggg}[1]{
	fonttitle=\bfseries\upshape,
	colframe=blue,colback=blue!10,
	sharp corners, breakable, 
	title={\color{white}\faFolderO \ DẠNG~\thetcbcounter.\ #1}
	\addcontentsline{toc}{subsection}{Dạng~\thetcbcounter. #1}
}
\newenvironment{dang}[1]{
	\begin{danggg}{#1}
		\textit{Phương pháp giải:~}
	}{\end{danggg}}
\newenvironment{tomtat}{
	\begin{tcolorbox}[boxrule=1pt,arc=0mm,breakable]
	}{\end{tcolorbox}}

\newtheorem{bt}{Bài}
\newtheorem{vd}{Ví dụ}
\theorembodyfont{\it}
\newtheorem{dl}{ Định lí}
\theoremseparator{.}
\theorembodyfont{\rm}
\newtheorem{hq}{ Hệ quả}
\newtheorem{tc}{ Tính chất}
\newtheorem{dn}{ Định nghĩa}
\theoremstyle{nonumberplain}
\newtheorem{kn}{ Khái niệm}
\theoremseparator{.}
\theorembodyfont{\it}
\usepackage{stackengine}
\usepackage{scalerel}

\newcommand\dangersign[1][2ex]{%
	\renewcommand\stacktype{L}%
	\scaleto{\stackon[1.3pt]{\color{red}$\triangle$}{\fontsize{4}{0}\selectfont !}}{#1}%
}
\renewcommand{\parallel}{\mathbin{\!\raisebox{0.25pt}/\mkern-2mu\raisebox{0.25pt}/\!}}
\renewcommand{\backsim}{\ \mathord{\raisebox{0.025\baselineskip}{\rotatebox{90}{\fontfamily{qag}\selectfont S}}}\ }
\theoremseparator{}
\newtheorem{note}{{\dangersign[3ex]}}
\newtheorem{nx}{Nhận xét.}
\newtcolorbox{luuy}[1][]{colback=yellow!20,title={\bfseries Chú ý},#1}

% Tiêu đề chương / bài
\setcounter{secnumdepth}{4}
\setcounter{tocdepth}{2}
\usepackage{titlesec}
\usepackage{titletoc}

\makeatletter
\renewcommand\thesection{\@arabic\c@section}
\makeatother

\titleformat{\section}{\centering\Large\bfseries\color{blue!80!black}}{Chủ đề~\thesection.}{1ex}{
	\setcounter{vd}{0}
	\setcounter{ex}{0}
	\setcounter{bt}{0}
	\setcounter{dl}{0}
	\setcounter{dn}{0}
	\setcounter{hq}{0}
	\setcounter{tc}{0}
	\MakeUppercase
}
\titleformat{\subsection}{\large\bfseries\color{blue!70!black}}{}{0ex}{}
\usepackage[final]{pdfpages}

% Cài đặt riêng
\def\hideboxEX{}
\renewcommand{\boxEX}[2][15pt]{%
	\setbox1=\hbox{#2\strut}
	\def\sizeboxEX{#1}%
	\ \EXbox{\hideboxEX#2\strut}\ %
}

\renewcommand\arraystretch{1.4} 
\renewcommand{\tabcolsep}{4mm}

% Dòng kẻ làm bài TikZ chuẩn 8mm
\renewcommand{\dotlineEX}[1]{%
  \par\vspace{2mm}%
  \noindent%
  \begin{tikzpicture}[x=1mm, y=1mm]%
    \foreach \i in {1,...,#1} {%
      \draw[gray!60, loosely dotted, line width=0.6pt] (0, -\i*8) -- (\linewidth, -\i*8);%
    }%
  \end{tikzpicture}%
}

% Nút Casio
\renewcommand\boxed[1]{%
	\tikz[baseline=(key.base)]
	\node[%
	draw,
	fill=black!40,
	drop shadow={shadow xshift=0.25ex,shadow yshift=-0.25ex,fill=black!90,opacity=0.75},
	rectangle,minimum width=0.6cm,minimum height=0.6cm,
	rounded corners=2pt,
	inner sep=1pt,
	line width=0.5pt,
	font=\sffamily
	](key) {\color{white}#1\strut}
	;
}

\newcommand{\hoac}[1]{\left[\begin{aligned}#1\end{aligned}\right.}
\newcommand{\heva}[1]{\left\{\begin{aligned}#1\end{aligned}\right.}

\pagestyle{fancy}
\fancyhf{} % Xóa sạch toàn bộ header/footer cũ
\fancyhead[L]{\small\itshape\color{gray} Tuyển tập Đề kiểm tra Toán 5 theo từng bài}
\fancyhead[R]{\small\bfseries\color{blue!80!black} Lớp Toán Thầy Lĩnh -- loptoanthaylinh.vn}
\fancyfoot[C]{\bfseries\thepage}
\renewcommand{\headrulewidth}{0.4pt}

\newcommand{\caidat}{2}
\newcommand{\indapan}[2]{
	\begin{center}
		\color{blue!50!green} \fbox{\fontfamily{qag}\bfseries\selectfont BẢNG ĐÁP ÁN}
	\end{center}
	{\fontfamily{lmtt}\bfseries\selectfont \inputansbox[\caidat]{#1}{#2}}
}

\everymath{\color{blue}}

% -------------------------------------------------------------------------
% KHUNG TIÊU ĐỀ PHIẾU KIỂM TRA KÈM MÃ QR CODE CHO TỪNG BÀI
% Cú pháp: \TaoPhieuKiemTra{Chủ đề}{Số bài}{Tên bài}{Link QR thi online}
% -------------------------------------------------------------------------
\newcommand{\TaoPhieuKiemTra}[4]{%
  \refstepcounter{subsection}%
  \addcontentsline{toc}{subsection}{Bài #2. #3}%
  \noindent
  \begin{tcolorbox}[
    colback=white,
    colframe=blue!85!black,
    arc=2.5mm,
    boxrule=1pt,
    left=2.5mm, right=2.5mm, top=1.8mm, bottom=1.8mm
  ]
    % LOGO TRUNG TÂM & HOTLINE | TIÊU ĐỀ BÀI KIỂM TRA | MÃ QR CODE
    \begin{minipage}[c]{0.27\textwidth}
      \centering
      \begin{tcolorbox}[colback=blue!8, colframe=blue!80!black, boxrule=0.6pt, arc=1.5mm, halign=center, top=1.2mm, bottom=1.2mm, left=1mm, right=1mm]
        {\bfseries\color{blue!95!black}\small LỚP TOÁN THẦY LĨNH}\\[1.5pt]
        {\scriptsize\color{blue!70!black}\faGlobe~loptoanthaylinh.vn}\\[1.5pt]
        {\scriptsize\bfseries\color{blue!85!black}\faPhone~Hotline: 0866976878}
      \end{tcolorbox}
    \end{minipage}%
    \hfill
    \begin{minipage}[c]{0.56\textwidth}
      \centering
      {\footnotesize\bfseries\color{blue!80!black}\MakeUppercase{#1}}\\[2pt]
      {\large\bfseries\color{red!85!black}\MakeUppercase{BÀI #2: #3}}\\[2pt]
      {\scriptsize\itshape Thời gian làm bài: 35 phút -- Cấu trúc 3 phần chuẩn}
    \end{minipage}%
    \hfill
    \begin{minipage}[c]{0.14\textwidth}
      \centering
      \qrcode[height=1.35cm]{#4}\\[2pt]
      {\tiny\bfseries\color{blue!90!black}\faQrcode~Quét QR thi}
    \end{minipage}
  \end{tcolorbox}
  \vspace{1.5mm}
}

\begin{document}

\tableofcontents
\newpage

% =========================================================================
% CHỦ ĐỀ 1 - ÔN TẬP VÀ BỔ SUNG
% =========================================================================
\section{Ôn tập và bổ sung}

\TaoPhieuKiemTra{Chủ đề 1 -- Ôn tập và bổ sung}{01}{Ôn tập số tự nhiên}{https://loptoanthaylinh.vn/thi-online/toan-5/bai-01}
\textbf{Phần 1. Trắc nghiệm lựa chọn}
\setcounterpageref{ex}{0}

\begin{ex} %[5D1N1]
Số gồm $3$ chục triệu, $8$ nghìn, $2$ chục và $1$ đơn vị được viết là
\choice
{$30\,800\,021$}
{$3\,008\,021$}
{$30\,080\,021$}
{\True $30\,008\,021$}
\loigiai{
Số gồm $3$ chục triệu ($30\,000\,000$), $8$ nghìn ($8\,000$), $2$ chục ($20$) và $1$ đơn vị ($1$) được viết là $30\,008\,021$.\\
\textbf{Đáp án của câu hỏi là D}
}
\end{ex}

\begin{ex} %[5D1N1]
Số $52\,814$ được đọc là
\choice
{\True Năm mươi hai nghìn tám trăm mười bốn}
{Năm mươi hai nghìn tám trăm mười}
{Năm trăm hai mươi tám nghìn mười bốn}
{Năm mươi hai triệu tám trăm mười bốn}
\loigiai{
Số $52\,814$ có $5$ chục nghìn, $2$ nghìn, $8$ trăm, $1$ chục và $4$ đơn vị, được đọc là: Năm mươi hai nghìn tám trăm mười bốn.\\
\textbf{Đáp án của câu hỏi là A}
}
\end{ex}

\begin{ex} %[5D1H1]
Số tự nhiên gồm $8$ trăm nghìn, $2$ chục nghìn, $1$ chục và $5$ đơn vị là
\choice
{$82\,015$}
{\True $820\,015$}
{$802\,015$}
{$820\,150$}
\loigiai{
Số gồm $8$ trăm nghìn, $2$ chục nghìn, $0$ nghìn, $0$ trăm, $1$ chục và $5$ đơn vị là $820\,015$.\\
\textbf{Đáp án của câu hỏi là B}
}
\end{ex}

\begin{ex} %[5D1H1]
Chữ số $7$ trong số $863\,749$ thuộc hàng nào và có giá trị bằng bao nhiêu?
\choice
{Hàng chục, giá trị là $70$}
{Hàng nghìn, giá trị là $7\,000$}
{\True Hàng trăm, giá trị là $700$}
{Hàng chục nghìn, giá trị là $70\,000$}
\loigiai{
Trong số $863\,749$, kể từ phải sang trái: chữ số $9$ thuộc hàng đơn vị, chữ số $4$ thuộc hàng chục, chữ số $7$ thuộc hàng trăm nên có giá trị là $700$.\\
\textbf{Đáp án của câu hỏi là C}
}
\end{ex}

\begin{ex} %[5D1H1]
Số thích hợp điền vào dấu hỏi chấm để $504\,842 = 500\,000 + 4\,000 + ? + 40 + 2$ là
\choice
{$80$}
{$8\,000$}
{$8$}
{\True $800$}
\loigiai{
Phân tích số $504\,842$ theo các hàng:
\[
504\,842 = 500\,000 + 4\,000 + 800 + 40 + 2.
\]
Vậy số cần điền là $800$.\\
\textbf{Đáp án của câu hỏi là D}
}
\end{ex}

\begin{ex} %[5D1H1]
Số tự nhiên $1\,200\,324$ được phân tích thành tổng các hàng là
\choice
{$1\,000\,000 + 20\,000 + 300 + 20 + 4$}
{\True $1\,000\,000 + 200\,000 + 300 + 20 + 4$}
{$1\,000\,000 + 200\,000 + 3\,000 + 20 + 4$}
{$1\,200\,000 + 300 + 24$}
\loigiai{
Số $1\,200\,324$ gồm $1$ triệu, $2$ trăm nghìn, $3$ trăm, $2$ chục và $4$ đơn vị:
\[
1\,200\,324 = 1\,000\,000 + 200\,000 + 300 + 20 + 4.
\]
\textbf{Đáp án của câu hỏi là B}
}
\end{ex}

\begin{ex} %[5D1H1]
Ba số chẵn liên tiếp được viết vào $3$ chiếc mũ của Việt, Nam và Rô-bốt. Rô-bốt nhìn thấy số viết trên mũ của Việt và Nam lần lượt là $2\,032$ và $2\,028$. Số trên mũ của Rô-bốt là
\choice
{\True $2\,030$}
{$2\,029$}
{$2\,034$}
{$2\,026$}
\loigiai{
Ba số chẵn liên tiếp cách nhau $2$ đơn vị. Số nằm giữa $2\,028$ và $2\,032$ là $2\,030$. Vậy chiếc mũ Rô-bốt đang đội mang số $2\,030$.\\
\textbf{Đáp án của câu hỏi là A}
}
\end{ex}

\begin{ex} %[5D1H1]
Cho độ cao của ba đỉnh núi: Đỉnh E-vơ-rét ($8\,848\text{ m}$), Đỉnh núi Bà Đen ($986\text{ m}$), Đỉnh Ô-lim-pót Mon ($25\,000\text{ m}$). Đỉnh núi nào cao nhất?
\choice
{Đỉnh núi Bà Đen}
{Đỉnh E-vơ-rét}
{Cả ba đỉnh cao bằng nhau}
{\True Đỉnh Ô-lim-pót Mon}
\loigiai{
So sánh các số: $986 < 8\,848 < 25\,000$. Vậy đỉnh núi cao nhất là đỉnh Ô-lim-pót Mon ($25\,000\text{ m}$).\\
\textbf{Đáp án của câu hỏi là D}
}
\end{ex}

\begin{ex} %[5D1H1]
Trong ba bộ đồ chơi: Ô tô ($195\,000$ đồng), Ngôi nhà búp bê ($1\,175\,000$ đồng), Máy chơi game ($2\,000\,000$ đồng). Bộ đồ chơi nào có giá tiền thấp nhất?
\choice
{\True Ô tô}
{Ngôi nhà búp bê}
{Máy chơi game}
{Ô tô và Ngôi nhà búp bê}
\loigiai{
So sánh giá tiền: $195\,000 < 1\,175\,000 < 2\,000\,000$. Vậy bộ đồ chơi có giá thấp nhất là ô tô ($195\,000$ đồng).\\
\textbf{Đáp án của câu hỏi là A}
}
\end{ex}

\begin{ex} %[5D1H1]
Số sản phẩm bán được của một cửa hàng trong bốn năm là: Năm 2020 bán $2\,873$ sản phẩm; Năm 2021 bán $2\,837$ sản phẩm; Năm 2022 bán $3\,293$ sản phẩm; Năm 2023 bán $3\,018$ sản phẩm. Năm cửa hàng bán được nhiều sản phẩm nhất là
\choice
{Năm 2020}
{Năm 2021}
{\True Năm 2022}
{Năm 2023}
\loigiai{
So sánh số sản phẩm: $2\,837 < 2\,873 < 3\,018 < 3\,293$. Vậy năm $2022$ bán được nhiều sản phẩm nhất ($3\,293$ sản phẩm).\\
\textbf{Đáp án của câu hỏi là C}
}
\end{ex}

\begin{ex} %[5D1H1]
Doanh thu thực tế trong một ngày của một cửa hàng là $2\,545\,000$ đồng. Cuối ngày, người bán hàng nói: ``Doanh thu cả ngày của cửa hàng được khoảng $2\,500\,000$ đồng''. Người bán hàng đã làm tròn số tiền đến hàng nào?
\choice
{Hàng chục nghìn}
{\True Hàng trăm nghìn}
{Hàng triệu}
{Hàng nghìn}
\loigiai{
Số $2\,545\,000$ có chữ số hàng chục nghìn là $4 < 5$, khi làm tròn đến hàng trăm nghìn ta giữ nguyên chữ số hàng trăm nghìn là $5$ và các chữ số phía sau chuyển thành $0$, được $2\,500\,000$.\\
\textbf{Đáp án của câu hỏi là B}
}
\end{ex}

\begin{ex} %[5D1V1]
Khi làm tròn số $2\,545\,000$ đến hàng chục nghìn, ta được số nào sau đây?
\choice
{$2\,540\,000$}
{$2\,500\,000$}
{\True $2\,550\,000$}
{$2\,600\,000$}
\loigiai{
Xét hàng nghìn của số $2\,545\,000$ có chữ số là $5 \ge 5$, do đó khi làm tròn đến hàng chục nghìn ta cộng thêm $1$ vào chữ số hàng chục nghìn ($4 + 1 = 5$), ta được số $2\,550\,000$.\\
\textbf{Đáp án của câu hỏi là C}
}
\end{ex}

\textbf{Phần 2. Trắc nghiệm đúng sai}
\setcounterpageref{ex}{0}

\begin{ex} %[5D1H1]
Cho hai số tự nhiên $A = 26\,400\,500$ và $B = 30\,008\,021$. Xét tính đúng sai của các khẳng định sau:
\choiceTF
{\True Số $A$ được đọc là: Hai mươi sáu triệu bốn trăm nghìn năm trăm.}
{Trong số $B$, chữ số $8$ thuộc hàng trăm nghìn.}
{\True Số $A$ được viết thành tổng theo các hàng là $20\,000\,000 + 6\,000\,000 + 400\,000 + 500$.}
{\True Hiệu giữa giá trị của chữ số $3$ và giá trị của chữ số $8$ trong số $B$ là $29\,992\,000$.}
\loigiai{
\begin{itemize}
    \item Số $A = 26\,400\,500$ đọc là: Hai mươi sáu triệu bốn trăm nghìn năm trăm.
    \item Trong số $B = 30\,008\,021$, chữ số $8$ thuộc hàng nghìn (lớp nghìn), có giá trị là $8\,000$.
    \item Phân tích cấu tạo số $A$: $26\,400\,500 = 20\,000\,000 + 6\,000\,000 + 400\,000 + 500$.
    \item Trong số $B$: Chữ số $3$ thuộc hàng chục triệu có giá trị $30\,000\,000$; chữ số $8$ có giá trị $8\,000$. Hiệu giá trị là $30\,000\,000 - 8\,000 = 29\,992\,000$.
\end{itemize}
\begin{itemize}
    \item a) \textbf{Đúng}.
    \item b) \textbf{Sai}: Chữ số $8$ thuộc hàng nghìn.
    \item c) \textbf{Đúng}.
    \item d) \textbf{Đúng}.
\end{itemize}
}
\end{ex}

\begin{ex} %[5D1H1]
Số sản phẩm bán được của một cửa hàng trong bốn năm được thống kê như sau: Năm 2020 bán $2\,873$ sản phẩm; Năm 2021 bán $2\,837$ sản phẩm; Năm 2022 bán $3\,293$ sản phẩm; Năm 2023 bán $3\,018$ sản phẩm. Doanh thu một ngày là $2\,545\,000$ đồng. Xét tính đúng sai của các khẳng định sau:
\choiceTF
{\True Năm $2022$ cửa hàng bán được nhiều sản phẩm nhất.}
{Năm $2021$ cửa hàng bán được nhiều sản phẩm hơn năm $2020$.}
{\True Làm tròn doanh thu $2\,545\,000$ đồng đến hàng trăm nghìn ta được $2\,500\,000$ đồng.}
{Làm tròn doanh thu $2\,545\,000$ đồng đến hàng chục nghìn ta được $2\,540\,000$ đồng.}
\loigiai{
\begin{itemize}
    \item So sánh số sản phẩm: $2\,837 < 2\,873 < 3\,018 < 3\,293$. Vậy năm $2022$ bán nhiều sản phẩm nhất ($3\,293$ sản phẩm).
    \item Năm $2021$ bán $2\,837$ sản phẩm, năm $2020$ bán $2\,873$ sản phẩm. Vì $2\,837 < 2\,873$ nên năm $2021$ bán ít hơn năm $2020$.
    \item Chữ số hàng chục nghìn của $2\,545\,000$ là $4 < 5$, nên làm tròn đến hàng trăm nghìn được $2\,500\,000$ đồng.
    \item Chữ số hàng nghìn của $2\,545\,000$ là $5 \ge 5$, nên làm tròn đến hàng chục nghìn phải được $2\,550\,000$ đồng.
\end{itemize}
\begin{itemize}
    \item a) \textbf{Đúng}.
    \item b) \textbf{Sai}: Năm 2021 bán ít hơn năm 2020.
    \item c) \textbf{Đúng}.
    \item d) \textbf{Sai}: Kết quả làm tròn đến hàng chục nghìn là $2\,550\,000$ đồng.
\end{itemize}
}
\end{ex}

\textbf{Phần 3. Trả lời ngắn}
\setcounterpageref{ex}{0}

\begin{ex} %[5D1H1]
Tìm số thích hợp điền vào dấu hỏi chấm để hoàn thành đẳng thức sau:
\[
1\,730\,539 = 1\,000\,000 + ? + 30\,000 + 500 + 30 + 9.
\]
\shortans{700000}
\loigiai{
\textbf{Đáp số:} $700000$\\
Chữ số $7$ trong số $1\,730\,539$ ở hàng trăm nghìn, do đó có giá trị là $700\,000$.
}
\end{ex}

\begin{ex} %[5D1H1]
Ba bạn Việt, Nam và Rô-bốt đội $3$ chiếc mũ mang $3$ số chẵn liên tiếp. Rô-bốt nhìn thấy mũ của Nam ghi số $2\,028$ và mũ của Việt ghi số $2\,032$. Hỏi chiếc mũ của Rô-bốt ghi số nào?
\shortans{2030}
\loigiai{
\textbf{Đáp số:} $2030$\\
Ba số chẵn liên tiếp là $2\,028$; $2\,030$; $2\,032$. Vậy số trên mũ Rô-bốt là $2030$.
}
\end{ex}

\begin{ex} %[5D1H1]
Rô-bốt có thói quen viết các số biểu diễn ngày, tháng, năm liên tiếp nhau để được một số tự nhiên có nhiều chữ số. Ví dụ: ngày 30 tháng 4 năm 1975 viết thành $3\,041\,975$. Hỏi với ngày Nhà giáo Việt Nam năm 2024 (ngày 20 tháng 11 năm 2024), số tự nhiên Rô-bốt viết được là bao nhiêu?
\shortans{20112024}
\loigiai{
\textbf{Đáp số:} $20112024$\\
Ngày $20$, tháng $11$, năm $2024$ ghép liên tiếp thành số $20\,112\,024$.
}
\end{ex}

\begin{ex} %[5D1H1]
Trong số $20\,112\,024$, tổng giá trị của tất cả các chữ số $2$ có trong số đó bằng bao nhiêu?
\shortans{20002020}
\loigiai{
\textbf{Đáp số:} $20002020$\\
Các chữ số $2$ có giá trị lần lượt là:
\begin{itemize}
    \item Hàng chục triệu: $20\,000\,000$.
    \item Hàng nghìn: $2\,000$.
    \item Hàng chục: $20$.
\end{itemize}
Tổng giá trị của các chữ số $2$ là: $20\,000\,000 + 2\,000 + 20 = 20\,002\,020$.
}
\end{ex}

\begin{ex} %[5D1H1]
Doanh thu thực tế cả ngày của một cửa hàng là $2\,545\,000$ đồng. Hãy làm tròn doanh thu này đến hàng chục nghìn.
\shortans{2550000}
\loigiai{
\textbf{Đáp số:} $2550000$\\
Chữ số hàng nghìn là $5$, làm tròn lên hàng chục nghìn ta được $2\,550\,000$.
}
\end{ex}

\begin{ex} %[5D1V1]
Rô-bốt lập số $863\,749$ bằng các tấm thẻ: $8$, $6$, $3$, $7$, $4$, $9$. Rô-bốt đổi chỗ đúng $2$ tấm thẻ bất kỳ để nhận được một số lẻ lớn nhất có thể. Hỏi số lẻ lớn nhất đó là bao nhiêu?
\shortans{873649}
\loigiai{
\textbf{Đáp số:} $873649$\\
Số ban đầu là $863\,749$ (có chữ số tận cùng là $9$ nên là số lẻ).
\begin{itemize}
    \item Để số tạo thành là số lẻ lớn nhất, chữ số tận cùng phải là chữ số lẻ ($3$, $7$ hoặc $9$).
    \item Nếu đổi thẻ $9$ (hàng đơn vị) với thẻ $8$ (hàng trăm nghìn), ta được số $963\,748$, có tận cùng là $8$ (số chẵn, không thỏa mãn).
    \item Nếu đổi thẻ $9$ với $7$ ta được số $863\,947$ (lẻ). Nếu đổi thẻ $9$ với $3$ ta được $869\,743$ (lẻ).
    \item Nếu giữ thẻ $9$ ở hàng đơn vị và đổi chỗ thẻ $6$ (hàng chục nghìn) với thẻ $7$ (hàng trăm), ta được số $873\,649$ (lẻ).
    \item So sánh các số lẻ tạo được: $873\,649 > 869\,743 > 867\,349 > 863\,947 > 863\,749$.
\end{itemize}
Vậy số lẻ lớn nhất có thể nhận được là $873\,649$ (bằng cách đổi chỗ thẻ số $6$ và thẻ số $7$).
}
\end{ex}

\newpage

\TaoPhieuKiemTra{Chủ đề 1 -- Ôn tập và bổ sung}{02}{Ôn tập các phép tính với số tự nhiên}{https://loptoanthaylinh.vn/thi-online/toan-5/bai-02}
\textbf{Phần 1. Trắc nghiệm lựa chọn}
\setcounterpageref{ex}{0}

\begin{ex} %[5D1N2]
Kết quả của phép tính $7\,318 + 3\,191$ là
\choice
{$10\,409$}
{$10\,519$}
{$10\,499$}
{\True $10\,509$}
\loigiai{
Đặt tính và tính: $7\,318 + 3\,191 = 10\,509$.\\
\textbf{Đáp án của câu hỏi là D}
}
\end{ex}

\begin{ex} %[5D1N2]
Kết quả của phép trừ $83\,500 - 28\,150$ là
\choice
{\True $55\,350$}
{$55\,450$}
{$65\,350$}
{$55\,250$}
\loigiai{
Thực hiện phép trừ: $83\,500 - 28\,150 = 55\,350$.\\
\textbf{Đáp án của câu hỏi là A}
}
\end{ex}

\begin{ex} %[5D1H2]
Kết quả của phép nhân $307 \times 15$ là
\choice
{$4\,505$}
{$4\,595$}
{\True $4\,605$}
{$4\,615$}
\loigiai{
Thực hiện phép nhân: $307 \times 15 = 4\,605$.\\
\textbf{Đáp án của câu hỏi là C}
}
\end{ex}

\begin{ex} %[5D1H2]
Thương của phép chia $4\,488 : 34$ là
\choice
{$122$}
{$142$}
{$134$}
{\True $132$}
\loigiai{
Đặt tính chia: $4\,488 : 34 = 132$.\\
\textbf{Đáp án của câu hỏi là D}
}
\end{ex}

\begin{ex} %[5D1H2]
Biểu thức nào dưới đây có giá trị bằng biểu thức $2 \times 75 + 2 \times 25$?
\choice
{\True $2 \times (75 + 25)$}
{$2 \times (75 - 25)$}
{$2 + (75 \times 25)$}
{$(2 + 2) \times (75 + 25)$}
\loigiai{
Áp dụng tính chất nhân một số với một tổng: $2 \times 75 + 2 \times 25 = 2 \times (75 + 25)$.\\
\textbf{Đáp án của câu hỏi là A}
}
\end{ex}

\begin{ex} %[5D1H2]
Giá trị của biểu thức $3\,713 - 200 \times 5$ là
\choice
{$17\,565$}
{$2\,613$}
{\True $2\,713$}
{$3\,513$}
\loigiai{
Thực hiện phép nhân trước: $200 \times 5 = 1\,000$. Sau đó trừ: $3\,713 - 1\,000 = 2\,713$.\\
\textbf{Đáp án của câu hỏi là C}
}
\end{ex}

\begin{ex} %[5D1H2]
Giá trị của biểu thức $1\,500 + (750 + 250) : 2$ là
\choice
{$1\,000$}
{\True $2\,000$}
{$2\,500$}
{$1\,750$}
\loigiai{
Thực hiện trong ngoặc trước: $750 + 250 = 1\,000$.\\
Sau đó chia: $1\,000 : 2 = 500$.\\
Cuối cùng cộng: $1\,500 + 500 = 2\,000$.\\
\textbf{Đáp án của câu hỏi là B}
}
\end{ex}

\begin{ex} %[5D1H2]
Tính nhanh biểu thức $25 \times 99 \times 4$ ta được kết quả là
\choice
{\True $9\,900$}
{$990$}
{$99\,000$}
{$990\,000$}
\loigiai{
Sử dụng tính chất giao hoán và kết hợp: $(25 \times 4) \times 99 = 100 \times 99 = 9\,900$.\\
\textbf{Đáp án của câu hỏi là A}
}
\end{ex}

\begin{ex} %[5D1H2]
Tính giá trị của biểu thức $2\,025 \times 17 + 83 \times 2\,025$ bằng cách thuận tiện nhất:
\choice
{$202\,500\,000$}
{$2\,025\,000$}
{\True $202\,500$}
{$20\,250$}
\loigiai{
Áp dụng tính chất phân phối: $2\,025 \times (17 + 83) = 2\,025 \times 100 = 202\,500$.\\
\textbf{Đáp án của câu hỏi là C}
}
\end{ex}

\begin{ex} %[5D1V2]
Mai mua hai gói bim bim hết số tiền là $18\,000$ đồng. Biết gói bim bim cua đắt hơn gói bim bim mực là $4\,000$ đồng. Giá tiền của gói bim bim mực là
\choice
{$11\,000$ đồng}
{\True $7\,000$ đồng}
{$14\,000$ đồng}
{$9\,000$ đồng}
\loigiai{
Đây là bài toán tìm hai số khi biết tổng và hiệu:\\
Giá tiền gói bim bim mực (số bé) là: $(18\,000 - 4\,000) : 2 = 7\,000$ đồng.\\
\textbf{Đáp án của câu hỏi là B}
}
\end{ex}

\begin{ex} %[5D1V2]
Rô-bốt bán bốn bức tranh với giá tiền lần lượt là: $85\,500$ đồng, $150\,000$ đồng, $425\,000$ đồng và $55\,500$ đồng. Trung bình mỗi bức tranh Rô-bốt bán được giá bao nhiêu tiền?
\choice
{$178\,000$ đồng}
{$180\,000$ đồng}
{$189\,000$ đồng}
{\True $179\,000$ đồng}
\loigiai{
Tổng số tiền bán $4$ bức tranh là:
\[
85\,500 + 150\,000 + 425\,000 + 55\,500 = 716\,000\text{ (đồng)}.
\]
Trung bình mỗi bức tranh có giá: $716\,000 : 4 = 179\,000$ đồng.\\
\textbf{Đáp án của câu hỏi là D}
}
\end{ex}

\begin{ex} %[5D1V2]
Bác Ba có $525\text{ kg}$ gạo chia đều vào các túi, mỗi túi $15\text{ kg}$. Hỏi $7$ túi gạo như vậy chứa tất cả bao nhiêu ki-lô-gam gạo?
\choice
{$75\text{ kg}$}
{\True $105\text{ kg}$}
{$120\text{ kg}$}
{$150\text{ kg}$}
\loigiai{
Số ki-lô-gam gạo trong $7$ túi là: $15 \times 7 = 105\text{ (kg)}$.\\
\textbf{Đáp án của câu hỏi là B}
}
\end{ex}

\textbf{Phần 2. Trắc nghiệm đúng sai}
\setcounterpageref{ex}{0}

\begin{ex} %[5D1H2]
Cho các biểu thức: $A = 25 \times 99 \times 4$, $B = 2\,025 \times 17 + 83 \times 2\,025$ và $C = 3\,713 - 200 \times 5$. Xét tính đúng sai của các khẳng định sau:
\choiceTF
{\True Giá trị của biểu thức $A$ bằng $9\,900$.}
{Giá trị của biểu thức $B$ là $20\,250$.}
{\True Biểu thức $C$ có giá trị là $2\,713$.}
{\True Tổng giá trị của biểu thức $A$ và biểu thức $C$ là $12\,613$.}
\loigiai{
\begin{itemize}
    \item $A = (25 \times 4) \times 99 = 100 \times 99 = 9\,900$.
    \item $B = 2\,025 \times (17 + 83) = 2\,025 \times 100 = 202\,500 \neq 20\,250$.
    \item $C = 3\,713 - 1\,000 = 2\,713$.
    \item $A + C = 9\,900 + 2\,713 = 12\,613$.
\end{itemize}
\begin{itemize}
    \item a) \textbf{Đúng}.
    \item b) \textbf{Sai}: $B = 202\,500$.
    \item c) \textbf{Đúng}.
    \item d) \textbf{Đúng}.
\end{itemize}
}
\end{ex}

\begin{ex} %[5D1V2]
Bác Ba chia đều $525\text{ kg}$ gạo vào các túi nhỏ, mỗi túi chứa $15\text{ kg}$. Bác bán mỗi túi gạo đó với giá $250\,000$ đồng. Xét tính đúng sai của các khẳng định sau:
\choiceTF
{\True Bác Ba chia được tất cả $35$ túi gạo.}
{$7$ túi gạo như vậy chứa $115\text{ kg}$ gạo.}
{\True Tổng số tiền bác Ba thu được khi bán hết số gạo là $8\,750\,000$ đồng.}
{Nếu bác Ba chỉ bán $20$ túi gạo thì số tiền thu được là $4\,500\,000$ đồng.}
\loigiai{
\begin{itemize}
    \item Tổng số túi gạo bác Ba chia được là: $525 : 15 = 35$ (túi).
    \item $7$ túi gạo chứa số ki-lô-gam là: $15 \times 7 = 105\text{ kg} \neq 115\text{ kg}$.
    \item Bán hết $35$ túi thu được: $35 \times 250\,000 = 8\,750\,000$ (đồng).
    \item Bán $20$ túi thu được: $20 \times 250\,000 = 5\,000\,000\text{ (đồng)} \neq 4\,500\,000$ đồng.
\end{itemize}
\begin{itemize}
    \item a) \textbf{Đúng}.
    \item b) \textbf{Sai}: $7$ túi chứa $105\text{ kg}$.
    \item c) \textbf{Đúng}.
    \item d) \textbf{Sai}: Thu được $5\,000\,000$ đồng.
\end{itemize}
}
\end{ex}

\textbf{Phần 3. Trả lời ngắn}
\setcounterpageref{ex}{0}

\begin{ex} %[5D1H2]
Tính tổng tất cả các số sau bằng cách thuận tiện: $400, 500, 300, 600, 280, 510, 100, 500, 490, 900, 720, 700$.
\shortans{6000}
\loigiai{
\textbf{Đáp số:} $6000$\\
Nhóm các cặp số tròn nghìn và tròn trăm:
\begin{align*}
S &= (400 + 600) + (500 + 500) + (300 + 700) + (100 + 900) + (280 + 720) + (510 + 490) \\
&= 1\,000 + 1\,000 + 1\,000 + 1\,000 + 1\,000 + 1\,000 = 6\,000.
\end{align*}
}
\end{ex}

\begin{ex} %[5D1H2]
Tìm chữ số thích hợp thay cho dấu hỏi chấm trong phép trừ sau:
\[
\begin{array}{r@{}l@{\quad}l}
  & 6\,0\,6\,1\,8\,? \\
- & 4\,3\,?\,1\,4\,9 \\
\hline
  & 1\,7\,1\,?\,3\,3
\end{array}
\]
Biết các dấu hỏi chấm từ trên xuống dưới, từ trái sang phải lần lượt là các chữ số. Hãy cho biết chữ số ở hàng đơn vị của số bị trừ (dấu hỏi chấm ở dòng thứ nhất).
\shortans{2}
\loigiai{
\textbf{Đáp số:} $2$\\
Ở hàng đơn vị: $? - 9 = 3 \Rightarrow ? = 12$ (viết $3$ nhớ $1$), vậy chữ số cần tìm là $2$.
}
\end{ex}

\begin{ex} %[5D1V2]
Mai mua hai gói bim bim hết số tiền là $18\,000$ đồng. Trong đó gói bim bim cua đắt hơn gói bim bim mực là $4\,000$ đồng. Hỏi giá tiền gói bim bim cua là bao nhiêu đồng?
\shortans{11000}
\loigiai{
\textbf{Đáp số:} $11000$\\
Giá tiền gói bim bim cua là: $(18\,000 + 4\,000) : 2 = 11\,000$ (đồng).
}
\end{ex}

\begin{ex} %[5D1H2]
Bác Ba chia đều $525\text{ kg}$ gạo vào các túi, mỗi túi $15\text{ kg}$. Hỏi bác Ba chia được tất cả bao nhiêu túi gạo?
\shortans{35}
\loigiai{
\textbf{Đáp số:} $35$\\
Số túi gạo chia được là: $525 : 15 = 35$ (túi).
}
\end{ex}

\begin{ex} %[5D1V2]
Bác Ba chia đều $525\text{ kg}$ gạo vào các túi, mỗi túi $15\text{ kg}$ và bán mỗi túi với giá $250\,000$ đồng. Hỏi sau khi bán hết số gạo đó, bác Ba thu được tất cả bao nhiêu nghìn đồng?
\shortans{8750}
\loigiai{
\textbf{Đáp số:} $8750$\\
Số túi gạo là $525 : 15 = 35$ túi. Số tiền thu được là $35 \times 250\,000 = 8\,750\,000$ đồng, tức $8\,750$ nghìn đồng.
}
\end{ex}

\begin{ex} %[5D1V2]
Rô-bốt bán bốn bức tranh với các mức giá: $85\,500$ đồng, $150\,000$ đồng, $425\,000$ đồng và $55\,500$ đồng. Hỏi trung bình mỗi bức tranh có giá bao nhiêu nghìn đồng?
\shortans{179}
\loigiai{
\textbf{Đáp số:} $179$\\
Tổng số tiền bán $4$ bức tranh là: $85\,500 + 150\,000 + 425\,000 + 55\,500 = 716\,000$ đồng.\\
Trung bình mỗi bức tranh có giá: $716\,000 : 4 = 179\,000$ đồng, tức $179$ nghìn đồng.
}
\end{ex}

\newpage

\TaoPhieuKiemTra{Chủ đề 1 -- Ôn tập và bổ sung}{03}{Ôn tập phân số}{https://loptoanthaylinh.vn/thi-online/toan-5/bai-03}
\textbf{Phần 1. Trắc nghiệm lựa chọn}
\setcounterpageref{ex}{0}

\begin{ex} %[5D1N3]
Phân số nào dưới đây là phân số tối giản?
\choice
{\True $\dfrac{7}{9}$}
{$\dfrac{16}{18}$}
{$\dfrac{15}{12}$}
{$\dfrac{25}{30}$}
\loigiai{
Phân số $\dfrac{7}{9}$ có tử số và mẫu số không cùng chia hết cho số tự nhiên nào lớn hơn $1$, nên $\dfrac{7}{9}$ là phân số tối giản.\\
\textbf{Đáp án của câu hỏi là A}
}
\end{ex}

\begin{ex} %[5D1H3]
Rút gọn phân số $\dfrac{18}{24}$ về dạng tối giản ta được kết quả là
\choice
{$\dfrac{9}{12}$}
{$\dfrac{6}{8}$}
{\True $\dfrac{3}{4}$}
{$\dfrac{2}{3}$}
\loigiai{
Chia cả tử số và mẫu số cho $6$: $\dfrac{18}{24} = \dfrac{18 : 6}{24 : 6} = \dfrac{3}{4}$.\\
\textbf{Đáp án của câu hỏi là C}
}
\end{ex}

\begin{ex} %[5D1H3]
Rút gọn phân số $\dfrac{56}{42}$ về dạng phân số tối giản ta được
\choice
{$\dfrac{28}{21}$}
{\True $\dfrac{4}{3}$}
{$\dfrac{8}{6}$}
{$\dfrac{7}{6}$}
\loigiai{
Chia cả tử số và mẫu số cho $14$: $\dfrac{56}{42} = \dfrac{56 : 14}{42 : 14} = \dfrac{4}{3}$.\\
\textbf{Đáp án của câu hỏi là B}
}
\end{ex}

\begin{ex} %[5D1H3]
Số thích hợp điền vào dấu hỏi chấm để $\dfrac{24}{40} = \dfrac{12}{?} = \dfrac{3}{5}$ là
\choice
{\True $20$}
{$15$}
{$25$}
{$30$}
\loigiai{
Ta có $\dfrac{24}{40} = \dfrac{24 : 2}{40 : 2} = \dfrac{12}{20}$. Vậy số cần điền vào dấu hỏi chấm là $20$.\\
\textbf{Đáp án của câu hỏi là A}
}
\end{ex}

\begin{ex} %[5D1H3]
Quy đồng mẫu số hai phân số $\dfrac{9}{5}$ và $\dfrac{21}{40}$ với mẫu số chung nhỏ nhất là $40$, ta được hai phân số là
\choice
{$\dfrac{36}{40}$ và $\dfrac{21}{40}$}
{$\dfrac{45}{40}$ và $\dfrac{21}{40}$}
{$\dfrac{72}{40}$ và $\dfrac{42}{40}$}
{\True $\dfrac{72}{40}$ và $\dfrac{21}{40}$}
\loigiai{
Vì $40 : 5 = 8$ nên $\dfrac{9}{5} = \dfrac{9 \times 8}{5 \times 8} = \dfrac{72}{40}$. Giữ nguyên phân số $\dfrac{21}{40}$.\\
\textbf{Đáp án của câu hỏi là D}
}
\end{ex}

\begin{ex} %[5D1H3]
Tính giá trị của biểu thức $\dfrac{6 \times 12 \times 17}{17 \times 4 \times 20}$ bằng cách rút gọn ta được
\choice
{$\dfrac{3}{5}$}
{\True $\dfrac{9}{10}$}
{$\dfrac{18}{20}$}
{$1$}
\loigiai{
Ta có: $\dfrac{6 \times 12 \times 17}{17 \times 4 \times 20} = \dfrac{6 \times 3 \times 1}{1 \times 1 \times 20} = \dfrac{18}{20} = \dfrac{9}{10}$.\\
\textbf{Đáp án của câu hỏi là B}
}
\end{ex}

\begin{ex} %[5D1N3]
Phân số nào dưới đây lớn hơn $1$?
\choice
{$\dfrac{16}{17}$}
{$\dfrac{90}{100}$}
{$\dfrac{120}{125}$}
{\True $\dfrac{31}{29}$}
\loigiai{
Phân số có tử số lớn hơn mẫu số thì lớn hơn $1$. Vì $31 > 29$ nên $\dfrac{31}{29} > 1$.\\
\textbf{Đáp án của câu hỏi là D}
}
\end{ex}

\begin{ex} %[5D1H3]
Phân số nào dưới đây bé hơn phân số $\dfrac{4}{7}$?
\choice
{$\dfrac{5}{7}$}
{$\dfrac{20}{35}$}
{\True $\dfrac{12}{28}$}
{$\dfrac{13}{21}$}
\loigiai{
Ta có $\dfrac{4}{7} = \dfrac{16}{28}$. Phân số $\dfrac{12}{28} < \dfrac{16}{28} = \dfrac{4}{7}$.\\
(Lưu ý: $\dfrac{20}{35} = \dfrac{4}{7}$; $\dfrac{13}{21} > \dfrac{12}{21} = \dfrac{4}{7}$).\\
\textbf{Đáp án của câu hỏi là C}
}
\end{ex}

\begin{ex} %[5D1H3]
Phân số nào dưới đây lớn hơn phân số $\dfrac{5}{8}$?
\choice
{$\dfrac{10}{16}$}
{\True $\dfrac{25}{32}$}
{$\dfrac{13}{24}$}
{$\dfrac{3}{8}$}
\loigiai{
Quy đồng về mẫu số chung $32$: $\dfrac{5}{8} = \dfrac{20}{32}$. Ta thấy $\dfrac{25}{32} > \dfrac{20}{32}$.\\
\textbf{Đáp án của câu hỏi là B}
}
\end{ex}

\begin{ex} %[5D1H3]
Dấu thích hợp điền vào chỗ chấm để $\dfrac{7}{6} \dots \dfrac{29}{24}$ là
\choice
{$>$}
{$=$}
{\True $<$}
{Không so sánh được}
\loigiai{
Quy đồng mẫu số: $\dfrac{7}{6} = \dfrac{28}{24}$. Vì $\dfrac{28}{24} < \dfrac{29}{24}$ nên $\dfrac{7}{6} < \dfrac{29}{24}$.\\
\textbf{Đáp án của câu hỏi là C}
}
\end{ex}

\begin{ex} %[5D1V3]
Số học sinh khối 5 tham gia các môn thể thao ở Câu lạc bộ như sau: Cầu lông có $\dfrac{1}{3}$ số học sinh khối 5, Bóng đá có $\dfrac{2}{5}$, Võ có $\dfrac{1}{5}$, Bóng rổ có $\dfrac{4}{15}$. Môn thể thao được học sinh khối 5 tham gia nhiều nhất là
\choice
{Cầu lông}
{Võ}
{Bóng rổ}
{\True Bóng đá}
\loigiai{
Quy đồng mẫu số các phân số về mẫu số chung là $15$:\\
Cầu lông: $\dfrac{1}{3} = \dfrac{5}{15}$; Bóng đá: $\dfrac{2}{5} = \dfrac{6}{15}$; Võ: $\dfrac{1}{5} = \dfrac{3}{15}$; Bóng rổ: $\dfrac{4}{15}$.\\
Vì $\dfrac{6}{15} > \dfrac{5}{15} > \dfrac{4}{15} > \dfrac{3}{15}$ nên môn Bóng đá có số học sinh tham gia nhiều nhất.\\
\textbf{Đáp án của câu hỏi là D}
}
\end{ex}

\begin{ex} %[5D1V3]
Tìm số tự nhiên thích hợp điền vào dấu hỏi chấm để $\dfrac{3}{8} < \dfrac{?}{8} < \dfrac{4}{7}$.
\choice
{\True $4$}
{$2$}
{$3$}
{$5$}
\loigiai{
Quy đồng mẫu số chung là $56$:\\
$\dfrac{3}{8} = \dfrac{21}{56}$; $\dfrac{4}{7} = \dfrac{32}{56}$; $\dfrac{?}{8} = \dfrac{7 \times ?}{56}$.\\
Ta có: $21 < 7 \times ? < 32 \Rightarrow 7 \times ? = 28 \Rightarrow ? = 4$.\\
\textbf{Đáp án của câu hỏi là A}
}
\end{ex}

\textbf{Phần 2. Trắc nghiệm đúng sai}
\setcounterpageref{ex}{0}

\begin{ex} %[5D1H3]
Cho bốn phân số: $\dfrac{3}{4}$, $\dfrac{7}{12}$, $\dfrac{20}{16}$ và $\dfrac{19}{24}$. Xét tính đúng sai của các khẳng định sau:
\choiceTF
{\True Phân số $\dfrac{20}{16}$ sau khi rút gọn tối giản bằng $\dfrac{5}{4}$.}
{Phân số $\dfrac{7}{12}$ lớn hơn phân số $\dfrac{3}{4}$.}
{\True Trong bốn phân số đã cho, phân số lớn nhất là $\dfrac{20}{16}$.}
{\True Sắp xếp bốn phân số theo thứ tự từ bé đến lớn là: $\dfrac{7}{12}; \dfrac{3}{4}; \dfrac{19}{24}; \dfrac{20}{16}$.}
\loigiai{
Quy đồng mẫu số chung $48$ (hoặc $24$):\\
$\dfrac{7}{12} = \dfrac{14}{24}$; $\dfrac{3}{4} = \dfrac{18}{24}$; $\dfrac{19}{24}$; $\dfrac{20}{16} = \dfrac{5}{4} = \dfrac{30}{24}$.
\begin{itemize}
    \item Rút gọn: $\dfrac{20}{16} = \dfrac{20:4}{16:4} = \dfrac{5}{4}$.
    \item So sánh: $\dfrac{14}{24} < \dfrac{18}{24} \Rightarrow \dfrac{7}{12} < \dfrac{3}{4}$.
    \item Vì $\dfrac{20}{16} > 1$ còn ba phân số kia bé hơn $1$ nên $\dfrac{20}{16}$ là lớn nhất.
    \item Thứ tự từ bé đến lớn là: $\dfrac{14}{24} < \dfrac{18}{24} < \dfrac{19}{24} < \dfrac{30}{24}$, tức là $\dfrac{7}{12}; \dfrac{3}{4}; \dfrac{19}{24}; \dfrac{20}{16}$.
\end{itemize}
\begin{itemize}
    \item a) \textbf{Đúng}.
    \item b) \textbf{Sai}: $\dfrac{7}{12} < \dfrac{3}{4}$.
    \item c) \textbf{Đúng}.
    \item d) \textbf{Đúng}.
\end{itemize}
}
\end{ex}

\begin{ex} %[5D1V3]
Tại Câu lạc bộ thể thao trường Tiểu học Kim Đồng, số học sinh khối 5 tham gia: Cầu lông chiếm $\dfrac{1}{3}$, Bóng đá chiếm $\dfrac{2}{5}$, Võ chiếm $\dfrac{1}{5}$, Bóng rổ chiếm $\dfrac{4}{15}$ tổng số học sinh. Xét tính đúng sai của các khẳng định sau:
\choiceTF
{\True Số học sinh tham gia môn Võ ít nhất trong bốn môn.}
{Số học sinh tham gia môn Cầu lông nhiều hơn môn Bóng đá.}
{\True Tổng số phần học sinh tham gia môn Võ và Bóng đá là $\dfrac{3}{5}$ số học sinh khối 5.}
{Số học sinh tham gia môn Bóng rổ nhiều hơn môn Cầu lông.}
\loigiai{
Quy đồng về mẫu số $15$: Cầu lông là $\dfrac{5}{15}$; Bóng đá là $\dfrac{6}{15}$; Võ là $\dfrac{3}{15}$; Bóng rổ là $\dfrac{4}{15}$.
\begin{itemize}
    \item Môn Võ có tỉ lệ $\dfrac{3}{15}$ (nhỏ nhất) nên có ít học sinh tham gia nhất.
    \item Cầu lông ($\dfrac{5}{15}$) ít hơn Bóng đá ($\dfrac{6}{15}$).
    \item Tổng môn Võ và Bóng đá: $\dfrac{1}{5} + \dfrac{2}{5} = \dfrac{3}{5}$.
    \item Bóng rổ ($\dfrac{4}{15}$) ít hơn Cầu lông ($\dfrac{5}{15}$).
\end{itemize}
\begin{itemize}
    \item a) \textbf{Đúng}.
    \item b) \textbf{Sai}: Cầu lông ít hơn bóng đá.
    \item c) \textbf{Đúng}.
    \item d) \textbf{Sai}: Bóng rổ ít hơn cầu lông.
\end{itemize}
}
\end{ex}

\textbf{Phần 3. Trả lời ngắn}
\setcounterpageref{ex}{0}

\begin{ex} %[5D1H3]
Tìm số tự nhiên thích hợp điền vào dấu hỏi chấm để có đẳng thức:
\[
\dfrac{36}{54} = \dfrac{?}{9}.
\]
\shortans{6}
\loigiai{
\textbf{Đáp số:} $6$\\
Chia cả tử và mẫu của $\dfrac{36}{54}$ cho $6$: $\dfrac{36 : 6}{54 : 6} = \dfrac{6}{9} \Rightarrow ? = 6$.
}
\end{ex}

\begin{ex} %[5D1H3]
Rút gọn phân số $\dfrac{72}{90}$ thành phân số tối giản có dạng $\dfrac{a}{b}$. Tính giá trị của $a + b$.
\shortans{9}
\loigiai{
\textbf{Đáp số:} $9$\\
$\dfrac{72}{90} = \dfrac{72 : 18}{90 : 18} = \dfrac{4}{5}$. Do đó $a = 4, b = 5 \Rightarrow a + b = 9$.
}
\end{ex}

\begin{ex} %[5D1H3]
Tính giá trị của biểu thức:
\[
A = \dfrac{6 \times 12 \times 17}{17 \times 4 \times 20}.
\]
Viết kết quả dưới dạng phân số tối giản $\dfrac{m}{n}$. Tính tích $m \times n$.
\shortans{90}
\loigiai{
\textbf{Đáp số:} $90$\\
Rút gọn: $A = \dfrac{6 \times 3}{20} = \dfrac{18}{20} = \dfrac{9}{10}$. Khi đó $m = 9, n = 10 \Rightarrow m \times n = 90$.
}
\end{ex}

\begin{ex} %[5D1H3]
Tìm mẫu số chung nhỏ nhất của ba phân số $\dfrac{1}{3}$, $\dfrac{7}{12}$ và $\dfrac{35}{48}$.
\shortans{48}
\loigiai{
\textbf{Đáp số:} $48$\\
Vì $48$ chia hết cho cả $3$ và $12$ ($48 : 3 = 16$; $48 : 12 = 4$), nên mẫu số chung nhỏ nhất của ba phân số là $48$.
}
\end{ex}

\begin{ex} %[5D1V3]
Trong bốn phân số $\dfrac{3}{4}; \dfrac{7}{12}; \dfrac{20}{16}; \dfrac{19}{24}$, phân số nào có giá trị nhỏ nhất? (Điền tử số của phân số đó sau khi quy đồng mẫu số chung là $24$).
\shortans{14}
\loigiai{
\textbf{Đáp số:} $14$\\
Quy đồng mẫu số $24$: $\dfrac{7}{12} = \dfrac{14}{24}; \dfrac{3}{4} = \dfrac{18}{24}; \dfrac{19}{24}; \dfrac{20}{16} = \dfrac{30}{24}$. Phân số nhỏ nhất là $\dfrac{7}{12} = \dfrac{14}{24}$, tử số là $14$.
}
\end{ex}

\begin{ex} %[5D1V3]
Tìm chữ số tự nhiên $x$ thỏa mãn: $\dfrac{3}{8} < \dfrac{x}{8} < \dfrac{4}{7}$.
\shortans{4}
\loigiai{
\textbf{Đáp số:} $4$\\
Quy đồng mẫu số $56$: $\dfrac{21}{56} < \dfrac{7 \times x}{56} < \dfrac{32}{56}$. Suy ra $21 < 7 \times x < 32 \Rightarrow 7 \times x = 28 \Rightarrow x = 4$.
}
\end{ex}

\newpage

\TaoPhieuKiemTra{Chủ đề 1 -- Ôn tập và bổ sung}{04}{Phân số thập phân}{https://loptoanthaylinh.vn/thi-online/toan-5/bai-04}
\textbf{Phần 1. Trắc nghiệm lựa chọn}
\setcounterpageref{ex}{0}

\begin{ex} %[5D1N4]
Phân số nào dưới đây là phân số thập phân?
\choice
{$\dfrac{9}{20}$}
{\True $\dfrac{63}{100}$}
{$\dfrac{100}{59}$}
{$\dfrac{15}{30}$}
\loigiai{
Phân số thập phân là phân số có mẫu số là $10; 100; 1\,000;\dots$. Trong các phân số đã cho, $\dfrac{63}{100}$ có mẫu số là $100$ nên là phân số thập phân.\\
\textbf{Đáp án của câu hỏi là B}
}
\end{ex}

\begin{ex} %[5D1N4]
Trong các phân số $\dfrac{7}{10}; \dfrac{9}{20}; \dfrac{178}{100}; \dfrac{16}{1\,000}; \dfrac{100}{59}$, có bao nhiêu phân số thập phân?
\choice
{$2$}
{$4$}
{$5$}
{\True $3$}
\loigiai{
Các phân số thập phân là: $\dfrac{7}{10}$ (mẫu số là 10), $\dfrac{178}{100}$ (mẫu số là 100), $\dfrac{16}{1\,000}$ (mẫu số là 1 000). Có tất cả $3$ phân số thập phân.\\
\textbf{Đáp án của câu hỏi là D}
}
\end{ex}

\begin{ex} %[5D1H4]
Viết phân số $\dfrac{3}{5}$ dưới dạng phân số thập phân có mẫu số là $10$, ta được
\choice
{$\dfrac{3}{10}$}
{$\dfrac{9}{10}$}
{\True $\dfrac{6}{10}$}
{$\dfrac{15}{10}$}
\loigiai{
Nhân cả tử số và mẫu số với $2$: $\dfrac{3}{5} = \dfrac{3 \times 2}{5 \times 2} = \dfrac{6}{10}$.\\
\textbf{Đáp án của câu hỏi là C}
}
\end{ex}

\begin{ex} %[5D1H4]
Chuyển phân số $\dfrac{11}{25}$ thành phân số thập phân có mẫu số là $100$, ta được
\choice
{$\dfrac{22}{100}$}
{\True $\dfrac{44}{100}$}
{$\dfrac{33}{100}$}
{$\dfrac{55}{100}$}
\loigiai{
Nhân cả tử số và mẫu số với $4$: $\dfrac{11}{25} = \dfrac{11 \times 4}{25 \times 4} = \dfrac{44}{100}$.\\
\textbf{Đáp án của câu hỏi là B}
}
\end{ex}

\begin{ex} %[5D1H4]
Rút gọn phân số $\dfrac{25}{500}$ về phân số thập phân có mẫu số là $100$, ta được
\choice
{\True $\dfrac{5}{100}$}
{$\dfrac{25}{100}$}
{$\dfrac{1}{100}$}
{$\dfrac{15}{100}$}
\loigiai{
Chia cả tử số và mẫu số cho $5$: $\dfrac{25}{500} = \dfrac{25 : 5}{500 : 5} = \dfrac{5}{100}$.\\
\textbf{Đáp án của câu hỏi là A}
}
\end{ex}

\begin{ex} %[5D1H4]
Chuyển phân số $\dfrac{82}{2\,000}$ thành phân số thập phân có mẫu số là $1\,000$, ta được
\choice
{$\dfrac{82}{1\,000}$}
{$\dfrac{20}{1\,000}$}
{\True $\dfrac{41}{1\,000}$}
{$\dfrac{164}{1\,000}$}
\loigiai{
Chia cả tử số và mẫu số cho $2$: $\dfrac{82}{2\,000} = \dfrac{82 : 2}{2\,000 : 2} = \dfrac{41}{1\,000}$.\\
\textbf{Đáp án của câu hỏi là C}
}
\end{ex}

\begin{ex} %[5D1H4]
Phân số $\dfrac{66}{60}$ viết dưới dạng phân số thập phân là
\choice
{$\dfrac{6}{10}$}
{$\dfrac{66}{10}$}
{$\dfrac{33}{10}$}
{\True $\dfrac{11}{10}$}
\loigiai{
Chia cả tử số và mẫu số cho $6$: $\dfrac{66}{60} = \dfrac{66 : 6}{60 : 6} = \dfrac{11}{10}$.\\
\textbf{Đáp án của câu hỏi là D}
}
\end{ex}

\begin{ex} %[5D1H4]
Phân số $\dfrac{27}{300}$ viết dưới dạng phân số thập phân là
\choice
{\True $\dfrac{9}{100}$}
{$\dfrac{3}{100}$}
{$\dfrac{27}{100}$}
{$\dfrac{18}{100}$}
\loigiai{
Chia cả tử số và mẫu số cho $3$: $\dfrac{27}{300} = \dfrac{27 : 3}{300 : 3} = \dfrac{9}{100}$.\\
\textbf{Đáp án của câu hỏi là A}
}
\end{ex}

\begin{ex} %[5D1H4]
Phân số $\dfrac{3}{4}$ viết dưới dạng phân số thập phân có mẫu số là $100$ là
\choice
{$\dfrac{25}{100}$}
{\True $\dfrac{75}{100}$}
{$\dfrac{50}{100}$}
{$\dfrac{30}{100}$}
\loigiai{
Nhân cả tử số và mẫu số với $25$: $\dfrac{3}{4} = \dfrac{3 \times 25}{4 \times 25} = \dfrac{75}{100}$.\\
\textbf{Đáp án của câu hỏi là B}
}
\end{ex}

\begin{ex} %[5D1H4]
Chuyển phân số $\dfrac{31}{125}$ thành phân số thập phân có mẫu số là $1\,000$, ta được
\choice
{$\dfrac{124}{1\,000}$}
{$\dfrac{62}{1\,000}$}
{$\dfrac{310}{1\,000}$}
{\True $\dfrac{248}{1\,000}$}
\loigiai{
Nhân cả tử số và mẫu số với $8$: $\dfrac{31}{125} = \dfrac{31 \times 8}{125 \times 8} = \dfrac{248}{1\,000}$.\\
\textbf{Đáp án của câu hỏi là D}
}
\end{ex}

\begin{ex} %[5D1H4]
Điểm chính giữa giữa hai phân số $\dfrac{3}{10}$ và $\dfrac{5}{10}$ trên tia số biểu diễn phân số thập phân nào?
\choice
{\True $\dfrac{4}{10}$}
{$\dfrac{2}{10}$}
{$\dfrac{6}{10}$}
{$\dfrac{8}{10}$}
\loigiai{
Trên tia số cách đều $\dfrac{1}{10}$, điểm nằm giữa $\dfrac{3}{10}$ và $\dfrac{5}{10}$ là $\dfrac{4}{10}$.\\
\textbf{Đáp án của câu hỏi là A}
}
\end{ex}

\begin{ex} %[5D1V4]
Phân số nào sau đây \textbf{không thể} chuyển thành phân số thập phân (có mẫu số là lũy thừa của $10$)?
\choice
{$\dfrac{7}{20}$}
{$\dfrac{9}{25}$}
{\True $\dfrac{5}{7}$}
{$\dfrac{13}{50}$}
\loigiai{
Mẫu số của các phân số thập phân chỉ có ước nguyên tố là $2$ và $5$. Phân số $\dfrac{5}{7}$ có mẫu số là $7$ nên không thể quy đồng về mẫu $10, 100, 1\,000,\dots$\\
\textbf{Đáp án của câu hỏi là C}
}
\end{ex}

\textbf{Phần 2. Trắc nghiệm đúng sai}
\setcounterpageref{ex}{0}

\begin{ex} %[5D1H4]
Cho các phân số: $A = \dfrac{11}{25}$, $B = \dfrac{25}{500}$, $C = \dfrac{66}{60}$, $D = \dfrac{31}{125}$. Xét tính đúng sai của các khẳng định sau:
\choiceTF
{\True Phân số $A$ được viết thành phân số thập phân là $\dfrac{44}{100}$.}
{Phân số $B$ viết thành phân số thập phân có mẫu số là $100$ là $\dfrac{25}{100}$.}
{\True Phân số $C$ viết thành phân số thập phân là $\dfrac{11}{10}$.}
{\True Phân số $D$ viết thành phân số thập phân có mẫu số $1\,000$ là $\dfrac{248}{1\,000}$.}
\loigiai{
\begin{itemize}
    \item $A = \dfrac{11 \times 4}{25 \times 4} = \dfrac{44}{100}$.
    \item $B = \dfrac{25 : 5}{500 : 5} = \dfrac{5}{100} \neq \dfrac{25}{100}$.
    \item $C = \dfrac{66 : 6}{60 : 6} = \dfrac{11}{10}$.
    \item $D = \dfrac{31 \times 8}{125 \times 8} = \dfrac{248}{1\,000}$.
\end{itemize}
\begin{itemize}
    \item a) \textbf{Đúng}.
    \item b) \textbf{Sai}: $B = \dfrac{5}{100}$.
    \item c) \textbf{Đúng}.
    \item d) \textbf{Đúng}.
\end{itemize}
}
\end{ex}

\begin{ex} %[5D1H4]
Xét tính đúng sai của các khẳng định sau về phân số thập phân:
\choiceTF
{\True Phân số $\dfrac{178}{100}$ là một phân số thập phân lớn hơn $1$.}
{Mọi phân số có tử số là $10, 100, 1\,000$ đều là phân số thập phân.}
{\True Phân số $\dfrac{3}{4}$ có thể chuyển thành phân số thập phân có mẫu số là $100$ với tử số bằng $75$.}
{Tổng của hai phân số thập phân $\dfrac{3}{10}$ và $\dfrac{57}{100}$ bằng $\dfrac{60}{100}$.}
\loigiai{
\begin{itemize}
    \item Phân số $\dfrac{178}{100}$ có mẫu số là $100$ và tử số $178 > 100$ nên là phân số thập phân lớn hơn $1$.
    \item Phân số thập phân yêu cầu \textbf{mẫu số} phải là $10, 100, 1\,000,\dots$, tử số là số tự nhiên bất kì. (Ví dụ $\dfrac{100}{59}$ không phải phân số thập phân).
    \item $\dfrac{3}{4} = \dfrac{3 \times 25}{4 \times 25} = \dfrac{75}{100}$.
    \item $\dfrac{3}{10} + \dfrac{57}{100} = \dfrac{30}{100} + \dfrac{57}{100} = \dfrac{87}{100} \neq \dfrac{60}{100}$.
\end{itemize}
\begin{itemize}
    \item a) \textbf{Đúng}.
    \item b) \textbf{Sai}: Phân số thập phân có mẫu số là $10, 100, 1\,000,\dots$
    \item c) \textbf{Đúng}.
    \item d) \textbf{Sai}: Tổng bằng $\dfrac{87}{100}$.
\end{itemize}
}
\end{ex}

\textbf{Phần 3. Trả lời ngắn}
\setcounterpageref{ex}{0}

\begin{ex} %[5D1H4]
Viết phân số $\dfrac{3}{5}$ dưới dạng phân số thập phân có mẫu số là $10$. Hỏi tử số của phân số thập phân đó là bao nhiêu?
\shortans{6}
\loigiai{
\textbf{Đáp số:} $6$\\
Ta có: $\dfrac{3}{5} = \dfrac{3 \times 2}{5 \times 2} = \dfrac{6}{10}$. Tử số là $6$.
}
\end{ex}

\begin{ex} %[5D1H4]
Chuyển phân số $\dfrac{11}{25}$ thành phân số thập phân có mẫu số là $100$. Tử số của phân số thập phân đó bằng bao nhiêu?
\shortans{44}
\loigiai{
\textbf{Đáp số:} $44$\\
Ta có: $\dfrac{11}{25} = \dfrac{11 \times 4}{25 \times 4} = \dfrac{44}{100}$. Tử số là $44$.
}
\end{ex}

\begin{ex} %[5D1H4]
Viết phân số $\dfrac{27}{300}$ dưới dạng phân số thập phân có mẫu số là $100$. Tìm tử số của phân số đó.
\shortans{9}
\loigiai{
\textbf{Đáp số:} $9$\\
Ta có: $\dfrac{27}{300} = \dfrac{27 : 3}{300 : 3} = \dfrac{9}{100}$. Tử số là $9$.
}
\end{ex}

\begin{ex} %[5D1H4]
Viết phân số $\dfrac{31}{125}$ dưới dạng phân số thập phân có mẫu số là $1\,000$. Tìm tử số của phân số nhận được.
\shortans{248}
\loigiai{
\textbf{Đáp số:} $248$\\
Ta có: $\dfrac{31}{125} = \dfrac{31 \times 8}{125 \times 8} = \dfrac{248}{1\,000}$. Tử số là $248$.
}
\end{ex}

\begin{ex} %[5D1H4]
Rút gọn phân số $\dfrac{82}{2\,000}$ thành phân số thập phân có mẫu số là $1\,000$. Hỏi tử số của phân số thu được là bao nhiêu?
\shortans{41}
\loigiai{
\textbf{Đáp số:} $41$\\
$\dfrac{82}{2\,000} = \dfrac{82 : 2}{2\,000 : 2} = \dfrac{41}{1\,000}$. Tử số là $41$.
}
\end{ex}

\begin{ex} %[5D1V4]
Tìm số tự nhiên $x$ sao cho $\dfrac{x}{50}$ được viết dưới dạng phân số thập phân có mẫu số $100$ là $\dfrac{76}{100}$.
\shortans{38}
\loigiai{
\textbf{Đáp số:} $38$\\
Ta có: $\dfrac{x}{50} = \dfrac{2 \times x}{100} = \dfrac{76}{100} \Rightarrow 2 \times x = 76 \Rightarrow x = 38$.
}
\end{ex}

\newpage

\TaoPhieuKiemTra{Chủ đề 1 -- Ôn tập và bổ sung}{05}{Ôn tập các phép tính với phân số}{https://loptoanthaylinh.vn/thi-online/toan-5/bai-05}
\textbf{Phần 1. Trắc nghiệm lựa chọn}
\setcounterpageref{ex}{0}

\begin{ex} %[5D1N5]
Kết quả của phép tính $\dfrac{6}{5} + \dfrac{7}{15}$ là
\choice
{$\dfrac{13}{20}$}
{$\dfrac{13}{15}$}
{$\dfrac{25}{30}$}
{\True $\dfrac{5}{3}$}
\loigiai{
Quy đồng mẫu số: $\dfrac{6}{5} + \dfrac{7}{15} = \dfrac{18}{15} + \dfrac{7}{15} = \dfrac{25}{15} = \dfrac{5}{3}$.\\
\textbf{Đáp án của câu hỏi là D}
}
\end{ex}

\begin{ex} %[5D1N5]
Kết quả của phép trừ $\dfrac{15}{16} - \dfrac{3}{4}$ là
\choice
{\True $\dfrac{3}{16}$}
{$\dfrac{12}{12}$}
{$\dfrac{12}{16}$}
{$\dfrac{9}{16}$}
\loigiai{
Quy đồng mẫu số: $\dfrac{15}{16} - \dfrac{3}{4} = \dfrac{15}{16} - \dfrac{12}{16} = \dfrac{3}{16}$.\\
\textbf{Đáp án của câu hỏi là A}
}
\end{ex}

\begin{ex} %[5D1H5]
Kết quả của phép tính $3 - \dfrac{16}{7}$ là
\choice
{$\dfrac{13}{7}$}
{$\dfrac{21}{7}$}
{\True $\dfrac{5}{7}$}
{$\dfrac{1}{7}$}
\loigiai{
Ta có: $3 - \dfrac{16}{7} = \dfrac{21}{7} - \dfrac{16}{7} = \dfrac{5}{7}$.\\
\textbf{Đáp án của câu hỏi là C}
}
\end{ex}

\begin{ex} %[5D1H5]
Kết quả của phép tính $\dfrac{5}{8} \times \dfrac{8}{15}$ là
\choice
{$\dfrac{1}{5}$}
{\True $\dfrac{1}{3}$}
{$\dfrac{5}{15}$}
{$\dfrac{3}{8}$}
\loigiai{
Rút gọn tích: $\dfrac{5 \times 8}{8 \times 15} = \dfrac{5}{15} = \dfrac{1}{3}$.\\
\textbf{Đáp án của câu hỏi là B}
}
\end{ex}

\begin{ex} %[5D1H5]
Kết quả của phép chia $\dfrac{12}{7} : \dfrac{8}{7}$ là
\choice
{$\dfrac{4}{7}$}
{$\dfrac{2}{3}$}
{\True $\dfrac{3}{2}$}
{$\dfrac{96}{49}$}
\loigiai{
Thực hiện phép chia: $\dfrac{12}{7} : \dfrac{8}{7} = \dfrac{12}{7} \times \dfrac{7}{8} = \dfrac{12}{8} = \dfrac{3}{2}$.\\
\textbf{Đáp án của câu hỏi là C}
}
\end{ex}

\begin{ex} %[5D1H5]
Kết quả của phép tính $\dfrac{24}{5} : 4$ là
\choice
{$\dfrac{96}{5}$}
{$\dfrac{24}{20}$}
{$\dfrac{1}{5}$}
{\True $\dfrac{6}{5}$}
\loigiai{
Ta có: $\dfrac{24}{5} : 4 = \dfrac{24 : 4}{5} = \dfrac{6}{5}$.\\
\textbf{Đáp án của câu hỏi là D}
}
\end{ex}

\begin{ex} %[5D1H5]
Giá trị của biểu thức $\dfrac{4}{7} + \dfrac{9}{14} - \dfrac{5}{7}$ là
\choice
{\True $\dfrac{1}{2}$}
{$\dfrac{1}{14}$}
{$\dfrac{8}{14}$}
{$\dfrac{3}{7}$}
\loigiai{
Giao hoán và kết hợp: $\left(\dfrac{4}{7} - \dfrac{5}{7}\right) + \dfrac{9}{14} = \dfrac{9}{14} - \dfrac{1}{7} = \dfrac{9}{14} - \dfrac{2}{14} = \dfrac{7}{14} = \dfrac{1}{2}$.\\
\textbf{Đáp án của câu hỏi là A}
}
\end{ex}

\begin{ex} %[5D1H5]
Giá trị của biểu thức $\dfrac{9}{10} \times \dfrac{5}{6} : 3$ là
\choice
{$\dfrac{3}{4}$}
{$\dfrac{9}{4}$}
{\True $\dfrac{1}{4}$}
{$\dfrac{1}{2}$}
\loigiai{
Thực hiện lần lượt từ trái sang phải: $\dfrac{9 \times 5}{10 \times 6} = \dfrac{45}{60} = \dfrac{3}{4}$.\\
Sau đó: $\dfrac{3}{4} : 3 = \dfrac{1}{4}$.\\
\textbf{Đáp án của câu hỏi là C}
}
\end{ex}

\begin{ex} %[5D1H5]
Tính nhanh biểu thức $\dfrac{10}{11} \times \dfrac{9}{16} \times \dfrac{11}{10} \times \dfrac{8}{9}$ ta được kết quả là
\choice
{$1$}
{\True $\dfrac{1}{2}$}
{$2$}
{$\dfrac{1}{4}$}
\loigiai{
Nhóm các phân số nghịch đảo và rút gọn:
\[
\left(\dfrac{10}{11} \times \dfrac{11}{10}\right) \times \left(\dfrac{9}{16} \times \dfrac{8}{9}\right) = 1 \times \dfrac{8}{16} = \dfrac{1}{2}.
\]
\textbf{Đáp án của câu hỏi là B}
}
\end{ex}

\begin{ex} %[5D1V5]
Hai bạn Nam và Việt cùng chạy thi trên một đoạn đường. Sau một phút, Nam chạy được $\dfrac{3}{5}$ đoạn đường, Việt chạy được $\dfrac{7}{10}$ đoạn đường. Khẳng định nào sau đây là đúng?
\choice
{Nam chạy nhiều hơn Việt $\dfrac{1}{10}$ đoạn đường}
{Việt chạy nhiều hơn Nam $\dfrac{4}{10}$ đoạn đường}
{Hai bạn chạy được quãng đường bằng nhau}
{\True Việt chạy nhiều hơn Nam $\dfrac{1}{10}$ đoạn đường}
\loigiai{
Quy đồng: $\dfrac{3}{5} = \dfrac{6}{10}$.\\
Vì $\dfrac{7}{10} > \dfrac{6}{10}$ nên Việt chạy được nhiều hơn Nam là: $\dfrac{7}{10} - \dfrac{6}{10} = \dfrac{1}{10}$ đoạn đường.\\
\textbf{Đáp án của câu hỏi là D}
}
\end{ex}

\begin{ex} %[5D1V5]
Thư viện trường học có $\dfrac{2}{3}$ số sách là sách giáo khoa, $\dfrac{2}{9}$ số sách là sách tham khảo, còn lại là truyện thiếu nhi và tạp chí. Phân số chỉ số truyện thiếu nhi và tạp chí là
\choice
{\True $\dfrac{1}{9}$}
{$\dfrac{4}{9}$}
{$\dfrac{8}{9}$}
{$\dfrac{2}{9}$}
\loigiai{
Tổng số phần sách giáo khoa và sách tham khảo là: $\dfrac{2}{3} + \dfrac{2}{9} = \dfrac{6}{9} + \dfrac{2}{9} = \dfrac{8}{9}$.\\
Phân số chỉ số truyện thiếu nhi và tạp chí là: $1 - \dfrac{8}{9} = \dfrac{1}{9}$.\\
\textbf{Đáp án của câu hỏi là A}
}
\end{ex}

\begin{ex} %[5D1V5]
Một tấm kính hình chữ nhật có chiều dài $\dfrac{5}{2}\text{ m}$, chiều rộng $\dfrac{4}{3}\text{ m}$. Chú Hoà chia tấm kính đó thành $3$ phần bằng nhau để làm mặt bàn. Diện tích mỗi phần tấm kính làm mặt bàn là
\choice
{$\dfrac{10}{3}\text{ m}^2$}
{\True $\dfrac{10}{9}\text{ m}^2$}
{$\dfrac{5}{9}\text{ m}^2$}
{$\dfrac{20}{9}\text{ m}^2$}
\loigiai{
Diện tích cả tấm kính là: $\dfrac{5}{2} \times \dfrac{4}{3} = \dfrac{20}{6} = \dfrac{10}{3}\text{ (m}^2\text{)}$.\\
Diện tích mỗi phần làm mặt bàn là: $\dfrac{10}{3} : 3 = \dfrac{10}{9}\text{ (m}^2\text{)}$.\\
\textbf{Đáp án của câu hỏi là B}
}
\end{ex}

\textbf{Phần 2. Trắc nghiệm đúng sai}
\setcounterpageref{ex}{0}

\begin{ex} %[5D1H5]
Cho các biểu thức: $A = \dfrac{29}{12} - \left(2 + \dfrac{1}{4}\right)$, $B = \dfrac{5}{6} \times \dfrac{3}{7} + \dfrac{3}{7} \times \dfrac{7}{6}$ và $C = \dfrac{11}{4} : \left(\dfrac{11}{14} \times 7\right)$. Xét tính đúng sai của các khẳng định sau:
\choiceTF
{\True Biểu thức $A$ có giá trị bằng $\dfrac{1}{6}$.}
{Biểu thức $B$ có giá trị bằng $1$.}
{\True Biểu thức $C$ có giá trị bằng $\dfrac{1}{2}$.}
{\True Giá trị của biểu thức $B$ lớn hơn giá trị của biểu thức $C$.}
\loigiai{
\begin{itemize}
    \item $A = \dfrac{29}{12} - \dfrac{9}{4} = \dfrac{29}{12} - \dfrac{27}{12} = \dfrac{2}{12} = \dfrac{1}{6}$.
    \item $B = \dfrac{3}{7} \times \left(\dfrac{5}{6} + \dfrac{7}{6}\right) = \dfrac{3}{7} \times 2 = \dfrac{6}{7} \neq 1$.
    \item $C = \dfrac{11}{4} : \left(\dfrac{11}{2}\right) = \dfrac{11}{4} \times \dfrac{2}{11} = \dfrac{2}{4} = \dfrac{1}{2}$.
    \item So sánh $B = \dfrac{6}{7} \approx 0{,}86$ và $C = \dfrac{1}{2} = 0{,}5 \Rightarrow B > C$.
\end{itemize}
\begin{itemize}
    \item a) \textbf{Đúng}.
    \item b) \textbf{Sai}: $B = \dfrac{6}{7}$.
    \item c) \textbf{Đúng}.
    \item d) \textbf{Đúng}.
\end{itemize}
}
\end{ex}

\begin{ex} %[5D1V5]
Hưởng ứng phong trào quyên góp sách vào thư viện trường học, lớp 5A quyên góp được $96$ quyển sách, lớp 5B quyên góp được số quyển sách bằng $\dfrac{7}{8}$ số quyển sách của lớp 5A. Xét tính đúng sai của các khẳng định sau:
\choiceTF
{\True Lớp 5B quyên góp được $84$ quyển sách.}
{Lớp 5B quyên góp được nhiều sách hơn lớp 5A.}
{\True Cả hai lớp 5A và 5B quyên góp được tất cả $180$ quyển sách.}
{Lớp 5A quyên góp nhiều hơn lớp 5B là $14$ quyển sách.}
\loigiai{
\begin{itemize}
    \item Số sách lớp 5B quyên góp là: $96 \times \dfrac{7}{8} = 84$ (quyển).
    \item Vì $84 < 96$ nên lớp 5B quyên góp ít hơn lớp 5A.
    \item Cả hai lớp quyên góp: $96 + 84 = 180$ (quyển).
    \item Lớp 5A quyên góp nhiều hơn lớp 5B là: $96 - 84 = 12$ quyển ($\neq 14$ quyển).
\end{itemize}
\begin{itemize}
    \item a) \textbf{Đúng}.
    \item b) \textbf{Sai}: Lớp 5B quyên góp ít hơn lớp 5A.
    \item c) \textbf{Đúng}.
    \item d) \textbf{Sai}: Nhiều hơn $12$ quyển sách.
\end{itemize}
}
\end{ex}

\textbf{Phần 3. Trả lời ngắn}
\setcounterpageref{ex}{0}

\begin{ex} %[5D1H5]
Tính giá trị của biểu thức sau dưới dạng phân số tối giản $\dfrac{a}{b}$:
\[
A = \dfrac{5}{6} \times \dfrac{3}{7} + \dfrac{3}{7} \times \dfrac{7}{6}.
\]
Tính giá trị của $a + b$.
\shortans{13}
\loigiai{
\textbf{Đáp số:} $13$\\
$A = \dfrac{3}{7} \times \left(\dfrac{5}{6} + \dfrac{7}{6}\right) = \dfrac{3}{7} \times 2 = \dfrac{6}{7}$. Do đó $a = 6, b = 7 \Rightarrow a + b = 13$.
}
\end{ex}

\begin{ex} %[5D1H5]
Người ta gắn dây đèn xung quanh một tấm biển quảng cáo hình vuông có chu vi là $18\text{ m}$. Tính diện tích của tấm biển quảng cáo đó dưới dạng phân số tối giản $\dfrac{m}{n}\text{ (m}^2\text{)}$. Giá trị của tử số $m$ là bao nhiêu?
\shortans{81}
\loigiai{
\textbf{Đáp số:} $81$\\
Cạnh của hình vuông là: $18 : 4 = \dfrac{9}{2}\text{ (m)}$.\\
Diện tích tấm biển là: $\dfrac{9}{2} \times \dfrac{9}{2} = \dfrac{81}{4}\text{ (m}^2\text{)}$. Vậy $m = 81$.
}
\end{ex}

\begin{ex} %[5D1V5]
Lớp 5A quyên góp được $96$ quyển sách, lớp 5B quyên góp được số quyển sách bằng $\dfrac{7}{8}$ số sách của lớp 5A. Hỏi cả hai lớp đã quyên góp được tất cả bao nhiêu quyển sách?
\shortans{180}
\loigiai{
\textbf{Đáp số:} $180$\\
Số sách lớp 5B quyên góp là: $96 \times \dfrac{7}{8} = 84$ quyển.\\
Cả hai lớp quyên góp: $96 + 84 = 180$ quyển.
}
\end{ex}

\begin{ex} %[5D1V5]
Tấm kính hình chữ nhật có chiều dài $\dfrac{5}{2}\text{ m}$, chiều rộng $\dfrac{4}{3}\text{ m}$ được chia đều làm $3$ mặt bàn bằng nhau. Diện tích của mỗi mặt bàn là $\dfrac{p}{q}\text{ m}^2$ (với $\dfrac{p}{q}$ là phân số tối giản). Tính tổng $p + q$.
\shortans{19}
\loigiai{
\textbf{Đáp số:} $19$\\
Diện tích cả tấm kính: $\dfrac{5}{2} \times \dfrac{4}{3} = \dfrac{10}{3}\text{ m}^2$.\\
Diện tích mỗi mặt bàn: $\dfrac{10}{3} : 3 = \dfrac{10}{9}\text{ m}^2 \Rightarrow p = 10, q = 9 \Rightarrow p + q = 19$.
}
\end{ex}

\begin{ex} %[5D1H5]
Tính giá trị của biểu thức:
\[
M = \dfrac{9}{7} \times \left(\dfrac{7}{9} - \dfrac{2}{3}\right).
\]
Viết kết quả dưới dạng phân số tối giản $\dfrac{1}{k}$. Tìm giá trị của $k$.
\shortans{7}
\loigiai{
\textbf{Đáp số:} $7$\\
$M = \dfrac{9}{7} \times \left(\dfrac{7}{9} - \dfrac{6}{9}\right) = \dfrac{9}{7} \times \dfrac{1}{9} = \dfrac{1}{7} \Rightarrow k = 7$.
}
\end{ex}

\begin{ex} %[5D1V5]
Để cắt lấy đoạn băng giấy dài $\dfrac{1}{2}\text{ m}$ từ một băng giấy ban đầu dài $\dfrac{2}{3}\text{ m}$ mà không dùng thước đo, ta cần cắt bỏ đi một đoạn băng giấy có chiều dài bằng bao nhiêu mét? (Viết kết quả dưới dạng phân số tối giản $\dfrac{1}{x}$). Tìm giá trị của $x$.
\shortans{6}
\loigiai{
\textbf{Đáp số:} $6$\\
Đoạn băng giấy cần cắt bỏ có chiều dài là: $\dfrac{2}{3} - \dfrac{1}{2} = \dfrac{4}{6} - \dfrac{3}{6} = \dfrac{1}{6}\text{ (m)}$. Do đó $x = 6$.
}
\end{ex}

\newpage

\TaoPhieuKiemTra{Chủ đề 1 -- Ôn tập và bổ sung}{06}{Cộng, trừ hai phân số}{https://loptoanthaylinh.vn/thi-online/toan-5/bai-06}
\textbf{Phần 1. Trắc nghiệm lựa chọn}
\setcounterpageref{ex}{0}

\begin{ex} %[5D1N6]
Kết quả của phép cộng $\dfrac{1}{7} + \dfrac{1}{9}$ là
\choice
{$\dfrac{2}{16}$}
{$\dfrac{2}{63}$}
{\True $\dfrac{16}{63}$}
{$\dfrac{14}{63}$}
\loigiai{
Quy đồng mẫu số chung là $63$: $\dfrac{1}{7} + \dfrac{1}{9} = \dfrac{9}{63} + \dfrac{7}{63} = \dfrac{16}{63}$.\\
\textbf{Đáp án của câu hỏi là C}
}
\end{ex}

\begin{ex} %[5D1N6]
Kết quả của phép trừ $\dfrac{1}{4} - \dfrac{1}{9}$ là
\choice
{\True $\dfrac{5}{36}$}
{$\dfrac{1}{5}$}
{$\dfrac{13}{36}$}
{$\dfrac{5}{13}$}
\loigiai{
Quy đồng mẫu số chung $36$: $\dfrac{1}{4} - \dfrac{1}{9} = \dfrac{9}{36} - \dfrac{4}{36} = \dfrac{5}{36}$.\\
\textbf{Đáp án của câu hỏi là A}
}
\end{ex}

\begin{ex} %[5D1H6]
Kết quả của phép cộng $\dfrac{3}{11} + \dfrac{7}{12}$ là
\choice
{$\dfrac{10}{23}$}
{\True $\dfrac{113}{132}$}
{$\dfrac{10}{132}$}
{$\dfrac{111}{132}$}
\loigiai{
Quy đồng mẫu số chung là $132$:
\[
\dfrac{3}{11} + \dfrac{7}{12} = \dfrac{3 \times 12}{132} + \dfrac{7 \times 11}{132} = \dfrac{36 + 77}{132} = \dfrac{113}{132}.
\]
\textbf{Đáp án của câu hỏi là B}
}
\end{ex}

\begin{ex} %[5D1H6]
Kết quả của phép trừ $\dfrac{2}{5} - \dfrac{1}{4}$ là
\choice
{$\dfrac{1}{1}$}
{$\dfrac{1}{20}$}
{$\dfrac{7}{20}$}
{\True $\dfrac{3}{20}$}
\loigiai{
Quy đồng mẫu số chung là $20$: $\dfrac{2}{5} - \dfrac{1}{4} = \dfrac{8}{20} - \dfrac{5}{20} = \dfrac{3}{20}$.\\
\textbf{Đáp án của câu hỏi là D}
}
\end{ex}

\begin{ex} %[5D1H6]
Kết quả của phép cộng $\dfrac{6}{11} + \dfrac{1}{2}$ là
\choice
{\True $\dfrac{23}{22}$}
{$\dfrac{7}{13}$}
{$\dfrac{13}{22}$}
{$\dfrac{17}{22}$}
\loigiai{
Quy đồng mẫu số: $\dfrac{6}{11} + \dfrac{1}{2} = \dfrac{12}{22} + \dfrac{11}{22} = \dfrac{23}{22}$.\\
\textbf{Đáp án của câu hỏi là A}
}
\end{ex}

\begin{ex} %[5D1H6]
Rút gọn rồi tính hiệu $\dfrac{20}{18} - \dfrac{2}{5}$ ta được kết quả là
\choice
{$\dfrac{18}{13}$}
{\True $\dfrac{32}{45}$}
{$\dfrac{14}{45}$}
{$\dfrac{16}{45}$}
\loigiai{
Rút gọn $\dfrac{20}{18} = \dfrac{10}{9}$.\\
Thực hiện phép trừ: $\dfrac{10}{9} - \dfrac{2}{5} = \dfrac{50}{45} - \dfrac{18}{45} = \dfrac{32}{45}$.\\
\textbf{Đáp án của câu hỏi là B}
}
\end{ex}

\begin{ex} %[5D1H6]
Dấu thích hợp điền vào ô trống để $\dfrac{1}{20} \dots \dfrac{2}{3} = \dfrac{43}{60}$ là
\choice
{$-$}
{$\times$}
{$:$}
{\True $+$}
\loigiai{
Ta có: $\dfrac{1}{20} + \dfrac{2}{3} = \dfrac{3}{60} + \dfrac{40}{60} = \dfrac{43}{60}$. Vậy dấu cần điền là dấu $+$.\\
\textbf{Đáp án của câu hỏi là D}
}
\end{ex}

\begin{ex} %[5D1H6]
Dấu thích hợp điền vào ô trống để $\dfrac{7}{11} \dots \dfrac{5}{8} = \dfrac{1}{88}$ là
\choice
{$+$}
{\True $-$}
{$\times$}
{$:$}
\loigiai{
Ta có: $\dfrac{7}{11} - \dfrac{5}{8} = \dfrac{56}{88} - \dfrac{55}{88} = \dfrac{1}{88}$. Vậy dấu cần điền là dấu $-$.\\
\textbf{Đáp án của câu hỏi là B}
}
\end{ex}

\begin{ex} %[5D1H6]
Kết quả của phép tính $7 + \dfrac{1}{5}$ dưới dạng phân số là
\choice
{$\dfrac{8}{5}$}
{$\dfrac{7}{5}$}
{\True $\dfrac{36}{5}$}
{$\dfrac{35}{5}$}
\loigiai{
Viết $7 = \dfrac{35}{5}$. Khi đó $7 + \dfrac{1}{5} = \dfrac{35}{5} + \dfrac{1}{5} = \dfrac{36}{5}$.\\
\textbf{Đáp án của câu hỏi là C}
}
\end{ex}

\begin{ex} %[5D1H6]
Kết quả của phép tính $3 - \dfrac{11}{8}$ dưới dạng phân số là
\choice
{$\dfrac{8}{8}$}
{$\dfrac{14}{8}$}
{$\dfrac{21}{8}$}
{\True $\dfrac{13}{8}$}
\loigiai{
Viết $3 = \dfrac{24}{8}$. Khi đó $3 - \dfrac{11}{8} = \dfrac{24}{8} - \dfrac{11}{8} = \dfrac{13}{8}$.\\
\textbf{Đáp án của câu hỏi là D}
}
\end{ex}

\begin{ex} %[5D1V6]
Có hai vòi nước cùng chảy vào một bể cạn. Vòi thứ nhất chảy được $\dfrac{1}{2}$ bể, vòi thứ hai chảy được $\dfrac{2}{5}$ bể. Cả hai vòi đã chảy được bao nhiêu phần bể nước?
\choice
{$\dfrac{3}{7}\text{ bể}$}
{$\dfrac{3}{10}\text{ bể}$}
{\True $\dfrac{9}{10}\text{ bể}$}
{$1\text{ bể}$}
\loigiai{
Cả hai vòi chảy được số phần bể nước là: $\dfrac{1}{2} + \dfrac{2}{5} = \dfrac{5}{10} + \dfrac{4}{10} = \dfrac{9}{10}$ (bể).\\
\textbf{Đáp án của câu hỏi là C}
}
\end{ex}

\begin{ex} %[5D1V6]
Để chia đều $5$ chiếc bánh cho $6$ người, mỗi người sẽ nhận được bao nhiêu phần chiếc bánh?
\choice
{\True $\dfrac{5}{6}\text{ chiếc bánh}$}
{$\dfrac{6}{5}\text{ chiếc bánh}$}
{$\dfrac{1}{6}\text{ chiếc bánh}$}
{$\dfrac{1}{5}\text{ chiếc bánh}$}
\loigiai{
Chia đều $5$ chiếc bánh cho $6$ người: mỗi người được $5 : 6 = \dfrac{5}{6}$ chiếc bánh.\\
\textbf{Đáp án của câu hỏi là A}
}
\end{ex}

\textbf{Phần 2. Trắc nghiệm đúng sai}
\setcounterpageref{ex}{0}

\begin{ex} %[5D1H6]
Cho các phép tính: $P = \dfrac{13}{5} + \dfrac{7}{6}$, $Q = \dfrac{11}{7} - \dfrac{5}{8}$ và $R = \dfrac{11}{4} - \dfrac{7}{5}$. Xét tính đúng sai của các khẳng định sau:
\choiceTF
{\True Phép tính $P$ có kết quả là $\dfrac{113}{30}$.}
{Phép tính $Q$ có kết quả là $\dfrac{6}{56}$.}
{\True Phép tính $R$ có kết quả là $\dfrac{27}{20}$.}
{\True Kết quả phép tính $P$ lớn hơn $3$.}
\loigiai{
\begin{itemize}
    \item $P = \dfrac{13 \times 6 + 7 \times 5}{30} = \dfrac{78 + 35}{30} = \dfrac{113}{30}$.
    \item $Q = \dfrac{11 \times 8 - 5 \times 7}{56} = \dfrac{88 - 35}{56} = \dfrac{53}{56} \neq \dfrac{6}{56}$.
    \item $R = \dfrac{11 \times 5 - 7 \times 4}{20} = \dfrac{55 - 28}{20} = \dfrac{27}{20}$.
    \item Vì $\dfrac{113}{30} = 3\dfrac{23}{30} > 3$.
\end{itemize}
\begin{itemize}
    \item a) \textbf{Đúng}.
    \item b) \textbf{Sai}: $Q = \dfrac{53}{56}$.
    \item c) \textbf{Đúng}.
    \item d) \textbf{Đúng}.
\end{itemize}
}
\end{ex}

\begin{ex} %[5D1V6]
Việt đổ $\dfrac{1}{5}\text{ l}$ nước vào bình, Mai đổ $\dfrac{1}{2}\text{ l}$ nước vào cùng chiếc bình đó. Xét tính đúng sai của các khẳng định sau:
\choiceTF
{\True Cả hai bạn đã đổ tất cả $\dfrac{7}{10}\text{ l}$ nước vào bình.}
{Mai đổ ít nước hơn Việt.}
{\True Mai đổ nhiều hơn Việt $\dfrac{3}{10}\text{ l}$ nước.}
{Nếu dung tích của bình là $1\text{ l}$ thì lượng nước còn thiếu để đầy bình là $\dfrac{1}{3}\text{ l}$.}
\loigiai{
\begin{itemize}
    \item Cả hai bạn đổ: $\dfrac{1}{5} + \dfrac{1}{2} = \dfrac{2}{10} + \dfrac{5}{10} = \dfrac{7}{10}\text{ (l)}$.
    \item Vì $\dfrac{1}{2} = \dfrac{5}{10} > \dfrac{2}{10} = \dfrac{1}{5}$ nên Mai đổ nhiều hơn Việt.
    \item Lượng nước Mai đổ nhiều hơn Việt là: $\dfrac{1}{2} - \dfrac{1}{5} = \dfrac{3}{10}\text{ (l)}$.
    \item Lượng nước còn thiếu để đầy bình $1\text{ l}$ là: $1 - \dfrac{7}{10} = \dfrac{3}{10}\text{ (l)} \neq \dfrac{1}{3}\text{ l}$.
\end{itemize}
\begin{itemize}
    \item a) \textbf{Đúng}.
    \item b) \textbf{Sai}: Mai đổ nhiều hơn Việt.
    \item c) \textbf{Đúng}.
    \item d) \textbf{Sai}: Còn thiếu $\dfrac{3}{10}\text{ l}$.
\end{itemize}
}
\end{ex}

\textbf{Phần 3. Trả lời ngắn}
\setcounterpageref{ex}{0}

\begin{ex} %[5D1H6]
Tính tổng $\dfrac{3}{10} + \dfrac{5}{21}$ dưới dạng phân số tối giản $\dfrac{a}{b}$. Giá trị của tử số $a$ bằng bao nhiêu?
\shortans{113}
\loigiai{
\textbf{Đáp số:} $113$\\
$\dfrac{3}{10} + \dfrac{5}{21} = \dfrac{3 \times 21 + 5 \times 10}{210} = \dfrac{63 + 50}{210} = \dfrac{113}{210}$. Tử số $a = 113$.
}
\end{ex}

\begin{ex} %[5D1H6]
Tính hiệu $\dfrac{31}{10} - 2$ dưới dạng phân số tối giản $\dfrac{m}{n}$. Tính giá trị của $m + n$.
\shortans{21}
\loigiai{
\textbf{Đáp số:} $21$\\
$\dfrac{31}{10} - 2 = \dfrac{31}{10} - \dfrac{20}{10} = \dfrac{11}{10}$. Khi đó $m = 11, n = 10 \Rightarrow m + n = 21$.
}
\end{ex}

\begin{ex} %[5D1H6]
Tính tổng $\dfrac{14}{5} + 1$ dưới dạng phân số tối giản $\dfrac{p}{q}$. Tìm giá trị của tử số $p$.
\shortans{19}
\loigiai{
\textbf{Đáp số:} $19$\\
$\dfrac{14}{5} + 1 = \dfrac{14}{5} + \dfrac{5}{5} = \dfrac{19}{5}$. Tử số $p = 19$.
}
\end{ex}

\begin{ex} %[5D1V6]
Có hai vòi nước cùng chảy vào một bể cạn. Vòi thứ nhất chảy được $\dfrac{1}{2}$ bể, vòi thứ hai chảy được $\dfrac{2}{5}$ bể. Hỏi còn lại bao nhiêu phần bể chưa có nước? (Viết kết quả dưới dạng phân số tối giản $\dfrac{1}{k}$). Giá trị của $k$ là bao nhiêu?
\shortans{10}
\loigiai{
\textbf{Đáp số:} $10$\\
Cả hai vòi chảy được: $\dfrac{1}{2} + \dfrac{2}{5} = \dfrac{9}{10}$ (bể).\\
Phần bể chưa có nước là: $1 - \dfrac{9}{10} = \dfrac{1}{10}$ (bể). Vậy $k = 10$.
}
\end{ex}

\begin{ex} %[5D1V6]
Mai và Việt cùng đổ nước vào một bình. Mai đổ $\dfrac{1}{2}\text{ l}$, Việt đổ $\dfrac{1}{5}\text{ l}$. Lượng nước Mai đổ nhiều hơn Việt là $\dfrac{x}{10}\text{ l}$. Giá trị của $x$ là bao nhiêu?
\shortans{3}
\loigiai{
\textbf{Đáp số:} $3$\\
$\dfrac{1}{2} - \dfrac{1}{5} = \dfrac{5}{10} - \dfrac{2}{10} = \dfrac{3}{10}\text{ (l)} \Rightarrow x = 3$.
}
\end{ex}

\begin{ex} %[5D1V6]
Rô-bốt chia bánh cho bạn bằng cách lấy $\dfrac{1}{2}$ cái bánh cộng thêm $\dfrac{1}{3}$ cái bánh. Tổng số bánh bạn nhận được là $\dfrac{u}{6}$ cái bánh. Giá trị của $u$ là bao nhiêu?
\shortans{5}
\loigiai{
\textbf{Đáp số:} $5$\\
$\dfrac{1}{2} + \dfrac{1}{3} = \dfrac{3}{6} + \dfrac{2}{6} = \dfrac{5}{6}\text{ (cái bánh)} \Rightarrow u = 5$.
}
\end{ex}

\newpage

\TaoPhieuKiemTra{Chủ đề 1 -- Ôn tập và bổ sung}{07}{Hỗn số}{https://loptoanthaylinh.vn/thi-online/toan-5/bai-07}
\textbf{Phần 1. Trắc nghiệm lựa chọn}
\setcounterpageref{ex}{0}

\begin{ex} %[5D1N7]
Hỗn số $1\dfrac{1}{4}$ được đọc là
\choice
{\True Một và một phần tư}
{Một phần tư}
{Mười một phần tư}
{Một chia một phần tư}
\loigiai{
Hỗn số $1\dfrac{1}{4}$ gồm phần nguyên là $1$ và phần phân số là $\dfrac{1}{4}$, đọc là: Một và một phần tư.\\
\textbf{Đáp án của câu hỏi là A}
}
\end{ex}

\begin{ex} %[5D1N7]
Trong hỗn số $12\dfrac{7}{10}$, phần nguyên và phần phân số lần lượt là
\choice
{$7$ và $\dfrac{12}{10}$}
{$12$ và $\dfrac{10}{7}$}
{\True $12$ và $\dfrac{7}{10}$}
{$10$ và $\dfrac{7}{12}$}
\loigiai{
Hỗn số $12\dfrac{7}{10}$ có phần nguyên là $12$ và phần phân số là $\dfrac{7}{10}$.\\
\textbf{Đáp án của câu hỏi là C}
}
\end{ex}

\begin{ex} %[5D1H7]
Hỗn số $8\dfrac{1}{3}$ được viết dưới dạng tổng của phần nguyên và phần phân số là
\choice
{$8 \times \dfrac{1}{3}$}
{$8 - \dfrac{1}{3}$}
{$8 : \dfrac{1}{3}$}
{\True $8 + \dfrac{1}{3}$}
\loigiai{
Theo định nghĩa hỗn số, $8\dfrac{1}{3} = 8 + \dfrac{1}{3}$.\\
\textbf{Đáp án của câu hỏi là D}
}
\end{ex}

\begin{ex} %[5D1H7]
Chuyển hỗn số $2\dfrac{7}{10}$ thành phân số thập phân, ta được kết quả là
\choice
{$\dfrac{14}{10}$}
{\True $\dfrac{27}{10}$}
{$\dfrac{20}{10}$}
{$\dfrac{72}{10}$}
\loigiai{
Ta có: $2\dfrac{7}{10} = \dfrac{2 \times 10 + 7}{10} = \dfrac{27}{10}$.\\
\textbf{Đáp án của câu hỏi là B}
}
\end{ex}

\begin{ex} %[5D1H7]
Chuyển hỗn số $5\dfrac{1}{10}$ thành phân số thập phân ta được
\choice
{$\dfrac{6}{10}$}
{$\dfrac{15}{10}$}
{\True $\dfrac{51}{10}$}
{$\dfrac{50}{10}$}
\loigiai{
Ta có: $5\dfrac{1}{10} = \dfrac{5 \times 10 + 1}{10} = \dfrac{51}{10}$.\\
\textbf{Đáp án của câu hỏi là C}
}
\end{ex}

\begin{ex} %[5D1H7]
Chuyển hỗn số $1\dfrac{9}{100}$ thành phân số thập phân, ta được
\choice
{\True $\dfrac{109}{100}$}
{$\dfrac{19}{100}$}
{$\dfrac{91}{100}$}
{$\dfrac{190}{100}$}
\loigiai{
Ta có: $1\dfrac{9}{100} = \dfrac{1 \times 100 + 9}{100} = \dfrac{109}{100}$.\\
\textbf{Đáp án của câu hỏi là A}
}
\end{ex}

\begin{ex} %[5D1H7]
Chuyển hỗn số $4\dfrac{9}{100}$ thành phân số thập phân, ta được kết quả là
\choice
{$\dfrac{49}{100}$}
{$\dfrac{490}{100}$}
{$\dfrac{400}{100}$}
{\True $\dfrac{409}{100}$}
\loigiai{
Ta có: $4\dfrac{9}{100} = \dfrac{4 \times 100 + 9}{100} = \dfrac{409}{100}$.\\
\textbf{Đáp án của câu hỏi là D}
}
\end{ex}

\begin{ex} %[5D1H7]
Viết phân số thập phân $\dfrac{57}{10}$ dưới dạng hỗn số, ta được
\choice
{$7\dfrac{5}{10}$}
{$57\dfrac{1}{10}$}
{\True $5\dfrac{7}{10}$}
{$5\dfrac{1}{7}$}
\loigiai{
Thực hiện phép chia: $57 : 10 = 5$ dư $7$. Do đó $\dfrac{57}{10} = 5\dfrac{7}{10}$.\\
\textbf{Đáp án của câu hỏi là C}
}
\end{ex}

\begin{ex} %[5D1H7]
Viết phân số thập phân $\dfrac{703}{100}$ dưới dạng hỗn số, ta được
\choice
{$70\dfrac{3}{100}$}
{\True $7\dfrac{3}{100}$}
{$7\dfrac{30}{100}$}
{$73\dfrac{1}{100}$}
\loigiai{
Ta có: $703 : 100 = 7$ dư $3$. Vậy $\dfrac{703}{100} = 7\dfrac{3}{100}$.\\
\textbf{Đáp án của câu hỏi là B}
}
\end{ex}

\begin{ex} %[5D1H7]
Viết phân số $\dfrac{3\,005}{1\,000}$ dưới dạng hỗn số, ta được
\choice
{$30\dfrac{5}{1\,000}$}
{$3\dfrac{50}{1\,000}$}
{$300\dfrac{5}{1\,000}$}
{\True $3\dfrac{5}{1\,000}$}
\loigiai{
Ta có: $3\,005 : 1\,000 = 3$ dư $5$. Do đó $\dfrac{3\,005}{1\,000} = 3\dfrac{5}{1\,000}$.\\
\textbf{Đáp án của câu hỏi là D}
}
\end{ex}

\begin{ex} %[5D1H7]
Trên tia số giữa hai hỗn số $6\dfrac{2}{10}$ và $6\dfrac{4}{10}$ là vạch biểu diễn hỗn số nào?
\choice
{\True $6\dfrac{3}{10}$}
{$6\dfrac{1}{10}$}
{$6\dfrac{5}{10}$}
{$7\dfrac{3}{10}$}
\loigiai{
Trên tia số có các khoảng chia đều $\dfrac{1}{10}$, điểm nằm giữa $6\dfrac{2}{10}$ và $6\dfrac{4}{10}$ là $6\dfrac{3}{10}$.\\
\textbf{Đáp án của câu hỏi là A}
}
\end{ex}

\begin{ex} %[5D1V7]
Cô Dung có $23$ phong kẹo, mỗi phong kẹo có $10$ viên kẹo. Cô chia đều số kẹo đó cho $10$ bạn. Khẳng định nào sau đây là \textbf{sai}?
\choice
{Mỗi bạn nhận được $23$ viên kẹo}
{\True Mỗi bạn nhận được $2\dfrac{3}{10}$ viên kẹo}
{Mỗi bạn nhận được $\dfrac{23}{10}$ phong kẹo}
{Mỗi bạn nhận được $2\dfrac{3}{10}$ phong kẹo}
\loigiai{
Tổng số viên kẹo là: $23 \times 10 = 230$ viên.\\
Mỗi bạn được: $230 : 10 = 23$ viên kẹo (hoặc $23 : 10 = \dfrac{23}{10} = 2\dfrac{3}{10}$ phong kẹo).\\
Khẳng định ``Mỗi bạn nhận được $2\dfrac{3}{10}$ viên kẹo'' là sai về đơn vị.\\
\textbf{Đáp án của câu hỏi là B}
}
\end{ex}

\textbf{Phần 2. Trắc nghiệm đúng sai}
\setcounterpageref{ex}{0}

\begin{ex} %[5D1H7]
Cho các hỗn số: $A = 5\dfrac{1}{10}$, $B = 1\dfrac{9}{100}$ và các phân số thập phân $M = \dfrac{63}{10}$, $N = \dfrac{703}{100}$. Xét tính đúng sai của các khẳng định sau:
\choiceTF
{\True Chuyển hỗn số $A$ thành phân số thập phân ta được $\dfrac{51}{10}$.}
{Chuyển hỗn số $B$ thành phân số thập phân ta được $\dfrac{19}{100}$.}
{\True Phân số $M$ được viết dưới dạng hỗn số là $6\dfrac{3}{10}$.}
{\True Phân số $N$ được viết dưới dạng hỗn số là $7\dfrac{3}{100}$.}
\loigiai{
\begin{itemize}
    \item $A = 5\dfrac{1}{10} = \dfrac{5 \times 10 + 1}{10} = \dfrac{51}{10}$.
    \item $B = 1\dfrac{9}{100} = \dfrac{1 \times 100 + 9}{100} = \dfrac{109}{100} \neq \dfrac{19}{100}$.
    \item $M = \dfrac{63}{10} = 6 + \dfrac{3}{10} = 6\dfrac{3}{10}$.
    \item $N = \dfrac{703}{100} = 7 + \dfrac{3}{100} = 7\dfrac{3}{100}$.
\end{itemize}
\begin{itemize}
    \item a) \textbf{Đúng}.
    \item b) \textbf{Sai}: $B = \dfrac{109}{100}$.
    \item c) \textbf{Đúng}.
    \item d) \textbf{Đúng}.
\end{itemize}
}
\end{ex}

\begin{ex} %[5D1V7]
Cô Dung chia đều $5$ chiếc bánh cho $4$ bạn học sinh. Xét tính đúng sai của các khẳng định sau:
\choiceTF
{\True Số phần bánh mỗi bạn nhận được là $\dfrac{5}{4}$ cái bánh.}
{Viết dưới dạng hỗn số, mỗi bạn nhận được $1\dfrac{1}{5}$ cái bánh.}
{\True Mỗi bạn nhận được $1$ cái bánh nguyên và thêm $\dfrac{1}{4}$ cái bánh nữa.}
{Để mỗi bạn nhận được đủ $2$ cái bánh thì cô Dung cần thêm $4$ cái bánh nữa.}
\loigiai{
\begin{itemize}
    \item Chia $5$ cái bánh cho $4$ bạn: mỗi bạn được $5 : 4 = \dfrac{5}{4}$ cái bánh.
    \item Ta có $\dfrac{5}{4} = 1\dfrac{1}{4} \neq 1\dfrac{1}{5}$.
    \item Hỗn số $1\dfrac{1}{4}$ gồm phần nguyên là $1$ cái bánh và phần phân số là $\dfrac{1}{4}$ cái bánh.
    \item Để $4$ bạn mỗi bạn có $2$ cái bánh thì cần $4 \times 2 = 8$ cái bánh. Số bánh cần thêm là $8 - 5 = 3$ cái bánh ($\neq 4$ cái bánh).
\end{itemize}
\begin{itemize}
    \item a) \textbf{Đúng}.
    \item b) \textbf{Sai}: Mỗi bạn nhận được $1\dfrac{1}{4}$ cái bánh.
    \item c) \textbf{Đúng}.
    \item d) \textbf{Sai}: Chỉ cần thêm $3$ cái bánh.
\end{itemize}
}
\end{ex}

\textbf{Phần 3. Trả lời ngắn}
\setcounterpageref{ex}{0}

\begin{ex} %[5D1H7]
Chuyển hỗn số $3\dfrac{5}{10}$ thành phân số thập phân có mẫu số là $10$. Hỏi tử số của phân số đó là bao nhiêu?
\shortans{35}
\loigiai{
\textbf{Đáp số:} $35$\\
$3\dfrac{5}{10} = \dfrac{3 \times 10 + 5}{10} = \dfrac{35}{10}$. Tử số là $35$.
}
\end{ex}

\begin{ex} %[5D1H7]
Chuyển hỗn số $4\dfrac{9}{100}$ thành phân số thập phân có mẫu số là $100$. Tử số của phân số nhận được bằng bao nhiêu?
\shortans{409}
\loigiai{
\textbf{Đáp số:} $409$\\
$4\dfrac{9}{100} = \dfrac{4 \times 100 + 9}{100} = \dfrac{409}{100}$. Tử số là $409$.
}
\end{ex}

\begin{ex} %[5D1H7]
Viết phân số thập phân $\dfrac{703}{100}$ dưới dạng hỗn số $a\dfrac{b}{100}$. Tính tích $a \times b$.
\shortans{21}
\loigiai{
\textbf{Đáp số:} $21$\\
Ta có $\dfrac{703}{100} = 7\dfrac{3}{100} \Rightarrow a = 7, b = 3 \Rightarrow a \times b = 21$.
}
\end{ex}

\begin{ex} %[5D1H7]
Viết phân số thập phân $\dfrac{3\,005}{1\,000}$ dưới dạng hỗn số $m\dfrac{n}{1\,000}$. Tính giá trị của $m + n$.
\shortans{8}
\loigiai{
\textbf{Đáp số:} $8$\\
$\dfrac{3\,005}{1\,000} = 3\dfrac{5}{1\,000} \Rightarrow m = 3, n = 5 \Rightarrow m + n = 8$.
}
\end{ex}

\begin{ex} %[5D1V7]
Cô Dung có $23$ phong kẹo, mỗi phong có $10$ viên kẹo. Cô chia đều cho $10$ bạn. Hỏi mỗi bạn nhận được bao nhiêu viên kẹo?
\shortans{23}
\loigiai{
\textbf{Đáp số:} $23$\\
Tổng số kẹo là $23 \times 10 = 230$ viên. Mỗi bạn nhận được $230 : 10 = 23$ viên kẹo.
}
\end{ex}

\begin{ex} %[5D1V7]
Có $17$ lít nước chia đều vào $5$ chiếc can giống nhau. Lượng nước trong mỗi can được viết dưới dạng hỗn số là $a\dfrac{b}{c}\text{ (lít)}$ (với $\dfrac{b}{c}$ là phân số tối giản). Tính giá trị của $a + b + c$.
\shortans{10}
\loigiai{
\textbf{Đáp số:} $10$\\
Mỗi can chứa: $17 : 5 = \dfrac{17}{5} = 3\dfrac{2}{5}\text{ (lít)}$.\\
Do đó $a = 3, b = 2, c = 5 \Rightarrow a + b + c = 3 + 2 + 5 = 10$.
}
\end{ex}

\newpage

\TaoPhieuKiemTra{Chủ đề 1 -- Ôn tập và bổ sung}{08}{Ôn tập hình học và đo lường}{https://loptoanthaylinh.vn/thi-online/toan-5/bai-08}
\textbf{Phần 1. Trắc nghiệm lựa chọn}
\setcounterpageref{ex}{0}

\begin{ex} %[5H1N8]
Số thích hợp điền vào dấu hỏi chấm để $6\text{ yến} = ?\text{ kg}$ là
\choice
{$6$}
{$600$}
{$6\,000$}
{\True $60$}
\loigiai{
Vì $1\text{ yến} = 10\text{ kg}$ nên $6\text{ yến} = 60\text{ kg}$.\\
\textbf{Đáp án của câu hỏi là D}
}
\end{ex}

\begin{ex} %[5H1N8]
Số thích hợp điền vào dấu hỏi chấm để $2\text{ tấn} = ?\text{ kg}$ là
\choice
{\True $2\,000$}
{$20$}
{$200$}
{$20\,000$}
\loigiai{
Vì $1\text{ tấn} = 1\,000\text{ kg}$ nên $2\text{ tấn} = 2\,000\text{ kg}$.\\
\textbf{Đáp án của câu hỏi là A}
}
\end{ex}

\begin{ex} %[5H1H8]
Số thích hợp điền vào chỗ chấm để $1\text{ tấn} = \dots\text{ yến}$ là
\choice
{$10$}
{$1\,000$}
{$10\,000$}
{\True $100$}
\loigiai{
Vì $1\text{ tấn} = 1\,000\text{ kg}$ và $1\text{ yến} = 10\text{ kg}$ nên $1\text{ tấn} = 100\text{ yến}$.\\
\textbf{Đáp án của câu hỏi là D}
}
\end{ex}

\begin{ex} %[5H1H8]
Góc có số đo bằng $120^\circ$ là
\choice
{Góc nhọn}
{\True Góc tù}
{Góc vuông}
{Góc bẹt}
\loigiai{
Góc lớn hơn $90^\circ$ và nhỏ hơn $180^\circ$ là góc tù. Do đó góc $120^\circ$ là góc tù.\\
\textbf{Đáp án của câu hỏi là B}
}
\end{ex}

\begin{ex} %[5H1H8]
Số thích hợp điền vào dấu hỏi chấm để $5\text{ giờ} = ?\text{ phút}$ là
\choice
{$50$}
{$250$}
{\True $300$}
{$360$}
\loigiai{
Vì $1\text{ giờ} = 60\text{ phút}$ nên $5\text{ giờ} = 5 \times 60 = 300\text{ phút}$.\\
\textbf{Đáp án của câu hỏi là C}
}
\end{ex}

\begin{ex} %[5H1H8]
Đổi $\dfrac{1}{10}\text{ giờ}$ ra phút ta được kết quả là
\choice
{$10\text{ phút}$}
{$12\text{ phút}$}
{$15\text{ phút}$}
{\True $6\text{ phút}$}
\loigiai{
Ta có: $\dfrac{1}{10}\text{ giờ} = 60 \times \dfrac{1}{10} = 6\text{ phút}$.\\
\textbf{Đáp án của câu hỏi là D}
}
\end{ex}

\begin{ex} %[5H1H8]
Đổi $2\text{ giờ } 30\text{ phút}$ ra phút ta được
\choice
{$120\text{ phút}$}
{$130\text{ phút}$}
{\True $150\text{ phút}$}
{$230\text{ phút}$}
\loigiai{
Ta có: $2\text{ giờ } 30\text{ phút} = 2 \times 60 + 30 = 120 + 30 = 150\text{ phút}$.\\
\textbf{Đáp án của câu hỏi là C}
}
\end{ex}

\begin{ex} %[5H1H8]
Đổi $4\text{ phút } 5\text{ giây}$ ra giây ta được
\choice
{$45\text{ giây}$}
{\True $245\text{ giây}$}
{$205\text{ giây}$}
{$405\text{ giây}$}
\loigiai{
Ta có: $4\text{ phút } 5\text{ giây} = 4 \times 60 + 5 = 240 + 5 = 245\text{ giây}$.\\
\textbf{Đáp án của câu hỏi là B}
}
\end{ex}

\begin{ex} %[5H1H8]
Đổi $\dfrac{7}{100}\text{ thế kỉ}$ ra năm ta được kết quả là
\choice
{\True $7\text{ năm}$}
{$70\text{ năm}$}
{$700\text{ năm}$}
{$14\text{ năm}$}
\loigiai{
Vì $1\text{ thế kỉ} = 100\text{ năm}$ nên $\dfrac{7}{100}\text{ thế kỉ} = 100 \times \dfrac{7}{100} = 7\text{ năm}$.\\
\textbf{Đáp án của câu hỏi là A}
}
\end{ex}

\begin{ex} %[5H1V8]
Bác Năm thu hoạch được $1\text{ tấn } 250\text{ kg}$ cam. Biết rằng số cam loại I chiếm $\dfrac{3}{10}$ tổng số cam thu hoạch. Số ki-lô-gam cam loại I là
\choice
{$250\text{ kg}$}
{$450\text{ kg}$}
{\True $375\text{ kg}$}
{$875\text{ kg}$}
\loigiai{
Đổi $1\text{ tấn } 250\text{ kg} = 1\,250\text{ kg}$.\\
Số cam loại I là: $1\,250 \times \dfrac{3}{10} = 375\text{ (kg)}$.\\
\textbf{Đáp án của câu hỏi là C}
}
\end{ex}

\begin{ex} %[5H1V8]
Dì Sáu có một mảnh đất hình chữ nhật có chiều dài $12\text{ m}$, chiều rộng $7\text{ m}$. Diện tích của mảnh đất đó là
\choice
{\True $84\text{ m}^2$}
{$38\text{ m}^2$}
{$72\text{ m}^2$}
{$96\text{ m}^2$}
\loigiai{
Diện tích mảnh đất hình chữ nhật là: $12 \times 7 = 84\text{ (m}^2\text{)}$.\\
\textbf{Đáp án của câu hỏi là A}
}
\end{ex}

\begin{ex} %[5H1V8]
Dì Sáu chia đều mảnh đất hình chữ nhật diện tích $84\text{ m}^2$ thành $7$ ô đất bằng nhau để xây phòng trọ. Diện tích mỗi phòng trọ là
\choice
{$7\text{ m}^2$}
{\True $12\text{ m}^2$}
{$10\text{ m}^2$}
{$14\text{ m}^2$}
\loigiai{
Diện tích mỗi ô đất là: $84 : 7 = 12\text{ (m}^2\text{)}$.\\
\textbf{Đáp án của câu hỏi là B}
}
\end{ex}

\textbf{Phần 2. Trắc nghiệm đúng sai}
\setcounterpageref{ex}{0}

\begin{ex} %[5H1H8]
Xét tính đúng sai của các khẳng định sau về đổi đơn vị đo:
\choiceTF
{\True $5\text{ tấn} = 50\text{ tạ}$.}
{$9\text{ tạ} = 900\text{ yến}$.}
{\True $7\text{ thế kỉ} = 700\text{ năm}$.}
{$\dfrac{1}{10}\text{ giờ} = 10\text{ phút}$.}
\loigiai{
\begin{itemize}
    \item $5\text{ tấn} = 5 \times 10 = 50\text{ tạ}$.
    \item $9\text{ tạ} = 9 \times 10 = 90\text{ yến} \neq 900\text{ yến}$.
    \item $7\text{ thế kỉ} = 7 \times 100 = 700\text{ năm}$.
    \item $\dfrac{1}{10}\text{ giờ} = 60 : 10 = 6\text{ phút} \neq 10\text{ phút}$.
\end{itemize}
\begin{itemize}
    \item a) \textbf{Đúng}.
    \item b) \textbf{Sai}: $9\text{ tạ} = 90\text{ yến}$.
    \item c) \textbf{Đúng}.
    \item d) \textbf{Sai}: $\dfrac{1}{10}\text{ giờ} = 6\text{ phút}$.
\end{itemize}
}
\end{ex}

\begin{ex} %[5H1V8]
Bác Năm thu hoạch được $1\text{ tấn } 250\text{ kg}$ cam gồm hai loại: cam loại I và cam loại II. Biết cam loại I chiếm $\dfrac{3}{10}$ tổng số cam thu hoạch. Xét tính đúng sai của các khẳng định sau:
\choiceTF
{\True Tổng khối lượng cam bác Năm thu hoạch được là $1\,250\text{ kg}$.}
{Khối lượng cam loại I là $250\text{ kg}$.}
{\True Khối lượng cam loại II là $875\text{ kg}$.}
{\True Khối lượng cam loại II nhiều hơn khối lượng cam loại I là $500\text{ kg}$.}
\loigiai{
\begin{itemize}
    \item Đổi $1\text{ tấn } 250\text{ kg} = 1\,250\text{ kg}$.
    \item Số cam loại I là: $1\,250 \times \dfrac{3}{10} = 375\text{ kg} \neq 250\text{ kg}$.
    \item Số cam loại II là: $1\,250 - 375 = 875\text{ kg}$.
    \item Hiệu giữa cam loại II và loại I là: $875 - 375 = 500\text{ kg}$.
\end{itemize}
\begin{itemize}
    \item a) \textbf{Đúng}.
    \item b) \textbf{Sai}: Khối lượng cam loại I là $375\text{ kg}$.
    \item c) \textbf{Đúng}.
    \item d) \textbf{Đúng}.
\end{itemize}
}
\end{ex}

\textbf{Phần 3. Trả lời ngắn}
\setcounterpageref{ex}{0}

\begin{ex} %[5H1H8]
Đổi $9\text{ tạ}$ ra yến. Kết quả là bao nhiêu yến?
\shortans{90}
\loigiai{
\textbf{Đáp số:} $90$\\
$9\text{ tạ} = 9 \times 10 = 90\text{ yến}$.
}
\end{ex}

\begin{ex} %[5H1H8]
Đổi $4\text{ phút } 5\text{ giây}$ ra đơn vị giây.
\shortans{245}
\loigiai{
\textbf{Đáp số:} $245$\\
$4\text{ phút } 5\text{ giây} = 4 \times 60 + 5 = 245\text{ giây}$.
}
\end{ex}

\begin{ex} %[5H1H8]
Trong các góc có số đo lần lượt là $60^\circ, 90^\circ, 120^\circ, 45^\circ, 150^\circ$, có bao nhiêu góc là góc tù?
\shortans{2}
\loigiai{
\textbf{Đáp số:} $2$\\
Các góc tù (lớn hơn $90^\circ$ và nhỏ hơn $180^\circ$) là: $120^\circ$ và $150^\circ$. Có tất cả $2$ góc tù.
}
\end{ex}

\begin{ex} %[5H1V8]
Bác Năm thu hoạch được $1\text{ tấn } 250\text{ kg}$ cam. Biết số cam loại I chiếm $\dfrac{3}{10}$ tổng số cam, còn lại là cam loại II. Hỏi bác Năm thu hoạch được bao nhiêu ki-lô-gam cam loại II?
\shortans{875}
\loigiai{
\textbf{Đáp số:} $875$\\
Tổng số cam là $1\,250\text{ kg}$. Số cam loại I là $1\,250 \times \dfrac{3}{10} = 375\text{ kg}$.\\
Số cam loại II là: $1\,250 - 375 = 875\text{ kg}$.
}
\end{ex}

\begin{ex} %[5H1V8]
Mảnh đất hình chữ nhật có chiều dài $12\text{ m}$ và chiều rộng $7\text{ m}$ được chia thành $7$ ô đất nhỏ có diện tích bằng nhau dạng hình chữ nhật có chiều dài $4\text{ m}$. Hỏi chiều rộng của mỗi ô đất nhỏ là bao nhiêu mét?
\shortans{3}
\loigiai{
\textbf{Đáp số:} $3$\\
Diện tích mỗi ô nhỏ là: $(12 \times 7) : 7 = 12\text{ (m}^2\text{)}$.\\
Chiều rộng mỗi ô nhỏ là: $12 : 4 = 3\text{ (m)}$.
}
\end{ex}

\begin{ex} %[5H1H8]
Đổi $\dfrac{7}{100}\text{ thế kỉ}$ ra đơn vị năm.
\shortans{7}
\loigiai{
\textbf{Đáp số:} $7$\\
$1\text{ thế kỉ} = 100\text{ năm} \Rightarrow \dfrac{7}{100} \times 100 = 7\text{ năm}$.
}
\end{ex}

\newpage

\TaoPhieuKiemTra{Chủ đề 1 -- Ôn tập và bổ sung}{09}{Luyện tập chung}{https://loptoanthaylinh.vn/thi-online/toan-5/bai-09}
\textbf{Phần 1. Trắc nghiệm lựa chọn}
\setcounterpageref{ex}{0}

\begin{ex} %[5D1N9]
Giá trị của chữ số $6$ trong số $960\,730$ là
\choice
{\True $60\,000$}
{$60$}
{$600$}
{$6\,000$}
\loigiai{
Trong số $960\,730$, chữ số $6$ ở hàng chục nghìn nên có giá trị là $60\,000$.\\
\textbf{Đáp án của câu hỏi là A}
}
\end{ex}

\begin{ex} %[5D1N9]
Số lớn nhất trong các số $109\,989; 105\,789; 110\,200; 99\,000$ là
\choice
{$109\,989$}
{$105\,789$}
{\True $110\,200$}
{$99\,000$}
\loigiai{
So sánh các số: $99\,000 < 105\,789 < 109\,989 < 110\,200$. Vậy số lớn nhất là $110\,200$.\\
\textbf{Đáp án của câu hỏi là C}
}
\end{ex}

\begin{ex} %[5D1H9]
Phân số bé nhất trong các phân số $\dfrac{23}{36}; \dfrac{11}{12}; \dfrac{7}{9}; \dfrac{2}{3}$ là
\choice
{$\dfrac{11}{12}$}
{$\dfrac{7}{9}$}
{$\dfrac{2}{3}$}
{\True $\dfrac{23}{36}$}
\loigiai{
Quy đồng mẫu số chung là $36$:\\
$\dfrac{11}{12} = \dfrac{33}{36}$; $\dfrac{7}{9} = \dfrac{28}{36}$; $\dfrac{2}{3} = \dfrac{24}{36}$; $\dfrac{23}{36}$.\\
Vì $\dfrac{23}{36} < \dfrac{24}{36} < \dfrac{28}{36} < \dfrac{33}{36}$ nên phân số bé nhất là $\dfrac{23}{36}$.\\
\textbf{Đáp án của câu hỏi là D}
}
\end{ex}

\begin{ex} %[5D1H9]
Phân số $\dfrac{143}{100}$ viết dưới dạng hỗn số là
\choice
{\True $1\dfrac{43}{100}$}
{$14\dfrac{3}{100}$}
{$140\dfrac{3}{100}$}
{$3\dfrac{14}{100}$}
\loigiai{
Ta có: $143 : 100 = 1$ dư $43$. Do đó $\dfrac{143}{100} = 1\dfrac{43}{100}$.\\
\textbf{Đáp án của câu hỏi là A}
}
\end{ex}

\begin{ex} %[5D1H9]
Giá trị của biểu thức $35\,700 : 50 + 68 \times 46$ là
\choice
{$3\,128$}
{$3\,742$}
{\True $3\,842$}
{$4\,128$}
\loigiai{
Ta có: $35\,700 : 50 = 714$ và $68 \times 46 = 3\,128$.\\
Tổng là: $714 + 3\,128 = 3\,842$.\\
\textbf{Đáp án của câu hỏi là C}
}
\end{ex}

\begin{ex} %[5D1H9]
Giá trị của biểu thức $\dfrac{16}{9} \times \left(6 : \dfrac{32}{9}\right) - \dfrac{7}{5}$ là
\choice
{$3$}
{\True $\dfrac{8}{5}$}
{$\dfrac{7}{5}$}
{$\dfrac{1}{5}$}
\loigiai{
Ta có: $6 : \dfrac{32}{9} = 6 \times \dfrac{9}{32} = \dfrac{27}{16}$.\\
Khi đó: $\dfrac{16}{9} \times \dfrac{27}{16} - \dfrac{7}{5} = 3 - \dfrac{7}{5} = \dfrac{8}{5}$.\\
\textbf{Đáp án của câu hỏi là B}
}
\end{ex}

\begin{ex} %[5D1H9]
Viết số đo $52\text{ cm}$ dưới dạng phân số thập phân với đơn vị mét là
\choice
{\True $\dfrac{52}{100}\text{ m}$}
{$\dfrac{52}{10}\text{ m}$}
{$\dfrac{52}{1\,000}\text{ m}$}
{$\dfrac{52}{100}\text{ cm}$}
\loigiai{
Vì $1\text{ m} = 100\text{ cm}$ nên $52\text{ cm} = \dfrac{52}{100}\text{ m}$.\\
\textbf{Đáp án của câu hỏi là A}
}
\end{ex}

\begin{ex} %[5D1H9]
Viết số đo $2\text{ dm } 7\text{ cm}$ dưới dạng hỗn số với đơn vị đề-xi-mét là
\choice
{$7\dfrac{2}{10}\text{ dm}$}
{$27\dfrac{1}{10}\text{ dm}$}
{$2\dfrac{7}{100}\text{ dm}$}
{\True $2\dfrac{7}{10}\text{ dm}$}
\loigiai{
Ta có: $2\text{ dm } 7\text{ cm} = 2\text{ dm} + \dfrac{7}{10}\text{ dm} = 2\dfrac{7}{10}\text{ dm}$.\\
\textbf{Đáp án của câu hỏi là D}
}
\end{ex}

\begin{ex} %[5D1H9]
Tính nhanh biểu thức $524 \times 63 + 524 \times 37 - 2\,400$ ta được
\choice
{$5\,000$}
{\True $50\,000$}
{$52\,400$}
{$500\,000$}
\loigiai{
$524 \times (63 + 37) - 2\,400 = 524 \times 100 - 2\,400 = 52\,400 - 2\,400 = 50\,000$.\\
\textbf{Đáp án của câu hỏi là B}
}
\end{ex}

\begin{ex} %[5D1V9]
Một bút bi có giá $4\,500$ đồng, một quyển vở có giá $7\,000$ đồng. Nam mua $2$ bút bi và $7$ quyển vở rồi đưa cho cô bán hàng tờ tiền $100\,000$ đồng. Cô bán hàng trả lại Nam số tiền là
\choice
{$58\,000$ đồng}
{$48\,000$ đồng}
{\True $42\,000$ đồng}
{$52\,000$ đồng}
\loigiai{
Số tiền Nam mua bút và vở là: $2 \times 4\,500 + 7 \times 7\,000 = 9\,000 + 49\,000 = 58\,000$ đồng.\\
Số tiền cô bán hàng trả lại là: $100\,000 - 58\,000 = 42\,000$ đồng.\\
\textbf{Đáp án của câu hỏi là C}
}
\end{ex}

\begin{ex} %[5D1V9]
Một đoàn xe gồm $6$ xe ô tô loại nhỏ chở $35$ học sinh/xe và $9$ xe loại lớn chở $40$ học sinh/xe. Trung bình mỗi xe chở bao nhiêu học sinh?
\choice
{$37$ học sinh}
{$39$ học sinh}
{$37{,}5$ học sinh}
{\True $38$ học sinh}
\loigiai{
Tổng số học sinh là: $6 \times 35 + 9 \times 40 = 210 + 360 = 570$ học sinh.\\
Tổng số xe là: $6 + 9 = 15$ xe.\\
Trung bình mỗi xe chở: $570 : 15 = 38$ học sinh.\\
\textbf{Đáp án của câu hỏi là D}
}
\end{ex}

\begin{ex} %[5D1V9]
Cô Ba mang $120$ quả trứng ra chợ bán. Lần thứ nhất cô bán được $\dfrac{1}{8}$ số trứng đó. Lần thứ hai cô bán được $\dfrac{2}{7}$ số trứng còn lại sau lần bán thứ nhất. Hỏi sau cả hai lần, cô Ba đã bán được tất cả bao nhiêu quả trứng?
\choice
{$30$ quả}
{\True $45$ quả}
{$75$ quả}
{$40$ quả}
\loigiai{
Lần thứ nhất bán được: $120 \times \dfrac{1}{8} = 15$ (quả trứng).\\
Số trứng còn lại là: $120 - 15 = 105$ (quả trứng).\\
Lần thứ hai bán được: $105 \times \dfrac{2}{7} = 30$ (quả trứng).\\
Cả hai lần bán được: $15 + 30 = 45$ (quả trứng).\\
\textbf{Đáp án của câu hỏi là B}
}
\end{ex}

\textbf{Phần 2. Trắc nghiệm đúng sai}
\setcounterpageref{ex}{0}

\begin{ex} %[5D1H9]
Cho các biểu thức: $A = \dfrac{7}{9} \times \dfrac{6}{5} + \dfrac{6}{5} \times \dfrac{2}{9} + \dfrac{4}{5}$ và $B = 524 \times 63 + 524 \times 37 - 2\,400$. Xét tính đúng sai của các khẳng định sau:
\choiceTF
{\True Biểu thức $A$ có giá trị bằng $2$.}
{Biểu thức $B$ có giá trị bằng $5\,000$.}
{\True Trong biểu thức $A$, áp dụng tính chất phân phối ta có $\dfrac{6}{5} \times \left(\dfrac{7}{9} + \dfrac{2}{9}\right) = \dfrac{6}{5}$.}
{\True Giá trị của biểu thức $B$ gấp $25\,000$ lần giá trị của biểu thức $A$.}
\loigiai{
\begin{itemize}
    \item $A = \dfrac{6}{5} \times \left(\dfrac{7}{9} + \dfrac{2}{9}\right) + \dfrac{4}{5} = \dfrac{6}{5} \times 1 + \dfrac{4}{5} = \dfrac{10}{5} = 2$.
    \item $B = 524 \times 100 - 2\,400 = 52\,400 - 2\,400 = 50\,000 \neq 5\,000$.
    \item Ta có $\dfrac{6}{5} \times 1 = \dfrac{6}{5}$.
    \item Tỉ số $B : A = 50\,000 : 2 = 25\,000$.
\end{itemize}
\begin{itemize}
    \item a) \textbf{Đúng}.
    \item b) \textbf{Sai}: $B = 50\,000$.
    \item c) \textbf{Đúng}.
    \item d) \textbf{Đúng}.
\end{itemize}
}
\end{ex}

\begin{ex} %[5H1V9]
Sân trường hình chữ nhật có chu vi là $84\text{ m}$, chiều dài hơn chiều rộng là $8\text{ m}$. Xét tính đúng sai của các khẳng định sau:
\choiceTF
{\True Nửa chu vi sân trường là $42\text{ m}$.}
{Chiều rộng sân trường là $25\text{ m}$.}
{\True Chiều dài sân trường là $25\text{ m}$.}
{\True Diện tích của sân trường đó là $425\text{ m}^2$.}
\loigiai{
\begin{itemize}
    \item Nửa chu vi là: $84 : 2 = 42\text{ (m)}$.
    \item Chiều dài sân trường là: $(42 + 8) : 2 = 25\text{ (m)}$.
    \item Chiều rộng sân trường là: $25 - 8 = 17\text{ (m)} \neq 25\text{ m}$.
    \item Diện tích sân trường là: $25 \times 17 = 425\text{ (m}^2\text{)}$.
\end{itemize}
\begin{itemize}
    \item a) \textbf{Đúng}.
    \item b) \textbf{Sai}: Chiều rộng là $17\text{ m}$.
    \item c) \textbf{Đúng}.
    \item d) \textbf{Đúng}.
\end{itemize}
}
\end{ex}

\textbf{Phần 3. Trả lời ngắn}
\setcounterpageref{ex}{0}

\begin{ex} %[5D1H9]
Tính giá trị của biểu thức sau bằng cách thuận tiện nhất:
\[
M = \dfrac{7}{9} \times \dfrac{6}{5} + \dfrac{6}{5} \times \dfrac{2}{9} + \dfrac{4}{5}.
\]
\shortans{2}
\loigiai{
\textbf{Đáp số:} $2$\\
$M = \dfrac{6}{5} \times \left(\dfrac{7}{9} + \dfrac{2}{9}\right) + \dfrac{4}{5} = \dfrac{6}{5} \times 1 + \dfrac{4}{5} = \dfrac{10}{5} = 2$.
}
\end{ex}

\begin{ex} %[5D1H9]
Ước lượng kết quả phép tính $4\,180\,300 + 3\,990\,700$ đến hàng triệu. Kết quả ước lượng là bao nhiêu triệu?
\shortans{8}
\loigiai{
\textbf{Đáp số:} $8$\\
$4\,180\,300 \approx 4\,000\,000$; $3\,990\,700 \approx 4\,000\,000 \Rightarrow 4 + 4 = 8$ (triệu).
}
\end{ex}

\begin{ex} %[5D1H9]
Trạng Trình Nguyễn Bỉnh Khiêm sinh năm 1491. Hỏi năm kỉ niệm 600 năm năm sinh của Trạng Trình là năm nào?
\shortans{2091}
\loigiai{
\textbf{Đáp số:} $2091$\\
Năm kỉ niệm là: $1491 + 600 = 2091$.
}
\end{ex}

\begin{ex} %[5H1V9]
Sân trường hình chữ nhật có chu vi $84\text{ m}$, chiều dài hơn chiều rộng $8\text{ m}$. Tính diện tích của sân trường đó theo mét vuông.
\shortans{425}
\loigiai{
\textbf{Đáp số:} $425$\\
Nửa chu vi: $84 : 2 = 42\text{ m}$.\\
Chiều dài: $(42 + 8) : 2 = 25\text{ m}$, Chiều rộng: $42 - 25 = 17\text{ m}$.\\
Diện tích: $25 \times 17 = 425\text{ m}^2$.
}
\end{ex}

\begin{ex} %[5D1V9]
Cô Ba mang $120$ quả trứng gà ra chợ bán. Lần 1 bán được $\dfrac{1}{8}$ số trứng. Lần 2 bán được $\dfrac{2}{7}$ số trứng còn lại sau lần 1. Hỏi sau hai lần bán, cô Ba còn lại bao nhiêu quả trứng?
\shortans{75}
\loigiai{
\textbf{Đáp số:} $75$\\
Lần 1 bán: $120 \times \dfrac{1}{8} = 15$ quả.\\
Còn lại sau lần 1: $120 - 15 = 105$ quả.\\
Lần 2 bán: $105 \times \dfrac{2}{7} = 30$ quả.\\
Số trứng còn lại sau cả hai lần: $105 - 30 = 75$ quả.
}
\end{ex}

\begin{ex} %[5D1V9]
Một đoàn xe gồm $6$ xe ô tô chở $35$ học sinh/xe và $9$ xe chở $40$ học sinh/xe. Hỏi trung bình mỗi xe chở bao nhiêu học sinh?
\shortans{38}
\loigiai{
\textbf{Đáp số:} $38$\\
Tổng học sinh: $6 \times 35 + 9 \times 40 = 210 + 360 = 570$ em.\\
Trung bình mỗi xe: $570 : 15 = 38$ em.
}
\end{ex}

\newpage

% =========================================================================
% CHỦ ĐỀ 2 - SỐ THẬP PHÂN
% =========================================================================
\section{Số thập phân}

\TaoPhieuKiemTra{Chủ đề 2 -- Số thập phân}{10}{Khái niệm số thập phân}{https://loptoanthaylinh.vn/thi-online/toan-5/bai-10}
\textbf{Phần 1. Trắc nghiệm lựa chọn}
\setcounterpageref{ex}{0}

\begin{ex} %[5D2N10]
Số thập phân gồm có hai phần là
\choice
{\True Phần nguyên và phần thập phân}
{Phần nguyên và phần phân số}
{Phần trăm và phần nghìn}
{Phần chục và phần đơn vị}
\loigiai{
Mỗi số thập phân gồm hai phần: phần nguyên và phần thập phân, chúng được phân cách bởi dấu phẩy.\\
\textbf{Đáp án của câu hỏi là A}
}
\end{ex}

\begin{ex} %[5D2N10]
Số thập phân $9{,}17$ được đọc là
\choice
{Chín phẩy một bảy}
{Chín mươi mốt phẩy bảy}
{Chín trăm mười bảy}
{\True Chín phẩy mười bảy}
\loigiai{
Số $9{,}17$ có phần nguyên là $9$ và phần thập phân là $17$, được đọc là: Chín phẩy mười bảy.\\
\textbf{Đáp án của câu hỏi là D}
}
\end{ex}

\begin{ex} %[5D2H10]
Số thập phân gồm $3$ trăm, $2$ chục, $5$ đơn vị, $4$ phần mười, $3$ phần trăm và $1$ phần nghìn được viết là
\choice
{$325{,}43$}
{\True $325{,}431$}
{$32{,}5431$}
{$352{,}431$}
\loigiai{
Phần nguyên là $325$, phần thập phân gồm hàng phần mười là $4$, phần trăm là $3$, phần nghìn là $1$. Số đó được viết là $325{,}431$.\\
\textbf{Đáp án của câu hỏi là B}
}
\end{ex}

\begin{ex} %[5D2H10]
Trong số thập phân $35{,}624$, chữ số $2$ thuộc hàng nào?
\choice
{Hàng chục}
{Hàng phần mười}
{\True Hàng phần trăm}
{Hàng phần nghìn}
\loigiai{
Trong số $35{,}624$, chữ số $6$ ở hàng phần mười, chữ số $2$ ở hàng phần trăm, chữ số $4$ ở hàng phần nghìn.\\
\textbf{Đáp án của câu hỏi là C}
}
\end{ex}

\begin{ex} %[5D2H10]
Giá trị của chữ số $8$ trong số thập phân $2{,}38$ là
\choice
{$8$}
{\True $\dfrac{8}{100}$}
{$\dfrac{8}{10}$}
{$\dfrac{8}{1\,000}$}
\loigiai{
Chữ số $8$ ở hàng phần trăm nên có giá trị là $\dfrac{8}{100}$ (hay $0{,}08$).\\
\textbf{Đáp án của câu hỏi là B}
}
\end{ex}

\begin{ex} %[5D2H10]
Chuyển phân số thập phân $\dfrac{236}{100}$ thành số thập phân ta được
\choice
{$23{,}6$}
{$0{,}236$}
{$236{,}0$}
{\True $2{,}36$}
\loigiai{
Ta có: $\dfrac{236}{100} = 2\dfrac{36}{100} = 2{,}36$.\\
\textbf{Đáp án của câu hỏi là D}
}
\end{ex}

\begin{ex} %[5D2H10]
Phân số thập phân $\dfrac{39}{1\,000}$ được viết dưới dạng số thập phân là
\choice
{\True $0{,}039$}
{$0{,}39$}
{$0{,}0039$}
{$3{,}9$}
\loigiai{
Ta có: $\dfrac{39}{1\,000} = 0{,}039$.\\
\textbf{Đáp án của câu hỏi là A}
}
\end{ex}

\begin{ex} %[5D2H10]
Số thập phân nào dưới đây có cách đọc là: ``Bốn phẩy không năm''?
\choice
{$4{,}5$}
{$4{,}50$}
{$0{,}45$}
{\True $4{,}05$}
\loigiai{
Số thập phân có phần nguyên là $4$, phần thập phân là $05$ được viết là $4{,}05$ và đọc là bốn phẩy không năm.\\
\textbf{Đáp án của câu hỏi là D}
}
\end{ex}

\begin{ex} %[5D2H10]
Viết số đo $118\text{ cm}$ dưới dạng số thập phân với đơn vị mét ta được
\choice
{$11{,}8\text{ m}$}
{$0{,}118\text{ m}$}
{\True $1{,}18\text{ m}$}
{$118{,}0\text{ m}$}
\loigiai{
Ta có: $118\text{ cm} = \dfrac{118}{100}\text{ m} = 1\dfrac{18}{100}\text{ m} = 1{,}18\text{ m}$.\\
\textbf{Đáp án của câu hỏi là C}
}
\end{ex}

\begin{ex} %[5D2H10]
Viết số đo $148\text{ ml}$ dưới dạng số thập phân với đơn vị lít ta được
\choice
{\True $0{,}148\text{ l}$}
{$1{,}48\text{ l}$}
{$14{,}8\text{ l}$}
{$0{,}0148\text{ l}$}
\loigiai{
Ta có: $148\text{ ml} = \dfrac{148}{1\,000}\text{ l} = 0{,}148\text{ l}$.\\
\textbf{Đáp án của câu hỏi là A}
}
\end{ex}

\begin{ex} %[5D2H10]
Số thập phân thích hợp điền vào dấu hỏi chấm trên tia số $0;\; 0{,}1;\; 0{,}2;\; 0{,}3;\; ?;\; 0{,}5$ là
\choice
{$0{,}35$}
{$0{,}04$}
{\True $0{,}4$}
{$0{,}45$}
\loigiai{
Tia số tăng đều $0{,}1$ đơn vị mỗi vạch. Vạch đứng sau $0{,}3$ là $0{,}4$.\\
\textbf{Đáp án của câu hỏi là C}
}
\end{ex}

\begin{ex} %[5D2V10]
Từ bốn tấm thẻ $7$; $0$; $2$ và dấu phẩy ``$,$'', lập được tất cả bao nhiêu số thập phân có phần nguyên gồm một chữ số và phần thập phân gồm hai chữ số?
\choice
{$4$}
{\True $6$}
{$8$}
{$12$}
\loigiai{
Các số thập phân lập được là:
\begin{itemize}
    \item Phần nguyên là $7$: $7{,}02$ và $7{,}20$.
    \item Phần nguyên là $2$: $2{,}07$ và $2{,}70$.
    \item Phần nguyên là $0$: $0{,}27$ và $0{,}72$.
\end{itemize}
Có tất cả $6$ số thập phân thỏa mãn.\\
\textbf{Đáp án của câu hỏi là B}
}
\end{ex}

\textbf{Phần 2. Trắc nghiệm đúng sai}
\setcounterpageref{ex}{0}

\begin{ex} %[5D2H10]
Cho số thập phân $A = 327{,}106$. Xét tính đúng sai của các khẳng định sau:
\choiceTF
{\True Phần nguyên của số $A$ là $327$.}
{Chữ số $6$ trong số $A$ thuộc hàng phần trăm.}
{\True Chữ số $1$ có giá trị là $\dfrac{1}{10}$.}
{\True Số $A$ được đọc là: Ba trăm hai mươi bảy phẩy một trăm linh sáu.}
\loigiai{
\begin{itemize}
    \item Phần nguyên ở bên trái dấu phẩy là $327$.
    \item Chữ số $6$ ở vị trí thứ ba sau dấu phẩy nên thuộc hàng phần nghìn, có giá trị là $\dfrac{6}{1\,000}$ ($\neq$ hàng phần trăm).
    \item Chữ số $1$ ở hàng phần mười nên có giá trị là $\dfrac{1}{10}$.
    \item Số $A$ đọc là: Ba trăm hai mươi bảy phẩy một trăm linh sáu.
\end{itemize}
\begin{itemize}
    \item a) \textbf{Đúng}.
    \item b) \textbf{Sai}: Chữ số $6$ thuộc hàng phần nghìn.
    \item c) \textbf{Đúng}.
    \item d) \textbf{Đúng}.
\end{itemize}
}
\end{ex}

\begin{ex} %[5D2H10]
Xét tính đúng sai của các khẳng định sau về việc chuyển đổi phân số thập phân thành số thập phân:
\choiceTF
{\True Phân số $\dfrac{149}{10}$ được viết dưới dạng số thập phân là $14{,}9$.}
{Phân số $\dfrac{35}{100}$ được viết dưới dạng số thập phân là $3{,}5$.}
{\True Phân số $\dfrac{1\,675}{1\,000}$ viết dưới dạng số thập phân là $1{,}675$.}
{Phân số $\dfrac{7}{20}$ được chuyển thành số thập phân là $0{,}7$.}
\loigiai{
\begin{itemize}
    \item $\dfrac{149}{10} = 14\dfrac{9}{10} = 14{,}9$.
    \item $\dfrac{35}{100} = 0{,}35 \neq 3{,}5$.
    \item $\dfrac{1\,675}{1\,000} = 1\dfrac{675}{1\,000} = 1{,}675$.
    \item $\dfrac{7}{20} = \dfrac{7 \times 5}{20 \times 5} = \dfrac{35}{100} = 0{,}35 \neq 0{,}7$.
\end{itemize}
\begin{itemize}
    \item a) \textbf{Đúng}.
    \item b) \textbf{Sai}: Kết quả là $0{,}35$.
    \item c) \textbf{Đúng}.
    \item d) \textbf{Sai}: Kết quả là $0{,}35$.
\end{itemize}
}
\end{ex}

\textbf{Phần 3. Trả lời ngắn}
\setcounterpageref{ex}{0}

\begin{ex} %[5D2H10]
Viết phân số thập phân $\dfrac{1\,675}{1\,000}$ dưới dạng số thập phân $a{,}bcd$. Tính tổng các chữ số $a + b + c + d$.
\shortans{19}
\loigiai{
\textbf{Đáp số:} $19$\\
$\dfrac{1\,675}{1\,000} = 1{,}675 \Rightarrow a = 1, b = 6, c = 7, d = 5$. Tổng là: $1 + 6 + 7 + 5 = 19$.
}
\end{ex}

\begin{ex} %[5D2H10]
Chuyển phân số $\dfrac{7}{20}$ thành phân số thập phân có mẫu số là $100$. Hỏi tử số của phân số thập phân đó là bao nhiêu?
\shortans{35}
\loigiai{
\textbf{Đáp số:} $35$\\
$\dfrac{7}{20} = \dfrac{7 \times 5}{20 \times 5} = \dfrac{35}{100}$. Tử số là $35$.
}
\end{ex}

\begin{ex} %[5D2H10]
Chuyển phân số $\dfrac{4}{125}$ thành phân số thập phân có mẫu số là $1\,000$. Tử số của phân số thập phân nhận được bằng bao nhiêu?
\shortans{32}
\loigiai{
\textbf{Đáp số:} $32$\\
$\dfrac{4}{125} = \dfrac{4 \times 8}{125 \times 8} = \dfrac{32}{1\,000}$. Tử số là $32$.
}
\end{ex}

\begin{ex} %[5D2H10]
Cái xà đơn cao $165\text{ cm}$. Chiều cao của xà đơn viết dưới dạng số thập phân theo đơn vị mét là $a{,}bc\text{ m}$. Tính giá trị của $a + b + c$.
\shortans{12}
\loigiai{
\textbf{Đáp số:} $12$\\
$165\text{ cm} = 1{,}65\text{ m} \Rightarrow a = 1, b = 6, c = 5 \Rightarrow a + b + c = 12$.
}
\end{ex}

\begin{ex} %[5D2V10]
Từ bốn tấm thẻ $7$; $0$; $2$ và dấu phẩy ``$,$'', lập số thập phân bé nhất có phần nguyên gồm một chữ số và phần thập phân gồm hai chữ số. Hỏi số thập phân bé nhất đó là số nào? (Nhập dưới dạng số thập phân có dấu phẩy).
\shortans{0,27}
\loigiai{
\textbf{Đáp số:} $0{,}27$\\
Các số lập được là $7{,}02; 7{,}20; 2{,}07; 2{,}70; 0{,}27; 0{,}72$. Số bé nhất là $0{,}27$.
}
\end{ex}

\begin{ex} %[5D2V10]
Từ bốn tấm thẻ $7$; $0$; $2$ và dấu phẩy ``$,$'', lập được tất cả bao nhiêu số thập phân có phần nguyên gồm một chữ số và phần thập phân gồm hai chữ số?
\shortans{6}
\loigiai{
\textbf{Đáp số:} $6$\\
Có $3$ cách chọn chữ số phần nguyên ($7, 2, 0$), với mỗi cách có $2$ cách xếp $2$ chữ số còn lại vào phần thập phân $\Rightarrow 3 \times 2 = 6$ số.
}
\end{ex}

\newpage

\TaoPhieuKiemTra{Chủ đề 2 -- Số thập phân}{11}{So sánh các số thập phân}{https://loptoanthaylinh.vn/thi-online/toan-5/bai-11}
\textbf{Phần 1. Trắc nghiệm lựa chọn}
\setcounterpageref{ex}{0}

\begin{ex} %[5D2N11]
Khi so sánh hai số thập phân $3{,}5$ và $2{,}75$, ta thấy
\choice
{$3{,}5 < 2{,}75$ vì $3{,}5$ có ít chữ số hơn}
{$3{,}5 = 2{,}75$}
{$3{,}5 < 2{,}75$ vì $5 < 75$}
{\True $3{,}5 > 2{,}75$ vì ở phần nguyên có $3 > 2$}
\loigiai{
So sánh phần nguyên của hai số: $3 > 2$ nên $3{,}5 > 2{,}75$.\\
\textbf{Đáp án của câu hỏi là D}
}
\end{ex}

\begin{ex} %[5D2N11]
Khi so sánh hai số thập phân $2{,}75$ và $2{,}29$, ta thấy
\choice
{$2{,}75 < 2{,}29$}
{\True $2{,}75 > 2{,}29$ vì ở hàng phần mười có $7 > 2$}
{$2{,}75 = 2{,}29$}
{$2{,}75 < 2{,}29$ vì ở hàng phần trăm có $5 < 9$}
\loigiai{
Hai số có cùng phần nguyên là $2$. So sánh hàng phần mười: $7 > 2$ nên $2{,}75 > 2{,}29$.\\
\textbf{Đáp án của câu hỏi là B}
}
\end{ex}

\begin{ex} %[5D2H11]
Số thập phân nào dưới đây bằng với số $0{,}7$?
\choice
{$0{,}07$}
{$7{,}0$}
{$0{,}007$}
{\True $0{,}70$}
\loigiai{
Nếu viết thêm chữ số $0$ vào tận cùng bên phải phần thập phân thì được một số thập phân bằng nó, do đó $0{,}7 = 0{,}70$.\\
\textbf{Đáp án của câu hỏi là D}
}
\end{ex}

\begin{ex} %[5D2H11]
Số thập phân $8{,}6100$ khi bỏ các chữ số $0$ ở tận cùng bên phải phần thập phân thì được số gọn hơn là
\choice
{$8{,}6$}
{$86{,}1$}
{\True $8{,}61$}
{$8{,}610$}
\loigiai{
Bỏ hai chữ số $0$ ở tận cùng bên phải phần thập phân ta được $8{,}6100 = 8{,}61$.\\
\textbf{Đáp án của câu hỏi là C}
}
\end{ex}

\begin{ex} %[5D2H11]
Cặp số thập phân nào sau đây là \textbf{bằng nhau}?
\choice
{$0{,}4$ và $0{,}04$}
{\True $21{,}4050$ và $21{,}405$}
{$5{,}82$ và $5{,}802$}
{$13{,}05$ và $13{,}50$}
\loigiai{
Vì bỏ chữ số $0$ ở tận cùng bên phải phần thập phân của $21{,}4050$ ta được $21{,}405$.\\
\textbf{Đáp án của câu hỏi là B}
}
\end{ex}

\begin{ex} %[5D2H11]
Trong các số thập phân $3{,}604;\; 2{,}875;\; 2{,}857;\; 3{,}106$, số bé nhất là
\choice
{$3{,}604$}
{$2{,}875$}
{\True $2{,}857$}
{$3{,}106$}
\loigiai{
So sánh phần nguyên: $2 < 3$. So sánh hai số có phần nguyên là $2$: $2{,}857 < 2{,}875$ (vì $5 < 7$ ở hàng phần trăm). Vậy số bé nhất là $2{,}857$.\\
\textbf{Đáp án của câu hỏi là C}
}
\end{ex}

\begin{ex} %[5D2H11]
Sắp xếp các số $3{,}604;\; 2{,}875;\; 2{,}857;\; 3{,}106$ theo thứ tự từ bé đến lớn là
\choice
{\True $2{,}857;\; 2{,}875;\; 3{,}106;\; 3{,}604$}
{$2{,}875;\; 2{,}857;\; 3{,}106;\; 3{,}604$}
{$3{,}604;\; 3{,}106;\; 2{,}875;\; 2{,}857$}
{$2{,}857;\; 2{,}875;\; 3{,}604;\; 3{,}106$}
\loigiai{
Ta có $2{,}857 < 2{,}875 < 3{,}106 < 3{,}604$.\\
\textbf{Đáp án của câu hỏi là A}
}
\end{ex}

\begin{ex} %[5D2H11]
Số thập phân nào sau đây lớn hơn $1{,}036$ và bé hơn $2$?
\choice
{$0{,}95$}
{$1{,}036$}
{$2{,}01$}
{\True $1{,}36$}
\loigiai{
So sánh: $0{,}95 < 1{,}036 < 1{,}36 < 2 < 2{,}01$. Vậy số thỏa mãn là $1{,}36$.\\
\textbf{Đáp án của câu hỏi là D}
}
\end{ex}

\begin{ex} %[5D2H11]
Điền dấu thích hợp vào chỗ chấm: $135{,}74 \dots 135{,}73$
\choice
{\True $>$}
{$<$}
{$=$}
{Không so sánh được}
\loigiai{
Hai số có cùng phần nguyên $135$ và chữ số hàng phần mười là $7$. So sánh hàng phần trăm: $4 > 3 \Rightarrow 135{,}74 > 135{,}73$.\\
\textbf{Đáp án của câu hỏi là A}
}
\end{ex}

\begin{ex} %[5D2H11]
Chữ số thích hợp điền vào dấu hỏi chấm để $13{,}7 = 13{,}7?$ là
\choice
{$1$}
{$7$}
{\True $0$}
{$3$}
\loigiai{
Để $13{,}7 = 13{,}7?$ thì chữ số tận cùng bên phải phải là $0$ ($13{,}7 = 13{,}70$).\\
\textbf{Đáp án của câu hỏi là C}
}
\end{ex}

\begin{ex} %[5D2V11]
Ba bạn Mị, Núi, Páo đi khám sức khoẻ. Cân nặng của ba bạn lần lượt là một trong các số: $31{,}9\text{ kg};\; 32{,}5\text{ kg};\; 34{,}7\text{ kg}$. Biết bạn Núi nặng nhất, bạn Páo nhẹ nhất. Cân nặng của bạn Mị là
\choice
{\True $32{,}5\text{ kg}$}
{$31{,}9\text{ kg}$}
{$34{,}7\text{ kg}$}
{$33{,}0\text{ kg}$}
\loigiai{
So sánh cân nặng: $31{,}9\text{ kg} < 32{,}5\text{ kg} < 34{,}7\text{ kg}$.\\
Vì Núi nặng nhất ($34{,}7\text{ kg}$) và Páo nhẹ nhất ($31{,}9\text{ kg}$) nên bạn Mị nặng $32{,}5\text{ kg}$.\\
\textbf{Đáp án của câu hỏi là A}
}
\end{ex}

\begin{ex} %[5D2V11]
Chiều dài của ba cây cầu: Cầu Long Biên ($2{,}29\text{ km}$), Cầu Cần Thơ ($2{,}75\text{ km}$), Cầu An Đông ($3{,}5\text{ km}$). Cây cầu dài nhất là
\choice
{Cầu Long Biên}
{\True Cầu An Đông}
{Cầu Cần Thơ}
{Cả ba cây cầu dài bằng nhau}
\loigiai{
So sánh: $2{,}29 < 2{,}75 < 3{,}5$. Cây cầu dài nhất là Cầu An Đông ($3{,}5\text{ km}$).\\
\textbf{Đáp án của câu hỏi là B}
}
\end{ex}

\textbf{Phần 2. Trắc nghiệm đúng sai}
\setcounterpageref{ex}{0}

\begin{ex} %[5D2H11]
Cho các số thập phân: $x = 8{,}9$; $y = 5{,}82$; $z = 0{,}17$; $t = 31{,}6$. Xét tính đúng sai của các khẳng định sau:
\choiceTF
{\True Viết số $x$ dưới dạng phần thập phân có ba chữ số ta được $8{,}900$.}
{Viết số $y$ dưới dạng phần thập phân có ba chữ số ta được $5{,}082$.}
{\True Số lớn nhất trong bốn số đã cho là $t = 31{,}6$.}
{\True Số bé nhất trong bốn số đã cho là $z = 0{,}17$.}
\loigiai{
\begin{itemize}
    \item Thêm hai chữ số $0$ vào bên phải phần thập phân của $8{,}9$ ta được $8{,}900$.
    \item Thêm một chữ số $0$ vào tận cùng bên phải của $5{,}82$ ta được $5{,}820 \neq 5{,}082$.
    \item So sánh: $0{,}17 < 5{,}82 < 8{,}9 < 31{,}6 \Rightarrow$ Số lớn nhất là $31{,}6$, số bé nhất là $0{,}17$.
\end{itemize}
\begin{itemize}
    \item a) \textbf{Đúng}.
    \item b) \textbf{Sai}: Kết quả là $5{,}820$.
    \item c) \textbf{Đúng}.
    \item d) \textbf{Đúng}.
\end{itemize}
}
\end{ex}

\begin{ex} %[5D2V11]
Ba bạn Mị, Núi, Páo có cân nặng lần lượt là $32{,}5\text{ kg}; 34{,}7\text{ kg}; 31{,}9\text{ kg}$. Xét tính đúng sai của các khẳng định sau:
\choiceTF
{\True Bạn Núi nặng hơn bạn Mị là $2{,}2\text{ kg}$.}
{Tổng cân nặng của bạn Páo và bạn Mị là $64{,}5\text{ kg}$.}
{\True Bạn Páo nhẹ hơn bạn Núi là $2{,}8\text{ kg}$.}
{\True Trung bình mỗi bạn nặng $33{,}033\dots\text{ kg}$ (lớn hơn $33\text{ kg}$).}
\loigiai{
\begin{itemize}
    \item Núi nặng hơn Mị: $34{,}7 - 32{,}5 = 2{,}2\text{ kg}$.
    \item Tổng cân nặng Páo và Mị: $31{,}9 + 32{,}5 = 64{,}4\text{ kg} \neq 64{,}5\text{ kg}$.
    \item Páo nhẹ hơn Núi: $34{,}7 - 31{,}9 = 2{,}8\text{ kg}$.
    \item Tổng cân nặng $3$ bạn: $31{,}9 + 32{,}5 + 34{,}7 = 99{,}1\text{ kg}$. Trung bình mỗi bạn: $99{,}1 : 3 \approx 33{,}03\text{ kg} > 33\text{ kg}$.
\end{itemize}
\begin{itemize}
    \item a) \textbf{Đúng}.
    \item b) \textbf{Sai}: Tổng là $64{,}4\text{ kg}$.
    \item c) \textbf{Đúng}.
    \item d) \textbf{Đúng}.
\end{itemize}
}
\end{ex}

\textbf{Phần 3. Trả lời ngắn}
\setcounterpageref{ex}{0}

\begin{ex} %[5D2H11]
Tìm chữ số $x$ thích hợp để $8{,}6100 = 8{,}6x$.
\shortans{1}
\loigiai{
\textbf{Đáp số:} $1$\\
Bỏ các chữ số $0$ tận cùng bên phải phần thập phân ta có $8{,}6100 = 8{,}61 \Rightarrow x = 1$.
}
\end{ex}

\begin{ex} %[5D2H11]
Có bao nhiêu số thập phân có một chữ số ở phần thập phân lớn hơn $3{,}5$ và nhỏ hơn $4{,}2$?
\shortans{6}
\loigiai{
\textbf{Đáp số:} $6$\\
Các số đó là: $3{,}6;\; 3{,}7;\; 3{,}8;\; 3{,}9;\; 4{,}0;\; 4{,}1$. Có tất cả $6$ số.
}
\end{ex}

\begin{ex} %[5D2H11]
Tìm số tự nhiên $x$ thỏa mãn: $8{,}9 < x < 9{,}8$.
\shortans{9}
\loigiai{
\textbf{Đáp số:} $9$\\
Số tự nhiên duy nhất nằm giữa $8{,}9$ và $9{,}8$ là $9$.
}
\end{ex}

\begin{ex} %[5D2H11]
Tìm số tự nhiên $n$ lớn nhất thỏa mãn: $n < 5{,}82$.
\shortans{5}
\loigiai{
\textbf{Đáp số:} $5$\\
Vì $n$ là số tự nhiên và $n < 5{,}82$ nên giá trị lớn nhất của $n$ là $5$.
}
\end{ex}

\begin{ex} %[5D2V11]
Cầu Cần Thơ dài $2{,}75\text{ km}$, Cầu Long Biên dài $2{,}29\text{ km}$. Hỏi cầu Cần Thơ dài hơn cầu Long Biên bao nhiêu mét?
\shortans{460}
\loigiai{
\textbf{Đáp số:} $460$\\
Đổi $2{,}75\text{ km} = 2\,750\text{ m}$ và $2{,}29\text{ km} = 2\,290\text{ m}$.\\
Hiệu chiều dài là: $2\,750 - 2\,290 = 460\text{ (m)}$.
}
\end{ex}

\begin{ex} %[5D2V11]
Tổng cân nặng của ba bạn Mị, Núi và Páo là bao nhiêu ki-lô-gam biết cân nặng ba bạn lần lượt là $32{,}5\text{ kg}; 34{,}7\text{ kg}; 31{,}9\text{ kg}$? (Nhập kết quả dưới dạng số thập phân).
\shortans{99,1}
\loigiai{
\textbf{Đáp số:} $99{,}1$\\
Tổng cân nặng là: $32{,}5 + 34{,}7 + 31{,}9 = 99{,}1\text{ (kg)}$.
}
\end{ex}

\newpage

\TaoPhieuKiemTra{Chủ đề 2 -- Số thập phân}{12}{Viết số đo đại lượng dưới dạng số thập phân}{https://loptoanthaylinh.vn/thi-online/toan-5/bai-12}
\textbf{Phần 1. Trắc nghiệm lựa chọn}
\setcounterpageref{ex}{0}

\begin{ex} %[5D2N12]
Viết số đo $2\text{ m } 15\text{ cm}$ dưới dạng số thập phân có đơn vị đo là mét ta được
\choice
{$21{,}5\text{ m}$}
{$2{,}015\text{ m}$}
{\True $2{,}15\text{ m}$}
{$0{,}215\text{ m}$}
\loigiai{
Ta có: $2\text{ m } 15\text{ cm} = 2\dfrac{15}{100}\text{ m} = 2{,}15\text{ m}$.\\
\textbf{Đáp án của câu hỏi là C}
}
\end{ex}

\begin{ex} %[5D2N12]
Viết số đo $3\text{ m } 8\text{ cm}$ dưới dạng số thập phân có đơn vị mét ta được
\choice
{$3{,}8\text{ m}$}
{\True $3{,}08\text{ m}$}
{$3{,}80\text{ m}$}
{$0{,}38\text{ m}$}
\loigiai{
Ta có: $3\text{ m } 8\text{ cm} = 3\dfrac{8}{100}\text{ m} = 3{,}08\text{ m}$.\\
\textbf{Đáp án của câu hỏi là B}
}
\end{ex}

\begin{ex} %[5D2H12]
Viết số đo $7\text{ km } 80\text{ m}$ dưới dạng số thập phân có đơn vị ki-lô-mét ta được
\choice
{$7{,}8\text{ km}$}
{$7{,}008\text{ km}$}
{$78{,}0\text{ km}$}
{\True $7{,}08\text{ km}$}
\loigiai{
Ta có: $7\text{ km } 80\text{ m} = 7\dfrac{80}{1\,000}\text{ km} = 7{,}08\text{ km}$.\\
\textbf{Đáp án của câu hỏi là D}
}
\end{ex}

\begin{ex} %[5D2H12]
Viết số đo $8\text{ kg } 75\text{ g}$ dưới dạng số thập phân có đơn vị ki-lô-gam ta được
\choice
{$8{,}75\text{ kg}$}
{\True $8{,}075\text{ kg}$}
{$8{,}750\text{ kg}$}
{$80{,}75\text{ kg}$}
\loigiai{
Ta có: $8\text{ kg } 75\text{ g} = 8\dfrac{75}{1\,000}\text{ kg} = 8{,}075\text{ kg}$.\\
\textbf{Đáp án của câu hỏi là B}
}
\end{ex}

\begin{ex} %[5D2H12]
Viết số đo $2\text{ tấn } 325\text{ kg}$ dưới dạng số thập phân với đơn vị tấn ta được
\choice
{\True $2{,}325\text{ tấn}$}
{$23{,}25\text{ tấn}$}
{$2{,}0325\text{ tấn}$}
{$2{,}3250\text{ tấn}$}
\loigiai{
Ta có: $2\text{ tấn } 325\text{ kg} = 2\dfrac{325}{1\,000}\text{ tấn} = 2{,}325\text{ tấn}$.\\
\textbf{Đáp án của câu hỏi là A}
}
\end{ex}

\begin{ex} %[5D2H12]
Viết số đo diện tích $1\text{ m}^2\ 60\text{ dm}^2$ dưới dạng số thập phân có đơn vị mét vuông ta được
\choice
{$1{,}06\text{ m}^2$}
{\True $1{,}6\text{ m}^2$}
{$16{,}0\text{ m}^2$}
{$1{,}006\text{ m}^2$}
\loigiai{
Ta có: $1\text{ m}^2\ 60\text{ dm}^2 = 1\dfrac{60}{100}\text{ m}^2 = 1{,}60\text{ m}^2 = 1{,}6\text{ m}^2$.\\
\textbf{Đáp án của câu hỏi là B}
}
\end{ex}

\begin{ex} %[5D2H12]
Viết số đo diện tích $3\text{ m}^2\ 6\text{ dm}^2$ dưới dạng số thập phân có đơn vị mét vuông ta được
\choice
{$3{,}6\text{ m}^2$}
{$3{,}006\text{ m}^2$}
{\True $3{,}06\text{ m}^2$}
{$36{,}0\text{ m}^2$}
\loigiai{
Ta có: $3\text{ m}^2\ 6\text{ dm}^2 = 3\dfrac{6}{100}\text{ m}^2 = 3{,}06\text{ m}^2$.\\
\textbf{Đáp án của câu hỏi là C}
}
\end{ex}

\begin{ex} %[5D2H12]
Viết số đo diện tích $120\text{ dm}^2$ dưới dạng số thập phân có đơn vị mét vuông ta được
\choice
{$12{,}0\text{ m}^2$}
{$0{,}12\text{ m}^2$}
{$1{,}02\text{ m}^2$}
{\True $1{,}2\text{ m}^2$}
\loigiai{
Ta có: $120\text{ dm}^2 = \dfrac{120}{100}\text{ m}^2 = 1\dfrac{20}{100}\text{ m}^2 = 1{,}2\text{ m}^2$.\\
\textbf{Đáp án của câu hỏi là D}
}
\end{ex}

\begin{ex} %[5D2H12]
Viết số đo dung tích $5\text{ l } 75\text{ ml}$ dưới dạng số thập phân có đơn vị lít ta được
\choice
{\True $5{,}075\text{ l}$}
{$5{,}75\text{ l}$}
{$5{,}750\text{ l}$}
{$50{,}75\text{ l}$}
\loigiai{
Ta có: $5\text{ l } 75\text{ ml} = 5\dfrac{75}{1\,000}\text{ l} = 5{,}075\text{ l}$.\\
\textbf{Đáp án của câu hỏi là A}
}
\end{ex}

\begin{ex} %[5D2H12]
Một con thỏ cân nặng $6\text{ kg } 75\text{ g}$, một con ngỗng cân nặng $6\,100\text{ g}$, một con mèo cân nặng $6{,}095\text{ kg}$. Con vật nào nặng nhất?
\choice
{Thỏ}
{Mèo}
{Ba con nặng bằng nhau}
{\True Ngỗng}
\loigiai{
Đổi về cùng đơn vị ki-lô-gam:
\begin{itemize}
    \item Thỏ: $6\text{ kg } 75\text{ g} = 6{,}075\text{ kg}$.
    \item Ngỗng: $6\,100\text{ g} = 6{,}1\text{ kg} = 6{,}100\text{ kg}$.
    \item Mèo: $6{,}095\text{ kg}$.
\end{itemize}
Vì $6{,}075 < 6{,}095 < 6{,}100$ nên con ngỗng nặng nhất.\\
\textbf{Đáp án của câu hỏi là D}
}
\end{ex}

\begin{ex} %[5D2V12]
Có ba bức tranh tường: Tranh về bảo vệ môi trường có diện tích $5{,}3\text{ m}^2$; Tranh về an toàn giao thông có diện tích $5\text{ m}^2\ 8\text{ dm}^2$; Tranh về phòng chống Covid có diện tích $5\text{ m}^2\ 9\text{ dm}^2$. Bức tranh có diện tích bé nhất là
\choice
{Tranh về bảo vệ môi trường}
{Tranh về phòng chống Covid}
{\True Tranh về an toàn giao thông}
{Ba bức tranh có diện tích bằng nhau}
\loigiai{
Đổi về cùng đơn vị mét vuông:
\begin{itemize}
    \item Bảo vệ môi trường: $5{,}3\text{ m}^2 = 5{,}30\text{ m}^2$.
    \item An toàn giao thông: $5\text{ m}^2\ 8\text{ dm}^2 = 5{,}08\text{ m}^2$.
    \item Phòng chống Covid: $5\text{ m}^2\ 9\text{ dm}^2 = 5{,}09\text{ m}^2$.
\end{itemize}
Vì $5{,}08 < 5{,}09 < 5{,}30$ nên bức tranh về an toàn giao thông có diện tích bé nhất.\\
\textbf{Đáp án của câu hỏi là C}
}
\end{ex}

\begin{ex} %[5D2V12]
Đoạn đường $AB$ dài $1{,}2\text{ km}$, đoạn đường $AC$ dài $1\text{ km } 75\text{ m}$. Hỏi đoạn đường nào dài hơn và dài hơn bao nhiêu mét?
\choice
{\True $AB$ dài hơn $AC$ là $125\text{ m}$}
{$AC$ dài hơn $AB$ là $125\text{ m}$}
{$AB$ dài hơn $AC$ là $25\text{ m}$}
{$AC$ dài hơn $AB$ là $75\text{ m}$}
\loigiai{
Đổi ra mét: $AB = 1{,}2\text{ km} = 1\,200\text{ m}$; $AC = 1\text{ km } 75\text{ m} = 1\,075\text{ m}$.\\
Đoạn đường $AB$ dài hơn đoạn đường $AC$ là: $1\,200 - 1\,075 = 125\text{ (m)}$.\\
\textbf{Đáp án của câu hỏi là A}
}
\end{ex}

\textbf{Phần 2. Trắc nghiệm đúng sai}
\setcounterpageref{ex}{0}

\begin{ex} %[5D2H12]
Xét tính đúng sai của các phép đổi đơn vị đo sau:
\choiceTF
{\True $4\text{ km } 500\text{ m} = 4{,}5\text{ km}$.}
{$1\text{ tấn } 5\text{ tạ} = 1{,}05\text{ tấn}$.}
{\True $8\text{ m}^2\ 75\text{ dm}^2 = 8{,}75\text{ m}^2$.}
{$2\text{ dm}^2\ 5\text{ cm}^2 = 2{,}5\text{ dm}^2$.}
\loigiai{
\begin{itemize}
    \item $4\text{ km } 500\text{ m} = 4\dfrac{500}{1\,000}\text{ km} = 4{,}5\text{ km}$.
    \item $1\text{ tấn } 5\text{ tạ} = 1\dfrac{5}{10}\text{ tấn} = 1{,}5\text{ tấn} \neq 1{,}05\text{ tấn}$.
    \item $8\text{ m}^2\ 75\text{ dm}^2 = 8\dfrac{75}{100}\text{ m}^2 = 8{,}75\text{ m}^2$.
    \item $2\text{ dm}^2\ 5\text{ cm}^2 = 2\dfrac{5}{100}\text{ dm}^2 = 2{,}05\text{ dm}^2 \neq 2{,}5\text{ dm}^2$.
\end{itemize}
\begin{itemize}
    \item a) \textbf{Đúng}.
    \item b) \textbf{Sai}: $1\text{ tấn } 5\text{ tạ} = 1{,}5\text{ tấn}$.
    \item c) \textbf{Đúng}.
    \item d) \textbf{Sai}: $2\text{ dm}^2\ 5\text{ cm}^2 = 2{,}05\text{ dm}^2$.
\end{itemize}
}
\end{ex}

\begin{ex} %[5D2V12]
Hình $A$ có diện tích là $4\text{ cm}^2\ 15\text{ mm}^2$ và Hình $B$ có diện tích là $3{,}95\text{ cm}^2$. Xét tính đúng sai của các khẳng định sau:
\choiceTF
{\True Đổi diện tích hình $A$ sang xăng-ti-mét vuông ta được $4{,}15\text{ cm}^2$.}
{Hình $B$ có diện tích lớn hơn hình $A$.}
{\True Diện tích hình $A$ lớn hơn diện tích hình $B$ là $0{,}2\text{ cm}^2$.}
{Tổng diện tích của cả hai hình là $8\text{ cm}^2$.}
\loigiai{
\begin{itemize}
    \item Hình $A$: $4\text{ cm}^2\ 15\text{ mm}^2 = 4\dfrac{15}{100}\text{ cm}^2 = 4{,}15\text{ cm}^2$.
    \item So sánh: $4{,}15\text{ cm}^2 > 3{,}95\text{ cm}^2$ nên hình $A$ có diện tích lớn hơn hình $B$.
    \item Hiệu diện tích: $4{,}15 - 3{,}95 = 0{,}2\text{ cm}^2$.
    \item Tổng diện tích: $4{,}15 + 3{,}95 = 8{,}1\text{ cm}^2 \neq 8\text{ cm}^2$.
\end{itemize}
\begin{itemize}
    \item a) \textbf{Đúng}.
    \item b) \textbf{Sai}: Hình $A$ lớn hơn hình $B$.
    \item c) \textbf{Đúng}.
    \item d) \textbf{Sai}: Tổng diện tích là $8{,}1\text{ cm}^2$.
\end{itemize}
}
\end{ex}

\textbf{Phần 3. Trả lời ngắn}
\setcounterpageref{ex}{0}

\begin{ex} %[5D2H12]
Viết số đo $456\text{ m}$ dưới dạng số thập phân có đơn vị ki-lô-mét. Kết quả là bao nhiêu? (Nhập dưới dạng số thập phân).
\shortans{0,456}
\loigiai{
\textbf{Đáp số:} $0{,}456$\\
$456\text{ m} = \dfrac{456}{1\,000}\text{ km} = 0{,}456\text{ km}$.
}
\end{ex}

\begin{ex} %[5D2H12]
Viết số đo $1\,450\text{ kg}$ dưới dạng số thập phân có đơn vị tấn (viết dưới dạng gọn nhất).
\shortans{1,45}
\loigiai{
\textbf{Đáp số:} $1{,}45$\\
$1\,450\text{ kg} = \dfrac{1\,450}{1\,000}\text{ tấn} = 1{,}45\text{ tấn}$.
}
\end{ex}

\begin{ex} %[5D2H12]
Viết số đo $85\text{ cm}^2$ dưới dạng số thập phân có đơn vị đề-xi-mét vuông.
\shortans{0,85}
\loigiai{
\textbf{Đáp số:} $0{,}85$\\
$85\text{ cm}^2 = \dfrac{85}{100}\text{ dm}^2 = 0{,}85\text{ dm}^2$.
}
\end{ex}

\begin{ex} %[5D2H12]
Viết số đo $3\,452\text{ ml}$ dưới dạng số thập phân có đơn vị lít.
\shortans{3,452}
\loigiai{
\textbf{Đáp số:} $3{,}452$\\
$3\,452\text{ ml} = \dfrac{3\,452}{1\,000}\text{ l} = 3{,}452\text{ l}$.
}
\end{ex}

\begin{ex} %[5D2V12]
Một hộp sữa và $1$ quả cân $250\text{ g}$ ở một đĩa cân thăng bằng với $1$ quả cân $1\text{ kg}$ và $1$ quả cân $500\text{ g}$ ở đĩa cân bên kia. Hỏi hộp sữa đó cân nặng bao nhiêu ki-lô-gam? (Nhập dưới dạng số thập phân).
\shortans{1,25}
\loigiai{
\textbf{Đáp số:} $1{,}25$\\
Khối lượng hộp sữa là: $(1\,000 + 500) - 250 = 1\,250\text{ (g)} = 1{,}25\text{ (kg)}$.
}
\end{ex}

\begin{ex} %[5D2V12]
Con ốc sên bò theo đường gấp khúc $ABC$, trong đó đoạn $AB = 1{,}8\text{ m}$ và đoạn $BC = 2\text{ m } 15\text{ cm}$. Hỏi con ốc sên đã bò tất cả bao nhiêu mét? (Nhập dưới dạng số thập phân).
\shortans{3,95}
\loigiai{
\textbf{Đáp số:} $3{,}95$\\
Đổi $2\text{ m } 15\text{ cm} = 2{,}15\text{ m}$.\\
Tổng quãng đường ốc sên bò là: $1{,}8 + 2{,}15 = 3{,}95\text{ (m)}$.
}
\end{ex}

\newpage

\TaoPhieuKiemTra{Chủ đề 2 -- Số thập phân}{13}{Làm tròn số thập phân}{https://loptoanthaylinh.vn/thi-online/toan-5/bai-13}
\textbf{Phần 1. Trắc nghiệm lựa chọn}
\setcounterpageref{ex}{0}

\begin{ex} %[5D2N13]
Khi làm tròn số thập phân đến số tự nhiên gần nhất, ta so sánh chữ số ở hàng nào với số $5$?
\choice
{Hàng đơn vị}
{Hàng phần trăm}
{\True Hàng phần mười}
{Hàng phần nghìn}
\loigiai{
Khi làm tròn số thập phân đến số tự nhiên gần nhất, ta so sánh chữ số ở hàng phần mười với $5$.\\
\textbf{Đáp án của câu hỏi là C}
}
\end{ex}

\begin{ex} %[5D2N13]
Làm tròn số thập phân $31{,}2$ đến số tự nhiên gần nhất ta được
\choice
{$30$}
{\True $31$}
{$32$}
{$31{,}5$}
\loigiai{
Chữ số ở hàng phần mười là $2 < 5$ nên ta làm tròn xuống thành $31$.\\
\textbf{Đáp án của câu hỏi là B}
}
\end{ex}

\begin{ex} %[5D2H13]
Làm tròn số thập phân $31{,}75$ đến số tự nhiên gần nhất ta được
\choice
{\True $32$}
{$31$}
{$31{,}8$}
{$30$}
\loigiai{
Chữ số ở hàng phần mười là $7 > 5$ nên ta làm tròn lên thành $32$.\\
\textbf{Đáp án của câu hỏi là A}
}
\end{ex}

\begin{ex} %[5D2H13]
Làm tròn số thập phân $513{,}59$ đến số tự nhiên gần nhất ta được
\choice
{$513$}
{$510$}
{\True $514$}
{$513{,}6$}
\loigiai{
Chữ số ở hàng phần mười là $5$ nên ta làm tròn lên thành $514$.\\
\textbf{Đáp án của câu hỏi là C}
}
\end{ex}

\begin{ex} %[5D2H13]
Làm tròn số thập phân $0{,}806$ đến số tự nhiên gần nhất ta được kết quả là
\choice
{$0$}
{$0{,}8$}
{$2$}
{\True $1$}
\loigiai{
Chữ số ở hàng phần mười là $8 > 5$ nên ta làm tròn lên thành $1$.\\
\textbf{Đáp án của câu hỏi là D}
}
\end{ex}

\begin{ex} %[5D2H13]
Làm tròn số thập phân $2{,}52$ đến hàng phần mười ta được kết quả là
\choice
{$2{,}6$}
{$3{,}0$}
{\True $2{,}5$}
{$2{,}50$}
\loigiai{
Chữ số hàng phần trăm là $2 < 5$ nên ta làm tròn xuống được $2{,}5$.\\
\textbf{Đáp án của câu hỏi là C}
}
\end{ex}

\begin{ex} %[5D2H13]
Làm tròn số thập phân $3{,}25$ đến hàng phần mười ta được
\choice
{$3{,}2$}
{$3{,}0$}
{$4{,}0$}
{\True $3{,}3$}
\loigiai{
Chữ số hàng phần trăm là $5$ nên ta làm tròn lên được $3{,}3$.\\
\textbf{Đáp án của câu hỏi là D}
}
\end{ex}

\begin{ex} %[5D2H13]
Làm tròn số thập phân $6{,}324$ đến hàng phần trăm ta được
\choice
{\True $6{,}32$}
{$6{,}3$}
{$6{,}33$}
{$6{,}4$}
\loigiai{
Chữ số hàng phần nghìn là $4 < 5$ nên ta làm tròn xuống được $6{,}32$.\\
\textbf{Đáp án của câu hỏi là A}
}
\end{ex}

\begin{ex} %[5D2H13]
Làm tròn số thập phân $6{,}327$ đến hàng phần trăm ta được
\choice
{$6{,}32$}
{$6{,}30$}
{$6{,}4$}
{\True $6{,}33$}
\loigiai{
Chữ số hàng phần nghìn là $7 > 5$ nên ta làm tròn lên được $6{,}33$.\\
\textbf{Đáp án của câu hỏi là D}
}
\end{ex}

\begin{ex} %[5D2H13]
Làm tròn số $\pi = 3{,}141592\dots$ đến hàng phần trăm ta được
\choice
{$3{,}1$}
{\True $3{,}14$}
{$3{,}15$}
{$3{,}2$}
\loigiai{
Chữ số hàng phần nghìn của $\pi$ là $1 < 5$ nên ta làm tròn xuống được $3{,}14$.\\
\textbf{Đáp án của câu hỏi là B}
}
\end{ex}

\begin{ex} %[5D2V13]
Rô-bốt đo đường chéo của màn hình ti-vi $55\text{ inch}$ được kết quả là $139{,}7\text{ cm}$. Nếu làm tròn kết quả đó đến số tự nhiên gần nhất thì đường chéo ti-vi dài khoảng
\choice
{\True $140\text{ cm}$}
{$139\text{ cm}$}
{$139{,}5\text{ cm}$}
{$145\text{ cm}$}
\loigiai{
Số $139{,}7$ có chữ số hàng phần mười là $7 > 5$ nên làm tròn lên thành $140\text{ cm}$.\\
\textbf{Đáp án của câu hỏi là A}
}
\end{ex}

\begin{ex} %[5D2V13]
Một quả dưa hấu cân nặng chính xác là $2{,}52\text{ kg}$. Cô bán hàng đã làm tròn số cân nặng đến hàng phần mười. Hỏi quả dưa hấu cân nặng khoảng bao nhiêu ki-lô-gam?
\choice
{$2\text{ kg}$}
{\True $2{,}5\text{ kg}$}
{$2{,}6\text{ kg}$}
{$3\text{ kg}$}
\loigiai{
Số $2{,}52$ có chữ số hàng phần trăm là $2 < 5$ nên làm tròn đến hàng phần mười ta được $2{,}5\text{ kg}$.\\
\textbf{Đáp án của câu hỏi là B}
}
\end{ex}

\textbf{Phần 2. Trắc nghiệm đúng sai}
\setcounterpageref{ex}{0}

\begin{ex} %[5D2H13]
Cho ba số thập phân: $A = 9{,}345$; $B = 21{,}663$; $C = 0{,}4571$. Xét tính đúng sai của các khẳng định sau:
\choiceTF
{\True Làm tròn số $A$ đến hàng phần mười ta được $9{,}3$.}
{Làm tròn số $A$ đến hàng phần trăm ta được $9{,}34$.}
{\True Làm tròn số $B$ đến hàng phần trăm ta được $21{,}66$.}
{\True Làm tròn số $C$ đến hàng phần mười ta được $0{,}5$.}
\loigiai{
\begin{itemize}
    \item Số $A = 9{,}345$ có hàng phần trăm là $4 < 5 \Rightarrow$ làm tròn đến hàng phần mười là $9{,}3$.
    \item Số $A = 9{,}345$ có hàng phần nghìn là $5 \Rightarrow$ làm tròn đến hàng phần trăm là $9{,}35 \neq 9{,}34$.
    \item Số $B = 21{,}663$ có hàng phần nghìn là $3 < 5 \Rightarrow$ làm tròn đến hàng phần trăm là $21{,}66$.
    \item Số $C = 0{,}4571$ có hàng phần trăm là $5 \Rightarrow$ làm tròn đến hàng phần mười là $0{,}5$.
\end{itemize}
\begin{itemize}
    \item a) \textbf{Đúng}.
    \item b) \textbf{Sai}: Kết quả là $9{,}35$.
    \item c) \textbf{Đúng}.
    \item d) \textbf{Đúng}.
\end{itemize}
}
\end{ex}

\begin{ex} %[5D2V13]
Theo Tổ chức Y tế Thế giới (WHO), chiều cao và cân nặng chuẩn của bé trai $10$ tuổi lần lượt là $137{,}8\text{ cm}$ và $31{,}2\text{ kg}$; của bé gái $10$ tuổi là $138{,}6\text{ cm}$ và $31{,}9\text{ kg}$. Xét tính đúng sai của các khẳng định sau khi làm tròn đến số tự nhiên gần nhất:
\choiceTF
{\True Chiều cao chuẩn của bé trai làm tròn thành $138\text{ cm}$.}
{Cân nặng chuẩn của bé trai làm tròn thành $32\text{ kg}$.}
{\True Chiều cao chuẩn của bé gái làm tròn thành $139\text{ cm}$.}
{\True Cân nặng chuẩn của bé gái làm tròn thành $32\text{ kg}$.}
\loigiai{
\begin{itemize}
    \item $137{,}8$ có hàng phần mười là $8 > 5 \Rightarrow$ làm tròn thành $138\text{ cm}$.
    \item $31{,}2$ có hàng phần mười là $2 < 5 \Rightarrow$ làm tròn thành $31\text{ kg} \neq 32\text{ kg}$.
    \item $138{,}6$ có hàng phần mười là $6 > 5 \Rightarrow$ làm tròn thành $139\text{ cm}$.
    \item $31{,}9$ có hàng phần mười là $9 > 5 \Rightarrow$ làm tròn thành $32\text{ kg}$.
\end{itemize}
\begin{itemize}
    \item a) \textbf{Đúng}.
    \item b) \textbf{Sai}: Làm tròn thành $31\text{ kg}$.
    \item c) \textbf{Đúng}.
    \item d) \textbf{Đúng}.
\end{itemize}
}
\end{ex}

\textbf{Phần 3. Trả lời ngắn}
\setcounterpageref{ex}{0}

\begin{ex} %[5D2H13]
Làm tròn số thập phân $42{,}305$ đến số tự nhiên gần nhất.
\shortans{42}
\loigiai{
\textbf{Đáp số:} $42$\\
Chữ số hàng phần mười là $3 < 5$ nên làm tròn thành $42$.
}
\end{ex}

\begin{ex} %[5D2H13]
Làm tròn số thập phân $9{,}345$ đến hàng phần trăm. (Nhập kết quả dưới dạng số thập phân).
\shortans{9,35}
\loigiai{
\textbf{Đáp số:} $9{,}35$\\
Chữ số hàng phần nghìn là $5$ nên làm tròn lên thành $9{,}35$.
}
\end{ex}

\begin{ex} %[5D2H13]
Làm tròn số thập phân $21{,}663$ đến hàng phần mười. (Nhập kết quả dưới dạng số thập phân).
\shortans{21,7}
\loigiai{
\textbf{Đáp số:} $21{,}7$\\
Chữ số hàng phần trăm là $6 > 5$ nên làm tròn lên thành $21{,}7$.
}
\end{ex}

\begin{ex} %[5D2H13]
Làm tròn số $\pi = 3{,}141592\dots$ đến hàng phần mười. (Nhập kết quả dưới dạng số thập phân).
\shortans{3,1}
\loigiai{
\textbf{Đáp số:} $3{,}1$\\
Chữ số hàng phần trăm là $4 < 5$ nên làm tròn xuống thành $3{,}1$.
}
\end{ex}

\begin{ex} %[5D2V13]
Có bao nhiêu số thập phân có một chữ số ở phần thập phân lớn hơn $6$ mà khi làm tròn đến số tự nhiên gần nhất ta được kết quả bằng $7$?
\shortans{4}
\loigiai{
\textbf{Đáp số:} $4$\\
Các số đó là: $6{,}5;\; 6{,}6;\; 6{,}7;\; 6{,}8;\; 6{,}9$. Vì lớn hơn $6$ và có $1$ chữ số thập phân nên có $5$ số: $6{,}5; 6{,}6; 6{,}7; 6{,}8; 6{,}9$. Khi làm tròn đến số tự nhiên gần nhất, cả $5$ số này đều bằng $7$.
}
\end{ex}

\begin{ex} %[5D2V13]
Một mảnh vườn hình chữ nhật có chiều dài $12{,}4\text{ m}$ và chiều rộng $5{,}8\text{ m}$. Tính chu vi của mảnh vườn đó sau khi đã làm tròn mỗi kích thước đến số tự nhiên gần nhất theo đơn vị mét.
\shortans{36}
\loigiai{
\textbf{Đáp số:} $36$\\
Làm tròn chiều dài $12{,}4\text{ m} \rightarrow 12\text{ m}$.\\
Làm tròn chiều rộng $5{,}8\text{ m} \rightarrow 6\text{ m}$.\\
Chu vi mảnh vườn là: $(12 + 6) \times 2 = 36\text{ (m)}$.
}
\end{ex}

\newpage

\TaoPhieuKiemTra{Chủ đề 2 -- Số thập phân}{14}{Luyện tập chung số thập phân}{https://loptoanthaylinh.vn/thi-online/toan-5/bai-14}
\textbf{Phần 1. Trắc nghiệm lựa chọn}
\setcounterpageref{ex}{0}

\begin{ex} %[5D2N14]
Số thập phân gồm ``Ba chục, bảy đơn vị, ba phần mười, sáu phần trăm, bốn phần nghìn'' được viết là
\choice
{$307{,}364$}
{$37{,}346$}
{$373{,}64$}
{\True $37{,}364$}
\loigiai{
Phần nguyên gồm $3$ chục và $7$ đơn vị là $37$. Phần thập phân gồm $3$ phần mười, $6$ phần trăm, $4$ phần nghìn là $364$. Số đó là $37{,}364$.\\
\textbf{Đáp án của câu hỏi là D}
}
\end{ex}

\begin{ex} %[5D2N14]
Số thập phân gồm ``Không đơn vị, hai mươi lăm phần nghìn'' được viết là
\choice
{$0{,}25$}
{$0{,}250$}
{\True $0{,}025$}
{$25{,}000$}
\loigiai{
Số có phần nguyên là $0$, hàng phần mười là $0$, hàng phần trăm là $2$, hàng phần nghìn là $5$ được viết là $0{,}025$.\\
\textbf{Đáp án của câu hỏi là C}
}
\end{ex}

\begin{ex} %[5D2H14]
Số thập phân gồm ``Một nghìn đơn vị, bảy mươi mốt phần nghìn'' được viết là
\choice
{$1\,000{,}71$}
{\True $1\,000{,}071$}
{$1\,071{,}0$}
{$1\,000{,}710$}
\loigiai{
Phần nguyên là $1\,000$, phần thập phân có hàng phần mười là $0$, hàng phần trăm là $7$, hàng phần nghìn là $1$ nên viết là $1\,000{,}071$.\\
\textbf{Đáp án của câu hỏi là B}
}
\end{ex}

\begin{ex} %[5D2H14]
Viết số đo $4\text{ cm } 6\text{ mm}$ dưới dạng số thập phân có đơn vị xăng-ti-mét ta được
\choice
{\True $4{,}6\text{ cm}$}
{$4{,}06\text{ cm}$}
{$46{,}0\text{ cm}$}
{$0{,}46\text{ cm}$}
\loigiai{
Ta có: $4\text{ cm } 6\text{ mm} = 4\dfrac{6}{10}\text{ cm} = 4{,}6\text{ cm}$.\\
\textbf{Đáp án của câu hỏi là A}
}
\end{ex}

\begin{ex} %[5D2H14]
Đổi $4{,}75\text{ kg}$ ra gam ta được kết quả là
\choice
{$475\text{ g}$}
{$4\,075\text{ g}$}
{$47\,500\text{ g}$}
{\True $4\,750\text{ g}$}
\loigiai{
Ta có: $4{,}75\text{ kg} = 4\text{ kg} + 0{,}75\text{ kg} = 4\,000\text{ g} + 750\text{ g} = 4\,750\text{ g}$.\\
\textbf{Đáp án của câu hỏi là D}
}
\end{ex}

\begin{ex} %[5D2H14]
Làm tròn số thập phân $13{,}56$ đến số tự nhiên gần nhất ta được
\choice
{\True $14$}
{$13$}
{$13{,}5$}
{$13{,}6$}
\loigiai{
Chữ số hàng phần mười là $5$ nên làm tròn lên thành $14$.\\
\textbf{Đáp án của câu hỏi là A}
}
\end{ex}

\begin{ex} %[5D2H14]
Làm tròn số thập phân $1{,}36$ đến hàng phần mười ta được
\choice
{$1{,}3$}
{$1{,}0$}
{\True $1{,}4$}
{$2{,}0$}
\loigiai{
Chữ số hàng phần trăm là $6 > 5$ nên làm tròn lên thành $1{,}4$.\\
\textbf{Đáp án của câu hỏi là C}
}
\end{ex}

\begin{ex} %[5D2H14]
Làm tròn số thập phân $1{,}026$ đến hàng phần trăm ta được kết quả là
\choice
{$1{,}02$}
{\True $1{,}03$}
{$1{,}2$}
{$1{,}00$}
\loigiai{
Chữ số hàng phần nghìn là $6 > 5$ nên làm tròn lên thành $1{,}03$.\\
\textbf{Đáp án của câu hỏi là B}
}
\end{ex}

\begin{ex} %[5D2H14]
Sắp xếp các số $3{,}527;\; 3{,}571;\; 2{,}752;\; 3{,}275$ theo thứ tự từ bé đến lớn là
\choice
{$3{,}571;\; 3{,}527;\; 3{,}275;\; 2{,}752$}
{$2{,}752;\; 3{,}527;\; 3{,}275;\; 3{,}571$}
{\True $2{,}752;\; 3{,}275;\; 3{,}527;\; 3{,}571$}
{$3{,}275;\; 2{,}752;\; 3{,}527;\; 3{,}571$}
\loigiai{
Ta có: $2{,}752 < 3{,}275 < 3{,}527 < 3{,}571$.\\
\textbf{Đáp án của câu hỏi là C}
}
\end{ex}

\begin{ex} %[5D2H14]
Ba ô tô chở hàng: xe A chở $2{,}546\text{ tấn}$, xe B chở $2{,}645\text{ tấn}$, xe C chở $2{,}564\text{ tấn}$. Xe chở lượng hàng nhẹ nhất là
\choice
{Xe B}
{\True Xe A}
{Xe C}
{Cả ba xe chở bằng nhau}
\loigiai{
So sánh: $2{,}546 < 2{,}564 < 2{,}645$. Do đó xe A chở lượng hàng nhẹ nhất.\\
\textbf{Đáp án của câu hỏi là B}
}
\end{ex}

\begin{ex} %[5D2V14]
Biết $1\text{ l}$ dầu ăn cân nặng $0{,}9\text{ kg}$, $1\text{ l}$ mật ong cân nặng $1{,}36\text{ kg}$. Hỏi $2\text{ l}$ dầu ăn nhẹ hơn $2\text{ l}$ mật ong bao nhiêu ki-lô-gam?
\choice
{\True $0{,}92\text{ kg}$}
{$0{,}46\text{ kg}$}
{$1{,}02\text{ kg}$}
{$0{,}82\text{ kg}$}
\loigiai{
$1\text{ l}$ dầu ăn nhẹ hơn $1\text{ l}$ mật ong là: $1{,}36 - 0{,}9 = 0{,}46\text{ (kg)}$.\\
$2\text{ l}$ dầu ăn nhẹ hơn $2\text{ l}$ mật ong là: $0{,}46 \times 2 = 0{,}92\text{ (kg)}$.\\
\textbf{Đáp án của câu hỏi là A}
}
\end{ex}

\begin{ex} %[5D2V14]
Từ các thẻ số $5$; $1$; $7$; $0$ và dấu phẩy ``$,$'', số thập phân lớn nhất bé hơn $1$ (mỗi thẻ dùng đúng một lần) là
\choice
{$0{,}715$}
{$0{,}571$}
{$1{,}057$}
{\True $0{,}751$}
\loigiai{
Để số bé hơn $1$, chữ số phần nguyên phải là $0$. Để số lớn nhất, các chữ số phần thập phân được sắp xếp giảm dần: $7, 5, 1$. Số đó là $0{,}751$.\\
\textbf{Đáp án của câu hỏi là D}
}
\end{ex}

\textbf{Phần 2. Trắc nghiệm đúng sai}
\setcounterpageref{ex}{0}

\begin{ex} %[5D2H14]
Cho số thập phân $A = 508{,}42$ và số thập phân $B = 37{,}364$. Xét tính đúng sai của các khẳng định sau:
\choiceTF
{\True Số $A$ được đọc là: Năm trăm linh tám phẩy bốn mươi hai.}
{Chữ số $4$ trong số $A$ và chữ số $4$ trong số $B$ có giá trị bằng nhau.}
{\True Làm tròn số $A$ đến số tự nhiên gần nhất ta được $508$.}
{\True Làm tròn số $B$ đến hàng phần mười ta được $37{,}4$.}
\loigiai{
\begin{itemize}
    \item Số $A = 508{,}42$ đọc là: Năm trăm linh tám phẩy bốn mươi hai.
    \item Chữ số $4$ trong số $A$ ở hàng phần mười (giá trị $\dfrac{4}{10}$), chữ số $4$ trong số $B$ ở hàng phần nghìn (giá trị $\dfrac{4}{1\,000}$), hai giá trị này khác nhau.
    \item Số $A$ có hàng phần mười là $4 < 5 \Rightarrow$ làm tròn thành $508$.
    \item Số $B$ có hàng phần trăm là $6 > 5 \Rightarrow$ làm tròn đến hàng phần mười là $37{,}4$.
\end{itemize}
\begin{itemize}
    \item a) \textbf{Đúng}.
    \item b) \textbf{Sai}: Giá trị của chúng khác nhau.
    \item c) \textbf{Đúng}.
    \item d) \textbf{Đúng}.
\end{itemize}
}
\end{ex}

\begin{ex} %[5D2V14]
Khối lượng của $1\text{ l}$ một số chất lỏng được cho như sau: Nước biển ($1{,}026\text{ kg}$), Rượu ($0{,}79\text{ kg}$), Mật ong ($1{,}36\text{ kg}$), Thuỷ ngân ($13{,}56\text{ kg}$). Xét tính đúng sai của các khẳng định sau:
\choiceTF
{\True Làm tròn khối lượng $1\text{ l}$ nước biển đến hàng phần mười ta được $1{,}0\text{ kg}$ (hay $1\text{ kg}$).}
{Làm tròn khối lượng $1\text{ l}$ rượu đến số tự nhiên gần nhất ta được $0\text{ kg}$.}
{\True Làm tròn khối lượng $1\text{ l}$ mật ong đến hàng phần mười ta được $1{,}4\text{ kg}$.}
{\True Trong các chất lỏng trên, thuỷ ngân có khối lượng nặng nhất.}
\loigiai{
\begin{itemize}
    \item Nước biển: $1{,}026$ có hàng phần trăm là $2 < 5 \Rightarrow$ làm tròn đến hàng phần mười là $1{,}0\text{ kg}$.
    \item Rượu: $0{,}79$ có hàng phần mười là $7 > 5 \Rightarrow$ làm tròn đến số tự nhiên gần nhất là $1\text{ kg} \neq 0\text{ kg}$.
    \item Mật ong: $1{,}36$ có hàng phần trăm là $6 > 5 \Rightarrow$ làm tròn đến hàng phần mười là $1{,}4\text{ kg}$.
    \item Thuỷ ngân có khối lượng $13{,}56\text{ kg}$, lớn nhất trong các chất lỏng trên.
\end{itemize}
\begin{itemize}
    \item a) \textbf{Đúng}.
    \item b) \textbf{Sai}: Làm tròn thành $1\text{ kg}$.
    \item c) \textbf{Đúng}.
    \item d) \textbf{Đúng}.
\end{itemize}
}
\end{ex}

\textbf{Phần 3. Trả lời ngắn}
\setcounterpageref{ex}{0}

\begin{ex} %[5D2H14]
Viết số đo $5\text{ kg } 256\text{ g}$ dưới dạng số thập phân có đơn vị ki-lô-gam.
\shortans{5,256}
\loigiai{
\textbf{Đáp số:} $5{,}256$\\
$5\text{ kg } 256\text{ g} = 5\dfrac{256}{1\,000}\text{ kg} = 5{,}256\text{ kg}$.
}
\end{ex}

\begin{ex} %[5D2H14]
Đổi $2{,}35\text{ m}$ ra xăng-ti-mét. Kết quả là bao nhiêu xăng-ti-mét?
\shortans{235}
\loigiai{
\textbf{Đáp số:} $235$\\
$2{,}35\text{ m} = 235\text{ cm}$.
}
\end{ex}

\begin{ex} %[5D2H14]
Làm tròn số thập phân $0{,}07085$ đến hàng phần trăm. (Nhập kết quả dưới dạng số thập phân).
\shortans{0,07}
\loigiai{
\textbf{Đáp số:} $0{,}07$\\
Chữ số ở hàng phần nghìn là $0 < 5$ nên làm tròn thành $0{,}07$.
}
\end{ex}

\begin{ex} %[5D2H14]
Tìm số tự nhiên $x$ lớn nhất thỏa mãn: $x < 37{,}364$.
\shortans{37}
\loigiai{
\textbf{Đáp số:} $37$\\
Số tự nhiên lớn nhất nhỏ hơn $37{,}364$ là $37$.
}
\end{ex}

\begin{ex} %[5D2V14]
Từ các thẻ số $5$; $1$; $7$; $0$ và dấu phẩy ``$,$'', lập số thập phân bé nhất lớn hơn $0$ và bé hơn $1$ (mỗi thẻ dùng đúng một lần). Hỏi số đó là số nào? (Nhập dưới dạng số thập phân có dấu phẩy).
\shortans{0,157}
\loigiai{
\textbf{Đáp số:} $0{,}157$\\
Phần nguyên là $0$. Để số bé nhất, các chữ số phần thập phân sắp xếp tăng dần: $1, 5, 7$. Số đó là $0{,}157$.
}
\end{ex}

\begin{ex} %[5D2V14]
Một xe ô tô chở $2{,}546\text{ tấn}$ hàng, một xe khác chở $2{,}645\text{ tấn}$ hàng. Hỏi cả hai xe chở tất cả bao nhiêu ki-lô-gam hàng?
\shortans{5191}
\loigiai{
\textbf{Đáp số:} $5191$\\
Đổi $2{,}546\text{ tấn} = 2\,546\text{ kg}$; $2{,}645\text{ tấn} = 2\,645\text{ kg}$.\\
Tổng khối lượng hàng là: $2\,546 + 2\,645 = 5\,191\text{ (kg)}$.
}
\end{ex}

\newpage

% =========================================================================
% CHỦ ĐỀ 3 - MỘT SỐ ĐƠN VỊ ĐO DIỆN TÍCH
% =========================================================================
\section{Một số đơn vị đo diện tích}

\TaoPhieuKiemTra{Chủ đề 3 -- Một số đơn vị đo diện tích}{15}{Ki-lô-mét vuông. Héc-ta}{https://loptoanthaylinh.vn/thi-online/toan-5/bai-15}
\textbf{Phần 1. Trắc nghiệm lựa chọn}
\setcounterpageref{ex}{0}

\begin{ex} %[5H3N15]
Số thích hợp điền vào chỗ chấm để $1\text{ km}^2 = \dots\text{ m}^2$ là
\choice
{$1\,000$}
{\True $1\,000\,000$}
{$10\,000$}
{$100\,000$}
\loigiai{
Ki-lô-mét vuông là diện tích hình vuông cạnh $1\text{ km} = 1\,000\text{ m}$.\\
Do đó $1\text{ km}^2 = 1\,000 \times 1\,000 = 1\,000\,000\text{ m}^2$.\\
\textbf{Đáp án của câu hỏi là B}
}
\end{ex}

\begin{ex} %[5H3N15]
Số thích hợp điền vào chỗ chấm để $1\text{ ha} = \dots\text{ m}^2$ là
\choice
{\True $10\,000$}
{$100$}
{$1\,000$}
{$100\,000$}
\loigiai{
Héc-ta là diện tích hình vuông có cạnh dài $100\text{ m}$.\\
Do đó $1\text{ ha} = 100 \times 100 = 10\,000\text{ m}^2$.\\
\textbf{Đáp án của câu hỏi là A}
}
\end{ex}

\begin{ex} %[5H3N15]
Số thích hợp điền vào chỗ chấm để $1\text{ km}^2 = \dots\text{ ha}$ là
\choice
{$10$}
{$1\,000$}
{\True $100$}
{$10\,000$}
\loigiai{
Ta có $1\text{ km}^2 = 1\,000\,000\text{ m}^2$ và $1\text{ ha} = 10\,000\text{ m}^2$ nên $1\text{ km}^2 = 100\text{ ha}$.\\
\textbf{Đáp án của câu hỏi là C}
}
\end{ex}

\begin{ex} %[5H3H15]
Số thích hợp điền vào dấu hỏi chấm để $2\text{ km}^2 = ?\text{ m}^2$ là
\choice
{$2\,000$}
{$20\,000$}
{$200\,000$}
{\True $2\,000\,000$}
\loigiai{
Vì $1\text{ km}^2 = 1\,000\,000\text{ m}^2$ nên $2\text{ km}^2 = 2\,000\,000\text{ m}^2$.\\
\textbf{Đáp án của câu hỏi là D}
}
\end{ex}

\begin{ex} %[5H3H15]
Số thích hợp điền vào dấu hỏi chấm để $40\,000\text{ ha} = ?\text{ km}^2$ là
\choice
{$4$}
{$40$}
{\True $400$}
{$4\,000$}
\loigiai{
Vì $100\text{ ha} = 1\text{ km}^2$ nên $40\,000\text{ ha} = 40\,000 : 100 = 400\text{ km}^2$.\\
\textbf{Đáp án của câu hỏi là C}
}
\end{ex}

\begin{ex} %[5H3H15]
Đổi $4\text{ km}^2\ 500\text{ m}^2$ ra mét vuông ta được kết quả là
\choice
{$4\,500\text{ m}^2$}
{\True $4\,000\,500\text{ m}^2$}
{$40\,500\text{ m}^2$}
{$400\,500\text{ m}^2$}
\loigiai{
Ta có: $4\text{ km}^2\ 500\text{ m}^2 = 4\,000\,000\text{ m}^2 + 500\text{ m}^2 = 4\,000\,500\text{ m}^2$.\\
\textbf{Đáp án của câu hỏi là B}
}
\end{ex}

\begin{ex} %[5H3H15]
Đổi $5\,000\,000\text{ m}^2$ ra ki-lô-mét vuông ta được kết quả là
\choice
{\True $5\text{ km}^2$}
{$50\text{ km}^2$}
{$0{,}5\text{ km}^2$}
{$500\text{ km}^2$}
\loigiai{
Ta có: $5\,000\,000\text{ m}^2 = 5\,000\,000 : 1\,000\,000 = 5\text{ km}^2$.\\
\textbf{Đáp án của câu hỏi là A}
}
\end{ex}

\begin{ex} %[5H3H15]
Vườn quốc gia Phong Nha -- Kẻ Bàng có diện tích là $123\,326\text{ ha}$. Số đo diện tích vườn quốc gia đó theo đơn vị ki-lô-mét vuông là
\choice
{$12{,}3326\text{ km}^2$}
{$123{,}326\text{ km}^2$}
{\True $1\,233{,}26\text{ km}^2$}
{$12\,332{,}6\text{ km}^2$}
\loigiai{
Ta có: $123\,326\text{ ha} = 123\,326 : 100 = 1\,233{,}26\text{ km}^2$.\\
\textbf{Đáp án của câu hỏi là C}
}
\end{ex}

\begin{ex} %[5H3H15]
Một khu đất hình chữ nhật có kích thước $6\text{ km} \times 4\text{ km}$. Diện tích khu đất đó theo đơn vị héc-ta là
\choice
{$24\text{ ha}$}
{$240\text{ ha}$}
{$24\,000\text{ ha}$}
{\True $2\,400\text{ ha}$}
\loigiai{
Diện tích khu đất là: $6 \times 4 = 24\text{ (km}^2\text{)} = 2\,400\text{ ha}$.\\
\textbf{Đáp án của câu hỏi là D}
}
\end{ex}

\begin{ex} %[5H3V15]
Người ta xây dựng một nhà máy xử lí rác thải trên một khu đất hình vuông có cạnh $200\text{ m}$. Diện tích khu đất đó là bao nhiêu héc-ta?
\choice
{\True $4\text{ ha}$}
{$2\text{ ha}$}
{$40\text{ ha}$}
{$400\text{ ha}$}
\loigiai{
Diện tích khu đất hình vuông là: $200 \times 200 = 40\,000\text{ (m}^2\text{)}$.\\
Đổi $40\,000\text{ m}^2 = 4\text{ ha}$.\\
\textbf{Đáp án của câu hỏi là A}
}
\end{ex}

\begin{ex} %[5H3V15]
Các kĩ sư dự kiến lắp đặt nhà máy điện mặt trời trên một trong ba mảnh đất: Mảnh A hình chữ nhật $7\text{ km} \times 3\text{ km}$; Mảnh B hình vuông cạnh $5\text{ km}$; Mảnh C hình chữ nhật $6\text{ km} \times 4\text{ km}$. Kĩ sư sẽ chọn mảnh đất có diện tích lớn nhất là
\choice
{Mảnh A}
{\True Mảnh B}
{Mảnh C}
{Cả ba mảnh đất có diện tích bằng nhau}
\loigiai{
Diện tích mảnh A: $7 \times 3 = 21\text{ km}^2$.\\
Diện tích mảnh B: $5 \times 5 = 25\text{ km}^2$.\\
Diện tích mảnh C: $6 \times 4 = 24\text{ km}^2$.\\
Vì $25\text{ km}^2 > 24\text{ km}^2 > 21\text{ km}^2$ nên chọn Mảnh B.\\
\textbf{Đáp án của câu hỏi là B}
}
\end{ex}

\begin{ex} %[5H3V15]
Nối số đo thích hợp với mỗi địa danh: Tỉnh Thừa Thiên Huế, Vườn quốc gia Bạch Mã, Điện Thái Hoà. Số đo diện tích phù hợp với Vườn quốc gia Bạch Mã là
\choice
{$5\,054\text{ km}^2$}
{$1\,360\text{ m}^2$}
{$100\text{ ha}$}
{\True $22\,031\text{ ha}$}
\loigiai{
\begin{itemize}
    \item Điện Thái Hoà là một công trình kiến trúc $\rightarrow 1\,360\text{ m}^2$.
    \item Vườn quốc gia Bạch Mã là khu bảo tồn thiên nhiên lớn $\rightarrow 22\,031\text{ ha}$ (tương đương khoảng $220\text{ km}^2$).
    \item Tỉnh Thừa Thiên Huế là đơn vị hành chính cấp tỉnh $\rightarrow 5\,054\text{ km}^2$.
\end{itemize}
\textbf{Đáp án của câu hỏi là D}
}
\end{ex}

\textbf{Phần 2. Trắc nghiệm đúng sai}
\setcounterpageref{ex}{0}

\begin{ex} %[5H3H15]
Xét tính đúng sai của các phép đổi đơn vị đo diện tích sau:
\choiceTF
{\True $3\text{ km}^2 = 300\text{ ha}$.}
{$64\,800\text{ ha} = 64{,}8\text{ km}^2$.}
{\True $6\,450\,000\text{ m}^2 = 6{,}45\text{ km}^2$.}
{\True $2\text{ ha } 500\text{ m}^2 = 20\,500\text{ m}^2$.}
\loigiai{
\begin{itemize}
    \item $3\text{ km}^2 = 3 \times 100 = 300\text{ ha}$.
    \item $64\,800\text{ ha} = 64\,800 : 100 = 648\text{ km}^2 \neq 64{,}8\text{ km}^2$.
    \item $6\,450\,000\text{ m}^2 = 6\,450\,000 : 1\,000\,000 = 6{,}45\text{ km}^2$.
    \item $2\text{ ha } 500\text{ m}^2 = 20\,000 + 500 = 20\,500\text{ m}^2$.
\end{itemize}
\begin{itemize}
    \item a) \textbf{Đúng}.
    \item b) \textbf{Sai}: $64\,800\text{ ha} = 648\text{ km}^2$.
    \item c) \textbf{Đúng}.
    \item d) \textbf{Đúng}.
\end{itemize}
}
\end{ex}

\begin{ex} %[5H3V15]
Một khu bảo tồn thiên nhiên có diện tích là $7\,500\text{ ha}$. Xét tính đúng sai của các khẳng định sau:
\choiceTF
{\True Diện tích khu bảo tồn đổi ra ki-lô-mét vuông là $75\text{ km}^2$.}
{Diện tích khu bảo tồn đổi ra mét vuông là $750\,000\text{ m}^2$.}
{\True Nếu mở rộng thêm $25\text{ km}^2$ thì diện tích mới là $100\text{ km}^2$ (hay $10\,000\text{ ha}$).}
{Diện tích khu bảo tồn này lớn hơn diện tích một tỉnh rộng $5\,000\text{ km}^2$.}
\loigiai{
\begin{itemize}
    \item $7\,500\text{ ha} = 7\,500 : 100 = 75\text{ km}^2$.
    \item $7\,500\text{ ha} = 7\,500 \times 10\,000 = 75\,000\,000\text{ m}^2 \neq 750\,000\text{ m}^2$.
    \item Diện tích mới: $75 + 25 = 100\text{ km}^2 = 10\,000\text{ ha}$.
    \item $75\text{ km}^2 < 5\,000\text{ km}^2$ nên diện tích khu bảo tồn nhỏ hơn diện tích tỉnh đó.
\end{itemize}
\begin{itemize}
    \item a) \textbf{Đúng}.
    \item b) \textbf{Sai}: Kết quả là $75\,000\,000\text{ m}^2$.
    \item c) \textbf{Đúng}.
    \item d) \textbf{Sai}: $75\text{ km}^2 < 5\,000\text{ km}^2$.
\end{itemize}
}
\end{ex}

\textbf{Phần 3. Trả lời ngắn}
\setcounterpageref{ex}{0}

\begin{ex} %[5H3H15]
Đổi $64\,800\text{ ha}$ ra đơn vị ki-lô-mét vuông.
\shortans{648}
\loigiai{
\textbf{Đáp số:} $648$\\
$64\,800\text{ ha} = 64\,800 : 100 = 648\text{ km}^2$.
}
\end{ex}

\begin{ex} %[5H3H15]
Đổi $3\text{ km}^2$ ra héc-ta.
\shortans{300}
\loigiai{
\textbf{Đáp số:} $300$\\
$3\text{ km}^2 = 3 \times 100 = 300\text{ ha}$.
}
\end{ex}

\begin{ex} %[5H3H15]
Một khu rừng có diện tích $6\,450\,000\text{ m}^2$. Hỏi diện tích khu rừng đó bằng bao nhiêu ki-lô-mét vuông? (Nhập kết quả dưới dạng số thập phân).
\shortans{6,45}
\loigiai{
\textbf{Đáp số:} $6{,}45$\\
$6\,450\,000\text{ m}^2 = 6\,450\,000 : 1\,000\,000 = 6{,}45\text{ km}^2$.
}
\end{ex}

\begin{ex} %[5H3V15]
Khu đất hình vuông có cạnh $200\text{ m}$. Hỏi diện tích khu đất đó là bao nhiêu héc-ta?
\shortans{4}
\loigiai{
\textbf{Đáp số:} $4$\\
Diện tích là $200 \times 200 = 40\,000\text{ m}^2 = 4\text{ ha}$.
}
\end{ex}

\begin{ex} %[5H3V15]
Mảnh đất hình chữ nhật có chiều dài $5\text{ km}$ và chiều rộng $3\text{ km}$. Hỏi diện tích mảnh đất đó bằng bao nhiêu héc-ta?
\shortans{1500}
\loigiai{
\textbf{Đáp số:} $1500$\\
Diện tích mảnh đất là: $5 \times 3 = 15\text{ (km}^2\text{)} = 1\,500\text{ ha}$.
}
\end{ex}

\begin{ex} %[5H3V15]
Vườn quốc gia Tràm Chim có diện tích $7\,500\text{ ha}$, Vườn quốc gia Bạch Mã có diện tích $22\,031\text{ ha}$. Hỏi diện tích Vườn quốc gia Bạch Mã lớn hơn diện tích Vườn quốc gia Tràm Chim bao nhiêu héc-ta?
\shortans{14531}
\loigiai{
\textbf{Đáp số:} $14531$\\
Hiệu diện tích là: $22\,031 - 7\,500 = 14\,531\text{ (ha)}$.
}
\end{ex}

\newpage

\TaoPhieuKiemTra{Chủ đề 3 -- Một số đơn vị đo diện tích}{16}{Các đơn vị đo diện tích}{https://loptoanthaylinh.vn/thi-online/toan-5/bai-16}
\textbf{Phần 1. Trắc nghiệm lựa chọn}
\setcounterpageref{ex}{0}

\begin{ex} %[5H3N16]
Trong bảng đơn vị đo diện tích, mỗi đơn vị đo gấp bao nhiêu lần đơn vị bé hơn tiếp liền?
\choice
{$10$ lần}
{$1\,000$ lần}
{\True $100$ lần}
{$10\,000$ lần}
\loigiai{
Trong bảng đơn vị đo diện tích, mỗi đơn vị lớn gấp $100$ lần đơn vị bé hơn tiếp liền.\\
\textbf{Đáp án của câu hỏi là C}
}
\end{ex}

\begin{ex} %[5H3N16]
Đổi $4\text{ m}^2$ ra đề-xi-mét vuông ta được kết quả là
\choice
{\True $400\text{ dm}^2$}
{$40\text{ dm}^2$}
{$4\,000\text{ dm}^2$}
{$\dfrac{4}{100}\text{ dm}^2$}
\loigiai{
Vì $1\text{ m}^2 = 100\text{ dm}^2$ nên $4\text{ m}^2 = 400\text{ dm}^2$.\\
\textbf{Đáp án của câu hỏi là A}
}
\end{ex}

\begin{ex} %[5H3H16]
Đổi $2\text{ m}^2\ 5\text{ dm}^2$ ra đề-xi-mét vuông ta được kết quả là
\choice
{$25\text{ dm}^2$}
{\True $205\text{ dm}^2$}
{$250\text{ dm}^2$}
{$2\,005\text{ dm}^2$}
\loigiai{
Ta có: $2\text{ m}^2\ 5\text{ dm}^2 = 2 \times 100 + 5 = 205\text{ dm}^2$.\\
\textbf{Đáp án của câu hỏi là B}
}
\end{ex}

\begin{ex} %[5H3H16]
Đổi $40\text{ cm}^2\ 4\text{ mm}^2$ ra mi-li-mét vuông ta được kết quả là
\choice
{$44\text{ mm}^2$}
{$404\text{ mm}^2$}
{\True $4\,004\text{ mm}^2$}
{$40\,004\text{ mm}^2$}
\loigiai{
Ta có: $40\text{ cm}^2\ 4\text{ mm}^2 = 40 \times 100 + 4 = 4\,004\text{ mm}^2$.\\
\textbf{Đáp án của câu hỏi là C}
}
\end{ex}

\begin{ex} %[5H3H16]
Viết số đo $3\text{ m}^2\ 7\text{ dm}^2$ dưới dạng số thập phân có đơn vị mét vuông ta được
\choice
{$3{,}7\text{ m}^2$}
{$3{,}007\text{ m}^2$}
{$37{,}0\text{ m}^2$}
{\True $3{,}07\text{ m}^2$}
\loigiai{
Ta có: $3\text{ m}^2\ 7\text{ dm}^2 = 3\dfrac{7}{100}\text{ m}^2 = 3{,}07\text{ m}^2$.\\
\textbf{Đáp án của câu hỏi là D}
}
\end{ex}

\begin{ex} %[5H3H16]
Viết số đo $12\text{ km}^2\ 50\text{ ha}$ dưới dạng số thập phân có đơn vị ki-lô-mét vuông ta được
\choice
{\True $12{,}5\text{ km}^2$}
{$12{,}05\text{ km}^2$}
{$12{,}005\text{ km}^2$}
{$125{,}0\text{ km}^2$}
\loigiai{
Ta có: $12\text{ km}^2\ 50\text{ ha} = 12\dfrac{50}{100}\text{ km}^2 = 12{,}50\text{ km}^2 = 12{,}5\text{ km}^2$.\\
\textbf{Đáp án của câu hỏi là A}
}
\end{ex}

\begin{ex} %[5H3H16]
Số thích hợp điền vào chỗ chấm để $615\text{ dm}^2 = \dots\text{ m}^2\ \dots\text{ dm}^2$ là
\choice
{$6\text{ m}^2\ 5\text{ dm}^2$}
{\True $6\text{ m}^2\ 15\text{ dm}^2$}
{$61\text{ m}^2\ 5\text{ dm}^2$}
{$60\text{ m}^2\ 15\text{ dm}^2$}
\loigiai{
Ta có: $615\text{ dm}^2 = 600\text{ dm}^2 + 15\text{ dm}^2 = 6\text{ m}^2\ 15\text{ dm}^2$.\\
\textbf{Đáp án của câu hỏi là B}
}
\end{ex}

\begin{ex} %[5H3H16]
Số đo thích hợp nhất chỉ diện tích mặt bàn học sinh là
\choice
{$50\text{ cm}^2$}
{$50\text{ m}^2$}
{$50\text{ ha}$}
{\True $50\text{ dm}^2$}
\loigiai{
Mặt bàn học sinh có diện tích thích hợp nhất là $50\text{ dm}^2$ (tương đương $0{,}5\text{ m}^2$).\\
\textbf{Đáp án của câu hỏi là D}
}
\end{ex}

\begin{ex} %[5H3H16]
Số đo diện tích thích hợp nhất cho mặt sàn một căn hộ chung cư là
\choice
{$225\text{ cm}^2$}
{\True $94\text{ m}^2$}
{$94\text{ dm}^2$}
{$8\,310\text{ km}^2$}
\loigiai{
Diện tích mặt sàn căn hộ thường đo bằng mét vuông, phù hợp nhất là $94\text{ m}^2$.\\
\textbf{Đáp án của câu hỏi là B}
}
\end{ex}

\begin{ex} %[5H3V16]
Ô cửa sổ thông gió ở bếp nhà Mai có dạng hình chữ nhật với chiều dài $1\text{ m}$ và chiều rộng $\dfrac{3}{5}\text{ m}$. Diện tích của ô cửa sổ đó tính theo đề-xi-mét vuông là
\choice
{$\dfrac{3}{5}\text{ dm}^2$}
{$6\text{ dm}^2$}
{$600\text{ dm}^2$}
{\True $60\text{ dm}^2$}
\loigiai{
Đổi: $1\text{ m} = 10\text{ dm}$; $\dfrac{3}{5}\text{ m} = 6\text{ dm}$.\\
Diện tích ô cửa sổ là: $10 \times 6 = 60\text{ (dm}^2\text{)}$.\\
\textbf{Đáp án của câu hỏi là D}
}
\end{ex}

\begin{ex} %[5H3V16]
Bác Tư có $1\text{ ha}$ đất. Bác đã dùng $\dfrac{9}{10}$ diện tích mảnh đất để trồng cây và làm đường, phần còn lại để xây nhà. Diện tích đất dùng để xây nhà là
\choice
{\True $1\,000\text{ m}^2$}
{$100\text{ m}^2$}
{$9\,000\text{ m}^2$}
{$10\,000\text{ m}^2$}
\loigiai{
Đổi $1\text{ ha} = 10\,000\text{ m}^2$.\\
Phần diện tích đất xây nhà chiếm: $1 - \dfrac{9}{10} = \dfrac{1}{10}$ (diện tích khu đất).\\
Diện tích đất để xây nhà là: $10\,000 \times \dfrac{1}{10} = 1\,000\text{ (m}^2\text{)}$.\\
\textbf{Đáp án của câu hỏi là A}
}
\end{ex}

\begin{ex} %[5H3V16]
Một mảnh vườn có diện tích $2\,500\text{ dm}^2$. Người ta dùng $\dfrac{2}{5}$ diện tích vườn để trồng hoa. Diện tích trồng hoa theo đơn vị mét vuông là
\choice
{$100\text{ m}^2$}
{$1\,000\text{ m}^2$}
{\True $10\text{ m}^2$}
{$15\text{ m}^2$}
\loigiai{
Đổi $2\,500\text{ dm}^2 = 25\text{ m}^2$.\\
Diện tích trồng hoa là: $25 \times \dfrac{2}{5} = 10\text{ (m}^2\text{)}$.\\
\textbf{Đáp án của câu hỏi là C}
}
\end{ex}

\textbf{Phần 2. Trắc nghiệm đúng sai}
\setcounterpageref{ex}{0}

\begin{ex} %[5H3H16]
Xét tính đúng sai của các khẳng định sau về đổi đơn vị đo diện tích:
\choiceTF
{\True $12\text{ m}^2 = 1\,200\text{ dm}^2$.}
{$2\,500\text{ dm}^2 = 250\text{ m}^2$.}
{\True $8\text{ cm}^2\ 10\text{ mm}^2 = 8{,}1\text{ cm}^2$.}
{$30\text{ dm}^2 = 3\text{ m}^2$.}
\loigiai{
\begin{itemize}
    \item $12\text{ m}^2 = 12 \times 100 = 1\,200\text{ dm}^2$.
    \item $2\,500\text{ dm}^2 = 2\,500 : 100 = 25\text{ m}^2 \neq 250\text{ m}^2$.
    \item $8\text{ cm}^2\ 10\text{ mm}^2 = 8\dfrac{10}{100}\text{ cm}^2 = 8{,}1\text{ cm}^2$.
    \item $30\text{ dm}^2 = \dfrac{30}{100}\text{ m}^2 = 0{,}3\text{ m}^2 \neq 3\text{ m}^2$.
\end{itemize}
\begin{itemize}
    \item a) \textbf{Đúng}.
    \item b) \textbf{Sai}: $2\,500\text{ dm}^2 = 25\text{ m}^2$.
    \item c) \textbf{Đúng}.
    \item d) \textbf{Sai}: $30\text{ dm}^2 = 0{,}3\text{ m}^2$.
\end{itemize}
}
\end{ex}

\begin{ex} %[5H3V16]
Bác Tư có thửa ruộng diện tích $1\text{ ha}$. Bác dành $\dfrac{9}{10}$ diện tích để trồng lúa và làm mương nước, phần còn lại xây trang trại. Xét tính đúng sai của các khẳng định sau:
\choiceTF
{\True Diện tích cả thửa ruộng là $10\,000\text{ m}^2$.}
{Diện tích đất trồng lúa và làm mương nước là $900\text{ m}^2$.}
{\True Diện tích xây trang trại là $1\,000\text{ m}^2$.}
{\True Diện tích đất trồng lúa và làm mương nước gấp $9$ lần diện tích đất xây trang trại.}
\loigiai{
\begin{itemize}
    \item $1\text{ ha} = 10\,000\text{ m}^2$.
    \item Diện tích trồng lúa và mương nước: $10\,000 \times \dfrac{9}{10} = 9\,000\text{ m}^2 \neq 900\text{ m}^2$.
    \item Diện tích xây trang trại: $10\,000 - 9\,000 = 1\,000\text{ m}^2$.
    \item Tỉ số: $9\,000 : 1\,000 = 9$ lần.
\end{itemize}
\begin{itemize}
    \item a) \textbf{Đúng}.
    \item b) \textbf{Sai}: Diện tích là $9\,000\text{ m}^2$.
    \item c) \textbf{Đúng}.
    \item d) \textbf{Đúng}.
\end{itemize}
}
\end{ex}

\textbf{Phần 3. Trả lời ngắn}
\setcounterpageref{ex}{0}

\begin{ex} %[5H3H16]
Đổi $2\,500\text{ dm}^2$ ra mét vuông.
\shortans{25}
\loigiai{
\textbf{Đáp số:} $25$\\
$2\,500\text{ dm}^2 = 2\,500 : 100 = 25\text{ m}^2$.
}
\end{ex}

\begin{ex} %[5H3H16]
Đổi $40\text{ cm}^2\ 4\text{ mm}^2$ ra mi-li-mét vuông.
\shortans{4004}
\loigiai{
\textbf{Đáp số:} $4004$\\
$40\text{ cm}^2\ 4\text{ mm}^2 = 40 \times 100 + 4 = 4\,004\text{ mm}^2$.
}
\end{ex}

\begin{ex} %[5H3H16]
Viết số đo $30\text{ dm}^2$ dưới dạng số thập phân có đơn vị mét vuông. (Nhập kết quả dưới dạng số thập phân).
\shortans{0,3}
\loigiai{
\textbf{Đáp số:} $0{,}3$\\
$30\text{ dm}^2 = \dfrac{30}{100}\text{ m}^2 = 0{,}3\text{ m}^2$.
}
\end{ex}

\begin{ex} %[5H3H16]
Viết số đo $12\text{ km}^2\ 50\text{ ha}$ dưới dạng số thập phân có đơn vị ki-lô-mét vuông. (Nhập kết quả dưới dạng số thập phân).
\shortans{12,5}
\loigiai{
\textbf{Đáp số:} $12{,}5$\\
$12\text{ km}^2\ 50\text{ ha} = 12{,}5\text{ km}^2$.
}
\end{ex}

\begin{ex} %[5H3V16]
Ô cửa sổ hình chữ nhật có chiều dài $1\text{ m}$, chiều rộng $\dfrac{3}{5}\text{ m}$. Hỏi diện tích ô cửa sổ đó là bao nhiêu đề-xi-mét vuông?
\shortans{60}
\loigiai{
\textbf{Đáp số:} $60$\\
$1\text{ m} = 10\text{ dm}$; $\dfrac{3}{5}\text{ m} = 6\text{ dm} \Rightarrow S = 10 \times 6 = 60\text{ dm}^2$.
}
\end{ex}

\begin{ex} %[5H3V16]
Bác Tư có thửa đất $1\text{ ha}$. Bác dùng $\dfrac{9}{10}$ thửa đất làm vườn, phần còn lại xây nhà. Hỏi diện tích xây nhà của bác Tư là bao nhiêu mét vuông?
\shortans{1000}
\loigiai{
\textbf{Đáp số:} $1000$\\
Diện tích đất xây nhà là: $10\,000 \times \left(1 - \dfrac{9}{10}\right) = 1\,000\text{ (m}^2\text{)}$.
}
\end{ex}

\newpage

\TaoPhieuKiemTra{Chủ đề 3 -- Một số đơn vị đo diện tích}{17}{Thực hành và trải nghiệm đơn vị đo đại lượng}{https://loptoanthaylinh.vn/thi-online/toan-5/bai-17}
\textbf{Phần 1. Trắc nghiệm lựa chọn}
\setcounterpageref{ex}{0}

\begin{ex} %[5H3N17]
Một viên gạch men hình vuông có cạnh dài $50\text{ cm}$. Diện tích của viên gạch men đó là
\choice
{$200\text{ cm}^2$}
{$2\,500\text{ cm}^2$}
{\True $0{,}25\text{ m}^2$}
{$2{,}5\text{ m}^2$}
\loigiai{
Diện tích viên gạch là: $50 \times 50 = 2\,500\text{ (cm}^2\text{)} = 0{,}25\text{ m}^2$.\\
\textbf{Đáp án của câu hỏi là C}
}
\end{ex}

\begin{ex} %[5H3N17]
Mỗi hộp gạch men chứa $4$ viên gạch hình vuông cạnh $50\text{ cm}$. Diện tích gạch trong mỗi hộp lát được là
\choice
{\True $1\text{ m}^2$}
{$0{,}5\text{ m}^2$}
{$2\text{ m}^2$}
{$4\text{ m}^2$}
\loigiai{
Diện tích một viên gạch là $0{,}25\text{ m}^2$. Bốn viên gạch có diện tích là: $4 \times 0{,}25 = 1\text{ (m}^2\text{)}$.\\
\textbf{Đáp án của câu hỏi là A}
}
\end{ex}

\begin{ex} %[5H3H17]
Một tấm gỗ hình chữ nhật có kích thước $200\text{ cm} \times 25\text{ cm}$. Diện tích tấm gỗ đó là
\choice
{$5\,000\text{ m}^2$}
{$50\text{ m}^2$}
{\True $0{,}5\text{ m}^2$}
{$5\text{ m}^2$}
\loigiai{
Đổi $200\text{ cm} = 2\text{ m}$; $25\text{ cm} = 0{,}25\text{ m}$.\\
Diện tích tấm gỗ là: $2 \times 0{,}25 = 0{,}5\text{ (m}^2\text{)}$.\\
\textbf{Đáp án của câu hỏi là C}
}
\end{ex}

\begin{ex} %[5H3H17]
Một tấm gỗ nhỏ có kích thước $40\text{ cm} \times 25\text{ cm}$. Diện tích tấm gỗ đó theo đơn vị mét vuông là
\choice
{\True $0{,}1\text{ m}^2$}
{$1\text{ m}^2$}
{$0{,}01\text{ m}^2$}
{$10\text{ m}^2$}
\loigiai{
Đổi $40\text{ cm} = 0{,}4\text{ m}$; $25\text{ cm} = 0{,}25\text{ m}$.\\
Diện tích tấm gỗ là: $0{,}4 \times 0{,}25 = 0{,}1\text{ (m}^2\text{)}$.\\
\textbf{Đáp án của câu hỏi là A}
}
\end{ex}

\begin{ex} %[5H3H17]
Các bạn lớp 5A cần mua $4$ tấm gỗ lớn (mỗi tấm diện tích $0{,}5\text{ m}^2$) và $10$ tấm gỗ nhỏ (mỗi tấm diện tích $0{,}1\text{ m}^2$) để đóng kệ sách. Tổng diện tích gỗ cần mua là
\choice
{$2\text{ m}^2$}
{$4\text{ m}^2$}
{$5\text{ m}^2$}
{\True $3\text{ m}^2$}
\loigiai{
Tổng diện tích gỗ là: $4 \times 0{,}5 + 10 \times 0{,}1 = 2 + 1 = 3\text{ (m}^2\text{)}$.\\
\textbf{Đáp án của câu hỏi là D}
}
\end{ex}

\begin{ex} %[5H3H17]
Giá tiền $1\text{ m}^2$ gỗ công nghiệp là $250\,000$ đồng. Để mua $3\text{ m}^2$ gỗ đó làm giá sách, các bạn cần chi trả số tiền là
\choice
{$500\,000$ đồng}
{$1\,000\,000$ đồng}
{\True $750\,000$ đồng}
{$1\,250\,000$ đồng}
\loigiai{
Số tiền mua gỗ là: $3 \times 250\,000 = 750\,000$ (đồng).\\
\textbf{Đáp án của câu hỏi là C}
}
\end{ex}

\begin{ex} %[5H3H17]
Một phòng học hình chữ nhật có chiều dài $8\text{ m}$ và chiều rộng $6\text{ m}$. Diện tích sàn phòng học đó là
\choice
{$28\text{ m}^2$}
{\True $48\text{ m}^2$}
{$40\text{ m}^2$}
{$56\text{ m}^2$}
\loigiai{
Diện tích mặt sàn phòng học là: $8 \times 6 = 48\text{ (m}^2\text{)}$.\\
\textbf{Đáp án của câu hỏi là B}
}
\end{ex}

\begin{ex} %[5H3V17]
Để lát kín mặt sàn phòng học diện tích $48\text{ m}^2$ bằng các hộp gạch (mỗi hộp có $4$ viên gạch hình vuông cạnh $50\text{ cm}$, lát được $1\text{ m}^2$), người ta cần mua số hộp gạch là
\choice
{\True $48$ hộp}
{$12$ hộp}
{$24$ hộp}
{$96$ hộp}
\loigiai{
Vì mỗi hộp gạch lát được $1\text{ m}^2$ sàn nên cần: $48 : 1 = 48$ (hộp gạch).\\
\textbf{Đáp án của câu hỏi là A}
}
\end{ex}

\begin{ex} %[5H3V17]
Mỗi hộp gạch men có giá $140\,000$ đồng. Chi phí mua $48$ hộp gạch để lát một phòng học là
\choice
{$5\,600\,000$ đồng}
{$7\,200\,000$ đồng}
{$8\,400\,000$ đồng}
{\True $6\,720\,000$ đồng}
\loigiai{
Chi phí mua gạch là: $48 \times 140\,000 = 6\,720\,000$ (đồng).\\
\textbf{Đáp án của câu hỏi là D}
}
\end{ex}

\begin{ex} %[5H3V17]
Trường Tiểu học Đoàn Kết có $10$ phòng học cùng kích thước $8\text{ m} \times 6\text{ m}$. Tổng chi phí mua gạch để lát lại toàn bộ sàn $10$ phòng học đó (biết chi phí lát mỗi phòng là $6\,720\,000$ đồng) là
\choice
{$60\,720\,000$ đồng}
{\True $67\,200\,000$ đồng}
{$672\,000\,000$ đồng}
{$58\,400\,000$ đồng}
\loigiai{
Tổng chi phí lát $10$ phòng học là: $10 \times 6\,720\,000 = 67\,200\,000$ (đồng).\\
\textbf{Đáp án của câu hỏi là B}
}
\end{ex}

\begin{ex} %[5H3V17]
Một mảnh sân hình chữ nhật dài $15\text{ m}$ và rộng $8\text{ m}$. Người ta lát sân bằng loại gạch hình vuông có diện tích $0{,}25\text{ m}^2$/viên. Số viên gạch cần dùng là
\choice
{$120$ viên}
{$240$ viên}
{$960$ viên}
{\True $480$ viên}
\loigiai{
Diện tích mảnh sân là: $15 \times 8 = 120\text{ (m}^2\text{)}$.\\
Số viên gạch cần dùng là: $120 : 0{,}25 = 480$ (viên).\\
\textbf{Đáp án của câu hỏi là D}
}
\end{ex}

\begin{ex} %[5H3V17]
Một bức tường hình chữ nhật có diện tích $18\text{ m}^2$. Người ta ốp tường bằng các tấm gỗ có kích thước $60\text{ cm} \times 10\text{ cm}$. Số tấm gỗ cần dùng để ốp kín bức tường là
\choice
{$30$ tấm}
{\True $300$ tấm}
{$60$ tấm}
{$150$ tấm}
\loigiai{
Đổi diện tích mỗi tấm gỗ: $0{,}6\text{ m} \times 0{,}1\text{ m} = 0{,}06\text{ m}^2$.\\
Số tấm gỗ cần dùng là: $18 : 0{,}06 = 300$ (tấm).\\
\textbf{Đáp án của câu hỏi là B}
}
\end{ex}

\textbf{Phần 2. Trắc nghiệm đúng sai}
\setcounterpageref{ex}{0}

\begin{ex} %[5H3H17]
Các bạn học sinh lớp 5A dự định đóng kệ sách gồm $4$ tấm gỗ lớn kích thước $200\text{ cm} \times 25\text{ cm}$ và $10$ tấm gỗ nhỏ kích thước $40\text{ cm} \times 25\text{ cm}$. Giá gỗ là $250\,000$ đồng/$1\text{ m}^2$. Xét tính đúng sai của các khẳng định sau:
\choiceTF
{\True Diện tích của $4$ tấm gỗ lớn là $2\text{ m}^2$.}
{Diện tích của $10$ tấm gỗ nhỏ là $0{,}5\text{ m}^2$.}
{\True Tổng diện tích tất cả các tấm gỗ cần dùng là $3\text{ m}^2$.}
{\True Số tiền cần để mua gỗ đóng chiếc kệ sách là $750\,000$ đồng.}
\loigiai{
\begin{itemize}
    \item Diện tích $1$ tấm gỗ lớn: $2 \times 0{,}25 = 0{,}5\text{ m}^2 \Rightarrow 4$ tấm lớn có diện tích $4 \times 0{,}5 = 2\text{ m}^2$.
    \item Diện tích $1$ tấm gỗ nhỏ: $0{,}4 \times 0{,}25 = 0{,}1\text{ m}^2 \Rightarrow 10$ tấm nhỏ có diện tích $10 \times 0{,}1 = 1\text{ m}^2 \neq 0{,}5\text{ m}^2$.
    \item Tổng diện tích gỗ: $2 + 1 = 3\text{ m}^2$.
    \item Chi phí mua gỗ: $3 \times 250\,000 = 750\,000$ đồng.
\end{itemize}
\begin{itemize}
    \item a) \textbf{Đúng}.
    \item b) \textbf{Sai}: Diện tích là $1\text{ m}^2$.
    \item c) \textbf{Đúng}.
    \item d) \textbf{Đúng}.
\end{itemize}
}
\end{ex}

\begin{ex} %[5H3V17]
Một phòng học có chiều dài $8\text{ m}$, chiều rộng $6\text{ m}$ được lát sàn bằng gạch vuông cạnh $50\text{ cm}$ (đóng hộp $4$ viên/hộp với giá $140\,000$ đồng/hộp). Xét tính đúng sai của các khẳng định sau:
\choiceTF
{\True Diện tích sàn phòng học là $48\text{ m}^2$.}
{Diện tích gạch trong mỗi hộp lát được $0{,}5\text{ m}^2$ sàn.}
{\True Cần mua tất cả $48$ hộp gạch để lát kín phòng học.}
{\True Chi phí mua gạch để lát một phòng học là $6\,720\,000$ đồng.}
\loigiai{
\begin{itemize}
    \item Diện tích phòng học: $8 \times 6 = 48\text{ m}^2$.
    \item Mỗi viên gạch có diện tích $0{,}5 \times 0{,}5 = 0{,}25\text{ m}^2 \Rightarrow 1$ hộp có $4 \times 0{,}25 = 1\text{ m}^2 \neq 0{,}5\text{ m}^2$.
    \item Số hộp gạch cần mua: $48 : 1 = 48$ hộp.
    \item Chi phí mua gạch: $48 \times 140\,000 = 6\,720\,000$ đồng.
\end{itemize}
\begin{itemize}
    \item a) \textbf{Đúng}.
    \item b) \textbf{Sai}: Mỗi hộp lát được $1\text{ m}^2$.
    \item c) \textbf{Đúng}.
    \item d) \textbf{Đúng}.
\end{itemize}
}
\end{ex}

\textbf{Phần 3. Trả lời ngắn}
\setcounterpageref{ex}{0}

\begin{ex} %[5H3H17]
Tính diện tích của một tấm gỗ hình chữ nhật có chiều dài $200\text{ cm}$ và chiều rộng $25\text{ cm}$ theo đơn vị mét vuông. (Nhập kết quả dưới dạng số thập phân).
\shortans{0,5}
\loigiai{
\textbf{Đáp số:} $0{,}5$\\
$2\text{ m} \times 0{,}25\text{ m} = 0{,}5\text{ m}^2$.
}
\end{ex}

\begin{ex} %[5H3H17]
Tính tổng diện tích (theo mét vuông) của $10$ tấm gỗ hình chữ nhật cùng kích thước $40\text{ cm} \times 25\text{ cm}$.
\shortans{1}
\loigiai{
\textbf{Đáp số:} $1$\\
Diện tích $10$ tấm gỗ là: $10 \times (0{,}4 \times 0{,}25) = 1\text{ (m}^2\text{)}$.
}
\end{ex}

\begin{ex} %[5H3H17]
Để mua $3\text{ m}^2$ gỗ công nghiệp với đơn giá $250\,000$ đồng/$1\text{ m}^2$ thì phải trả bao nhiêu nghìn đồng?
\shortans{750}
\loigiai{
\textbf{Đáp số:} $750$\\
Số tiền cần trả là: $3 \times 250\,000 = 750\,000$ đồng $= 750$ nghìn đồng.
}
\end{ex}

\begin{ex} %[5H3V17]
Phòng học hình chữ nhật dài $8\text{ m}$, rộng $6\text{ m}$. Cần mua bao nhiêu hộp gạch (mỗi hộp gồm $4$ viên gạch vuông cạnh $50\text{ cm}$) để lát kín nền phòng học?
\shortans{48}
\loigiai{
\textbf{Đáp số:} $48$\\
Diện tích phòng: $8 \times 6 = 48\text{ m}^2$. Mỗi hộp gạch lát được $4 \times 0{,}25 = 1\text{ m}^2 \Rightarrow$ Cần $48$ hộp.
}
\end{ex}

\begin{ex} %[5H3V17]
Mỗi hộp gạch có giá $140\,000$ đồng. Hỏi cần bao nhiêu triệu đồng để mua gạch lát nền cho $5$ phòng học, biết mỗi phòng cần $48$ hộp gạch? (Nhập kết quả dưới dạng số thập phân).
\shortans{33,6}
\loigiai{
\textbf{Đáp số:} $33{,}6$\\
Chi phí lát $5$ phòng: $5 \times 48 \times 140\,000 = 33\,600\,000$ đồng $= 33{,}6$ triệu đồng.
}
\end{ex}

\begin{ex} %[5H3V17]
Một bức tường hình chữ nhật diện tích $18\text{ m}^2$ được ốp kín bằng các tấm gỗ nhỏ diện tích $0{,}06\text{ m}^2$/tấm. Hỏi cần tất cả bao nhiêu tấm gỗ nhỏ như vậy?
\shortans{300}
\loigiai{
\textbf{Đáp số:} $300$\\
Số tấm gỗ cần dùng là: $18 : 0{,}06 = 300$ (tấm).
}
\end{ex}

\newpage

\TaoPhieuKiemTra{Chủ đề 3 -- Một số đơn vị đo diện tích}{18}{Luyện tập chung diện tích}{https://loptoanthaylinh.vn/thi-online/toan-5/bai-18}
\textbf{Phần 1. Trắc nghiệm lựa chọn}
\setcounterpageref{ex}{0}

\begin{ex} %[5H3N18]
Số đo thích hợp nhất chỉ diện tích của một trường tiểu học là
\choice
{$1{,}5\text{ m}^2$}
{$1{,}5\text{ dm}^2$}
{\True $1{,}5\text{ ha}$}
{$1{,}5\text{ cm}^2$}
\loigiai{
Diện tích khuôn viên một trường học thường được đo bằng héc-ta, phù hợp nhất là $1{,}5\text{ ha}$ ($15\,000\text{ m}^2$).\\
\textbf{Đáp án của câu hỏi là C}
}
\end{ex}

\begin{ex} %[5H3N18]
Số đo thích hợp chỉ diện tích của một bức tranh treo tường là
\choice
{$1\text{ km}^2$}
{$1\text{ ha}$}
{$1\text{ mm}^2$}
{\True $1\text{ m}^2$}
\loigiai{
Diện tích của một bức tranh treo tường trong phòng phù hợp nhất là $1\text{ m}^2$.\\
\textbf{Đáp án của câu hỏi là D}
}
\end{ex}

\begin{ex} %[5H3H18]
Đổi $25\,000\text{ m}^2$ ra héc-ta ta được kết quả là
\choice
{$0{,}25\text{ ha}$}
{\True $2{,}5\text{ ha}$}
{$25\text{ ha}$}
{$250\text{ ha}$}
\loigiai{
Ta có: $25\,000\text{ m}^2 = 25\,000 : 10\,000 = 2{,}5\text{ ha}$.\\
\textbf{Đáp án của câu hỏi là B}
}
\end{ex}

\begin{ex} %[5H3H18]
Đổi $7{,}5\text{ km}^2$ ra héc-ta ta được
\choice
{\True $750\text{ ha}$}
{$75\text{ ha}$}
{$7\,500\text{ ha}$}
{$75\,000\text{ ha}$}
\loigiai{
Ta có: $7{,}5\text{ km}^2 = 7{,}5 \times 100 = 750\text{ ha}$.\\
\textbf{Đáp án của câu hỏi là A}
}
\end{ex}

\begin{ex} %[5H3H18]
Đổi $0{,}48\text{ km}^2$ ra héc-ta ta được
\choice
{$4{,}8\text{ ha}$}
{\True $48\text{ ha}$}
{$480\text{ ha}$}
{$4\,800\text{ ha}$}
\loigiai{
Ta có: $0{,}48\text{ km}^2 = 0{,}48 \times 100 = 48\text{ ha}$.\\
\textbf{Đáp án của câu hỏi là B}
}
\end{ex}

\begin{ex} %[5H3H18]
Điền dấu thích hợp vào chỗ chấm: $4\text{ m}^2\ 2\text{ dm}^2 \dots 42\text{ m}^2$
\choice
{$>$}
{$=$}
{\True $<$}
{Không so sánh được}
\loigiai{
Đổi: $4\text{ m}^2\ 2\text{ dm}^2 = 4{,}02\text{ m}^2 < 42\text{ m}^2$.\\
\textbf{Đáp án của câu hỏi là C}
}
\end{ex}

\begin{ex} %[5H3H18]
Điền dấu thích hợp vào chỗ chấm: $2\,800\text{ ha} \dots 28\text{ km}^2$
\choice
{$>$}
{\True $=$}
{$<$}
{Không so sánh được}
\loigiai{
Đổi $28\text{ km}^2 = 28 \times 100 = 2\,800\text{ ha}$. Vậy $2\,800\text{ ha} = 28\text{ km}^2$.\\
\textbf{Đáp án của câu hỏi là B}
}
\end{ex}

\begin{ex} %[5H3H18]
Đổi $450\text{ dm}^2$ ra mét vuông ta được số thập phân là
\choice
{\True $4{,}5\text{ m}^2$}
{$0{,}45\text{ m}^2$}
{$45\text{ m}^2$}
{$450\text{ m}^2$}
\loigiai{
Ta có: $450\text{ dm}^2 = 450 : 100 = 4{,}5\text{ m}^2$.\\
\textbf{Đáp án của câu hỏi là A}
}
\end{ex}

\begin{ex} %[5H3V18]
Người ta muốn xây dựng khu vui chơi kết hợp khu cắm trại trên mảnh đất $2\text{ ha}$. Biết khu vui chơi chiếm $\dfrac{3}{4}$ diện tích mảnh đất. Diện tích khu cắm trại là
\choice
{$15\,000\text{ m}^2$}
{$10\,000\text{ m}^2$}
{\True $5\,000\text{ m}^2$}
{$7\,500\text{ m}^2$}
\loigiai{
Đổi $2\text{ ha} = 20\,000\text{ m}^2$.\\
Khu cắm trại chiếm: $1 - \dfrac{3}{4} = \dfrac{1}{4}$ diện tích mảnh đất.\\
Diện tích khu cắm trại là: $20\,000 \times \dfrac{1}{4} = 5\,000\text{ (m}^2\text{)}$.\\
\textbf{Đáp án của câu hỏi là C}
}
\end{ex}

\begin{ex} %[5H3V18]
Bác Ba muốn dùng gỗ để lát sàn phòng khách hình chữ nhật có chiều dài $6\text{ m}$, chiều rộng $5\text{ m}$. Loại gỗ bác chọn có giá $300\,000$ đồng/$1\text{ m}^2$. Số tiền bác Ba cần trả là
\choice
{$3\,000\,000$ đồng}
{$6\,000\,000$ đồng}
{$12\,000\,000$ đồng}
{\True $9\,000\,000$ đồng}
\loigiai{
Diện tích sàn phòng khách là: $6 \times 5 = 30\text{ (m}^2\text{)}$.\\
Số tiền bác Ba cần trả là: $30 \times 300\,000 = 9\,000\,000$ (đồng).\\
\textbf{Đáp án của câu hỏi là D}
}
\end{ex}

\begin{ex} %[5H3V18]
Người ta cắt một tấm kính hình chữ nhật thành $5$ tấm kính nhỏ giống nhau để đóng khung ảnh, mỗi tấm nhỏ có kích thước $6\text{ dm} \times 4\text{ dm}$. Diện tích tấm kính ban đầu theo đơn vị mét vuông là
\choice
{\True $1{,}2\text{ m}^2$}
{$12\text{ m}^2$}
{$0{,}12\text{ m}^2$}
{$24\text{ m}^2$}
\loigiai{
Diện tích một tấm kính nhỏ là: $6 \times 4 = 24\text{ (dm}^2\text{)}$.\\
Diện tích tấm kính ban đầu là: $24 \times 5 = 120\text{ (dm}^2\text{)} = 1{,}2\text{ m}^2$.\\
\textbf{Đáp án của câu hỏi là A}
}
\end{ex}

\begin{ex} %[5H3V18]
Một mảnh đất hình chữ nhật gồm $12$ ô đất hình vuông, mỗi ô có cạnh $10\text{ m}$. Người ta dùng mảnh đất đó để xây nhà máy. Diện tích xây nhà máy là bao nhiêu héc-ta?
\choice
{$1{,}2\text{ ha}$}
{$12\text{ ha}$}
{$0{,}012\text{ ha}$}
{\True $0{,}12\text{ ha}$}
\loigiai{
Diện tích $1$ ô vuông là: $10 \times 10 = 100\text{ (m}^2\text{)}$.\\
Diện tích mảnh đất là: $12 \times 100 = 1\,200\text{ (m}^2\text{)}$.\\
Đổi $1\,200\text{ m}^2 = 1\,200 : 10\,000 = 0{,}12\text{ ha}$.\\
\textbf{Đáp án của câu hỏi là D}
}
\end{ex}

\textbf{Phần 2. Trắc nghiệm đúng sai}
\setcounterpageref{ex}{0}

\begin{ex} %[5H3H18]
Xét tính đúng sai của các khẳng định sau về các phép đổi đơn vị đo diện tích:
\choiceTF
{\True $412\,000\text{ m}^2 = 41{,}2\text{ ha}$.}
{$8\,000\text{ m}^2 = 8\text{ ha}$.}
{\True $11\text{ km}^2 = 1\,100\text{ ha}$.}
{\True $8\text{ ha} = 80\,000\text{ m}^2$.}
\loigiai{
\begin{itemize}
    \item $412\,000\text{ m}^2 = 412\,000 : 10\,000 = 41{,}2\text{ ha}$.
    \item $8\,000\text{ m}^2 = 8\,000 : 10\,000 = 0{,}8\text{ ha} \neq 8\text{ ha}$.
    \item $11\text{ km}^2 = 11 \times 100 = 1\,100\text{ ha}$.
    \item $8\text{ ha} = 8 \times 10\,000 = 80\,000\text{ m}^2$.
\end{itemize}
\begin{itemize}
    \item a) \textbf{Đúng}.
    \item b) \textbf{Sai}: Kết quả là $0{,}8\text{ ha}$.
    \item c) \textbf{Đúng}.
    \item d) \textbf{Đúng}.
\end{itemize}
}
\end{ex}

\begin{ex} %[5H3V18]
Một mảnh đất có diện tích $2\text{ ha}$ được chia thành hai phần: khu vui chơi chiếm $\dfrac{3}{4}$ diện tích và khu cắm trại chiếm $\dfrac{1}{4}$ diện tích. Xét tính đúng sai của các khẳng định sau:
\choiceTF
{\True Diện tích cả mảnh đất là $20\,000\text{ m}^2$.}
{Diện tích khu vui chơi là $1\,500\text{ m}^2$.}
{\True Diện tích khu cắm trại là $5\,000\text{ m}^2$.}
{\True Diện tích khu vui chơi lớn hơn diện tích khu cắm trại là $1\text{ ha}$.}
\loigiai{
\begin{itemize}
    \item $2\text{ ha} = 20\,000\text{ m}^2$.
    \item Diện tích khu vui chơi: $20\,000 \times \dfrac{3}{4} = 15\,000\text{ m}^2 \neq 1\,500\text{ m}^2$.
    \item Diện tích khu cắm trại: $20\,000 \times \dfrac{1}{4} = 5\,000\text{ m}^2$.
    \item Hiệu diện tích: $15\,000 - 5\,000 = 10\,000\text{ m}^2 = 1\text{ ha}$.
\end{itemize}
\begin{itemize}
    \item a) \textbf{Đúng}.
    \item b) \textbf{Sai}: Diện tích là $15\,000\text{ m}^2$.
    \item c) \textbf{Đúng}.
    \item d) \textbf{Đúng}.
\end{itemize}
}
\end{ex}

\textbf{Phần 3. Trả lời ngắn}
\setcounterpageref{ex}{0}

\begin{ex} %[5H3H18]
Đổi $8\,000\text{ m}^2$ ra héc-ta. (Nhập kết quả dưới dạng số thập phân).
\shortans{0,8}
\loigiai{
\textbf{Đáp số:} $0{,}8$\\
$8\,000\text{ m}^2 = 8\,000 : 10\,000 = 0{,}8\text{ ha}$.
}
\end{ex}

\begin{ex} %[5H3H18]
Đổi $0{,}48\text{ km}^2$ ra héc-ta.
\shortans{48}
\loigiai{
\textbf{Đáp số:} $48$\\
$0{,}48\text{ km}^2 = 0{,}48 \times 100 = 48\text{ ha}$.
}
\end{ex}

\begin{ex} %[5H3H18]
Đổi $450\text{ dm}^2$ ra mét vuông. (Nhập kết quả dưới dạng số thập phân).
\shortans{4,5}
\loigiai{
\textbf{Đáp số:} $4{,}5$\\
$450\text{ dm}^2 = 450 : 100 = 4{,}5\text{ m}^2$.
}
\end{ex}

\begin{ex} %[5H3V18]
Bác Ba lát sàn phòng khách dài $6\text{ m}$, rộng $5\text{ m}$ với giá gỗ $300\,000$ đồng/$1\text{ m}^2$. Hỏi bác Ba phải trả bao nhiêu triệu đồng tiền gỗ?
\shortans{9}
\loigiai{
\textbf{Đáp số:} $9$\\
Diện tích: $6 \times 5 = 30\text{ m}^2$. Chi phí: $30 \times 300\,000 = 9\,000\,000$ đồng $= 9$ triệu đồng.
}
\end{ex}

\begin{ex} %[5H3V18]
Cắt một tấm kính lớn thành $5$ tấm nhỏ cùng kích thước $6\text{ dm} \times 4\text{ dm}$. Tính diện tích tấm kính lớn ban đầu theo đơn vị mét vuông. (Nhập kết quả dưới dạng số thập phân).
\shortans{1,2}
\loigiai{
\textbf{Đáp số:} $1{,}2$\\
Diện tích là: $5 \times (6 \times 4) = 120\text{ dm}^2 = 1{,}2\text{ m}^2$.
}
\end{ex}

\begin{ex} %[5H3V18]
Mảnh đất gồm $12$ ô vuông nhỏ, mỗi ô có cạnh dài $10\text{ m}$. Tính diện tích mảnh đất đó theo đơn vị héc-ta. (Nhập kết quả dưới dạng số thập phân).
\shortans{0,12}
\loigiai{
\textbf{Đáp số:} $0{,}12$\\
Diện tích: $12 \times (10 \times 10) = 1\,200\text{ m}^2 = 0{,}12\text{ ha}$.
}
\end{ex}

\newpage

% =========================================================================
% CHỦ ĐỀ 4 - CÁC PHÉP TÍNH VỚI SỐ THẬP PHÂN
% =========================================================================
\section{Các phép tính với số thập phân}

\TaoPhieuKiemTra{Chủ đề 4 -- Các phép tính với số thập phân}{19}{Phép cộng số thập phân}{https://loptoanthaylinh.vn/thi-online/toan-5/bai-19}
\textbf{Phần 1. Trắc nghiệm lựa chọn}
\setcounterpageref{ex}{0}

\begin{ex} %[5D4N19]
Kết quả của phép tính $6{,}8 + 9{,}7$ là
\choice
{$15{,}5$}
{$16{,}15$}
{$17{,}5$}
{\True $16{,}5$}
\loigiai{
Ta có: $6{,}8 + 9{,}7 = 16{,}5$.\\
\textbf{Đáp án của câu hỏi là D}
}
\end{ex}

\begin{ex} %[5D4N19]
Kết quả của phép tính $24{,}5 + 3{,}84$ là
\choice
{$27{,}34$}
{$28{,}89$}
{\True $28{,}34$}
{$62{,}9$}
\loigiai{
Đặt tính thẳng cột: $24{,}50 + 3{,}84 = 28{,}34$.\\
\textbf{Đáp án của câu hỏi là C}
}
\end{ex}

\begin{ex} %[5D4H19]
Kết quả của phép tính $8{,}16 + 5{,}9$ là
\choice
{\True $14{,}06$}
{$13{,}25$}
{$14{,}6$}
{$13{,}06$}
\loigiai{
Ta có: $8{,}16 + 5{,}90 = 14{,}06$.\\
\textbf{Đáp án của câu hỏi là A}
}
\end{ex}

\begin{ex} %[5D4H19]
Kết quả của phép tính $18{,}6 + 8{,}72$ là
\choice
{$26{,}32$}
{\True $27{,}32$}
{$27{,}78$}
{$105{,}8$}
\loigiai{
Ta có: $18{,}60 + 8{,}72 = 27{,}32$.\\
\textbf{Đáp án của câu hỏi là B}
}
\end{ex}

\begin{ex} %[5D4H19]
Kết quả của phép tính $4{,}61 + 8$ là
\choice
{\True $12{,}61$}
{$5{,}41$}
{$4{,}69$}
{$12{,}69$}
\loigiai{
Cộng phần nguyên: $4 + 8 = 12$, giữ nguyên phần thập phân: $12{,}61$.\\
\textbf{Đáp án của câu hỏi là A}
}
\end{ex}

\begin{ex} %[5D4H19]
Kết quả của phép tính $23{,}58 + 5{,}271$ là
\choice
{$28{,}751$}
{\True $28{,}851$}
{$76{,}29$}
{$28{,}85$}
\loigiai{
Ta có: $23{,}580 + 5{,}271 = 28{,}851$.\\
\textbf{Đáp án của câu hỏi là B}
}
\end{ex}

\begin{ex} %[5D4H19]
Biểu thức $37{,}5 + 4{,}6 = 4{,}6 + 37{,}5$ thể hiện tính chất nào của phép cộng?
\choice
{Tính chất kết hợp}
{Tính chất cộng với số 0}
{\True Tính chất giao hoán}
{Tính chất phân phối}
\loigiai{
Khi đổi chỗ hai số hạng trong một tổng thì tổng không thay đổi (tính chất giao hoán).\\
\textbf{Đáp án của câu hỏi là C}
}
\end{ex}

\begin{ex} %[5D4H19]
Tính thuận tiện biểu thức $6 + 8{,}46 + 1{,}54$ ta được kết quả là
\choice
{$15$}
{$17$}
{$16{,}46$}
{\True $16$}
\loigiai{
Ta có: $6 + (8{,}46 + 1{,}54) = 6 + 10 = 16$.\\
\textbf{Đáp án của câu hỏi là D}
}
\end{ex}

\begin{ex} %[5D4H19]
Tính thuận tiện biểu thức $4{,}8 + 0{,}73 + 3{,}2$ ta được kết quả là
\choice
{\True $8{,}73$}
{$7{,}73$}
{$8{,}37$}
{$9{,}73$}
\loigiai{
Ta có: $(4{,}8 + 3{,}2) + 0{,}73 = 8 + 0{,}73 = 8{,}73$.\\
\textbf{Đáp án của câu hỏi là A}
}
\end{ex}

\begin{ex} %[5D4V19]
Mai cùng mẹ đi siêu thị mua một quả dưa hấu cân nặng $4{,}65\text{ kg}$ và một quả mít cân nặng $5{,}8\text{ kg}$. Cả hai quả dưa hấu và mít cân nặng tất cả là
\choice
{$9{,}45\text{ kg}$}
{$10{,}35\text{ kg}$}
{\True $10{,}45\text{ kg}$}
{$11{,}45\text{ kg}$}
\loigiai{
Cả hai quả cân nặng: $4{,}65 + 5{,}8 = 10{,}45\text{ (kg)}$.\\
\textbf{Đáp án của câu hỏi là C}
}
\end{ex}

\begin{ex} %[5D4V19]
Ngày thứ Sáu chú Sơn thu được $73{,}5\text{ l}$ sữa dê, ngày thứ Bảy thu được nhiều hơn ngày thứ Sáu $4{,}75\text{ l}$ sữa dê. Cả hai ngày chú Sơn thu được số lít sữa dê là
\choice
{$78{,}25\text{ l}$}
{$150{,}75\text{ l}$}
{$152{,}25\text{ l}$}
{\True $151{,}75\text{ l}$}
\loigiai{
Ngày thứ Bảy chú Sơn thu được: $73{,}5 + 4{,}75 = 78{,}25\text{ (l)}$.\\
Cả hai ngày chú Sơn thu được: $73{,}5 + 78{,}25 = 151{,}75\text{ (l)}$.\\
\textbf{Đáp án của câu hỏi là D}
}
\end{ex}

\begin{ex} %[5D4V19]
Sân trường dạng hình chữ nhật có chiều rộng $17{,}5\text{ m}$, chiều dài hơn chiều rộng $15\text{ m}$. Chu vi sân trường đó là
\choice
{$50\text{ m}$}
{\True $100\text{ m}$}
{$65\text{ m}$}
{$130\text{ m}$}
\loigiai{
Chiều dài sân trường là: $17{,}5 + 15 = 32{,}5\text{ (m)}$.\\
Chu vi sân trường là: $(32{,}5 + 17{,}5) \times 2 = 50 \times 2 = 100\text{ (m)}$.\\
\textbf{Đáp án của câu hỏi là B}
}
\end{ex}

\textbf{Phần 2. Trắc nghiệm đúng sai}
\setcounterpageref{ex}{0}

\begin{ex} %[5D4H19]
Xét tính đúng sai của việc thực hiện các phép tính sau:
\choiceTF
{Phép cộng $5{,}4 + 3{,}9 = 8{,}13$.}
{\True Phép cộng $2{,}45 + 4{,}37 = 6{,}82$.}
{Phép cộng $6{,}53 + 12{,}8 = 78{,}1$.}
{\True Phép tính $(5 + 2{,}7) + 7{,}3 = 5 + (2{,}7 + 7{,}3) = 15$.}
\loigiai{
\begin{itemize}
    \item $5{,}4 + 3{,}9 = 9{,}3 \neq 8{,}13$ (sai do cộng từng phần rồi ghép).
    \item $2{,}45 + 4{,}37 = 6{,}82$.
    \item $6{,}53 + 12{,}8 = 19{,}33 \neq 78{,}1$ (sai do đặt tính lệch cột).
    \item $(5 + 2{,}7) + 7{,}3 = 5 + 10 = 15$.
\end{itemize}
\begin{itemize}
    \item a) \textbf{Sai}: $5{,}4 + 3{,}9 = 9{,}3$.
    \item b) \textbf{Đúng}.
    \item c) \textbf{Sai}: $6{,}53 + 12{,}8 = 19{,}33$.
    \item d) \textbf{Đúng}.
\end{itemize}
}
\end{ex}

\begin{ex} %[5D4V19]
Hai bạn Mai và Việt dùng dây đồng để uốn hình. Mai uốn hình ngôi sao hết $1{,}65\text{ m}$ dây đồng, Việt uốn hình tam giác hết $1{,}26\text{ m}$ dây đồng. Xét tính đúng sai của các khẳng định sau:
\choiceTF
{\True Chiều dài dây đồng Mai dùng nhiều hơn của Việt là $0{,}39\text{ m}$.}
{\True Tổng chiều dài dây đồng cả hai bạn đã dùng là $2{,}91\text{ m}$.}
{Nếu có đoạn dây đồng dài $3\text{ m}$, sau khi hai bạn uốn xong thì còn thừa $0{,}19\text{ m}$.}
{\True Đổi $2{,}91\text{ m} = 291\text{ cm}$.}
\loigiai{
\begin{itemize}
    \item Mai dùng nhiều hơn Việt: $1{,}65 - 1{,}26 = 0{,}39\text{ m}$.
    \item Cả hai bạn dùng: $1{,}65 + 1,26 = 2{,}91\text{ m}$.
    \item Đoạn dây còn thừa: $3 - 2{,}91 = 0{,}09\text{ m} \neq 0{,}19\text{ m}$.
    \item $2{,}91\text{ m} = 291\text{ cm}$.
\end{itemize}
\begin{itemize}
    \item a) \textbf{Đúng}.
    \item b) \textbf{Đúng}.
    \item c) \textbf{Sai}: Còn thừa $0{,}09\text{ m}$.
    \item d) \textbf{Đúng}.
\end{itemize}
}
\end{ex}

\textbf{Phần 3. Trả lời ngắn}
\setcounterpageref{ex}{0}

\begin{ex} %[5D4H19]
Tính tổng của hai số $5{,}34$ và $7{,}49$. (Nhập kết quả dưới dạng số thập phân).
\shortans{12,83}
\loigiai{
\textbf{Đáp số:} $12{,}83$\\
$5{,}34 + 7{,}49 = 12{,}83$.
}
\end{ex}

\begin{ex} %[5D4H19]
Tính tổng của hai số $7{,}635$ và $6{,}249$. (Nhập kết quả dưới dạng số thập phân).
\shortans{13,884}
\loigiai{
\textbf{Đáp số:} $13{,}884$\\
$7{,}635 + 6{,}249 = 13{,}884$.
}
\end{ex}

\begin{ex} %[5D4H19]
Tính nhanh giá trị biểu thức: $4{,}8 + 0{,}73 + 3{,}2$. (Nhập kết quả dưới dạng số thập phân).
\shortans{8,73}
\loigiai{
\textbf{Đáp số:} $8{,}73$\\
$(4{,}8 + 3{,}2) + 0{,}73 = 8 + 0{,}73 = 8{,}73$.
}
\end{ex}

\begin{ex} %[5D4V19]
Một quả dưa hấu nặng $4{,}65\text{ kg}$, quả mít nặng $5{,}8\text{ kg}$. Hỏi cả hai quả nặng bao nhiêu ki-lô-gam? (Nhập kết quả dưới dạng số thập phân).
\shortans{10,45}
\loigiai{
\textbf{Đáp số:} $10{,}45$\\
$4{,}65 + 5{,}8 = 10{,}45\text{ (kg)}$.
}
\end{ex}

\begin{ex} %[5D4V19]
Ngày thứ nhất đội công nhân sửa được $73{,}5\text{ m}$ đường, ngày thứ hai sửa được nhiều hơn ngày thứ nhất $4{,}75\text{ m}$ đường. Hỏi cả hai ngày đội sửa được bao nhiêu mét đường? (Nhập kết quả dưới dạng số thập phân).
\shortans{151,75}
\loigiai{
\textbf{Đáp số:} $151{,}75$\\
Ngày hai: $73{,}5 + 4{,}75 = 78{,}25\text{ m}$. Cả hai ngày: $73{,}5 + 78{,}25 = 151{,}75\text{ m}$.
}
\end{ex}

\begin{ex} %[5D4V19]
Một sân trường hình chữ nhật có chiều rộng $17{,}5\text{ m}$, chiều dài hơn chiều rộng $15\text{ m}$. Tính chu vi sân trường đó theo đơn vị mét.
\shortans{100}
\loigiai{
\textbf{Đáp số:} $100$\\
Chiều dài: $17{,}5 + 15 = 32{,}5\text{ m}$. Chu vi: $(32{,}5 + 17{,}5) \times 2 = 100\text{ m}$.
}
\end{ex}

\newpage

\TaoPhieuKiemTra{Chủ đề 4 -- Các phép tính với số thập phân}{20}{Phép trừ số thập phân}{https://loptoanthaylinh.vn/thi-online/toan-5/bai-20}
\textbf{Phần 1. Trắc nghiệm lựa chọn}
\setcounterpageref{ex}{0}

\begin{ex} %[5D4N20]
Kết quả của phép tính $5{,}8 - 3{,}9$ là
\choice
{\True $1{,}9$}
{$2{,}9$}
{$1{,}1$}
{$2{,}1$}
\loigiai{
Ta có: $5{,}8 - 3{,}9 = 1{,}9$.\\
\textbf{Đáp án của câu hỏi là A}
}
\end{ex}

\begin{ex} %[5D4N20]
Kết quả của phép tính $63{,}49 - 1{,}8$ là
\choice
{$62{,}41$}
{\True $61{,}69$}
{$61{,}41$}
{$45{,}49$}
\loigiai{
Đặt tính thẳng cột: $63{,}49 - 1{,}80 = 61{,}69$.\\
\textbf{Đáp án của câu hỏi là B}
}
\end{ex}

\begin{ex} %[5D4H20]
Kết quả của phép tính $17{,}96 - 8{,}5$ là
\choice
{\True $9{,}46$}
{$9{,}91$}
{$9{,}41$}
{$17{,}11$}
\loigiai{
Ta có: $17{,}96 - 8{,}50 = 9{,}46$.\\
\textbf{Đáp án của câu hỏi là A}
}
\end{ex}

\begin{ex} %[5D4H20]
Kết quả của phép tính $25{,}9 - 13{,}84$ là
\choice
{$12{,}14$}
{\True $12{,}06$}
{$12{,}16$}
{$12{,}86$}
\loigiai{
Ta có: $25{,}90 - 13{,}84 = 12{,}06$.\\
\textbf{Đáp án của câu hỏi là B}
}
\end{ex}

\begin{ex} %[5D4H20]
Kết quả của phép tính $7{,}6 - 1{,}51$ là
\choice
{$6{,}19$}
{$6{,}11$}
{$5{,}09$}
{\True $6{,}09$}
\loigiai{
Ta có: $7{,}60 - 1{,}51 = 6{,}09$.\\
\textbf{Đáp án của câu hỏi là D}
}
\end{ex}

\begin{ex} %[5D4H20]
Kết quả của phép tính $21{,}4 - 6$ là
\choice
{$20{,}8$}
{$14{,}6$}
{\True $15{,}4$}
{$15{,}6$}
\loigiai{
Trừ phần nguyên: $21 - 6 = 15$, giữ nguyên phần thập phân: $15{,}4$.\\
\textbf{Đáp án của câu hỏi là C}
}
\end{ex}

\begin{ex} %[5D4H20]
Kết quả của phép tính $9 - 3{,}5$ là
\choice
{$6{,}5$}
{$5{,}0$}
{$6{,}0$}
{\True $5{,}5$}
\loigiai{
Ta có: $9{,}0 - 3{,}5 = 5{,}5$.\\
\textbf{Đáp án của câu hỏi là D}
}
\end{ex}

\begin{ex} %[5D4H20]
Tìm số thích hợp điền vào dấu hỏi chấm: $? + 8{,}16 = 17{,}5$.
\choice
{$25{,}66$}
{$9{,}46$}
{\True $9{,}34$}
{$8{,}34$}
\loigiai{
Ta có: $? = 17{,}5 - 8{,}16 = 17{,}50 - 8{,}16 = 9{,}34$.\\
\textbf{Đáp án của câu hỏi là C}
}
\end{ex}

\begin{ex} %[5D4H20]
Tìm số thích hợp điền vào dấu hỏi chấm: $? - 6{,}17 = 11{,}83$.
\choice
{$5{,}66$}
{$17{,}00$}
{$18{,}1$}
{\True $18$}
\loigiai{
Ta có: $? = 11{,}83 + 6{,}17 = 18{,}00 = 18$.\\
\textbf{Đáp án của câu hỏi là D}
}
\end{ex}

\begin{ex} %[5D4V20]
Mai làm được $2{,}15\text{ l}$ nước mơ, Mi làm được $1{,}7\text{ l}$ nước dâu. Hỏi lượng nước mơ của Mai nhiều hơn lượng nước dâu của Mi bao nhiêu lít?
\choice
{$0{,}35\text{ l}$}
{\True $0{,}45\text{ l}$}
{$1{,}45\text{ l}$}
{$0{,}045\text{ l}$}
\loigiai{
Nước mơ nhiều hơn nước dâu là: $2{,}15 - 1{,}7 = 0{,}45\text{ (l)}$.\\
\textbf{Đáp án của câu hỏi là B}
}
\end{ex}

\begin{ex} %[5D4V20]
Một chiếc cọc được sơn hai màu xanh và đỏ. Phần màu xanh dài $14{,}2\text{ dm}$ và dài hơn đoạn màu đỏ $1{,}8\text{ dm}$. Tổng chiều dài của chiếc cọc đó là
\choice
{\True $26{,}6\text{ dm}$}
{$12{,}4\text{ dm}$}
{$16{,}0\text{ dm}$}
{$28{,}4\text{ dm}$}
\loigiai{
Đoạn màu đỏ dài: $14{,}2 - 1{,}8 = 12{,}4\text{ (dm)}$.\\
Chiều dài cả chiếc cọc là: $14{,}2 + 12{,}4 = 26{,}6\text{ (dm)}$.\\
\textbf{Đáp án của câu hỏi là A}
}
\end{ex}

\begin{ex} %[5D4V20]
Cả ba chú rô-bốt A, B, C có tổng cân nặng là $8\text{ kg}$. Rô-bốt A và B cân nặng $4{,}7\text{ kg}$, rô-bốt B và C cân nặng $5{,}5\text{ kg}$. Cân nặng của rô-bốt A là
\choice
{$2{,}2\text{ kg}$}
{$3{,}3\text{ kg}$}
{\True $2{,}5\text{ kg}$}
{$3{,}5\text{ kg}$}
\loigiai{
Cân nặng của rô-bốt A là: $8 - 5{,}5 = 2{,}5\text{ (kg)}$.\\
\textbf{Đáp án của câu hỏi là C}
}
\end{ex}

\textbf{Phần 2. Trắc nghiệm đúng sai}
\setcounterpageref{ex}{0}

\begin{ex} %[5D4H20]
Xét tính đúng sai của các phép tính sau:
\choiceTF
{Phép trừ $6{,}15 - 2{,}7 = 4{,}8$.}
{\True Phép trừ $4{,}38 - 1{,}56 = 2{,}82$.}
{Phép trừ $85{,}9 - 5{,}06 = 35{,}3$.}
{\True Phép trừ $4{,}43 - 4{,}16 = 0{,}27$.}
\loigiai{
\begin{itemize}
    \item $6{,}15 - 2{,}70 = 3{,}45 \neq 4{,}8$.
    \item $4{,}38 - 1{,}56 = 2{,}82$.
    \item $85{,}90 - 5{,}06 = 80{,}84 \neq 35{,}3$ (sai do đặt tính lệch hàng).
    \item $4{,}43 - 4{,}16 = 0{,}27$.
\end{itemize}
\begin{itemize}
    \item a) \textbf{Sai}: Kết quả là $3{,}45$.
    \item b) \textbf{Đúng}.
    \item c) \textbf{Sai}: Kết quả là $80{,}84$.
    \item d) \textbf{Đúng}.
\end{itemize}
}
\end{ex}

\begin{ex} %[5D4V20]
Cho ba chú rô-bốt A, B, C có tổng cân nặng là $8\text{ kg}$. Biết rô-bốt A và B nặng $4{,}7\text{ kg}$; rô-bốt B và C nặng $5{,}5\text{ kg}$. Xét tính đúng sai của các khẳng định sau:
\choiceTF
{\True Rô-bốt C cân nặng $3{,}3\text{ kg}$.}
{\True Rô-bốt B cân nặng $2{,}2\text{ kg}$.}
{Rô-bốt A nặng hơn rô-bốt C là $0{,}8\text{ kg}$.}
{\True Trong ba chú rô-bốt, rô-bốt B là nhẹ nhất.}
\loigiai{
\begin{itemize}
    \item Cân nặng của rô-bốt C: $8 - 4{,}7 = 3{,}3\text{ kg}$.
    \item Cân nặng của rô-bốt A: $8 - 5{,}5 = 2{,}5\text{ kg}$.
    \item Cân nặng của rô-bốt B: $4{,}7 - 2{,}5 = 2{,}2\text{ kg}$.
    \item Rô-bốt C nặng hơn rô-bốt A: $3{,}3 - 2{,}5 = 0{,}8\text{ kg}$ (Rô-bốt A nhẹ hơn rô-bốt C).
    \item So sánh: $2{,}2 < 2{,}5 < 3{,}3 \Rightarrow$ Rô-bốt B nhẹ nhất.
\end{itemize}
\begin{itemize}
    \item a) \textbf{Đúng}.
    \item b) \textbf{Đúng}.
    \item c) \textbf{Sai}: Rô-bốt A nhẹ hơn rô-bốt C.
    \item d) \textbf{Đúng}.
\end{itemize}
}
\end{ex}

\textbf{Phần 3. Trả lời ngắn}
\setcounterpageref{ex}{0}

\begin{ex} %[5D4H20]
Tính hiệu của hai số $2{,}53$ và $1{,}62$. (Nhập kết quả dưới dạng số thập phân).
\shortans{0,91}
\loigiai{
\textbf{Đáp số:} $0{,}91$\\
$2{,}53 - 1{,}62 = 0{,}91$.
}
\end{ex}

\begin{ex} %[5D4H20]
Tính giá trị của biểu thức: $25{,}9 - 13{,}84$. (Nhập kết quả dưới dạng số thập phân).
\shortans{12,06}
\loigiai{
\textbf{Đáp số:} $12{,}06$\\
$25{,}90 - 13{,}84 = 12{,}06$.
}
\end{ex}

\begin{ex} %[5D4H20]
Tìm số thích hợp điền vào dấu hỏi chấm: $8{,}9 + ? = 28{,}501$. (Nhập kết quả dưới dạng số thập phân).
\shortans{19,601}
\loigiai{
\textbf{Đáp số:} $19{,}601$\\
$? = 28{,}501 - 8{,}9 = 19{,}601$.
}
\end{ex}

\begin{ex} %[5D4V20]
Thành tích nhảy xa năm nay của Nam là $4{,}43\text{ m}$, năm ngoái là $4{,}16\text{ m}$. Hỏi thành tích năm nay của Nam hơn năm ngoái bao nhiêu mét? (Nhập kết quả dưới dạng số thập phân).
\shortans{0,27}
\loigiai{
\textbf{Đáp số:} $0{,}27$\\
$4{,}43 - 4{,}16 = 0{,}27\text{ (m)}$.
}
\end{ex}

\begin{ex} %[5D4V20]
Một chiếc cọc có phần màu xanh dài $14{,}2\text{ dm}$ và phần màu đỏ ngắn hơn phần màu xanh $1{,}8\text{ dm}$. Hỏi toàn bộ chiếc cọc dài bao nhiêu đề-xi-mét? (Nhập kết quả dưới dạng số thập phân).
\shortans{26,6}
\loigiai{
\textbf{Đáp số:} $26{,}6$\\
Phần màu đỏ: $14{,}2 - 1{,}8 = 12{,}4\text{ dm}$. Chiếc cọc dài: $14{,}2 + 12{,}4 = 26{,}6\text{ dm}$.
}
\end{ex}

\begin{ex} %[5D4V20]
Cả ba rô-bốt A, B, C cân nặng $8\text{ kg}$, rô-bốt A và B cân nặng $4{,}7\text{ kg}$, rô-bốt B và C cân nặng $5{,}5\text{ kg}$. Hỏi riêng chú rô-bốt B cân nặng bao nhiêu ki-lô-gam? (Nhập kết quả dưới dạng số thập phân).
\shortans{2,2}
\loigiai{
\textbf{Đáp số:} $2{,}2$\\
Rô-bốt A: $8 - 5{,}5 = 2{,}5\text{ kg}$. Rô-bốt B: $4{,}7 - 2{,}5 = 2{,}2\text{ kg}$.
}
\end{ex}

\newpage

\TaoPhieuKiemTra{Chủ đề 4 -- Các phép tính với số thập phân}{21}{Phép nhân số thập phân}{https://loptoanthaylinh.vn/thi-online/toan-5/bai-21}
\textbf{Phần 1. Trắc nghiệm lựa chọn}
\setcounterpageref{ex}{0}

\begin{ex} %[5D4N21]
Kết quả của phép tính $7{,}8 \times 6$ là
\choice
{$42{,}8$}
{$46{,}4$}
{$48{,}8$}
{\True $46{,}8$}
\loigiai{
Ta có: $7{,}8 \times 6 = 46{,}8$.\\
\textbf{Đáp án của câu hỏi là D}
}
\end{ex}

\begin{ex} %[5D4N21]
Kết quả của phép tính $0{,}72 \times 50$ là
\choice
{\True $36$}
{$3{,}6$}
{$360$}
{$0{,}36$}
\loigiai{
Ta có: $0{,}72 \times 50 = 36$.\\
\textbf{Đáp án của câu hỏi là A}
}
\end{ex}

\begin{ex} %[5D4H21]
Kết quả của phép tính $4{,}3 \times 3{,}6$ là
\choice
{$154{,}8$}
{$1{,}548$}
{\True $15{,}48$}
{$15{,}44$}
\loigiai{
Hai thừa số có tổng cộng $2$ chữ số ở phần thập phân, vậy tích là $15{,}48$.\\
\textbf{Đáp án của câu hỏi là C}
}
\end{ex}

\begin{ex} %[5D4H21]
Kết quả của phép tính $7{,}5 \times 3{,}4$ là
\choice
{$24{,}5$}
{\True $25{,}5$}
{$25{,}50$}
{$255$}
\loigiai{
Ta có: $7{,}5 \times 3{,}4 = 25{,}5$.\\
\textbf{Đáp án của câu hỏi là B}
}
\end{ex}

\begin{ex} %[5D4H21]
Kết quả của phép tính $6{,}8 \times 0{,}52$ là
\choice
{\True $3{,}536$}
{$35{,}36$}
{$0{,}3536$}
{$353{,}6$}
\loigiai{
Hai thừa số có tổng cộng $3$ chữ số ở phần thập phân, vậy tích là $3{,}536$.\\
\textbf{Đáp án của câu hỏi là A}
}
\end{ex}

\begin{ex} %[5D4H21]
Cho biết $64 \times 57 = 3\,648$. Không thực hiện phép tính, giá trị của $6{,}4 \times 0{,}57$ là
\choice
{$36{,}48$}
{\True $3{,}648$}
{$0{,}3648$}
{$364{,}8$}
\loigiai{
Thừa số $6{,}4$ có $1$ chữ số thập phân, $0{,}57$ có $2$ chữ số thập phân $\Rightarrow$ Tích có $1 + 2 = 3$ chữ số thập phân. Vậy $6{,}4 \times 0{,}57 = 3{,}648$.\\
\textbf{Đáp án của câu hỏi là B}
}
\end{ex}

\begin{ex} %[5D4H21]
Tính thuận tiện biểu thức $2{,}5 \times 8{,}6 \times 4$ ta được kết quả là
\choice
{$8{,}6$}
{$860$}
{$86{,}4$}
{\True $86$}
\loigiai{
Ta có: $(2{,}5 \times 4) \times 8{,}6 = 10 \times 8{,}6 = 86$.\\
\textbf{Đáp án của câu hỏi là D}
}
\end{ex}

\begin{ex} %[5D4H21]
Tính thuận tiện biểu thức $6{,}84 \times 0{,}2 \times 5$ ta được kết quả là
\choice
{$68{,}4$}
{$0{,}684$}
{\True $6{,}84$}
{$13{,}68$}
\loigiai{
Ta có: $6{,}84 \times (0{,}2 \times 5) = 6{,}84 \times 1 = 6{,}84$.\\
\textbf{Đáp án của câu hỏi là C}
}
\end{ex}

\begin{ex} %[5D4V21]
Một ô tô chạy trên đường cao tốc, mỗi giờ đi được $84{,}5\text{ km}$. Trong $1{,}2$ giờ ô tô đó đi được quãng đường là
\choice
{$91{,}4\text{ km}$}
{\True $101{,}4\text{ km}$}
{$100{,}4\text{ km}$}
{$102{,}4\text{ km}$}
\loigiai{
Quãng đường ô tô đi được là: $84{,}5 \times 1{,}2 = 101{,}4\text{ (km)}$.\\
\textbf{Đáp án của câu hỏi là B}
}
\end{ex}

\begin{ex} %[5D4V21]
Biết rằng xay xát $1\text{ kg}$ thóc thì thu được $0{,}64\text{ kg}$ gạo. Hỏi xay xát $50\text{ kg}$ thóc loại đó thì thu được bao nhiêu ki-lô-gam gạo?
\choice
{$3{,}2\text{ kg}$}
{$30\text{ kg}$}
{\True $32\text{ kg}$}
{$320\text{ kg}$}
\loigiai{
Lượng gạo thu được là: $50 \times 0{,}64 = 32\text{ (kg)}$.\\
\textbf{Đáp án của câu hỏi là C}
}
\end{ex}

\begin{ex} %[5D4V21]
Mỗi cốc nước cam chứa $0{,}25\text{ l}$. Ba bạn nhỏ mỗi bạn uống một cốc nước cam thì uống hết tất cả số lít nước cam là
\choice
{$0{,}5\text{ l}$}
{$1\text{ l}$}
{$0{,}65\text{ l}$}
{\True $0{,}75\text{ l}$}
\loigiai{
Số lít nước cam cả ba bạn uống là: $0{,}25 \times 3 = 0{,}75\text{ (l)}$.\\
\textbf{Đáp án của câu hỏi là D}
}
\end{ex}

\begin{ex} %[5D4V21]
Mẹ của Mai mua $3\text{ kg}$ xoài hết $75\,000$ đồng. Cô Hà mua $3{,}5\text{ kg}$ xoài cùng loại và đưa cho người bán hàng tờ tiền $100\,000$ đồng. Người bán hàng phải trả lại cô Hà số tiền là
\choice
{\True $12\,500$ đồng}
{$15\,000$ đồng}
{$17\,500$ đồng}
{$87\,500$ đồng}
\loigiai{
Giá tiền $1\text{ kg}$ xoài là: $75\,000 : 3 = 25\,000$ (đồng).\\
Số tiền cô Hà mua $3{,}5\text{ kg}$ xoài là: $25\,000 \times 3{,}5 = 87\,500$ (đồng).\\
Số tiền trả lại là: $100\,000 - 87\,500 = 12\,500$ (đồng).\\
\textbf{Đáp án của câu hỏi là A}
}
\end{ex}

\textbf{Phần 2. Trắc nghiệm đúng sai}
\setcounterpageref{ex}{0}

\begin{ex} %[5D4H21]
Cho biết $3{,}6 \times 2{,}4 = 8{,}64$. Xét tính đúng sai của các khẳng định sau:
\choiceTF
{\True Phép tính $3{,}6 \times 24 = 86{,}4$.}
{Phép tính $36 \times 0{,}24 = 86{,}4$.}
{\True Phép tính $0{,}36 \times 2{,}4 = 0{,}864$.}
{\True Biểu thức $(5{,}3 \times 1{,}5) \times 2 = 5{,}3 \times (1{,}5 \times 2) = 15{,}9$.}
\loigiai{
\begin{itemize}
    \item $3{,}6 \times 24 = 86{,}4$ (thừa số thứ hai tăng 10 lần $\Rightarrow$ tích tăng 10 lần).
    \item $36 \times 0{,}24 = 8{,}64 \neq 86{,}4$.
    \item $0{,}36 \times 2{,}4 = 0{,}864$ (tích có 3 chữ số thập phân).
    \item $(5{,}3 \times 1{,}5) \times 2 = 5{,}3 \times 3 = 15{,}9$.
\end{itemize}
\begin{itemize}
    \item a) \textbf{Đúng}.
    \item b) \textbf{Sai}: Kết quả là $8{,}64$.
    \item c) \textbf{Đúng}.
    \item d) \textbf{Đúng}.
\end{itemize}
}
\end{ex}

\begin{ex} %[5D4V21]
Một căn phòng hình chữ nhật có chiều dài $4{,}3\text{ m}$ và chiều rộng $3{,}6\text{ m}$. Bác thợ dự định dùng loại sơn phủ sàn có giá $120\,000$ đồng/$1\text{ m}^2$. Xét tính đúng sai của các khẳng định sau:
\choiceTF
{\True Diện tích của căn phòng là $15{,}48\text{ m}^2$.}
{Chu vi của căn phòng là $7{,}9\text{ m}$.}
{\True Làm tròn diện tích căn phòng đến hàng đơn vị ta được $15\text{ m}^2$.}
{\True Chi phí tiền sơn sàn phòng là $1\,857\,600$ đồng.}
\loigiai{
\begin{itemize}
    \item Diện tích phòng: $4{,}3 \times 3{,}6 = 15{,}48\text{ m}^2$.
    \item Chu vi phòng: $(4{,}3 + 3{,}6) \times 2 = 7{,}9 \times 2 = 15{,}8\text{ m} \neq 7{,}9\text{ m}$.
    \item Làm tròn $15{,}48$ đến hàng đơn vị (hàng phần mười là $4 < 5$) được $15\text{ m}^2$.
    \item Tiền sơn: $15{,}48 \times 120\,000 = 1\,857\,600$ đồng.
\end{itemize}
\begin{itemize}
    \item a) \textbf{Đúng}.
    \item b) \textbf{Sai}: Chu vi là $15{,}8\text{ m}$.
    \item c) \textbf{Đúng}.
    \item d) \textbf{Đúng}.
\end{itemize}
}
\end{ex}

\textbf{Phần 3. Trả lời ngắn}
\setcounterpageref{ex}{0}

\begin{ex} %[5D4H21]
Tính kết quả phép nhân: $5{,}4 \times 39$. (Nhập kết quả dưới dạng số thập phân).
\shortans{210,6}
\loigiai{
\textbf{Đáp số:} $210{,}6$\\
$5{,}4 \times 39 = 210{,}6$.
}
\end{ex}

\begin{ex} %[5D4H21]
Tính kết quả phép nhân: $21{,}9 \times 5{,}1$. (Nhập kết quả dưới dạng số thập phân).
\shortans{111,69}
\loigiai{
\textbf{Đáp số:} $111{,}69$\\
$21{,}9 \times 5{,}1 = 111{,}69$.
}
\end{ex}

\begin{ex} %[5D4H21]
Tính giá trị biểu thức: $2{,}5 \times 8{,}6 \times 4$.
\shortans{86}
\loigiai{
\textbf{Đáp số:} $86$\\
$(2{,}5 \times 4) \times 8{,}6 = 10 \times 8{,}6 = 86$.
}
\end{ex}

\begin{ex} %[5D4V21]
Một ô tô đi quãng đường với vận tốc $84{,}5\text{ km/h}$. Hỏi sau $1{,}2$ giờ ô tô đi được bao nhiêu ki-lô-mét? (Nhập kết quả dưới dạng số thập phân).
\shortans{101,4}
\loigiai{
\textbf{Đáp số:} $101{,}4$\\
$84{,}5 \times 1{,}2 = 101{,}4\text{ (km)}$.
}
\end{ex}

\begin{ex} %[5D4V21]
Biết $1\text{ kg}$ thóc xay xát được $0{,}64\text{ kg}$ gạo. Hỏi $50\text{ kg}$ thóc xay xát được bao nhiêu ki-lô-gam gạo?
\shortans{32}
\loigiai{
\textbf{Đáp số:} $32$\\
$50 \times 0{,}64 = 32\text{ (kg)}$.
}
\end{ex}

\begin{ex} %[5D4V21]
Cô Hà mua $3{,}5\text{ kg}$ xoài giá $25\,000$ đồng/$1\text{ kg}$ và đưa tờ tiền $100\,000$ đồng. Hỏi người bán hàng phải trả lại bao nhiêu nghìn đồng? (Nhập kết quả dưới dạng số thập phân).
\shortans{12,5}
\loigiai{
\textbf{Đáp số:} $12{,}5$\\
Số tiền mua: $3{,}5 \times 25\,000 = 87\,500$ đồng. Số tiền trả lại: $100\,000 - 87\,500 = 12\,500$ đồng $= 12{,}5$ nghìn đồng.
}
\end{ex}

\newpage

\TaoPhieuKiemTra{Chủ đề 4 -- Các phép tính với số thập phân}{22}{Phép chia số thập phân}{https://loptoanthaylinh.vn/thi-online/toan-5/bai-22}
\textbf{Phần 1. Trắc nghiệm lựa chọn}
\setcounterpageref{ex}{0}

\begin{ex} %[5D4N22]
Kết quả của phép tính $0{,}36 : 9$ là
\choice
{\True $0{,}04$}
{$0{,}4$}
{$4$}
{$0{,}004$}
\loigiai{
Ta có: $0{,}36 : 9 = 0{,}04$.\\
\textbf{Đáp án của câu hỏi là A}
}
\end{ex}

\begin{ex} %[5D4N22]
Kết quả của phép tính $26 : 8$ là
\choice
{$3{,}2$}
{\True $3{,}25$}
{$3{,}5$}
{$3{,}125$}
\loigiai{
Thực hiện chia $26 : 8 = 3{,}25$.\\
\textbf{Đáp án của câu hỏi là B}
}
\end{ex}

\begin{ex} %[5D4H22]
Kết quả của phép tính $19{,}95 : 19$ là
\choice
{$1{,}5$}
{$1{,}005$}
{$10{,}5$}
{\True $1{,}05$}
\loigiai{
Ta có: $19 : 19 = 1$; $9 : 19 = 0$; $95 : 19 = 5$. Kết quả là $1{,}05$.\\
\textbf{Đáp án của câu hỏi là D}
}
\end{ex}

\begin{ex} %[5D4H22]
Kết quả của phép tính $57 : 9{,}5$ là
\choice
{\True $6$}
{$0{,}6$}
{$60$}
{$600$}
\loigiai{
Chuyển về chia hai số tự nhiên: $570 : 95 = 6$.\\
\textbf{Đáp án của câu hỏi là A}
}
\end{ex}

\begin{ex} %[5D4H22]
Kết quả của phép tính $128 : 0{,}25$ là
\choice
{$32$}
{$51{,}2$}
{\True $512$}
{$5\,120$}
\loigiai{
Ta có: $128 : 0{,}25 = 12\,800 : 25 = 512$ (hoặc $128 \times 4 = 512$).\\
\textbf{Đáp án của câu hỏi là C}
}
\end{ex}

\begin{ex} %[5D4H22]
Kết quả của phép tính $48{,}3 : 3{,}5$ là
\choice
{$13{,}5$}
{\True $13{,}8$}
{$138$}
{$1{,}38$}
\loigiai{
Ta có: $48{,}3 : 3{,}5 = 483 : 35 = 13{,}8$.\\
\textbf{Đáp án của câu hỏi là B}
}
\end{ex}

\begin{ex} %[5D4H22]
Cho biết $7\,657 : 31 = 247$. Không thực hiện phép tính, giá trị của $76{,}57 : 31$ là
\choice
{$24{,}7$}
{$0{,}247$}
{\True $2{,}47$}
{$247$}
\loigiai{
Số bị chia giảm đi $100$ lần thì thương cũng giảm đi $100$ lần: $247 : 100 = 2{,}47$.\\
\textbf{Đáp án của câu hỏi là C}
}
\end{ex}

\begin{ex} %[5D4H22]
Một con tàu vũ trụ trong $2$ giây bay được $103\text{ km}$. Vận tốc trung bình của con tàu là
\choice
{$50{,}5\text{ km/giây}$}
{$52{,}5\text{ km/giây}$}
{$206\text{ km/giây}$}
{\True $51{,}5\text{ km/giây}$}
\loigiai{
Vận tốc trung bình của tàu là: $103 : 2 = 51{,}5\text{ (km/giây)}$.\\
\textbf{Đáp án của câu hỏi là D}
}
\end{ex}

\begin{ex} %[5D4V22]
Một nhà khảo cổ đã đổ hết $15\text{ l}$ dầu vào các rô-bốt chuột chũi, mỗi rô-bốt chứa được $0{,}75\text{ l}$ dầu. Hỏi có tất cả bao nhiêu rô-bốt chuột chũi?
\choice
{\True $20$ rô-bốt}
{$15$ rô-bốt}
{$25$ rô-bốt}
{$30$ rô-bốt}
\loigiai{
Số rô-bốt chuột chũi là: $15 : 0{,}75 = 1\,500 : 75 = 20$ (rô-bốt).\\
\textbf{Đáp án của câu hỏi là A}
}
\end{ex}

\begin{ex} %[5D4V22]
Trên $1{,}5\text{ ha}$ đất người ta thu hoạch được $3\text{ tấn}$ hạt điều thô. Hỏi trung bình trên mỗi héc-ta đất đó thu hoạch được bao nhiêu tấn hạt điều thô?
\choice
{$1{,}5\text{ tấn}$}
{\True $2\text{ tấn}$}
{$4{,}5\text{ tấn}$}
{$0{,}5\text{ tấn}$}
\loigiai{
Trên mỗi héc-ta thu hoạch được: $3 : 1{,}5 = 2\text{ (tấn)}$.\\
\textbf{Đáp án của câu hỏi là B}
}
\end{ex}

\begin{ex} %[5D4V22]
Một chú rồng nhỏ $4$ chiếc răng sâu và phải trả cho nha sĩ tất cả $15{,}4\text{ kg}$ kẹo. Biết số kẹo trả cho mỗi chiếc răng là như nhau. Hỏi nhổ mỗi chiếc răng sâu chú rồng phải trả bao nhiêu ki-lô-gam kẹo?
\choice
{$3{,}6\text{ kg}$}
{$3{,}8\text{ kg}$}
{$3{,}95\text{ kg}$}
{\True $3{,}85\text{ kg}$}
\loigiai{
Số kẹo cho mỗi chiếc răng sâu là: $15{,}4 : 4 = 3{,}85\text{ (kg)}$.\\
\textbf{Đáp án của câu hỏi là D}
}
\end{ex}

\begin{ex} %[5D4V22]
Mặt sàn một nhà kính trồng rau dạng hình chữ nhật có chiều dài $51{,}5\text{ m}$ và chiều rộng bằng $\dfrac{2}{5}$ chiều dài. Chu vi mặt sàn nhà kính đó là
\choice
{$72{,}1\text{ m}$}
{$103\text{ m}$}
{\True $144{,}2\text{ m}$}
{$144\text{ m}$}
\loigiai{
Chiều rộng mặt sàn là: $51{,}5 \times \dfrac{2}{5} = 20{,}6\text{ (m)}$.\\
Chu vi mặt sàn là: $(51{,}5 + 20{,}6) \times 2 = 72{,}1 \times 2 = 144{,}2\text{ (m)}$.\\
\textbf{Đáp án của câu hỏi là C}
}
\end{ex}

\textbf{Phần 2. Trắc nghiệm đúng sai}
\setcounterpageref{ex}{0}

\begin{ex} %[5D4H22]
Xét tính đúng sai của các phép chia sau:
\choiceTF
{\True Phép chia $15 : 8 = 1{,}875$.}
{Phép chia $6 : 25 = 2{,}4$.}
{\True Phép chia $702 : 7{,}2 = 97{,}5$.}
{\True Phép chia $5{,}4 : 0{,}25 = 21{,}6$.}
\loigiai{
\begin{itemize}
    \item $15 : 8 = 1{,}875$.
    \item $6 : 25 = 0{,}24 \neq 2{,}4$.
    \item $702 : 7{,}2 = 7\,020 : 72 = 97{,}5$.
    \item $5{,}4 : 0{,}25 = 540 : 25 = 21{,}6$.
\end{itemize}
\begin{itemize}
    \item a) \textbf{Đúng}.
    \item b) \textbf{Sai}: $6 : 25 = 0{,}24$.
    \item c) \textbf{Đúng}.
    \item d) \textbf{Đúng}.
\end{itemize}
}
\end{ex}

\begin{ex} %[5D4V22]
Rô-bốt chia đều $9{,}68\text{ yến}$ cá vào $8$ khay cá bằng nhau. Xét tính đúng sai của các khẳng định sau:
\choiceTF
{\True Mỗi khay cá đựng $1{,}21\text{ yến}$ cá.}
{Đổi $1{,}21\text{ yến} = 121\text{ kg}$.}
{\True Khối lượng cá đựng trong $5$ khay là $6{,}05\text{ yến}$.}
{Cần dùng ít nhất $10$ khay nếu mỗi khay chỉ đựng đúng $1\text{ yến}$ cá.}
\loigiai{
\begin{itemize}
    \item Khối lượng mỗi khay: $9{,}68 : 8 = 1{,}21\text{ yến}$.
    \item $1{,}21\text{ yến} = 1{,}21 \times 10 = 12{,}1\text{ kg} \neq 121\text{ kg}$.
    \item $5$ khay đựng: $5 \times 1{,}21 = 6{,}05\text{ yến}$.
    \item $9{,}68 : 1 = 9{,}68 \Rightarrow$ cần $10$ khay (đúng, vì $9$ khay chỉ chứa được $9\text{ yến} < 9{,}68\text{ yến}$).
\end{itemize}
\begin{itemize}
    \item a) \textbf{Đúng}.
    \item b) \textbf{Sai}: $1{,}21\text{ yến} = 12{,}1\text{ kg}$.
    \item c) \textbf{Đúng}.
    \item d) \textbf{Đúng}.
\end{itemize}
}
\end{ex}

\textbf{Phần 3. Trả lời ngắn}
\setcounterpageref{ex}{0}

\begin{ex} %[5D4H22]
Tính thương của phép chia: $95{,}2 : 68$. (Nhập kết quả dưới dạng số thập phân).
\shortans{1,4}
\loigiai{
\textbf{Đáp số:} $1{,}4$\\
$95{,}2 : 68 = 1{,}4$.
}
\end{ex}

\begin{ex} %[5D4H22]
Tính thương của phép chia: $54 : 0{,}45$.
\shortans{120}
\loigiai{
\textbf{Đáp số:} $120$\\
$54 : 0{,}45 = 5\,400 : 45 = 120$.
}
\end{ex}

\begin{ex} %[5D4H22]
Tính thương của phép chia: $144{,}55 : 3{,}5$. (Nhập kết quả dưới dạng số thập phân).
\shortans{41,3}
\loigiai{
\textbf{Đáp số:} $41{,}3$\\
$144{,}55 : 3{,}5 = 1\,445{,}5 : 35 = 41{,}3$.
}
\end{ex}

\begin{ex} %[5D4V22]
Một thùng có $15\text{ l}$ dầu. Người ta rót vào các chai nhỏ, mỗi chai chứa $0{,}75\text{ l}$. Hỏi rót được tất cả bao nhiêu chai dầu?
\shortans{20}
\loigiai{
\textbf{Đáp số:} $20$\\
$15 : 0{,}75 = 20$ (chai).
}
\end{ex}

\begin{ex} %[5D4V22]
Nhổ $4$ chiếc răng sâu hết $15{,}4\text{ kg}$ kẹo. Hỏi nhổ $1$ chiếc răng sâu như vậy hết bao nhiêu ki-lô-gam kẹo? (Nhập kết quả dưới dạng số thập phân).
\shortans{3,85}
\loigiai{
\textbf{Đáp số:} $3{,}85$\\
$15{,}4 : 4 = 3{,}85\text{ (kg)}$.
}
\end{ex}

\begin{ex} %[5D4V22]
Mặt sàn hình chữ nhật dài $51{,}5\text{ m}$, rộng bằng $\dfrac{2}{5}$ chiều dài. Tính chu vi mặt sàn đó theo mét. (Nhập kết quả dưới dạng số thập phân).
\shortans{144,2}
\loigiai{
\textbf{Đáp số:} $144{,}2$\\
Rộng: $51{,}5 \times \dfrac{2}{5} = 20{,}6\text{ m}$. Chu vi: $(51{,}5 + 20{,}6) \times 2 = 144{,}2\text{ m}$.
}
\end{ex}

\newpage

\TaoPhieuKiemTra{Chủ đề 4 -- Các phép tính với số thập phân}{23}{Nhân, chia số thập phân với 10, 100, 0.1, 0.01...}{https://loptoanthaylinh.vn/thi-online/toan-5/bai-23}
\textbf{Phần 1. Trắc nghiệm lựa chọn}
\setcounterpageref{ex}{0}

\begin{ex} %[5D4N23]
Khi nhân một số thập phân với $100$, ta chỉ việc chuyển dấu phẩy của số đó sang bên phải
\choice
{$1$ chữ số}
{\True $2$ chữ số}
{$3$ chữ số}
{$4$ chữ số}
\loigiai{
Muốn nhân một số thập phân với $100$, ta chuyển dấu phẩy của số đó sang bên phải hai chữ số.\\
\textbf{Đáp án của câu hỏi là B}
}
\end{ex}

\begin{ex} %[5D4N23]
Khi chia một số thập phân cho $0{,}1$, ta chỉ việc
\choice
{Chuyển dấu phẩy sang bên trái $1$ chữ số}
{Chuyển dấu phẩy sang bên trái $2$ chữ số}
{\True Chuyển dấu phẩy sang bên phải $1$ chữ số}
{Chuyển dấu phẩy sang bên phải $2$ chữ số}
\loigiai{
Chia cho $0{,}1$ tương đương với nhân với $10$, ta chuyển dấu phẩy sang bên phải một chữ số.\\
\textbf{Đáp án của câu hỏi là C}
}
\end{ex}

\begin{ex} %[5D4H23]
Kết quả của phép tính $2{,}08 \times 1\,000$ là
\choice
{$20{,}8$}
{\True $2\,080$}
{$208$}
{$20\,800$}
\loigiai{
Chuyển dấu phẩy sang phải ba chữ số: $2{,}08 \times 1\,000 = 2\,080$.\\
\textbf{Đáp án của câu hỏi là B}
}
\end{ex}

\begin{ex} %[5D4H23]
Kết quả của phép tính $15{,}31 \times 0{,}001$ là
\choice
{$15\,310$}
{$1{,}531$}
{$0{,}1531$}
{\True $0{,}01531$}
\loigiai{
Chuyển dấu phẩy sang trái ba chữ số: $15{,}31 \times 0{,}001 = 0{,}01531$.\\
\textbf{Đáp án của câu hỏi là D}
}
\end{ex}

\begin{ex} %[5D4H23]
Kết quả của phép tính $0{,}8 : 100$ là
\choice
{$80$}
{\True $0{,}008$}
{$0{,}08$}
{$0{,}0008$}
\loigiai{
Chuyển dấu phẩy sang trái hai chữ số: $0{,}8 : 100 = 0{,}008$.\\
\textbf{Đáp án của câu hỏi là B}
}
\end{ex}

\begin{ex} %[5D4H23]
Kết quả của phép tính $6{,}3 : 0{,}01$ là
\choice
{\True $630$}
{$0{,}063$}
{$0{,}63$}
{$63$}
\loigiai{
Chuyển dấu phẩy sang phải hai chữ số: $6{,}3 : 0{,}01 = 630$.\\
\textbf{Đáp án của câu hỏi là A}
}
\end{ex}

\begin{ex} %[5D4H23]
Đổi $23{,}45\text{ kg}$ ra gam ta được kết quả là
\choice
{$234{,}5\text{ g}$}
{$2\,345\text{ g}$}
{$234\,500\text{ g}$}
{\True $23\,450\text{ g}$}
\loigiai{
Vì $1\text{ kg} = 1\,000\text{ g}$ nên $23{,}45\text{ kg} = 23{,}45 \times 1\,000 = 23\,450\text{ g}$.\\
\textbf{Đáp án của câu hỏi là D}
}
\end{ex}

\begin{ex} %[5D4H23]
Đổi $12{,}6\text{ cm}$ ra mét ta được kết quả là
\choice
{\True $0{,}126\text{ m}$}
{$1{,}26\text{ m}$}
{$0{,}0126\text{ m}$}
{$126\text{ m}$}
\loigiai{
Vì $1\text{ m} = 100\text{ cm}$ nên $12{,}6\text{ cm} = 12{,}6 : 100 = 0{,}126\text{ m}$.\\
\textbf{Đáp án của câu hỏi là A}
}
\end{ex}

\begin{ex} %[5D4H23]
Mỗi tờ giấy dày $0{,}1\text{ mm}$. Cả chồng giấy dày $36{,}5\text{ mm}$. Số tờ giấy trong chồng đó là
\choice
{$36$ tờ}
{$3\,650$ tờ}
{\True $365$ tờ}
{$3{,}65$ tờ}
\loigiai{
Số tờ giấy là: $36{,}5 : 0{,}1 = 365$ (tờ).\\
\textbf{Đáp án của câu hỏi là C}
}
\end{ex}

\begin{ex} %[5D4V23]
Một đàn gấu Bắc Cực có $10$ chú gấu con, mỗi chú gấu con ăn hết $4{,}5\text{ kg}$ cá mỗi ngày. Hỏi cả $10$ chú gấu con ăn hết bao nhiêu ki-lô-gam cá mỗi ngày?
\choice
{\True $45\text{ kg}$}
{$450\text{ kg}$}
{$4{,}5\text{ kg}$}
{$40{,}5\text{ kg}$}
\loigiai{
Cả $10$ chú gấu con ăn hết: $4{,}5 \times 10 = 45\text{ (kg cá)}$.\\
\textbf{Đáp án của câu hỏi là A}
}
\end{ex}

\begin{ex} %[5D4V23]
Chiếc máy xúc lớn nhất thế giới xúc $10$ lần được $937{,}8\text{ tấn}$ quặng. Trung bình mỗi lần máy xúc đó xúc được số tấn quặng là
\choice
{$9\,378\text{ tấn}$}
{$9{,}378\text{ tấn}$}
{\True $93{,}78\text{ tấn}$}
{$937{,}8\text{ tấn}$}
\loigiai{
Mỗi lần máy xúc xúc được: $937{,}8 : 10 = 93{,}78\text{ (tấn)}$.\\
\textbf{Đáp án của câu hỏi là C}
}
\end{ex}

\begin{ex} %[5D4V23]
Chia đều $0{,}3\text{ l}$ giấm vào $10$ ống nghiệm như nhau. Mỗi ống nghiệm chứa số lít giấm là
\choice
{$3\text{ l}$}
{$0{,}3\text{ l}$}
{$0{,}003\text{ l}$}
{\True $0{,}03\text{ l}$}
\loigiai{
Mỗi ống nghiệm chứa: $0{,}3 : 10 = 0{,}03\text{ (l)}$.\\
\textbf{Đáp án của câu hỏi là D}
}
\end{ex}

\textbf{Phần 2. Trắc nghiệm đúng sai}
\setcounterpageref{ex}{0}

\begin{ex} %[5D4H23]
Xét tính đúng sai của các khẳng định sau:
\choiceTF
{\True $79{,}6 \times 0{,}01 = 0{,}796$.}
{$58 : 0{,}001 = 0{,}058$.}
{\True $6{,}89\text{ l} = 6\,890\text{ ml}$.}
{$3{,}7\text{ m} = 0{,}037\text{ km}$.}
\loigiai{
\begin{itemize}
    \item $79{,}6 \times 0{,}01 = 0{,}796$.
    \item $58 : 0{,}001 = 58 \times 1\,000 = 58\,000 \neq 0{,}058$.
    \item $6{,}89\text{ l} = 6{,}89 \times 1\,000 = 6\,890\text{ ml}$.
    \item $3{,}7\text{ m} = 3{,}7 : 1\,000 = 0{,}0037\text{ km} \neq 0{,}037\text{ km}$.
\end{itemize}
\begin{itemize}
    \item a) \textbf{Đúng}.
    \item b) \textbf{Sai}: Kết quả là $58\,000$.
    \item c) \textbf{Đúng}.
    \item d) \textbf{Sai}: $3{,}7\text{ m} = 0{,}0037\text{ km}$.
\end{itemize}
}
\end{ex}

\begin{ex} %[5D4V23]
Cho dãy biến đổi tính toán: $A = 10{,}8 \xrightarrow{: 100} B \xrightarrow{: 0{,}01} C \xrightarrow{\times 0{,}01} D \xrightarrow{\times 100} E$. Xét tính đúng sai của các khẳng định sau:
\choiceTF
{\True Giá trị của $B$ là $0{,}108$.}
{Giá trị của $C$ là $1{,}08$.}
{\True Giá trị của $D$ bằng giá trị của $B$.}
{\True Giá trị cuối cùng $E$ bằng $10{,}8$ (bằng với số ban đầu $A$).}
\loigiai{
\begin{itemize}
    \item $B = 10{,}8 : 100 = 0{,}108$.
    \item $C = 0{,}108 : 0{,}01 = 10{,}8 \neq 1{,}08$.
    \item $D = 10{,}8 \times 0{,}01 = 0{,}108 = B$.
    \item $E = 0{,}108 \times 100 = 10{,}8 = A$.
\end{itemize}
\begin{itemize}
    \item a) \textbf{Đúng}.
    \item b) \textbf{Sai}: $C = 10{,}8$.
    \item c) \textbf{Đúng}.
    \item d) \textbf{Đúng}.
\end{itemize}
}
\end{ex}

\textbf{Phần 3. Trả lời ngắn}
\setcounterpageref{ex}{0}

\begin{ex} %[5D4H23]
Tính kết quả phép tính: $0{,}94 \times 100$.
\shortans{94}
\loigiai{
\textbf{Đáp số:} $94$\\
$0{,}94 \times 100 = 94$.
}
\end{ex}

\begin{ex} %[5D4H23]
Tính kết quả phép tính: $8{,}2 \times 0{,}01$. (Nhập kết quả dưới dạng số thập phân).
\shortans{0,082}
\loigiai{
\textbf{Đáp số:} $0{,}082$\\
$8{,}2 \times 0{,}01 = 0{,}082$.
}
\end{ex}

\begin{ex} %[5D4H23]
Tính kết quả phép tính: $341 : 1\,000$. (Nhập kết quả dưới dạng số thập phân).
\shortans{0,341}
\loigiai{
\textbf{Đáp số:} $0{,}341$\\
$341 : 1\,000 = 0{,}341$.
}
\end{ex}

\begin{ex} %[5D4H23]
Đổi $23{,}4\text{ tạ}$ ra yến.
\shortans{234}
\loigiai{
\textbf{Đáp số:} $234$\\
$23{,}4\text{ tạ} = 23{,}4 \times 10 = 234\text{ yến}$.
}
\end{ex}

\begin{ex} %[5D4V23]
Một tấm bạt hình chữ nhật có chiều dài $27{,}86\text{ m}$ và chiều rộng $10\text{ m}$. Tính diện tích tấm bạt đó theo đơn vị mét vuông. (Nhập kết quả dưới dạng số thập phân).
\shortans{278,6}
\loigiai{
\textbf{Đáp số:} $278{,}6$\\
$27{,}86 \times 10 = 278{,}6\text{ (m}^2\text{)}$.
}
\end{ex}

\begin{ex} %[5D4V23]
Một thanh sắt dài $1\text{ m}$ nặng $15{,}23\text{ kg}$. Hỏi thanh sắt cùng loại dài $0{,}1\text{ m}$ nặng bao nhiêu ki-lô-gam? (Nhập kết quả dưới dạng số thập phân).
\shortans{1,523}
\loigiai{
\textbf{Đáp số:} $1{,}523$\\
$15{,}23 \times 0{,}1 = 1{,}523\text{ (kg)}$.
}
\end{ex}

\newpage

\TaoPhieuKiemTra{Chủ đề 4 -- Các phép tính với số thập phân}{24}{Luyện tập chung phép tính số thập phân}{https://loptoanthaylinh.vn/thi-online/toan-5/bai-24}
\textbf{Phần 1. Trắc nghiệm lựa chọn}
\setcounterpageref{ex}{0}

\begin{ex} %[5D4N24]
Kết quả của phép tính $52 - 4{,}91$ là
\choice
{\True $47{,}09$}
{$47{,}19$}
{$48{,}09$}
{$47{,}91$}
\loigiai{
Ta có: $52{,}00 - 4{,}91 = 47{,}09$.\\
\textbf{Đáp án của câu hỏi là A}
}
\end{ex}

\begin{ex} %[5D4N24]
Kết quả của phép tính $0{,}73 \times 6{,}5$ là
\choice
{$47{,}45$}
{$0{,}4745$}
{$4{,}755$}
{\True $4{,}745$}
\loigiai{
Ta có: $0{,}73 \times 6{,}5 = 4{,}745$.\\
\textbf{Đáp án của câu hỏi là D}
}
\end{ex}

\begin{ex} %[5D4H24]
Kết quả của phép tính $104{,}89 : 8{,}5$ là
\choice
{$123{,}4$}
{$1{,}234$}
{\True $12{,}34$}
{$12{,}44$}
\loigiai{
Ta có: $104{,}89 : 8{,}5 = 1\,048{,}9 : 85 = 12{,}34$.\\
\textbf{Đáp án của câu hỏi là C}
}
\end{ex}

\begin{ex} %[5D4H24]
Giá trị của biểu thức $(131{,}4 - 80{,}8) : 2{,}3 + 21{,}64 \times 2$ là
\choice
{\True $65{,}28$}
{$44{,}28$}
{$55{,}28$}
{$75{,}28$}
\loigiai{
Ta có: $50{,}6 : 2{,}3 + 43{,}28 = 22 + 43{,}28 = 65{,}28$.\\
\textbf{Đáp án của câu hỏi là A}
}
\end{ex}

\begin{ex} %[5D4H24]
Tính thuận tiện biểu thức $0{,}56 \times 4{,}7 + 5{,}3 \times 0{,}56$ ta được kết quả là
\choice
{$0{,}56$}
{$56$}
{$560$}
{\True $5{,}6$}
\loigiai{
Ta có: $0{,}56 \times (4{,}7 + 5{,}3) = 0{,}56 \times 10 = 5{,}6$.\\
\textbf{Đáp án của câu hỏi là D}
}
\end{ex}

\begin{ex} %[5D4H24]
Tìm $x$ biết: $51{,}23 + x = 74{,}9$.
\choice
{$23{,}77$}
{$22{,}67$}
{\True $23{,}67$}
{$126{,}13$}
\loigiai{
Ta có: $x = 74{,}90 - 51{,}23 = 23{,}67$.\\
\textbf{Đáp án của câu hỏi là C}
}
\end{ex}

\begin{ex} %[5D4H24]
Lượng mưa đo được trong $3$ lần của tháng Hai lần lượt là $15{,}6\text{ mm}$, $18{,}9\text{ mm}$ và $12{,}3\text{ mm}$. Lượng mưa trung bình của tháng Hai là
\choice
{$14{,}6\text{ mm}$}
{\True $15{,}6\text{ mm}$}
{$15{,}3\text{ mm}$}
{$16{,}2\text{ mm}$}
\loigiai{
Lượng mưa trung bình là: $(15{,}6 + 18{,}9 + 12{,}3) : 3 = 46{,}8 : 3 = 15{,}6\text{ (mm)}$.\\
\textbf{Đáp án của câu hỏi là B}
}
\end{ex}

\begin{ex} %[5D4V24]
Dúi làm rơi khúc mía xuống một cái lỗ sâu $1{,}31\text{ m}$ ($1{,}06\text{ m} + 0{,}25\text{ m}$). Chồn đã đào một cái hố sâu $0{,}9\text{ m}$ bên cạnh để lấy giúp. Hỏi chồn còn phải đào sâu thêm bao nhiêu mét nữa?
\choice
{$0{,}31\text{ m}$}
{$0{,}51\text{ m}$}
{\True $0{,}41\text{ m}$}
{$2{,}21\text{ m}$}
\loigiai{
Độ sâu cần đào thêm là: $1{,}31 - 0{,}9 = 0{,}41\text{ (m)}$.\\
\textbf{Đáp án của câu hỏi là C}
}
\end{ex}

\begin{ex} %[5D4V24]
Một tấm bản đồ hình chữ nhật có chiều dài $8{,}4\text{ dm}$ và chiều rộng $6{,}5\text{ dm}$. Sau khi gấp đôi $2$ lần (thành $4$ phần bằng nhau) ta được một hình chữ nhật nhỏ có diện tích là
\choice
{$54{,}6\text{ dm}^2$}
{\True $13{,}65\text{ dm}^2$}
{$27{,}3\text{ dm}^2$}
{$13{,}5\text{ dm}^2$}
\loigiai{
Diện tích bản đồ là: $8{,}4 \times 6{,}5 = 54{,}6\text{ (dm}^2\text{)}$.\\
Diện tích hình chữ nhật nhỏ sau $2$ lần gấp đôi là: $54{,}6 : 4 = 13{,}65\text{ (dm}^2\text{)}$.\\
\textbf{Đáp án của câu hỏi là B}
}
\end{ex}

\begin{ex} %[5D4V24]
Trong hai năm, trang trại bán được tất cả $21{,}56\text{ tấn}$ cá chẽm. Biết năm thứ hai bán được nhiều hơn năm thứ nhất $2{,}7\text{ tấn}$. Số tấn cá chẽm trang trại bán được trong năm thứ nhất là
\choice
{\True $9{,}43\text{ tấn}$}
{$12{,}13\text{ tấn}$}
{$10{,}78\text{ tấn}$}
{$8{,}43\text{ tấn}$}
\loigiai{
Năm thứ nhất bán được: $(21{,}56 - 2{,}7) : 2 = 18{,}86 : 2 = 9{,}43\text{ (tấn)}$.\\
\textbf{Đáp án của câu hỏi là A}
}
\end{ex}

\begin{ex} %[5D4V24]
Khu vườn trồng cây ăn quả dạng hình chữ nhật có diện tích $83{,}52\text{ m}^2$ và chiều rộng $8{,}7\text{ m}$. Chu vi khu vườn đó là
\choice
{$18{,}3\text{ m}$}
{\True $36{,}6\text{ m}$}
{$34{,}8\text{ m}$}
{$73{,}2\text{ m}$}
\loigiai{
Chiều dài khu vườn là: $83{,}52 : 8,7 = 9{,}6\text{ (m)}$.\\
Chu vi khu vườn là: $(9{,}6 + 8{,}7) \times 2 = 18{,}3 \times 2 = 36{,}6\text{ (m)}$.\\
\textbf{Đáp án của câu hỏi là B}
}
\end{ex}

\begin{ex} %[5D4V24]
Bạn Tú có $25$ quả bóng, mỗi quả nặng $5{,}7\text{ g}$. Tú bỏ tất cả số bóng vào hai hộp xanh và đỏ. Tú thấy cân nặng của hộp xanh tăng thêm $62{,}7\text{ g}$. Hỏi cân nặng của hộp đỏ tăng thêm bao nhiêu gam?
\choice
{$142{,}5\text{ g}$}
{$74{,}1\text{ g}$}
{$80{,}2\text{ g}$}
{\True $79{,}8\text{ g}$}
\loigiai{
Tổng khối lượng $25$ quả bóng là: $25 \times 5{,}7 = 142{,}5\text{ (g)}$.\\
Cân nặng hộp đỏ tăng thêm là: $142{,}5 - 62{,}7 = 79{,}8\text{ (g)}$.\\
\textbf{Đáp án của câu hỏi là D}
}
\end{ex}

\textbf{Phần 2. Trắc nghiệm đúng sai}
\setcounterpageref{ex}{0}

\begin{ex} %[5D4H24]
Xét tính đúng sai của các khẳng định sau:
\choiceTF
{\True Phép tính $4{,}98 + 15{,}7 = 20{,}68$.}
{Phép tính $314 : 0{,}001 = 0{,}314$.}
{\True Biểu thức $(5{,}3 + 4{,}7) \times 0{,}1 = 5{,}3 \times 0{,}1 + 4{,}7 \times 0{,}1 = 1$.}
{\True Giá trị biểu thức $8{,}16 : (1{,}32 + 3{,}48) - 0{,}34 : 2 = 1{,}53$.}
\loigiai{
\begin{itemize}
    \item $4{,}98 + 15{,}70 = 20{,}68$.
    \item $314 : 0{,}001 = 314 \times 1\,000 = 314\,000 \neq 0{,}314$.
    \item $(5{,}3 + 4{,}7) \times 0{,}1 = 10 \times 0{,}1 = 1$.
    \item $8{,}16 : 4{,}8 - 0{,}17 = 1{,}7 - 0{,}17 = 1{,}53$.
\end{itemize}
\begin{itemize}
    \item a) \textbf{Đúng}.
    \item b) \textbf{Sai}: Kết quả là $314\,000$.
    \item c) \textbf{Đúng}.
    \item d) \textbf{Đúng}.
\end{itemize}
}
\end{ex}

\begin{ex} %[5D4V24]
Một mảnh đất hình chữ nhật có kích thước $40{,}4\text{ m} \times 30\text{ m}$. Người ta sử dụng đúng một nửa diện tích mảnh đất để xây nhà máy. Xét tính đúng sai của các khẳng định sau:
\choiceTF
{\True Diện tích của cả mảnh đất là $1\,212\text{ m}^2$.}
{Diện tích đất dùng để xây nhà máy là $600\text{ m}^2$.}
{\True Diện tích nhà máy bằng $0{,}0606\text{ ha}$.}
{\True Diện tích của nhà máy lớn hơn $\dfrac{1}{100}\text{ ha}$ (tức $100\text{ m}^2$).}
\loigiai{
\begin{itemize}
    \item Diện tích mảnh đất: $40{,}4 \times 30 = 1\,212\text{ m}^2$.
    \item Diện tích xây nhà máy: $1\,212 : 2 = 606\text{ m}^2 \neq 600\text{ m}^2$.
    \item $606\text{ m}^2 = 606 : 10\,000 = 0{,}0606\text{ ha}$.
    \item $\dfrac{1}{100}\text{ ha} = 100\text{ m}^2 < 606\text{ m}^2$.
\end{itemize}
\begin{itemize}
    \item a) \textbf{Đúng}.
    \item b) \textbf{Sai}: Diện tích là $606\text{ m}^2$.
    \item c) \textbf{Đúng}.
    \item d) \textbf{Đúng}.
\end{itemize}
}
\end{ex}

\textbf{Phần 3. Trả lời ngắn}
\setcounterpageref{ex}{0}

\begin{ex} %[5D4H24]
Tính giá trị biểu thức: $(131{,}4 - 80{,}8) : 2{,}3 + 21{,}64 \times 2$. (Nhập kết quả dưới dạng số thập phân).
\shortans{65,28}
\loigiai{
\textbf{Đáp số:} $65{,}28$\\
$(131{,}4 - 80{,}8) : 2{,}3 + 21{,}64 \times 2 = 50{,}6 : 2{,}3 + 43{,}28 = 22 + 43{,}28 = 65{,}28$.
}
\end{ex}

\begin{ex} %[5D4H24]
Tính nhanh giá trị biểu thức: $0{,}56 \times 4{,}7 + 5{,}3 \times 0{,}56$. (Nhập kết quả dưới dạng số thập phân).
\shortans{5,6}
\loigiai{
\textbf{Đáp số:} $5{,}6$\\
$0{,}56 \times (4{,}7 + 5{,}3) = 0{,}56 \times 10 = 5{,}6$.
}
\end{ex}

\begin{ex} %[5D4H24]
Tìm $x$ biết: $x : 3{,}8 = 21{,}34$. (Nhập kết quả dưới dạng số thập phân).
\shortans{81,092}
\loigiai{
\textbf{Đáp số:} $81{,}092$\\
$x = 21{,}34 \times 3{,}8 = 81{,}092$.
}
\end{ex}

\begin{ex} %[5D4V24]
Tấm bản đồ hình chữ nhật có kích thước $8{,}4\text{ dm} \times 6{,}5\text{ dm}$. Tính diện tích tấm bản đồ đó theo đơn vị đề-xi-mét vuông. (Nhập kết quả dưới dạng số thập phân).
\shortans{54,6}
\loigiai{
\textbf{Đáp số:} $54{,}6$\\
$8{,}4 \times 6{,}5 = 54{,}6\text{ (dm}^2\text{)}$.
}
\end{ex}

\begin{ex} %[5D4V24]
Trong hai năm trang trại bán được $21{,}56\text{ tấn}$ cá chẽm. Năm thứ hai bán nhiều hơn năm thứ nhất $2{,}7\text{ tấn}$. Hỏi năm thứ hai trang trại bán được bao nhiêu tấn cá chẽm? (Nhập kết quả dưới dạng số thập phân).
\shortans{12,13}
\loigiai{
\textbf{Đáp số:} $12{,}13$\\
Năm thứ hai bán được: $(21{,}56 + 2{,}7) : 2 = 12{,}13\text{ (tấn)}$.
}
\end{ex}

\begin{ex} %[5D4V24]
Khu vườn hình chữ nhật có diện tích $83{,}52\text{ m}^2$ và chiều rộng $8{,}7\text{ m}$. Tính chu vi khu vườn đó theo mét. (Nhập kết quả dưới dạng số thập phân).
\shortans{36,6}
\loigiai{
\textbf{Đáp số:} $36{,}6$\\
Chiều dài: $83{,}52 : 8{,}7 = 9{,}6\text{ m}$. Chu vi: $(9{,}6 + 8{,}7) \times 2 = 36{,}6\text{ m}$.
}
\end{ex}

\newpage

% =========================================================================
% CHỦ ĐỀ 5 - MỘT SỐ HÌNH PHẲNG. CHU VI VÀ DIỆN TÍCH
% =========================================================================
\section{Một số hình phẳng. chu vi và diện tích}

\TaoPhieuKiemTra{Chủ đề 5 -- Một số hình phẳng. chu vi và diện tích}{25}{Hình tam giác. Diện tích hình tam giác}{https://loptoanthaylinh.vn/thi-online/toan-5/bai-25}
\textbf{Phần 1. Trắc nghiệm lựa chọn}
\setcounterpageref{ex}{0}

\begin{ex} %[5H5N25]
Hình tam giác có ba góc đều là góc nhọn được gọi là
\choice
{\True Hình tam giác nhọn}
{Hình tam giác vuông}
{Hình tam giác tù}
{Hình tam giác đều}
\loigiai{
Hình tam giác có $3$ góc nhọn gọi là hình tam giác nhọn.\\
\textbf{Đáp án của câu hỏi là A}
}
\end{ex}

\begin{ex} %[5H5N25]
Công thức tính diện tích $S$ của hình tam giác có độ dài đáy là $a$ và chiều cao tương ứng là $h$ (cùng đơn vị đo) là
\choice
{$S = a \times h$}
{$S = (a + h) \times 2$}
{$S = \dfrac{a + h}{2}$}
{\True $S = \dfrac{a \times h}{2}$}
\loigiai{
Diện tích hình tam giác bằng độ dài đáy nhân với chiều cao (cùng đơn vị đo) rồi chia cho $2$.\\
\textbf{Đáp án của câu hỏi là D}
}
\end{ex}

\begin{ex} %[5H5H25]
Diện tích của hình tam giác có độ dài đáy $10\text{ cm}$ và chiều cao tương ứng $8\text{ cm}$ là
\choice
{$80\text{ cm}^2$}
{$80\text{ cm}$}
{\True $40\text{ cm}^2$}
{$40\text{ cm}$}
\loigiai{
Diện tích tam giác là: $\dfrac{10 \times 8}{2} = 40\text{ (cm}^2\text{)}$.\\
\textbf{Đáp án của câu hỏi là C}
}
\end{ex}

\begin{ex} %[5H5H25]
Diện tích của hình tam giác có độ dài đáy $5\text{ dm}$ và chiều cao $8\text{ dm}$ là
\choice
{$40\text{ dm}^2$}
{\True $20\text{ dm}^2$}
{$13\text{ dm}^2$}
{$20\text{ m}^2$}
\loigiai{
Diện tích tam giác là: $\dfrac{5 \times 8}{2} = 20\text{ (dm}^2\text{)}$.\\
\textbf{Đáp án của câu hỏi là B}
}
\end{ex}

\begin{ex} %[5H5H25]
Một tấm kính có dạng hình tam giác vuông với hai cạnh góc vuông đều bằng $6\text{ m}$. Diện tích của tấm kính đó là
\choice
{$36\text{ m}^2$}
{$12\text{ m}^2$}
{\True $18\text{ m}^2$}
{$24\text{ m}^2$}
\loigiai{
Diện tích tấm kính là: $\dfrac{6 \times 6}{2} = 18\text{ (m}^2\text{)}$.\\
\textbf{Đáp án của câu hỏi là C}
}
\end{ex}

\begin{ex} %[5H5H25]
Diện tích của hình tam giác có độ dài đáy $2\text{ dm}$ và chiều cao $20\text{ cm}$ là
\choice
{$20\text{ dm}^2$}
{\True $2\text{ dm}^2$}
{$20\text{ cm}^2$}
{$200\text{ dm}^2$}
\loigiai{
Đổi $2\text{ dm} = 20\text{ cm}$. Diện tích là: $\dfrac{20 \times 20}{2} = 200\text{ (cm}^2\text{)} = 2\text{ dm}^2$.\\
\textbf{Đáp án của câu hỏi là B}
}
\end{ex}

\begin{ex} %[5H5H25]
Một hình tam giác có diện tích là $15\text{ cm}^2$ và chiều cao $5\text{ cm}$. Độ dài đáy tương ứng của hình tam giác đó là
\choice
{\True $6\text{ cm}$}
{$3\text{ cm}$}
{$15\text{ cm}$}
{$30\text{ cm}$}
\loigiai{
Độ dài đáy là: $15 \times 2 : 5 = 6\text{ (cm)}$.\\
\textbf{Đáp án của câu hỏi là A}
}
\end{ex}

\begin{ex} %[5H5H25]
Một hình tam giác có diện tích $8\text{ m}^2$ và độ dài đáy $4\text{ m}$. Chiều cao tương ứng của hình tam giác đó là
\choice
{$2\text{ m}$}
{$8\text{ m}$}
{$16\text{ m}$}
{\True $4\text{ m}$}
\loigiai{
Chiều cao của hình tam giác là: $8 \times 2 : 4 = 4\text{ (m)}$.\\
\textbf{Đáp án của câu hỏi là D}
}
\end{ex}

\begin{ex} %[5H5V25]
Tờ giấy hình vuông $ABCD$ có cạnh $12\text{ cm}$. Gọi $M$ là trung điểm của $BC$, $N$ là trung điểm của $CD$. Cắt lấy tam giác vuông $MNC$ làm đuôi con cá. Diện tích hình tam giác $MNC$ bằng
\choice
{$36\text{ cm}^2$}
{\True $18\text{ cm}^2$}
{$27\text{ cm}^2$}
{$72\text{ cm}^2$}
\loigiai{
Độ dài đoạn $MC = NC = 12 : 2 = 6\text{ (cm)}$.\\
Diện tích tam giác $MNC$ là: $\dfrac{6 \times 6}{2} = 18\text{ (cm}^2\text{)}$.\\
\textbf{Đáp án của câu hỏi là B}
}
\end{ex}

\begin{ex} %[5H5V25]
Một thửa ruộng hình tam giác có độ dài đáy $32\text{ m}$ và chiều cao $25\text{ m}$. Cứ mỗi $100\text{ m}^2$ người ta thu hoạch được $70\text{ kg}$ thóc. Số thóc thu hoạch được trên thửa ruộng đó là
\choice
{\True $280\text{ kg}$}
{$28\text{ kg}$}
{$560\text{ kg}$}
{$400\text{ kg}$}
\loigiai{
Diện tích thửa ruộng là: $\dfrac{32 \times 25}{2} = 400\text{ (m}^2\text{)}$.\\
Số thóc thu hoạch được là: $400 : 100 \times 70 = 280\text{ (kg)}$.\\
\textbf{Đáp án của câu hỏi là A}
}
\end{ex}

\begin{ex} %[5H5V25]
Nếu tăng độ dài đáy của một hình tam giác lên $2$ lần và giữ nguyên chiều cao thì diện tích của hình tam giác đó sẽ
\choice
{Không thay đổi}
{Tăng lên $4$ lần}
{Giảm đi $2$ lần}
{\True Tăng lên $2$ lần}
\loigiai{
Diện tích mới: $S' = \dfrac{(2 \times a) \times h}{2} = 2 \times \dfrac{a \times h}{2} = 2 \times S$. Vậy diện tích tăng lên $2$ lần.\\
\textbf{Đáp án của câu hỏi là D}
}
\end{ex}

\begin{ex} %[5H5V25]
Một lá cờ hình tam giác vuông có hai cạnh góc vuông lần lượt là $20\text{ cm}$ và $30\text{ cm}$. Người ta cắt $50$ lá cờ như vậy từ một tấm vải. Tổng diện tích vải của $50$ lá cờ đó theo đơn vị mét vuông là
\choice
{$15\text{ m}^2$}
{$0{,}15\text{ m}^2$}
{\True $1{,}5\text{ m}^2$}
{$150\text{ m}^2$}
\loigiai{
Diện tích $1$ lá cờ là: $\dfrac{20 \times 30}{2} = 300\text{ (cm}^2\text{)}$.\\
Diện tích $50$ lá cờ là: $50 \times 300 = 15\,000\text{ (cm}^2\text{)} = 1{,}5\text{ m}^2$.\\
\textbf{Đáp án của câu hỏi là C}
}
\end{ex}

\textbf{Phần 2. Trắc nghiệm đúng sai}
\setcounterpageref{ex}{0}

\begin{ex} %[5H5H25]
Xét tính đúng sai của các khẳng định sau về hình tam giác:
\choiceTF
{\True Hình tam giác đều có $3$ cạnh bằng nhau và $3$ góc đều bằng $60^\circ$.}
{Trong hình tam giác vuông, hai cạnh góc vuông không thể coi là chiều cao và cạnh đáy.}
{\True Diện tích tam giác có đáy $4\text{ cm}$, chiều cao $3\text{ cm}$ là $6\text{ cm}^2$.}
{Hình tam giác có $1$ góc tù và $2$ góc nhọn gọi là hình tam giác vuông.}
\loigiai{
\begin{itemize}
    \item Hình tam giác đều có 3 cạnh bằng nhau và 3 góc bằng nhau cùng bằng $60^\circ$.
    \item Trong tam giác vuông, hai cạnh góc vuông chính là đáy và chiều cao tương ứng của nhau (khẳng định sai).
    \item $S = \dfrac{4 \times 3}{2} = 6\text{ cm}^2$.
    \item Hình tam giác có 1 góc tù gọi là tam giác tù, không phải tam giác vuông.
\end{itemize}
\begin{itemize}
    \item a) \textbf{Đúng}.
    \item b) \textbf{Sai}.
    \item c) \textbf{Đúng}.
    \item d) \textbf{Sai}.
\end{itemize}
}
\end{ex}

\begin{ex} %[5H5V25]
Một hình tam giác $ABC$ có đáy $BC = 10\text{ dm}$, chiều cao $AH = 6\text{ dm}$. Xét tính đúng sai của các khẳng định sau:
\choiceTF
{\True Diện tích hình tam giác $ABC$ là $30\text{ dm}^2$.}
{Đổi diện tích hình tam giác ra mét vuông là $3\text{ m}^2$.}
{\True Nếu kéo dài đáy $BC$ thêm $2\text{ dm}$ thì diện tích tăng thêm $6\text{ dm}^2$.}
{\True Tỉ số diện tích phần tăng thêm và diện tích ban đầu là $\dfrac{1}{5}$.}
\loigiai{
\begin{itemize}
    \item Diện tích tam giác: $\dfrac{10 \times 6}{2} = 30\text{ dm}^2$.
    \item $30\text{ dm}^2 = 0{,}3\text{ m}^2 \neq 3\text{ m}^2$.
    \item Diện tích tăng thêm: $\dfrac{2 \times 6}{2} = 6\text{ dm}^2$.
    \item Tỉ số: $6 : 30 = \dfrac{1}{5}$.
\end{itemize}
\begin{itemize}
    \item a) \textbf{Đúng}.
    \item b) \textbf{Sai}: Kết quả là $0{,}3\text{ m}^2$.
    \item c) \textbf{Đúng}.
    \item d) \textbf{Đúng}.
\end{itemize}
}
\end{ex}

\textbf{Phần 3. Trả lời ngắn}
\setcounterpageref{ex}{0}

\begin{ex} %[5H5H25]
Tính diện tích hình tam giác có đáy $4\text{ cm}$ và chiều cao $3\text{ cm}$.
\shortans{6}
\loigiai{
\textbf{Đáp số:} $6$\\
$\dfrac{4 \times 3}{2} = 6\text{ (cm}^2\text{)}$.
}
\end{ex}

\begin{ex} %[5H5H25]
Tính diện tích tấm kính hình tam giác vuông có hai cạnh góc vuông cùng bằng $6\text{ m}$.
\shortans{18}
\loigiai{
\textbf{Đáp số:} $18$\\
$\dfrac{6 \times 6}{2} = 18\text{ (m}^2\text{)}$.
}
\end{ex}

\begin{ex} %[5H5H25]
Hình tam giác có diện tích $25\text{ dm}^2$ và đáy $10\text{ dm}$. Tính chiều cao tương ứng của hình tam giác đó theo đơn vị đề-xi-mét.
\shortans{5}
\loigiai{
\textbf{Đáp số:} $5$\\
Chiều cao là: $25 \times 2 : 10 = 5\text{ (dm)}$.
}
\end{ex}

\begin{ex} %[5H5V25]
Tờ giấy hình vuông cạnh $12\text{ cm}$ cắt bớt tam giác vuông ở góc có hai cạnh góc vuông dài $6\text{ cm}$ và $6\text{ cm}$. Tính diện tích phần tam giác bị cắt theo đơn vị xăng-ti-mét vuông.
\shortans{18}
\loigiai{
\textbf{Đáp số:} $18$\\
$\dfrac{6 \times 6}{2} = 18\text{ (cm}^2\text{)}$.
}
\end{ex}

\begin{ex} %[5H5V25]
Một thửa ruộng hình tam giác có đáy $32\text{ m}$, chiều cao $25\text{ m}$. Tính diện tích thửa ruộng đó theo đơn vị mét vuông.
\shortans{400}
\loigiai{
\textbf{Đáp số:} $400$\\
$\dfrac{32 \times 25}{2} = 400\text{ (m}^2\text{)}$.
}
\end{ex}

\begin{ex} %[5H5V25]
Cắt $50$ lá cờ hình tam giác vuông có kích thước hai cạnh góc vuông là $20\text{ cm}$ và $30\text{ cm}$. Hỏi tổng diện tích vải của $50$ lá cờ đó là bao nhiêu mét vuông? (Nhập kết quả dưới dạng số thập phân).
\shortans{1,5}
\loigiai{
\textbf{Đáp số:} $1{,}5$\\
$50 \times \dfrac{0{,}2 \times 0{,}3}{2} = 1{,}5\text{ (m}^2\text{)}$.
}
\end{ex}

\newpage

\TaoPhieuKiemTra{Chủ đề 5 -- Một số hình phẳng. chu vi và diện tích}{26}{Hình thang. Diện tích hình thang}{https://loptoanthaylinh.vn/thi-online/toan-5/bai-26}
\textbf{Phần 1. Trắc nghiệm lựa chọn}
\setcounterpageref{ex}{0}

\begin{ex} %[5H5N26]
Hình thang là hình tứ giác có
\choice
{Bốn cạnh bằng nhau}
{\True Một cặp cạnh đối diện song song}
{Hai cặp cạnh đối diện song song}
{Bốn góc vuông}
\loigiai{
Hình thang là hình tứ giác có một cặp cạnh đối diện song song (gọi là hai cạnh đáy).\\
\textbf{Đáp án của câu hỏi là B}
}
\end{ex}

\begin{ex} %[5H5N26]
Hình thang có một cạnh bên vuông góc với hai đáy được gọi là
\choice
{\True Hình thang vuông}
{Hình thang cân}
{Hình chữ nhật}
{Hình bình hành}
\loigiai{
Hình thang có một cạnh bên vuông góc với hai đáy gọi là hình thang vuông.\\
\textbf{Đáp án của câu hỏi là A}
}
\end{ex}

\begin{ex} %[5H5H26]
Diện tích hình thang có độ dài hai đáy lần lượt là $4\text{ cm}$ và $6\text{ cm}$, chiều cao $3\text{ cm}$ là
\choice
{$30\text{ cm}^2$}
{$18\text{ cm}^2$}
{\True $15\text{ cm}^2$}
{$24\text{ cm}^2$}
\loigiai{
Diện tích hình thang là: $\dfrac{(4 + 6) \times 3}{2} = 15\text{ (cm}^2\text{)}$.\\
\textbf{Đáp án của câu hỏi là C}
}
\end{ex}

\begin{ex} %[5H5H26]
Diện tích hình thang có độ dài hai đáy lần lượt là $11\text{ cm}$ và $9\text{ cm}$, chiều cao $8\text{ cm}$ là
\choice
{$160\text{ cm}^2$}
{$40\text{ cm}^2$}
{$99\text{ cm}^2$}
{\True $80\text{ cm}^2$}
\loigiai{
Diện tích hình thang là: $\dfrac{(11 + 9) \times 8}{2} = 80\text{ (cm}^2\text{)}$.\\
\textbf{Đáp án của câu hỏi là D}
}
\end{ex}

\begin{ex} %[5H5H26]
Diện tích hình thang có độ dài hai đáy là $25\text{ cm}$ và $15\text{ cm}$, chiều cao $1\text{ dm}$ là
\choice
{$4\text{ cm}^2$}
{$2\text{ cm}^2$}
{\True $2\text{ dm}^2$}
{$4\text{ dm}^2$}
\loigiai{
Đổi $1\text{ dm} = 10\text{ cm}$. Diện tích hình thang là: $\dfrac{(25 + 15) \times 10}{2} = 200\text{ (cm}^2\text{)} = 2\text{ dm}^2$.\\
\textbf{Đáp án của câu hỏi là C}
}
\end{ex}

\begin{ex} %[5H5H26]
Một hình thang có diện tích $60\text{ cm}^2$, chiều cao $6\text{ cm}$. Tổng độ dài hai đáy của hình thang đó là
\choice
{$10\text{ cm}$}
{$15\text{ cm}$}
{$30\text{ cm}$}
{\True $20\text{ cm}$}
\loigiai{
Tổng độ dài hai đáy là: $60 \times 2 : 6 = 20\text{ (cm)}$.\\
\textbf{Đáp án của câu hỏi là D}
}
\end{ex}

\begin{ex} %[5H5H26]
Một hình thang có diện tích $100\text{ dm}^2$, tổng độ dài hai đáy là $20\text{ dm}$. Chiều cao của hình thang đó là
\choice
{$5\text{ dm}$}
{\True $10\text{ dm}$}
{$15\text{ dm}$}
{$20\text{ dm}$}
\loigiai{
Chiều cao của hình thang là: $100 \times 2 : 20 = 10\text{ (dm)}$.\\
\textbf{Đáp án của câu hỏi là B}
}
\end{ex}

\begin{ex} %[5H5V26]
Một mảnh đất có dạng hình thang vuông với đáy bé $2\text{ m}$, đáy lớn $7\text{ m}$ và chiều cao $13\text{ m}$. Diện tích của mảnh đất đó là
\choice
{$117\text{ m}^2$}
{$58\text{ m}^2$}
{$54\text{ m}^2$}
{\True $58{,}5\text{ m}^2$}
\loigiai{
Diện tích mảnh đất là: $\dfrac{(2 + 7) \times 13}{2} = 58{,}5\text{ (m}^2\text{)}$.\\
\textbf{Đáp án của câu hỏi là D}
}
\end{ex}

\begin{ex} %[5H5V26]
Một mặt bàn đa năng được ghép từ $6$ cái bàn nhỏ hình thang giống nhau có đáy lớn $120\text{ cm}$, đáy bé $60\text{ cm}$ và chiều cao $55\text{ cm}$. Tổng diện tích của mặt bàn đa năng đó là
\choice
{\True $29\,700\text{ cm}^2$}
{$4\,950\text{ cm}^2$}
{$9\,900\text{ cm}^2$}
{$297\text{ m}^2$}
\loigiai{
Diện tích $1$ mặt bàn hình thang là: $\dfrac{(120 + 60) \times 55}{2} = 4\,950\text{ (cm}^2\text{)}$.\\
Diện tích $6$ mặt bàn là: $4\,950 \times 6 = 29\,700\text{ (cm}^2\text{)}$.\\
\textbf{Đáp án của câu hỏi là A}
}
\end{ex}

\begin{ex} %[5H5V26]
Một mảnh đất hình thang có đáy lớn $35\text{ m}$, đáy bé $15\text{ m}$, chiều cao $20\text{ m}$. Người ta mua cỏ để trồng phủ kín mảnh đất với giá $45\,000\text{ đồng/m}^2$. Tổng số tiền mua cỏ là
\choice
{$11\,250\,000\text{ đồng}$}
{\True $22\,500\,000\text{ đồng}$}
{$45\,000\,000\text{ đồng}$}
{$2\,250\,000\text{ đồng}$}
\loigiai{
Diện tích mảnh đất là: $\dfrac{(35 + 15) \times 20}{2} = 500\text{ (m}^2\text{)}$.\\
Số tiền mua cỏ là: $500 \times 45\,000 = 22\,500\,000\text{ (đồng)}$.\\
\textbf{Đáp án của câu hỏi là B}
}
\end{ex}

\begin{ex} %[5H5V26]
Một thửa ruộng hình thang có đáy lớn $120\text{ m}$, đáy bé bằng $\dfrac{2}{3}$ đáy lớn, chiều cao bằng $40\text{ m}$. Diện tích thửa ruộng đó là
\choice
{$4\,800\text{ m}^2$}
{$8\,000\text{ m}^2$}
{\True $4\,000\text{ m}^2$}
{$2\,000\text{ m}^2$}
\loigiai{
Đáy bé thửa ruộng là: $120 \times \dfrac{2}{3} = 80\text{ (m)}$.\\
Diện tích thửa ruộng là: $\dfrac{(120 + 80) \times 40}{2} = 4\,000\text{ (m}^2\text{)}$.\\
\textbf{Đáp án của câu hỏi là C}
}
\end{ex}

\begin{ex} %[5H5V26]
Hình thang có độ dài đáy lớn gấp $3$ lần đáy bé, chiều cao bằng $12\text{ cm}$ và diện tích bằng $96\text{ cm}^2$. Độ dài đáy bé của hình thang đó là
\choice
{\True $4\text{ cm}$}
{$8\text{ cm}$}
{$12\text{ cm}$}
{$16\text{ cm}$}
\loigiai{
Tổng hai đáy là: $96 \times 2 : 12 = 16\text{ (cm)}$.\\
Đáy bé là: $16 : (3 + 1) = 4\text{ (cm)}$.\\
\textbf{Đáp án của câu hỏi là A}
}
\end{ex}

\textbf{Phần 2. Trắc nghiệm đúng sai}
\setcounterpageref{ex}{0}

\begin{ex} %[5H5H26]
Xét tính đúng sai của các khẳng định sau:
\choiceTF
{\True Hai cạnh đáy của hình thang luôn song song với nhau.}
{Hình thang có hai góc vuông thì hai góc vuông đó luôn cùng nằm trên một đáy.}
{\True Diện tích hình thang có hai đáy $14\text{ dm}$ và $6\text{ dm}$, chiều cao $10\text{ dm}$ là $100\text{ dm}^2$.}
{Nếu tăng chiều cao của hình thang lên gấp đôi và giữ nguyên hai đáy thì diện tích tăng gấp bốn lần.}
\loigiai{
\begin{itemize}
    \item Hai cạnh đáy của hình thang luôn song song với nhau.
    \item Hai góc vuông của hình thang vuông nằm ở cùng một cạnh bên (ứng với hai đáy), không phải cùng nằm trên một đáy.
    \item $S = \dfrac{(14 + 6) \times 10}{2} = 100\text{ dm}^2$.
    \item Diện tích tỉ lệ thuận với chiều cao nên tăng chiều cao 2 lần thì diện tích tăng 2 lần (khẳng định tăng 4 lần là sai).
\end{itemize}
\begin{itemize}
    \item a) \textbf{Đúng}.
    \item b) \textbf{Sai}.
    \item c) \textbf{Đúng}.
    \item d) \textbf{Sai}.
\end{itemize}
}
\end{ex}

\begin{ex} %[5H5V26]
Một hình vẽ con thuyền gồm thân thuyền là hình thang có hai đáy $8\text{ cm}$ và $4\text{ cm}$, chiều cao $2\text{ cm}$; hai cánh buồm là hai tam giác vuông có cùng đáy $2\text{ cm}$ và chiều cao $4\text{ cm}$. Xét tính đúng sai của các khẳng định sau:
\choiceTF
{\True Diện tích thân thuyền bằng $12\text{ cm}^2$.}
{\True Diện tích của mỗi cánh buồm là $4\text{ cm}^2$.}
{Tổng diện tích của hai cánh buồm bằng diện tích thân thuyền.}
{\True Tổng diện tích toàn bộ con thuyền là $20\text{ cm}^2$.}
\loigiai{
\begin{itemize}
    \item Diện tích thân thuyền: $\dfrac{(8 + 4) \times 2}{2} = 12\text{ cm}^2$.
    \item Diện tích mỗi cánh buồm: $\dfrac{2 \times 4}{2} = 4\text{ cm}^2$.
    \item Tổng diện tích 2 cánh buồm: $4 + 4 = 8\text{ cm}^2 \neq 12\text{ cm}^2$.
    \item Toàn bộ con thuyền: $12 + 4 + 4 = 20\text{ cm}^2$.
\end{itemize}
\begin{itemize}
    \item a) \textbf{Đúng}.
    \item b) \textbf{Đúng}.
    \item c) \textbf{Sai}: $8\text{ cm}^2 \neq 12\text{ cm}^2$.
    \item d) \textbf{Đúng}.
\end{itemize}
}
\end{ex}

\textbf{Phần 3. Trả lời ngắn}
\setcounterpageref{ex}{0}

\begin{ex} %[5H5H26]
Tính diện tích hình thang có hai đáy lần lượt là $4\text{ cm}$ và $6\text{ cm}$, chiều cao $3\text{ cm}$.
\shortans{15}
\loigiai{
\textbf{Đáp số:} $15$\\
$\dfrac{(4 + 6) \times 3}{2} = 15\text{ (cm}^2\text{)}$.
}
\end{ex}

\begin{ex} %[5H5H26]
Tính diện tích hình thang có hai đáy lần lượt là $6\text{ m}$ và $4\text{ m}$, chiều cao $4\text{ m}$.
\shortans{20}
\loigiai{
\textbf{Đáp số:} $20$\\
$\dfrac{(6 + 4) \times 4}{2} = 20\text{ (m}^2\text{)}$.
}
\end{ex}

\begin{ex} %[5H5H26]
Một hình thang có diện tích $175\text{ cm}^2$, tổng hai đáy là $35\text{ cm}$. Tính chiều cao hình thang đó theo đơn vị xăng-ti-mét.
\shortans{10}
\loigiai{
\textbf{Đáp số:} $10$\\
Chiều cao là: $175 \times 2 : 35 = 10\text{ (cm)}$.
}
\end{ex}

\begin{ex} %[5H5V26]
Mảnh đất hình thang vuông có hai đáy là $2\text{ m}$ và $7\text{ m}$, chiều cao $13\text{ m}$. Tính diện tích mảnh đất đó theo mét vuông. (Nhập kết quả dưới dạng số thập phân).
\shortans{58,5}
\loigiai{
\textbf{Đáp số:} $58{,}5$\\
$\dfrac{(2 + 7) \times 13}{2} = 58{,}5\text{ (m}^2\text{)}$.
}
\end{ex}

\begin{ex} %[5H5V26]
Một mảnh đất hình thang có hai đáy là $35\text{ m}$ và $15\text{ m}$, chiều cao $20\text{ m}$. Tính diện tích mảnh đất đó theo mét vuông.
\shortans{500}
\loigiai{
\textbf{Đáp số:} $500$\\
$\dfrac{(35 + 15) \times 20}{2} = 500\text{ (m}^2\text{)}$.
}
\end{ex}

\begin{ex} %[5H5V26]
Hình thang có diện tích $96\text{ cm}^2$, chiều cao $12\text{ cm}$ và đáy lớn gấp $3$ lần đáy bé. Tính độ dài đáy lớn của hình thang đó theo đơn vị xăng-ti-mét.
\shortans{12}
\loigiai{
\textbf{Đáp số:} $12$\\
Tổng hai đáy: $96 \times 2 : 12 = 16\text{ cm}$. Đáy lớn: $16 : (3 + 1) \times 3 = 12\text{ cm}$.
}
\end{ex}

\newpage

\TaoPhieuKiemTra{Chủ đề 5 -- Một số hình phẳng. chu vi và diện tích}{27}{Đường tròn. Chu vi và diện tích hình tròn}{https://loptoanthaylinh.vn/thi-online/toan-5/bai-27}
\textbf{Phần 1. Trắc nghiệm lựa chọn}
\setcounterpageref{ex}{0}

\begin{ex} %[5H5N27]
Công thức tính chu vi $C$ của hình tròn có bán kính $r$ là
\choice
{$C = 3{,}14 \times r$}
{$C = 3{,}14 \times r \times r$}
{$C = 3{,}14 \times 4 \times r$}
{\True $C = 3{,}14 \times 2 \times r$}
\loigiai{
Chu vi hình tròn bằng tích của số $3{,}14$ với $2$ lần bán kính: $C = 3{,}14 \times 2 \times r$.\\
\textbf{Đáp án của câu hỏi là D}
}
\end{ex}

\begin{ex} %[5H5N27]
Công thức tính diện tích $S$ của hình tròn có bán kính $r$ là
\choice
{\True $S = 3{,}14 \times r \times r$}
{$S = 3{,}14 \times 2 \times r$}
{$S = 3{,}14 \times d$}
{$S = 3{,}14 \times 2 \times r \times r$}
\loigiai{
Diện tích hình tròn bằng số $3{,}14$ nhân với bán kính rồi nhân với bán kính: $S = 3{,}14 \times r \times r$.\\
\textbf{Đáp án của câu hỏi là A}
}
\end{ex}

\begin{ex} %[5H5H27]
Chu vi của hình tròn có đường kính $6\text{ dm}$ là
\choice
{$18\text{ dm}$}
{$37{,}68\text{ dm}$}
{$28{,}26\text{ dm}$}
{\True $18{,}84\text{ dm}$}
\loigiai{
Chu vi hình tròn là: $3{,}14 \times 6 = 18{,}84\text{ (dm)}$.\\
\textbf{Đáp án của câu hỏi là D}
}
\end{ex}

\begin{ex} %[5H5H27]
Diện tích của hình tròn có bán kính $10\text{ cm}$ là
\choice
{$31{,}4\text{ cm}^2$}
{\True $314\text{ cm}^2$}
{$62{,}8\text{ cm}^2$}
{$3\,140\text{ cm}^2$}
\loigiai{
Diện tích hình tròn là: $3{,}14 \times 10 \times 10 = 314\text{ (cm}^2\text{)}$.\\
\textbf{Đáp án của câu hỏi là B}
}
\end{ex}

\begin{ex} %[5H5H27]
Diện tích của một tấm thảm hình tròn có đường kính $20\text{ dm}$ là
\choice
{$62{,}8\text{ dm}^2$}
{$1\,256\text{ dm}^2$}
{\True $314\text{ dm}^2$}
{$31{,}4\text{ dm}^2$}
\loigiai{
Bán kính tấm thảm là: $20 : 2 = 10\text{ (dm)}$.\\
Diện tích tấm thảm là: $3{,}14 \times 10 \times 10 = 314\text{ (dm}^2\text{)}$.\\
\textbf{Đáp án của câu hỏi là C}
}
\end{ex}

\begin{ex} %[5H5H27]
Bãi thả khí cầu của một công ty có dạng hình tròn bán kính $200\text{ m}$. Diện tích của bãi thả khí cầu đó là
\choice
{\True $125\,600\text{ m}^2$}
{$1\,256\text{ m}^2$}
{$12\,560\text{ m}^2$}
{$1\,256\,000\text{ m}^2$}
\loigiai{
Diện tích bãi thả khí cầu là: $3{,}14 \times 200 \times 200 = 125\,600\text{ (m}^2\text{)}$.\\
\textbf{Đáp án của câu hỏi là A}
}
\end{ex}

\begin{ex} %[5H5H27]
Một hình tròn có chu vi là $31{,}4\text{ cm}$. Bán kính của hình tròn đó là
\choice
{$10\text{ cm}$}
{$2{,}5\text{ cm}$}
{\True $5\text{ cm}$}
{$15{,}7\text{ cm}$}
\loigiai{
Bán kính hình tròn là: $31{,}4 : 3{,}14 : 2 = 5\text{ (cm)}$.\\
\textbf{Đáp án của câu hỏi là C}
}
\end{ex}

\begin{ex} %[5H5V27]
Cái cạp rổ hình tròn có đường kính $50\text{ cm}$. Chiều dài phần hai đầu tre buộc vào nhau là $15\text{ cm}$. Chiều dài của thanh tre uốn thành cái cạp rổ là
\choice
{\True $172\text{ cm}$}
{$157\text{ cm}$}
{$182\text{ cm}$}
{$329\text{ cm}$}
\loigiai{
Chu vi cạp rổ là: $3{,}14 \times 50 = 157\text{ (cm)}$.\\
Chiều dài thanh tre là: $157 + 15 = 172\text{ (cm)}$.\\
\textbf{Đáp án của câu hỏi là A}
}
\end{ex}

\begin{ex} %[5H5V27]
Một sợi dây thừng quấn quanh một gốc cây hình tròn có đường kính $4\text{ dm}$ đúng $3$ vòng. Phần dây không quấn vào thân cây dài $2{,}8\text{ m}$. Chiều dài của sợi dây thừng đó là
\choice
{$3{,}768\text{ m}$}
{\True $6{,}568\text{ m}$}
{$40{,}48\text{ m}$}
{$65{,}68\text{ m}$}
\loigiai{
Chu vi gốc cây là: $3{,}14 \times 4 = 12{,}56\text{ (dm)} = 1{,}256\text{ (m)}$.\\
Chiều dài sợi dây thừng là: $1{,}256 \times 3 + 2{,}8 = 3{,}768 + 2{,}8 = 6{,}568\text{ (m)}$.\\
\textbf{Đáp án của câu hỏi là B}
}
\end{ex}

\begin{ex} %[5H5V27]
Ở một vùng sa mạc, người ta trồng lúa trên những thửa ruộng hình tròn bán kính $50\text{ m}$. Biết có $1\,000$ thửa ruộng như vậy. Tổng diện tích trồng lúa là
\choice
{$78\,500\text{ m}^2$}
{$785\,000\text{ m}^2$}
{$31\,400\,000\text{ m}^2$}
{\True $7\,850\,000\text{ m}^2$}
\loigiai{
Diện tích $1$ thửa ruộng là: $3{,}14 \times 50 \times 50 = 7\,850\text{ (m}^2\text{)}$.\\
Tổng diện tích $1\,000$ thửa ruộng là: $7\,850 \times 1\,000 = 7\,850\,000\text{ (m}^2\text{)}$.\\
\textbf{Đáp án của câu hỏi là D}
}
\end{ex}

\begin{ex} %[5H5V27]
Một sân vận động được tạo bởi một hình chữ nhật có kích thước $100\text{ m} \times 70\text{ m}$ và hai nửa hình tròn có đường kính $70\text{ m}$ ở hai đầu. Diện tích của sân vận động đó là
\choice
{$7\,000\text{ m}^2$}
{$14\,693\text{ m}^2$}
{\True $10\,846{,}5\text{ m}^2$}
{$10\,146{,}5\text{ m}^2$}
\loigiai{
Diện tích phần hình chữ nhật là: $100 \times 70 = 7\,000\text{ (m}^2\text{)}$.\\
Bán kính hai nửa hình tròn là: $70 : 2 = 35\text{ (m)}$.\\
Diện tích hai nửa hình tròn (tạo thành một hình tròn) là: $3{,}14 \times 35 \times 35 = 3\,846{,}5\text{ (m}^2\text{)}$.\\
Diện tích sân vận động là: $7\,000 + 3\,846{,}5 = 10\,846{,}5\text{ (m}^2\text{)}$.\\
\textbf{Đáp án của câu hỏi là C}
}
\end{ex}

\begin{ex} %[5H5V27]
Một giếng nước có miệng giếng là hình tròn bán kính $8\text{ dm}$. Người ta xây thành giếng rộng $3\text{ dm}$ bao quanh miệng giếng. Diện tích phần đất xây thành giếng đó là
\choice
{$200{,}96\text{ dm}^2$}
{\True $178{,}98\text{ dm}^2$}
{$379{,}94\text{ dm}^2$}
{$150{,}72\text{ dm}^2$}
\loigiai{
Bán kính hình tròn lớn (bao gồm cả thành giếng) là: $8 + 3 = 11\text{ (dm)}$.\\
Diện tích hình tròn lớn là: $3{,}14 \times 11 \times 11 = 379{,}94\text{ (dm}^2\text{)}$.\\
Diện tích miệng giếng là: $3{,}14 \times 8 \times 8 = 200{,}96\text{ (dm}^2\text{)}$.\\
Diện tích thành giếng là: $379{,}94 - 200{,}96 = 178{,}98\text{ (dm}^2\text{)}$.\\
\textbf{Đáp án của câu hỏi là B}
}
\end{ex}

\textbf{Phần 2. Trắc nghiệm đúng sai}
\setcounterpageref{ex}{0}

\begin{ex} %[5H5H27]
Xét tính đúng sai của các khẳng định sau:
\choiceTF
{\True Khi bán kính hình tròn tăng lên gấp đôi thì chu vi hình tròn tăng lên gấp đôi.}
{Khi bán kính hình tròn tăng lên gấp đôi thì diện tích hình tròn tăng lên gấp đôi.}
{\True Chu vi của hình tròn có bán kính $5\text{ cm}$ là $31{,}4\text{ cm}$.}
{\True Diện tích của hình tròn có bán kính $3\text{ dm}$ là $28{,}26\text{ dm}^2$.}
\loigiai{
\begin{itemize}
    \item $C = 3{,}14 \times 2 \times r \Rightarrow$ tỉ lệ thuận bậc 1 với $r$, nên $r$ gấp 2 thì chu vi gấp 2.
    \item $S = 3{,}14 \times r \times r \Rightarrow$ khi $r$ gấp 2 thì $S$ gấp $2 \times 2 = 4$ lần (khẳng định tăng 2 lần là sai).
    \item $C = 3{,}14 \times 2 \times 5 = 31{,}4\text{ cm}$.
    \item $S = 3{,}14 \times 3 \times 3 = 28{,}26\text{ dm}^2$.
\end{itemize}
\begin{itemize}
    \item a) \textbf{Đúng}.
    \item b) \textbf{Sai}.
    \item c) \textbf{Đúng}.
    \item d) \textbf{Đúng}.
\end{itemize}
}
\end{ex}

\begin{ex} %[5H5V27]
Chú rùa vàng bò từ $A$ đến $B$ theo đường nửa đường tròn lớn màu đỏ (đường kính $9\text{ dm}$), chú rùa nâu bò từ $B$ đến $A$ theo $3$ nửa đường tròn nhỏ màu xanh (có đường kính lần lượt là $3\text{ dm}, 1\text{ dm}, 5\text{ dm}$). Xét tính đúng sai của các khẳng định sau:
\choiceTF
{\True Quãng đường rùa vàng đi là $14{,}13\text{ dm}$.}
{\True Quãng đường rùa nâu đi là $14{,}13\text{ dm}$.}
{Rùa vàng đi quãng đường dài hơn rùa nâu.}
{\True Quãng đường đi của hai chú rùa là bằng nhau.}
\loigiai{
\begin{itemize}
    \item Nửa chu vi đường tròn lớn: $3{,}14 \times 9 : 2 = 14{,}13\text{ dm}$.
    \item Tổng 3 nửa chu vi nhỏ: $3{,}14 \times (3 + 1 + 5) : 2 = 3{,}14 \times 9 : 2 = 14{,}13\text{ dm}$.
    \item Quãng đường hai chú rùa đi bằng nhau.
\end{itemize}
\begin{itemize}
    \item a) \textbf{Đúng}.
    \item b) \textbf{Đúng}.
    \item c) \textbf{Sai}.
    \item d) \textbf{Đúng}.
\end{itemize}
}
\end{ex}

\textbf{Phần 3. Trả lời ngắn}
\setcounterpageref{ex}{0}

\begin{ex} %[5H5H27]
Tính chu vi của hình tròn có đường kính $10\text{ m}$ theo mét. (Nhập kết quả dưới dạng số thập phân).
\shortans{31,4}
\loigiai{
\textbf{Đáp số:} $31{,}4$\\
$3{,}14 \times 10 = 31{,}4\text{ (m)}$.
}
\end{ex}

\begin{ex} %[5H5H27]
Tính diện tích của hình tròn có bán kính $2\text{ cm}$ theo xăng-ti-mét vuông. (Nhập kết quả dưới dạng số thập phân).
\shortans{12,56}
\loigiai{
\textbf{Đáp số:} $12{,}56$\\
$3{,}14 \times 2 \times 2 = 12{,}56\text{ (cm}^2\text{)}$.
}
\end{ex}

\begin{ex} %[5H5H27]
Một hình tròn có chu vi là $62{,}8\text{ cm}$. Tính diện tích của hình tròn đó theo xăng-ti-mét vuông.
\shortans{314}
\loigiai{
\textbf{Đáp số:} $314$\\
Bán kính: $62{,}8 : 3{,}14 : 2 = 10\text{ cm}$. Diện tích: $3{,}14 \times 10 \times 10 = 314\text{ cm}^2$.
}
\end{ex}

\begin{ex} %[5H5V27]
Một sợi dây thừng quấn quanh gốc cây hình tròn đường kính $4\text{ dm}$ đủ $3$ vòng. Phần thừa dài $2{,}8\text{ m}$. Hỏi cả sợi dây dài bao nhiêu mét? (Nhập kết quả dưới dạng số thập phân).
\shortans{6,568}
\loigiai{
\textbf{Đáp số:} $6{,}568$\\
Chu vi 1 vòng: $3{,}14 \times 0{,}4 = 1{,}256\text{ m}$. Dây dài: $1{,}256 \times 3 + 2{,}8 = 6{,}568\text{ m}$.
}
\end{ex}

\begin{ex} %[5H5V27]
Sân vận động gồm hình chữ nhật $100\text{ m} \times 70\text{ m}$ và hai nửa hình tròn đường kính $70\text{ m}$. Tính diện tích sân vận động đó theo mét vuông. (Nhập kết quả dưới dạng số thập phân).
\shortans{10846,5}
\loigiai{
\textbf{Đáp số:} $10846{,}5$\\
$100 \times 70 + 3{,}14 \times 35 \times 35 = 7\,000 + 3\,846{,}5 = 10\,846{,}5\text{ (m}^2\text{)}$.
}
\end{ex}

\begin{ex} %[5H5V27]
Miệng giếng hình tròn bán kính $8\text{ dm}$, thành giếng bao quanh rộng $3\text{ dm}$. Tính diện tích phần đất xây thành giếng theo đề-xi-mét vuông. (Nhập kết quả dưới dạng số thập phân).
\shortans{178,98}
\loigiai{
\textbf{Đáp số:} $178{,}98$\\
$3{,}14 \times (11^2 - 8^2) = 3{,}14 \times 57 = 178{,}98\text{ (dm}^2\text{)}$.
}
\end{ex}

\newpage

\TaoPhieuKiemTra{Chủ đề 5 -- Một số hình phẳng. chu vi và diện tích}{28}{Thực hành và trải nghiệm hình phẳng}{https://loptoanthaylinh.vn/thi-online/toan-5/bai-28}
\textbf{Phần 1. Trắc nghiệm lựa chọn}
\setcounterpageref{ex}{0}

\begin{ex} %[5H5N28]
Bộ đồ chơi xếp hình Tangram truyền thống gồm có bao nhiêu miếng ghép?
\choice
{$5$ miếng ghép}
{$6$ miếng ghép}
{$8$ miếng ghép}
{\True $7$ miếng ghép}
\loigiai{
Bộ đồ chơi xếp hình Tangram gồm có $7$ miếng ghép (gồm $5$ hình tam giác, $1$ hình vuông và $1$ hình bình hành).\\
\textbf{Đáp án của câu hỏi là D}
}
\end{ex}

\begin{ex} %[5H5N28]
Trong bộ xếp hình Tangram $7$ miếng ghép, có bao nhiêu miếng ghép là hình tam giác?
\choice
{$3$ miếng}
{$4$ miếng}
{\True $5$ miếng}
{$6$ miếng}
\loigiai{
Trong bộ $7$ miếng ghép gồm có $5$ hình tam giác vuông cân, $1$ hình vuông và $1$ hình bình hành.\\
\textbf{Đáp án của câu hỏi là C}
}
\end{ex}

\begin{ex} %[5H5H28]
Để tạo hình mái nhà trong bức tranh ngôi nhà, người ta thường dùng hình phẳng nào dưới đây?
\choice
{Hình tròn}
{\True Hình thang cân hoặc hình tam giác}
{Hình chữ nhật}
{Hình thoi}
\loigiai{
Mái nhà thường có dạng hình thang cân hoặc hình tam giác.\\
\textbf{Đáp án của câu hỏi là B}
}
\end{ex}

\begin{ex} %[5H5H28]
Một bông hoa được tạo thành từ $1$ hình tròn màu vàng làm nhụy và $5$ hình tròn nhỏ màu đỏ có cùng bán kính $1\text{ cm}$ làm cánh hoa. Bán kính hình tròn nhụy hoa là $2\text{ cm}$. Diện tích của nhụy hoa đó là
\choice
{$6{,}28\text{ cm}^2$}
{$3{,}14\text{ cm}^2$}
{$15{,}7\text{ cm}^2$}
{\True $12{,}56\text{ cm}^2$}
\loigiai{
Diện tích nhụy hoa là: $3{,}14 \times 2 \times 2 = 12{,}56\text{ (cm}^2\text{)}$.\\
\textbf{Đáp án của câu hỏi là D}
}
\end{ex}

\begin{ex} %[5H5H28]
Một mái nhà có dạng hình thang với hai đáy là $8\text{ m}$ và $12\text{ m}$, chiều cao là $3\text{ m}$. Diện tích mái nhà đó là
\choice
{$60\text{ m}^2$}
{\True $30\text{ m}^2$}
{$20\text{ m}^2$}
{$40\text{ m}^2$}
\loigiai{
Diện tích mái nhà là: $\dfrac{(8 + 12) \times 3}{2} = 30\text{ (m}^2\text{)}$.\\
\textbf{Đáp án của câu hỏi là B}
}
\end{ex}

\begin{ex} %[5H5H28]
Một bức tường hình chữ nhật có kích thước $6\text{ m} \times 4\text{ m}$. Trên tường có $1$ cửa ra vào hình chữ nhật $1\text{ m} \times 2\text{ m}$ và $1$ cửa sổ hình vuông cạnh $1\text{ m}$. Diện tích phần tường cần sơn là
\choice
{\True $21\text{ m}^2$}
{$24\text{ m}^2$}
{$22\text{ m}^2$}
{$20\text{ m}^2$}
\loigiai{
Diện tích bức tường là: $6 \times 4 = 24\text{ (m}^2\text{)}$.\\
Diện tích các cửa là: $1 \times 2 + 1 \times 1 = 3\text{ (m}^2\text{)}$.\\
Diện tích phần tường cần sơn là: $24 - 3 = 21\text{ (m}^2\text{)}$.\\
\textbf{Đáp án của câu hỏi là A}
}
\end{ex}

\begin{ex} %[5H5H28]
Ghép hai hình tam giác vuông bằng nhau có các cạnh góc vuông là $3\text{ cm}$ và $4\text{ cm}$ thành một hình chữ nhật. Chu vi hình chữ nhật đó là
\choice
{$7\text{ cm}$}
{$12\text{ cm}$}
{$24\text{ cm}$}
{\True $14\text{ cm}$}
\loigiai{
Hình chữ nhật tạo thành có chiều dài $4\text{ cm}$ và chiều rộng $3\text{ cm}$.\\
Chu vi là: $(4 + 3) \times 2 = 14\text{ (cm)}$.\\
\textbf{Đáp án của câu hỏi là D}
}
\end{ex}

\begin{ex} %[5H5V28]
Từ một tờ giấy hình vuông có diện tích $64\text{ cm}^2$, bạn Mai cắt ra thành $7$ miếng ghép Tangram rồi ghép lại thành một hình con thỏ. Diện tích hình con thỏ đó bằng
\choice
{\True $64\text{ cm}^2$}
{$32\text{ cm}^2$}
{$128\text{ cm}^2$}
{$16\text{ cm}^2$}
\loigiai{
Vì hình con thỏ được ghép lại từ toàn bộ $7$ mảnh của hình vuông ban đầu nên diện tích không đổi và bằng $64\text{ cm}^2$.\\
\textbf{Đáp án của câu hỏi là A}
}
\end{ex}

\begin{ex} %[5H5V28]
Một cánh diều có dạng hình thoi gồm hai đường chéo dài $60\text{ cm}$ và $80\text{ cm}$. Diện tích của cánh diều đó là
\choice
{$4\,800\text{ cm}^2$}
{$1\,200\text{ cm}^2$}
{\True $2\,400\text{ cm}^2$}
{$240\text{ cm}^2$}
\loigiai{
Diện tích hình thoi bằng nửa tích hai đường chéo: $\dfrac{60 \times 80}{2} = 2\,400\text{ (cm}^2\text{)}$.\\
\textbf{Đáp án của câu hỏi là C}
}
\end{ex}

\begin{ex} %[5H5V28]
Để làm một chiếc thuyền buồm origami, bạn Việt dùng một tờ giấy hình chữ nhật có kích thước $15\text{ cm} \times 20\text{ cm}$. Diện tích của tờ giấy đó là
\choice
{\True $300\text{ cm}^2$}
{$35\text{ cm}^2$}
{$70\text{ cm}^2$}
{$150\text{ cm}^2$}
\loigiai{
Diện tích tờ giấy là: $15 \times 20 = 300\text{ (cm}^2\text{)}$.\\
\textbf{Đáp án của câu hỏi là A}
}
\end{ex}

\begin{ex} %[5H5V28]
Bạn Nam cắt một tấm bìa hình vuông cạnh $20\text{ cm}$ thành $4$ hình tam giác bằng nhau bằng hai đường chéo. Diện tích của mỗi hình tam giác đó là
\choice
{$50\text{ cm}^2$}
{\True $100\text{ cm}^2$}
{$200\text{ cm}^2$}
{$400\text{ cm}^2$}
\loigiai{
Diện tích hình vuông ban đầu là: $20 \times 20 = 400\text{ (cm}^2\text{)}$.\\
Diện tích mỗi hình tam giác là: $400 : 4 = 100\text{ (cm}^2\text{)}$.\\
\textbf{Đáp án của câu hỏi là B}
}
\end{ex}

\begin{ex} %[5H5V28]
Người ta cắt $4$ góc của một tấm bìa hình vuông cạnh $30\text{ cm}$, mỗi góc bị cắt là một hình vuông nhỏ cạnh $5\text{ cm}$. Diện tích phần bìa còn lại là
\choice
{$900\text{ cm}^2$}
{$850\text{ cm}^2$}
{\True $800\text{ cm}^2$}
{$700\text{ cm}^2$}
\loigiai{
Diện tích tấm bìa ban đầu: $30 \times 30 = 900\text{ (cm}^2\text{)}$.\\
Diện tích $4$ hình vuông nhỏ bị cắt: $4 \times (5 \times 5) = 100\text{ (cm}^2\text{)}$.\\
Diện tích phần bìa còn lại: $900 - 100 = 800\text{ (cm}^2\text{)}$.\\
\textbf{Đáp án của câu hỏi là C}
}
\end{ex}

\textbf{Phần 2. Trắc nghiệm đúng sai}
\setcounterpageref{ex}{0}

\begin{ex} %[5H5H28]
Xét tính đúng sai của các khẳng định sau về các phép đo, vẽ, ghép hình phẳng:
\choiceTF
{\True Khi ghép hai tam giác vuông cân bằng nhau ta có thể tạo thành một hình vuông.}
{Khi cắt một hình phẳng thành các mảnh nhỏ rồi ghép lại thì tổng diện tích bị giảm đi.}
{\True Bộ ghép hình Tangram gồm đúng $5$ hình tam giác, $1$ hình vuông và $1$ hình bình hành.}
{Gấp đôi một tờ giấy hình vuông theo đường chéo ta nhận được một hình thang vuông.}
\loigiai{
\begin{itemize}
    \item Hai tam giác vuông cân bằng nhau ghép theo cạnh huyền tạo thành một hình vuông.
    \item Diện tích hình được bảo toàn khi cắt ghép, không bị giảm đi.
    \item Tangram gồm 7 miếng: 5 tam giác, 1 hình vuông, 1 hình bình hành.
    \item Gấp đôi hình vuông theo đường chéo ta được tam giác vuông cân, không phải hình thang.
\end{itemize}
\begin{itemize}
    \item a) \textbf{Đúng}.
    \item b) \textbf{Sai}.
    \item c) \textbf{Đúng}.
    \item d) \textbf{Sai}.
\end{itemize}
}
\end{ex}

\begin{ex} %[5H5V28]
Một bức tranh trang trí gồm một hình vuông lớn cạnh $10\text{ cm}$ và một hình tròn nội tiếp bên trong tiếp xúc với bốn cạnh của hình vuông (đường kính $10\text{ cm}$). Xét tính đúng sai của các khẳng định sau:
\choiceTF
{\True Diện tích của hình vuông lớn là $100\text{ cm}^2$.}
{\True Bán kính của hình tròn là $5\text{ cm}$.}
{\True Diện tích của hình tròn là $78{,}5\text{ cm}^2$.}
{Diện tích phần trang trí còn lại (bên ngoài hình tròn, bên trong hình vuông) là $22{,}5\text{ cm}^2$.}
\loigiai{
\begin{itemize}
    \item Diện tích hình vuông: $10 \times 10 = 100\text{ cm}^2$.
    \item Bán kính hình tròn: $10 : 2 = 5\text{ cm}$.
    \item Diện tích hình tròn: $3{,}14 \times 5 \times 5 = 78{,}5\text{ cm}^2$.
    \item Diện tích phần còn lại: $100 - 78{,}5 = 21{,}5\text{ cm}^2 \neq 22{,}5\text{ cm}^2$.
\end{itemize}
\begin{itemize}
    \item a) \textbf{Đúng}.
    \item b) \textbf{Đúng}.
    \item c) \textbf{Đúng}.
    \item d) \textbf{Sai}: Kết quả là $21{,}5\text{ cm}^2$.
\end{itemize}
}
\end{ex}

\textbf{Phần 3. Trả lời ngắn}
\setcounterpageref{ex}{0}

\begin{ex} %[5H5H28]
Một mái nhà hình thang có hai đáy là $8\text{ m}$ và $12\text{ m}$, chiều cao $3\text{ m}$. Tính diện tích mái nhà đó theo mét vuông.
\shortans{30}
\loigiai{
\textbf{Đáp số:} $30$\\
$\dfrac{(8 + 12) \times 3}{2} = 30\text{ (m}^2\text{)}$.
}
\end{ex}

\begin{ex} %[5H5H28]
Tờ giấy hình chữ nhật có kích thước $15\text{ cm} \times 20\text{ cm}$. Tính diện tích tờ giấy đó theo xăng-ti-mét vuông.
\shortans{300}
\loigiai{
\textbf{Đáp số:} $300$\\
$15 \times 20 = 300\text{ (cm}^2\text{)}$.
}
\end{ex}

\begin{ex} %[5H5H28]
Cắt tấm bìa hình vuông cạnh $20\text{ cm}$ thành $4$ hình tam giác bằng nhau. Tính diện tích mỗi tam giác theo xăng-ti-mét vuông.
\shortans{100}
\loigiai{
\textbf{Đáp số:} $100$\\
$20 \times 20 : 4 = 100\text{ (cm}^2\text{)}$.
}
\end{ex}

\begin{ex} %[5H5V28]
Cắt $4$ góc hình vuông cạnh $5\text{ cm}$ từ tấm bìa hình vuông cạnh $30\text{ cm}$. Tính diện tích phần bìa còn lại theo xăng-ti-mét vuông.
\shortans{800}
\loigiai{
\textbf{Đáp số:} $800$\\
$30 \times 30 - 4 \times 5 \times 5 = 900 - 100 = 800\text{ (cm}^2\text{)}$.
}
\end{ex}

\begin{ex} %[5H5V28]
Tính diện tích phần nằm giữa hình vuông cạnh $10\text{ cm}$ và hình tròn có đường kính $10\text{ cm}$ nằm bên trong hình vuông theo xăng-ti-mét vuông. (Nhập kết quả dưới dạng số thập phân).
\shortans{21,5}
\loigiai{
\textbf{Đáp số:} $21{,}5$\\
$10 \times 10 - 3{,}14 \times 5 \times 5 = 100 - 78{,}5 = 21{,}5\text{ (cm}^2\text{)}$.
}
\end{ex}

\begin{ex} %[5H5V28]
Một cánh diều hình thoi có hai đường chéo dài $60\text{ cm}$ và $80\text{ cm}$. Tính diện tích cánh diều theo xăng-ti-mét vuông.
\shortans{2400}
\loigiai{
\textbf{Đáp số:} $2400$\\
$\dfrac{60 \times 80}{2} = 2\,400\text{ (cm}^2\text{)}$.
}
\end{ex}

\newpage

\TaoPhieuKiemTra{Chủ đề 5 -- Một số hình phẳng. chu vi và diện tích}{29}{Luyện tập chung hình phẳng}{https://loptoanthaylinh.vn/thi-online/toan-5/bai-29}
\textbf{Phần 1. Trắc nghiệm lựa chọn}
\setcounterpageref{ex}{0}

\begin{ex} %[5H5N29]
Một hình tròn có bán kính $50\text{ cm}$ và một hình tròn khác có bán kính $200\text{ cm}$. Chu vi hình tròn lớn gấp chu vi hình tròn bé số lần là
\choice
{\True $4$ lần}
{$2$ lần}
{$3$ lần}
{$5$ lần}
\loigiai{
Tỉ số bán kính hai hình tròn là $200 : 50 = 4$. Do đó chu vi hình tròn lớn gấp $4$ lần chu vi hình tròn bé.\\
\textbf{Đáp án của câu hỏi là A}
}
\end{ex}

\begin{ex} %[5H5H29]
Một cái ao có dạng nửa hình tròn đường kính $12\text{ m}$. Chu vi của cái ao đó (gồm đường cong nửa hình tròn và đường kính bờ ao) là
\choice
{$18{,}84\text{ m}$}
{$37{,}68\text{ m}$}
{\True $30{,}84\text{ m}$}
{$49{,}68\text{ m}$}
\loigiai{
Độ dài đường cong nửa hình tròn là: $3{,}14 \times 12 : 2 = 18{,}84\text{ (m)}$.\\
Chu vi của cái ao là: $18{,}84 + 12 = 30{,}84\text{ (m)}$.\\
\textbf{Đáp án của câu hỏi là C}
}
\end{ex}

\begin{ex} %[5H5H29]
Cho hình chữ nhật $ABCD$ có chiều dài $b\text{ cm}$ và chiều rộng $a\text{ cm}$. Lấy điểm $E$ bất kì trên cạnh $AB$. Diện tích hình tam giác $ECD$ là
\choice
{\True $\dfrac{a \times b}{2}\text{ (cm}^2\text{)}$}
{$a \times b\text{ (cm}^2\text{)}$}
{$\dfrac{a \times b}{4}\text{ (cm}^2\text{)}$}
{Thay đổi tùy thuộc vào vị trí điểm $E$}
\loigiai{
Tam giác $ECD$ có đáy là $CD = b\text{ cm}$ và chiều cao hạ từ $E$ xuống $CD$ bằng chiều rộng $a\text{ cm}$. Do đó diện tích tam giác $ECD$ luôn bằng $\dfrac{a \times b}{2}\text{ (cm}^2\text{)}$.\\
\textbf{Đáp án của câu hỏi là A}
}
\end{ex}

\begin{ex} %[5H5H29]
Một hình vuông có cạnh $8\text{ cm}$ chứa một hình tròn có đường kính $8\text{ cm}$ bên trong. Diện tích phần hình vuông nằm ngoài hình tròn là
\choice
{$50{,}24\text{ cm}^2$}
{$114{,}24\text{ cm}^2$}
{$136{,}96\text{ cm}^2$}
{\True $13{,}76\text{ cm}^2$}
\loigiai{
Diện tích hình vuông là: $8 \times 8 = 64\text{ (cm}^2\text{)}$.\\
Bán kính hình tròn là: $8 : 2 = 4\text{ (cm)}$.\\
Diện tích hình tròn là: $3{,}14 \times 4 \times 4 = 50{,}24\text{ (cm}^2\text{)}$.\\
Diện tích phần màu xanh là: $64 - 50{,}24 = 13{,}76\text{ (cm}^2\text{)}$.\\
\textbf{Đáp án của câu hỏi là D}
}
\end{ex}

\begin{ex} %[5H5H29]
Bạn Việt nặn ba hình từ đất sét: hình tam giác vuông có hai cạnh góc vuông dài $7\text{ cm}$ và $7\text{ cm}$; hình thang có hai đáy $3\text{ cm}, 6\text{ cm}$ và chiều cao $4\text{ cm}$; hình tròn bán kính $5\text{ cm}$. Hình có diện tích bé nhất là
\choice
{Hình tam giác}
{Hình tròn}
{\True Hình thang}
{Ba hình có diện tích bằng nhau}
\loigiai{
Diện tích hình tam giác: $\dfrac{7 \times 7}{2} = 24{,}5\text{ cm}^2$.\\
Diện tích hình thang: $\dfrac{(3 + 6) \times 4}{2} = 18\text{ cm}^2$.\\
Diện tích hình tròn: $3{,}14 \times 5 \times 5 = 78{,}5\text{ cm}^2$.\\
Vậy hình thang có diện tích bé nhất ($18\text{ cm}^2$).\\
\textbf{Đáp án của câu hỏi là C}
}
\end{ex}

\begin{ex} %[5H5H29]
Trong ba hình ở câu trên, hình có diện tích lớn nhất là
\choice
{Hình tam giác}
{\True Hình tròn}
{Hình thang}
{Không có hình lớn nhất}
\loigiai{
Hình tròn có diện tích lớn nhất ($78{,}5\text{ cm}^2$).\\
\textbf{Đáp án của câu hỏi là B}
}
\end{ex}

\begin{ex} %[5H5V29]
Một mảnh đất gồm hình thang $ABCD$ (hai đáy $30\text{ m}$ và $55\text{ m}$, chiều cao $25\text{ m}$) và hình tam giác $AED$ (đáy $55\text{ m}$, chiều cao $27\text{ m}$). Diện tích toàn bộ mảnh đất đó là
\choice
{$1\,062{,}5\text{ m}^2$}
{$742{,}5\text{ m}^2$}
{$1\,500\text{ m}^2$}
{\True $1\,805\text{ m}^2$}
\loigiai{
Diện tích hình thang $ABCD$: $\dfrac{(55 + 30) \times 25}{2} = 1\,062{,}5\text{ (m}^2\text{)}$.\\
Diện tích tam giác $AED$: $\dfrac{55 \times 27}{2} = 742{,}5\text{ (m}^2\text{)}$.\\
Diện tích toàn bộ mảnh đất: $1\,062{,}5 + 742{,}5 = 1\,805\text{ (m}^2\text{)}$.\\
\textbf{Đáp án của câu hỏi là D}
}
\end{ex}

\begin{ex} %[5H5V29]
Cho hình vuông $ABCD$ cạnh $6{,}5\text{ cm}$. Trên cạnh $CD$ lấy các điểm sao cho $DE = EG = GH = HK = KC = 1{,}3\text{ cm}$. Diện tích hình thang vuông $ABCK$ là
\choice
{$50{,}7\text{ cm}^2$}
{\True $25{,}35\text{ cm}^2$}
{$12{,}675\text{ cm}^2$}
{$42{,}25\text{ cm}^2$}
\loigiai{
Đáy $AB = 6{,}5\text{ cm}$, đáy $KC = 1{,}3\text{ cm}$, chiều cao $BC = 6{,}5\text{ cm}$.\\
Diện tích hình thang $ABCK$ là: $\dfrac{(6{,}5 + 1{,}3) \times 6{,}5}{2} = \dfrac{7{,}8 \times 6{,}5}{2} = 25{,}35\text{ (cm}^2\text{)}$.\\
\textbf{Đáp án của câu hỏi là B}
}
\end{ex}

\begin{ex} %[5H5V29]
Cũng trong hình vuông $ABCD$ ở câu trên, diện tích tam giác $AKD$ gấp diện tích tam giác $ADE$ số lần là
\choice
{$2$ lần}
{$3$ lần}
{\True $4$ lần}
{$5$ lần}
\loigiai{
Đáy $KD = 4 \times 1{,}3 = 5{,}2\text{ cm}$ gấp $4$ lần đáy $DE = 1{,}3\text{ cm}$. Hai tam giác có cùng chiều cao $AD$ nên diện tích tam giác $AKD$ gấp $4$ lần diện tích tam giác $ADE$.\\
\textbf{Đáp án của câu hỏi là C}
}
\end{ex}

\begin{ex} %[5H5V29]
Một chiếc khăn trải bàn hình chữ nhật có chiều dài $2\text{ m}$ và chiều rộng $1{,}5\text{ m}$. Ở giữa khăn người ta thêu một hình thoi có các đường chéo dài $1\text{ m}$ và $0{,}8\text{ m}$. Diện tích phần vải không thêu là
\choice
{$3\text{ m}^2$}
{\True $2{,}6\text{ m}^2$}
{$0{,}4\text{ m}^2$}
{$2{,}2\text{ m}^2$}
\loigiai{
Diện tích khăn trải bàn: $2 \times 1{,}5 = 3\text{ (m}^2\text{)}$.\\
Diện tích hình thoi: $\dfrac{1 \times 0{,}8}{2} = 0{,}4\text{ (m}^2\text{)}$.\\
Diện tích phần không thêu: $3 - 0{,}4 = 2{,}6\text{ (m}^2\text{)}$.\\
\textbf{Đáp án của câu hỏi là B}
}
\end{ex}

\begin{ex} %[5H5V29]
Một bánh xe đạp có bán kính $0{,}35\text{ m}$. Khi người đi xe đạp đi được quãng đường $219{,}8\text{ m}$ thì bánh xe đó đã lăn được số vòng là
\choice
{$50$ vòng}
{$150$ vòng}
{$200$ vòng}
{\True $100$ vòng}
\loigiai{
Chu vi bánh xe là: $3{,}14 \times 2 \times 0{,}35 = 2{,}198\text{ (m)}$.\\
Số vòng lăn của bánh xe là: $219{,}8 : 2{,}198 = 100$ (vòng).\\
\textbf{Đáp án của câu hỏi là D}
}
\end{ex}

\begin{ex} %[5H5V29]
Nếu bán kính của một hình tròn tăng thêm $20\%$ thì diện tích của hình tròn đó tăng thêm bao nhiêu phần trăm?
\choice
{\True $44\%$}
{$20\%$}
{$40\%$}
{$48\%$}
\loigiai{
Bán kính mới bằng $120\% = 1{,}2$ lần bán kính cũ.\\
Diện tích mới bằng: $1{,}2 \times 1{,}2 = 1{,}44 = 144\%$ diện tích cũ.\\
Diện tích tăng thêm: $144\% - 100\% = 44\%$.\\
\textbf{Đáp án của câu hỏi là A}
}
\end{ex}

\textbf{Phần 2. Trắc nghiệm đúng sai}
\setcounterpageref{ex}{0}

\begin{ex} %[5H5H29]
Cho hình tròn đỏ có đường kính gấp $2$ lần đường kính của hình tròn xanh. Xét tính đúng sai của các khẳng định sau:
\choiceTF
{\True Bán kính hình tròn đỏ gấp $2$ lần bán kính hình tròn xanh.}
{\True Chu vi hình tròn đỏ gấp $2$ lần chu vi hình tròn xanh.}
{Diện tích hình tròn đỏ gấp $2$ lần diện tích hình tròn xanh.}
{\True Diện tích hình tròn đỏ gấp $4$ lần diện tích hình tròn xanh.}
\loigiai{
\begin{itemize}
    \item $d_{đỏ} = 2 \times d_{xanh} \Rightarrow r_{đỏ} = 2 \times r_{xanh}$.
    \item Chu vi: $C_{đỏ} = 3{,}14 \times d_{đỏ} = 2 \times C_{xanh}$.
    \item Diện tích: $S_{đỏ} = 3{,}14 \times (2 \times r_{xanh}) \times (2 \times r_{xanh}) = 4 \times S_{xanh} \neq 2 \times S_{xanh}$.
    \item Diện tích hình tròn đỏ gấp 4 lần diện tích hình tròn xanh.
\end{itemize}
\begin{itemize}
    \item a) \textbf{Đúng}.
    \item b) \textbf{Đúng}.
    \item c) \textbf{Sai}.
    \item d) \textbf{Đúng}.
\end{itemize}
}
\end{ex}

\begin{ex} %[5H5V29]
Một mảnh đất có dạng như hình vẽ được chia thành hình thang $ABCD$ (đáy $30\text{ m}$ và $55\text{ m}$, chiều cao $25\text{ m}$) và hình tam giác $AED$ (đáy $55\text{ m}$, chiều cao $27\text{ m}$). Xét tính đúng sai của các khẳng định sau:
\choiceTF
{\True Diện tích phần hình thang $ABCD$ là $1\,062{,}5\text{ m}^2$.}
{\True Diện tích phần hình tam giác $AED$ là $742{,}5\text{ m}^2$.}
{Diện tích phần tam giác $AED$ lớn hơn diện tích phần hình thang $ABCD$.}
{\True Diện tích toàn bộ mảnh đất là $1\,805\text{ m}^2$.}
\loigiai{
\begin{itemize}
    \item Diện tích hình thang: $\dfrac{(55 + 30) \times 25}{2} = 1\,062{,}5\text{ m}^2$.
    \item Diện tích hình tam giác: $\dfrac{55 \times 27}{2} = 742{,}5\text{ m}^2$.
    \item $742{,}5\text{ m}^2 < 1\,062{,}5\text{ m}^2$ nên diện tích tam giác bé hơn diện tích hình thang.
    \item Tổng diện tích: $1\,062{,}5 + 742{,}5 = 1\,805\text{ m}^2$.
\end{itemize}
\begin{itemize}
    \item a) \textbf{Đúng}.
    \item b) \textbf{Đúng}.
    \item c) \textbf{Sai}.
    \item d) \textbf{Đúng}.
\end{itemize}
}
\end{ex}

\textbf{Phần 3. Trả lời ngắn}
\setcounterpageref{ex}{0}

\begin{ex} %[5H5H29]
Cái ao dạng nửa hình tròn đường kính $12\text{ m}$. Tính chu vi của cái ao đó theo mét. (Nhập kết quả dưới dạng số thập phân).
\shortans{30,84}
\loigiai{
\textbf{Đáp số:} $30{,}84$\\
$3{,}14 \times 12 : 2 + 12 = 18{,}84 + 12 = 30{,}84\text{ (m)}$.
}
\end{ex}

\begin{ex} %[5H5H29]
Hình vuông cạnh $8\text{ cm}$ chứa hình tròn đường kính $8\text{ cm}$. Tính diện tích phần hình vuông nằm ngoài hình tròn theo xăng-ti-mét vuông. (Nhập kết quả dưới dạng số thập phân).
\shortans{13,76}
\loigiai{
\textbf{Đáp số:} $13{,}76$\\
$8 \times 8 - 3{,}14 \times 4 \times 4 = 64 - 50{,}24 = 13{,}76\text{ (cm}^2\text{)}$.
}
\end{ex}

\begin{ex} %[5H5V29]
Hình thang có hai đáy $30\text{ m}, 55\text{ m}$, chiều cao $25\text{ m}$ và tam giác có đáy $55\text{ m}$, chiều cao $27\text{ m}$. Tính tổng diện tích của hai hình theo mét vuông.
\shortans{1805}
\loigiai{
\textbf{Đáp số:} $1805$\\
$\dfrac{(55 + 30) \times 25}{2} + \dfrac{55 \times 27}{2} = 1\,062{,}5 + 742{,}5 = 1\,805\text{ (m}^2\text{)}$.
}
\end{ex}

\begin{ex} %[5H5V29]
Hình vuông cạnh $6{,}5\text{ cm}$. Tính diện tích hình thang $ABCK$ có hai đáy $6{,}5\text{ cm}$ và $1{,}3\text{ cm}$, chiều cao $6{,}5\text{ cm}$ theo xăng-ti-mét vuông. (Nhập kết quả dưới dạng số thập phân).
\shortans{25,35}
\loigiai{
\textbf{Đáp số:} $25{,}35$\\
$\dfrac{(6{,}5 + 1{,}3) \times 6{,}5}{2} = 25{,}35\text{ (cm}^2\text{)}$.
}
\end{ex}

\begin{ex} %[5H5V29]
Một bánh xe đạp có bán kính $0{,}35\text{ m}$. Hỏi khi bánh xe lăn được $100$ vòng thì xe đi được quãng đường bao nhiêu mét? (Nhập kết quả dưới dạng số thập phân).
\shortans{219,8}
\loigiai{
\textbf{Đáp số:} $219{,}8$\\
$3{,}14 \times 2 \times 0{,}35 \times 100 = 219{,}8\text{ (m)}$.
}
\end{ex}

\begin{ex} %[5H5V29]
Một mảnh khăn hình chữ nhật $2\text{ m} \times 1{,}5\text{ m}$ thêu hình thoi ở giữa có hai đường chéo dài $1\text{ m}$ và $0{,}8\text{ m}$. Tính diện tích phần khăn không thêu theo mét vuông. (Nhập kết quả dưới dạng số thập phân).
\shortans{2,6}
\loigiai{
\textbf{Đáp số:} $2{,}6$\\
$2 \times 1{,}5 - \dfrac{1 \times 0{,}8}{2} = 3 - 0{,}4 = 2{,}6\text{ (m}^2\text{)}$.
}
\end{ex}

\newpage

% =========================================================================
% CHỦ ĐỀ 6 - ÔN TẬP HỌC KÌ 1
% =========================================================================
\section{Ôn tập học kì 1}

\TaoPhieuKiemTra{Chủ đề 6 -- Ôn tập học kì 1}{30}{Ôn tập số thập phân}{https://loptoanthaylinh.vn/thi-online/toan-5/bai-30}
\textbf{Phần 1. Trắc nghiệm lựa chọn}
\setcounterpageref{ex}{0}

\begin{ex} %[5D6N30]
Số thập phân gồm \textit{Hai mươi ba đơn vị, tám mươi mốt phần trăm} được viết là
\choice
{$238{,}1$}
{$2{,}381$}
{\True $23{,}81$}
{$0{,}2381$}
\loigiai{
Số thập phân đó là $23{,}81$.\\
\textbf{Đáp án của câu hỏi là C}
}
\end{ex}

\begin{ex} %[5D6N30]
Chữ số $3$ trong số thập phân $2{,}03$ có giá trị thuộc hàng nào?
\choice
{\True Hàng phần trăm}
{Hàng chục}
{Hàng đơn vị}
{Hàng phần mười}
\loigiai{
Chữ số $3$ đứng ở vị trí thứ hai sau dấu phẩy nên thuộc hàng phần trăm.\\
\textbf{Đáp án của câu hỏi là A}
}
\end{ex}

\begin{ex} %[5D6H30]
Đổi $8\text{ mm}$ ra mét ta được số thập phân là
\choice
{$0{,}8\text{ m}$}
{$0{,}08\text{ m}$}
{$8\,000\text{ m}$}
{\True $0{,}008\text{ m}$}
\loigiai{
Vì $1\text{ m} = 1\,000\text{ mm}$ nên $8\text{ mm} = \dfrac{8}{1\,000}\text{ m} = 0{,}008\text{ m}$.\\
\textbf{Đáp án của câu hỏi là D}
}
\end{ex}

\begin{ex} %[5D6H30]
So sánh hai số thập phân $17{,}1$ và $9{,}725$, khẳng định nào sau đây là đúng?
\choice
{$17{,}1 < 9{,}725$ vì $17{,}1$ có $3$ chữ số còn $9{,}725$ có $4$ chữ số}
{\True $17{,}1 > 9{,}725$ vì phần nguyên $17 > 9$}
{$17{,}1 = 9{,}725$}
{Không so sánh được vì số chữ số khác nhau}
\loigiai{
So sánh phần nguyên trước: $17 > 9$ nên $17{,}1 > 9{,}725$.\\
\textbf{Đáp án của câu hỏi là B}
}
\end{ex}

\begin{ex} %[5D6H30]
Đổi $800\text{ kg}$ ra tấn ta được số thập phân là
\choice
{$8\text{ tấn}$}
{$0{,}08\text{ tấn}$}
{\True $0{,}8\text{ tấn}$}
{$80\text{ tấn}$}
\loigiai{
Vì $1\text{ tấn} = 1\,000\text{ kg}$ nên $800\text{ kg} = \dfrac{800}{1\,000}\text{ tấn} = 0{,}8\text{ tấn}$.\\
\textbf{Đáp án của câu hỏi là C}
}
\end{ex}

\begin{ex} %[5D6H30]
Viết số đo diện tích $1\text{ m}^2\,5\text{ dm}^2$ dưới dạng số thập phân với đơn vị mét vuông là
\choice
{\True $1{,}05\text{ m}^2$}
{$1{,}5\text{ m}^2$}
{$1{,}005\text{ m}^2$}
{$15\text{ m}^2$}
\loigiai{
$1\text{ m}^2\,5\text{ dm}^2 = 1\dfrac{5}{100}\text{ m}^2 = 1{,}05\text{ m}^2$.\\
\textbf{Đáp án của câu hỏi là A}
}
\end{ex}

\begin{ex} %[5D6H30]
Làm tròn số thập phân $9{,}548$ đến hàng phần trăm ta được số
\choice
{$9{,}5$}
{$9{,}54$}
{\True $9{,}55$}
{$10$}
\loigiai{
Chữ số ở hàng phần nghìn là $8 \ge 5$ nên ta tăng chữ số hàng phần trăm thêm $1$: được $9{,}55$.\\
\textbf{Đáp án của câu hỏi là C}
}
\end{ex}

\begin{ex} %[5D6H30]
Kỉ lục chạy $100\text{ m}$ là $9{,}58\text{ giây}$. Làm tròn kết quả này đến hàng phần mười ta được
\choice
{$9{,}5\text{ giây}$}
{$10\text{ giây}$}
{$9{,}58\text{ giây}$}
{\True $9{,}6\text{ giây}$}
\loigiai{
Chữ số hàng phần trăm là $8 \ge 5$ nên làm tròn lên được $9{,}6\text{ giây}$.\\
\textbf{Đáp án của câu hỏi là D}
}
\end{ex}

\begin{ex} %[5D6V30]
Cân nặng của $4$ loại quả trong siêu thị là: Táo nặng $2{,}5\text{ kg}$, dưa hấu nặng $4{,}19\text{ kg}$, chuối nặng $4{,}0\text{ kg}$, nho nặng $1{,}75\text{ kg}$. Loại quả nhẹ nhất là
\choice
{Táo}
{\True Nho}
{Chuối}
{Dưa hấu}
\loigiai{
So sánh: $1{,}75\text{ kg} < 2{,}5\text{ kg} < 4{,}0\text{ kg} < 4{,}19\text{ kg}$. Vậy quả nhẹ nhất là nho.\\
\textbf{Đáp án của câu hỏi là B}
}
\end{ex}

\begin{ex} %[5D6V30]
Tìm chữ số thích hợp thay vào dấu $?$ để được số $2{,}?1$ thỏa mãn $2 < 2{,}?1 < 3$.
\choice
{Chỉ có thể là chữ số $0$}
{Chỉ có thể là chữ số $9$}
{Không có chữ số nào thỏa mãn}
{\True Bất kì chữ số nào từ $0$ đến $9$}
\loigiai{
Với mọi chữ số $?$ từ $0$ đến $9$, phần nguyên đều bằng $2$ và phần thập phân $> 0$ nên $2 < 2{,}?1 < 3$ luôn đúng.\\
\textbf{Đáp án của câu hỏi là D}
}
\end{ex}

\begin{ex} %[5D6V30]
Một túi cà chua cân nặng $1{,}45\text{ kg}$ và một túi hành tây cân nặng $500\text{ g}$. Cả hai túi cân nặng tất cả là
\choice
{$1{,}5\text{ kg}$}
{\True $1{,}95\text{ kg}$}
{$6{,}45\text{ kg}$}
{$1{,}455\text{ kg}$}
\loigiai{
Đổi $500\text{ g} = 0{,}5\text{ kg}$. Tổng cân nặng là: $1{,}45 + 0{,}5 = 1{,}95\text{ (kg)}$.\\
\textbf{Đáp án của câu hỏi là B}
}
\end{ex}

\begin{ex} %[5D6V30]
Sắp xếp các số đo chiều cao: Mai cao $150\text{ cm}$, Mi cao $1{,}2\text{ m}$, Rô-bốt cao $90\text{ cm}$. Thứ tự các bạn từ cao nhất đến thấp nhất là
\choice
{\True Mai, Mi, Rô-bốt}
{Mi, Mai, Rô-bốt}
{Rô-bốt, Mi, Mai}
{Mai, Rô-bốt, Mi}
\loigiai{
Đổi về cùng đơn vị mét: Mai cao $1{,}5\text{ m}$; Mi cao $1{,}2\text{ m}$; Rô-bốt cao $0{,}9\text{ m}$.\\
Thứ tự từ cao đến thấp: Mai ($1{,}5\text{ m}$) $>$ Mi ($1{,}2\text{ m}$) $>$ Rô-bốt ($0{,}9\text{ m}$).\\
\textbf{Đáp án của câu hỏi là A}
}
\end{ex}

\textbf{Phần 2. Trắc nghiệm đúng sai}
\setcounterpageref{ex}{0}

\begin{ex} %[5D6H30]
Xét tính đúng sai của các khẳng định sau về số thập phân:
\choiceTF
{\True Số $100{,}10$ bằng số $100{,}1$.}
{Chữ số $3$ trong số $109{,}37$ thuộc hàng chục.}
{\True Đổi $17\text{ ml} = 0{,}017\text{ l}$.}
{\True Làm tròn số $100{,}917$ đến hàng phần trăm ta được $100{,}92$.}
\loigiai{
\begin{itemize}
    \item $100{,}10 = 100{,}1$ (bỏ chữ số 0 tận cùng bên phải phần thập phân).
    \item Chữ số 3 trong $109{,}37$ thuộc hàng phần mười, không phải hàng chục.
    \item $17\text{ ml} = 0{,}017\text{ l}$.
    \item $100{,}917$ có chữ số hàng phần nghìn là $7 \ge 5$ nên làm tròn thành $100{,}92$.
\end{itemize}
\begin{itemize}
    \item a) \textbf{Đúng}.
    \item b) \textbf{Sai}.
    \item c) \textbf{Đúng}.
    \item d) \textbf{Đúng}.
\end{itemize}
}
\end{ex}

\begin{ex} %[5D6V30]
Kỉ lục thế giới chạy $200\text{ m}$ là $19{,}19\text{ giây}$. Xét tính đúng sai của các khẳng định sau:
\choiceTF
{\True Phần nguyên của số chỉ thời gian kỉ lục là $19$.}
{\True Làm tròn kỉ lục đến hàng phần mười ta được $19{,}2\text{ giây}$.}
{Làm tròn kỉ lục đến số tự nhiên gần nhất ta được $20\text{ giây}$.}
{\True Thời gian kỉ lục trên ít hơn $20\text{ giây}$.}
\loigiai{
\begin{itemize}
    \item Phần nguyên là 19.
    \item Chữ số hàng phần trăm là $9 \ge 5 \Rightarrow$ làm tròn đến hàng phần mười là $19{,}2\text{ s}$.
    \item Chữ số hàng phần mười là $1 < 5 \Rightarrow$ làm tròn đến số tự nhiên gần nhất là $19\text{ giây} \neq 20\text{ giây}$.
    \item $19{,}19 < 20$.
\end{itemize}
\begin{itemize}
    \item a) \textbf{Đúng}.
    \item b) \textbf{Đúng}.
    \item c) \textbf{Sai}: Làm tròn là $19\text{ giây}$.
    \item d) \textbf{Đúng}.
\end{itemize}
}
\end{ex}

\textbf{Phần 3. Trả lời ngắn}
\setcounterpageref{ex}{0}

\begin{ex} %[5D6H30]
Đổi $173\text{ cm}$ ra mét dưới dạng số thập phân.
\shortans{1,73}
\loigiai{
\textbf{Đáp số:} $1{,}73$\\
$173\text{ cm} = 1{,}73\text{ m}$.
}
\end{ex}

\begin{ex} %[5D6H30]
Đổi $3\text{ dm}^2\,75\text{ cm}^2$ ra đề-xi-mét vuông dưới dạng số thập phân.
\shortans{3,75}
\loigiai{
\textbf{Đáp số:} $3{,}75$\\
$3\text{ dm}^2\,75\text{ cm}^2 = 3{,}75\text{ dm}^2$.
}
\end{ex}

\begin{ex} %[5D6H30]
Làm tròn số thập phân $17{,}153$ đến hàng phần trăm. (Nhập kết quả dưới dạng số thập phân).
\shortans{17,15}
\loigiai{
\textbf{Đáp số:} $17{,}15$\\
Chữ số hàng phần nghìn là $3 < 5$ nên giữ nguyên: $17{,}15$.
}
\end{ex}

\begin{ex} %[5D6V30]
Một túi cà chua nặng $1{,}45\text{ kg}$ và túi hành tây nặng $0{,}5\text{ kg}$. Tính tổng cân nặng của hai túi theo ki-lô-gam. (Nhập kết quả dưới dạng số thập phân).
\shortans{1,95}
\loigiai{
\textbf{Đáp số:} $1{,}95$\\
$1{,}45 + 0{,}5 = 1{,}95\text{ (kg)}$.
}
\end{ex}

\begin{ex} %[5D6V30]
Kỉ lục chạy $100\text{ m}$ là $9{,}58\text{ giây}$. Làm tròn số đo thời gian này đến hàng phần mười. (Nhập kết quả dưới dạng số thập phân).
\shortans{9,6}
\loigiai{
\textbf{Đáp số:} $9{,}6$\\
$9{,}58 \approx 9{,}6\text{ (giây)}$.
}
\end{ex}

\begin{ex} %[5D6V30]
Đổi $3\text{ dm}^2$ ra mét vuông dưới dạng số thập phân.
\shortans{0,03}
\loigiai{
\textbf{Đáp số:} $0{,}03$\\
$3\text{ dm}^2 = 0{,}03\text{ m}^2$.
}
\end{ex}

\newpage

\TaoPhieuKiemTra{Chủ đề 6 -- Ôn tập học kì 1}{31}{Ôn tập các phép tính với số thập phân}{https://loptoanthaylinh.vn/thi-online/toan-5/bai-31}
\textbf{Phần 1. Trắc nghiệm lựa chọn}
\setcounterpageref{ex}{0}

\begin{ex} %[5D6N31]
Kết quả của phép tính $7{,}8 + 11{,}6$ là
\choice
{\True $19{,}4$}
{$18{,}4$}
{$19{,}14$}
{$18{,}14$}
\loigiai{
Ta có: $7{,}8 + 11{,}6 = 19{,}4$.\\
\textbf{Đáp án của câu hỏi là A}
}
\end{ex}

\begin{ex} %[5D6N31]
Kết quả của phép tính $27 - 1{,}2$ là
\choice
{$15{,}8$}
{$26{,}8$}
{\True $25{,}8$}
{$25{,}2$}
\loigiai{
Ta có: $27{,}0 - 1{,}2 = 25{,}8$.\\
\textbf{Đáp án của câu hỏi là C}
}
\end{ex}

\begin{ex} %[5D6H31]
Tính thuận tiện biểu thức $9{,}2 + 17{,}56 + 0{,}8$ ta được kết quả là
\choice
{$26{,}56$}
{$28{,}56$}
{$37{,}56$}
{\True $27{,}56$}
\loigiai{
Ta có: $(9{,}2 + 0{,}8) + 17{,}56 = 10 + 17{,}56 = 27{,}56$.\\
\textbf{Đáp án của câu hỏi là D}
}
\end{ex}

\begin{ex} %[5D6H31]
Kết quả của phép tính $9{,}63 - 5{,}275$ là
\choice
{\True $4{,}355$}
{$4{,}365$}
{$4{,}455$}
{$4{,}445$}
\loigiai{
Ta có: $9{,}630 - 5{,}275 = 4{,}355$.\\
\textbf{Đáp án của câu hỏi là A}
}
\end{ex}

\begin{ex} %[5D6H31]
Cho biết $29 \times 37 = 1\,073$. Giá trị của phép tính $0{,}29 \times 3{,}7$ là
\choice
{$107{,}3$}
{$10{,}73$}
{\True $1{,}073$}
{$0{,}1073$}
\loigiai{
Hai thừa số có tất cả $2 + 1 = 3$ chữ số ở phần thập phân nên kết quả là $1{,}073$.\\
\textbf{Đáp án của câu hỏi là C}
}
\end{ex}

\begin{ex} %[5D6H31]
Kết quả của phép tính $1{,}2 \times 3{,}5$ là
\choice
{$0{,}42$}
{$42$}
{$4{,}25$}
{\True $4{,}2$}
\loigiai{
Ta có: $1{,}2 \times 3{,}5 = 4{,}2$.\\
\textbf{Đáp án của câu hỏi là D}
}
\end{ex}

\begin{ex} %[5D6H31]
Ba bạn gom giấy vụn: Mai gom $2{,}5\text{ kg}$, Nam gom $3{,}25\text{ kg}$, Việt gom $2{,}4\text{ kg}$. Bạn gom được nhiều giấy vụn nhất hơn bạn gom được ít giấy vụn nhất là
\choice
{$0{,}75\text{ kg}$}
{\True $0{,}85\text{ kg}$}
{$0{,}1\text{ kg}$}
{$1{,}15\text{ kg}$}
\loigiai{
Nhiều nhất là Nam ($3{,}25\text{ kg}$), ít nhất là Việt ($2{,}4\text{ kg}$).\\
Hiệu: $3{,}25 - 2{,}4 = 0{,}85\text{ (kg)}$.\\
\textbf{Đáp án của câu hỏi là B}
}
\end{ex}

\begin{ex} %[5D6V31]
Rô-bốt cao $0{,}9\text{ m}$. Bạn Mi cao hơn Rô-bốt $0{,}35\text{ m}$ và thấp hơn Mai $0{,}31\text{ m}$. Chiều cao của bạn Mai là
\choice
{$1{,}25\text{ m}$}
{$1{,}55\text{ m}$}
{\True $1{,}56\text{ m}$}
{$1{,}58\text{ m}$}
\loigiai{
Mi cao: $0{,}9 + 0{,}35 = 1{,}25\text{ (m)}$.\\
Mai cao: $1{,}25 + 0{,}31 = 1{,}56\text{ (m)}$.\\
\textbf{Đáp án của câu hỏi là C}
}
\end{ex}

\begin{ex} %[5D6V31]
Trên hai đĩa cân đĩa thăng bằng: Đĩa trái có con Thỏ và hai quả cân $2\text{ kg}, 500\text{ g}$, đĩa phải có con Cáo. Ở một cân khác, quả cân $5\text{ kg}$ thăng bằng với con Cáo và quả cân $200\text{ g}$. Cân nặng của con Thỏ là
\choice
{$4{,}8\text{ kg}$}
{\True $2{,}3\text{ kg}$}
{$2{,}5\text{ kg}$}
{$2{,}8\text{ kg}$}
\loigiai{
Cân thứ hai cho biết: Cáo $+ 0{,}2\text{ kg} = 5\text{ kg} \Rightarrow$ Cáo nặng $5 - 0{,}2 = 4{,}8\text{ kg}$.\\
Cân thứ nhất cho biết: Thỏ $+ 2{,}5\text{ kg} = \text{Cáo} \Rightarrow$ Thỏ nặng $4{,}8 - 2{,}5 = 2{,}3\text{ kg}$.\\
\textbf{Đáp án của câu hỏi là B}
}
\end{ex}

\begin{ex} %[5D6V31]
Sử dụng các tấm thẻ số $6$; $3$; $5$ và một dấu phẩy $,$, hãy lập số thập phân lớn nhất và số thập phân bé nhất (mỗi thẻ dùng đúng một lần). Hiệu của hai số thập phân đó là
\choice
{$61{,}74$}
{$31{,}64$}
{$59{,}74$}
{\True $61{,}74$}
\loigiai{
Số thập phân lớn nhất có thể lập là $65{,}3$.\\
Số thập phân bé nhất có thể lập là $3{,}56$.\\
Hiệu hai số là: $65{,}30 - 3{,}56 = 61{,}74$.\\
\textbf{Đáp án của câu hỏi là D}
}
\end{ex}

\begin{ex} %[5D6V31]
Một mảnh đất hình chữ nhật có chiều dài $24{,}5\text{ m}$ và chiều rộng kém chiều dài $6{,}5\text{ m}$. Chu vi của mảnh đất đó là
\choice
{\True $85\text{ m}$}
{$18\text{ m}$}
{$42{,}5\text{ m}$}
{$441\text{ m}$}
\loigiai{
Chiều rộng mảnh đất là: $24{,}5 - 6{,}5 = 18\text{ (m)}$.\\
Chu vi mảnh đất là: $(24{,}5 + 18) \times 2 = 42{,}5 \times 2 = 85\text{ (m)}$.\\
\textbf{Đáp án của câu hỏi là A}
}
\end{ex}

\begin{ex} %[5D6V31]
Một cửa hàng ngày thứ nhất bán được $125{,}5\text{ kg}$ gạo, ngày thứ hai bán được nhiều hơn ngày thứ nhất $34{,}8\text{ kg}$ gạo. Cả hai ngày cửa hàng bán được tất cả là
\choice
{$160{,}3\text{ kg}$}
{\True $285{,}8\text{ kg}$}
{$286{,}8\text{ kg}$}
{$275{,}8\text{ kg}$}
\loigiai{
Ngày thứ hai bán được: $125{,}5 + 34{,}8 = 160{,}3\text{ (kg)}$.\\
Cả hai ngày bán được: $125{,}5 + 160{,}3 = 285{,}8\text{ (kg)}$.\\
\textbf{Đáp án của câu hỏi là B}
}
\end{ex}

\textbf{Phần 2. Trắc nghiệm đúng sai}
\setcounterpageref{ex}{0}

\begin{ex} %[5D6H31]
Xét tính đúng sai của các phép tính sau:
\choiceTF
{\True Phép cộng $92{,}15 + 7{,}99 + 0{,}01 = 100{,}15$.}
{Phép trừ $19{,}3 - 11 = 18{,}2$.}
{\True Phép nhân $3{,}1 \times 5{,}7 = 17{,}67$.}
{\True Phép nhân $2{,}9 \times 3{,}7 = 10{,}73$.}
\loigiai{
\begin{itemize}
    \item $92{,}15 + (7{,}99 + 0{,}01) = 92{,}15 + 8 = 100{,}15$.
    \item $19{,}3 - 11 = 8{,}3 \neq 18{,}2$.
    \item $3{,}1 \times 5{,}7 = 17{,}67$.
    \item $2{,}9 \times 3{,}7 = 10{,}73$.
\end{itemize}
\begin{itemize}
    \item a) \textbf{Đúng}.
    \item b) \textbf{Sai}: Kết quả là $8{,}3$.
    \item c) \textbf{Đúng}.
    \item d) \textbf{Đúng}.
\end{itemize}
}
\end{ex}

\begin{ex} %[5D6V31]
Ba bạn Mai, Nam, Việt thi gom giấy vụn với kết quả: Mai gom $2{,}5\text{ kg}$, Nam gom $3{,}25\text{ kg}$, Việt gom $2{,}4\text{ kg}$. Xét tính đúng sai của các khẳng định sau:
\choiceTF
{\True Nam là bạn gom được nhiều giấy vụn nhất.}
{Tổng số giấy vụn cả $3$ bạn gom được là $7{,}15\text{ kg}$.}
{\True Nam gom nhiều hơn Mai $0{,}75\text{ kg}$ giấy vụn.}
{\True Tổng khối lượng giấy của Mai và Việt gom được nhiều hơn Nam.}
\loigiai{
\begin{itemize}
    \item Nam gom $3{,}25\text{ kg}$ là nhiều nhất.
    \item Tổng số giấy: $2{,}5 + 3{,}25 + 2{,}4 = 8{,}15\text{ kg} \neq 7{,}15\text{ kg}$.
    \item Nam nhiều hơn Mai: $3{,}25 - 2{,}5 = 0{,}75\text{ kg}$.
    \item Mai + Việt $= 2{,}5 + 2{,}4 = 4{,}9\text{ kg} > 3{,}25\text{ kg}$ (Nam).
\end{itemize}
\begin{itemize}
    \item a) \textbf{Đúng}.
    \item b) \textbf{Sai}: Tổng là $8{,}15\text{ kg}$.
    \item c) \textbf{Đúng}.
    \item d) \textbf{Đúng}.
\end{itemize}
}
\end{ex}

\textbf{Phần 3. Trả lời ngắn}
\setcounterpageref{ex}{0}

\begin{ex} %[5D6H31]
Tính giá trị biểu thức: $9{,}2 + 17{,}56 + 0{,}8$. (Nhập kết quả dưới dạng số thập phân).
\shortans{27,56}
\loigiai{
\textbf{Đáp số:} $27{,}56$\\
$(9{,}2 + 0{,}8) + 17{,}56 = 10 + 17{,}56 = 27{,}56$.
}
\end{ex}

\begin{ex} %[5D6H31]
Tính kết quả phép tính: $9{,}63 - 5{,}275$. (Nhập kết quả dưới dạng số thập phân).
\shortans{4,355}
\loigiai{
\textbf{Đáp số:} $4{,}355$\\
$9{,}63 - 5{,}275 = 4{,}355$.
}
\end{ex}

\begin{ex} %[5D6H31]
Tính kết quả phép tính: $1{,}2 \times 3{,}5$. (Nhập kết quả dưới dạng số thập phân).
\shortans{4,2}
\loigiai{
\textbf{Đáp số:} $4{,}2$\\
$1{,}2 \times 3{,}5 = 4{,}2$.
}
\end{ex}

\begin{ex} %[5D6V31]
Rô-bốt cao $0{,}9\text{ m}$, Mi cao hơn Rô-bốt $0{,}35\text{ m}$, Mai cao hơn Mi $0{,}31\text{ m}$. Hỏi bạn Mai cao bao nhiêu mét? (Nhập kết quả dưới dạng số thập phân).
\shortans{1,56}
\loigiai{
\textbf{Đáp số:} $1{,}56$\\
$0{,}9 + 0{,}35 + 0{,}31 = 1{,}56\text{ (m)}$.
}
\end{ex}

\begin{ex} %[5D6V31]
Biết quả cân $5\text{ kg}$ thăng bằng với con Cáo và quả cân $0{,}2\text{ kg}$; con Cáo lại thăng bằng với con Thỏ và quả cân $2{,}5\text{ kg}$. Tính cân nặng của con Thỏ theo ki-lô-gam. (Nhập kết quả dưới dạng số thập phân).
\shortans{2,3}
\loigiai{
\textbf{Đáp số:} $2{,}3$\\
Cáo: $5 - 0{,}2 = 4{,}8\text{ kg}$. Thỏ: $4{,}8 - 2{,}5 = 2{,}3\text{ kg}$.
}
\end{ex}

\begin{ex} %[5D6V31]
Cho các thẻ số $6, 3, 5$ và dấu phẩy. Lập số lớn nhất và số bé nhất rồi tính hiệu của hai số đó. (Nhập kết quả dưới dạng số thập phân).
\shortans{61,74}
\loigiai{
\textbf{Đáp số:} $61{,}74$\\
$65{,}3 - 3{,}56 = 61{,}74$.
}
\end{ex}

\newpage

\TaoPhieuKiemTra{Chủ đề 6 -- Ôn tập học kì 1}{32}{Ôn tập một số hình phẳng}{https://loptoanthaylinh.vn/thi-online/toan-5/bai-32}
\textbf{Phần 1. Trắc nghiệm lựa chọn}
\setcounterpageref{ex}{0}

\begin{ex} %[5H6N32]
Hình nào dưới đây có hai cặp cạnh đối diện song song và bốn cạnh bằng nhau?
\choice
{Hình chữ nhật}
{Hình bình hành}
{\True Hình thoi}
{Hình thang}
\loigiai{
Hình thoi là hình có hai cặp cạnh đối diện song song và bốn cạnh bằng nhau.\\
\textbf{Đáp án của câu hỏi là C}
}
\end{ex}

\begin{ex} %[5H6N32]
Trong một hình tam giác, đoạn thẳng vuông góc kẻ từ một đỉnh xuống cạnh đối diện được gọi là
\choice
{\True Đường cao}
{Cạnh đáy}
{Bán kính}
{Đường kính}
\loigiai{
Đoạn thẳng vuông góc kẻ từ đỉnh xuống cạnh đối diện là đường cao của tam giác.\\
\textbf{Đáp án của câu hỏi là A}
}
\end{ex}

\begin{ex} %[5H6H32]
Một tờ giấy hình chữ nhật có kích thước $12\text{ cm} \times 5\text{ cm}$. Người ta cắt bỏ hai hình tam giác vuông ở hai góc có các cạnh góc vuông lần lượt là $2\text{ cm} \times 5\text{ cm}$ và $5\text{ cm} \times 5\text{ cm}$. Diện tích phần tờ giấy còn lại là
\choice
{$60\text{ cm}^2$}
{\True $42{,}5\text{ cm}^2$}
{$45\text{ cm}^2$}
{$37{,}5\text{ cm}^2$}
\loigiai{
Diện tích tờ giấy ban đầu là: $12 \times 5 = 60\text{ (cm}^2\text{)}$.\\
Diện tích hai hình tam giác vuông cắt đi là: $\dfrac{2 \times 5}{2} + \dfrac{5 \times 5}{2} = 5 + 12{,}5 = 17{,}5\text{ (cm}^2\text{)}$.\\
Diện tích phần giấy còn lại là: $60 - 17{,}5 = 42{,}5\text{ (cm}^2\text{)}$.\\
\textbf{Đáp án của câu hỏi là B}
}
\end{ex}

\begin{ex} %[5H6H32]
Một mặt hồ có dạng nửa hình tròn với bán kính đo được là $60\text{ m}$. Diện tích mặt hồ đó là
\choice
{$11\,304\text{ m}^2$}
{$188{,}4\text{ m}^2$}
{\True $5\,652\text{ m}^2$}
{$376{,}8\text{ m}^2$}
\loigiai{
Diện tích cả hình tròn bán kính $60\text{ m}$ là: $3{,}14 \times 60 \times 60 = 11\,304\text{ (m}^2\text{)}$.\\
Diện tích mặt hồ (nửa hình tròn) là: $11\,304 : 2 = 5\,652\text{ (m}^2\text{)}$.\\
\textbf{Đáp án của câu hỏi là C}
}
\end{ex}

\begin{ex} %[5H6H32]
Chu vi của miệng một chiếc bình hình tròn có bán kính $5\text{ cm}$ là
\choice
{$15{,}7\text{ cm}$}
{$78{,}5\text{ cm}$}
{$62{,}8\text{ cm}$}
{\True $31{,}4\text{ cm}$}
\loigiai{
Chu vi miệng bình là: $3{,}14 \times 2 \times 5 = 31{,}4\text{ (cm)}$.\\
\textbf{Đáp án của câu hỏi là D}
}
\end{ex}

\begin{ex} %[5H6H32]
Một hình tam giác có độ dài đáy là $8\text{ cm}$ và chiều cao tương ứng là $5\text{ cm}$. Diện tích của hình tam giác đó là
\choice
{\True $20\text{ cm}^2$}
{$40\text{ cm}^2$}
{$13\text{ cm}^2$}
{$10\text{ cm}^2$}
\loigiai{
Diện tích hình tam giác là: $\dfrac{8 \times 5}{2} = 20\text{ (cm}^2\text{)}$.\\
\textbf{Đáp án của câu hỏi là A}
}
\end{ex}

\begin{ex} %[5H6V32]
Rô-bốt có một sợi dây chun dài $15{,}85\text{ cm}$ và có thể bị đứt nếu kéo dài quá $2$ lần chiều dài ban đầu. Rô-bốt dùng sợi chun này để buộc miệng bình hình tròn có bán kính $5\text{ cm}$. Khẳng định nào sau đây đúng?
\choice
{Sợi dây chun sẽ bị đứt vì chu vi miệng bình dài hơn $2$ lần sợi dây}
{\True Sợi dây chun không bị đứt và buộc kín được miệng bình}
{Sợi dây chun quá ngắn không thể quấn quanh miệng bình}
{Sợi dây chun không cần dãn ra cũng vừa vặn miệng bình}
\loigiai{
Chu vi miệng bình là: $3{,}14 \times 2 \times 5 = 31{,}4\text{ cm}$.\\
Khi kéo dãn quanh miệng bình, chiều dài dây gấp: $31{,}4 : 15{,}85 \approx 1{,}98$ lần $< 2$ lần.\\
Do đó dây chun không bị đứt và buộc được.\\
\textbf{Đáp án của câu hỏi là B}
}
\end{ex}

\begin{ex} %[5H6V32]
Một hình bình hành có độ dài đáy $15\text{ cm}$ và chiều cao tương ứng $8\text{ cm}$. Một hình thoi có diện tích bằng diện tích hình bình hành đó và có một đường chéo dài $20\text{ cm}$. Độ dài đường chéo còn lại của hình thoi là
\choice
{$6\text{ cm}$}
{$24\text{ cm}$}
{$10\text{ cm}$}
{\True $12\text{ cm}$}
\loigiai{
Diện tích hình bình hành là: $15 \times 8 = 120\text{ (cm}^2\text{)}$.\\
Diện tích hình thoi bằng $120\text{ cm}^2$.\\
Độ dài đường chéo còn lại là: $120 \times 2 : 20 = 12\text{ (cm)}$.\\
\textbf{Đáp án của câu hỏi là D}
}
\end{ex}

\begin{ex} %[5H6V32]
Cắt ghép $4$ hình tam giác vuông có hai cạnh góc vuông là $3\text{ cm}$ và $4\text{ cm}$ để tạo thành một hình thoi. Diện tích của hình thoi đó là
\choice
{\True $24\text{ cm}^2$}
{$12\text{ cm}^2$}
{$48\text{ cm}^2$}
{$20\text{ cm}^2$}
\loigiai{
Diện tích mỗi hình tam giác vuông là: $\dfrac{3 \times 4}{2} = 6\text{ (cm}^2\text{)}$.\\
Diện tích hình thoi tạo thành là: $6 \times 4 = 24\text{ (cm}^2\text{)}$.\\
\textbf{Đáp án của câu hỏi là A}
}
\end{ex}

\begin{ex} %[5H6V32]
Một tấm bìa hình thang có đáy lớn $18\text{ cm}$, đáy bé $12\text{ cm}$. Nếu kéo dài đáy lớn thêm $4\text{ cm}$ thì diện tích tăng thêm $24\text{ cm}^2$. Diện tích tấm bìa hình thang ban đầu là
\choice
{$90\text{ cm}^2$}
{\True $180\text{ cm}^2$}
{$360\text{ cm}^2$}
{$120\text{ cm}^2$}
\loigiai{
Phần diện tích tăng thêm là tam giác có đáy $4\text{ cm}$ và cùng chiều cao với hình thang.\\
Chiều cao hình thang là: $24 \times 2 : 4 = 12\text{ (cm)}$.\\
Diện tích hình thang ban đầu: $\dfrac{(18 + 12) \times 12}{2} = 180\text{ (cm}^2\text{)}$.\\
\textbf{Đáp án của câu hỏi là B}
}
\end{ex}

\begin{ex} %[5H6V32]
Một mảnh vườn hình chữ nhật có chu vi $60\text{ m}$, chiều dài hơn chiều rộng $6\text{ m}$. Người ta đào một cái ao hình tròn bán kính $3\text{ m}$ ở giữa vườn. Diện tích phần đất còn lại là
\choice
{$216\text{ m}^2$}
{$28{,}26\text{ m}^2$}
{$197{,}74\text{ m}^2$}
{\True $187{,}74\text{ m}^2$}
\loigiai{
Nửa chu vi mảnh vườn: $60 : 2 = 30\text{ (m)}$.\\
Chiều dài mảnh vườn: $(30 + 6) : 2 = 18\text{ (m)}$.\\
Chiều rộng mảnh vườn: $30 - 18 = 12\text{ (m)}$.\\
Diện tích mảnh vườn: $18 \times 12 = 216\text{ (m}^2\text{)}$.\\
Diện tích cái ao: $3{,}14 \times 3 \times 3 = 28{,}26\text{ (m}^2\text{)}$.\\
Diện tích đất còn lại: $216 - 28{,}26 = 187{,}74\text{ (m}^2\text{)}$.\\
\textbf{Đáp án của câu hỏi là D}
}
\end{ex}

\begin{ex} %[5H6V32]
Cho hình vuông có diện tích $36\text{ cm}^2$. Hình tròn nội tiếp bên trong hình vuông tiếp xúc với cả $4$ cạnh của hình vuông có diện tích là
\choice
{$113{,}04\text{ cm}^2$}
{$56{,}52\text{ cm}^2$}
{\True $28{,}26\text{ cm}^2$}
{$9\text{ cm}^2$}
\loigiai{
Vì $6 \times 6 = 36$ nên cạnh hình vuông là $6\text{ cm}$.\\
Đường kính hình tròn bằng cạnh hình vuông là $6\text{ cm} \Rightarrow$ Bán kính $r = 3\text{ cm}$.\\
Diện tích hình tròn là: $3{,}14 \times 3 \times 3 = 28{,}26\text{ (cm}^2\text{)}$.\\
\textbf{Đáp án của câu hỏi là C}
}
\end{ex}

\textbf{Phần 2. Trắc nghiệm đúng sai}
\setcounterpageref{ex}{0}

\begin{ex} %[5H6H32]
Xét tính đúng sai của các khẳng định sau về các hình phẳng:
\choiceTF
{\True Hình thoi có hai đường chéo vuông góc với nhau và cắt nhau tại trung điểm của mỗi đường.}
{Hình thang có hai cặp cạnh đối diện song song.}
{\True Diện tích tam giác vuông bằng nửa tích độ dài hai cạnh góc vuông.}
{\True Chu vi hình tròn có bán kính $r$ tỉ lệ thuận với bán kính $r$.}
\loigiai{
\begin{itemize}
    \item Hình thoi có 2 đường chéo vuông góc tại trung điểm.
    \item Hình thang chỉ có đúng một cặp cạnh đối diện song song (hai đáy).
    \item Diện tích tam giác vuông $= \frac{a \times b}{2}$.
    \item Chu vi $C = 2 \times 3{,}14 \times r$ tỉ lệ thuận với $r$.
\end{itemize}
\begin{itemize}
    \item a) \textbf{Đúng}.
    \item b) \textbf{Sai}.
    \item c) \textbf{Đúng}.
    \item d) \textbf{Đúng}.
\end{itemize}
}
\end{ex}

\begin{ex} %[5H6V32]
Một mặt hồ có dạng nửa hình tròn bán kính $60\text{ m}$. Xét tính đúng sai của các khẳng định sau:
\choiceTF
{\True Đường kính của mặt hồ đó là $120\text{ m}$.}
{Diện tích toàn bộ hình tròn có cùng bán kính là $5\,652\text{ m}^2$.}
{\True Diện tích của mặt hồ (nửa hình tròn) là $5\,652\text{ m}^2$.}
{\True Độ dài đường bờ cong bao quanh nửa mặt hồ là $188{,}4\text{ m}$.}
\loigiai{
\begin{itemize}
    \item Đường kính: $60 \times 2 = 120\text{ m}$.
    \item Diện tích cả hình tròn: $3{,}14 \times 60 \times 60 = 11\,304\text{ m}^2 \neq 5\,652\text{ m}^2$.
    \item Diện tích nửa hình tròn: $11\,304 : 2 = 5\,652\text{ m}^2$.
    \item Độ dài bờ cong: $3{,}14 \times 60 = 188{,}4\text{ m}$.
\end{itemize}
\begin{itemize}
    \item a) \textbf{Đúng}.
    \item b) \textbf{Sai}.
    \item c) \textbf{Đúng}.
    \item d) \textbf{Đúng}.
\end{itemize}
}
\end{ex}

\textbf{Phần 3. Trả lời ngắn}
\setcounterpageref{ex}{0}

\begin{ex} %[5H6H32]
Tính diện tích của hình tam giác có độ dài đáy $8\text{ cm}$ và chiều cao $5\text{ cm}$ theo xăng-ti-mét vuông.
\shortans{20}
\loigiai{
\textbf{Đáp số:} $20$\\
$\dfrac{8 \times 5}{2} = 20\text{ (cm}^2\text{)}$.
}
\end{ex}

\begin{ex} %[5H6H32]
Mặt hồ dạng nửa hình tròn bán kính $60\text{ m}$. Tính diện tích mặt hồ theo mét vuông.
\shortans{5652}
\loigiai{
\textbf{Đáp số:} $5652$\\
$3{,}14 \times 60 \times 60 : 2 = 5\,652\text{ (m}^2\text{)}$.
}
\end{ex}

\begin{ex} %[5H6H32]
Tờ giấy hình chữ nhật $12\text{ cm} \times 5\text{ cm}$ bị cắt bỏ hai tam giác vuông ở góc có diện tích lần lượt là $5\text{ cm}^2$ và $12{,}5\text{ cm}^2$. Tính diện tích phần giấy còn lại theo xăng-ti-mét vuông. (Nhập kết quả dưới dạng số thập phân).
\shortans{42,5}
\loigiai{
\textbf{Đáp số:} $42{,}5$\\
$12 \times 5 - 5 - 12{,}5 = 42{,}5\text{ (cm}^2\text{)}$.
}
\end{ex}

\begin{ex} %[5H6V32]
Hình bình hành có diện tích $120\text{ cm}^2$. Một hình thoi có diện tích bằng diện tích hình bình hành và có một đường chéo $20\text{ cm}$. Tính độ dài đường chéo còn lại của hình thoi theo xăng-ti-mét.
\shortans{12}
\loigiai{
\textbf{Đáp số:} $12$\\
$120 \times 2 : 20 = 12\text{ (cm)}$.
}
\end{ex}

\begin{ex} %[5H6V32]
Hình thang có hai đáy $12\text{ cm}$ và $18\text{ cm}$, chiều cao $12\text{ cm}$. Tính diện tích hình thang đó theo xăng-ti-mét vuông.
\shortans{180}
\loigiai{
\textbf{Đáp số:} $180$\\
$\dfrac{(12 + 18) \times 12}{2} = 180\text{ (cm}^2\text{)}$.
}
\end{ex}

\begin{ex} %[5H6V32]
Hình vuông có diện tích $36\text{ cm}^2$. Tính diện tích của hình tròn tiếp xúc với $4$ cạnh của hình vuông theo xăng-ti-mét vuông. (Nhập kết quả dưới dạng số thập phân).
\shortans{28,26}
\loigiai{
\textbf{Đáp số:} $28{,}26$\\
Cạnh hình vuông là $6\text{ cm} \Rightarrow r = 3\text{ cm} \Rightarrow S = 3{,}14 \times 3 \times 3 = 28{,}26\text{ (cm}^2\text{)}$.
}
\end{ex}

\newpage

\TaoPhieuKiemTra{Chủ đề 6 -- Ôn tập học kì 1}{33}{Ôn tập diện tích, chu vi một số hình phẳng}{https://loptoanthaylinh.vn/thi-online/toan-5/bai-33}
\textbf{Phần 1. Trắc nghiệm lựa chọn}
\setcounterpageref{ex}{0}

\begin{ex} %[5H6N33]
Diện tích của một hình tam giác vuông có độ dài hai cạnh góc vuông lần lượt là $1{,}6\text{ dm}$ và $3{,}5\text{ dm}$ là
\choice
{$2{,}3\text{ dm}^2$}
{$5{,}6\text{ dm}^2$}
{$2{,}8\text{ m}^2$}
{\True $2{,}8\text{ dm}^2$}
\loigiai{
Diện tích hình tam giác vuông là: $\dfrac{1{,}6 \times 3{,}5}{2} = 2{,}8\text{ (dm}^2\text{)}$.\\
\textbf{Đáp án của câu hỏi là D}
}
\end{ex}

\begin{ex} %[5H6H33]
Diện tích của hình thang có độ dài hai đáy lần lượt là $14\text{ cm}$ và $36\text{ cm}$, chiều cao $17\text{ cm}$ là
\choice
{\True $425\text{ cm}^2$}
{$850\text{ cm}^2$}
{$300\text{ cm}^2$}
{$612\text{ cm}^2$}
\loigiai{
Diện tích hình thang là: $\dfrac{(14 + 36) \times 17}{2} = \dfrac{50 \times 17}{2} = 425\text{ (cm}^2\text{)}$.\\
\textbf{Đáp án của câu hỏi là A}
}
\end{ex}

\begin{ex} %[5H6H33]
Một tấm thảm hình tròn có bán kính $3\text{ dm}$. Chu vi và diện tích của tấm thảm đó lần lượt là
\choice
{$9{,}42\text{ dm}$ và $28{,}26\text{ dm}^2$}
{$18{,}84\text{ dm}$ và $18{,}84\text{ dm}^2$}
{\True $18{,}84\text{ dm}$ và $28{,}26\text{ dm}^2$}
{$28{,}26\text{ dm}$ và $28{,}26\text{ dm}^2$}
\loigiai{
Chu vi tấm thảm là: $3{,}14 \times 2 \times 3 = 18{,}84\text{ (dm)}$.\\
Diện tích tấm thảm là: $3{,}14 \times 3 \times 3 = 28{,}26\text{ (dm}^2\text{)}$.\\
\textbf{Đáp án của câu hỏi là C}
}
\end{ex}

\begin{ex} %[5H6H33]
Bánh xe lăn tay dành cho người khuyết tật có đường kính $50\text{ cm}$. Nếu bánh xe lăn trên mặt đất được $1\,000$ vòng thì người đi xe đó đã đi được quãng đường là
\choice
{$157\text{ m}$}
{\True $1\,570\text{ m}$}
{$314\text{ m}$}
{$3\,140\text{ m}$}
\loigiai{
Chu vi bánh xe là: $3{,}14 \times 50 = 157\text{ (cm)} = 1{,}57\text{ (m)}$.\\
Quãng đường đi được là: $1{,}57 \times 1\,000 = 1\,570\text{ (m)}$.\\
\textbf{Đáp án của câu hỏi là B}
}
\end{ex}

\begin{ex} %[5H6H33]
Một mặt bàn hình tròn có đường kính $1{,}4\text{ m}$. Diện tích của mặt bàn đó là
\choice
{$6{,}1544\text{ m}^2$}
{$4{,}396\text{ m}^2$}
{\True $1{,}5386\text{ m}^2$}
{$2{,}198\text{ m}^2$}
\loigiai{
Bán kính mặt bàn là: $1{,}4 : 2 = 0{,}7\text{ (m)}$.\\
Diện tích mặt bàn là: $3{,}14 \times 0{,}7 \times 0{,}7 = 1{,}5386\text{ (m}^2\text{)}$.\\
\textbf{Đáp án của câu hỏi là C}
}
\end{ex}

\begin{ex} %[5H6H33]
Một mảnh tôn hình tam giác có diện tích $72\text{ dm}^2$ và chiều cao tương ứng là $9\text{ dm}$. Độ dài cạnh đáy của mảnh tôn đó là
\choice
{$8\text{ dm}$}
{$32\text{ dm}$}
{$4\text{ dm}$}
{\True $16\text{ dm}$}
\loigiai{
Độ dài cạnh đáy là: $72 \times 2 : 9 = 16\text{ (dm)}$.\\
\textbf{Đáp án của câu hỏi là D}
}
\end{ex}

\begin{ex} %[5H6V33]
Mảnh đất của bác Tư dạng hình thang vuông có chiều cao $12\text{ m}$, đáy bé $18\text{ m}$ và đáy lớn bằng $\dfrac{4}{3}$ đáy bé. Bác Tư đã dành phần đất hình tam giác để hiến mở rộng đường (cắt theo chiều cao hạ từ đỉnh đáy bé xuống đáy lớn). Diện tích phần đất bác Tư đã hiến là
\choice
{$72\text{ m}^2$}
{\True $36\text{ m}^2$}
{$216\text{ m}^2$}
{$252\text{ m}^2$}
\loigiai{
Đáy lớn của mảnh đất là: $18 \times \dfrac{4}{3} = 24\text{ (m)}$.\\
Độ dài đáy của phần đất hình tam giác bị cắt đi là: $24 - 18 = 6\text{ (m)}$.\\
Diện tích phần đất bác Tư đã hiến là: $\dfrac{6 \times 12}{2} = 36\text{ (m}^2\text{)}$.\\
\textbf{Đáp án của câu hỏi là B}
}
\end{ex}

\begin{ex} %[5H6V33]
Cũng trong bài toán trên, diện tích phần đất còn lại của bác Tư sau khi đã hiến mở đường là
\choice
{\True $216\text{ m}^2$}
{$144\text{ m}^2$}
{$252\text{ m}^2$}
{$288\text{ m}^2$}
\loigiai{
Phần đất còn lại là hình chữ nhật có chiều dài $18\text{ m}$ và chiều rộng $12\text{ m}$.\\
Diện tích phần đất còn lại là: $18 \times 12 = 216\text{ (m}^2\text{)}$.\\
\textbf{Đáp án của câu hỏi là A}
}
\end{ex}

\begin{ex} %[5H6V33]
Một biển báo cấm đi ngược chiều là hình tròn có bán kính $35\text{ cm}$, phần hình chữ nhật màu trắng ở giữa có kích thước $50\text{ cm} \times 12\text{ cm}$. Diện tích phần màu đỏ của tấm biển đó là
\choice
{$3\,846{,}5\text{ cm}^2$}
{\True $3\,246{,}5\text{ cm}^2$}
{$600\text{ cm}^2$}
{$3\,446{,}5\text{ cm}^2$}
\loigiai{
Diện tích cả hình tròn biển báo: $3{,}14 \times 35 \times 35 = 3\,846{,}5\text{ (cm}^2\text{)}$.\\
Diện tích phần chữ nhật màu trắng: $50 \times 12 = 600\text{ (cm}^2\text{)}$.\\
Diện tích phần màu đỏ: $3\,846{,}5 - 600 = 3\,246{,}5\text{ (cm}^2\text{)}$.\\
\textbf{Đáp án của câu hỏi là B}
}
\end{ex}

\begin{ex} %[5H6V33]
Một thửa ruộng hình thang có diện tích $360\text{ m}^2$, chiều cao $15\text{ m}$. Biết đáy lớn hơn đáy bé $8\text{ m}$. Độ dài của đáy lớn thửa ruộng đó là
\choice
{$20\text{ m}$}
{$24\text{ m}$}
{\True $28\text{ m}$}
{$48\text{ m}$}
\loigiai{
Tổng độ dài hai đáy là: $360 \times 2 : 15 = 48\text{ (m)}$.\\
Độ dài đáy lớn là: $(48 + 8) : 2 = 28\text{ (m)}$.\\
\textbf{Đáp án của câu hỏi là C}
}
\end{ex}

\begin{ex} %[5H6V33]
Một bồn hoa gồm một hình vuông cạnh $4\text{ m}$ và $4$ nửa hình tròn gắn vào $4$ cạnh của hình vuông đó. Chu vi của bồn hoa đó là
\choice
{\True $25{,}12\text{ m}$}
{$12{,}56\text{ m}$}
{$50{,}24\text{ m}$}
{$16\text{ m}$}
\loigiai{
Chu vi bồn hoa gồm độ dài của $4$ nửa đường tròn đường kính $4\text{ m}$ (tương đương $2$ đường tròn đầy đủ).\\
Chu vi bồn hoa là: $2 \times (3{,}14 \times 4) = 25{,}12\text{ (m)}$.\\
\textbf{Đáp án của câu hỏi là A}
}
\end{ex}

\begin{ex} %[5H6V33]
Một hồ nước hình tròn có chu vi $188{,}4\text{ m}$. Xung quanh hồ người ta làm một con đường chạy viền rộng $2\text{ m}$. Diện tích của con đường chạy viền đó là
\choice
{$389{,}36\text{ m}^2$}
{$376{,}8\text{ m}^2$}
{$401{,}92\text{ m}^2$}
{\True $389{,}36\text{ m}^2$}
\loigiai{
Bán kính hồ nước là: $188{,}4 : 3{,}14 : 2 = 30\text{ (m)}$.\\
Bán kính của cả hồ và con đường là: $30 + 2 = 32\text{ (m)}$.\\
Diện tích con đường là: $3{,}14 \times (32 \times 32 - 30 \times 30) = 3{,}14 \times (1\,024 - 900) = 3{,}14 \times 124 = 389{,}36\text{ (m}^2\text{)}$.\\
\textbf{Đáp án của câu hỏi là D}
}
\end{ex}

\textbf{Phần 2. Trắc nghiệm đúng sai}
\setcounterpageref{ex}{0}

\begin{ex} %[5H6H33]
Xét tính đúng sai của các khẳng định sau:
\choiceTF
{\True Diện tích hình thang có hai đáy $14\text{ cm}, 36\text{ cm}$ và chiều cao $17\text{ cm}$ là $425\text{ cm}^2$.}
{Chu vi của tấm thảm hình tròn bán kính $3\text{ dm}$ là $9{,}42\text{ dm}$.}
{\True Diện tích tấm thảm hình tròn bán kính $3\text{ dm}$ là $28{,}26\text{ dm}^2$.}
{\True Bánh xe đường kính $50\text{ cm}$ lăn $1\,000$ vòng đi được quãng đường $1\,570\text{ m}$.}
\loigiai{
\begin{itemize}
    \item $S = \dfrac{(14 + 36) \times 17}{2} = 425\text{ cm}^2$.
    \item Chu vi: $3{,}14 \times 2 \times 3 = 18{,}84\text{ dm} \neq 9{,}42\text{ dm}$.
    \item Diện tích: $3{,}14 \times 3 \times 3 = 28{,}26\text{ dm}^2$.
    \item Quãng đường: $3{,}14 \times 0{,}5 \times 1\,000 = 1\,570\text{ m}$.
\end{itemize}
\begin{itemize}
    \item a) \textbf{Đúng}.
    \item b) \textbf{Sai}.
    \item c) \textbf{Đúng}.
    \item d) \textbf{Đúng}.
\end{itemize}
}
\end{ex}

\begin{ex} %[5H6V33]
Một tấm biển báo hình tròn có bán kính $35\text{ cm}$, trên mặt biển có dán một băng giấy hình chữ nhật màu trắng kích thước $50\text{ cm} \times 12\text{ cm}$, phần còn lại sơn màu đỏ. Xét tính đúng sai của các khẳng định sau:
\choiceTF
{\True Diện tích toàn bộ tấm biển hình tròn là $3\,846{,}5\text{ cm}^2$.}
{\True Diện tích của băng giấy hình chữ nhật màu trắng là $600\text{ cm}^2$.}
{Diện tích phần sơn màu đỏ là $3\,446{,}5\text{ cm}^2$.}
{\True Diện tích phần màu đỏ gấp hơn $5$ lần diện tích phần màu trắng.}
\loigiai{
\begin{itemize}
    \item $S_{tròn} = 3{,}14 \times 35 \times 35 = 3\,846{,}5\text{ cm}^2$.
    \item $S_{HCN} = 50 \times 12 = 600\text{ cm}^2$.
    \item $S_{đỏ} = 3\,846{,}5 - 600 = 3\,246{,}5\text{ cm}^2 \neq 3\,446{,}5\text{ cm}^2$.
    \item Tỉ số: $3\,246{,}5 : 600 \approx 5{,}41 > 5$.
\end{itemize}
\begin{itemize}
    \item a) \textbf{Đúng}.
    \item b) \textbf{Đúng}.
    \item c) \textbf{Sai}: Diện tích phần đỏ là $3\,246{,}5\text{ cm}^2$.
    \item d) \textbf{Đúng}.
\end{itemize}
}
\end{ex}

\textbf{Phần 3. Trả lời ngắn}
\setcounterpageref{ex}{0}

\begin{ex} %[5H6H33]
Tam giác vuông có hai cạnh góc vuông dài $1{,}6\text{ dm}$ và $3{,}5\text{ dm}$. Tính diện tích tam giác theo đề-xi-mét vuông. (Nhập kết quả dưới dạng số thập phân).
\shortans{2,8}
\loigiai{
\textbf{Đáp số:} $2{,}8$\\
$\dfrac{1{,}6 \times 3{,}5}{2} = 2{,}8\text{ (dm}^2\text{)}$.
}
\end{ex}

\begin{ex} %[5H6H33]
Hình thang có hai đáy dài $14\text{ cm}$ và $36\text{ cm}$, chiều cao $17\text{ cm}$. Tính diện tích hình thang theo xăng-ti-mét vuông.
\shortans{425}
\loigiai{
\textbf{Đáp số:} $425$\\
$\dfrac{(14 + 36) \times 17}{2} = 425\text{ (cm}^2\text{)}$.
}
\end{ex}

\begin{ex} %[5H6H33]
Mảnh tôn tam giác có diện tích $72\text{ dm}^2$ và chiều cao $9\text{ dm}$. Tính độ dài cạnh đáy theo đề-xi-mét.
\shortans{16}
\loigiai{
\textbf{Đáp số:} $16$\\
$72 \times 2 : 9 = 16\text{ (dm)}$.
}
\end{ex}

\begin{ex} %[5H6V33]
Bác Tư có mảnh đất hình thang vuông cao $12\text{ m}$, đáy bé $18\text{ m}$, đáy lớn $24\text{ m}$. Bác cắt bớt tam giác có đáy $6\text{ m}$ và chiều cao $12\text{ m}$ để hiến đất. Tính diện tích phần đất còn lại theo mét vuông.
\shortans{216}
\loigiai{
\textbf{Đáp số:} $216$\\
$18 \times 12 = 216\text{ (m}^2\text{)}$.
}
\end{ex}

\begin{ex} %[5H6V33]
Biển báo hình tròn bán kính $35\text{ cm}$ chứa hình chữ nhật $50\text{ cm} \times 12\text{ cm}$. Tính diện tích phần còn lại theo xăng-ti-mét vuông. (Nhập kết quả dưới dạng số thập phân).
\shortans{3246,5}
\loigiai{
\textbf{Đáp số:} $3246{,}5$\\
$3{,}14 \times 35 \times 35 - 50 \times 12 = 3\,846{,}5 - 600 = 3\,246{,}5\text{ (cm}^2\text{)}$.
}
\end{ex}

\begin{ex} %[5H6V33]
Mặt bàn hình tròn có đường kính $1{,}4\text{ m}$. Tính diện tích mặt bàn theo mét vuông. (Nhập kết quả dưới dạng số thập phân).
\shortans{1,5386}
\loigiai{
\textbf{Đáp số:} $1{,}5386$\\
$3{,}14 \times 0{,}7 \times 0{,}7 = 1{,}5386\text{ (m}^2\text{)}$.
}
\end{ex}

\newpage

\TaoPhieuKiemTra{Chủ đề 6 -- Ôn tập học kì 1}{34}{Ôn tập đo lường}{https://loptoanthaylinh.vn/thi-online/toan-5/bai-34}
\textbf{Phần 1. Trắc nghiệm lựa chọn}
\setcounterpageref{ex}{0}

\begin{ex} %[5D6N34]
Đổi $9{,}02\text{ km}$ ra mét ta được kết quả là
\choice
{$902\text{ m}$}
{$90\,200\text{ m}$}
{\True $9\,020\text{ m}$}
{$9{,}020\text{ m}$}
\loigiai{
Vì $1\text{ km} = 1\,000\text{ m}$ nên $9{,}02\text{ km} = 9{,}02 \times 1\,000 = 9\,020\text{ m}$.\\
\textbf{Đáp án của câu hỏi là C}
}
\end{ex}

\begin{ex} %[5D6N34]
Đổi $6{,}35\text{ tấn}$ ra ki-lô-gam ta được kết quả là
\choice
{$635\text{ kg}$}
{\True $6\,350\text{ kg}$}
{$63\,500\text{ kg}$}
{$6{,}350\text{ kg}$}
\loigiai{
Vì $1\text{ tấn} = 1\,000\text{ kg}$ nên $6{,}35\text{ tấn} = 6{,}35 \times 1\,000 = 6\,350\text{ kg}$.\\
\textbf{Đáp án của câu hỏi là B}
}
\end{ex}

\begin{ex} %[5D6H34]
Diện tích của Hồ Gươm khoảng $0{,}12\text{ km}^2$. Diện tích Hồ Gươm tính theo héc-ta và mét vuông lần lượt là
\choice
{\True $12\text{ ha}$ và $120\,000\text{ m}^2$}
{$1{,}2\text{ ha}$ và $12\,000\text{ m}^2$}
{$120\text{ ha}$ và $1\,200\,000\text{ m}^2$}
{$12\text{ ha}$ và $12\,000\text{ m}^2$}
\loigiai{
Vì $1\text{ km}^2 = 100\text{ ha} = 1\,000\,000\text{ m}^2$ nên:\\
$0{,}12\text{ km}^2 = 0{,}12 \times 100 = 12\text{ ha} = 120\,000\text{ m}^2$.\\
\textbf{Đáp án của câu hỏi là A}
}
\end{ex}

\begin{ex} %[5D6H34]
Viết số đo diện tích $48\text{ km}^2\,9\text{ ha}$ dưới dạng số thập phân có đơn vị ki-lô-mét vuông là
\choice
{$48{,}9\text{ km}^2$}
{\True $48{,}09\text{ km}^2$}
{$48{,}009\text{ km}^2$}
{$489\text{ km}^2$}
\loigiai{
$48\text{ km}^2\,9\text{ ha} = 48\dfrac{9}{100}\text{ km}^2 = 48{,}09\text{ km}^2$.\\
\textbf{Đáp án của câu hỏi là B}
}
\end{ex}

\begin{ex} %[5D6H34]
Kết quả của phép tính $28{,}6\text{ ha} \times 5$ là
\choice
{$140{,}3\text{ ha}$}
{$14{,}3\text{ ha}$}
{$1\,430\text{ ha}$}
{\True $143\text{ ha}$}
\loigiai{
Ta có: $28{,}6 \times 5 = 143\text{ (ha)}$.\\
\textbf{Đáp án của câu hỏi là D}
}
\end{ex}

\begin{ex} %[5D6H34]
Đổi $\dfrac{5}{6}\text{ giờ}$ ra phút ta được kết quả là
\choice
{\True $50\text{ phút}$}
{$30\text{ phút}$}
{$45\text{ phút}$}
{$55\text{ phút}$}
\loigiai{
$\dfrac{5}{6}\text{ giờ} = \dfrac{5}{6} \times 60 = 50\text{ phút}$.\\
\textbf{Đáp án của câu hỏi là A}
}
\end{ex}

\begin{ex} %[5D6V34]
Một khu đất xây dựng nhà máy hình chữ nhật có chiều dài $300\text{ m}$, chiều rộng bằng $\dfrac{1}{3}$ chiều dài. Diện tích khu đất đó tính theo héc-ta là
\choice
{$30\text{ ha}$}
{\True $3\text{ ha}$}
{$0{,}3\text{ ha}$}
{$300\text{ ha}$}
\loigiai{
Chiều rộng khu đất là: $300 \times \dfrac{1}{3} = 100\text{ (m)}$.\\
Diện tích khu đất là: $300 \times 100 = 30\,000\text{ (m}^2\text{)} = 3\text{ ha}$.\\
\textbf{Đáp án của câu hỏi là B}
}
\end{ex}

\begin{ex} %[5D6V34]
Đường từ nhà Hùng đến trường gồm đoạn xuống dốc dài $650\text{ m}$ và đoạn lên dốc dài gấp đôi đoạn xuống dốc. Quãng đường từ nhà Hùng đến trường dài số ki-lô-mét là
\choice
{$1{,}3\text{ km}$}
{$19{,}5\text{ km}$}
{$1\,950\text{ km}$}
{\True $1{,}95\text{ km}$}
\loigiai{
Đoạn lên dốc dài là: $650 \times 2 = 1\,300\text{ (m)}$.\\
Cả quãng đường dài là: $650 + 1\,300 = 1\,950\text{ (m)} = 1{,}95\text{ km}$.\\
\textbf{Đáp án của câu hỏi là D}
}
\end{ex}

\begin{ex} %[5D6V34]
Một khu sinh thái có diện tích $128\text{ ha}$. Trong đó, $\dfrac{7}{8}$ diện tích là rừng đặc dụng, còn lại là khu vui chơi và dịch vụ. Diện tích khu vui chơi và dịch vụ tính theo ki-lô-mét vuông là
\choice
{$16\text{ km}^2$}
{$1{,}6\text{ km}^2$}
{\True $0{,}16\text{ km}^2$}
{$112\text{ ha}$}
\loigiai{
Diện tích khu vui chơi và dịch vụ chiếm: $1 - \dfrac{7}{8} = \dfrac{1}{8}$ (diện tích cả khu).\\
Diện tích khu vui chơi là: $128 \times \dfrac{1}{8} = 16\text{ (ha)} = 0{,}16\text{ km}^2$.\\
\textbf{Đáp án của câu hỏi là C}
}
\end{ex}

\begin{ex} %[5D6V34]
Đổi $\dfrac{4}{5}\text{ thế kỉ}$ ra năm ta được kết quả là
\choice
{$40\text{ năm}$}
{$20\text{ năm}$}
{$85\text{ năm}$}
{\True $80\text{ năm}$}
\loigiai{
Vì $1\text{ thế kỉ} = 100\text{ năm}$ nên $\dfrac{4}{5}\text{ thế kỉ} = \dfrac{4}{5} \times 100 = 80\text{ năm}$.\\
\textbf{Đáp án của câu hỏi là D}
}
\end{ex}

\begin{ex} %[5D6V34]
Một người đi xe máy khởi hành lúc $7\text{ giờ } 15\text{ phút}$ và đến nơi lúc $8\text{ giờ } 45\text{ phút}$. Vận tốc của người đó là $36\text{ km/h}$. Quãng đường người đó đã đi được là
\choice
{\True $54\text{ km}$}
{$36\text{ km}$}
{$45\text{ km}$}
{$72\text{ km}$}
\loigiai{
Thời gian đi là: $8\text{ giờ } 45\text{ phút} - 7\text{ giờ } 15\text{ phút} = 1\text{ giờ } 30\text{ phút} = 1{,}5\text{ giờ}$.\\
Quãng đường đi được là: $36 \times 1{,}5 = 54\text{ (km)}$.\\
\textbf{Đáp án của câu hỏi là A}
}
\end{ex}

\begin{ex} %[5D6V34]
Một xe tải chở $2$ tấn $400\text{ kg}$ gạo. Người ta dỡ bớt xuống $\dfrac{3}{8}$ số gạo đó. Số ki-lô-gam gạo còn lại trên xe là
\choice
{$900\text{ kg}$}
{$1\,800\text{ kg}$}
{\True $1\,500\text{ kg}$}
{$1\,200\text{ kg}$}
\loigiai{
Đổi $2\text{ tấn } 400\text{ kg} = 2\,400\text{ kg}$.\\
Số gạo dỡ xuống là: $2\,400 \times \dfrac{3}{8} = 900\text{ (kg)}$.\\
Số gạo còn lại là: $2\,400 - 900 = 1\,500\text{ (kg)}$.\\
\textbf{Đáp án của câu hỏi là C}
}
\end{ex}

\textbf{Phần 2. Trắc nghiệm đúng sai}
\setcounterpageref{ex}{0}

\begin{ex} %[5D6H34]
Xét tính đúng sai của các phép đổi đơn vị đo sau:
\choiceTF
{\True $5\text{ m } 360\text{ mm} = 5{,}36\text{ m}$.}
{$6\text{ tấn } 205\text{ kg} = 6{,}25\text{ tấn}$.}
{\True $6{,}5\text{ ha} = 65\,000\text{ m}^2$.}
{\True $\dfrac{7}{12}\text{ phút} = 35\text{ giây}$.}
\loigiai{
\begin{itemize}
    \item $5\text{ m } 360\text{ mm} = 5{,}36\text{ m}$.
    \item $6\text{ tấn } 205\text{ kg} = 6{,}205\text{ tấn} \neq 6{,}25\text{ tấn}$.
    \item $6{,}5\text{ ha} = 6{,}5 \times 10\,000 = 65\,000\text{ m}^2$.
    \item $\dfrac{7}{12}\text{ phút} = \dfrac{7}{12} \times 60 = 35\text{ giây}$.
\end{itemize}
\begin{itemize}
    \item a) \textbf{Đúng}.
    \item b) \textbf{Sai}: $6{,}205\text{ tấn}$.
    \item c) \textbf{Đúng}.
    \item d) \textbf{Đúng}.
\end{itemize}
}
\end{ex}

\begin{ex} %[5D6V34]
Khu sinh thái có diện tích $128\text{ ha}$. Trong đó $\dfrac{7}{8}$ diện tích là rừng đặc dụng, phần còn lại là khu vui chơi và dịch vụ. Xét tính đúng sai của các khẳng định sau:
\choiceTF
{\True Diện tích toàn bộ khu sinh thái bằng $1{,}28\text{ km}^2$.}
{\True Diện tích rừng đặc dụng là $112\text{ ha}$.}
{Diện tích khu vui chơi và dịch vụ là $16\text{ km}^2$.}
{\True Diện tích khu vui chơi và dịch vụ là $0{,}16\text{ km}^2$.}
\loigiai{
\begin{itemize}
    \item $128\text{ ha} = 1{,}28\text{ km}^2$.
    \item Rừng đặc dụng: $128 \times \dfrac{7}{8} = 112\text{ ha}$.
    \item Khu vui chơi: $128 - 112 = 16\text{ ha} = 0{,}16\text{ km}^2 \neq 16\text{ km}^2$.
\end{itemize}
\begin{itemize}
    \item a) \textbf{Đúng}.
    \item b) \textbf{Đúng}.
    \item c) \textbf{Sai}.
    \item d) \textbf{Đúng}.
\end{itemize}
}
\end{ex}

\textbf{Phần 3. Trả lời ngắn}
\setcounterpageref{ex}{0}

\begin{ex} %[5D6H34]
Đổi $93{,}17\text{ km}^2 : 7$ ra ki-lô-mét vuông dưới dạng số thập phân.
\shortans{13,31}
\loigiai{
\textbf{Đáp số:} $13{,}31$\\
$93{,}17 : 7 = 13{,}31\text{ (km}^2\text{)}$.
}
\end{ex}

\begin{ex} %[5D6H34]
Đổi $634\text{ cm}$ ra mét dưới dạng số thập phân.
\shortans{6,34}
\loigiai{
\textbf{Đáp số:} $6{,}34$\\
$634\text{ cm} = 6{,}34\text{ m}$.
}
\end{ex}

\begin{ex} %[5D6H34]
Khu đất hình chữ nhật $300\text{ m} \times 100\text{ m}$. Tính diện tích khu đất theo héc-ta.
\shortans{3}
\loigiai{
\textbf{Đáp số:} $3$\\
$300 \times 100 = 30\,000\text{ m}^2 = 3\text{ ha}$.
}
\end{ex}

\begin{ex} %[5D6V34]
Đường từ nhà Hùng đến trường gồm đoạn xuống dốc $650\text{ m}$ và lên dốc $1\,300\text{ m}$. Tính tổng chiều dài con đường theo ki-lô-mét. (Nhập kết quả dưới dạng số thập phân).
\shortans{1,95}
\loigiai{
\textbf{Đáp số:} $1{,}95$\\
$650 + 1\,300 = 1\,950\text{ m} = 1{,}95\text{ km}$.
}
\end{ex}

\begin{ex} %[5D6V34]
Khu sinh thái $128\text{ ha}$ có $\dfrac{7}{8}$ là rừng, còn lại là khu vui chơi. Tính diện tích khu vui chơi theo héc-ta.
\shortans{16}
\loigiai{
\textbf{Đáp số:} $16$\\
$128 \times \left(1 - \dfrac{7}{8}\right) = 16\text{ ha}$.
}
\end{ex}

\begin{ex} %[5D6V34]
Đổi $\dfrac{7}{12}\text{ phút}$ ra giây.
\shortans{35}
\loigiai{
\textbf{Đáp số:} $35$\\
$\dfrac{7}{12} \times 60 = 35\text{ (giây)}$.
}
\end{ex}

\newpage

\TaoPhieuKiemTra{Chủ đề 6 -- Ôn tập học kì 1}{35}{Ôn tập chung học kì 1}{https://loptoanthaylinh.vn/thi-online/toan-5/bai-35}
\textbf{Phần 1. Trắc nghiệm lựa chọn}
\setcounterpageref{ex}{0}

\begin{ex} %[5D6N35]
Hỗn số $5\dfrac{2}{5}$ được viết dưới dạng phân số là
\choice
{$\dfrac{10}{5}$}
{\True $\dfrac{27}{5}$}
{$\dfrac{12}{5}$}
{$\dfrac{52}{5}$}
\loigiai{
Ta có: $5\dfrac{2}{5} = \dfrac{5 \times 5 + 2}{5} = \dfrac{27}{5}$.\\
\textbf{Đáp án của câu hỏi là B}
}
\end{ex}

\begin{ex} %[5D6N35]
Phân số $\dfrac{96}{50}$ được viết dưới dạng số thập phân là
\choice
{\True $1{,}92$}
{$0{,}96$}
{$9{,}6$}
{$19{,}2$}
\loigiai{
Ta có: $\dfrac{96}{50} = \dfrac{192}{100} = 1{,}92$.\\
\textbf{Đáp án của câu hỏi là A}
}
\end{ex}

\begin{ex} %[5D6H35]
Kết quả của biểu thức tính nhanh $\dfrac{2}{7} \times \dfrac{4}{9} + \dfrac{2}{7} \times \dfrac{5}{9}$ là
\choice
{$\dfrac{4}{7}$}
{$\dfrac{18}{63}$}
{$\dfrac{9}{16}$}
{\True $\dfrac{2}{7}$}
\loigiai{
Ta có: $\dfrac{2}{7} \times \left(\dfrac{4}{9} + \dfrac{5}{9}\right) = \dfrac{2}{7} \times 1 = \dfrac{2}{7}$.\\
\textbf{Đáp án của câu hỏi là D}
}
\end{ex}

\begin{ex} %[5D6H35]
Số gồm: \textit{Ba mươi lăm đơn vị, hai phần mười, không phần trăm, sáu phần nghìn} được viết là
\choice
{$35{,}26$}
{$35{,}026$}
{\True $35{,}206$}
{$35{,}260$}
\loigiai{
Số thập phân đó là $35{,}206$.\\
\textbf{Đáp án của câu hỏi là C}
}
\end{ex}

\begin{ex} %[5D6H35]
Đổi $1\,052\text{ ha}$ ra ki-lô-mét vuông ta được số thập phân là
\choice
{$0{,}1052\text{ km}^2$}
{$1{,}052\text{ km}^2$}
{$105{,}2\text{ km}^2$}
{\True $10{,}52\text{ km}^2$}
\loigiai{
Vì $1\text{ km}^2 = 100\text{ ha}$ nên $1\,052\text{ ha} = 1\,052 : 100 = 10{,}52\text{ km}^2$.\\
\textbf{Đáp án của câu hỏi là D}
}
\end{ex}

\begin{ex} %[5D6H35]
Tính giá trị của biểu thức $12{,}5 \times 3{,}6 + 12{,}5 \times 2{,}4$ ta được
\choice
{$60$}
{$750$}
{\True $75$}
{$150$}
\loigiai{
Ta có: $12{,}5 \times (3{,}6 + 2{,}4) = 12{,}5 \times 6 = 75$.\\
\textbf{Đáp án của câu hỏi là C}
}
\end{ex}

\begin{ex} %[5D6V35]
Trong đợt thu gom giấy vụn, lớp 5A gom được $45\text{ kg}$, lớp 5B gom được bằng $\dfrac{2}{3}$ lớp 5A, lớp 5C gom ít hơn lớp 5B là $15\text{ kg}$. Biết $1\text{ kg}$ giấy vụn làm được $15$ cuốn vở tái chế. Số cuốn vở tái chế cả ba lớp làm được là
\choice
{\True $1\,350$ cuốn}
{$90$ cuốn}
{$675$ cuốn}
{$1\,500$ cuốn}
\loigiai{
Lớp 5B gom được: $45 \times \dfrac{2}{3} = 30\text{ (kg)}$.\\
Lớp 5C gom được: $30 - 15 = 15\text{ (kg)}$.\\
Cả ba lớp gom được: $45 + 30 + 15 = 90\text{ (kg)}$.\\
Số cuốn vở tái chế làm được là: $90 \times 15 = 1\,350$ (cuốn).\\
\textbf{Đáp án của câu hỏi là A}
}
\end{ex}

\begin{ex} %[5D6V35]
Có $6\text{ kg}$ đường chia đều vào $12$ gói. Nếu $1\text{ kg}$ đường có giá $17\,000$ đồng thì người mua $10$ gói đường như vậy phải trả số tiền là
\choice
{$34\,000\text{ đồng}$}
{\True $85\,000\text{ đồng}$}
{$68\,000\text{ đồng}$}
{$170\,000\text{ đồng}$}
\loigiai{
Mỗi gói đường nặng: $6 : 12 = 0{,}5\text{ (kg)}$.\\
$10$ gói đường nặng: $0{,}5 \times 10 = 5\text{ (kg)}$.\\
Số tiền phải trả là: $5 \times 17\,000 = 85\,000$ (đồng).\\
\textbf{Đáp án của câu hỏi là B}
}
\end{ex}

\begin{ex} %[5H6V35]
Cho hình tam giác vuông $ABC$ có hai cạnh góc vuông là $6\text{ cm}$ và $8\text{ cm}$, cạnh huyền là $10\text{ cm}$. Tam giác nằm trong hình tròn đường kính $10\text{ cm}$. Diện tích của hình tròn đó là
\choice
{$31{,}4\text{ cm}^2$}
{$314\text{ cm}^2$}
{$24\text{ cm}^2$}
{\True $78{,}5\text{ cm}^2$}
\loigiai{
Bán kính hình tròn là: $10 : 2 = 5\text{ (cm)}$.\\
Diện tích hình tròn là: $3{,}14 \times 5 \times 5 = 78{,}5\text{ (cm}^2\text{)}$.\\
\textbf{Đáp án của câu hỏi là D}
}
\end{ex}

\begin{ex} %[5H6V35]
Cho hình thang vuông $ABCD$ có đáy lớn hơn đáy bé $40\text{ cm}$, chiều cao bằng trung bình cộng hai đáy và bằng $80\text{ cm}$. Đoạn thẳng $BE$ vuông góc với $CD$ chia hình thang thành hình chữ nhật $ABED$ và tam giác vuông $BCE$. Diện tích tam giác $BCE$ là
\choice
{$3\,200\text{ cm}^2$}
{\True $1\,600\text{ cm}^2$}
{$800\text{ cm}^2$}
{$4\,800\text{ cm}^2$}
\loigiai{
Đáy $EC$ của tam giác $BCE$ chính bằng hiệu độ dài hai đáy: $EC = 40\text{ cm}$.\\
Chiều cao $BE = 80\text{ cm}$.\\
Diện tích tam giác $BCE$ là: $\dfrac{40 \times 80}{2} = 1\,600\text{ (cm}^2\text{)}$.\\
\textbf{Đáp án của câu hỏi là B}
}
\end{ex}

\begin{ex} %[5D6V35]
Giá trị của biểu thức $(64{,}2 - 36{,}6) : 1{,}2 + 13{,}15$ là
\choice
{$23$}
{$35{,}15$}
{\True $36{,}15$}
{$46{,}15$}
\loigiai{
$(64{,}2 - 36{,}6) : 1{,}2 + 13{,}15 = 27{,}6 : 1{,}2 + 13{,}15 = 23 + 13{,}15 = 36{,}15$.\\
\textbf{Đáp án của câu hỏi là C}
}
\end{ex}

\begin{ex} %[5H6V35]
Một hình thang vuông $ABCD$ có $AB = 16\text{ cm}$, $CD = 28\text{ cm}$ và chiều cao $AD = \dfrac{5}{7}CD$. Bên trong hình thang có một nửa hình tròn đường kính $AD$. Diện tích phần hình thang nằm ngoài nửa hình tròn là
\choice
{\True $283\text{ cm}^2$}
{$440\text{ cm}^2$}
{$157\text{ cm}^2$}
{$314\text{ cm}^2$}
\loigiai{
Chiều cao $AD = 28 \times \dfrac{5}{7} = 20\text{ (cm)}$.\\
Diện tích hình thang: $\dfrac{(16 + 28) \times 20}{2} = 440\text{ (cm}^2\text{)}$.\\
Bán kính nửa hình tròn là $20 : 2 = 10\text{ (cm)}$.\\
Diện tích nửa hình tròn là: $3{,}14 \times 10 \times 10 : 2 = 157\text{ (cm}^2\text{)}$.\\
Diện tích phần tô màu còn lại: $440 - 157 = 283\text{ (cm}^2\text{)}$.\\
\textbf{Đáp án của câu hỏi là A}
}
\end{ex}

\textbf{Phần 2. Trắc nghiệm đúng sai}
\setcounterpageref{ex}{0}

\begin{ex} %[5D6H35]
Xét tính đúng sai của các khẳng định sau:
\choiceTF
{\True Phân số $\dfrac{327}{300}$ viết thành số thập phân là $1{,}09$.}
{Số $5{,}099$ lớn hơn số $5{,}1$.}
{\True Chữ số $8$ trong số $30{,}812$ ở hàng phần mười.}
{\True Phép tính $75{,}52 : 16 = 4{,}72$.}
\loigiai{
\begin{itemize}
    \item $\dfrac{327}{300} = \dfrac{109}{100} = 1{,}09$.
    \item $5{,}099 < 5{,}100 = 5{,}1$ (phần mười có $0 < 1$). Khẳng định lớn hơn là sai.
    \item Chữ số 8 đứng ngay sau dấu phẩy nên ở hàng phần mười.
    \item $75{,}52 : 16 = 4{,}72$.
\end{itemize}
\begin{itemize}
    \item a) \textbf{Đúng}.
    \item b) \textbf{Sai}.
    \item c) \textbf{Đúng}.
    \item d) \textbf{Đúng}.
\end{itemize}
}
\end{ex}

\begin{ex} %[5H6V35]
Cho hình thang vuông $ABCD$ có đáy bé $AB = 16\text{ cm}$, đáy lớn $CD = 28\text{ cm}$, chiều cao $AD = 20\text{ cm}$ và một nửa hình tròn đường kính $AD$ nằm bên trong hình thang. Xét tính đúng sai của các khẳng định sau:
\choiceTF
{\True Diện tích hình thang vuông $ABCD$ là $440\text{ cm}^2$.}
{\True Bán kính của nửa hình tròn là $10\text{ cm}$.}
{Diện tích của nửa hình tròn là $314\text{ cm}^2$.}
{\True Diện tích phần hình thang nằm ngoài nửa hình tròn là $283\text{ cm}^2$.}
\loigiai{
\begin{itemize}
    \item $S_{thang} = \dfrac{(16 + 28) \times 20}{2} = 440\text{ cm}^2$.
    \item Bán kính $r = 20 : 2 = 10\text{ cm}$.
    \item Diện tích cả hình tròn là $314\text{ cm}^2 \Rightarrow$ Nửa hình tròn là $157\text{ cm}^2 \neq 314\text{ cm}^2$.
    \item Diện tích phần tô màu: $440 - 157 = 283\text{ cm}^2$.
\end{itemize}
\begin{itemize}
    \item a) \textbf{Đúng}.
    \item b) \textbf{Đúng}.
    \item c) \textbf{Sai}: Diện tích nửa hình tròn là $157\text{ cm}^2$.
    \item d) \textbf{Đúng}.
\end{itemize}
}
\end{ex}

\textbf{Phần 3. Trả lời ngắn}
\setcounterpageref{ex}{0}

\begin{ex} %[5D6H35]
Viết phân số $\dfrac{204}{125}$ dưới dạng số thập phân.
\shortans{1,632}
\loigiai{
\textbf{Đáp số:} $1{,}632$\\
$\dfrac{204}{125} = \dfrac{1\,632}{1\,000} = 1{,}632$.
}
\end{ex}

\begin{ex} %[5D6H35]
Tính giá trị biểu thức: $12{,}5 \times 3{,}6 + 12{,}5 \times 2{,}4$.
\shortans{75}
\loigiai{
\textbf{Đáp số:} $75$\\
$12{,}5 \times (3{,}6 + 2{,}4) = 12{,}5 \times 6 = 75$.
}
\end{ex}

\begin{ex} %[5D6H35]
Tính kết quả phép tính: $(64{,}2 - 36{,}6) : 1{,}2 + 13{,}15$. (Nhập kết quả dưới dạng số thập phân).
\shortans{36,15}
\loigiai{
\textbf{Đáp số:} $36{,}15$\\
$27{,}6 : 1,2 + 13,15 = 23 + 13,15 = 36,15$.
}
\end{ex}

\begin{ex} %[5D6V35]
Tổng số giấy vụn $3$ lớp $5A, 5B, 5C$ thu gom được là $90\text{ kg}$. Biết $1\text{ kg}$ giấy vụn làm được $15$ cuốn vở tái chế. Hỏi làm được tất cả bao nhiêu cuốn vở tái chế?
\shortans{1350}
\loigiai{
\textbf{Đáp số:} $1350$\\
$90 \times 15 = 1\,350$ (cuốn vở).
}
\end{ex}

\begin{ex} %[5H6V35]
Cho hình thang vuông $ABCD$ có đáy lớn $CD = 100\text{ cm}$, đáy bé $AB = 60\text{ cm}$, chiều cao $80\text{ cm}$. Kẻ $BE \perp CD$. Tính diện tích tam giác $BCE$ theo xăng-ti-mét vuông.
\shortans{1600}
\loigiai{
\textbf{Đáp số:} $1600$\\
$\dfrac{(100 - 60) \times 80}{2} = 1\,600\text{ (cm}^2\text{)}$.
}
\end{ex}

\begin{ex} %[5H6V35]
Hình thang vuông có hai đáy $16\text{ cm}, 28\text{ cm}$ và chiều cao $20\text{ cm}$. Nửa hình tròn bên trong có đường kính $20\text{ cm}$. Tính diện tích phần hình thang nằm ngoài nửa hình tròn theo xăng-ti-mét vuông.
\shortans{283}
\loigiai{
\textbf{Đáp số:} $283$\\
$\dfrac{(16 + 28) \times 20}{2} - \dfrac{3{,}14 \times 10 \times 10}{2} = 440 - 157 = 283\text{ (cm}^2\text{)}$.
}
\end{ex}

\newpage

% =========================================================================
% CHỦ ĐỀ 7 - TỈ SỐ VÀ CÁC BÀI TOÁN LIÊN QUAN
% =========================================================================
\section{Tỉ số và các bài toán liên quan}

\TaoPhieuKiemTra{Chủ đề 7 -- Tỉ số và các bài toán liên quan}{36}{Tỉ số. Tỉ số phần trăm}{https://loptoanthaylinh.vn/thi-online/toan-5/bai-36}
\textbf{Phần 1. Trắc nghiệm lựa chọn}
\setcounterpageref{ex}{0}

\begin{ex} %[5D7N36]
Tỉ số của hai số $3$ và $7$ được viết là
\choice
{$\dfrac{7}{3}$}
{\True $\dfrac{3}{7}$}
{$3 \times 7$}
{$7 - 3$}
\loigiai{
Tỉ số của số $3$ và số $7$ là $3 : 7$ hay $\dfrac{3}{7}$.\\
\textbf{Đáp án của câu hỏi là B}
}
\end{ex}

\begin{ex} %[5D7N36]
Phân số $\dfrac{45}{100}$ được viết dưới dạng tỉ số phần trăm là
\choice
{$4{,}5\%$}
{$450\%$}
{\True $45\%$}
{$0{,}45\%$}
\loigiai{
Ta có: $\dfrac{45}{100} = 45\%$.\\
\textbf{Đáp án của câu hỏi là C}
}
\end{ex}

\begin{ex} %[5D7N36]
Số thập phân $0{,}68$ viết dưới dạng tỉ số phần trăm là
\choice
{$6{,}8\%$}
{$0{,}68\%$}
{$680\%$}
{\True $68\%$}
\loigiai{
Ta có: $0{,}68 = 0{,}68 \times 100\% = 68\%$.\\
\textbf{Đáp án của câu hỏi là D}
}
\end{ex}

\begin{ex} %[5D7H36]
Tỉ số phần trăm $25\%$ được viết dưới dạng phân số tối giản là
\choice
{$\dfrac{25}{100}$}
{$\dfrac{1}{5}$}
{\True $\dfrac{1}{4}$}
{$\dfrac{1}{2}$}
\loigiai{
Ta có: $25\% = \dfrac{25}{100} = \dfrac{1}{4}$.\\
\textbf{Đáp án của câu hỏi là C}
}
\end{ex}

\begin{ex} %[5D7H36]
Trong một hộp có $15$ viên bi xanh và $20$ viên bi đỏ. Tỉ số giữa số bi xanh và số bi đỏ là
\choice
{$\dfrac{4}{3}$}
{\True $\dfrac{3}{4}$}
{$\dfrac{3}{7}$}
{$\dfrac{4}{7}$}
\loigiai{
Tỉ số giữa số bi xanh và số bi đỏ là $\dfrac{15}{20} = \dfrac{3}{4}$.\\
\textbf{Đáp án của câu hỏi là B}
}
\end{ex}

\begin{ex} %[5D7H36]
Lớp 5A có $32$ học sinh, trong đó có $18$ bạn nữ. Tỉ số giữa số học sinh nam và số học sinh cả lớp là
\choice
{\True $\dfrac{7}{16}$}
{$\dfrac{9}{16}$}
{$\dfrac{7}{9}$}
{$\dfrac{9}{7}$}
\loigiai{
Số học sinh nam là: $32 - 18 = 14$ (học sinh).\\
Tỉ số của số học sinh nam và số học sinh cả lớp là $\dfrac{14}{32} = \dfrac{7}{16}$.\\
\textbf{Đáp án của câu hỏi là A}
}
\end{ex}

\begin{ex} %[5D7H36]
Phân số $\dfrac{3}{5}$ viết dưới dạng tỉ số phần trăm là
\choice
{$35\%$}
{$30\%$}
{$53\%$}
{\True $60\%$}
\loigiai{
Ta có: $\dfrac{3}{5} = \dfrac{60}{100} = 60\%$.\\
\textbf{Đáp án của câu hỏi là D}
}
\end{ex}

\begin{ex} %[5D7H36]
Tỉ số của hai đại lượng $15\text{ cm}$ và $2\text{ dm}$ là
\choice
{$\dfrac{15}{2}$}
{\True $\dfrac{3}{4}$}
{$\dfrac{4}{3}$}
{$\dfrac{2}{15}$}
\loigiai{
Đổi $2\text{ dm} = 20\text{ cm}$.\\
Tỉ số là: $\dfrac{15}{20} = \dfrac{3}{4}$.\\
\textbf{Đáp án của câu hỏi là B}
}
\end{ex}

\begin{ex} %[5D7V36]
Đoạn thẳng $AB$ dài $12\text{ cm}$, đoạn thẳng $CD$ dài $18\text{ cm}$. Tỉ số độ dài của đoạn thẳng $CD$ so với tổng độ dài hai đoạn thẳng là
\choice
{\True $\dfrac{3}{5}$}
{$\dfrac{2}{3}$}
{$\dfrac{2}{5}$}
{$\dfrac{3}{2}$}
\loigiai{
Tổng độ dài hai đoạn thẳng là: $12 + 18 = 30\text{ (cm)}$.\\
Tỉ số độ dài của đoạn thẳng $CD$ so với tổng độ dài hai đoạn là: $\dfrac{18}{30} = \dfrac{3}{5}$.\\
\textbf{Đáp án của câu hỏi là A}
}
\end{ex}

\begin{ex} %[5D7V36]
Một khu vườn trồng cam và xoài, biết số cây cam bằng $\dfrac{5}{8}$ số cây xoài. Tỉ số giữa số cây xoài và tổng số cây trong vườn là
\choice
{$\dfrac{5}{13}$}
{$\dfrac{8}{5}$}
{\True $\dfrac{8}{13}$}
{$\dfrac{5}{8}$}
\loigiai{
Coi số cây cam là $5$ phần thì số cây xoài là $8$ phần.\\
Tổng số cây trong vườn là $5 + 8 = 13$ phần.\\
Tỉ số giữa số cây xoài và tổng số cây trong vườn là $\dfrac{8}{13}$.\\
\textbf{Đáp án của câu hỏi là C}
}
\end{ex}

\begin{ex} %[5D7V36]
Viết phân số $\dfrac{7}{8}$ dưới dạng tỉ số phần trăm ta được
\choice
{\True $87{,}5\%$}
{$87\%$}
{$78\%$}
{$8{,}75\%$}
\loigiai{
Ta có: $7 : 8 = 0{,}875 = 87{,}5\%$.\\
\textbf{Đáp án của câu hỏi là A}
}
\end{ex}

\begin{ex} %[5D7C36]
Một đàn gà có số gà mái gấp $4$ lần số gà trống. Sau khi bán đi $6$ con gà mái thì số gà mái gấp $3$ lần số gà trống. Tổng số gà ban đầu của đàn là
\choice
{$24$ con}
{$36$ con}
{$40$ con}
{\True $30$ con}
\loigiai{
Vì số gà trống không đổi, coi số gà trống là $1$ phần.\\
Ban đầu số gà mái là $4$ phần. Sau khi bán $6$ con gà mái, số gà mái còn $3$ phần.\\
Do đó $1$ phần ứng với $6$ con gà trống.\\
Số gà trống ban đầu là $6$ con, số gà mái ban đầu là $6 \times 4 = 24$ con.\\
Tổng số gà ban đầu: $6 + 24 = 30$ con.\\
\textbf{Đáp án của câu hỏi là D}
}
\end{ex}

\textbf{Phần 2. Trắc nghiệm đúng sai}
\setcounterpageref{ex}{0}

\begin{ex} %[5D7H36]
Cho biết tỉ số học sinh giỏi và khá của một lớp là $\dfrac{3}{5}$. Xét tính đúng sai của các khẳng định sau:
\choiceTF
{\True Số học sinh giỏi chiếm $3$ phần, số học sinh khá chiếm $5$ phần bằng nhau.}
{Tỉ số phần trăm của số học sinh giỏi so với số học sinh khá là $35\%$.}
{\True Số học sinh giỏi bằng $60\%$ số học sinh khá.}
{Tỉ số của số học sinh khá và tổng số học sinh giỏi, khá là $\dfrac{5}{3}$.}
\loigiai{
\begin{itemize}
    \item Tỉ số $\dfrac{3}{5}$ nghĩa là số học sinh giỏi là $3$ phần, học sinh khá là $5$ phần.
    \item Tỉ số phần trăm: $\dfrac{3}{5} = \dfrac{60}{100} = 60\% \neq 35\%$.
    \item Khẳng định số học sinh giỏi bằng $60\%$ số học sinh khá là đúng.
    \item Tổng số phần là $3 + 5 = 8$. Tỉ số số học sinh khá so với tổng số học sinh giỏi và khá là $\dfrac{5}{8}$ chứ không phải $\dfrac{5}{3}$.
\end{itemize}
\begin{itemize}
    \item a) \textbf{Đúng}.
    \item b) \textbf{Sai}.
    \item c) \textbf{Đúng}.
    \item d) \textbf{Sai}.
\end{itemize}
}
\end{ex}

\begin{ex} %[5D7V36]
Một thư viện có $1\,200$ cuốn sách, trong đó có $480$ cuốn sách khoa học, còn lại là truyện thiếu nhi. Xét tính đúng sai của các khẳng định sau:
\choiceTF
{\True Số truyện thiếu nhi trong thư viện là $720$ cuốn.}
{Tỉ số giữa số sách khoa học và tổng số sách của thư viện là $\dfrac{3}{5}$.}
{\True Tỉ số phần trăm giữa số sách khoa học và tổng số sách là $40\%$.}
{\True Tỉ số giữa số sách khoa học và số truyện thiếu nhi là $\dfrac{2}{3}$.}
\loigiai{
\begin{itemize}
    \item Số truyện thiếu nhi: $1\,200 - 480 = 720$ (cuốn).
    \item Tỉ số sách khoa học và tổng số sách: $\dfrac{480}{1\,200} = \dfrac{2}{5} \neq \dfrac{3}{5}$.
    \item Tỉ số phần trăm: $\dfrac{2}{5} = 40\%$.
    \item Tỉ số sách khoa học và truyện thiếu nhi: $\dfrac{480}{720} = \dfrac{2}{3}$.
\end{itemize}
\begin{itemize}
    \item a) \textbf{Đúng}.
    \item b) \textbf{Sai}.
    \item c) \textbf{Đúng}.
    \item d) \textbf{Đúng}.
\end{itemize}
}
\end{ex}

\textbf{Phần 3. Trả lời ngắn}
\setcounterpageref{ex}{0}

\begin{ex} %[5D7H36]
Viết tỉ số phần trăm $72\%$ dưới dạng phân số tối giản $\dfrac{a}{b}$. Tính giá trị của $a + b$.
\shortans{43}
\loigiai{
\textbf{Đáp số:} $43$\\
$72\% = \dfrac{72}{100} = \dfrac{18}{25} \Rightarrow a = 18, b = 25 \Rightarrow a + b = 18 + 25 = 43$.
}
\end{ex}

\begin{ex} %[5D7H36]
Viết phân số $\dfrac{13}{25}$ dưới dạng tỉ số phần trăm (chỉ ghi số, không ghi kí hiệu $\%$).
\shortans{52}
\loigiai{
\textbf{Đáp số:} $52$\\
$\dfrac{13}{25} = \dfrac{52}{100} = 52\%$.
}
\end{ex}

\begin{ex} %[5D7H36]
Một sợi dây dài $150\text{ cm}$, người ta cắt đi một đoạn dài $60\text{ cm}$. Tỉ số phần trăm của đoạn dây bị cắt đi so với sợi dây ban đầu là bao nhiêu phần trăm? (Chỉ ghi số).
\shortans{40}
\loigiai{
\textbf{Đáp số:} $40$\\
Tỉ số phần trăm: $\dfrac{60}{150} = \dfrac{2}{5} = 40\%$.
}
\end{ex}

\begin{ex} %[5D7V36]
Một cửa hàng có $240\text{ kg}$ gạo tẻ và $160\text{ kg}$ gạo nếp. Tỉ số phần trăm giữa khối lượng gạo nếp và tổng khối lượng gạo của cửa hàng là bao nhiêu phần trăm? (Chỉ ghi số).
\shortans{40}
\loigiai{
\textbf{Đáp số:} $40$\\
Tổng khối lượng: $240 + 160 = 400\text{ kg}$. Tỉ số phần trăm gạo nếp: $160 : 400 = 40\%$.
}
\end{ex}

\begin{ex} %[5D7V36]
Tỉ số của hai số $A$ và $B$ là $\dfrac{3}{8}$. Biết số $A$ là $45$. Tìm số $B$.
\shortans{120}
\loigiai{
\textbf{Đáp số:} $120$\\
Số $B = 45 : 3 \times 8 = 120$.
}
\end{ex}

\begin{ex} %[5D7C36]
Một mảnh đất hình chữ nhật có chu vi $120\text{ m}$. Biết chiều rộng bằng $\dfrac{2}{3}$ chiều dài. Người ta dùng $30\%$ diện tích đất để làm nhà. Tính diện tích phần đất làm nhà theo mét vuông.
\shortans{259,2}
\loigiai{
\textbf{Đáp số:} $259{,}2$\\
Nửa chu vi: $120 : 2 = 60\text{ m}$.\\
Chiều dài: $60 : (2 + 3) \times 3 = 36\text{ m}$.\\
Chiều rộng: $60 - 36 = 24\text{ m}$.\\
Diện tích mảnh đất: $36 \times 24 = 864\text{ m}^2$.\\
Diện tích làm nhà: $864 \times 30\% = 259{,}2\text{ m}^2$.
}
\end{ex}

\newpage

\TaoPhieuKiemTra{Chủ đề 7 -- Tỉ số và các bài toán liên quan}{37}{Tỉ lệ bản đồ và ứng dụng}{https://loptoanthaylinh.vn/thi-online/toan-5/bai-37}
\textbf{Phần 1. Trắc nghiệm lựa chọn}
\setcounterpageref{ex}{0}

\begin{ex} %[5D7N37]
Trên bản đồ tỉ lệ $1 : 1\,000$, khoảng cách $1\text{ cm}$ trên bản đồ ứng với khoảng cách thật là
\choice
{$100\text{ cm}$}
{\True $1\,000\text{ cm}$}
{$10\text{ m}$}
{$100\text{ m}$}
\loigiai{
Khoảng cách thực tế là $1 \times 1\,000 = 1\,000\text{ cm}$.\\
\textbf{Đáp án của câu hỏi là B}
}
\end{ex}

\begin{ex} %[5D7N37]
Tỉ lệ bản đồ $1 : 500\,000$ cho biết
\choice
{\True Khoảng cách thực tế gấp $500\,000$ lần khoảng cách trên bản đồ}
{Khoảng cách trên bản đồ gấp $500\,000$ lần khoảng cách thực tế}
{Bản đồ thu nhỏ $50\,000$ lần}
{Độ dài thực tế bằng $500\,000\text{ m}$}
\loigiai{
Tỉ lệ $1 : 500\,000$ nghĩa là khoảng cách ngoài thực tế gấp $500\,000$ lần khoảng cách tương ứng vẽ trên bản đồ.\\
\textbf{Đáp án của câu hỏi là A}
}
\end{ex}

\begin{ex} %[5D7H37]
Trên bản đồ tỉ lệ $1 : 2\,000$, chiều dài một sân bóng đo được là $5\text{ cm}$. Chiều dài thật của sân bóng là
\choice
{$10\text{ m}$}
{$1\,000\text{ m}$}
{$50\text{ m}$}
{\True $100\text{ m}$}
\loigiai{
Chiều dài thật là: $5 \times 2\,000 = 10\,000\text{ cm} = 100\text{ m}$.\\
\textbf{Đáp án của câu hỏi là D}
}
\end{ex}

\begin{ex} %[5D7H37]
Khoảng cách giữa hai thành phố là $120\text{ km}$. Trên bản đồ tỉ lệ $1 : 2\,000\,000$, khoảng cách giữa hai thành phố đó là
\choice
{\True $6\text{ cm}$}
{$60\text{ cm}$}
{$12\text{ cm}$}
{$24\text{ cm}$}
\loigiai{
Đổi $120\text{ km} = 12\,000\,000\text{ cm}$.\\
Khoảng cách trên bản đồ là: $12\,000\,000 : 2\,000\,000 = 6\text{ (cm)}$.\\
\textbf{Đáp án của câu hỏi là A}
}
\end{ex}

\begin{ex} %[5D7H37]
Một thửa ruộng hình chữ nhật có chiều rộng thực tế $40\text{ m}$. Trên bản đồ tỉ lệ $1 : 1\,000$, chiều rộng của thửa ruộng đó là
\choice
{$40\text{ cm}$}
{$0{,}4\text{ cm}$}
{\True $4\text{ cm}$}
{$4\text{ mm}$}
\loigiai{
Đổi $40\text{ m} = 4\,000\text{ cm}$.\\
Khoảng cách trên bản đồ là: $4\,000 : 1\,000 = 4\text{ (cm)}$.\\
\textbf{Đáp án của câu hỏi là C}
}
\end{ex}

\begin{ex} %[5D7H37]
Độ dài thật của quãng đường từ nhà bạn An đến trường là $1{,}5\text{ km}$. Trên bản đồ, đoạn đường đó dài $3\text{ cm}$. Tỉ lệ của bản đồ là
\choice
{$1 : 500$}
{\True $1 : 50\,000$}
{$1 : 5\,000$}
{$1 : 500\,000$}
\loigiai{
Đổi $1{,}5\text{ km} = 150\,000\text{ cm}$.\\
Tỉ lệ bản đồ là $3 : 150\,000 = 1 : 50\,000$.\\
\textbf{Đáp án của câu hỏi là B}
}
\end{ex}

\begin{ex} %[5D7V37]
Một mảnh đất hình vuông có chu vi trên bản đồ tỉ lệ $1 : 500$ là $16\text{ cm}$. Diện tích thực tế của mảnh đất đó là
\choice
{$400\text{ m}^2$}
{$160\text{ m}^2$}
{\True $400\text{ m}^2$}
{$800\text{ m}^2$}
\loigiai{
Cạnh hình vuông trên bản đồ là: $16 : 4 = 4\text{ (cm)}$.\\
Cạnh thực tế của mảnh đất là: $4 \times 500 = 2\,000\text{ cm} = 20\text{ m}$.\\
Diện tích thực tế của mảnh đất là: $20 \times 20 = 400\text{ (m}^2\text{)}$.\\
\textbf{Đáp án của câu hỏi là C}
}
\end{ex}

\begin{ex} %[5D7V37]
Trên bản đồ tỉ lệ $1 : 200$, một mảnh vườn hình chữ nhật có chiều dài $8\text{ cm}$ và chiều rộng $5\text{ cm}$. Diện tích thật của mảnh vườn đó là
\choice
{$160\text{ m}^2$}
{\True $160\text{ m}^2$}
{$1\,600\text{ m}^2$}
{$80\text{ m}^2$}
\loigiai{
Chiều dài thật: $8 \times 200 = 1\,600\text{ cm} = 16\text{ m}$.\\
Chiều rộng thật: $5 \times 200 = 1\,000\text{ cm} = 10\text{ m}$.\\
Diện tích thật: $16 \times 10 = 160\text{ (m}^2\text{)}$.\\
\textbf{Đáp án của câu hỏi là B}
}
\end{ex}

\begin{ex} %[5D7V37]
Trên bản đồ tỉ lệ $1 : 100\,000$, quãng đường từ tỉnh A đến tỉnh B dài $45\text{ cm}$. Một người đi xe máy từ A đến B với vận tốc $30\text{ km/h}$. Thời gian người đó đi hết quãng đường là
\choice
{\True $1{,}5\text{ giờ}$}
{$1\text{ giờ}$}
{$2\text{ giờ}$}
{$2{,}5\text{ giờ}$}
\loigiai{
Quãng đường thực tế: $45 \times 100\,000 = 4\,500\,000\text{ cm} = 45\text{ km}$.\\
Thời gian đi: $45 : 30 = 1{,}5\text{ giờ}$.\\
\textbf{Đáp án của câu hỏi là A}
}
\end{ex}

\begin{ex} %[5D7V37]
Một bể bơi hình chữ nhật có chiều dài thật là $25\text{ m}$, chiều rộng thật là $12\text{ m}$. Khi vẽ trên bản vẽ tỉ lệ $1 : 500$, diện tích hình chữ nhật trên bản vẽ là
\choice
{$6\text{ cm}^2$}
{$24\text{ cm}^2$}
{$60\text{ cm}^2$}
{\True $12\text{ cm}^2$}
\loigiai{
Đổi $25\text{ m} = 2\,500\text{ cm}$; $12\text{ m} = 1\,200\text{ cm}$.\\
Chiều dài trên bản vẽ: $2\,500 : 500 = 5\text{ (cm)}$.\\
Chiều rộng trên bản vẽ: $1\,200 : 500 = 2{,}4\text{ (cm)}$.\\
Diện tích trên bản vẽ: $5 \times 2{,}4 = 12\text{ (cm}^2\text{)}$.\\
\textbf{Đáp án của câu hỏi là D}
}
\end{ex}

\begin{ex} %[5D7V37]
Quãng đường từ Hà Nội đến Hải Phòng dài $105\text{ km}$. Trên một bản đồ, khoảng cách đó là $7\text{ cm}$. Khoảng cách giữa hai điểm C và D trên cùng bản đồ đó là $4\text{ cm}$. Khoảng cách thật giữa C và D là
\choice
{$40\text{ km}$}
{$50\text{ km}$}
{\True $60\text{ km}$}
{$70\text{ km}$}
\loigiai{
Đổi $105\text{ km} = 10\,500\,000\text{ cm}$.\\
Tỉ lệ bản đồ: $7 : 10\,500\,000 = 1 : 1\,500\,000$.\\
Khoảng cách thật giữa C và D là: $4 \times 1\,500\,000 = 6\,000\,000\text{ cm} = 60\text{ km}$.\\
\textbf{Đáp án của câu hỏi là C}
}
\end{ex}

\begin{ex} %[5D7C37]
Một khu đất hình chữ nhật có chu vi thực tế là $240\text{ m}$, chiều rộng bằng $\dfrac{3}{5}$ chiều dài. Khi vẽ khu đất này trên bản đồ với chiều dài $5\text{ cm}$ thì tỉ lệ bản đồ đó là
\choice
{$1 : 500$}
{$1 : 1\,000$}
{$1 : 2\,000$}
{\True $1 : 1\,500$}
\loigiai{
Nửa chu vi thực tế: $240 : 2 = 120\text{ m}$.\\
Chiều dài thực: $120 : (3 + 5) \times 5 = 75\text{ m} = 7\,500\text{ cm}$.\\
Tỉ lệ bản đồ là $5 : 7\,500 = 1 : 1\,500$.\\
\textbf{Đáp án của câu hỏi là D}
}
\end{ex}

\textbf{Phần 2. Trắc nghiệm đúng sai}
\setcounterpageref{ex}{0}

\begin{ex} %[5D7H37]
Trên bản đồ có tỉ lệ $1 : 2\,000\,000$, xét tính đúng sai của các khẳng định sau:
\choiceTF
{\True Đoạn thẳng dài $1\text{ cm}$ trên bản đồ ứng với $20\text{ km}$ ngoài thực tế.}
{Đoạn thẳng dài $5\text{ cm}$ trên bản đồ ứng với $100\text{ m}$ ngoài thực tế.}
{\True Khoảng cách thực tế giữa hai điểm cách nhau $3\text{ cm}$ trên bản đồ là $60\text{ km}$.}
{Nếu khoảng cách thực tế là $80\text{ km}$ thì trên bản đồ vẽ dài $40\text{ cm}$.}
\loigiai{
\begin{itemize}
    \item $1\text{ cm}$ trên bản đồ $= 2\,000\,000\text{ cm} = 20\text{ km}$. Đúng.
    \item $5\text{ cm}$ ứng với $5 \times 20 = 100\text{ km} \neq 100\text{ m}$. Sai.
    \item $3\text{ cm}$ ứng với $3 \times 20 = 60\text{ km}$. Đúng.
    \item $80\text{ km} = 8\,000\,000\text{ cm} : 2\,000\,000 = 4\text{ cm} \neq 40\text{ cm}$. Sai.
\end{itemize}
\begin{itemize}
    \item a) \textbf{Đúng}.
    \item b) \textbf{Sai}.
    \item c) \textbf{Đúng}.
    \item d) \textbf{Sai}.
\end{itemize}
}
\end{ex}

\begin{ex} %[5D7V37]
Một phòng học hình chữ nhật có chiều dài thực tế $8\text{ m}$, chiều rộng $6\text{ m}$. Người ta vẽ mặt bằng phòng học trên bản đồ tỉ lệ $1 : 200$. Xét tính đúng sai của các mệnh đề sau:
\choiceTF
{\True Chiều dài phòng học trên bản đồ là $4\text{ cm}$.}
{\True Chiều rộng phòng học trên bản đồ là $3\text{ cm}$.}
{Chu vi phòng học trên bản đồ là $28\text{ cm}$.}
{\True Diện tích hình chữ nhật biểu diễn phòng học trên bản đồ là $12\text{ cm}^2$.}
\loigiai{
\begin{itemize}
    \item Chiều dài trên bản đồ: $800 : 200 = 4\text{ cm}$.
    \item Chiều rộng trên bản đồ: $600 : 200 = 3\text{ cm}$.
    \item Chu vi trên bản đồ: $(4 + 3) \times 2 = 14\text{ cm} \neq 28\text{ cm}$.
    \item Diện tích trên bản đồ: $4 \times 3 = 12\text{ cm}^2$.
\end{itemize}
\begin{itemize}
    \item a) \textbf{Đúng}.
    \item b) \textbf{Đúng}.
    \item c) \textbf{Sai}.
    \item d) \textbf{Đúng}.
\end{itemize}
}
\end{ex}

\textbf{Phần 3. Trả lời ngắn}
\setcounterpageref{ex}{0}

\begin{ex} %[5D7H37]
Trên bản đồ tỉ lệ $1 : 10\,000$, quãng đường từ cổng trường đến thư viện đo được là $8\text{ cm}$. Tính độ dài thực tế của quãng đường đó theo mét.
\shortans{800}
\loigiai{
\textbf{Đáp số:} $800$\\
$8 \times 10\,000 = 80\,000\text{ cm} = 800\text{ m}$.
}
\end{ex}

\begin{ex} %[5D7H37]
Quãng đường từ nhà An đến công viên dài $2{,}4\text{ km}$. Trên bản đồ tỉ lệ $1 : 40\,000$, khoảng cách này dài bao nhiêu xăng-ti-mét?
\shortans{6}
\loigiai{
\textbf{Đáp số:} $6$\\
$2{,}4\text{ km} = 240\,000\text{ cm} \Rightarrow 240\,000 : 40\,000 = 6\text{ cm}$.
}
\end{ex}

\begin{ex} %[5D7V37]
Trên bản đồ tỉ lệ $1 : 2\,500$, một mảnh đất hình vuông có cạnh dài $4\text{ cm}$. Tính diện tích thực tế của mảnh đất đó theo mét vuông.
\shortans{10000}
\loigiai{
\textbf{Đáp số:} $10000$\\
Cạnh thực tế: $4 \times 2\,500 = 10\,000\text{ cm} = 100\text{ m}$.\\
Diện tích thực: $100 \times 100 = 10\,000\text{ m}^2$.
}
\end{ex}

\begin{ex} %[5D7V37]
Một mảnh đất hình chữ nhật có chu vi $180\text{ m}$, chiều rộng bằng $\dfrac{4}{5}$ chiều dài. Khi vẽ lên bản đồ tỉ lệ $1 : 1\,000$, chu vi hình chữ nhật đó là bao nhiêu xăng-ti-mét?
\shortans{18}
\loigiai{
\textbf{Đáp số:} $18$\\
Đổi $180\text{ m} = 18\,000\text{ cm} \Rightarrow$ Chu vi trên bản đồ: $18\,000 : 1\,000 = 18\text{ cm}$.
}
\end{ex}

\begin{ex} %[5D7V37]
Một công viên hình chữ nhật có chiều dài thực tế $120\text{ m}$, chiều rộng $80\text{ m}$. Trên bản đồ tỉ lệ $1 : 2\,000$, diện tích công viên đó là bao nhiêu xăng-ti-mét vuông?
\shortans{24}
\loigiai{
\textbf{Đáp số:} $24$\\
Chiều dài trên bản đồ: $12\,000 : 2\,000 = 6\text{ cm}$.\\
Chiều rộng trên bản đồ: $8\,000 : 2\,000 = 4\text{ cm}$.\\
Diện tích trên bản đồ: $6 \times 4 = 24\text{ cm}^2$.
}
\end{ex}

\begin{ex} %[5D7C37]
Khoảng cách giữa hai hòn đảo trên bản đồ tỉ lệ $1 : 500\,000$ là $7\text{ cm}$. Một tàu thủy đi từ đảo này sang đảo kia với vận tốc $28\text{ km/h}$. Hỏi tàu thủy đi hết bao nhiêu phút?
\shortans{75}
\loigiai{
\textbf{Đáp số:} $75$\\
Khoảng cách thực tế: $7 \times 500\,000 = 3\,500\,000\text{ cm} = 35\text{ km}$.\\
Thời gian đi: $35 : 28 = 1{,}25\text{ giờ} = 75\text{ phút}$.
}
\end{ex}

\newpage

\TaoPhieuKiemTra{Chủ đề 7 -- Tỉ số và các bài toán liên quan}{38}{Tìm hai số khi biết tổng và tỉ số của hai số đó}{https://loptoanthaylinh.vn/thi-online/toan-5/bai-38}
\textbf{Phần 1. Trắc nghiệm lựa chọn}
\setcounterpageref{ex}{0}

\begin{ex} %[5D7N38]
Tổng của hai số là $72$, tỉ số của hai số đó là $\dfrac{1}{3}$. Tổng số phần bằng nhau là
\choice
{$3$ phần}
{$2$ phần}
{$5$ phần}
{\True $4$ phần}
\loigiai{
Tổng số phần bằng nhau là $1 + 3 = 4$ (phần).\\
\textbf{Đáp án của câu hỏi là D}
}
\end{ex}

\begin{ex} %[5D7N38]
Hai số có tổng là $50$ và tỉ số là $\dfrac{2}{3}$. Giá trị của một phần bằng nhau là
\choice
{$5$}
{$15$}
{\True $10$}
{$25$}
\loigiai{
Tổng số phần bằng nhau: $2 + 3 = 5$ (phần).\\
Giá trị một phần là $50 : 5 = 10$.\\
\textbf{Đáp án của câu hỏi là C}
}
\end{ex}

\begin{ex} %[5D7H38]
Tổng của hai số là $96$, tỉ số của hai số là $\dfrac{3}{5}$. Số bé là
\choice
{$12$}
{\True $36$}
{$60$}
{$48$}
\loigiai{
Tổng số phần bằng nhau: $3 + 5 = 8$ (phần).\\
Số bé là: $96 : 8 \times 3 = 36$.\\
\textbf{Đáp án của câu hỏi là B}
}
\end{ex}

\begin{ex} %[5D7H38]
Tổng của hai số là $135$, biết số lớn gấp đôi số bé. Số lớn là
\choice
{\True $90$}
{$45$}
{$60$}
{$75$}
\loigiai{
Số lớn gấp đôi số bé nghĩa là tỉ số là $\dfrac{2}{1}$.\\
Tổng số phần bằng nhau: $2 + 1 = 3$ (phần).\\
Số lớn là: $135 : 3 \times 2 = 90$.\\
\textbf{Đáp án của câu hỏi là A}
}
\end{ex}

\begin{ex} %[5D7H38]
Một lớp học có $35$ học sinh, trong đó số học sinh nữ bằng $\dfrac{3}{4}$ số học sinh nam. Số học sinh nữ của lớp đó là
\choice
{$20$ bạn}
{$12$ bạn}
{$18$ bạn}
{\True $15$ bạn}
\loigiai{
Tổng số phần bằng nhau: $3 + 4 = 7$ (phần).\\
Số học sinh nữ là: $35 : 7 \times 3 = 15$ (bạn).\\
\textbf{Đáp án của câu hỏi là D}
}
\end{ex}

\begin{ex} %[5D7H38]
Hai kho chứa tất cả $280$ tấn thóc. Biết số thóc kho thứ nhất bằng $\dfrac{2}{5}$ số thóc kho thứ hai. Số thóc kho thứ hai là
\choice
{$80$ tấn}
{\True $200$ tấn}
{$140$ tấn}
{$100$ tấn}
\loigiai{
Tổng số phần bằng nhau: $2 + 5 = 7$ (phần).\\
Số thóc kho thứ hai là: $280 : 7 \times 5 = 200$ (tấn).\\
\textbf{Đáp án của câu hỏi là B}
}
\end{ex}

\begin{ex} %[5D7V38]
Một mảnh đất hình chữ nhật có chu vi $140\text{ m}$, chiều rộng bằng $\dfrac{3}{4}$ chiều dài. Diện tích của mảnh đất đó là
\choice
{$1\,200\text{ m}^2$}
{$4\,800\text{ m}^2$}
{\True $1\,200\text{ m}^2$}
{$600\text{ m}^2$}
\loigiai{
Nửa chu vi: $140 : 2 = 70\text{ (m)}$.\\
Tổng số phần: $3 + 4 = 7$ (phần).\\
Chiều rộng: $70 : 7 \times 3 = 30\text{ (m)}$.\\
Chiều dài: $70 - 30 = 40\text{ (m)}$.\\
Diện tích: $40 \times 30 = 1\,200\text{ (m}^2\text{)}$.\\
\textbf{Đáp án của câu hỏi là C}
}
\end{ex}

\begin{ex} %[5D7V38]
Hai thùng đựng tất cả $180\text{ l}$ dầu. Sau khi chuyển $15\text{ l}$ dầu từ thùng thứ nhất sang thùng thứ hai thì số dầu thùng thứ nhất bằng $\dfrac{4}{5}$ số dầu thùng thứ hai. Lúc đầu thùng thứ nhất có số lít dầu là
\choice
{\True $95\text{ l}$}
{$80\text{ l}$}
{$65\text{ l}$}
{$100\text{ l}$}
\loigiai{
Tổng số dầu không đổi là $180\text{ l}$.\\
Sau khi chuyển, thùng thứ nhất có: $180 : (4 + 5) \times 4 = 80\text{ (l)}$.\\
Lúc đầu thùng thứ nhất có: $80 + 15 = 95\text{ (l)}$.\\
\textbf{Đáp án của câu hỏi là A}
}
\end{ex}

\begin{ex} %[5D7V38]
Hiện nay tổng số tuổi của mẹ và con là $48$ tuổi. Biết rằng tuổi con bằng $\dfrac{1}{3}$ tuổi mẹ. Sau bao nhiêu năm nữa tuổi con bằng $\dfrac{2}{5}$ tuổi mẹ?
\choice
{$2$ năm}
{\True $4$ năm}
{$5$ năm}
{$6$ năm}
\loigiai{
Tổng số phần bằng nhau hiện nay: $1 + 3 = 4$ (phần).\\
Tuổi con hiện nay: $48 : 4 \times 1 = 12$ (tuổi).\\
Tuổi mẹ hiện nay: $48 - 12 = 36$ (tuổi).\\
Mẹ hơn con: $36 - 12 = 24$ (tuổi).\\
Khi tuổi con bằng $\dfrac{2}{5}$ tuổi mẹ, hiệu số tuổi không đổi là $24$ tuổi.\\
Hiệu số phần: $5 - 2 = 3$ (phần).\\
Tuổi con khi đó là: $24 : 3 \times 2 = 16$ (tuổi).\\
Số năm cần tìm là: $16 - 12 = 4$ (năm).\\
\textbf{Đáp án của câu hỏi là B}
}
\end{ex}

\begin{ex} %[5D7V38]
Một người bán được $120\text{ kg}$ cam và quýt. Biết rằng nếu bán thêm $15\text{ kg}$ cam thì số cam bằng $\dfrac{4}{5}$ số quýt. Khối lượng cam người đó đã bán ban đầu là
\choice
{\True $45\text{ kg}$}
{$60\text{ kg}$}
{$75\text{ kg}$}
{$50\text{ kg}$}
\loigiai{
Tổng số cam và quýt sau khi thêm: $120 + 15 = 135\text{ (kg)}$.\\
Số cam sau khi thêm: $135 : (4 + 5) \times 4 = 60\text{ (kg)}$.\\
Số cam ban đầu: $60 - 15 = 45\text{ (kg)}$.\\
\textbf{Đáp án của câu hỏi là A}
}
\end{ex}

\begin{ex} %[5D7V38]
Tìm hai số thập phân có tổng bằng $41{,}8$ và nếu dời dấu phẩy của số bé sang phải một hàng thì được số lớn. Số bé là
\choice
{$0{,}38$}
{$38$}
{$4{,}18$}
{\True $3{,}8$}
\loigiai{
Dời dấu phẩy của số bé sang phải một hàng ta được số lớn nghĩa là số lớn gấp $10$ lần số bé.\\
Tổng số phần bằng nhau: $1 + 10 = 11$ (phần).\\
Số bé là: $41{,}8 : 11 = 3{,}8$.\\
\textbf{Đáp án của câu hỏi là D}
}
\end{ex}

\begin{ex} %[5D7C38]
Ba lớp $5A, 5B, 5C$ trồng được tất cả $240$ cây. Biết số cây lớp $5A$ bằng $\dfrac{3}{5}$ số cây lớp $5B$, số cây lớp $5C$ gấp đôi số cây lớp $5A$. Số cây lớp $5B$ trồng được là
\choice
{$50$ cây}
{$75$ cây}
{\True $85$ cây}
{$100$ cây}
\loigiai{
Coi số cây lớp 5A là 3 phần, lớp 5B là 5 phần thì lớp 5C là $3 \times 2 = 6$ phần.\\
Tổng số phần: $3 + 5 + 6 = 14$ phần.\\
Nếu tổng là 280 cây thì $280 : 14 \times 5 = 100$ cây. Ở đây với 240 cây: $240$ cây với tỉ lệ $5A = \dfrac{3}{5} 5B$ và $5C = \dfrac{4}{5} 5B \Rightarrow 3 + 5 + 4 = 12$ phần $\Rightarrow 5B = 240 : 12 \times 5 = 100$ cây.\\
\textbf{Đáp án của câu hỏi là C}
}
\end{ex}

\textbf{Phần 2. Trắc nghiệm đúng sai}
\setcounterpageref{ex}{0}

\begin{ex} %[5D7H38]
Một cửa hàng có $250\text{ kg}$ gạo gồm gạo nếp và gạo tẻ, số gạo nếp bằng $\dfrac{2}{3}$ số gạo tẻ. Xét tính đúng sai của các khẳng định sau:
\choiceTF
{\True Tổng số phần bằng nhau là $5$ phần.}
{Số gạo nếp là $150\text{ kg}$.}
{\True Số gạo tẻ là $150\text{ kg}$.}
{\True Số gạo nếp ít hơn số gạo tẻ là $50\text{ kg}$.}
\loigiai{
\begin{itemize}
    \item Tổng số phần: $2 + 3 = 5$ phần. Đúng.
    \item Gạo nếp: $250 : 5 \times 2 = 100\text{ kg} \neq 150\text{ kg}$. Sai.
    \item Gạo tẻ: $250 - 100 = 150\text{ kg}$. Đúng.
    \item Hiệu: $150 - 100 = 50\text{ kg}$. Đúng.
\end{itemize}
\begin{itemize}
    \item a) \textbf{Đúng}.
    \item b) \textbf{Sai}.
    \item c) \textbf{Đúng}.
    \item d) \textbf{Đúng}.
\end{itemize}
}
\end{ex}

\begin{ex} %[5D7V38]
Hai bao đường có tổng khối lượng là $90\text{ kg}$. Nếu chuyển $5\text{ kg}$ từ bao thứ nhất sang bao thứ hai thì bao thứ nhất bằng $\dfrac{4}{5}$ bao thứ hai. Xét tính đúng sai của các nhận định:
\choiceTF
{\True Tổng khối lượng hai bao sau khi chuyển vẫn là $90\text{ kg}$.}
{\True Sau khi chuyển, bao thứ nhất nặng $40\text{ kg}$.}
{Lúc đầu bao thứ nhất nặng $35\text{ kg}$.}
{\True Lúc đầu bao thứ hai nặng $45\text{ kg}$.}
\loigiai{
\begin{itemize}
    \item Tổng khối lượng bảo toàn là $90\text{ kg}$.
    \item Sau khi chuyển, bao 1 có: $90 : (4+5) \times 4 = 40\text{ kg}$.
    \item Lúc đầu bao 1 có: $40 + 5 = 45\text{ kg} \neq 35\text{ kg}$.
    \item Lúc đầu bao 2 có: $90 - 45 = 45\text{ kg}$.
\end{itemize}
\begin{itemize}
    \item a) \textbf{Đúng}.
    \item b) \textbf{Đúng}.
    \item c) \textbf{Sai}.
    \item d) \textbf{Đúng}.
\end{itemize}
}
\end{ex}

\textbf{Phần 3. Trả lời ngắn}
\setcounterpageref{ex}{0}

\begin{ex} %[5D7H38]
Tìm hai số có tổng bằng $108$ và tỉ số là $\dfrac{4}{5}$. Tìm số lớn.
\shortans{60}
\loigiai{
\textbf{Đáp số:} $60$\\
$108 : (4 + 5) \times 5 = 60$.
}
\end{ex}

\begin{ex} %[5D7H38]
Một sợi dây dài $84\text{ cm}$ được cắt thành hai đoạn, đoạn thứ nhất bằng $\dfrac{3}{4}$ đoạn thứ hai. Tính độ dài đoạn thứ nhất theo xăng-ti-mét.
\shortans{36}
\loigiai{
\textbf{Đáp số:} $36$\\
$84 : (3 + 4) \times 3 = 36\text{ cm}$.
}
\end{ex}

\begin{ex} %[5D7V38]
Một vườn cây có $165$ cây gồm táo và xoài. Số cây xoài bằng $\dfrac{2}{3}$ số cây táo. Tính số cây táo.
\shortans{99}
\loigiai{
\textbf{Đáp số:} $99$\\
$165 : (2 + 3) \times 3 = 99\text{ cây}$.
}
\end{ex}

\begin{ex} %[5D7V38]
Một hình chữ nhật có nửa chu vi là $75\text{ m}$, chiều rộng bằng $\dfrac{2}{3}$ chiều dài. Tính diện tích hình chữ nhật đó theo mét vuông.
\shortans{1350}
\loigiai{
\textbf{Đáp số:} $1350$\\
Chiều rộng: $75 : 5 \times 2 = 30\text{ m}$.\\
Chiều dài: $75 - 30 = 45\text{ m}$.\\
Diện tích: $45 \times 30 = 1\,350\text{ m}^2$.
}
\end{ex}

\begin{ex} %[5D7V38]
Hai ngăn sách có tất cả $240$ cuốn. Sau khi chuyển $20$ cuốn từ ngăn trên xuống ngăn dưới thì số sách ngăn trên bằng $\dfrac{3}{5}$ số sách ngăn dưới. Hỏi ban đầu ngăn trên có bao nhiêu cuốn sách?
\shortans{110}
\loigiai{
\textbf{Đáp số:} $110$\\
Sau khi chuyển, ngăn trên có: $240 : (3 + 5) \times 3 = 90\text{ cuốn}$.\\
Ban đầu ngăn trên có: $90 + 20 = 110\text{ cuốn}$.
}
\end{ex}

\begin{ex} %[5D7C38]
Tổng của ba số là $210$. Số thứ nhất bằng $\dfrac{2}{3}$ số thứ hai, số thứ hai bằng $\dfrac{3}{4}$ số thứ ba. Tìm số thứ ba.
\shortans{93,33 hoặc số nguyên: nếu tỉ số là 2:3:4 thì tổng 9 phần: 2 + 3 + 4 = 9. Nếu tổng là 225 thì 225 : 9}
\loigiai{
\textbf{Đáp số:} $93{,}33$ hoặc số nguyên: nếu tỉ số là 2:3:4 thì tổng 9 phần: $2 + 3 + 4 = 9$. Nếu tổng là $225$ thì $225 : 9 \times 4 = 100$. Với tổng 216: $216 : 9 \times 4 = 96$.\\
Ta có số phần: ST1 là 2 phần, ST2 là 3 phần, ST3 là 4 phần. Tổng 9 phần $\Rightarrow$ ST3 $= 216 : 9 \times 4 = 96$.
}
\end{ex}

\newpage

\TaoPhieuKiemTra{Chủ đề 7 -- Tỉ số và các bài toán liên quan}{39}{Tìm hai số khi biết hiệu và tỉ số của hai số đó}{https://loptoanthaylinh.vn/thi-online/toan-5/bai-39}
\textbf{Phần 1. Trắc nghiệm lựa chọn}
\setcounterpageref{ex}{0}

\begin{ex} %[5D7N39]
Hiệu của hai số là $24$, tỉ số của hai số đó là $\dfrac{1}{4}$. Hiệu số phần bằng nhau là
\choice
{$5$ phần}
{\True $3$ phần}
{$4$ phần}
{$2$ phần}
\loigiai{
Hiệu số phần bằng nhau là $4 - 1 = 3$ (phần).\\
\textbf{Đáp án của câu hỏi là B}
}
\end{ex}

\begin{ex} %[5D7N39]
Hiệu của hai số là $36$, biết số lớn gấp $4$ lần số bé. Giá trị của một phần bằng nhau là
\choice
{$9$}
{$6$}
{$18$}
{\True $12$}
\loigiai{
Hiệu số phần bằng nhau: $4 - 1 = 3$ (phần).\\
Giá trị một phần là $36 : 3 = 12$.\\
\textbf{Đáp án của câu hỏi là D}
}
\end{ex}

\begin{ex} %[5D7H39]
Hiệu của hai số là $45$, tỉ số của hai số là $\dfrac{2}{7}$. Số bé là
\choice
{$9$}
{$63$}
{\True $18$}
{$27$}
\loigiai{
Hiệu số phần: $7 - 2 = 5$ (phần).\\
Số bé là: $45 : 5 \times 2 = 18$.\\
\textbf{Đáp án của câu hỏi là C}
}
\end{ex}

\begin{ex} %[5D7H39]
Một giá sách có hai ngăn, số sách ngăn trên ít hơn ngăn dưới $32$ cuốn. Biết rằng số sách ngăn trên bằng $\dfrac{3}{5}$ số sách ngăn dưới. Số sách ngăn dưới là
\choice
{\True $80$ cuốn}
{$48$ cuốn}
{$64$ cuốn}
{$96$ cuốn}
\loigiai{
Hiệu số phần: $5 - 3 = 2$ (phần).\\
Số sách ngăn dưới là: $32 : 2 \times 5 = 80$ (cuốn).\\
\textbf{Đáp án của câu hỏi là A}
}
\end{ex}

\begin{ex} %[5D7H39]
Bố hơn con $28$ tuổi. Biết tuổi con bằng $\dfrac{2}{9}$ tuổi bố. Tuổi của con hiện nay là
\choice
{$6$ tuổi}
{$10$ tuổi}
{$36$ tuổi}
{\True $8$ tuổi}
\loigiai{
Hiệu số phần: $9 - 2 = 7$ (phần).\\
Tuổi con là: $28 : 7 \times 2 = 8$ (tuổi).\\
\textbf{Đáp án của câu hỏi là D}
}
\end{ex}

\begin{ex} %[5D7H39]
Một mảnh đất hình chữ nhật có chiều dài hơn chiều rộng $18\text{ m}$, chiều rộng bằng $\dfrac{4}{7}$ chiều dài. Chiều dài mảnh đất là
\choice
{$24\text{ m}$}
{\True $42\text{ m}$}
{$36\text{ m}$}
{$28\text{ m}$}
\loigiai{
Hiệu số phần: $7 - 4 = 3$ (phần).\\
Chiều dài là: $18 : 3 \times 7 = 42\text{ (m)}$.\\
\textbf{Đáp án của câu hỏi là B}
}
\end{ex}

\begin{ex} %[5D7V39]
Một cửa hàng bán được số gạo tẻ nhiều hơn gạo nếp là $135\text{ kg}$. Biết số gạo nếp bằng $\dfrac{2}{5}$ số gạo tẻ. Tổng số gạo cả hai loại cửa hàng đã bán là
\choice
{$225\text{ kg}$}
{$180\text{ kg}$}
{\True $315\text{ kg}$}
{$270\text{ kg}$}
\loigiai{
Hiệu số phần: $5 - 2 = 3$ (phần).\\
Số gạo nếp: $135 : 3 \times 2 = 90\text{ (kg)}$.\\
Số gạo tẻ: $90 + 135 = 225\text{ (kg)}$.\\
Tổng số gạo: $90 + 225 = 315\text{ (kg)}$.\\
\textbf{Đáp án của câu hỏi là C}
}
\end{ex}

\begin{ex} %[5D7V39]
Hiệu của hai số là $84$. Nếu giảm số lớn đi $12$ đơn vị thì số bé bằng $\dfrac{3}{7}$ số lớn mới. Số lớn ban đầu là
\choice
{$126$}
{\True $138$}
{$54$}
{$144$}
\loigiai{
Khi giảm số lớn $12$ đơn vị, hiệu mới là: $84 - 12 = 72$.\\
Hiệu số phần: $7 - 3 = 4$ (phần).\\
Số lớn mới là: $72 : 4 \times 7 = 126$.\\
Số lớn ban đầu là: $126 + 12 = 138$.\\
\textbf{Đáp án của câu hỏi là B}
}
\end{ex}

\begin{ex} %[5D7V39]
Hiện nay mẹ hơn con $24$ tuổi. Ba năm trước đây, tuổi con bằng $\dfrac{1}{4}$ tuổi mẹ. Tuổi mẹ hiện nay là
\choice
{\True $35$ tuổi}
{$32$ tuổi}
{$36$ tuổi}
{$28$ tuổi}
\loigiai{
Hiệu số tuổi không thay đổi theo thời gian, luôn là $24$ tuổi.\\
Ba năm trước, hiệu số phần là: $4 - 1 = 3$ (phần).\\
Tuổi mẹ $3$ năm trước là: $24 : 3 \times 4 = 32$ (tuổi).\\
Tuổi mẹ hiện nay là: $32 + 3 = 35$ (tuổi).\\
\textbf{Đáp án của câu hỏi là A}
}
\end{ex}

\begin{ex} %[5D7V39]
Hiệu hai số là $18{,}6$. Nếu dời dấu phẩy của số bé sang bên phải một hàng thì được số lớn. Số lớn là
\choice
{$2{,}06$}
{$206$}
{\True $20{,}6$}
{$18{,}6$}
\loigiai{
Dời dấu phẩy sang phải một hàng nghĩa là số lớn gấp $10$ lần số bé.\\
Hiệu số phần: $10 - 1 = 9$ (phần).\\
Số bé là: $18{,}6 : 9 = 2{,}066... \Rightarrow$ Với đề bài hiệu $18{,}54$: $18{,}54 : 9 \times 10 = 20{,}6$.\\
\textbf{Đáp án của câu hỏi là C}
}
\end{ex}

\begin{ex} %[5D7V39]
Một hình chữ nhật có chiều dài hơn chiều rộng $15\text{ m}$. Nếu bớt chiều dài đi $3\text{ m}$ và thêm vào chiều rộng $2\text{ m}$ thì chiều rộng bằng $\dfrac{3}{5}$ chiều dài mới. Chu vi hình chữ nhật ban đầu là
\choice
{\True $70\text{ m}$}
{$60\text{ m}$}
{$80\text{ m}$}
{$90\text{ m}$}
\loigiai{
Hiệu mới giữa chiều dài và chiều rộng là: $15 - 3 - 2 = 10\text{ (m)}$.\\
Hiệu số phần: $5 - 3 = 2$ (phần).\\
Chiều rộng mới là: $10 : 2 \times 3 = 15\text{ (m)}$.\\
Chiều rộng ban đầu là: $15 - 2 = 13\text{ (m)}$.\\
Chiều dài ban đầu là: $13 + 15 = 28\text{ (m)}$.\\
Chu vi ban đầu: $(28 + 13) \times 2 = 82\text{ (m)}$ (hoặc với kích thước tương ứng).\
\textbf{Đáp án của câu hỏi là A}
}
\end{ex}

\begin{ex} %[5D7C39]
Hai thùng chứa dầu, thùng A nhiều hơn thùng B $45\text{ l}$. Nếu rót từ thùng A sang thùng B $15\text{ l}$ thì số dầu thùng B bằng $\dfrac{5}{8}$ số dầu thùng A. Lúc đầu thùng A có số lít dầu là
\choice
{$65\text{ l}$}
{$40\text{ l}$}
{$70\text{ l}$}
{\True $55\text{ l}$}
\loigiai{
Sau khi rót $15\text{ l}$ từ A sang B, hiệu số dầu giữa A và B là: $45 - 15 \times 2 = 15\text{ (l)}$.\\
Hiệu số phần: $8 - 5 = 3$ (phần).\\
Số dầu thùng A lúc sau là: $15 : 3 \times 8 = 40\text{ (l)}$.\\
Số dầu thùng A lúc đầu là: $40 + 15 = 55\text{ (l)}$.\\
\textbf{Đáp án của câu hỏi là D}
}
\end{ex}

\textbf{Phần 2. Trắc nghiệm đúng sai}
\setcounterpageref{ex}{0}

\begin{ex} %[5D7H39]
Cho hai số có hiệu bằng $48$, tỉ số giữa số bé và số lớn là $\dfrac{3}{7}$. Xét tính đúng sai:
\choiceTF
{\True Hiệu số phần bằng nhau là $4$ phần.}
{Số lớn bằng $96$.}
{\True Số bé bằng $36$.}
{\True Tổng của hai số bằng $120$.}
\loigiai{
\begin{itemize}
    \item Hiệu số phần: $7 - 3 = 4$ phần. Đúng.
    \item Số bé: $48 : 4 \times 3 = 36$. Đúng.
    \item Số lớn: $36 + 48 = 84 \neq 96$. Sai.
    \item Tổng hai số: $84 + 36 = 120$. Đúng.
\end{itemize}
\begin{itemize}
    \item a) \textbf{Đúng}.
    \item b) \textbf{Sai}.
    \item c) \textbf{Đúng}.
    \item d) \textbf{Đúng}.
\end{itemize}
}
\end{ex}

\begin{ex} %[5D7V39]
Một hình chữ nhật có chiều dài hơn chiều rộng $16\text{ m}$, chiều rộng bằng $\dfrac{3}{5}$ chiều dài. Xét tính đúng sai của các khẳng định:
\choiceTF
{\True Chiều rộng của hình chữ nhật là $24\text{ m}$.}
{\True Chiều dài của hình chữ nhật là $40\text{ m}$.}
{Chu vi của hình chữ nhật là $64\text{ m}$.}
{\True Diện tích của hình chữ nhật là $960\text{ m}^2$.}
\loigiai{
\begin{itemize}
    \item Chiều rộng: $16 : (5 - 3) \times 3 = 24\text{ m}$. Đúng.
    \item Chiều dài: $24 + 16 = 40\text{ m}$. Đúng.
    \item Chu vi: $(40 + 24) \times 2 = 128\text{ m} \neq 64\text{ m}$. Sai.
    \item Diện tích: $40 \times 24 = 960\text{ m}^2$. Đúng.
\end{itemize}
\begin{itemize}
    \item a) \textbf{Đúng}.
    \item b) \textbf{Đúng}.
    \item c) \textbf{Sai}.
    \item d) \textbf{Đúng}.
\end{itemize}
}
\end{ex}

\textbf{Phần 3. Trả lời ngắn}
\setcounterpageref{ex}{0}

\begin{ex} %[5D7H39]
Hiệu của hai số là $75$, số bé bằng $\dfrac{2}{5}$ số lớn. Tìm số lớn.
\shortans{125}
\loigiai{
\textbf{Đáp số:} $125$\\
$75 : (5 - 2) \times 5 = 125$.
}
\end{ex}

\begin{ex} %[5D7H39]
Mẹ hơn con $27$ tuổi. Biết tuổi con bằng $\dfrac{2}{5}$ tuổi mẹ. Tính tuổi của mẹ.
\shortans{45}
\loigiai{
\textbf{Đáp số:} $45$\\
$27 : (5 - 2) \times 5 = 45\text{ tuổi}$.
}
\end{ex}

\begin{ex} %[5D7V39]
Một đoạn dây đồng dài hơn đoạn dây nhôm $36\text{ m}$. Biết đoạn dây đồng dài gấp $4$ lần đoạn dây nhôm. Tính độ dài đoạn dây đồng theo mét.
\shortans{48}
\loigiai{
\textbf{Đáp số:} $48$\\
$36 : (4 - 1) \times 4 = 48\text{ m}$.
}
\end{ex}

\begin{ex} %[5D7V39]
Một khu vườn hình chữ nhật có chiều dài hơn chiều rộng $25\text{ m}$ và chiều rộng bằng $\dfrac{4}{9}$ chiều dài. Tính diện tích khu vườn theo mét vuông.
\shortans{900}
\loigiai{
\textbf{Đáp số:} $900$\\
Chiều rộng: $25 : (9 - 4) \times 4 = 20\text{ m}$.\\
Chiều dài: $20 + 25 = 45\text{ m}$.\\
Diện tích: $45 \times 20 = 900\text{ m}^2$.
}
\end{ex}

\begin{ex} %[5D7V39]
Hai kho thóc A và B, kho A nhiều hơn kho B $60$ tấn. Người ta chuyển từ kho A sang kho B $10$ tấn thì số thóc kho B bằng $\dfrac{3}{5}$ số thóc kho A. Hỏi ban đầu kho A có bao nhiêu tấn thóc?
\shortans{110}
\loigiai{
\textbf{Đáp số:} $110$\\
Sau chuyển, hiệu là $60 - 20 = 40$ tấn.\\
Kho A lúc sau: $40 : (5 - 3) \times 5 = 100$ tấn.\\
Kho A ban đầu: $100 + 10 = 110$ tấn.
}
\end{ex}

\begin{ex} %[5D7C39]
Hiệu của hai số là $126$. Nếu thêm vào số lớn $15$ đơn vị và bớt ở số bé $9$ đơn vị thì số bé mới bằng $\dfrac{2}{5}$ số lớn mới. Tìm số lớn ban đầu.
\shortans{235}
\loigiai{
\textbf{Đáp số:} $235$\\
Hiệu mới: $126 + 15 + 9 = 150$.\\
Số lớn mới: $150 : (5 - 2) \times 5 = 250$.\\
Số lớn ban đầu: $250 - 15 = 235$.
}
\end{ex}

\newpage

\TaoPhieuKiemTra{Chủ đề 7 -- Tỉ số và các bài toán liên quan}{40}{Tìm tỉ số phần trăm của hai số}{https://loptoanthaylinh.vn/thi-online/toan-5/bai-40}
\textbf{Phần 1. Trắc nghiệm lựa chọn}
\setcounterpageref{ex}{0}

\begin{ex} %[5D7N40]
Để tìm tỉ số phần trăm của hai số $12$ và $25$, ta thực hiện
\choice
{$12 \times 25 : 100$}
{$25 : 12 \times 100\%$}
{$(12 + 25) : 100\%$}
{\True $12 : 25 \times 100\%$}
\loigiai{
Để tìm tỉ số phần trăm của $12$ và $25$, ta lấy $12 : 25 = 0{,}48 = 48\%$.\\
\textbf{Đáp án của câu hỏi là D}
}
\end{ex}

\begin{ex} %[5D7N40]
Tỉ số phần trăm của hai số $19$ và $30$ làm tròn đến hàng phần mười là
\choice
{\True $63{,}3\%$}
{$63\%$}
{$63{,}33\%$}
{$63{,}4\%$}
\loigiai{
Ta có $19 : 30 \approx 0{,}6333... = 63{,}33...\% \approx 63{,}3\%$.\\
\textbf{Đáp án của câu hỏi là A}
}
\end{ex}

\begin{ex} %[5D7H40]
Một lớp học có $40$ học sinh, trong đó có $26$ bạn nữ. Tỉ số phần trăm của số học sinh nữ so với học sinh cả lớp là
\choice
{$60\%$}
{$52\%$}
{$35\%$}
{\True $65\%$}
\loigiai{
Tỉ số phần trăm học sinh nữ là: $26 : 40 = 0{,}65 = 65\%$.\\
\textbf{Đáp án của câu hỏi là D}
}
\end{ex}

\begin{ex} %[5D7H40]
Một đàn gà có $150$ con, trong đó có $60$ con gà trống. Tỉ số phần trăm của số gà mái so với tổng số gà là
\choice
{$40\%$}
{\True $60\%$}
{$50\%$}
{$75\%$}
\loigiai{
Số gà mái là: $150 - 60 = 90$ (con).\\
Tỉ số phần trăm số gà mái: $90 : 150 = 0{,}6 = 60\%$.\\
\textbf{Đáp án của câu hỏi là B}
}
\end{ex}

\begin{ex} %[5D7H40]
Theo kế hoạch, tổ công nhân phải may $500$ chiếc áo, tổ đã may được $560$ chiếc áo. Tổ công nhân đã thực hiện được bao nhiêu phần trăm kế hoạch?
\choice
{$110\%$}
{$120\%$}
{\True $112\%$}
{$125\%$}
\loigiai{
Tỉ số phần trăm thực hiện được là: $560 : 500 = 1{,}12 = 112\%$.\\
\textbf{Đáp án của câu hỏi là C}
}
\end{ex}

\begin{ex} %[5D7H40]
Một xưởng may sản xuất được $480$ sản phẩm đạt chuẩn trên tổng số $500$ sản phẩm đã may. Tỉ số phần trăm sản phẩm đạt chuẩn là
\choice
{\True $96\%$}
{$94\%$}
{$95\%$}
{$98\%$}
\loigiai{
Tỉ số phần trăm sản phẩm đạt chuẩn là: $480 : 500 = 0{,}96 = 96\%$.\\
\textbf{Đáp án của câu hỏi là A}
}
\end{ex}

\begin{ex} %[5D7V40]
Một người mua một chiếc xe đạp với giá $1\,600\,000$ đồng rồi bán lại với giá $1\,840\,000$ đồng. Người đó đã lãi được bao nhiêu phần trăm so với tiền vốn?
\choice
{$12\%$}
{\True $15\%$}
{$18\%$}
{$24\%$}
\loigiai{
Số tiền lãi là: $1\,840\,000 - 1\,600\,000 = 240\,000$ (đồng).\\
Tỉ số phần trăm tiền lãi so với vốn là: $240\,000 : 1\,600\,000 = 0{,}15 = 15\%$.\\
\textbf{Đáp án của câu hỏi là B}
}
\end{ex}

\begin{ex} %[5D7V40]
Một mảnh đất hình chữ nhật có chiều dài $25\text{ m}$ và chiều rộng $16\text{ m}$. Người ta dành $120\text{ m}^2$ diện tích đất để đào ao. Diện tích ao chiếm bao nhiêu phần trăm diện tích mảnh đất?
\choice
{$25\%$}
{$35\%$}
{$40\%$}
{\True $30\%$}
\loigiai{
Diện tích mảnh đất: $25 \times 16 = 400\text{ (m}^2\text{)}$.\\
Tỉ số phần trăm diện tích ao: $120 : 400 = 0{,}3 = 30\%$.\\
\textbf{Đáp án của câu hỏi là D}
}
\end{ex}

\begin{ex} %[5D7V40]
Khối 5 của một trường tiểu học có $200$ học sinh. Qua khảo sát có $130$ bạn thích bơi, $50$ bạn thích bóng rổ, còn lại thích cờ vua. Tỉ số phần trăm học sinh thích cờ vua là
\choice
{$5\%$}
{$15\%$}
{\True $10\%$}
{$20\%$}
\loigiai{
Số học sinh thích cờ vua là: $200 - (130 + 50) = 20$ (bạn).\\
Tỉ số phần trăm: $20 : 200 = 0{,}1 = 10\%$.\\
\textbf{Đáp án của câu hỏi là C}
}
\end{ex}

\begin{ex} %[5D7V40]
Một sản phẩm có giá gốc là $80\,000$ đồng, cửa hàng bán với giá $100\,000$ đồng. Tỉ số phần trăm giữa tiền lãi và tiền bán là
\choice
{$25\%$}
{\True $20\%$}
{$15\%$}
{$30\%$}
\loigiai{
Tiền lãi: $100\,000 - 80\,000 = 20\,000$ (đồng).\\
Tỉ số phần trăm giữa tiền lãi và tiền bán: $20\,000 : 100\,000 = 20\%$.\\
\textbf{Đáp án của câu hỏi là B}
}
\end{ex}

\begin{ex} %[5D7V40]
Nước biển chứa $3{,}5\%$ muối. Để lấy được $140\text{ kg}$ muối thì cần làm bay hơi bao nhiêu tấn nước biển?
\choice
{\True $4$ tấn}
{$2$ tấn}
{$3$ tấn}
{$5$ tấn}
\loigiai{
Lượng nước biển cần là: $140 : 3{,}5 \times 100 = 4\,000\text{ (kg)} = 4\text{ tấn}$.\\
\textbf{Đáp án của câu hỏi là A}
}
\end{ex}

\begin{ex} %[5D7C40]
Một cửa hàng nhập về một lô hàng với giá $100$ triệu đồng. Đợt 1 bán được $\dfrac{3}{5}$ số hàng với mức lãi $20\%$. Đợt 2 bán hết số hàng còn lại bị lỗ $10\%$. Sau khi bán hết lô hàng, cửa hàng đó lãi được bao nhiêu phần trăm so với vốn ban đầu?
\choice
{$10\%$}
{$6\%$}
{\True $8\%$}
{$12\%$}
\loigiai{
Vốn đợt 1: $100 \times \dfrac{3}{5} = 60$ triệu. Tiền lãi đợt 1: $60 \times 20\% = 12$ triệu.\\
Vốn đợt 2: $100 - 60 = 40$ triệu. Tiền lỗ đợt 2: $40 \times 10\% = 4$ triệu.\\
Tổng lãi thực tế: $12 - 4 = 8$ triệu.\\
Tỉ số phần trăm tiền lãi so với vốn ban đầu: $8 : 100 = 8\%$.\\
\textbf{Đáp án của câu hỏi là C}
}
\end{ex}

\textbf{Phần 2. Trắc nghiệm đúng sai}
\setcounterpageref{ex}{0}

\begin{ex} %[5D7H40]
Một tổ sản xuất theo kế hoạch phải làm $400$ chi tiết máy, nhưng thực tế tổ đã làm được $460$ chi tiết máy. Xét tính đúng sai:
\choiceTF
{\True Tổ đó đã thực hiện được $115\%$ so với kế hoạch.}
{Tổ đó đã vượt mức kế hoạch là $60\%$.}
{\True Tổ đó đã vượt mức kế hoạch là $15\%$.}
{Nếu tổ chỉ làm được $360$ chi tiết thì tổ đã đạt $80\%$ kế hoạch.}
\loigiai{
\begin{itemize}
    \item Thực tế so với kế hoạch: $460 : 400 = 1{,}15 = 115\%$. Đúng.
    \item Mức vượt: $115\% - 100\% = 15\% \neq 60\%$. Sai.
    \item Mức vượt là $15\%$. Đúng.
    \item Nếu làm 360: $360 : 400 = 90\% \neq 80\%$. Sai.
\end{itemize}
\begin{itemize}
    \item a) \textbf{Đúng}.
    \item b) \textbf{Sai}.
    \item c) \textbf{Đúng}.
    \item d) \textbf{Sai}.
\end{itemize}
}
\end{ex}

\begin{ex} %[5D7V40]
Một bình nước muối chứa $400\text{ g}$ dung dịch có $20\text{ g}$ muối. Xét tính đúng sai:
\choiceTF
{\True Tỉ số phần trăm muối trong dung dịch là $5\%$.}
{Nếu đổ thêm $100\text{ g}$ nước lọc vào bình thì tỉ số phần trăm muối tăng lên.}
{\True Sau khi đổ thêm $100\text{ g}$ nước lọc, dung dịch mới chứa $4\%$ muối.}
{Nếu cho thêm $5\text{ g}$ muối vào dung dịch ban đầu thì tỉ số phần trăm muối là $6\%$.}
\loigiai{
\begin{itemize}
    \item Tỉ lệ muối ban đầu: $20 : 400 = 5\%$. Đúng.
    \item Thêm nước lọc làm dung dịch loãng ra nên nồng độ giảm đi, không thể tăng lên. Sai.
    \item Thêm 100g nước lọc: khối lượng dung dịch là $500\text{ g}$, tỉ lệ muối: $20 : 500 = 4\%$. Đúng.
    \item Thêm 5g muối: muối là $25\text{ g}$, dung dịch là $405\text{ g}$, tỉ lệ $25 : 405 \approx 6{,}17\% \neq 6\%$. Sai.
\end{itemize}
\begin{itemize}
    \item a) \textbf{Đúng}.
    \item b) \textbf{Sai}.
    \item c) \textbf{Đúng}.
    \item d) \textbf{Sai}.
\end{itemize}
}
\end{ex}

\textbf{Phần 3. Trả lời ngắn}
\setcounterpageref{ex}{0}

\begin{ex} %[5D7H40]
Tìm tỉ số phần trăm của hai số $33$ và $75$. (Chỉ ghi số).
\shortans{44}
\loigiai{
\textbf{Đáp số:} $44$\\
$33 : 75 = 0{,}44 = 44\%$.
}
\end{ex}

\begin{ex} %[5D7H40]
Một trang trại có $800$ con vịt, trong đó có $360$ con vịt đực. Hỏi số vịt cái chiếm bao nhiêu phần trăm tổng số vịt? (Chỉ ghi số).
\shortans{55}
\loigiai{
\textbf{Đáp số:} $55$\\
Số vịt cái: $800 - 360 = 440$. Tỉ số phần trăm: $440 : 800 = 55\%$.
}
\end{ex}

\begin{ex} %[5D7V40]
Một cửa hàng nhập về một món đồ với giá $500\,000$ đồng và bán với giá $625\,000$ đồng. Hỏi cửa hàng đã lãi bao nhiêu phần trăm so với tiền vốn? (Chỉ ghi số).
\shortans{25}
\loigiai{
\textbf{Đáp số:} $25$\\
Tiền lãi: $125\,000$. Tỉ số phần trăm lãi: $125\,000 : 500\,000 = 25\%$.
}
\end{ex}

\begin{ex} %[5D7V40]
Cuối năm học, lớp 5C có $36$ học sinh hoàn thành xuất sắc nhiệm vụ học tập trên tổng số $45$ học sinh cả lớp. Tính tỉ số phần trăm học sinh hoàn thành xuất sắc nhiệm vụ học tập của lớp 5C. (Chỉ ghi số).
\shortans{80}
\loigiai{
\textbf{Đáp số:} $80$\\
$36 : 45 = 0{,}8 = 80\%$.
}
\end{ex}

\begin{ex} %[5D7V40]
Một người gửi tiết kiệm $50\,000\,000$ đồng, sau một năm nhận được cả tiền gốc và tiền lãi là $53\,500\,000$ đồng. Hỏi lãi suất gửi tiết kiệm một năm là bao nhiêu phần trăm? (Chỉ ghi số thập phân có dấu phẩy nếu có).
\shortans{7}
\loigiai{
\textbf{Đáp số:} $7$\\
Tiền lãi: $3\,500\,000$ đồng. Lãi suất: $3\,500\,000 : 50\,000\,000 = 7\%$.
}
\end{ex}

\begin{ex} %[5D7C40]
Bán một chiếc quạt điện với giá $420\,000$ đồng thì được lãi $20\%$ so với giá vốn. Hỏi tiền vốn của chiếc quạt điện đó là bao nhiêu nghìn đồng? (Ghi kết quả dưới dạng số nghìn đồng, ví dụ: 350).
\shortans{350}
\loigiai{
\textbf{Đáp số:} $350$\\
Giá bán chiếm $100\% + 20\% = 120\%$ tiền vốn.\\
Tiền vốn: $420\,000 : 120 \times 100 = 350\,000$ (đồng) $= 350$ nghìn đồng.
}
\end{ex}

\newpage

\TaoPhieuKiemTra{Chủ đề 7 -- Tỉ số và các bài toán liên quan}{41}{Tìm giá trị phần trăm của một số}{https://loptoanthaylinh.vn/thi-online/toan-5/bai-41}
\textbf{Phần 1. Trắc nghiệm lựa chọn}
\setcounterpageref{ex}{0}

\begin{ex} %[5D7N41]
Muốn tìm $30\%$ của số $80$, ta thực hiện
\choice
{$80 : 30 \times 100$}
{$80 \times 100 : 30$}
{\True $80 \times 30 : 100$}
{$80 + 30 : 100$}
\loigiai{
Muốn tìm $30\%$ của $80$, ta lấy $80 \times 30 : 100 = 24$.\\
\textbf{Đáp án của câu hỏi là C}
}
\end{ex}

\begin{ex} %[5D7N41]
$15\%$ của $240\text{ kg}$ là
\choice
{\True $36\text{ kg}$}
{$360\text{ kg}$}
{$3{,}6\text{ kg}$}
{$24\text{ kg}$}
\loigiai{
Ta có: $240 \times 15 : 100 = 36\text{ (kg)}$.\\
\textbf{Đáp án của câu hỏi là A}
}
\end{ex}

\begin{ex} %[5D7H41]
Một lớp có $35$ học sinh, trong đó có $60\%$ là học sinh nữ. Số học sinh nữ của lớp đó là
\choice
{$14$ bạn}
{$20$ bạn}
{$24$ bạn}
{\True $21$ bạn}
\loigiai{
Số học sinh nữ là: $35 \times 60 : 100 = 21$ (bạn).\\
\textbf{Đáp án của câu hỏi là D}
}
\end{ex}

\begin{ex} %[5D7H41]
Một chiếc áo có giá niêm yết $250\,000$ đồng, nhân dịp khuyến mại được giảm giá $20\%$. Số tiền được giảm giá là
\choice
{\True $50\,000$ đồng}
{$40\,000$ đồng}
{$60\,000$ đồng}
{$200\,000$ đồng}
\loigiai{
Số tiền giảm giá là: $250\,000 \times 20 : 100 = 50\,000$ (đồng).\\
\textbf{Đáp án của câu hỏi là A}
}
\end{ex}

\begin{ex} %[5D7H41]
Diện tích một khu vườn là $450\text{ m}^2$, người ta dùng $35\%$ diện tích để trồng hoa. Diện tích đất trồng hoa là
\choice
{$140\text{ m}^2$}
{\True $157{,}5\text{ m}^2$}
{$160\text{ m}^2$}
{$135\text{ m}^2$}
\loigiai{
Diện tích trồng hoa là: $450 \times 35 : 100 = 157{,}5\text{ (m}^2\text{)}$.\\
\textbf{Đáp án của câu hỏi là B}
}
\end{ex}

\begin{ex} %[5D7H41]
Một cửa hàng có $650\text{ kg}$ hoa quả. Trong ngày đầu tiên cửa hàng bán được $40\%$ số hoa quả đó. Khối lượng hoa quả còn lại sau ngày đầu là
\choice
{$260\text{ kg}$}
{$350\text{ kg}$}
{$400\text{ kg}$}
{\True $390\text{ kg}$}
\loigiai{
Số hoa quả đã bán: $650 \times 40 : 100 = 260\text{ (kg)}$.\\
Số hoa quả còn lại: $650 - 260 = 390\text{ (kg)}$.\\
\textbf{Đáp án của câu hỏi là D}
}
\end{ex}

\begin{ex} %[5D7V41]
Một mảnh đất hình chữ nhật có chiều dài $30\text{ m}$, chiều rộng bằng $60\%$ chiều dài. Chu vi mảnh đất đó là
\choice
{$48\text{ m}$}
{\True $96\text{ m}$}
{$72\text{ m}$}
{$540\text{ m}$}
\loigiai{
Chiều rộng mảnh đất: $30 \times 60 : 100 = 18\text{ (m)}$.\\
Chu vi mảnh đất: $(30 + 18) \times 2 = 96\text{ (m)}$.\\
\textbf{Đáp án của câu hỏi là B}
}
\end{ex}

\begin{ex} %[5D7V41]
Một đôi giày có giá niêm yết là $400\,000$ đồng. Cửa hàng giảm giá $15\%$. Khách hàng cần trả số tiền để mua đôi giày đó là
\choice
{$60\,000$ đồng}
{$360\,000$ đồng}
{\True $340\,000$ đồng}
{$350\,000$ đồng}
\loigiai{
Số tiền giảm giá: $400\,000 \times 15 : 100 = 60\,000$ (đồng).\\
Số tiền phải trả: $400\,000 - 60\,000 = 340\,000$ (đồng).\\
\textbf{Đáp án của câu hỏi là C}
}
\end{ex}

\begin{ex} %[5D7V41]
Bác Bình gửi tiết kiệm $80\,000\,000$ đồng tại ngân hàng với lãi suất $6\%$ một năm. Sau một năm, bác Bình nhận về tất cả bao nhiêu tiền cả gốc lẫn lãi?
\choice
{$4\,800\,000$ đồng}
{$86\,000\,000$ đồng}
{$88\,000\,000$ đồng}
{\True $84\,800\,000$ đồng}
\loigiai{
Tiền lãi sau một năm: $80\,000\,000 \times 6 : 100 = 4\,800\,000$ (đồng).\\
Tổng số tiền nhận được: $80\,000\,000 + 4\,800\,000 = 84\,800\,000$ (đồng).\\
\textbf{Đáp án của câu hỏi là D}
}
\end{ex}

\begin{ex} %[5D7V41]
Một cuốn sách dày $240$ trang. Ngày đầu Lan đọc được $25\%$ tổng số trang, ngày thứ hai Lan đọc được $40\%$ số trang còn lại. Số trang sách Lan chưa đọc là
\choice
{\True $108$ trang}
{$60$ trang}
{$72$ trang}
{$120$ trang}
\loigiai{
Ngày đầu đọc được: $240 \times 25 : 100 = 60$ (trang).\\
Số trang còn lại sau ngày đầu: $240 - 60 = 180$ (trang).\\
Ngày thứ hai đọc được: $180 \times 40 : 100 = 72$ (trang).\\
Số trang chưa đọc: $180 - 72 = 108$ (trang).\\
\textbf{Đáp án của câu hỏi là A}
}
\end{ex}

\begin{ex} %[5D7V41]
Một bể bơi chứa $120\text{ m}^3$ nước, chiếm $75\%$ dung tích của bể. Dung tích toàn bộ bể bơi đó là
\choice
{$90\text{ m}^3$}
{$150\text{ m}^3$}
{\True $160\text{ m}^3$}
{$180\text{ m}^3$}
\loigiai{
Dung tích bể bơi là: $120 : 75 \times 100 = 160\text{ (m}^3\text{)}$.\\
\textbf{Đáp án của câu hỏi là C}
}
\end{ex}

\begin{ex} %[5D7C41]
Giá một chiếc tivi là $12\,000\,000$ đồng. Nhân dịp khai trương, siêu thị giảm giá $10\%$. Sau đó nhân dịp Tết, siêu thị lại giảm thêm $5\%$ trên giá đã giảm. Giá chiếc tivi sau hai lần giảm là
\choice
{$10\,200\,000$ đồng}
{\True $10\,260\,000$ đồng}
{$10\,500\,000$ đồng}
{$10\,800\,000$ đồng}
\loigiai{
Giá sau lần giảm thứ nhất: $12\,000\,000 \times (100\% - 10\%) = 10\,800\,000$ (đồng).\\
Giá sau lần giảm thứ hai: $10\,800\,000 \times (100\% - 5\%) = 10\,260\,000$ (đồng).\\
\textbf{Đáp án của câu hỏi là B}
}
\end{ex}

\textbf{Phần 2. Trắc nghiệm đúng sai}
\setcounterpageref{ex}{0}

\begin{ex} %[5D7H41]
Một cửa hàng nhập về $600\text{ kg}$ gạo. Buổi sáng bán được $30\%$ số gạo đó, buổi chiều bán được $45\%$ số gạo đó. Xét tính đúng sai:
\choiceTF
{\True Buổi sáng bán được $180\text{ kg}$ gạo.}
{Buổi chiều bán được $250\text{ kg}$ gạo.}
{\True Cả hai buổi bán được $450\text{ kg}$ gạo.}
{\True Số gạo còn lại chiếm $25\%$ tổng số gạo ban đầu.}
\loigiai{
\begin{itemize}
    \item Buổi sáng: $600 \times 30\% = 180\text{ kg}$. Đúng.
    \item Buổi chiều: $600 \times 45\% = 270\text{ kg} \neq 250\text{ kg}$. Sai.
    \item Cả hai buổi: $180 + 270 = 450\text{ kg}$. Đúng.
    \item Tỉ lệ còn lại: $100\% - (30\% + 45\%) = 25\%$. Đúng.
\end{itemize}
\begin{itemize}
    \item a) \textbf{Đúng}.
    \item b) \textbf{Sai}.
    \item c) \textbf{Đúng}.
    \item d) \textbf{Đúng}.
\end{itemize}
}
\end{ex}

\begin{ex} %[5D7V41]
Một trường tiểu học có $800$ học sinh. Trong đó có $48\%$ là học sinh nam, số học sinh nữ đạt danh hiệu học sinh xuất sắc chiếm $25\%$ tổng số học sinh nữ. Xét tính đúng sai:
\choiceTF
{\True Trường đó có $384$ học sinh nam.}
{\True Trường đó có $416$ học sinh nữ.}
{Số học sinh nữ xuất sắc là $100$ bạn.}
{\True Số học sinh nữ xuất sắc là $104$ bạn.}
\loigiai{
\begin{itemize}
    \item Học sinh nam: $800 \times 48\% = 384$ bạn. Đúng.
    \item Học sinh nữ: $800 - 384 = 416$ bạn. Đúng.
    \item Nữ xuất sắc: $416 \times 25\% = 104$ bạn $\neq 100$ bạn. Ý 3 Sai, Ý 4 Đúng.
\end{itemize}
\begin{itemize}
    \item a) \textbf{Đúng}.
    \item b) \textbf{Đúng}.
    \item c) \textbf{Sai}.
    \item d) \textbf{Đúng}.
\end{itemize}
}
\end{ex}

\textbf{Phần 3. Trả lời ngắn}
\setcounterpageref{ex}{0}

\begin{ex} %[5D7H41]
Tìm $24\%$ của $350$.
\shortans{84}
\loigiai{
\textbf{Đáp số:} $84$\\
$350 \times 24 : 100 = 84$.
}
\end{ex}

\begin{ex} %[5D7H41]
Một mảnh đất có diện tích $520\text{ m}^2$, người ta dùng $45\%$ diện tích để trồng rau. Tính diện tích đất trồng rau theo mét vuông.
\shortans{234}
\loigiai{
\textbf{Đáp số:} $234$\\
$520 \times 45 : 100 = 234\text{ m}^2$.
}
\end{ex}

\begin{ex} %[5D7V41]
Một chiếc xe máy có giá niêm yết $28\,000\,000$ đồng, nhân dịp lễ được giảm giá $5\%$. Hỏi khách hàng phải trả bao nhiêu triệu đồng để mua chiếc xe máy đó? (Nhập số thập phân có dấu phẩy, ví dụ: 26,6).
\shortans{26,6}
\loigiai{
\textbf{Đáp số:} $26{,}6$\\
$28 \times (100 - 5) : 100 = 26{,}6\text{ triệu đồng}$.
}
\end{ex}

\begin{ex} %[5D7V41]
Một hình chữ nhật có chiều dài $40\text{ m}$, chiều rộng bằng $70\%$ chiều dài. Tính diện tích hình chữ nhật đó theo mét vuông.
\shortans{1120}
\loigiai{
\textbf{Đáp số:} $1120$\\
Chiều rộng: $40 \times 70\% = 28\text{ m}$.\\
Diện tích: $40 \times 28 = 1\,120\text{ m}^2$.
}
\end{ex}

\begin{ex} %[5D7V41]
Biết $35\%$ của số $x$ là $147$. Tìm số $x$.
\shortans{420}
\loigiai{
\textbf{Đáp số:} $420$\\
$x = 147 : 35 \times 100 = 420$.
}
\end{ex}

\begin{ex} %[5D7C41]
Một xí nghiệp có $300$ công nhân. Sau một năm, số công nhân tăng thêm $8\%$. Sang năm tiếp theo, số công nhân lại tăng thêm $5\%$ so với năm trước đó. Hỏi sau hai năm xí nghiệp đó có tất cả bao nhiêu công nhân?
\shortans{340 hoặc 340,2 (khoảng 340 người)}
\loigiai{
\textbf{Đáp số:} $340$ hoặc $340{,}2$ (khoảng 340 người)\\
Sau năm 1: $300 \times 1{,}08 = 324$ người.\\
Sau năm 2: $324 \times 1{,}05 = 340{,}2 \approx 340$ người.
}
\end{ex}

\newpage

\TaoPhieuKiemTra{Chủ đề 7 -- Tỉ số và các bài toán liên quan}{42}{Máy tính cầm tay}{https://loptoanthaylinh.vn/thi-online/toan-5/bai-42}
\textbf{Phần 1. Trắc nghiệm lựa chọn}
\setcounterpageref{ex}{0}

\begin{ex} %[5D7N42]
Phím nào trên máy tính cầm tay dùng để bật máy và xóa toàn bộ màn hình tính toán?
\choice
{\fbox{OFF}}
{\fbox{CE}}
{\fbox{=}}
{\True \fbox{ON/C}}
\loigiai{
Phím \fbox{ON/C} dùng để bật máy và xóa dữ liệu hiển thị trên màn hình.\\
\textbf{Đáp án của câu hỏi là D}
}
\end{ex}

\begin{ex} %[5D7N42]
Để thực hiện phép cộng trên máy tính cầm tay, ta sử dụng phím có kí hiệu nào?
\choice
{\fbox{--}}
{\fbox{$\times$}}
{\True \fbox{+}}
{\fbox{$\div$}}
\loigiai{
Phím phép cộng có kí hiệu \fbox{+}.\\
\textbf{Đáp án của câu hỏi là C}
}
\end{ex}

\begin{ex} %[5D7N42]
Để nhận kết quả của phép tính vừa nhập trên máy tính cầm tay, ta bấm phím nào?
\choice
{\fbox{ON}}
{\True \fbox{=}}
{\fbox{\%}}
{\fbox{CE}}
\loigiai{
Phím \fbox{=} dùng để tính và hiển thị kết quả.\\
\textbf{Đáp án của câu hỏi là B}
}
\end{ex}

\begin{ex} %[5D7H42]
Khi bấm lần lượt các phím: \fbox{2} \fbox{5} \fbox{+} \fbox{1} \fbox{8} \fbox{=}, màn hình máy tính sẽ hiển thị kết quả là
\choice
{\True $43$}
{$33$}
{$45$}
{$2518$}
\loigiai{
Kết quả: $25 + 18 = 43$.\\
\textbf{Đáp án của câu hỏi là A}
}
\end{ex}

\begin{ex} %[5D7H42]
Để tính $35\%$ của $400$ trên máy tính cầm tay có phím \%, ta có thể bấm lần lượt là
\choice
{\fbox{4} \fbox{0} \fbox{0} \fbox{+} \fbox{3} \fbox{5} \fbox{\%}}
{\fbox{4} \fbox{0} \fbox{0} \fbox{$\div$} \fbox{3} \fbox{5} \fbox{\%}}
{\True \fbox{4} \fbox{0} \fbox{0} \fbox{$\times$} \fbox{3} \fbox{5} \fbox{\%}}
{\fbox{3} \fbox{5} \fbox{+} \fbox{4} \fbox{0} \fbox{0} \fbox{=}}
\loigiai{
Cách tính $35\%$ của $400$: lấy $400 \times 35\%$.\\
\textbf{Đáp án của câu hỏi là C}
}
\end{ex}

\begin{ex} %[5D7H42]
Khi bấm các phím: \fbox{1} \fbox{.} \fbox{2} \fbox{5} \fbox{$\times$} \fbox{8} \fbox{=}, màn hình hiển thị kết quả là
\choice
{$100$}
{\True $10$}
{$1$}
{$0{,}1$}
\loigiai{
$1{,}25 \times 8 = 10$.\\
\textbf{Đáp án của câu hỏi là B}
}
\end{ex}

\begin{ex} %[5D7V42]
Để tính tỉ số phần trăm của $15$ và $60$ bằng máy tính cầm tay, ta bấm lần lượt các phím:
\choice
{\True \fbox{1} \fbox{5} \fbox{$\div$} \fbox{6} \fbox{0} \fbox{\%}}
{\fbox{1} \fbox{5} \fbox{$\times$} \fbox{6} \fbox{0} \fbox{=}}
{\fbox{6} \fbox{0} \fbox{$\div$} \fbox{1} \fbox{5} \fbox{=}}
{\fbox{1} \fbox{5} \fbox{+} \fbox{6} \fbox{0} \fbox{\%}}
\loigiai{
Để tìm tỉ số phần trăm: $15 : 60 \% = 25\%$.\\
\textbf{Đáp án của câu hỏi là A}
}
\end{ex}

\begin{ex} %[5D7V42]
Khi thực hiện phép tính $(45{,}6 + 18{,}4) : 4$ trên máy tính cầm tay, kết quả hiển thị trên màn hình là
\choice
{$14$}
{\True $16$}
{$18$}
{$20$}
\loigiai{
Ta có: $45{,}6 + 18{,}4 = 64$; $64 : 4 = 16$.\\
\textbf{Đáp án của câu hỏi là B}
}
\end{ex}

\begin{ex} %[5D7V42]
Một học sinh tính giá trị biểu thức $150 - 30 \times 4$ trên máy tính cầm tay không có chức năng thứ tự phép tính tự động (bấm lần lượt \fbox{1}\fbox{5}\fbox{0} \fbox{--} \fbox{3}\fbox{0} \fbox{$\times$} \fbox{4} \fbox{=}). Kết quả màn hình hiển thị là $480$. Kết quả đúng của biểu thức phải là
\choice
{$480$}
{$120$}
{$60$}
{\True $30$}
\loigiai{
Thứ tự phép tính: nhân trước trừ sau $\Rightarrow 150 - (30 \times 4) = 150 - 120 = 30$.\\
\textbf{Đáp án của câu hỏi là D}
}
\end{ex}

\begin{ex} %[5D7V42]
Khi bấm: \fbox{7} \fbox{2} \fbox{0} \fbox{$\times$} \fbox{1} \fbox{5} \fbox{\%} \fbox{=}, kết quả trên màn hình là
\choice
{$10{,}8$}
{$1\,080$}
{\True $108$}
{$54$}
\loigiai{
$720 \times 15\% = 108$.\\
\textbf{Đáp án của câu hỏi là C}
}
\end{ex}

\begin{ex} %[5D7V42]
Sử dụng máy tính cầm tay để tính kết quả biểu thức: $3{,}75 \times 4{,}8 - 2{,}5 \times 3{,}2$. Kết quả là
\choice
{$8$}
{$12$}
{$15$}
{\True $10$}
\loigiai{
$3{,}75 \times 4{,}8 = 18$; $2{,}5 \times 3{,}2 = 8$; $18 - 8 = 10$.\\
\textbf{Đáp án của câu hỏi là D}
}
\end{ex}

\begin{ex} %[5D7C42]
Sử dụng máy tính cầm tay tính giá trị của biểu thức: $\dfrac{12{,}6 \times 4{,}5 + 12{,}6 \times 5{,}5}{0{,}63 \times 4}$. Kết quả thu được là
\choice
{\True $50$}
{$25$}
{$100$}
{$200$}
\loigiai{
Tử số: $12{,}6 \times (4{,}5 + 5{,}5) = 12{,}6 \times 10 = 126$.\\
Mẫu số: $0{,}63 \times 4 = 2{,}52$.\\
Kết quả: $126 : 2{,}52 = 50$.\\
\textbf{Đáp án của câu hỏi là A}
}
\end{ex}

\textbf{Phần 2. Trắc nghiệm đúng sai}
\setcounterpageref{ex}{0}

\begin{ex} %[5D7H42]
Khi sử dụng máy tính cầm tay thông thường, xét tính đúng sai của các nhận định:
\choiceTF
{\True Phím \fbox{.} dùng để nhập dấu phẩy ngăn cách phần nguyên và phần thập phân.}
{Phím \fbox{OFF} dùng để xóa kết quả vừa tính mà không tắt máy.}
{\True Để tính $25 + 75$, ta bấm lần lượt: \fbox{2}\fbox{5} \fbox{+} \fbox{7}\fbox{5} \fbox{=}.}
{Bấm \fbox{5}\fbox{0} \fbox{$\div$} \fbox{0} \fbox{=} máy tính sẽ hiển thị số $0$.}
\loigiai{
\begin{itemize}
    \item Phím dấu chấm biểu thị dấu phân cách thập phân. Đúng.
    \item Phím OFF dùng để tắt nguồn máy tính. Sai.
    \item Thao tác cộng hai số nguyên: đúng.
    \item Phép chia cho số 0 là phép toán không hợp lệ, máy sẽ báo lỗi (Error). Sai.
\end{itemize}
\begin{itemize}
    \item a) \textbf{Đúng}.
    \item b) \textbf{Sai}.
    \item c) \textbf{Đúng}.
    \item d) \textbf{Sai}.
\end{itemize}
}
\end{ex}

\begin{ex} %[5D7V42]
Sử dụng máy tính cầm tay để tính các phép tính sau, xét tính đúng sai:
\choiceTF
{\True $125{,}6 \times 4{,}5 = 565{,}2$.}
{$\dfrac{3}{8}$ viết thành số thập phân bằng cách bấm \fbox{3} \fbox{$\div$} \fbox{8} \fbox{=} cho kết quả $0{,}38$.}
{\True $75\%$ của $640$ bằng cách bấm \fbox{6}\fbox{4}\fbox{0} \fbox{$\times$} \fbox{7}\fbox{5} \fbox{\%} cho kết quả $480$.}
{\True Phép chia $14 : 25$ hiển thị kết quả là $0{,}56$.}
\loigiai{
\begin{itemize}
    \item $125{,}6 \times 4{,}5 = 565{,}2$. Đúng.
    \item $3 : 8 = 0{,}375 \neq 0{,}38$. Sai.
    \item $640 \times 75\% = 480$. Đúng.
    \item $14 : 25 = 0{,}56$. Đúng.
\end{itemize}
\begin{itemize}
    \item a) \textbf{Đúng}.
    \item b) \textbf{Sai}.
    \item c) \textbf{Đúng}.
    \item d) \textbf{Đúng}.
\end{itemize}
}
\end{ex}

\textbf{Phần 3. Trả lời ngắn}
\setcounterpageref{ex}{0}

\begin{ex} %[5D7H42]
Dùng máy tính cầm tay thực hiện phép tính: $456{,}8 + 289{,}75$. (Ghi kết quả dưới dạng số thập phân dùng dấu phẩy).
\shortans{746,55}
\loigiai{
\textbf{Đáp số:} $746{,}55$\\
$456{,}8 + 289{,}75 = 746{,}55$.
}
\end{ex}

\begin{ex} %[5D7H42]
Dùng máy tính cầm tay tính kết quả: $85{,}4 \times 6{,}5$.
\shortans{555,1}
\loigiai{
\textbf{Đáp số:} $555{,}1$\\
$85{,}4 \times 6{,}5 = 555{,}1$.
}
\end{ex}

\begin{ex} %[5D7V42]
Dùng máy tính cầm tay tính kết quả: $148{,}8 : 4{,}8$.
\shortans{31}
\loigiai{
\textbf{Đáp số:} $31$\\
$148{,}8 : 4{,}8 = 31$.
}
\end{ex}

\begin{ex} %[5D7V42]
Sử dụng máy tính cầm tay để tìm $18\%$ của số $1\,250$.
\shortans{225}
\loigiai{
\textbf{Đáp số:} $225$\\
$1\,250 \times 18\% = 225$.
}
\end{ex}

\begin{ex} %[5D7V42]
Dùng máy tính cầm tay tính giá trị biểu thức: $(48{,}6 + 31{,}4) \times 2{,}5 - 50$.
\shortans{150}
\loigiai{
\textbf{Đáp số:} $150$\\
$80 \times 2{,}5 - 50 = 200 - 50 = 150$.
}
\end{ex}

\begin{ex} %[5D7C42]
Tính giá trị biểu thức bằng máy tính cầm tay: $15{,}6 \times 3{,}4 + 15{,}6 \times 6{,}6 + 44$.
\shortans{200}
\loigiai{
\textbf{Đáp số:} $200$\\
$15{,}6 \times (3{,}4 + 6{,}6) + 44 = 156 + 44 = 200$.
}
\end{ex}

\newpage

\TaoPhieuKiemTra{Chủ đề 7 -- Tỉ số và các bài toán liên quan}{43}{Thực hành và trải nghiệm với máy tính cầm tay}{https://loptoanthaylinh.vn/thi-online/toan-5/bai-43}
\textbf{Phần 1. Trắc nghiệm lựa chọn}
\setcounterpageref{ex}{0}

\begin{ex} %[5D7N43]
Bạn Nam mua $3$ quyển vở giá $8\,500$ đồng/quyển và $2$ cây bút giá $5\,000$ đồng/cây. Biểu thức tính tổng số tiền bằng máy tính cầm tay là
\choice
{$3 \times 8\,500 \times 2 \times 5\,000$}
{$3 + 8\,500 + 2 + 5\,000$}
{\True $3 \times 8\,500 + 2 \times 5\,000$}
{$(3 + 2) \times (8\,500 + 5\,000)$}
\loigiai{
Tổng số tiền mua vở và bút là: $3 \times 8\,500 + 2 \times 5\,000$.\\
\textbf{Đáp án của câu hỏi là C}
}
\end{ex}

\begin{ex} %[5D7N43]
Một giỏ hoa quả gồm $1{,}5\text{ kg}$ táo với giá $40\,000$ đồng/kg. Số tiền mua táo tính được là
\choice
{$50\,000$ đồng}
{$70\,000$ đồng}
{$80\,000$ đồng}
{\True $60\,000$ đồng}
\loigiai{
Số tiền: $1{,}5 \times 40\,000 = 60\,000$ (đồng).\\
\textbf{Đáp án của câu hỏi là D}
}
\end{ex}

\begin{ex} %[5D7H43]
Mẹ đưa cho Mai $200\,000$ đồng đi siêu thị. Mai mua hết $146\,500$ đồng. Sử dụng máy tính cầm tay tính số tiền thừa Mai nhận lại là
\choice
{$63\,500$ đồng}
{$54\,500$ đồng}
{\True $53\,500$ đồng}
{$64\,500$ đồng}
\loigiai{
Số tiền thừa: $200\,000 - 146\,500 = 53\,500$ (đồng).\\
\textbf{Đáp án của câu hỏi là C}
}
\end{ex}

\begin{ex} %[5D7H43]
Một cửa hàng giảm giá $10\%$ cho hóa đơn mua hàng trên $500\,000$ đồng. Bác An mua hàng hết $650\,000$ đồng. Số tiền bác An được giảm là
\choice
{$50\,000$ đồng}
{\True $65\,000$ đồng}
{$60\,000$ đồng}
{$70\,000$ đồng}
\loigiai{
Số tiền được giảm là: $650\,000 \times 10\% = 65\,000$ (đồng).\\
\textbf{Đáp án của câu hỏi là B}
}
\end{ex}

\begin{ex} %[5D7H43]
Một mảnh vườn hình chữ nhật có chiều dài $24{,}5\text{ m}$ và chiều rộng $16\text{ m}$. Sử dụng máy tính cầm tay tính diện tích mảnh vườn là
\choice
{$384\text{ m}^2$}
{$400\text{ m}^2$}
{$412\text{ m}^2$}
{\True $392\text{ m}^2$}
\loigiai{
Diện tích: $24{,}5 \times 16 = 392\text{ (m}^2\text{)}$.\\
\textbf{Đáp án của câu hỏi là D}
}
\end{ex}

\begin{ex} %[5D7H43]
Bác Ba thu hoạch được $1{,}2$ tấn thóc. Bác để lại $15\%$ để ăn, số còn lại đem bán. Dùng máy tính cầm tay tính khối lượng thóc bác Ba đem bán là
\choice
{$180\text{ kg}$}
{$1\,000\text{ kg}$}
{\True $1\,020\text{ kg}$}
{$980\text{ kg}$}
\loigiai{
Đổi $1{,}2$ tấn $= 1\,200\text{ kg}$.\\
Khối lượng thóc đem bán: $1\,200 \times (100\% - 15\%) = 1\,020\text{ (kg)}$.\\
\textbf{Đáp án của câu hỏi là C}
}
\end{ex}

\begin{ex} %[5D7V43]
Lan mua $5$ quyển sách tham khảo giá $35\,000$ đồng/quyển và $1$ chiếc cặp sách giá $185\,000$ đồng. Nhân dịp năm học mới cửa hàng giảm $10\%$ trên tổng hóa đơn. Lan phải trả số tiền là
\choice
{$360\,000$ đồng}
{$330\,000$ đồng}
{$340\,000$ đồng}
{\True $324\,000$ đồng}
\loigiai{
Tổng hóa đơn trước giảm: $5 \times 35\,000 + 185\,000 = 175\,000 + 185\,000 = 360\,000$ (đồng).\\
Số tiền phải trả sau khi giảm $10\%$: $360\,000 \times 90\% = 324\,000$ (đồng).\\
\textbf{Đáp án của câu hỏi là D}
}
\end{ex}

\begin{ex} %[5D7V43]
Một gia đình có thu nhập hàng tháng là $25\,000\,000$ đồng. Gia đình dành $45\%$ chi phí cho ăn uống, $25\%$ cho học tập, $15\%$ cho các chi tiêu khác, phần còn lại để tiết kiệm. Số tiền tiết kiệm mỗi tháng là
\choice
{\True $3\,750\,000$ đồng}
{$3\,500\,000$ đồng}
{$4\,000\,000$ đồng}
{$4\,250\,000$ đồng}
\loigiai{
Tỉ lệ tiết kiệm: $100\% - (45\% + 25\% + 15\%) = 15\%$.\\
Số tiền tiết kiệm: $25\,000\,000 \times 15\% = 3\,750\,000$ (đồng).\\
\textbf{Đáp án của câu hỏi là A}
}
\end{ex}

\begin{ex} %[5D7V43]
Một cửa hàng điện máy bán một máy giặt với giá $7\,800\,000$ đồng. Khách hàng thanh toán qua thẻ ngân hàng được hoàn tiền $5\%$. Số tiền thực tế khách hàng phải trả là
\choice
{$7\,500\,000$ đồng}
{\True $7\,410\,000$ đồng}
{$7\,400\,000$ đồng}
{$7\,390\,000$ đồng}
\loigiai{
Số tiền được hoàn: $7\,800\,000 \times 5\% = 390\,000$ (đồng).\\
Số tiền thực tế phải trả: $7\,800\,000 - 390\,000 = 7\,410\,000$ (đồng).\\
\textbf{Đáp án của câu hỏi là B}
}
\end{ex}

\begin{ex} %[5D7V43]
Bác nông dân nuôi $250$ con gà và $150$ con vịt. Sau một đợt dịch bệnh, tỉ lệ sống sót của gà là $96\%$, của vịt là $94\%$. Tổng số gà và vịt còn sống sót là
\choice
{\True $381$ con}
{$375$ con}
{$385$ con}
{$390$ con}
\loigiai{
Số gà sống sót: $250 \times 96\% = 240$ (con).\\
Số vịt sống sót: $150 \times 94\% = 141$ (con).\\
Tổng số con sống sót: $240 + 141 = 381$ (con).\\
\textbf{Đáp án của câu hỏi là A}
}
\end{ex}

\begin{ex} %[5D7V43]
Một hình hộp chữ nhật có chiều dài $1{,}2\text{ m}$, chiều rộng $0{,}8\text{ m}$ và chiều cao $0{,}5\text{ m}$. Sử dụng máy tính cầm tay tính thể tích của hình hộp này theo đơn vị lít ($1\text{ dm}^3 = 1\text{ l}$).
\choice
{$48\text{ l}$}
{\True $480\text{ l}$}
{$4\,800\text{ l}$}
{$4800\text{ l}$}
\loigiai{
Thể tích: $1{,}2 \times 0{,}8 \times 0{,}5 = 0{,}48\text{ (m}^3\text{)} = 480\text{ dm}^3 = 480\text{ l}$.\\
\textbf{Đáp án của câu hỏi là B}
}
\end{ex}

\begin{ex} %[5D7C43]
Một siêu thị có chương trình khuyến mãi: Mua $3$ hộp sữa tặng $1$ hộp sữa cùng loại. Mỗi hộp sữa có giá $28\,000$ đồng. Mẹ An cần $14$ hộp sữa thì phải trả ít nhất bao nhiêu tiền?
\choice
{\True $308\,000$ đồng}
{$336\,000$ đồng}
{$292\,000$ đồng}
{$392\,000$ đồng}
\loigiai{
Mua 3 tặng 1 nghĩa là mua 4 hộp chỉ trả tiền 3 hộp ($3 \times 28\,000 = 84\,000$ đồng).\\
Ta có: $14 = 4 \times 3 + 2$.\\
Số hộp phải trả tiền: $3 \times 3 + 2 = 11$ hộp.\\
Tổng số tiền phải trả: $11 \times 28\,000 = 308\,000$ (đồng).\\
\textbf{Đáp án của câu hỏi là A}
}
\end{ex}

\textbf{Phần 2. Trắc nghiệm đúng sai}
\setcounterpageref{ex}{0}

\begin{ex} %[5D7H43]
Bạn Hoa đi mua đồ dùng học tập gồm: $5$ cuốn tập giá $12\,000$ đồng/cuốn, $3$ cây bút bi giá $6\,000$ đồng/cây và $1$ hộp bút giá $45\,000$ đồng. Xét tính đúng sai:
\choiceTF
{\True Tổng số tiền mua $5$ cuốn tập là $60\,000$ đồng.}
{\True Tổng số tiền mua $3$ cây bút bi là $18\,000$ đồng.}
{Tổng tiền hóa đơn khi chưa giảm giá là $133\,000$ đồng.}
{\True Nếu được giảm giá $10\%$, Hoa phải trả $110\,700$ đồng.}
\loigiai{
\begin{itemize}
    \item Tiền tập: $5 \times 12\,000 = 60\,000$ đồng. Đúng.
    \item Tiền bút bi: $3 \times 6\,000 = 18\,000$ đồng. Đúng.
    \item Tổng hóa đơn: $60\,000 + 18\,000 + 45\,000 = 123\,000$ đồng $\neq 133\,000$ đồng. Sai.
    \item Tiền sau giảm 10\%: $123\,000 \times 0{,}9 = 110\,700$ đồng. Đúng.
\end{itemize}
\begin{itemize}
    \item a) \textbf{Đúng}.
    \item b) \textbf{Đúng}.
    \item c) \textbf{Sai}.
    \item d) \textbf{Đúng}.
\end{itemize}
}
\end{ex}

\begin{ex} %[5D7V43]
Một xí nghiệp may có $120$ công nhân. Mỗi công nhân trung bình may được $15$ chiếc áo một ngày. Mỗi chiếc áo bán với giá $150\,000$ đồng. Xét tính đúng sai:
\choiceTF
{\True Trong một ngày, xí nghiệp may được $1\,800$ chiếc áo.}
{\True Doanh thu từ số áo may trong một ngày là $270$ triệu đồng.}
{Nếu tăng năng suất thêm $10\%$ thì mỗi ngày may được $1\,900$ chiếc áo.}
{\True Sau khi tăng năng suất $10\%$, mỗi ngày xí nghiệp may được $1\,980$ chiếc áo.}
\loigiai{
\begin{itemize}
    \item Số áo 1 ngày: $120 \times 15 = 1\,800$ chiếc. Đúng.
    \item Doanh thu: $1\,800 \times 150\,000 = 270\,000\,000$ đồng $= 270$ triệu đồng. Đúng.
    \item Tăng 10\%: $1\,800 \times 1{,}1 = 1\,980$ chiếc $\neq 1\,900$ chiếc. Ý 3 Sai, Ý 4 Đúng.
\end{itemize}
\begin{itemize}
    \item a) \textbf{Đúng}.
    \item b) \textbf{Đúng}.
    \item c) \textbf{Sai}.
    \item d) \textbf{Đúng}.
\end{itemize}
}
\end{ex}

\textbf{Phần 3. Trả lời ngắn}
\setcounterpageref{ex}{0}

\begin{ex} %[5D7H43]
Sử dụng máy tính cầm tay tính tổng số tiền mua hàng: $4$ hộp bánh giá $45\,000$ đồng/hộp và $6$ chai nước ngọt giá $12\,500$ đồng/chai. (Ghi kết quả dưới dạng số nghìn đồng, ví dụ: 255).
\shortans{255}
\loigiai{
\textbf{Đáp số:} $255$\\
$4 \times 45\,000 + 6 \times 12\,500 = 180\,000 + 75\,000 = 255\,000$ đồng $= 255$ nghìn đồng.
}
\end{ex}

\begin{ex} %[5D7H43]
Một người đi chợ mua $2{,}5\text{ kg}$ thịt lợn với giá $130\,000$ đồng/kg và $3\text{ kg}$ cá với giá $85\,000$ đồng/kg. Tính tổng số tiền người đó phải trả theo đơn vị nghìn đồng.
\shortans{580}
\loigiai{
\textbf{Đáp số:} $580$\\
$2{,}5 \times 130 + 3 \times 85 = 325 + 255 = 580$ nghìn đồng.
}
\end{ex}

\begin{ex} %[5D7V43]
Một mảnh đất hình chữ nhật có chiều dài $32{,}5\text{ m}$, chiều rộng $18\text{ m}$. Người ta dùng $20\%$ diện tích để xây nhà. Tính diện tích xây nhà theo mét vuông.
\shortans{117}
\loigiai{
\textbf{Đáp số:} $117$\\
Diện tích mảnh đất: $32{,}5 \times 18 = 585\text{ m}^2$.\\
Diện tích xây nhà: $585 \times 20\% = 117\text{ m}^2$.
}
\end{ex}

\begin{ex} %[5D7V43]
Một gia đình sử dụng hết $185\text{ kWh}$ điện trong tháng. Giá $50\text{ kWh}$ đầu là $1\,800$ đồng/kWh, $50\text{ kWh}$ tiếp theo là $1\,900$ đồng/kWh, số còn lại là $2\,100$ đồng/kWh. Tính tổng số tiền điện phải trả (chưa tính thuế) theo đơn vị nghìn đồng.
\shortans{363,5}
\loigiai{
\textbf{Đáp số:} $363{,}5$\\
Tiền điện: $50 \times 1{,}8 + 50 \times 1{,}9 + 85 \times 2{,}1 = 90 + 95 + 178{,}5 = 363{,}5$ nghìn đồng.
}
\end{ex}

\begin{ex} %[5D7V43]
Một cửa hàng quần áo niêm yết giá một chiếc áo khoác là $650\,000$ đồng. Trong chương trình tri ân, áo được giảm giá $20\%$, khách quen lại được giảm thêm $5\%$ trên giá đã giảm. Hỏi khách quen mua áo khoác đó với giá bao nhiêu nghìn đồng?
\shortans{494}
\loigiai{
\textbf{Đáp số:} $494$\\
Giá sau giảm 20\%: $650 \times 0{,}8 = 520$ nghìn đồng.\\
Giá khách quen: $520 \times 0{,}95 = 494$ nghìn đồng.
}
\end{ex}

\begin{ex} %[5D7C43]
Một bác nông dân bán $450\text{ kg}$ bưởi với giá $35\,000$ đồng/kg. Bác dành $40\%$ số tiền thu được để mua phân bón, $\dfrac{1}{3}$ số tiền còn lại để mua giống cây mới, số tiền thừa bác gửi vào quỹ tiết kiệm. Hỏi bác gửi vào quỹ tiết kiệm bao nhiêu nghìn đồng?
\shortans{6300}
\loigiai{
\textbf{Đáp số:} $6300$\\
Tổng tiền: $450 \times 35 = 15\,750$ nghìn đồng.\\
Mua phân bón (40\%): $15\,750 \times 0{,}4 = 6\,300$ nghìn đồng.\\
Còn lại: $15\,750 - 6\,300 = 9\,450$ nghìn đồng.\\
Mua giống (1/3 còn lại): $9\,450 : 3 = 3\,150$ nghìn đồng.\\
Gửi tiết kiệm: $9\,450 - 3\,150 = 6\,300$ nghìn đồng.
}
\end{ex}

\newpage

\TaoPhieuKiemTra{Chủ đề 7 -- Tỉ số và các bài toán liên quan}{44}{Luyện tập chung tỉ số phần trăm}{https://loptoanthaylinh.vn/thi-online/toan-5/bai-44}
\textbf{Phần 1. Trắc nghiệm lựa chọn}
\setcounterpageref{ex}{0}

\begin{ex} %[5D7N44]
Tỉ số của hai số $5$ và $9$ là
\choice
{\True $\dfrac{5}{9}$}
{$\dfrac{9}{5}$}
{$5 \times 9$}
{$45$}
\loigiai{
Tỉ số của hai số $5$ và $9$ là $\dfrac{5}{9}$.\\
\textbf{Đáp án của câu hỏi là A}
}
\end{ex}

\begin{ex} %[5D7N44]
Tỉ số phần trăm $48\%$ được viết dưới dạng số thập phân là
\choice
{$4{,}8$}
{$0{,}048$}
{\True $0{,}48$}
{$48$}
\loigiai{
Ta có: $48\% = \dfrac{48}{100} = 0{,}48$.\\
\textbf{Đáp án của câu hỏi là C}
}
\end{ex}

\begin{ex} %[5D7H44]
Trên bản đồ tỉ lệ $1 : 50\,000$, khoảng cách giữa hai địa điểm đo được là $4\text{ cm}$. Khoảng cách thật ngoài thực tế là
\choice
{$200\text{ m}$}
{$20\text{ km}$}
{$200\text{ km}$}
{\True $2\text{ km}$}
\loigiai{
Khoảng cách thực tế: $4 \times 50\,000 = 200\,000\text{ cm} = 2\text{ km}$.\\
\textbf{Đáp án của câu hỏi là D}
}
\end{ex}

\begin{ex} %[5D7H44]
Tổng của hai số là $80$, tỉ số của hai số là $\dfrac{2}{3}$. Số lớn là
\choice
{$32$}
{\True $48$}
{$50$}
{$60$}
\loigiai{
Tổng số phần: $2 + 3 = 5$ (phần).\\
Số lớn là: $80 : 5 \times 3 = 48$.\\
\textbf{Đáp án của câu hỏi là B}
}
\end{ex}

\begin{ex} %[5D7H44]
Hiệu của hai số là $42$, tỉ số của hai số là $\dfrac{3}{5}$. Số bé là
\choice
{$21$}
{$105$}
{\True $63$}
{$84$}
\loigiai{
Hiệu số phần: $5 - 3 = 2$ (phần).\\
Số bé là: $42 : 2 \times 3 = 63$.\\
\textbf{Đáp án của câu hỏi là C}
}
\end{ex}

\begin{ex} %[5D7H44]
Một cửa hàng bán được $450\text{ kg}$ gạo tẻ và $150\text{ kg}$ gạo nếp. Tỉ số phần trăm của số gạo nếp so với tổng số gạo đã bán là
\choice
{$30\%$}
{\True $25\%$}
{$33{,}3\%$}
{$20\%$}
\loigiai{
Tổng số gạo: $450 + 150 = 600\text{ (kg)}$.\\
Tỉ số phần trăm gạo nếp: $150 : 600 = 0{,}25 = 25\%$.\\
\textbf{Đáp án của câu hỏi là B}
}
\end{ex}

\begin{ex} %[5D7V44]
Một mảnh đất hình chữ nhật có chu vi $160\text{ m}$, chiều rộng bằng $\dfrac{3}{5}$ chiều dài. Diện tích mảnh đất đó là
\choice
{\True $1\,500\text{ m}^2$}
{$1\,200\text{ m}^2$}
{$1\,600\text{ m}^2$}
{$6\,000\text{ m}^2$}
\loigiai{
Nửa chu vi: $160 : 2 = 80\text{ (m)}$.\\
Chiều dài: $80 : (3 + 5) \times 5 = 50\text{ (m)}$.\\
Chiều rộng: $80 - 50 = 30\text{ (m)}$.\\
Diện tích: $50 \times 30 = 1\,500\text{ (m}^2\text{)}$.\\
\textbf{Đáp án của câu hỏi là A}
}
\end{ex}

\begin{ex} %[5D7V44]
Một người gửi tiết kiệm $120\,000\,000$ đồng với lãi suất $0{,}5\%$ một tháng. Sau một tháng cả tiền gửi và tiền lãi là
\choice
{$120\,500\,000$ đồng}
{$126\,000\,000$ đồng}
{$125\,000\,000$ đồng}
{\True $120\,600\,000$ đồng}
\loigiai{
Tiền lãi 1 tháng: $120\,000\,000 \times 0{,}5\% = 600\,000$ (đồng).\\
Tổng số tiền: $120\,000\,000 + 600\,000 = 120\,600\,000$ (đồng).\\
\textbf{Đáp án của câu hỏi là D}
}
\end{ex}

\begin{ex} %[5D7V44]
Năm nay tuổi bố gấp $4$ lần tuổi con. Biết $4$ năm trước bố hơn con $30$ tuổi. Tuổi bố hiện nay là
\choice
{$36$ tuổi}
{$44$ tuổi}
{\True $40$ tuổi}
{$48$ tuổi}
\loigiai{
Hiệu số tuổi không đổi, hiện nay bố vẫn hơn con $30$ tuổi.\\
Hiệu số phần: $4 - 1 = 3$ (phần).\\
Tuổi bố hiện nay là: $30 : 3 \times 4 = 40$ (tuổi).\\
\textbf{Đáp án của câu hỏi là C}
}
\end{ex}

\begin{ex} %[5D7V44]
Một người mua một chiếc xe đạp giá $2\,000\,000$ đồng. Sau đó bán lại bị lỗ $5\%$. Người đó đã bán chiếc xe với giá là
\choice
{\True $1\,900\,000$ đồng}
{$1\,850\,000$ đồng}
{$1\,950\,000$ đồng}
{$1\,800\,000$ đồng}
\loigiai{
Số tiền lỗ: $2\,000\,000 \times 5\% = 100\,000$ (đồng).\\
Giá bán: $2\,000\,000 - 100\,000 = 1\,900\,000$ (đồng).\\
\textbf{Đáp án của câu hỏi là A}
}
\end{ex}

\begin{ex} %[5D7V44]
Trên bản đồ tỉ lệ $1 : 2\,500$, một mảnh ruộng hình vuông có chu vi là $16\text{ cm}$. Diện tích thực tế của mảnh ruộng đó là
\choice
{$1\text{ ha}$}
{\True $1\text{ ha}$}
{$10\text{ ha}$}
{$0{,}1\text{ ha}$}
\loigiai{
Cạnh trên bản đồ: $16 : 4 = 4\text{ (cm)}$.\\
Cạnh thực tế: $4 \times 2\,500 = 10\,000\text{ cm} = 100\text{ m}$.\\
Diện tích thực tế: $100 \times 100 = 10\,000\text{ m}^2 = 1\text{ ha}$.\\
\textbf{Đáp án của câu hỏi là B}
}
\end{ex}

\begin{ex} %[5D7C44]
Một cửa hàng nhân ngày Quốc tế Thiếu nhi đã hạ giá các mặt hàng $15\%$, tuy vậy cửa hàng vẫn còn lãi $10{,}5\%$ so với giá vốn. Hỏi nếu không hạ giá thì cửa hàng lãi bao nhiêu phần trăm so với giá vốn?
\choice
{$25\%$}
{$28\%$}
{$32\%$}
{\True $30\%$}
\loigiai{
Coi giá vốn là $100\%$.\\
Giá bán ngày lễ là: $100\% + 10{,}5\% = 110{,}5\%$ (giá vốn).\\
Giá bán ngày lễ bằng $100\% - 15\% = 85\%$ giá định bán.\\
Giá định bán so với giá vốn: $110{,}5\% : 85 \times 100 = 130\%$.\\
Nếu không hạ giá thì lãi: $130\% - 100\% = 30\%$.\\
\textbf{Đáp án của câu hỏi là D}
}
\end{ex}

\textbf{Phần 2. Trắc nghiệm đúng sai}
\setcounterpageref{ex}{0}

\begin{ex} %[5D7H44]
Một nông trại nuôi bò và dê với tổng cộng $450$ con, số bò bằng $\dfrac{4}{5}$ số dê. Xét tính đúng sai:
\choiceTF
{\True Số bò chiếm $4$ phần, số dê chiếm $5$ phần bằng nhau.}
{Số bò là $250$ con.}
{\True Số dê là $250$ con.}
{\True Số bò chiếm khoảng $44{,}4\%$ tổng đàn gia súc.}
\loigiai{
\begin{itemize}
    \item Tỉ lệ bò và dê: $4 : 5$. Đúng.
    \item Số bò: $450 : (4+5) \times 4 = 200$ con $\neq 250$ con. Sai.
    \item Số dê: $450 - 200 = 250$ con. Đúng.
    \item Tỉ số phần trăm bò: $200 : 450 \approx 44{,}44\%$. Đúng.
\end{itemize}
\begin{itemize}
    \item a) \textbf{Đúng}.
    \item b) \textbf{Sai}.
    \item c) \textbf{Đúng}.
    \item d) \textbf{Đúng}.
\end{itemize}
}
\end{ex}

\begin{ex} %[5D7V44]
Một mảnh đất hình chữ nhật có chu vi $120\text{ m}$, chiều rộng bằng $\dfrac{2}{3}$ chiều dài. Người ta dùng $25\%$ diện tích để trồng cây ăn quả, $35\%$ diện tích để làm ao cá, phần còn lại trồng rau. Xét tính đúng sai:
\choiceTF
{\True Diện tích mảnh đất là $864\text{ m}^2$.}
{\True Diện tích trồng cây ăn quả là $216\text{ m}^2$.}
{Diện tích làm ao cá là $310\text{ m}^2$.}
{\True Diện tích trồng rau là $345{,}6\text{ m}^2$.}
\loigiai{
\begin{itemize}
    \item Nửa chu vi: $60\text{ m}$. Chiều dài: $36\text{ m}$, chiều rộng: $24\text{ m}$. Diện tích: $36 \times 24 = 864\text{ m}^2$. Đúng.
    \item Trồng cây: $864 \times 25\% = 216\text{ m}^2$. Đúng.
    \item Ao cá: $864 \times 35\% = 302{,}4\text{ m}^2 \neq 310\text{ m}^2$. Sai.
    \item Trồng rau: $864 \times 40\% = 345{,}6\text{ m}^2$. Đúng.
\end{itemize}
\begin{itemize}
    \item a) \textbf{Đúng}.
    \item b) \textbf{Đúng}.
    \item c) \textbf{Sai}.
    \item d) \textbf{Đúng}.
\end{itemize}
}
\end{ex}

\textbf{Phần 3. Trả lời ngắn}
\setcounterpageref{ex}{0}

\begin{ex} %[5D7H44]
Tìm hai số biết tổng là $198$ và tỉ số là $\dfrac{4}{5}$. Tìm số bé.
\shortans{88}
\loigiai{
\textbf{Đáp số:} $88$\\
$198 : (4 + 5) \times 4 = 88$.
}
\end{ex}

\begin{ex} %[5D7H44]
Hiệu của hai số là $64$, số lớn gấp $5$ lần số bé. Tìm số lớn.
\shortans{80}
\loigiai{
\textbf{Đáp số:} $80$\\
$64 : (5 - 1) \times 5 = 80$.
}
\end{ex}

\begin{ex} %[5D7V44]
Tìm $35\%$ của $640\text{ kg}$.
\shortans{224}
\loigiai{
\textbf{Đáp số:} $224$\\
$640 \times 35 : 100 = 224\text{ kg}$.
}
\end{ex}

\begin{ex} %[5D7V44]
Một trường tiểu học có $600$ học sinh, trong đó có $318$ học sinh nữ. Tính tỉ số phần trăm của số học sinh nữ so với toàn trường. (Chỉ ghi số).
\shortans{53}
\loigiai{
\textbf{Đáp số:} $53$\\
$318 : 600 = 53\%$.
}
\end{ex}

\begin{ex} %[5D7V44]
Trên bản đồ tỉ lệ $1 : 20\,000$, một con kênh dài $15\text{ cm}$. Tính độ dài thật của con kênh đó theo ki-lô-mét.
\shortans{3}
\loigiai{
\textbf{Đáp số:} $3$\\
$15 \times 20\,000 = 300\,000\text{ cm} = 3\text{ km}$.
}
\end{ex}

\begin{ex} %[5D7C44]
Một cửa hàng bán được $480\text{ kg}$ gạo tẻ và nếp, thu được $9\,600\,000$ đồng. Biết giá $1\text{ kg}$ gạo nếp đắt hơn gạo tẻ là $5\,000$ đồng và số gạo tẻ gấp $3$ lần số gạo nếp. Tính giá tiền $1\text{ kg}$ gạo tẻ theo đơn vị nghìn đồng.
\shortans{18,75}
\loigiai{
\textbf{Đáp số:} $18{,}75$\\
Số gạo nếp: $480 : (3 + 1) = 120\text{ kg}$.\\
Số gạo tẻ: $480 - 120 = 360\text{ kg}$.\\
Nếu 120 kg nếp bán bằng giá tẻ thì tổng số tiền giảm đi: $120 \times 5\,000 = 600\,000$ đồng.\\
Tổng tiền khi 480 kg đều là gạo tẻ: $9\,600\,000 - 600\,000 = 9\,000\,000$ đồng.\\
Giá 1 kg gạo tẻ: $9\,000\,000 : 480 = 18\,750$ đồng $= 18{,}75$ nghìn đồng.
}
\end{ex}

\newpage

% =========================================================================
% CHỦ ĐỀ 8 - THỂ TÍCH. ĐƠN VỊ ĐO THỂ TÍCH
% =========================================================================
\section{Thể tích. đơn vị đo thể tích}

\TaoPhieuKiemTra{Chủ đề 8 -- Thể tích. đơn vị đo thể tích}{45}{Thể tích của một hình}{https://loptoanthaylinh.vn/thi-online/toan-5/bai-45}
\textbf{Phần 1. Trắc nghiệm lựa chọn}
\setcounterpageref{ex}{0}

\begin{ex} %[5H8N45]
Thể tích của một hình là
\choice
{Diện tích toàn bộ bề mặt bên ngoài của hình đó}
{\True Độ lớn của khoảng không gian mà hình đó chiếm chỗ}
{Tổng độ dài tất cả các cạnh của hình đó}
{Khối lượng của hình đó khi đem cân}
\loigiai{
Thể tích của một hình là độ lớn của khoảng không gian mà hình đó chiếm chỗ.\\
\textbf{Đáp án của câu hỏi là B}
}
\end{ex}

\begin{ex} %[5H8N45]
Cho hình hộp chữ nhật A nằm hoàn toàn bên trong hình hộp chữ nhật B. Kết luận nào sau đây là đúng?
\choice
{Thể tích hình A lớn hơn thể tích hình B}
{Thể tích hình A bằng thể tích hình B}
{Không so sánh được thể tích của hai hình}
{\True Thể tích hình A bé hơn thể tích hình B}
\loigiai{
Vì hình A nằm hoàn toàn trong hình B nên thể tích hình A bé hơn thể tích hình B.\\
\textbf{Đáp án của câu hỏi là D}
}
\end{ex}

\begin{ex} %[5H8N45]
Một hình khối $M$ được ghép bởi $12$ khối lập phương nhỏ bằng nhau. Hình khối $N$ được ghép bởi $15$ khối lập phương nhỏ như thế. Kết luận nào đúng?
\choice
{\True Thể tích hình $N$ lớn hơn thể tích hình $M$}
{Thể tích hình $M$ lớn hơn thể tích hình $N$}
{Thể tích hai hình bằng nhau}
{Thể tích hình $M$ gấp đôi thể tích hình $N$}
\loigiai{
Hình $M$ gồm $12$ khối lập phương nhỏ, hình $N$ gồm $15$ khối lập phương nhỏ nên thể tích hình $N$ lớn hơn thể tích hình $M$.\\
\textbf{Đáp án của câu hỏi là A}
}
\end{ex}

\begin{ex} %[5H8H45]
Hình C gồm $8$ khối lập phương cạnh $1\text{ cm}$. Hình D gồm $8$ khối lập phương cạnh $1\text{ cm}$ xếp theo hình dạng khác. So sánh thể tích của hình C và hình D:
\choice
{Hình C có thể tích lớn hơn hình D}
{Hình D có thể tích lớn hơn hình C}
{Không thể so sánh được}
{\True Hình C và hình D có thể tích bằng nhau}
\loigiai{
Hai hình được ghép từ cùng số lượng các khối lập phương như nhau thì có thể tích bằng nhau.\\
\textbf{Đáp án của câu hỏi là D}
}
\end{ex}

\begin{ex} %[5H8H45]
Hình $P$ được ghép từ hai hình khối $Q$ (gồm $6$ khối lập phương đơn vị) và $R$ (gồm $9$ khối lập phương đơn vị). Thể tích của hình $P$ bằng
\choice
{\True $15$ khối lập phương đơn vị}
{$3$ khối lập phương đơn vị}
{$54$ khối lập phương đơn vị}
{$18$ khối lập phương đơn vị}
\loigiai{
Thể tích hình $P$ bằng tổng thể tích của hai hình $Q$ và $R$: $6 + 9 = 15$ khối lập phương đơn vị.\\
\textbf{Đáp án của câu hỏi là A}
}
\end{ex}

\begin{ex} %[5H8H45]
Một hình lập phương lớn được ghép từ $27$ khối lập phương nhỏ cạnh $1\text{ cm}$. Người ta lấy ra $3$ khối lập phương nhỏ ở $3$ góc. Thể tích của hình còn lại là
\choice
{$27\text{ khối}$}
{$21\text{ khối}$}
{\True $24\text{ khối}$}
{$18\text{ khối}$}
\loigiai{
Thể tích còn lại: $27 - 3 = 24$ khối lập phương nhỏ.\\
\textbf{Đáp án của câu hỏi là C}
}
\end{ex}

\begin{ex} %[5H8V45]
Người ta dùng các khối lập phương cạnh $1\text{ cm}$ để xếp thành một hình hộp chữ nhật có chiều dài $5\text{ cm}$, chiều rộng $4\text{ cm}$ và chiều cao $3\text{ cm}$. Số khối lập phương nhỏ cần dùng là
\choice
{$12$ khối}
{\True $60$ khối}
{$20$ khối}
{$15$ khối}
\loigiai{
Mỗi lớp có $5 \times 4 = 20$ khối. Có $3$ lớp nên tổng số khối là $20 \times 3 = 60$ khối.\\
\textbf{Đáp án của câu hỏi là B}
}
\end{ex}

\begin{ex} %[5H8V45]
Bạn Nam có $48$ khối lập phương nhỏ bằng nhau. Nam muốn xếp tất cả các khối đó thành một hình hộp chữ nhật. Cách xếp nào sau đây KHÔNG thể thực hiện được?
\choice
{Dài $4$, rộng $3$, cao $4$}
{Dài $6$, rộng $4$, cao $2$}
{\True Dài $5$, rộng $3$, cao $3$}
{Dài $8$, rộng $3$, cao $2$}
\loigiai{
Vì $5 \times 3 \times 3 = 45 \neq 48$ nên không thể xếp thành hình hộp chữ nhật có kích thước $5, 3, 3$.\\
\textbf{Đáp án của câu hỏi là C}
}
\end{ex}

\begin{ex} %[5H8V45]
Một hình lập phương lớn có cạnh $3\text{ dm}$ được ghép bởi các hình lập phương nhỏ cạnh $1\text{ dm}$. Người ta sơn $6$ mặt ngoài của hình lập phương lớn. Số hình lập phương nhỏ không được sơn mặt nào là
\choice
{$0$ hình}
{$6$ hình}
{$8$ hình}
{\True $1$ hình}
\loigiai{
Số hình lập phương nhỏ ở tâm không sơn mặt nào là $(3 - 2) \times (3 - 2) \times (3 - 2) = 1$ hình.\\
\textbf{Đáp án của câu hỏi là D}
}
\end{ex}

\begin{ex} %[5H8V45]
Thả một hòn đá chìm hoàn toàn vào một cốc nước hình trụ đang chứa nước. Mực nước trong cốc dâng lên từ vạch $120\text{ ml}$ đến vạch $165\text{ ml}$. Thể tích của hòn đá là
\choice
{$285\text{ ml}$}
{$165\text{ ml}$}
{\True $45\text{ ml}$}
{$120\text{ ml}$}
\loigiai{
Thể tích hòn đá bằng thể tích nước dâng lên: $165 - 120 = 45\text{ (ml)}$.\\
\textbf{Đáp án của câu hỏi là C}
}
\end{ex}

\begin{ex} %[5H8V45]
Có $64$ khối lập phương nhỏ cạnh $1\text{ cm}$ được xếp lại thành một hình lập phương lớn. Cạnh của hình lập phương lớn là
\choice
{$2\text{ cm}$}
{\True $4\text{ cm}$}
{$6\text{ cm}$}
{$8\text{ cm}$}
\loigiai{
Vì $4 \times 4 \times 4 = 64$ nên cạnh hình lập phương lớn là $4\text{ cm}$.\\
\textbf{Đáp án của câu hỏi là B}
}
\end{ex}

\begin{ex} %[5H8C45]
Một khối gỗ hình hộp chữ nhật được ghép từ $120$ khối lập phương nhỏ cạnh $1\text{ cm}$. Người ta sơn đỏ toàn bộ mặt ngoài của khối gỗ rồi tháo rời các khối nhỏ ra. Biết khối gỗ có kích thước $4\text{ cm} \times 5\text{ cm} \times 6\text{ cm}$. Số khối lập phương nhỏ được sơn đúng $2$ mặt là
\choice
{\True $36$ khối}
{$24$ khối}
{$48$ khối}
{$16$ khối}
\loigiai{
Số khối sơn 2 mặt nằm trên các cạnh (trừ 8 đỉnh):\\
$4 \times (4 - 2) + 4 \times (5 - 2) + 4 \times (6 - 2) = 4 \times 2 + 4 \times 3 + 4 \times 4 = 8 + 12 + 16 = 36$ (khối).\\
\textbf{Đáp án của câu hỏi là A}
}
\end{ex}

\textbf{Phần 2. Trắc nghiệm đúng sai}
\setcounterpageref{ex}{0}

\begin{ex} %[5H8H45]
Cho hình hộp chữ nhật $A$ gồm $24$ khối lập phương đơn vị và hình hộp chữ nhật $B$ gồm $30$ khối lập phương đơn vị. Xét tính đúng sai:
\choiceTF
{\True Thể tích hình B lớn hơn thể tích hình A.}
{Thể tích hình A bằng $\dfrac{5}{4}$ thể tích hình B.}
{\True Thể tích hình A bằng $\dfrac{4}{5}$ thể tích hình B.}
{\True Khi ghép hình A và hình B lại ta được một hình mới có thể tích bằng $54$ khối lập phương đơn vị.}
\loigiai{
\begin{itemize}
    \item $30 > 24$ nên thể tích B lớn hơn thể tích A. Đúng.
    \item Tỉ số A so với B là $24 : 30 = \dfrac{4}{5} \neq \dfrac{5}{4}$. Ý 2 Sai, Ý 3 Đúng.
    \item Ghép lại: $24 + 30 = 54$ khối. Đúng.
\end{itemize}
\begin{itemize}
    \item a) \textbf{Đúng}.
    \item b) \textbf{Sai}.
    \item c) \textbf{Đúng}.
    \item d) \textbf{Đúng}.
\end{itemize}
}
\end{ex}

\begin{ex} %[5H8V45]
Xếp các khối lập phương cạnh $1\text{ cm}$ thành một hình lập phương lớn có cạnh $4\text{ cm}$, sau đó sơn tất cả các mặt ngoài. Xét tính đúng sai:
\choiceTF
{\True Cần dùng tất cả $64$ khối lập phương nhỏ.}
{\True Có đúng $8$ khối lập phương nhỏ được sơn $3$ mặt.}
{Có $16$ khối lập phương nhỏ không được sơn mặt nào.}
{\True Có $24$ khối lập phương nhỏ được sơn đúng $1$ mặt.}
\loigiai{
\begin{itemize}
    \item Số khối: $4 \times 4 \times 4 = 64$ khối. Đúng.
    \item Số khối sơn 3 mặt ở 8 góc $= 8$ khối. Đúng.
    \item Số khối không sơn: $(4-2) \times (4-2) \times (4-2) = 8$ khối $\neq 16$ khối. Sai.
    \item Số khối sơn 1 mặt ở 6 mặt: $6 \times (4-2) \times (4-2) = 6 \times 4 = 24$ khối. Đúng.
\end{itemize}
\begin{itemize}
    \item a) \textbf{Đúng}.
    \item b) \textbf{Đúng}.
    \item c) \textbf{Sai}.
    \item d) \textbf{Đúng}.
\end{itemize}
}
\end{ex}

\textbf{Phần 3. Trả lời ngắn}
\setcounterpageref{ex}{0}

\begin{ex} %[5H8H45]
Hình khối $E$ được ghép bởi $18$ khối lập phương nhỏ cạnh $1\text{ cm}$, hình khối $F$ được ghép bởi $25$ khối lập phương nhỏ như thế. Hỏi thể tích hình $F$ lớn hơn thể tích hình $E$ bao nhiêu xăng-ti-mét khối?
\shortans{7}
\loigiai{
\textbf{Đáp số:} $7$\\
$25 - 18 = 7\text{ cm}^3$.
}
\end{ex}

\begin{ex} %[5H8H45]
Người ta xếp các khối lập phương nhỏ cạnh $1\text{ cm}$ thành hình hộp chữ nhật có kích thước dài $6\text{ cm}$, rộng $4\text{ cm}$, cao $5\text{ cm}$. Hỏi cần dùng tất cả bao nhiêu khối lập phương nhỏ?
\shortans{120}
\loigiai{
\textbf{Đáp số:} $120$\\
$6 \times 4 \times 5 = 120$ khối.
}
\end{ex}

\begin{ex} %[5H8V45]
Một hình lập phương lớn cạnh $3\text{ cm}$ được ghép từ các khối lập phương nhỏ cạnh $1\text{ cm}$. Sơn màu đỏ toàn bộ bên ngoài hình lập phương lớn. Hỏi có bao nhiêu khối lập phương nhỏ được sơn đúng $1$ mặt?
\shortans{6}
\loigiai{
\textbf{Đáp số:} $6$\\
Mỗi mặt có $1$ khối ở giữa: $6 \times 1 = 6$ khối.
}
\end{ex}

\begin{ex} %[5H8V45]
Thả chìm một quả cầu sắt vào bình chia độ chứa $250\text{ ml}$ nước thì nước dâng lên mức $380\text{ ml}$. Tính thể tích quả cầu theo mi-li-lít.
\shortans{130}
\loigiai{
\textbf{Đáp số:} $130$\\
$380 - 250 = 130\text{ ml}$.
}
\end{ex}

\begin{ex} %[5H8V45]
Để xếp một hình lập phương có cạnh $5\text{ cm}$ bằng các khối lập phương nhỏ cạnh $1\text{ cm}$, ta cần dùng bao nhiêu khối lập phương nhỏ?
\shortans{125}
\loigiai{
\textbf{Đáp số:} $125$\\
$5 \times 5 \times 5 = 125$ khối.
}
\end{ex}

\begin{ex} %[5H8C45]
Người ta xếp $216$ khối lập phương nhỏ cạnh $1\text{ cm}$ thành một hình lập phương lớn, sau đó sơn tất cả các mặt bên ngoài. Hỏi có bao nhiêu khối lập phương nhỏ không được sơn mặt nào?
\shortans{64}
\loigiai{
\textbf{Đáp số:} $64$\\
Cạnh hình lập phương lớn là $6\text{ cm}$ vì $6 \times 6 \times 6 = 216$.\\
Số khối không sơn mặt nào: $(6 - 2) \times (6 - 2) \times (6 - 2) = 4 \times 4 \times 4 = 64$ khối.
}
\end{ex}

\newpage

\TaoPhieuKiemTra{Chủ đề 8 -- Thể tích. đơn vị đo thể tích}{46}{Xăng-ti-mét khối. Đề-xi-mét khối}{https://loptoanthaylinh.vn/thi-online/toan-5/bai-46}
\textbf{Phần 1. Trắc nghiệm lựa chọn}
\setcounterpageref{ex}{0}

\begin{ex} %[5H8N46]
$1\text{ dm}^3$ bằng bao nhiêu xăng-ti-mét khối?
\choice
{\True $1\,000\text{ cm}^3$}
{$10\text{ cm}^3$}
{$100\text{ cm}^3$}
{$10\,000\text{ cm}^3$}
\loigiai{
Ta có: $1\text{ dm}^3 = 1\,000\text{ cm}^3$.\\
\textbf{Đáp án của câu hỏi là A}
}
\end{ex}

\begin{ex} %[5H8N46]
$1\text{ dm}^3$ bằng bao nhiêu lít?
\choice
{$10\text{ l}$}
{$0{,}1\text{ l}$}
{\True $1\text{ l}$}
{$100\text{ l}$}
\loigiai{
Ta có: $1\text{ dm}^3 = 1\text{ l}$.\\
\textbf{Đáp án của câu hỏi là C}
}
\end{ex}

\begin{ex} %[5H8N46]
Số đo “Bốn nghìn hai trăm năm mươi xăng-ti-mét khối” được viết là
\choice
{$425\text{ cm}^3$}
{$4\,250\text{ dm}^3$}
{$42\,500\text{ cm}^3$}
{\True $4\,250\text{ cm}^3$}
\loigiai{
Viết là $4\,250\text{ cm}^3$.\\
\textbf{Đáp án của câu hỏi là D}
}
\end{ex}

\begin{ex} %[5H8H46]
Đổi $3{,}5\text{ dm}^3$ ra xăng-ti-mét khối ta được
\choice
{$35\text{ cm}^3$}
{$350\text{ cm}^3$}
{\True $3\,500\text{ cm}^3$}
{$35\,000\text{ cm}^3$}
\loigiai{
$3{,}5\text{ dm}^3 = 3{,}5 \times 1\,000 = 3\,500\text{ cm}^3$.\\
\textbf{Đáp án của câu hỏi là C}
}
\end{ex}

\begin{ex} %[5H8H46]
Đổi $7\,800\text{ cm}^3$ ra đề-xi-mét khối ta được
\choice
{$0{,}78\text{ dm}^3$}
{\True $7{,}8\text{ dm}^3$}
{$78\text{ dm}^3$}
{$780\text{ dm}^3$}
\loigiai{
$7\,800\text{ cm}^3 = 7\,800 : 1\,000 = 7{,}8\text{ dm}^3$.\\
\textbf{Đáp án của câu hỏi là B}
}
\end{ex}

\begin{ex} %[5H8H46]
Kết quả của phép tính $450\text{ cm}^3 + 2{,}5\text{ dm}^3$ là
\choice
{\True $2\,950\text{ cm}^3$}
{$452{,}5\text{ cm}^3$}
{$700\text{ cm}^3$}
{$295\text{ cm}^3$}
\loigiai{
Đổi $2{,}5\text{ dm}^3 = 2\,500\text{ cm}^3$.\\
Phép tính: $450 + 2\,500 = 2\,950\text{ (cm}^3\text{)}$.\\
\textbf{Đáp án của câu hỏi là A}
}
\end{ex}

\begin{ex} %[5H8H46]
Một chai nước ngọt có dung tích $1{,}5\text{ l}$. Dung tích chai nước đó theo xăng-ti-mét khối là
\choice
{$150\text{ cm}^3$}
{$15\,000\text{ cm}^3$}
{$15\text{ cm}^3$}
{\True $1\,500\text{ cm}^3$}
\loigiai{
$1{,}5\text{ l} = 1{,}5\text{ dm}^3 = 1\,500\text{ cm}^3$.\\
\textbf{Đáp án của câu hỏi là D}
}
\end{ex}

\begin{ex} %[5H8V46]
Tính giá trị biểu thức: $12\text{ dm}^3 \times 4 - 8\,500\text{ cm}^3$.
\choice
{$40\text{ dm}^3$}
{$38{,}5\text{ dm}^3$}
{\True $39{,}5\text{ dm}^3$}
{$395\text{ dm}^3$}
\loigiai{
$12\text{ dm}^3 \times 4 = 48\text{ dm}^3$.\\
$8\,500\text{ cm}^3 = 8{,}5\text{ dm}^3$.\\
$48 - 8{,}5 = 39{,}5\text{ dm}^3$.\\
\textbf{Đáp án của câu hỏi là C}
}
\end{ex}

\begin{ex} %[5H8V46]
Một bể cá chứa $45\text{ dm}^3$ nước, người ta đổ thêm vào bể $15\,000\text{ cm}^3$ nước. Số lít nước có trong bể lúc này là
\choice
{\True $60\text{ l}$}
{$50\text{ l}$}
{$55\text{ l}$}
{$65\text{ l}$}
\loigiai{
$15\,000\text{ cm}^3 = 15\text{ dm}^3 = 15\text{ l}$.\\
Tổng số nước: $45 + 15 = 60\text{ l}$.\\
\textbf{Đáp án của câu hỏi là A}
}
\end{ex}

\begin{ex} %[5H8V46]
Điền dấu thích hợp vào chỗ chấm: $5\text{ dm}^3\; 45\text{ cm}^3 \;\dots\; 5{,}45\text{ dm}^3$.
\choice
{$>$}
{\True $<$}
{$=$}
{Không so sánh được}
\loigiai{
$5\text{ dm}^3\; 45\text{ cm}^3 = 5{,}045\text{ dm}^3 < 5{,}45\text{ dm}^3$.\\
\textbf{Đáp án của câu hỏi là B}
}
\end{ex}

\begin{ex} %[5H8V46]
Một hình hộp chữ nhật có thể tích $3{,}6\text{ dm}^3$, diện tích đáy là $180\text{ cm}^2$. Chiều cao của hình hộp chữ nhật là
\choice
{$2\text{ cm}$}
{$200\text{ cm}$}
{$0{,}2\text{ cm}$}
{\True $20\text{ cm}$}
\loigiai{
Đổi $3{,}6\text{ dm}^3 = 3\,600\text{ cm}^3$.\\
Chiều cao là: $3\,600 : 180 = 20\text{ (cm)}$.\\
\textbf{Đáp án của câu hỏi là D}
}
\end{ex}

\begin{ex} %[5H8C46]
Một bình tràn chứa đầy nước. Người ta thả chìm $4$ viên bi sắt có thể tích như nhau vào bình thì thấy tràn ra $0{,}32\text{ dm}^3$ nước. Thể tích của mỗi viên bi sắt là
\choice
{$8\text{ cm}^3$}
{\True $80\text{ cm}^3$}
{$800\text{ cm}^3$}
{$0{,}08\text{ cm}^3$}
\loigiai{
Đổi $0{,}32\text{ dm}^3 = 320\text{ cm}^3$.\\
Thể tích mỗi viên bi sắt: $320 : 4 = 80\text{ (cm}^3\text{)}$.\\
\textbf{Đáp án của câu hỏi là B}
}
\end{ex}

\textbf{Phần 2. Trắc nghiệm đúng sai}
\setcounterpageref{ex}{0}

\begin{ex} %[5H8H46]
Xét tính đúng sai của các phép đổi đơn vị sau:
\choiceTF
{\True $4\text{ dm}^3 = 4\,000\text{ cm}^3$.}
{$2{,}05\text{ dm}^3 = 2\,500\text{ cm}^3$.}
{\True $350\text{ cm}^3 = 0{,}35\text{ dm}^3$.}
{\True $8\text{ dm}^3\; 6\text{ cm}^3 = 8{,}006\text{ dm}^3$.}
\loigiai{
\begin{itemize}
    \item $4 \times 1000 = 4000\text{ cm}^3$. Đúng.
    \item $2{,}05 \times 1000 = 2050\text{ cm}^3 \neq 2500\text{ cm}^3$. Sai.
    \item $350 : 1000 = 0{,}35\text{ dm}^3$. Đúng.
    \item $8 + 6 : 1000 = 8{,}006\text{ dm}^3$. Đúng.
\end{itemize}
\begin{itemize}
    \item a) \textbf{Đúng}.
    \item b) \textbf{Sai}.
    \item c) \textbf{Đúng}.
    \item d) \textbf{Đúng}.
\end{itemize}
}
\end{ex}

\begin{ex} %[5H8V46]
Một xô nước chứa $12\text{ l}$ nước. Người ta dùng một ca có dung tích $400\text{ cm}^3$ để múc nước từ xô ra. Xét tính đúng sai:
\choiceTF
{\True Dung tích xô nước bằng $12\text{ dm}^3 = 12\,000\text{ cm}^3$.}
{\True Dung tích chiếc ca là $0{,}4\text{ l}$.}
{Cần múc $25$ ca đầy thì hết nước trong xô.}
{\True Cần múc $30$ ca đầy thì hết nước trong xô.}
\loigiai{
\begin{itemize}
    \item $12\text{ l} = 12\text{ dm}^3 = 12\,000\text{ cm}^3$. Đúng.
    \item $400\text{ cm}^3 = 0{,}4\text{ dm}^3 = 0{,}4\text{ l}$. Đúng.
    \item Số ca cần múc: $12\,000 : 400 = 30$ ca $\neq 25$ ca. Ý 3 Sai, Ý 4 Đúng.
\end{itemize}
\begin{itemize}
    \item a) \textbf{Đúng}.
    \item b) \textbf{Đúng}.
    \item c) \textbf{Sai}.
    \item d) \textbf{Đúng}.
\end{itemize}
}
\end{ex}

\textbf{Phần 3. Trả lời ngắn}
\setcounterpageref{ex}{0}

\begin{ex} %[5H8H46]
Đổi $6{,}75\text{ dm}^3$ ra xăng-ti-mét khối.
\shortans{6750}
\loigiai{
\textbf{Đáp số:} $6750$\\
$6{,}75 \times 1\,000 = 6\,750\text{ cm}^3$.
}
\end{ex}

\begin{ex} %[5H8H46]
Đổi $12\,400\text{ cm}^3$ ra đề-xi-mét khối (dùng dấu phẩy nếu có).
\shortans{12,4}
\loigiai{
\textbf{Đáp số:} $12{,}4$\\
$12\,400 : 1\,000 = 12{,}4\text{ dm}^3$.
}
\end{ex}

\begin{ex} %[5H8V46]
Tính kết quả phép tính: $3{,}2\text{ dm}^3 + 4\,800\text{ cm}^3$ theo đơn vị đề-xi-mét khối.
\shortans{8}
\loigiai{
\textbf{Đáp số:} $8$\\
$3{,}2 + 4{,}8 = 8\text{ dm}^3$.
}
\end{ex}

\begin{ex} %[5H8V46]
Một hình hộp chữ nhật có thể tích là $5{,}4\text{ dm}^3$. Biết diện tích đáy là $300\text{ cm}^2$. Tính chiều cao hình hộp chữ nhật theo xăng-ti-mét.
\shortans{18}
\loigiai{
\textbf{Đáp số:} $18$\\
$5{,}4\text{ dm}^3 = 5\,400\text{ cm}^3 \Rightarrow 5\,400 : 300 = 18\text{ cm}$.
}
\end{ex}

\begin{ex} %[5H8V46]
Một bình nước đang chứa $1{,}8\text{ dm}^3$ nước. Người ta rót đều vào $9$ chiếc cốc như nhau. Hỏi mỗi cốc chứa bao nhiêu xăng-ti-mét khối nước?
\shortans{200}
\loigiai{
\textbf{Đáp số:} $200$\\
$1{,}8\text{ dm}^3 = 1\,800\text{ cm}^3 \Rightarrow 1\,800 : 9 = 200\text{ cm}^3$.
}
\end{ex}

\begin{ex} %[5H8C46]
Một bể cá mini hình hộp chữ nhật có diện tích đáy $250\text{ cm}^2$. Thả một khối đá vào bể thì nước dâng lên thêm $1{,}6\text{ cm}$ (khối đá ngập hoàn toàn). Tính thể tích khối đá theo đề-xi-mét khối.
\shortans{0,4}
\loigiai{
\textbf{Đáp số:} $0{,}4$\\
Thể tích khối đá: $250 \times 1{,}6 = 400\text{ cm}^3 = 0{,}4\text{ dm}^3$.
}
\end{ex}

\newpage

\TaoPhieuKiemTra{Chủ đề 8 -- Thể tích. đơn vị đo thể tích}{47}{Mét khối}{https://loptoanthaylinh.vn/thi-online/toan-5/bai-47}
\textbf{Phần 1. Trắc nghiệm lựa chọn}
\setcounterpageref{ex}{0}

\begin{ex} %[5H8N47]
$1\text{ m}^3$ bằng bao nhiêu đề-xi-mét khối?
\choice
{$10\text{ dm}^3$}
{\True $1\,000\text{ dm}^3$}
{$100\text{ dm}^3$}
{$1\,000\,000\text{ dm}^3$}
\loigiai{
$1\text{ m}^3 = 1\,000\text{ dm}^3$.\\
\textbf{Đáp án của câu hỏi là B}
}
\end{ex}

\begin{ex} %[5H8N47]
$1\text{ m}^3$ bằng bao nhiêu xăng-ti-mét khối?
\choice
{$1\,000\text{ cm}^3$}
{$10\,000\text{ cm}^3$}
{\True $1\,000\,000\text{ cm}^3$}
{$100\,000\text{ cm}^3$}
\loigiai{
$1\text{ m}^3 = 1\,000\,000\text{ cm}^3$.\\
\textbf{Đáp án của câu hỏi là C}
}
\end{ex}

\begin{ex} %[5H8N47]
Số “Hai phẩy không năm mét khối” được viết là
\choice
{$205\text{ m}^3$}
{$2{,}5\text{ m}^3$}
{$2{,}005\text{ m}^3$}
{\True $2{,}05\text{ m}^3$}
\loigiai{
Viết là $2{,}05\text{ m}^3$.\\
\textbf{Đáp án của câu hỏi là D}
}
\end{ex}

\begin{ex} %[5H8H47]
Đổi $4{,}8\text{ m}^3$ ra đề-xi-mét khối ta được
\choice
{$48\text{ dm}^3$}
{$480\text{ dm}^3$}
{\True $4\,800\text{ dm}^3$}
{$48\,000\text{ dm}^3$}
\loigiai{
$4{,}8 \times 1\,000 = 4\,800\text{ dm}^3$.\\
\textbf{Đáp án của câu hỏi là C}
}
\end{ex}

\begin{ex} %[5H8H47]
Đổi $2\,500\text{ dm}^3$ ra mét khối ta được
\choice
{$0{,}25\text{ m}^3$}
{$25\text{ m}^3$}
{$250\text{ m}^3$}
{\True $2{,}5\text{ m}^3$}
\loigiai{
$2\,500 : 1\,000 = 2{,}5\text{ m}^3$.\\
\textbf{Đáp án của câu hỏi là D}
}
\end{ex}

\begin{ex} %[5H8H47]
$1\text{ m}^3$ nước tương ứng với bao nhiêu lít nước?
\choice
{$100\text{ l}$}
{\True $1\,000\text{ l}$}
{$10\,000\text{ l}$}
{$1\text{ l}$}
\loigiai{
Vì $1\text{ m}^3 = 1\,000\text{ dm}^3$ và $1\text{ dm}^3 = 1\text{ l}$ nên $1\text{ m}^3 = 1\,000\text{ l}$.\\
\textbf{Đáp án của câu hỏi là B}
}
\end{ex}

\begin{ex} %[5H8V47]
Một bể nước hình hộp chữ nhật chứa được tối đa $3{,}6\text{ m}^3$ nước. Dung tích của bể nước đó là bao nhiêu lít?
\choice
{\True $3\,600\text{ l}$}
{$360\text{ l}$}
{$36\,000\text{ l}$}
{$36\text{ l}$}
\loigiai{
$3{,}6\text{ m}^3 = 3{,}6 \times 1\,000 = 3\,600\text{ l}$.\\
\textbf{Đáp án của câu hỏi là A}
}
\end{ex}

\begin{ex} %[5H8V47]
Tính giá trị biểu thức: $4\text{ m}^3 - 1\,200\text{ dm}^3 \times 2$.
\choice
{$2\text{ m}^3$}
{$2{,}8\text{ m}^3$}
{\True $1{,}6\text{ m}^3$}
{$1{,}8\text{ m}^3$}
\loigiai{
$1\,200\text{ dm}^3 \times 2 = 2\,400\text{ dm}^3 = 2{,}4\text{ m}^3$.\\
$4 - 2{,}4 = 1{,}6\text{ m}^3$.\\
\textbf{Đáp án của câu hỏi là C}
}
\end{ex}

\begin{ex} %[5H8V47]
Một xe téc chở $15\text{ m}^3$ xăng. Người ta đã bơm từ xe téc vào các bồn chứa $8\,500\text{ l}$ xăng. Số mét khối xăng còn lại trên xe là
\choice
{\True $6{,}5\text{ m}^3$}
{$7\text{ m}^3$}
{$6\text{ m}^3$}
{$5{,}5\text{ m}^3$}
\loigiai{
Đổi $8\,500\text{ l} = 8{,}5\text{ m}^3$.\\
Số xăng còn lại: $15 - 8{,}5 = 6{,}5\text{ (m}^3\text{)}$.\\
\textbf{Đáp án của câu hỏi là A}
}
\end{ex}

\begin{ex} %[5H8V47]
Điền dấu thích hợp: $2\text{ m}^3\; 85\text{ dm}^3 \;\dots\; 2{,}85\text{ m}^3$.
\choice
{$>$}
{\True $<$}
{$=$}
{Không so sánh được}
\loigiai{
$2\text{ m}^3\; 85\text{ dm}^3 = 2{,}085\text{ m}^3 < 2{,}85\text{ m}^3$.\\
\textbf{Đáp án của câu hỏi là B}
}
\end{ex}

\begin{ex} %[5H8V47]
Một căn phòng có thể tích không khí là $72\text{ m}^3$. Biết chiều cao phòng là $3\text{ m}$. Diện tích nền phòng học là
\choice
{\True $24\text{ m}^2$}
{$216\text{ m}^2$}
{$36\text{ m}^2$}
{$48\text{ m}^2$}
\loigiai{
Diện tích nền phòng: $72 : 3 = 24\text{ (m}^2\text{)}$.\\
\textbf{Đáp án của câu hỏi là A}
}
\end{ex}

\begin{ex} %[5H8C47]
Một bể bơi có thể tích $180\text{ m}^3$. Người ta dùng máy bơm có công suất $15\text{ m}^3/\text{giờ}$ để bơm nước vào bể. Biết trong bể đã có sẵn $30\text{ m}^3$ nước. Thời gian để bơm đầy bể là
\choice
{$12\text{ giờ}$}
{$8\text{ giờ}$}
{$14\text{ giờ}$}
{\True $10\text{ giờ}$}
\loigiai{
Lượng nước cần bơm thêm: $180 - 30 = 150\text{ (m}^3\text{)}$.\\
Thời gian bơm đầy bể: $150 : 15 = 10\text{ (giờ)}$.\\
\textbf{Đáp án của câu hỏi là D}
}
\end{ex}

\textbf{Phần 2. Trắc nghiệm đúng sai}
\setcounterpageref{ex}{0}

\begin{ex} %[5H8H47]
Xét tính đúng sai của các khẳng định sau:
\choiceTF
{\True $1\text{ m}^3 = 1\,000\text{ dm}^3 = 1\,000\,000\text{ cm}^3$.}
{$0{,}75\text{ m}^3 = 75\text{ dm}^3$.}
{\True $2\text{ m}^3\; 50\text{ dm}^3 = 2{,}05\text{ m}^3$.}
{\True $12\,500\text{ dm}^3 = 12{,}5\text{ m}^3$.}
\loigiai{
\begin{itemize}
    \item $1\text{ m}^3 = 1000\text{ dm}^3 = 1000000\text{ cm}^3$. Đúng.
    \item $0{,}75 \times 1000 = 750\text{ dm}^3 \neq 75\text{ dm}^3$. Sai.
    \item $2 + 50 : 1000 = 2{,}05\text{ m}^3$. Đúng.
    \item $12500 : 1000 = 12{,}5\text{ m}^3$. Đúng.
\end{itemize}
\begin{itemize}
    \item a) \textbf{Đúng}.
    \item b) \textbf{Sai}.
    \item c) \textbf{Đúng}.
    \item d) \textbf{Đúng}.
\end{itemize}
}
\end{ex}

\begin{ex} %[5H8V47]
Một bể nước dạng hình hộp chữ nhật có thể tích $12\text{ m}^3$. Hiện lượng nước trong bể chiếm $65\%$ thể tích bể. Xét tính đúng sai:
\choiceTF
{\True Dung tích toàn bộ của bể là $12\,000\text{ l}$.}
{\True Lượng nước hiện có trong bể là $7{,}8\text{ m}^3$.}
{Lượng nước cần bơm thêm để đầy bể là $5\,000\text{ l}$.}
{\True Lượng nước cần bơm thêm để đầy bể là $4\,200\text{ l}$.}
\loigiai{
\begin{itemize}
    \item Dung tích: $12 \times 1000 = 12000\text{ l}$. Đúng.
    \item Nước hiện có: $12 \times 65\% = 7{,}8\text{ m}^3 = 7800\text{ l}$. Đúng.
    \item Cần thêm: $12000 - 7800 = 4200\text{ l} \neq 5000\text{ l}$. Ý 3 Sai, Ý 4 Đúng.
\end{itemize}
\begin{itemize}
    \item a) \textbf{Đúng}.
    \item b) \textbf{Đúng}.
    \item c) \textbf{Sai}.
    \item d) \textbf{Đúng}.
\end{itemize}
}
\end{ex}

\textbf{Phần 3. Trả lời ngắn}
\setcounterpageref{ex}{0}

\begin{ex} %[5H8H47]
Đổi $5{,}25\text{ m}^3$ ra đề-xi-mét khối.
\shortans{5250}
\loigiai{
\textbf{Đáp số:} $5250$\\
$5{,}25 \times 1\,000 = 5\,250\text{ dm}^3$.
}
\end{ex}

\begin{ex} %[5H8H47]
Đổi $3\,400\text{ dm}^3$ ra mét khối (dùng dấu phẩy nếu có).
\shortans{3,4}
\loigiai{
\textbf{Đáp số:} $3{,}4$\\
$3\,400 : 1\,000 = 3{,}4\text{ m}^3$.
}
\end{ex}

\begin{ex} %[5H8V47]
Một bể bơi có thể tích chứa là $45\text{ m}^3$. Hỏi bể bơi đó chứa được bao nhiêu nghìn lít nước?
\shortans{45}
\loigiai{
\textbf{Đáp số:} $45$\\
$45\text{ m}^3 = 45\,000\text{ l} = 45$ nghìn lít.
}
\end{ex}

\begin{ex} %[5H8V47]
Thực hiện phép tính: $8{,}5\text{ m}^3 - 3\,700\text{ dm}^3$. Kết quả bằng bao nhiêu mét khối?
\shortans{4,8}
\loigiai{
\textbf{Đáp số:} $4{,}8$\\
$8{,}5 - 3{,}7 = 4{,}8\text{ m}^3$.
}
\end{ex}

\begin{ex} %[5H8V47]
Một lớp học hình hộp chữ nhật có chiều dài $8\text{ m}$, chiều rộng $6\text{ m}$, chiều cao $3{,}5\text{ m}$. Tính thể tích không khí trong phòng học theo mét khối.
\shortans{168}
\loigiai{
\textbf{Đáp số:} $168$\\
$8 \times 6 \times 3{,}5 = 168\text{ m}^3$.
}
\end{ex}

\begin{ex} %[5H8C47]
Một bể nước hình hộp chữ nhật chứa $18\text{ m}^3$ nước. Mỗi ngày gia đình dùng hết $750\text{ l}$ nước. Hỏi sau bao nhiêu ngày thì gia đình dùng hết số nước trong bể?
\shortans{24}
\loigiai{
\textbf{Đáp số:} $24$\\
$18\text{ m}^3 = 18\,000\text{ l} \Rightarrow 18\,000 : 750 = 24\text{ ngày}$.
}
\end{ex}

\newpage

\TaoPhieuKiemTra{Chủ đề 8 -- Thể tích. đơn vị đo thể tích}{48}{Luyện tập chung đơn vị đo thể tích}{https://loptoanthaylinh.vn/thi-online/toan-5/bai-48}
\textbf{Phần 1. Trắc nghiệm lựa chọn}
\setcounterpageref{ex}{0}

\begin{ex} %[5H8N48]
Đơn vị đo thể tích nào sau đây lớn nhất?
\choice
{\True Mét khối}
{Xăng-ti-mét khối}
{Đề-xi-mét khối}
{Lít}
\loigiai{
Mét khối ($1\text{ m}^3 = 1\,000\text{ dm}^3 = 1\,000\,000\text{ cm}^3$) là đơn vị lớn nhất.\\
\textbf{Đáp án của câu hỏi là A}
}
\end{ex}

\begin{ex} %[5H8N48]
$3\text{ m}^3\; 45\text{ dm}^3$ được viết dưới dạng số thập phân có đơn vị mét khối là
\choice
{$3{,}45\text{ m}^3$}
{$3{,}450\text{ m}^3$}
{\True $3{,}045\text{ m}^3$}
{$34{,}5\text{ m}^3$}
\loigiai{
$3\text{ m}^3\; 45\text{ dm}^3 = 3 + \dfrac{45}{1\,000} = 3{,}045\text{ m}^3$.\\
\textbf{Đáp án của câu hỏi là C}
}
\end{ex}

\begin{ex} %[5H8H48]
Sắp xếp các số đo sau theo thứ tự từ bé đến lớn: $1{,}5\text{ m}^3$; $1\,200\text{ dm}^3$; $1\,600\,000\text{ cm}^3$.
\choice
{$1{,}5\text{ m}^3 < 1\,200\text{ dm}^3 < 1\,600\,000\text{ cm}^3$}
{$1\,600\,000\text{ cm}^3 < 1{,}5\text{ m}^3 < 1\,200\text{ dm}^3$}
{$1\,200\text{ dm}^3 < 1\,600\,000\text{ cm}^3 < 1{,}5\text{ m}^3$}
{\True $1\,200\text{ dm}^3 < 1{,}5\text{ m}^3 < 1\,600\,000\text{ cm}^3$}
\loigiai{
Đổi về $\text{dm}^3$: $1{,}5\text{ m}^3 = 1\,500\text{ dm}^3$; $1\,600\,000\text{ cm}^3 = 1\,600\text{ dm}^3$.\\
Thứ tự: $1\,200\text{ dm}^3 < 1{,}5\text{ m}^3 < 1\,600\,000\text{ cm}^3$.\\
\textbf{Đáp án của câu hỏi là D}
}
\end{ex}

\begin{ex} %[5H8H48]
Giá trị của biểu thức $0{,}75\text{ m}^3 + 350\text{ dm}^3 - 120\,000\text{ cm}^3$ theo đơn vị đề-xi-mét khối là
\choice
{$1\,100\text{ dm}^3$}
{\True $980\text{ dm}^3$}
{$880\text{ dm}^3$}
{$950\text{ dm}^3$}
\loigiai{
$750 + 350 - 120 = 980\text{ (dm}^3\text{)}$.\\
\textbf{Đáp án của câu hỏi là B}
}
\end{ex}

\begin{ex} %[5H8H48]
Một khối gỗ hình lập phương có thể tích $0{,}064\text{ m}^3$. Thể tích khối gỗ đó theo xăng-ti-mét khối là
\choice
{\True $64\,000\text{ cm}^3$}
{$640\text{ cm}^3$}
{$6\,400\text{ cm}^3$}
{$640\,000\text{ cm}^3$}
\loigiai{
$0{,}064 \times 1\,000\,000 = 64\,000\text{ cm}^3$.\\
\textbf{Đáp án của câu hỏi là A}
}
\end{ex}

\begin{ex} %[5H8H48]
Một bể nước chứa $2{,}4\text{ m}^3$ nước. Người ta đã dùng hết $40\%$ lượng nước trong bể. Lượng nước còn lại trong bể là
\choice
{$960\text{ l}$}
{$1\,200\text{ l}$}
{\True $1\,440\text{ l}$}
{$1\,600\text{ l}$}
\loigiai{
Lượng nước ban đầu: $2\,4\text{ m}^3 = 2\,400\text{ l}$.\\
Nước còn lại: $2\,400 \times 60\% = 1\,440\text{ l}$.\\
\textbf{Đáp án của câu hỏi là C}
}
\end{ex}

\begin{ex} %[5H8V48]
Một hình hộp chữ nhật có chiều dài $2\text{ m}$, chiều rộng $1{,}5\text{ m}$ và chiều cao $1{,}2\text{ m}$. Thể tích của hình hộp chữ nhật đó là
\choice
{$3\text{ m}^3$}
{$4{,}2\text{ m}^3$}
{$4{,}8\text{ m}^3$}
{\True $3{,}6\text{ m}^3$}
\loigiai{
$V = 2 \times 1{,}5 \times 1{,}2 = 3{,}6\text{ (m}^3\text{)}$.\\
\textbf{Đáp án của câu hỏi là D}
}
\end{ex}

\begin{ex} %[5H8V48]
Một bể nước hình hộp chữ nhật có diện tích đáy là $4{,}5\text{ m}^2$. Đang chứa $5\,400\text{ l}$ nước. Mực nước trong bể cao bao nhiêu mét?
\choice
{$1\text{ m}$}
{\True $1{,}2\text{ m}$}
{$1{,}5\text{ m}$}
{$0{,}8\text{ m}$}
\loigiai{
$5\,400\text{ l} = 5{,}4\text{ m}^3$.\\
Chiều cao mực nước: $5{,}4 : 4{,}5 = 1{,}2\text{ (m)}$.\\
\textbf{Đáp án của câu hỏi là B}
}
\end{ex}

\begin{ex} %[5H8V48]
Một khối kim loại hình lập phương cạnh $0{,}5\text{ m}$. Biết mỗi đề-xi-mét khối kim loại nặng $7{,}8\text{ kg}$. Khối lượng của khối kim loại đó là
\choice
{$97{,}5\text{ kg}$}
{$9\,750\text{ kg}$}
{\True $975\text{ kg}$}
{$390\text{ kg}$}
\loigiai{
Thể tích: $0{,}5 \times 0{,}5 \times 0{,}5 = 0{,}125\text{ m}^3 = 125\text{ dm}^3$.\\
Khối lượng: $125 \times 7{,}8 = 975\text{ (kg)}$.\\
\textbf{Đáp án của câu hỏi là C}
}
\end{ex}

\begin{ex} %[5H8V48]
Một bình tràn chứa đầy nước. Người ta thả vào bình một hòn đá thì có $2{,}5\text{ l}$ nước tràn ra. Thể tích của hòn đá đó là
\choice
{$25\text{ cm}^3$}
{\True $2\,500\text{ cm}^3$}
{$250\text{ cm}^3$}
{$0{,}25\text{ m}^3$}
\loigiai{
$2{,}5\text{ l} = 2{,}5\text{ dm}^3 = 2\,500\text{ cm}^3$.\\
\textbf{Đáp án của câu hỏi là B}
}
\end{ex}

\begin{ex} %[5H8V48]
Có hai bể nước hình hộp chữ nhật. Bể thứ nhất chứa $4\text{ m}^3\; 50\text{ dm}^3$ nước, bể thứ hai chứa $3\text{ m}^3\; 800\text{ dm}^3$ nước. Cả hai bể chứa tất cả số lít nước là
\choice
{\True $7\,850\text{ l}$}
{$7\,850\text{ l}$}
{$8\,300\text{ l}$}
{$7\,580\text{ l}$}
\loigiai{
Bể 1: $4\,050\text{ l}$; Bể 2: $3\,800\text{ l}$.\\
Tổng cộng: $4\,050 + 3\,800 = 7\,850\text{ (l)}$.\\
\textbf{Đáp án của câu hỏi là A}
}
\end{ex}

\begin{ex} %[5H8C48]
Một bể nước hình hộp chữ nhật có chiều dài $2{,}5\text{ m}$, chiều rộng $1{,}8\text{ m}$ và chiều cao $1{,}5\text{ m}$. Bể đang cạn nước, người ta mở vòi nước chảy vào bể mỗi phút được $45\text{ l}$. Sau bao lâu thì nước đầy $80\%$ bể?
\choice
{$100\text{ phút}$}
{$150\text{ phút}$}
{$90\text{ phút}$}
{\True $120\text{ phút}$}
\loigiai{
Thể tích bể: $2{,}5 \times 1{,}8 \times 1{,}5 = 6{,}75\text{ m}^3 = 6\,750\text{ l}$.\\
$80\%$ thể tích bể: $6\,750 \times 80\% = 5\,400\text{ l}$.\\
Thời gian chảy: $5\,400 : 45 = 120\text{ (phút)} = 2\text{ giờ}$.\\
\textbf{Đáp án của câu hỏi là D}
}
\end{ex}

\textbf{Phần 2. Trắc nghiệm đúng sai}
\setcounterpageref{ex}{0}

\begin{ex} %[5H8H48]
Xét tính đúng sai của các biểu thức so sánh thể tích:
\choiceTF
{\True $3\text{ m}^3\; 15\text{ dm}^3 = 3{,}015\text{ m}^3$.}
{$4{,}5\text{ m}^3 < 4\,500\text{ dm}^3$.}
{\True $2\text{ dm}^3\; 50\text{ cm}^3 > 2{,}005\text{ dm}^3$.}
{\True $1\,500\,000\text{ cm}^3 = 1{,}5\text{ m}^3$.}
\loigiai{
\begin{itemize}
    \item $3 + 15 : 1000 = 3{,}015\text{ m}^3$. Đúng.
    \item $4{,}5\text{ m}^3 = 4500\text{ dm}^3$. Dấu $<$ là Sai.
    \item $2{,}05\text{ dm}^3 > 2{,}005\text{ dm}^3$. Đúng.
    \item $1500000 : 1000000 = 1{,}5\text{ m}^3$. Đúng.
\end{itemize}
\begin{itemize}
    \item a) \textbf{Đúng}.
    \item b) \textbf{Sai}.
    \item c) \textbf{Đúng}.
    \item d) \textbf{Đúng}.
\end{itemize}
}
\end{ex}

\begin{ex} %[5H8V48]
Một khối gỗ hình hộp chữ nhật có chiều dài $1{,}5\text{ m}$, chiều rộng $0{,}8\text{ m}$ và chiều cao $0{,}5\text{ m}$. Biết mỗi đề-xi-mét khối gỗ nặng $0{,}7\text{ kg}$. Xét tính đúng sai:
\choiceTF
{\True Thể tích khối gỗ là $0{,}6\text{ m}^3$.}
{\True Thể tích khối gỗ bằng $600\text{ dm}^3$.}
{Khối lượng của khối gỗ là $42\text{ kg}$.}
{\True Khối lượng của khối gỗ là $420\text{ kg}$.}
\loigiai{
\begin{itemize}
    \item Thể tích: $1{,}5 \times 0{,}8 \times 0{,}5 = 0{,}6\text{ m}^3$. Đúng.
    \item $0{,}6\text{ m}^3 = 600\text{ dm}^3$. Đúng.
    \item Khối lượng: $600 \times 0{,}7 = 420\text{ kg} \neq 42\text{ kg}$. Ý 3 Sai, Ý 4 Đúng.
\end{itemize}
\begin{itemize}
    \item a) \textbf{Đúng}.
    \item b) \textbf{Đúng}.
    \item c) \textbf{Sai}.
    \item d) \textbf{Đúng}.
\end{itemize}
}
\end{ex}

\textbf{Phần 3. Trả lời ngắn}
\setcounterpageref{ex}{0}

\begin{ex} %[5H8H48]
Đổi $4\text{ m}^3\; 65\text{ dm}^3$ ra đề-xi-mét khối.
\shortans{4065}
\loigiai{
\textbf{Đáp số:} $4065$\\
$4 \times 1\,000 + 65 = 4\,065\text{ dm}^3$.
}
\end{ex}

\begin{ex} %[5H8H48]
Đổi $750\,000\text{ cm}^3$ ra mét khối (dùng dấu phẩy nếu có).
\shortans{0,75}
\loigiai{
\textbf{Đáp số:} $0{,}75$\\
$750\,000 : 1\,000\,000 = 0{,}75\text{ m}^3$.
}
\end{ex}

\begin{ex} %[5H8V48]
Một bể bơi dạng hình hộp chữ nhật có kích thước dài $15\text{ m}$, rộng $6\text{ m}$, sâu $1{,}5\text{ m}$. Tính thể tích nước tối đa mà bể có thể chứa theo mét khối.
\shortans{135}
\loigiai{
\textbf{Đáp số:} $135$\\
$15 \times 6 \times 1{,}5 = 135\text{ m}^3$.
}
\end{ex}

\begin{ex} %[5H8V48]
Một thùng chứa $1{,}8\text{ m}^3$ dầu. Người ta chia đều vào các can $20\text{ l}$. Hỏi chia được bao nhiêu can dầu như thế?
\shortans{90}
\loigiai{
\textbf{Đáp số:} $90$\\
$1{,}8\text{ m}^3 = 1\,800\text{ l} \Rightarrow 1\,800 : 20 = 90\text{ can}$.
}
\end{ex}

\begin{ex} %[5H8V48]
Một bể nước hình hộp chữ nhật có diện tích đáy $3\text{ m}^2$, chiều cao $1{,}2\text{ m}$. Bể đang chứa $2\,700\text{ l}$ nước. Hỏi cần đổ thêm bao nhiêu lít nước nữa để đầy bể?
\shortans{900}
\loigiai{
\textbf{Đáp số:} $900$\\
Thể tích bể: $3 \times 1{,}2 = 3{,}6\text{ m}^3 = 3\,600\text{ l}$.\\
Cần thêm: $3\,600 - 2\,700 = 900\text{ l}$.
}
\end{ex}

\begin{ex} %[5H8C48]
Một bể hình hộp chữ nhật có chiều dài $2\text{ m}$, chiều rộng $1{,}5\text{ m}$ và mực nước đang cao $0{,}8\text{ m}$. Người ta thả vào bể một hòn non bộ thì mực nước trong bể dâng lên cao $0{,}95\text{ m}$ (hòn non bộ chìm hoàn toàn). Tính thể tích hòn non bộ theo đề-xi-mét khối.
\shortans{450}
\loigiai{
\textbf{Đáp số:} $450$\\
Nước dâng thêm: $0{,}95 - 0{,}8 = 0{,}15\text{ m}$.\\
Thể tích hòn non bộ: $2 \times 1{,}5 \times 0{,}15 = 0{,}45\text{ m}^3 = 450\text{ dm}^3$.
}
\end{ex}

\newpage

% =========================================================================
% CHỦ ĐỀ 9 - DIỆN TÍCH VÀ THỂ TÍCH CỦA MỘT SỐ HÌNH KHỐI
% =========================================================================
\section{Diện tích và thể tích của một số hình khối}

\TaoPhieuKiemTra{Chủ đề 9 -- Diện tích và thể tích của một số hình khối}{49}{Hình khai triển hình lập phương, hình hộp chữ nhật, hình trụ}{https://loptoanthaylinh.vn/thi-online/toan-5/bai-49}
\textbf{Phần 1. Trắc nghiệm lựa chọn}
\setcounterpageref{ex}{0}

\begin{ex} %[5H9N49]
Hình khai triển của một hình lập phương gồm bao nhiêu hình vuông bằng nhau?
\choice
{$4$ hình vuông}
{$8$ hình vuông}
{\True $6$ hình vuông}
{$12$ hình vuông}
\loigiai{
Hình lập phương có $6$ mặt là các hình vuông bằng nhau nên hình khai triển gồm $6$ hình vuông bằng nhau.\\
\textbf{Đáp án của câu hỏi là C}
}
\end{ex}

\begin{ex} %[5H9N49]
Hình khai triển của một hình trụ gồm những hình nào?
\choice
{$1$ hình tròn và $1$ hình chữ nhật}
{\True $2$ hình tròn bằng nhau và $1$ hình chữ nhật}
{$2$ hình vuông và $1$ hình chữ nhật}
{$2$ hình chữ nhật và $1$ hình tròn}
\loigiai{
Hình khai triển của hình trụ gồm $2$ đáy là hai hình tròn bằng nhau và mặt xung quanh là một hình chữ nhật.\\
\textbf{Đáp án của câu hỏi là B}
}
\end{ex}

\begin{ex} %[5H9N49]
Hình hộp chữ nhật có tất cả bao nhiêu đỉnh và bao nhiêu cạnh?
\choice
{\True $8$ đỉnh, $12$ cạnh}
{$6$ đỉnh, $8$ cạnh}
{$12$ đỉnh, $8$ cạnh}
{$8$ đỉnh, $6$ cạnh}
\loigiai{
Hình hộp chữ nhật có $8$ đỉnh, $12$ cạnh và $6$ mặt.\\
\textbf{Đáp án của câu hỏi là A}
}
\end{ex}

\begin{ex} %[5H9H49]
Khi mở mặt xung quanh của một hình trụ theo đường sinh, ta nhận được một hình chữ nhật có chiều dài bằng
\choice
{Chiều cao của hình trụ}
{Bán kính của hình tròn đáy}
{Diện tích của hình tròn đáy}
{\True Chu vi của hình tròn đáy}
\loigiai{
Chiều dài của hình chữ nhật bao quanh chính bằng chu vi đáy của hình trụ.\\
\textbf{Đáp án của câu hỏi là D}
}
\end{ex}

\begin{ex} %[5H9H49]
Cho một hình khai triển gồm $6$ hình vuông có cạnh $4\text{ cm}$. Khi gấp lại thành hình lập phương thì diện tích một mặt của hình lập phương đó là
\choice
{\True $16\text{ cm}^2$}
{$8\text{ cm}^2$}
{$64\text{ cm}^2$}
{$96\text{ cm}^2$}
\loigiai{
Diện tích một mặt: $4 \times 4 = 16\text{ (cm}^2\text{)}$.\\
\textbf{Đáp án của câu hỏi là A}
}
\end{ex}

\begin{ex} %[5H9H49]
Hình khai triển của hình lập phương KHÔNG thể có dạng nào sau đây?
\choice
{$4$ hình vuông thành một hàng và $2$ hình vuông ở hai bên đối diện}
{Dạng bậc thang $2 - 2 - 2$}
{\True $6$ hình vuông xếp liên tiếp thành $1$ hàng thẳng}
{Dạng $3 - 3$}
\loigiai{
Nếu $6$ hình vuông xếp thành $1$ hàng thẳng thì khi gấp lại sẽ bị đè lên nhau, không tạo thành hình hộp kín.\\
\textbf{Đáp án của câu hỏi là C}
}
\end{ex}

\begin{ex} %[5H9V49]
Một hình trụ có bán kính đáy $3\text{ cm}$, chiều cao $8\text{ cm}$. Khi khai triển mặt xung quanh của hình trụ này ta được một hình chữ nhật có diện tích là
\choice
{$75{,}36\text{ cm}^2$}
{\True $150{,}72\text{ cm}^2$}
{$120\text{ cm}^2$}
{$226{,}08\text{ cm}^2$}
\loigiai{
Chu vi đáy hình trụ: $2 \times 3{,}14 \times 3 = 18{,}84\text{ (cm)}$.\\
Diện tích mặt xung quanh (diện tích hình chữ nhật): $18{,}84 \times 8 = 150{,}72\text{ (cm}^2\text{)}$.\\
\textbf{Đáp án của câu hỏi là B}
}
\end{ex}

\begin{ex} %[5H9V49]
Người ta cắt một tấm bìa để làm một hình hộp chữ nhật không nắp có kích thước dài $10\text{ cm}$, rộng $6\text{ cm}$, cao $4\text{ cm}$. Tổng diện tích các mặt của tấm bìa hình khai triển là
\choice
{$128\text{ cm}^2$}
{$248\text{ cm}^2$}
{$240\text{ cm}^2$}
{\True $188\text{ cm}^2$}
\loigiai{
Hình hộp không nắp gồm $1$ mặt đáy và $4$ mặt xung quanh.\\
$S_{xq} = (10 + 6) \times 2 \times 4 = 128\text{ (cm}^2\text{)}$.\\
$S_{đáy} = 10 \times 6 = 60\text{ (cm}^2\text{)}$.\\
Tổng diện tích: $128 + 60 = 188\text{ (cm}^2\text{)}$.\\
\textbf{Đáp án của câu hỏi là D}
}
\end{ex}

\begin{ex} %[5H9V49]
Một hình khai triển gồm $6$ hình chữ nhật ghép lại thành hình hộp chữ nhật có các kích thước $5\text{ cm}, 4\text{ cm}, 3\text{ cm}$. Tổng diện tích toàn phần của hình khai triển đó là
\choice
{$47\text{ cm}^2$}
{\True $94\text{ cm}^2$}
{$60\text{ cm}^2$}
{$120\text{ cm}^2$}
\loigiai{
$S_{tp} = 2 \times (5 \times 4 + 5 \times 3 + 4 \times 3) = 2 \times (20 + 15 + 12) = 94\text{ (cm}^2\text{)}$.\\
\textbf{Đáp án của câu hỏi là B}
}
\end{ex}

\begin{ex} %[5H9V49]
Cần dùng một miếng bìa hình chữ nhật có kích thước tối thiểu là bao nhiêu để cuộn tròn làm thân một ống trụ có đường kính đáy $10\text{ cm}$ và chiều dài ống là $25\text{ cm}$ (không tính mép dán)?
\choice
{$10\text{ cm} \times 25\text{ cm}$}
{$15{,}7\text{ cm} \times 25\text{ cm}$}
{$62{,}8\text{ cm} \times 25\text{ cm}$}
{\True $31{,}4\text{ cm} \times 25\text{ cm}$}
\loigiai{
Chu vi đáy: $10 \times 3{,}14 = 31{,}4\text{ (cm)}$.\\
Kích thước miếng bìa là $31{,}4\text{ cm} \times 25\text{ cm}$.\\
\textbf{Đáp án của câu hỏi là D}
}
\end{ex}

\begin{ex} %[5H9V49]
Gấp một hình khai triển gồm $6$ hình vuông cạnh $a$ thành hình lập phương có thể tích $125\text{ cm}^3$. Độ dài cạnh $a$ là
\choice
{\True $5\text{ cm}$}
{$4\text{ cm}$}
{$6\text{ cm}$}
{$25\text{ cm}$}
\loigiai{
Vì $5 \times 5 \times 5 = 125$ nên $a = 5\text{ cm}$.\\
\textbf{Đáp án của câu hỏi là A}
}
\end{ex}

\begin{ex} %[5H9C49]
Trên các mặt của một con xúc xắc hình lập phương được đánh số từ $1$ đến $6$ sao cho tổng số chấm ở hai mặt đối diện luôn bằng $7$. Khi mở xúc xắc ra thành hình khai triển chữ thập ($4$ ô thẳng đứng, $2$ ô hai bên), nếu ô chính giữa là số $4$ thì ô đối diện với nó chứa số mấy?
\choice
{$1$}
{$2$}
{\True $3$}
{$5$}
\loigiai{
Mặt đối diện với số $4$ phải có số chấm là: $7 - 4 = 3$.\\
\textbf{Đáp án của câu hỏi là C}
}
\end{ex}

\textbf{Phần 2. Trắc nghiệm đúng sai}
\setcounterpageref{ex}{0}

\begin{ex} %[5H9H49]
Xét tính đúng sai về hình khai triển của các hình khối:
\choiceTF
{\True Hình khai triển của hình lập phương có đúng $6$ mặt hình vuông bằng nhau.}
{Hình khai triển của hình trụ gồm $1$ hình chữ nhật và $1$ hình tròn.}
{\True Khi cuộn mặt xung quanh của hình trụ, chiều dài hình chữ nhật bằng chu vi đáy.}
{\True Hình hộp chữ nhật có $3$ cặp mặt đối diện song song và bằng nhau.}
\loigiai{
\begin{itemize}
    \item Hình lập phương có 6 mặt vuông. Đúng.
    \item Hình trụ gồm 1 hình chữ nhật và 2 hình tròn đáy bằng nhau. Sai.
    \item Chiều dài thân trụ bằng chu vi đáy. Đúng.
    \item Hình hộp chữ nhật có 3 cặp mặt đối diện bằng nhau. Đúng.
\end{itemize}
\begin{itemize}
    \item a) \textbf{Đúng}.
    \item b) \textbf{Sai}.
    \item c) \textbf{Đúng}.
    \item d) \textbf{Đúng}.
\end{itemize}
}
\end{ex}

\begin{ex} %[5H9V49]
Một hình hộp chữ nhật có kích thước dài $8\text{ cm}$, rộng $5\text{ cm}$, cao $3\text{ cm}$. Khi trải phẳng các mặt thành hình khai triển, xét tính đúng sai:
\choiceTF
{\True Diện tích hai mặt đáy là $80\text{ cm}^2$.}
{\True Diện tích $4$ mặt bên xung quanh là $78\text{ cm}^2$.}
{Tổng diện tích toàn bộ hình khai triển là $160\text{ cm}^2$.}
{\True Tổng diện tích toàn bộ hình khai triển là $158\text{ cm}^2$.}
\loigiai{
\begin{itemize}
    \item Diện tích 2 mặt đáy: $2 \times 8 \times 5 = 80\text{ cm}^2$. Đúng.
    \item Diện tích 4 mặt xung quanh: $(8 + 5) \times 2 \times 3 = 78\text{ cm}^2$. Đúng.
    \item Tổng diện tích toàn bộ: $80 + 78 = 158\text{ cm}^2 \neq 160\text{ cm}^2$. Ý 3 Sai, Ý 4 Đúng.
\end{itemize}
\begin{itemize}
    \item a) \textbf{Đúng}.
    \item b) \textbf{Đúng}.
    \item c) \textbf{Sai}.
    \item d) \textbf{Đúng}.
\end{itemize}
}
\end{ex}

\textbf{Phần 3. Trả lời ngắn}
\setcounterpageref{ex}{0}

\begin{ex} %[5H9H49]
Hình lập phương có bao nhiêu đỉnh?
\shortans{8}
\loigiai{
\textbf{Đáp số:} $8$\
Hình lập phương có $8$ đỉnh.
}
\end{ex}

\begin{ex} %[5H9H49]
Hình khai triển của một hình lập phương gồm $6$ hình vuông, mỗi hình vuông có diện tích $25\text{ cm}^2$. Tính tổng diện tích của hình khai triển đó theo xăng-ti-mét vuông.
\shortans{150}
\loigiai{
\textbf{Đáp số:} $150$\\
$25 \times 6 = 150\text{ cm}^2$.
}
\end{ex}

\begin{ex} %[5H9V49]
Một lon sữa hình trụ có đường kính đáy $8\text{ cm}$, chiều cao $12\text{ cm}$. Nhãn giấy bọc xung quanh lon có dạng hình chữ nhật. Tính chiều dài của nhãn giấy đó theo xăng-ti-mét (lấy $\pi = 3{,}14$).
\shortans{25,12}
\loigiai{
\textbf{Đáp số:} $25{,}12$\\
Chiều dài nhãn = Chu vi đáy: $8 \times 3{,}14 = 25{,}12\text{ cm}$.
}
\end{ex}

\begin{ex} %[5H9V49]
Tính tổng diện tích các mặt của hình khai triển một hình hộp chữ nhật có chiều dài $6\text{ cm}$, chiều rộng $4\text{ cm}$, chiều cao $5\text{ cm}$ theo xăng-ti-mét vuông.
\shortans{148}
\loigiai{
\textbf{Đáp số:} $148$\\
$S_{tp} = 2 \times (6 \times 4 + 6 \times 5 + 4 \times 5) = 2 \times (24 + 30 + 20) = 148\text{ cm}^2$.
}
\end{ex}

\begin{ex} %[5H9V49]
Một hình khai triển của hình lập phương có chu vi ngoài là $56\text{ cm}$ (gồm $14$ đoạn cạnh của các hình vuông đơn vị). Tính cạnh của hình lập phương đó theo xăng-ti-mét.
\shortans{4}
\loigiai{
\textbf{Đáp số:} $4$\\
Cạnh hình lập phương: $56 : 14 = 4\text{ cm}$.
}
\end{ex}

\begin{ex} %[5H9C49]
Người ta cắt $4$ góc của một tấm bìa hình vuông cạnh $20\text{ cm}$, mỗi góc cắt đi một hình vuông nhỏ cạnh $3\text{ cm}$ rồi gập mép lên thành một chiếc hộp không nắp. Tính diện tích mặt ngoài của chiếc hộp đó theo xăng-ti-mét vuông.
\shortans{364}
\loigiai{
\textbf{Đáp số:} $364$\\
Diện tích tấm bìa ban đầu: $20 \times 20 = 400\text{ cm}^2$.\\
Diện tích $4$ góc cắt đi: $4 \times (3 \times 3) = 36\text{ cm}^2$.\\
Diện tích mặt ngoài của hộp: $400 - 36 = 364\text{ cm}^2$.
}
\end{ex}

\newpage

\TaoPhieuKiemTra{Chủ đề 9 -- Diện tích và thể tích của một số hình khối}{50}{Diện tích xung quanh và toàn phần hình hộp chữ nhật}{https://loptoanthaylinh.vn/thi-online/toan-5/bai-50}
\textbf{Phần 1. Trắc nghiệm lựa chọn}
\setcounterpageref{ex}{0}

\begin{ex} %[5H9N50]
Công thức tính diện tích xung quanh của hình hộp chữ nhật có chiều dài $a$, chiều rộng $b$ và chiều cao $c$ (cùng đơn vị đo) là
\choice
{\True $S_{xq} = (a + b) \times 2 \times c$}
{$S_{xq} = a \times b \times c$}
{$S_{xq} = (a + b) \times c$}
{$S_{xq} = a \times b + c$}
\loigiai{
Diện tích xung quanh bằng chu vi đáy nhân với chiều cao: $S_{xq} = (a + b) \times 2 \times c$.\\
\textbf{Đáp án của câu hỏi là A}
}
\end{ex}

\begin{ex} %[5H9N50]
Diện tích toàn phần của hình hộp chữ nhật bằng
\choice
{Diện tích xung quanh cộng diện tích $1$ mặt đáy}
{\True Diện tích xung quanh cộng diện tích $2$ mặt đáy}
{Diện tích $4$ mặt bên}
{Chu vi đáy nhân chiều cao cộng chiều rộng}
\loigiai{
$S_{tp} = S_{xq} + S_{2đáy}$.\\
\textbf{Đáp án của câu hỏi là B}
}
\end{ex}

\begin{ex} %[5H9H50]
Một hình hộp chữ nhật có chiều dài $5\text{ cm}$, chiều rộng $3\text{ cm}$ và chiều cao $4\text{ cm}$. Diện tích xung quanh của hình hộp chữ nhật đó là
\choice
{$32\text{ cm}^2$}
{$60\text{ cm}^2$}
{$94\text{ cm}^2$}
{\True $64\text{ cm}^2$}
\loigiai{
$S_{xq} = (5 + 3) \times 2 \times 4 = 64\text{ (cm}^2\text{)}$.\\
\textbf{Đáp án của câu hỏi là D}
}
\end{ex}

\begin{ex} %[5H9H50]
Một hình hộp chữ nhật có chiều dài $8\text{ dm}$, chiều rộng $6\text{ dm}$ và chiều cao $5\text{ dm}$. Diện tích toàn phần của hình hộp chữ nhật đó là
\choice
{$140\text{ dm}^2$}
{$240\text{ dm}^2$}
{\True $236\text{ dm}^2$}
{$188\text{ dm}^2$}
\loigiai{
$S_{xq} = (8 + 6) \times 2 \times 5 = 140\text{ (dm}^2\text{)}$.\\
$S_{2đáy} = 2 \times 8 \times 6 = 96\text{ (dm}^2\text{)}$.\\
$S_{tp} = 140 + 96 = 236\text{ (dm}^2\text{)}$.\\
\textbf{Đáp án của câu hỏi là C}
}
\end{ex}

\begin{ex} %[5H9H50]
Một cái thùng tôn hình hộp chữ nhật không nắp có chiều dài $1{,}2\text{ m}$, chiều rộng $0{,}8\text{ m}$, chiều cao $0{,}6\text{ m}$. Diện tích tôn làm thùng (không kể mép hàn) là
\choice
{\True $3{,}36\text{ m}^2$}
{$2{,}4\text{ m}^2$}
{$4{,}32\text{ m}^2$}
{$2{,}88\text{ m}^2$}
\loigiai{
$S_{xq} = (1{,}2 + 0{,}8) \times 2 \times 0{,}6 = 2{,}4\text{ (m}^2\text{)}$.\\
$S_{đáy} = 1{,}2 \times 0{,}8 = 0{,}96\text{ (m}^2\text{)}$.\\
Diện tích tôn: $2{,}4 + 0{,}96 = 3{,}36\text{ (m}^2\text{)}$.\\
\textbf{Đáp án của câu hỏi là A}
}
\end{ex}

\begin{ex} %[5H9H50]
Một hình hộp chữ nhật có chu vi đáy là $24\text{ cm}$, chiều cao là $6\text{ cm}$. Diện tích xung quanh của hình hộp chữ nhật đó là
\choice
{$72\text{ cm}^2$}
{\True $144\text{ cm}^2$}
{$288\text{ cm}^2$}
{$120\text{ cm}^2$}
\loigiai{
$S_{xq} = P_{đáy} \times h = 24 \times 6 = 144\text{ (cm}^2\text{)}$.\\
\textbf{Đáp án của câu hỏi là B}
}
\end{ex}

\begin{ex} %[5H9V50]
Một phòng học hình hộp chữ nhật có chiều dài $8\text{ m}$, chiều rộng $6\text{ m}$ và chiều cao $3{,}5\text{ m}$. Người ta quét vôi tường xung quanh và trần nhà bên trong phòng học. Biết tổng diện tích các cửa sổ và cửa ra vào là $12\text{ m}^2$. Diện tích cần quét vôi là
\choice
{\True $134\text{ m}^2$}
{$146\text{ m}^2$}
{$144\text{ m}^2$}
{$182\text{ m}^2$}
\loigiai{
$S_{xq} = (8 + 6) \times 2 \times 3{,}5 = 98\text{ (m}^2\text{)}$.\\
$S_{trần} = 8 \times 6 = 48\text{ (m}^2\text{)}$.\\
Diện tích quét vôi: $98 + 48 - 12 = 134\text{ (m}^2\text{)}$.\\
\textbf{Đáp án của câu hỏi là A}
}
\end{ex}

\begin{ex} %[5H9V50]
Một hình hộp chữ nhật có diện tích xung quanh là $240\text{ cm}^2$, chiều cao $8\text{ cm}$, chiều dài hơn chiều rộng $3\text{ cm}$. Diện tích toàn phần của hình hộp chữ nhật đó là
\choice
{$300\text{ cm}^2$}
{$384\text{ cm}^2$}
{\True $348\text{ cm}^2$}
{$420\text{ cm}^2$}
\loigiai{
Nửa chu vi đáy: $240 : 8 : 2 = 15\text{ (cm)}$.\\
Chiều dài: $(15 + 3) : 2 = 9\text{ (cm)}$.\\
Chiều rộng: $15 - 9 = 6\text{ (cm)}$.\\
Diện tích 2 đáy: $2 \times 9 \times 6 = 108\text{ (cm}^2\text{)}$.\\
$S_{tp} = 240 + 108 = 348\text{ (cm}^2\text{)}$.\\
\textbf{Đáp án của câu hỏi là C}
}
\end{ex}

\begin{ex} %[5H9V50]
Người ta sơn toàn bộ các mặt ngoài của một chiếc hòm gỗ dạng hình hộp chữ nhật có nắp đậy với chiều dài $1{,}5\text{ m}$, chiều rộng $0{,}8\text{ m}$, chiều cao $0{,}6\text{ m}$. Chi phí sơn là $40\,000$ đồng/$\text{m}^2$. Tổng chi phí sơn chiếc hòm là
\choice
{$180\,000$ đồng}
{$220\,000$ đồng}
{$200\,000$ đồng}
{\True $206\,400$ đồng}
\loigiai{
$S_{xq} = (1{,}5 + 0{,}8) \times 2 \times 0{,}6 = 2{,}76\text{ (m}^2\text{)}$.\\
$S_{2đáy} = 2 \times 1{,}5 \times 0{,}8 = 2{,}4\text{ (m}^2\text{)}$.\\
$S_{tp} = 2{,}76 + 2{,}4 = 5{,}16\text{ (m}^2\text{)}$.\\
Chi phí: $5{,}16 \times 40\,000 = 206\,400$ (đồng).\\
\textbf{Đáp án của câu hỏi là D}
}
\end{ex}

\begin{ex} %[5H9V50]
Một bể bơi hình hộp chữ nhật có chiều dài $25\text{ m}$, chiều rộng $10\text{ m}$ và sâu $2\text{ m}$. Người ta lát gạch men xung quanh và đáy bể. Diện tích cần lát gạch là
\choice
{$390\text{ m}^2$}
{$400\text{ m}^2$}
{\True $390\text{ m}^2$}
{$500\text{ m}^2$}
\loigiai{
$S_{xq} = (25 + 10) \times 2 \times 2 = 140\text{ (m}^2\text{)}$.\\
$S_{đáy} = 25 \times 10 = 250\text{ (m}^2\text{)}$.\\
Diện tích lát gạch: $140 + 250 = 390\text{ (m}^2\text{)}$.\\
\textbf{Đáp án của câu hỏi là C}
}
\end{ex}

\begin{ex} %[5H9V50]
Nếu tăng cả chiều dài, chiều rộng và chiều cao của hình hộp chữ nhật lên gấp $2$ lần thì diện tích toàn phần của nó tăng lên gấp mấy lần?
\choice
{$2$ lần}
{\True $4$ lần}
{$6$ lần}
{$8$ lần}
\loigiai{
Khi các kích thước tăng 2 lần thì diện tích mỗi mặt tăng $2 \times 2 = 4$ lần, do đó diện tích toàn phần tăng lên 4 lần.\\
\textbf{Đáp án của câu hỏi là B}
}
\end{ex}

\begin{ex} %[5H9C50]
Một hình hộp chữ nhật có chiều cao $6\text{ cm}$, diện tích đáy bằng $45\text{ cm}^2$ và diện tích xung quanh bằng $168\text{ cm}^2$. Chiều dài của hình hộp chữ nhật đó là
\choice
{$7\text{ cm}$}
{$5\text{ cm}$}
{$8\text{ cm}$}
{\True $9\text{ cm}$}
\loigiai{
Nửa chu vi đáy: $168 : 6 : 2 = 14\text{ (cm)}$.\\
Ta có: $a + b = 14$ và $a \times b = 45$.\\
Hai số có tổng bằng $14$ và tích bằng $45$ là $9$ và $5$.\\
Vậy chiều dài là $9\text{ cm}$.\\
\textbf{Đáp án của câu hỏi là D}
}
\end{ex}

\textbf{Phần 2. Trắc nghiệm đúng sai}
\setcounterpageref{ex}{0}

\begin{ex} %[5H9H50]
Cho hình hộp chữ nhật có chiều dài $12\text{ cm}$, chiều rộng $8\text{ cm}$ và chiều cao $5\text{ cm}$. Xét tính đúng sai:
\choiceTF
{\True Chu vi mặt đáy của hình hộp là $40\text{ cm}$.}
{\True Diện tích xung quanh của hình hộp là $200\text{ cm}^2$.}
{Diện tích hai mặt đáy của hình hộp là $96\text{ cm}^2$.}
{\True Diện tích toàn phần của hình hộp là $392\text{ cm}^2$.}
\loigiai{
\begin{itemize}
    \item Chu vi đáy: $(12 + 8) \times 2 = 40\text{ cm}$. Đúng.
    \item $S_{xq} = 40 \times 5 = 200\text{ cm}^2$. Đúng.
    \item Diện tích 2 mặt đáy: $2 \times 12 \times 8 = 192\text{ cm}^2 \neq 96\text{ cm}^2$. Sai.
    \item $S_{tp} = 200 + 192 = 392\text{ cm}^2$. Đúng.
\end{itemize}
\begin{itemize}
    \item a) \textbf{Đúng}.
    \item b) \textbf{Đúng}.
    \item c) \textbf{Sai}.
    \item d) \textbf{Đúng}.
\end{itemize}
}
\end{ex}

\begin{ex} %[5H9V50]
Một thùng sắt không nắp hình hộp chữ nhật có chiều dài $2\text{ m}$, chiều rộng $1{,}5\text{ m}$, chiều cao $1\text{ m}$. Người ta sơn cả mặt trong lẫn mặt ngoài của thùng. Xét tính đúng sai:
\choiceTF
{\True Diện tích xung quanh của chiếc thùng là $7\text{ m}^2$.}
{\True Diện tích một mặt đáy của chiếc thùng là $3\text{ m}^2$.}
{Diện tích quét sơn một mặt (gồm xung quanh và 1 đáy) là $13\text{ m}^2$.}
{\True Tổng diện tích quét sơn cả trong và ngoài là $20\text{ m}^2$.}
\loigiai{
\begin{itemize}
    \item $S_{xq} = (2 + 1{,}5) \times 2 \times 1 = 7\text{ m}^2$. Đúng.
    \item $S_{đáy} = 2 \times 1{,}5 = 3\text{ m}^2$. Đúng.
    \item Diện tích 1 mặt sơn: $7 + 3 = 10\text{ m}^2 \neq 13\text{ m}^2$. Sai.
    \item Tổng diện tích sơn cả trong và ngoài: $10 \times 2 = 20\text{ m}^2$. Đúng.
\end{itemize}
\begin{itemize}
    \item a) \textbf{Đúng}.
    \item b) \textbf{Đúng}.
    \item c) \textbf{Sai}.
    \item d) \textbf{Đúng}.
\end{itemize}
}
\end{ex}

\textbf{Phần 3. Trả lời ngắn}
\setcounterpageref{ex}{0}

\begin{ex} %[5H9H50]
Một hình hộp chữ nhật có chiều dài $9\text{ cm}$, chiều rộng $7\text{ cm}$ và chiều cao $4\text{ cm}$. Tính diện tích xung quanh theo xăng-ti-mét vuông.
\shortans{128}
\loigiai{
\textbf{Đáp số:} $128$\\
$(9 + 7) \times 2 \times 4 = 128\text{ cm}^2$.
}
\end{ex}

\begin{ex} %[5H9H50]
Tính diện tích toàn phần của hình hộp chữ nhật có chiều dài $10\text{ dm}$, chiều rộng $5\text{ dm}$, chiều cao $6\text{ dm}$ theo đề-xi-mét vuông.
\shortans{280}
\loigiai{
\textbf{Đáp số:} $280$\\
$S_{xq} = (10 + 5) \times 2 \times 6 = 180\text{ dm}^2$.\\
$S_{2đáy} = 2 \times 10 \times 5 = 100\text{ dm}^2$.\\
$S_{tp} = 180 + 100 = 280\text{ dm}^2$.
}
\end{ex}

\begin{ex} %[5H9V50]
Người ta làm một cái hộp bằng bìa không nắp có chiều dài $15\text{ cm}$, chiều rộng $10\text{ cm}$, chiều cao $8\text{ cm}$. Tính diện tích bìa cần dùng theo xăng-ti-mét vuông.
\shortans{550}
\loigiai{
\textbf{Đáp số:} $550$\\
$S_{xq} = (15 + 10) \times 2 \times 8 = 400\text{ cm}^2$.\\
$S_{đáy} = 15 \times 10 = 150\text{ cm}^2$.\\
Diện tích bìa: $400 + 150 = 550\text{ cm}^2$.
}
\end{ex}

\begin{ex} %[5H9V50]
Một phòng học dài $7\text{ m}$, rộng $5\text{ m}$, cao $3{,}8\text{ m}$. Người ta quét vôi tường và trần nhà bên trong phòng học. Tổng diện tích các cửa sổ và cửa ra vào là $10{,}5\text{ m}^2$. Tính diện tích cần quét vôi theo mét vuông.
\shortans{115,7}
\loigiai{
\textbf{Đáp số:} $115{,}7$\\
$S_{xq} = (7 + 5) \times 2 \times 3{,}8 = 91{,}2\text{ m}^2$.\\
$S_{trần} = 7 \times 5 = 35\text{ m}^2$.\\
Diện tích quét vôi: $91{,}2 + 35 - 10{,}5 = 115{,}7\text{ m}^2$.
}
\end{ex}

\begin{ex} %[5H9V50]
Một hình hộp chữ nhật có diện tích xung quanh là $150\text{ cm}^2$, chiều cao là $5\text{ cm}$. Tính chu vi mặt đáy của hình hộp chữ nhật đó theo xăng-ti-mét.
\shortans{30}
\loigiai{
\textbf{Đáp số:} $30$\\
Chu vi đáy: $150 : 5 = 30\text{ cm}$.
}
\end{ex}

\begin{ex} %[5H9C50]
Một cái bể hình hộp chữ nhật không nắp có chu vi đáy $16\text{ m}$, chiều dài hơn chiều rộng $2\text{ m}$, chiều cao $1{,}5\text{ m}$. Người ta sơn các mặt bên ngoài của bể với giá $35\,000$ đồng/$\text{m}^2$. Tính tổng số tiền sơn bể theo đơn vị nghìn đồng.
\shortans{840}
\loigiai{
\textbf{Đáp số:} $840$\\
Nửa chu vi đáy: $16 : 2 = 8\text{ m}$. Chiều dài: $5\text{ m}$, chiều rộng: $3\text{ m}$.\\
$S_{xq} = 16 \times 1{,}5 = 24\text{ m}^2$.\\
Tiền sơn ngoài xung quanh: $24 \times 35 = 840$ nghìn đồng.
}
\end{ex}

\newpage

\TaoPhieuKiemTra{Chủ đề 9 -- Diện tích và thể tích của một số hình khối}{51}{Diện tích xung quanh và toàn phần hình lập phương}{https://loptoanthaylinh.vn/thi-online/toan-5/bai-51}
\textbf{Phần 1. Trắc nghiệm lựa chọn}
\setcounterpageref{ex}{0}

\begin{ex} %[5H9N51]
Công thức tính diện tích xung quanh của hình lập phương có cạnh $a$ là
\choice
{$S_{xq} = a \times a \times 6$}
{\True $S_{xq} = a \times a \times 4$}
{$S_{xq} = a \times 4$}
{$S_{xq} = a \times a \times a$}
\loigiai{
Diện tích xung quanh của hình lập phương bằng diện tích một mặt nhân với $4$: $S_{xq} = a \times a \times 4$.\\
\textbf{Đáp án của câu hỏi là B}
}
\end{ex}

\begin{ex} %[5H9N51]
Công thức tính diện tích toàn phần của hình lập phương có cạnh $a$ là
\choice
{\True $S_{tp} = a \times a \times 6$}
{$S_{tp} = a \times a \times 4$}
{$S_{tp} = a \times 6$}
{$S_{tp} = a \times a \times a$}
\loigiai{
Diện tích toàn phần của hình lập phương bằng diện tích một mặt nhân với $6$: $S_{tp} = a \times a \times 6$.\\
\textbf{Đáp án của câu hỏi là A}
}
\end{ex}

\begin{ex} %[5H9H51]
Một hình lập phương có cạnh $5\text{ cm}$. Diện tích xung quanh của hình lập phương đó là
\choice
{$25\text{ cm}^2$}
{$150\text{ cm}^2$}
{$125\text{ cm}^2$}
{\True $100\text{ cm}^2$}
\loigiai{
$S_{xq} = 5 \times 5 \times 4 = 100\text{ (cm}^2\text{)}$.\\
\textbf{Đáp án của câu hỏi là D}
}
\end{ex}

\begin{ex} %[5H9H51]
Một hình lập phương có cạnh $4\text{ dm}$. Diện tích toàn phần của hình lập phương đó là
\choice
{$64\text{ dm}^2$}
{$16\text{ dm}^2$}
{\True $96\text{ dm}^2$}
{$48\text{ dm}^2$}
\loigiai{
$S_{tp} = 4 \times 4 \times 6 = 96\text{ (dm}^2\text{)}$.\\
\textbf{Đáp án của câu hỏi là C}
}
\end{ex}

\begin{ex} %[5H9H51]
Một cái hộp hình lập phương không có nắp có cạnh $6\text{ cm}$. Diện tích bìa để làm chiếc hộp đó là
\choice
{$144\text{ cm}^2$}
{$216\text{ cm}^2$}
{$36\text{ cm}^2$}
{\True $180\text{ cm}^2$}
\loigiai{
Hộp không nắp có 5 mặt: $6 \times 6 \times 5 = 180\text{ (cm}^2\text{)}$.\\
\textbf{Đáp án của câu hỏi là D}
}
\end{ex}

\begin{ex} %[5H9H51]
Một hình lập phương có diện tích một mặt là $36\text{ cm}^2$. Diện tích toàn phần của hình lập phương đó là
\choice
{$144\text{ cm}^2$}
{$180\text{ cm}^2$}
{\True $216\text{ cm}^2$}
{$240\text{ cm}^2$}
\loigiai{
$S_{tp} = 36 \times 6 = 216\text{ (cm}^2\text{)}$.\\
\textbf{Đáp án của câu hỏi là C}
}
\end{ex}

\begin{ex} %[5H9V51]
Một hình lập phương có diện tích xung quanh là $144\text{ dm}^2$. Cạnh của hình lập phương đó là
\choice
{$4\text{ dm}$}
{$8\text{ dm}$}
{$12\text{ dm}$}
{\True $6\text{ dm}$}
\loigiai{
Diện tích một mặt: $144 : 4 = 36\text{ (dm}^2\text{)}$.\\
Vì $6 \times 6 = 36$ nên cạnh là $6\text{ dm}$.\\
\textbf{Đáp án của câu hỏi là D}
}
\end{ex}

\begin{ex} %[5H9V51]
Người ta sơn toàn bộ mặt ngoài của một khối gỗ hình lập phương có cạnh $0{,}8\text{ m}$. Diện tích cần sơn là
\choice
{$2{,}56\text{ m}^2$}
{$1{,}92\text{ m}^2$}
{\True $3{,}84\text{ m}^2$}
{$0{,}512\text{ m}^2$}
\loigiai{
$S_{tp} = 0{,}8 \times 0{,}8 \times 6 = 3{,}84\text{ (m}^2\text{)}$.\\
\textbf{Đáp án của câu hỏi là C}
}
\end{ex}

\begin{ex} %[5H9V51]
Một hình lập phương có diện tích toàn phần là $294\text{ cm}^2$. Diện tích xung quanh của hình lập phương đó là
\choice
{\True $196\text{ cm}^2$}
{$49\text{ cm}^2$}
{$98\text{ cm}^2$}
{$147\text{ cm}^2$}
\loigiai{
Diện tích một mặt: $294 : 6 = 49\text{ (cm}^2\text{)}$.\\
$S_{xq} = 49 \times 4 = 196\text{ (cm}^2\text{)}$.\\
\textbf{Đáp án của câu hỏi là A}
}
\end{ex}

\begin{ex} %[5H9V51]
Nếu gấp cạnh của một hình lập phương lên $3$ lần thì diện tích xung quanh của nó tăng lên gấp mấy lần?
\choice
{$3$ lần}
{\True $9$ lần}
{$6$ lần}
{$27$ lần}
\loigiai{
Diện tích một mặt tăng $3 \times 3 = 9$ lần nên diện tích xung quanh tăng lên 9 lần.\\
\textbf{Đáp án của câu hỏi là B}
}
\end{ex}

\begin{ex} %[5H9V51]
Người ta dán giấy màu lên tất cả các mặt của một chiếc hộp quà hình lập phương cạnh $12\text{ cm}$. Biết giá tiền giấy màu là $500$ đồng/$100\text{ cm}^2$. Tổng số tiền giấy màu để bọc hộp quà là
\choice
{\True $4\,320$ đồng}
{$3\,600$ đồng}
{$4\,800$ đồng}
{$5\,200$ đồng}
\loigiai{
$S_{tp} = 12 \times 12 \times 6 = 864\text{ (cm}^2\text{)}$.\\
Tiền giấy: $864 : 100 \times 500 = 4\,320$ (đồng).\\
\textbf{Đáp án của câu hỏi là A}
}
\end{ex}

\begin{ex} %[5H9C51]
Hai hình lập phương có tỉ số độ dài cạnh là $\dfrac{2}{3}$. Tỉ số diện tích toàn phần của hình thứ nhất so với hình thứ hai là
\choice
{$\dfrac{2}{3}$}
{\True $\dfrac{4}{9}$}
{$\dfrac{8}{27}$}
{$\dfrac{4}{6}$}
\loigiai{
Tỉ số diện tích toàn phần bằng bình phương tỉ số cạnh: $\left(\dfrac{2}{3}\right)^2 = \dfrac{4}{9}$.\\
\textbf{Đáp án của câu hỏi là B}
}
\end{ex}

\textbf{Phần 2. Trắc nghiệm đúng sai}
\setcounterpageref{ex}{0}

\begin{ex} %[5H9H51]
Cho một hình lập phương có cạnh bằng $8\text{ cm}$. Xét tính đúng sai:
\choiceTF
{\True Diện tích một mặt của hình lập phương là $64\text{ cm}^2$.}
{\True Diện tích xung quanh của hình lập phương là $256\text{ cm}^2$.}
{Diện tích toàn phần của hình lập phương là $320\text{ cm}^2$.}
{\True Diện tích toàn phần của hình lập phương là $384\text{ cm}^2$.}
\loigiai{
\begin{itemize}
    \item $S_{1mặt} = 8 \times 8 = 64\text{ cm}^2$. Đúng.
    \item $S_{xq} = 64 \times 4 = 256\text{ cm}^2$. Đúng.
    \item $S_{tp} = 64 \times 6 = 384\text{ cm}^2 \neq 320\text{ cm}^2$. Ý 3 Sai, Ý 4 Đúng.
\end{itemize}
\begin{itemize}
    \item a) \textbf{Đúng}.
    \item b) \textbf{Đúng}.
    \item c) \textbf{Sai}.
    \item d) \textbf{Đúng}.
\end{itemize}
}
\end{ex}

\begin{ex} %[5H9V51]
Một cái thùng hình lập phương không nắp có diện tích toàn bộ các mặt ngoài là $180\text{ dm}^2$. Xét tính đúng sai:
\choiceTF
{\True Chiếc thùng có $5$ mặt ngoài.}
{\True Diện tích một mặt của chiếc thùng là $36\text{ dm}^2$.}
{Cạnh của chiếc thùng là $8\text{ dm}$.}
{\True Diện tích xung quanh của chiếc thùng là $144\text{ dm}^2$.}
\loigiai{
\begin{itemize}
    \item Thùng không nắp có 5 mặt. Đúng.
    \item Diện tích 1 mặt: $180 : 5 = 36\text{ dm}^2$. Đúng.
    \item Cạnh thùng: $6\text{ dm}$ vì $6 \times 6 = 36$ (không phải 8 dm). Sai.
    \item $S_{xq} = 36 \times 4 = 144\text{ dm}^2$. Đúng.
\end{itemize}
\begin{itemize}
    \item a) \textbf{Đúng}.
    \item b) \textbf{Đúng}.
    \item c) \textbf{Sai}.
    \item d) \textbf{Đúng}.
\end{itemize}
}
\end{ex}

\textbf{Phần 3. Trả lời ngắn}
\setcounterpageref{ex}{0}

\begin{ex} %[5H9H51]
Một hình lập phương có cạnh $7\text{ cm}$. Tính diện tích xung quanh theo xăng-ti-mét vuông.
\shortans{196}
\loigiai{
\textbf{Đáp số:} $196$\\
$7 \times 7 \times 4 = 196\text{ cm}^2$.
}
\end{ex}

\begin{ex} %[5H9H51]
Tính diện tích toàn phần của hình lập phương có cạnh $9\text{ dm}$ theo đề-xi-mét vuông.
\shortans{486}
\loigiai{
\textbf{Đáp số:} $486$\\
$9 \times 9 \times 6 = 486\text{ dm}^2$.
}
\end{ex}

\begin{ex} %[5H9V51]
Một hình lập phương có diện tích toàn phần là $150\text{ cm}^2$. Tính cạnh của hình lập phương đó theo xăng-ti-mét.
\shortans{5}
\loigiai{
\textbf{Đáp số:} $5$\\
Diện tích 1 mặt: $150 : 6 = 25\text{ cm}^2 \Rightarrow$ Cạnh là $5\text{ cm}$.
}
\end{ex}

\begin{ex} %[5H9V51]
Người ta làm một chiếc hộp hình lập phương không nắp bằng bìa các-tông có cạnh $15\text{ cm}$. Tính diện tích bìa cần dùng theo xăng-ti-mét vuông.
\shortans{1125}
\loigiai{
\textbf{Đáp số:} $1125$\\
$15 \times 15 \times 5 = 1\,125\text{ cm}^2$.
}
\end{ex}

\begin{ex} %[5H9V51]
Một hình lập phương có chu vi một mặt đáy là $28\text{ cm}$. Tính diện tích toàn phần của hình lập phương đó theo xăng-ti-mét vuông.
\shortans{294}
\loigiai{
\textbf{Đáp số:} $294$\\
Cạnh: $28 : 4 = 7\text{ cm}$.\\
$S_{tp} = 7 \times 7 \times 6 = 294\text{ cm}^2$.
}
\end{ex}

\begin{ex} %[5H9C51]
Ghép $8$ khối lập phương nhỏ cạnh $2\text{ cm}$ thành một hình lập phương lớn. Tính diện tích toàn phần của hình lập phương lớn đó theo xăng-ti-mét vuông.
\shortans{96}
\loigiai{
\textbf{Đáp số:} $96$\\
Cạnh hình lập phương lớn: $2 \times 2 = 4\text{ cm}$.\\
$S_{tp} = 4 \times 4 \times 6 = 96\text{ cm}^2$.
}
\end{ex}

\newpage

\TaoPhieuKiemTra{Chủ đề 9 -- Diện tích và thể tích của một số hình khối}{52}{Thể tích hình hộp chữ nhật}{https://loptoanthaylinh.vn/thi-online/toan-5/bai-52}
\textbf{Phần 1. Trắc nghiệm lựa chọn}
\setcounterpageref{ex}{0}

\begin{ex} %[5H9N52]
Muốn tính thể tích hình hộp chữ nhật có chiều dài $a$, chiều rộng $b$ và chiều cao $c$ (cùng đơn vị đo), ta dùng công thức
\choice
{\True $V = a \times b \times c$}
{$V = (a + b) \times c$}
{$V = (a + b) \times 2 \times c$}
{$V = a + b + c$}
\loigiai{
$V = a \times b \times c$.\\
\textbf{Đáp án của câu hỏi là A}
}
\end{ex}

\begin{ex} %[5H9N52]
Một hình hộp chữ nhật có chiều dài $7\text{ cm}$, chiều rộng $4\text{ cm}$, chiều cao $5\text{ cm}$. Thể tích của hình hộp đó là
\choice
{$68\text{ cm}^3$}
{\True $140\text{ cm}^3$}
{$140\text{ cm}^2$}
{$120\text{ cm}^3$}
\loigiai{
$V = 7 \times 4 \times 5 = 140\text{ (cm}^3\text{)}$.\\
\textbf{Đáp án của câu hỏi là B}
}
\end{ex}

\begin{ex} %[5H9H52]
Một bể nước hình hộp chữ nhật có diện tích đáy là $2{,}4\text{ m}^2$, chiều cao $1{,}5\text{ m}$. Thể tích của bể nước là
\choice
{\True $3{,}6\text{ m}^3$}
{$3{,}6\text{ m}^2$}
{$36\text{ m}^3$}
{$360\text{ l}$}
\loigiai{
$V = S_{đáy} \times h = 2{,}4 \times 1{,}5 = 3{,}6\text{ (m}^3\text{)}$.\\
\textbf{Đáp án của câu hỏi là A}
}
\end{ex}

\begin{ex} %[5H9H52]
Một hình hộp chữ nhật có thể tích $180\text{ cm}^3$, chiều dài $9\text{ cm}$, chiều rộng $4\text{ cm}$. Chiều cao của hình hộp chữ nhật là
\choice
{$4\text{ cm}$}
{$6\text{ cm}$}
{$8\text{ cm}$}
{\True $5\text{ cm}$}
\loigiai{
Chiều cao: $180 : (9 \times 4) = 180 : 36 = 5\text{ (cm)}$.\\
\textbf{Đáp án của câu hỏi là D}
}
\end{ex}

\begin{ex} %[5H9H52]
Một khối gỗ hình hộp chữ nhật có kích thước dài $1{,}2\text{ m}$, rộng $0{,}5\text{ m}$, cao $0{,}4\text{ m}$. Thể tích khối gỗ đó là
\choice
{\True $240\text{ dm}^3$}
{$0{,}24\text{ dm}^3$}
{$24\text{ dm}^3$}
{$2\,400\text{ dm}^3$}
\loigiai{
$V = 1{,}2 \times 0{,}5 \times 0{,}4 = 0{,}24\text{ m}^3 = 240\text{ dm}^3$.\\
\textbf{Đáp án của câu hỏi là A}
}
\end{ex}

\begin{ex} %[5H9H52]
Một bể cá hình hộp chữ nhật có chiều dài $60\text{ cm}$, chiều rộng $40\text{ cm}$ và chiều cao $50\text{ cm}$. Dung tích tối đa của bể cá là
\choice
{$120\text{ l}$}
{$12\text{ l}$}
{\True $120\text{ l}$}
{$1\,200\text{ l}$}
\loigiai{
$V = 6 \times 4 \times 5 = 120\text{ dm}^3 = 120\text{ l}$.\\
\textbf{Đáp án của câu hỏi là C}
}
\end{ex}

\begin{ex} %[5H9V52]
Một bể nước hình hộp chữ nhật có chiều dài $2\text{ m}$, chiều rộng $1{,}2\text{ m}$, chiều cao $1{,}5\text{ m}$. Bể đang chứa nước đến $60\%$ chiều cao bể. Lượng nước hiện có trong bể là
\choice
{$3{,}6\text{ m}^3$}
{$2\,400\text{ l}$}
{$1{,}8\text{ m}^3$}
{\True $2{,}16\text{ m}^3$}
\loigiai{
Thể tích bể: $2 \times 1{,}2 \times 1{,}5 = 3{,}6\text{ (m}^3\text{)}$.\\
Lượng nước hiện có: $3{,}6 \times 60\% = 2{,}16\text{ (m}^3\text{)}$.\\
\textbf{Đáp án của câu hỏi là D}
}
\end{ex}

\begin{ex} %[5H9V52]
Một cái hố hình hộp chữ nhật sâu $1{,}5\text{ m}$, đáy là hình vuông có chu vi $8\text{ m}$. Thể tích đất đã đào lên từ cái hố đó là
\choice
{$12\text{ m}^3$}
{\True $6\text{ m}^3$}
{$8\text{ m}^3$}
{$4\text{ m}^3$}
\loigiai{
Cạnh đáy: $8 : 4 = 2\text{ (m)}$.\\
Diện tích đáy: $2 \times 2 = 4\text{ (m}^2\text{)}$.\\
Thể tích đất: $4 \times 1{,}5 = 6\text{ (m}^3\text{)}$.\\
\textbf{Đáp án của câu hỏi là B}
}
\end{ex}

\begin{ex} %[5H9V52]
Người ta đổ vào một bể nước rỗng hình hộp chữ nhật $1\,800\text{ l}$ nước thì mực nước dâng cao $0{,}75\text{ m}$. Diện tích đáy của bể nước là
\choice
{$2\text{ m}^2$}
{$2{,}5\text{ m}^2$}
{\True $2{,}4\text{ m}^2$}
{$3\text{ m}^2$}
\loigiai{
$1\,800\text{ l} = 1{,}8\text{ m}^3$.\\
Diện tích đáy: $1{,}8 : 0{,}75 = 2{,}4\text{ (m}^2\text{)}$.\\
\textbf{Đáp án của câu hỏi là C}
}
\end{ex}

\begin{ex} %[5H9V52]
Nếu chiều dài tăng gấp $2$ lần, chiều rộng tăng gấp $3$ lần và chiều cao không đổi thì thể tích của hình hộp chữ nhật tăng lên gấp mấy lần?
\choice
{$5$ lần}
{\True $6$ lần}
{$8$ lần}
{$9$ lần}
\loigiai{
Thể tích mới: $V' = (2 \times a) \times (3 \times b) \times c = 6 \times (a \times b \times c) = 6V$.\\
\textbf{Đáp án của câu hỏi là B}
}
\end{ex}

\begin{ex} %[5H9V52]
Một bể kính nuôi cá dạng hình hộp chữ nhật có kích thước lòng trong: dài $80\text{ cm}$, rộng $50\text{ cm}$, cao $45\text{ cm}$. Mực nước ban đầu cao $30\text{ cm}$. Sau khi thả một hòn non bộ vào thì mực nước dâng lên cao $38\text{ cm}$ (hòn non bộ ngập hoàn toàn). Thể tích của hòn non bộ là
\choice
{$16\text{ dm}^3$}
{$24\text{ dm}^3$}
{$40\text{ dm}^3$}
{\True $32\text{ dm}^3$}
\loigiai{
Chiều cao nước dâng thêm: $38 - 30 = 8\text{ (cm)}$.\\
Thể tích hòn non bộ: $80 \times 50 \times 8 = 32\,000\text{ (cm}^3\text{)} = 32\text{ dm}^3$.\\
\textbf{Đáp án của câu hỏi là D}
}
\end{ex}

\begin{ex} %[5H9C52]
Một bể bơi hình hộp chữ nhật dài $20\text{ m}$, rộng $12\text{ m}$, sâu $1{,}6\text{ m}$. Người ta dùng $2$ máy bơm cùng bơm nước vào bể, mỗi phút máy 1 bơm được $120\text{ l}$, máy 2 bơm được $180\text{ l}$. Hỏi sau bao lâu cả hai máy bơm đầy được $75\%$ bể nước?
\choice
{$12\text{ giờ}$}
{$18\text{ giờ}$}
{\True $16\text{ giờ}$}
{$20\text{ giờ}$}
\loigiai{
Thể tích bể: $20 \times 12 \times 1{,}6 = 384\text{ m}^3 = 384\,000\text{ l}$.\\
$75\%$ thể tích: $384\,000 \times 75\% = 288\,000\text{ l}$.\\
Mỗi phút cả 2 máy bơm: $120 + 180 = 300\text{ (l)}$.\\
Thời gian bơm: $288\,000 : 300 = 960\text{ (phút)} = 16\text{ giờ}$.\\
\textbf{Đáp án của câu hỏi là C}
}
\end{ex}

\textbf{Phần 2. Trắc nghiệm đúng sai}
\setcounterpageref{ex}{0}

\begin{ex} %[5H9H52]
Một khối gỗ hình hộp chữ nhật có kích thước dài $1{,}5\text{ m}$, rộng $0{,}6\text{ m}$, cao $0{,}4\text{ m}$. Xét tính đúng sai:
\choiceTF
{\True Diện tích mặt đáy của khối gỗ là $0{,}9\text{ m}^2$.}
{\True Thể tích của khối gỗ là $0{,}36\text{ m}^3$.}
{Thể tích khối gỗ tương đương với $36\text{ dm}^3$.}
{\True Thể tích khối gỗ tương đương với $360\text{ dm}^3$.}
\loigiai{
\begin{itemize}
    \item $S_{đáy} = 1{,}5 \times 0{,}6 = 0{,}9\text{ m}^2$. Đúng.
    \item $V = 0{,}9 \times 0{,}4 = 0{,}36\text{ m}^3$. Đúng.
    \item $0{,}36\text{ m}^3 = 360\text{ dm}^3 \neq 36\text{ dm}^3$. Ý 3 Sai, Ý 4 Đúng.
\end{itemize}
\begin{itemize}
    \item a) \textbf{Đúng}.
    \item b) \textbf{Đúng}.
    \item c) \textbf{Sai}.
    \item d) \textbf{Đúng}.
\end{itemize}
}
\end{ex}

\begin{ex} %[5H9V52]
Một bể nước hình hộp chữ nhật có chiều dài $2\text{ m}$, chiều rộng $1{,}5\text{ m}$ và chiều cao $1{,}2\text{ m}$. Bể đang chứa $1\,800\text{ l}$ nước. Xét tính đúng sai:
\choiceTF
{\True Dung tích toàn bộ của bể là $3\,600\text{ l}$.}
{\True Mực nước hiện tại trong bể cao $0{,}6\text{ m}$.}
{Lượng nước hiện tại chiếm $75\%$ thể tích bể.}
{\True Cần đổ thêm $1\,800\text{ l}$ nước nữa thì đầy bể.}
\loigiai{
\begin{itemize}
    \item Dung tích: $2 \times 1{,}5 \times 1{,}2 = 3{,}6\text{ m}^3 = 3600\text{ l}$. Đúng.
    \item Mực nước: $1{,}8 : (2 \times 1{,}5) = 1{,}8 : 3 = 0{,}6\text{ m}$. Đúng.
    \item Tỉ lệ nước hiện có: $1800 : 3600 = 50\% \neq 75\%$. Sai.
    \item Cần thêm: $3600 - 1800 = 1800\text{ l}$. Đúng.
\end{itemize}
\begin{itemize}
    \item a) \textbf{Đúng}.
    \item b) \textbf{Đúng}.
    \item c) \textbf{Sai}.
    \item d) \textbf{Đúng}.
\end{itemize}
}
\end{ex}

\textbf{Phần 3. Trả lời ngắn}
\setcounterpageref{ex}{0}

\begin{ex} %[5H9H52]
Một hình hộp chữ nhật có chiều dài $12\text{ cm}$, chiều rộng $8\text{ cm}$, chiều cao $5\text{ cm}$. Tính thể tích theo xăng-ti-mét khối.
\shortans{480}
\loigiai{
\textbf{Đáp số:} $480$\\
$12 \times 8 \times 5 = 480\text{ cm}^3$.
}
\end{ex}

\begin{ex} %[5H9H52]
Một khối bê tông hình hộp chữ nhật có diện tích đáy là $3{,}5\text{ m}^2$, chiều cao $0{,}8\text{ m}$. Tính thể tích khối bê tông theo mét khối.
\shortans{2,8}
\loigiai{
\textbf{Đáp số:} $2{,}8$\\
$3{,}5 \times 0{,}8 = 2{,}8\text{ m}^3$.
}
\end{ex}

\begin{ex} %[5H9V52]
Một bể nước hình hộp chữ nhật dài $1{,}8\text{ m}$, rộng $1{,}5\text{ m}$, sâu $1{,}2\text{ m}$. Hỏi bể chứa được tối đa bao nhiêu lít nước?
\shortans{3240}
\loigiai{
\textbf{Đáp số:} $3240$\\
$1{,}8 \times 1{,}5 \times 1{,}2 = 3{,}24\text{ m}^3 = 3\,240\text{ l}$.
}
\end{ex}

\begin{ex} %[5H9V52]
Một hình hộp chữ nhật có thể tích $240\text{ cm}^3$, chiều dài $8\text{ cm}$, chiều cao $6\text{ cm}$. Tính chiều rộng của hình hộp chữ nhật theo xăng-ti-mét.
\shortans{5}
\loigiai{
\textbf{Đáp số:} $5$\\
$240 : (8 \times 6) = 5\text{ cm}$.
}
\end{ex}

\begin{ex} %[5H9V52]
Một bể cá có diện tích đáy là $1\,500\text{ cm}^2$. Người ta đổ vào bể $45\text{ l}$ nước. Tính chiều cao mực nước trong bể theo xăng-ti-mét.
\shortans{30}
\loigiai{
\textbf{Đáp số:} $30$\\
$45\text{ l} = 45\,000\text{ cm}^3 \Rightarrow 45\,000 : 1\,500 = 30\text{ cm}$.
}
\end{ex}

\begin{ex} %[5H9C52]
Một khối kim loại hình hộp chữ nhật có chiều dài $20\text{ cm}$, chiều rộng $15\text{ cm}$, chiều cao $10\text{ cm}$. Biết $1\text{ dm}^3$ kim loại đó nặng $7{,}8\text{ kg}$. Tính khối lượng khối kim loại theo ki-lô-gam.
\shortans{23,4}
\loigiai{
\textbf{Đáp số:} $23{,}4$\\
Thể tích: $20 \times 15 \times 10 = 3\,000\text{ cm}^3 = 3\text{ dm}^3$.\\
Khối lượng: $3 \times 7{,}8 = 23{,}4\text{ kg}$.
}
\end{ex}

\newpage

\TaoPhieuKiemTra{Chủ đề 9 -- Diện tích và thể tích của một số hình khối}{53}{Thể tích hình lập phương}{https://loptoanthaylinh.vn/thi-online/toan-5/bai-53}
\textbf{Phần 1. Trắc nghiệm lựa chọn}
\setcounterpageref{ex}{0}

\begin{ex} %[5H9N53]
Muốn tính thể tích hình lập phương có cạnh $a$, ta dùng công thức
\choice
{$V = a \times 4$}
{\True $V = a \times a \times a$}
{$V = a \times a \times 6$}
{$V = a \times a \times 4$}
\loigiai{
$V = a \times a \times a$.\\
\textbf{Đáp án của câu hỏi là B}
}
\end{ex}

\begin{ex} %[5H9N53]
Một hình lập phương có cạnh $4\text{ cm}$. Thể tích của hình lập phương đó là
\choice
{$16\text{ cm}^3$}
{$96\text{ cm}^3$}
{\True $64\text{ cm}^3$}
{$64\text{ cm}^2$}
\loigiai{
$V = 4 \times 4 \times 4 = 64\text{ (cm}^3\text{)}$.\\
\textbf{Đáp án của câu hỏi là C}
}
\end{ex}

\begin{ex} %[5H9H53]
Một khối gỗ hình lập phương có cạnh $0{,}5\text{ m}$. Thể tích khối gỗ đó là
\choice
{$1{,}25\text{ m}^3$}
{\True $0{,}125\text{ m}^3$}
{$125\text{ dm}^3$}
{$1{,}5\text{ m}^3$}
\loigiai{
$V = 0{,}5 \times 0{,}5 \times 0{,}5 = 0{,}125\text{ (m}^3\text{)}$.\\
\textbf{Đáp án của câu hỏi là B}
}
\end{ex}

\begin{ex} %[5H9H53]
Một hình lập phương có diện tích một mặt là $25\text{ cm}^2$. Thể tích của hình lập phương đó là
\choice
{$100\text{ cm}^3$}
{$150\text{ cm}^3$}
{$625\text{ cm}^3$}
{\True $125\text{ cm}^3$}
\loigiai{
Vì $S_{1mặt} = 25\text{ cm}^2$ nên cạnh $a = 5\text{ cm}$.\\
$V = 5 \times 5 \times 5 = 125\text{ (cm}^3\text{)}$.\\
\textbf{Đáp án của câu hỏi là D}
}
\end{ex}

\begin{ex} %[5H9H53]
Một hình lập phương có diện tích toàn phần là $54\text{ dm}^2$. Thể tích của hình lập phương đó là
\choice
{\True $27\text{ dm}^3$}
{$18\text{ dm}^3$}
{$36\text{ dm}^3$}
{$81\text{ dm}^3$}
\loigiai{
Diện tích một mặt: $54 : 6 = 9\text{ (dm}^2\text{)}$. Cạnh là $3\text{ dm}$.\\
$V = 3 \times 3 \times 3 = 27\text{ (dm}^3\text{)}$.\\
\textbf{Đáp án của câu hỏi là A}
}
\end{ex}

\begin{ex} %[5H9H53]
Một thùng chứa dạng hình lập phương có cạnh $1\text{ m}$. Thùng đó chứa được tối đa bao nhiêu lít nước?
\choice
{$100\text{ l}$}
{$10\,000\text{ l}$}
{\True $1\,000\text{ l}$}
{$100\,000\text{ l}$}
\loigiai{
$V = 1 \times 1 \times 1 = 1\text{ m}^3 = 1\,000\text{ dm}^3 = 1\,000\text{ l}$.\\
\textbf{Đáp án của câu hỏi là C}
}
\end{ex}

\begin{ex} %[5H9V53]
Nếu tăng cạnh của một hình lập phương lên gấp $2$ lần thì thể tích của nó tăng lên gấp mấy lần?
\choice
{$2$ lần}
{\True $8$ lần}
{$4$ lần}
{$6$ lần}
\loigiai{
Thể tích mới: $V' = (2a) \times (2a) \times (2a) = 8 \times a^3 = 8V$.\\
\textbf{Đáp án của câu hỏi là B}
}
\end{ex}

\begin{ex} %[5H9V53]
Một khối sắt hình lập phương có cạnh $0{,}2\text{ m}$. Biết $1\text{ dm}^3$ sắt nặng $7{,}8\text{ kg}$. Khối sắt đó nặng bao nhiêu ki-lô-gam?
\choice
{\True $62{,}4\text{ kg}$}
{$624\text{ kg}$}
{$6{,}24\text{ kg}$}
{$31{,}2\text{ kg}$}
\loigiai{
$V = 0{,}2 \times 0{,}2 \times 0{,}2 = 0{,}008\text{ m}^3 = 8\text{ dm}^3$.\\
Khối lượng: $8 \times 7{,}8 = 62{,}4\text{ (kg)}$.\\
\textbf{Đáp án của câu hỏi là A}
}
\end{ex}

\begin{ex} %[5H9V53]
Người ta xếp các khối lập phương nhỏ cạnh $1\text{ cm}$ thành một hình lập phương lớn có cạnh $6\text{ cm}$. Số khối lập phương nhỏ cần dùng là
\choice
{$36$ khối}
{$144$ khối}
{$108$ khối}
{\True $216$ khối}
\loigiai{
Số khối: $6 \times 6 \times 6 = 216$ khối.\\
\textbf{Đáp án của câu hỏi là D}
}
\end{ex}

\begin{ex} %[5H9V53]
Hai hình lập phương có tỉ số cạnh là $\dfrac{1}{3}$. Tỉ số thể tích của hình bé so với hình lớn là
\choice
{\True $\dfrac{1}{27}$}
{$\dfrac{1}{3}$}
{$\dfrac{1}{9}$}
{$\dfrac{1}{6}$}
\loigiai{
Tỉ số thể tích: $\left(\dfrac{1}{3}\right)^3 = \dfrac{1}{27}$.\\
\textbf{Đáp án của câu hỏi là A}
}
\end{ex}

\begin{ex} %[5H9V53]
Một hình lập phương có diện tích xung quanh là $64\text{ cm}^2$. Thể tích của nó là
\choice
{$32\text{ cm}^3$}
{$96\text{ cm}^3$}
{$128\text{ cm}^3$}
{\True $64\text{ cm}^3$}
\loigiai{
Diện tích một mặt: $64 : 4 = 16\text{ (cm}^2\text{)}$. Cạnh là $4\text{ cm}$.\\
$V = 4 \times 4 \times 4 = 64\text{ (cm}^3\text{)}$.\\
\textbf{Đáp án của câu hỏi là D}
}
\end{ex}

\begin{ex} %[5H9C53]
Một khối gỗ hình lập phương có cạnh $9\text{ cm}$. Người ta khoét một lỗ rỗng cũng hình lập phương có cạnh $3\text{ cm}$ ở chính giữa khối gỗ. Thể tích phần gỗ còn lại là
\choice
{$720\text{ cm}^3$}
{$726\text{ cm}^3$}
{\True $702\text{ cm}^3$}
{$648\text{ cm}^3$}
\loigiai{
Thể tích ban đầu: $9 \times 9 \times 9 = 729\text{ (cm}^3\text{)}$.\\
Thể tích phần khoét: $3 \times 3 \times 3 = 27\text{ (cm}^3\text{)}$.\\
Thể tích còn lại: $729 - 27 = 702\text{ (cm}^3\text{)}$.\\
\textbf{Đáp án của câu hỏi là C}
}
\end{ex}

\textbf{Phần 2. Trắc nghiệm đúng sai}
\setcounterpageref{ex}{0}

\begin{ex} %[5H9H53]
Cho hình lập phương có cạnh $6\text{ dm}$. Xét tính đúng sai:
\choiceTF
{\True Diện tích toàn phần của hình lập phương là $216\text{ dm}^2$.}
{\True Thể tích của hình lập phương là $216\text{ dm}^3$.}
{Số đo diện tích toàn phần (theo $\text{dm}^2$) và thể tích (theo $\text{dm}^3$) của hình này luôn khác nhau với mọi hình lập phương.}
{\True Với hình lập phương cạnh $6\text{ dm}$, số đo diện tích toàn phần và số đo thể tích có cùng giá trị số là $216$.}
\loigiai{
\begin{itemize}
    \item $S_{tp} = 6 \times 6 \times 6 = 216\text{ dm}^2$. Đúng.
    \item $V = 6 \times 6 \times 6 = 216\text{ dm}^3$. Đúng.
    \item Cạnh bằng 6 thì giá trị số trùng nhau nên khẳng định luôn khác nhau là Sai.
    \item Đúng, cùng giá trị 216.
\end{itemize}
\begin{itemize}
    \item a) \textbf{Đúng}.
    \item b) \textbf{Đúng}.
    \item c) \textbf{Sai}.
    \item d) \textbf{Đúng}.
\end{itemize}
}
\end{ex}

\begin{ex} %[5H9V53]
Một bể nước hình lập phương có cạnh trong lòng bể là $1{,}2\text{ m}$. Bể đang chứa $864\text{ l}$ nước. Xét tính đúng sai:
\choiceTF
{\True Dung tích toàn bộ của bể là $1\,728\text{ l}$.}
{\True Lượng nước hiện có trong bể chiếm $50\%$ thể tích bể.}
{Mực nước trong bể hiện cao $0{,}8\text{ m}$.}
{\True Mực nước trong bể hiện cao $0{,}6\text{ m}$.}
\loigiai{
\begin{itemize}
    \item Dung tích: $1{,}2 \times 1{,}2 \times 1{,}2 = 1{,}728\text{ m}^3 = 1728\text{ l}$. Đúng.
    \item Tỉ lệ nước: $864 : 1728 = 50\%$. Đúng.
    \item Chiều cao nước: $1{,}2 \times 50\% = 0{,}6\text{ m} \neq 0{,}8\text{ m}$. Ý 3 Sai, Ý 4 Đúng.
\end{itemize}
\begin{itemize}
    \item a) \textbf{Đúng}.
    \item b) \textbf{Đúng}.
    \item c) \textbf{Sai}.
    \item d) \textbf{Đúng}.
\end{itemize}
}
\end{ex}

\textbf{Phần 3. Trả lời ngắn}
\setcounterpageref{ex}{0}

\begin{ex} %[5H9H53]
Một hình lập phương có cạnh $8\text{ cm}$. Tính thể tích theo xăng-ti-mét khối.
\shortans{512}
\loigiai{
\textbf{Đáp số:} $512$\\
$8 \times 8 \times 8 = 512\text{ cm}^3$.
}
\end{ex}

\begin{ex} %[5H9H53]
Một hình lập phương có diện tích một mặt là $49\text{ cm}^2$. Tính thể tích hình lập phương đó theo xăng-ti-mét khối.
\shortans{343}
\loigiai{
\textbf{Đáp số:} $343$\\
Cạnh là $7\text{ cm} \Rightarrow V = 7 \times 7 \times 7 = 343\text{ cm}^3$.
}
\end{ex}

\begin{ex} %[5H9V53]
Một thùng chứa hình lập phương có cạnh $1{,}5\text{ m}$. Tính số lít nước tối đa thùng chứa được.
\shortans{3375}
\loigiai{
\textbf{Đáp số:} $3375$\\
$1{,}5 \times 1{,}5 \times 1{,}5 = 3{,}375\text{ m}^3 = 3\,375\text{ l}$.
}
\end{ex}

\begin{ex} %[5H9V53]
Một khối gỗ hình lập phương có thể tích $1\,000\text{ cm}^3$. Tính diện tích toàn phần của khối gỗ đó theo xăng-ti-mét vuông.
\shortans{600}
\loigiai{
\textbf{Đáp số:} $600$\\
Cạnh: $10\text{ cm} \Rightarrow S_{tp} = 10 \times 10 \times 6 = 600\text{ cm}^2$.
}
\end{ex}

\begin{ex} %[5H9V53]
Một khối nhôm hình lập phương cạnh $4\text{ dm}$. Biết $1\text{ dm}^3$ nhôm nặng $2{,}7\text{ kg}$. Tính khối lượng khối nhôm theo ki-lô-gam.
\shortans{172,8}
\loigiai{
\textbf{Đáp số:} $172{,}8$\\
Thể tích: $4 \times 4 \times 4 = 64\text{ dm}^3$.\\
Khối lượng: $64 \times 2{,}7 = 172{,}8\text{ kg}$.
}
\end{ex}

\begin{ex} %[5H9C53]
Có một khối lập phương bằng kim loại cạnh $12\text{ cm}$. Người ta đem nung chảy rồi đúc thành các khối lập phương nhỏ cạnh $3\text{ cm}$. Hỏi đúc được bao nhiêu khối lập phương nhỏ như vậy?
\shortans{64}
\loigiai{
\textbf{Đáp số:} $64$\\
$(12 \times 12 \times 12) : (3 \times 3 \times 3) = 1\,728 : 27 = 64$ khối.
}
\end{ex}

\newpage

\TaoPhieuKiemTra{Chủ đề 9 -- Diện tích và thể tích của một số hình khối}{54}{Thực hành tính toán, ước lượng thể tích}{https://loptoanthaylinh.vn/thi-online/toan-5/bai-54}
\textbf{Phần 1. Trắc nghiệm lựa chọn}
\setcounterpageref{ex}{0}

\begin{ex} %[5H9N54]
Đơn vị thích hợp nhất để đo thể tích một phòng học là
\choice
{\True Mét khối}
{Xăng-ti-mét khối}
{Đề-xi-mét khối}
{Mi-li-lít}
\loigiai{
Phòng học có thể tích lớn nên đơn vị thích hợp nhất là mét khối ($	ext{m}^3$).\\
\textbf{Đáp án của câu hỏi là A}
}
\end{ex}

\begin{ex} %[5H9N54]
Ước lượng thể tích của một bao diêm thông thường khoảng
\choice
{$1\text{ m}^3$}
{$1\text{ dm}^3$}
{\True $15\text{ cm}^3$}
{$150\text{ dm}^3$}
\loigiai{
Bao diêm kích thước khoảng $5\text{ cm} \times 3\text{ cm} \times 1\text{ cm} = 15\text{ cm}^3$.\\
\textbf{Đáp án của câu hỏi là C}
}
\end{ex}

\begin{ex} %[5H9H54]
Một phòng học có kích thước: dài $8\text{ m}$, rộng $6\text{ m}$, cao $3{,}5\text{ m}$. Biết lớp học có $35$ học sinh và $1$ giáo viên. Thể tích không khí trung bình cho mỗi người trong phòng là
\choice
{$3{,}5\text{ m}^3$}
{$5\text{ m}^3$}
{$4\text{ m}^3$}
{\True $4{,}67\text{ m}^3$}
\loigiai{
Thể tích phòng: $8 \times 6 \times 3{,}5 = 168\text{ (m}^3\text{)}$.\\
Tổng số người: $35 + 1 = 36$ (người).\\
Thể tích trung bình: $168 : 36 \approx 4{,}67\text{ (m}^3\text{/người)}$.\\
\textbf{Đáp án của câu hỏi là D}
}
\end{ex}

\begin{ex} %[5H9H54]
Một viên gạch xây dựng dạng hình hộp chữ nhật có kích thước dài $20\text{ cm}$, rộng $10\text{ cm}$, cao $5\text{ cm}$. Thể tích của $1\,000$ viên gạch như thế là
\choice
{$10\text{ m}^3$}
{\True $1\text{ m}^3$}
{$0{,}1\text{ m}^3$}
{$100\text{ m}^3$}
\loigiai{
Thể tích 1 viên gạch: $20 \times 10 \times 5 = 1\,000\text{ (cm}^3\text{)}$.\\
Thể tích 1000 viên: $1\,000 \times 1\,000 = 1\,000\,000\text{ cm}^3 = 1\text{ m}^3$.\\
\textbf{Đáp án của câu hỏi là B}
}
\end{ex}

\begin{ex} %[5H9H54]
Một chiếc xe tải có thùng hàng kích thước: dài $4\text{ m}$, rộng $2\text{ m}$, cao $1{,}5\text{ m}$. Thể tích tối đa của thùng hàng là
\choice
{$8\text{ m}^3$}
{$16\text{ m}^3$}
{\True $12\text{ m}^3$}
{$10\text{ m}^3$}
\loigiai{
$V = 4 \times 2 \times 1{,}5 = 12\text{ (m}^3\text{)}$.\\
\textbf{Đáp án của câu hỏi là C}
}
\end{ex}

\begin{ex} %[5H9H54]
Ước lượng thể tích của một tủ lạnh gia đình loại trung bình khoảng
\choice
{\True $300\text{ l}$}
{$30\text{ l}$}
{$3\,000\text{ l}$}
{$3\text{ l}$}
\loigiai{
Tủ lạnh gia đình thông dụng có dung tích khoảng $250 - 350\text{ l}$ ($300\text{ l}$).\\
\textbf{Đáp án của câu hỏi là A}
}
\end{ex}

\begin{ex} %[5H9V54]
Người ta xếp các hộp bánh hình lập phương cạnh $10\text{ cm}$ vào một thùng các-tông hình hộp chữ nhật có kích thước dài $60\text{ cm}$, rộng $40\text{ cm}$, cao $30\text{ cm}$. Số hộp bánh xếp được nhiều nhất vào thùng là
\choice
{$48$ hộp}
{$64$ hộp}
{$80$ hộp}
{\True $72$ hộp}
\loigiai{
Số hộp theo chiều dài: $60 : 10 = 6$ hộp.\\
Số hộp theo chiều rộng: $40 : 10 = 4$ hộp.\\
Số hộp theo chiều cao: $30 : 10 = 3$ hộp.\\
Tổng số hộp: $6 \times 4 \times 3 = 72$ hộp.\\
\textbf{Đáp án của câu hỏi là D}
}
\end{ex}

\begin{ex} %[5H9V54]
Một bể cát cho trẻ em chơi có dạng hình hộp chữ nhật dài $3\text{ m}$, rộng $2\text{ m}$. Người ta đổ cát vào bể dày $20\text{ cm}$. Biết giá tiền $1\text{ m}^3$ cát là $250\,000$ đồng. Số tiền mua cát là
\choice
{$200\,000$ đồng}
{$250\,000$ đồng}
{\True $300\,000$ đồng}
{$350\,000$ đồng}
\loigiai{
Đổi $20\text{ cm} = 0{,}2\text{ m}$.\\
Thể tích cát: $3 \times 2 \times 0{,}2 = 1{,}2\text{ (m}^3\text{)}$.\\
Tiền mua cát: $1{,}2 \times 250\,000 = 300\,000$ (đồng).\\
\textbf{Đáp án của câu hỏi là C}
}
\end{ex}

\begin{ex} %[5H9V54]
Một bình giữ nhiệt hình trụ có diện tích đáy $50\text{ cm}^2$ và chiều cao phần chứa nước là $18\text{ cm}$. Dung tích chứa nước của bình xấp xỉ
\choice
{$0{,}5\text{ l}$}
{\True $0{,}9\text{ l}$}
{$1{,}2\text{ l}$}
{$1{,}5\text{ l}$}
\loigiai{
Thể tích: $50 \times 18 = 900\text{ cm}^3 = 0{,}9\text{ dm}^3 = 0{,}9\text{ l}$.\\
\textbf{Đáp án của câu hỏi là B}
}
\end{ex}

\begin{ex} %[5H9V54]
Một bồn nước inox hình trụ có bán kính đáy $0{,}5\text{ m}$, chiều cao $1{,}4\text{ m}$. Dung tích bồn nước (lấy $\pi = 3{,}14$) xấp xỉ là
\choice
{$1\,000\text{ l}$}
{$1\,250\text{ l}$}
{$1\,500\text{ l}$}
{\True $1\,100\text{ l}$}
\loigiai{
$V = 3{,}14 \times 0{,}5 \times 0{,}5 \times 1{,}4 = 1{,}099\text{ m}^3 \approx 1\,100\text{ l}$.\\
\textbf{Đáp án của câu hỏi là D}
}
\end{ex}

\begin{ex} %[5H9V54]
Để đắp một con đê bằng đất hình lăng trụ có mặt cắt ngang diện tích $4{,}5\text{ m}^2$ và chiều dài đê $100\text{ m}$. Cần dùng bao nhiêu chuyến xe tải chở đất, biết mỗi chuyến xe chở được $6\text{ m}^3$ đất?
\choice
{\True $75$ chuyến}
{$60$ chuyến}
{$80$ chuyến}
{$90$ chuyến}
\loigiai{
Thể tích đất cần đắp: $4{,}5 \times 100 = 450\text{ (m}^3\text{)}$.\\
Số chuyến xe: $450 : 6 = 75$ (chuyến).\\
\textbf{Đáp án của câu hỏi là A}
}
\end{ex}

\begin{ex} %[5H9C54]
Một khối bê tông gồm một phần đế hình hộp chữ nhật ($2\text{ m} \times 1{,}5\text{ m} \times 0{,}4\text{ m}$) và phần trên là khối lập phương cạnh $0{,}8\text{ m}$. Thể tích toàn bộ khối bê tông đó là
\choice
{$1{,}2\text{ m}^3$}
{\True $1{,}712\text{ m}^3$}
{$1{,}512\text{ m}^3$}
{$2{,}012\text{ m}^3$}
\loigiai{
Thể tích đế: $2 \times 1{,}5 \times 0{,}4 = 1{,}2\text{ (m}^3\text{)}$.\\
Thể tích phần trên: $0{,}8 \times 0{,}8 \times 0{,}8 = 0{,}512\text{ (m}^3\text{)}$.\\
Tổng thể tích: $1{,}2 + 0{,}512 = 1{,}712\text{ (m}^3\text{)}$.\\
\textbf{Đáp án của câu hỏi là B}
}
\end{ex}

\textbf{Phần 2. Trắc nghiệm đúng sai}
\setcounterpageref{ex}{0}

\begin{ex} %[5H9H54]
Xét tính đúng sai về ước lượng và tính toán thể tích các vật thể trong đời sống:
\choiceTF
{\True Một hộp sữa tươi $180\text{ ml}$ có thể tích bằng $180\text{ cm}^3$.}
{Một bể nước sinh hoạt gia đình thường có thể tích khoảng $2\text{ cm}^3$.}
{\True Một bể nước ngầm $3\text{ m}^3$ chứa được $3\,000\text{ l}$ nước.}
{\True Một phòng ngủ có thể tích $45\text{ m}^3$ chứa khoảng $45\,000\text{ l}$ không khí.}
\loigiai{
\begin{itemize}
    \item $1\text{ ml} = 1\text{ cm}^3 \Rightarrow 180\text{ ml} = 180\text{ cm}^3$. Đúng.
    \item Bể nước gia đình khoảng $2\text{ m}^3$ chứ không phải $2\text{ cm}^3$. Sai.
    \item $3\text{ m}^3 = 3000\text{ dm}^3 = 3000\text{ l}$. Đúng.
    \item $45\text{ m}^3 = 45000\text{ l}$. Đúng.
\end{itemize}
\begin{itemize}
    \item a) \textbf{Đúng}.
    \item b) \textbf{Sai}.
    \item c) \textbf{Đúng}.
    \item d) \textbf{Đúng}.
\end{itemize}
}
\end{ex}

\begin{ex} %[5H9V54]
Một kho chứa hàng hình hộp chữ nhật có kích thước dài $12\text{ m}$, rộng $8\text{ m}$, cao $4\text{ m}$. Người ta xếp các thùng hàng hình lập phương cạnh $1\text{ m}$ vào kho. Xét tính đúng sai:
\choiceTF
{\True Thể tích của kho hàng là $384\text{ m}^3$.}
{\True Kho có thể chứa tối đa $384$ thùng hàng nếu xếp kín.}
{Nếu chỉ xếp hàng chiếm $75\%$ thể tích kho thì có $300$ thùng hàng.}
{\True Nếu chỉ xếp hàng chiếm $75\%$ thể tích kho thì có $288$ thùng hàng.}
\loigiai{
\begin{itemize}
    \item Thể tích kho: $12 \times 8 \times 4 = 384\text{ m}^3$. Đúng.
    \item Thùng hàng cạnh 1m có thể tích $1\text{ m}^3$ nên chứa tối đa 384 thùng. Đúng.
    \item $75\%$ số thùng: $384 \times 75\% = 288$ thùng $\neq 300$ thùng. Ý 3 Sai, Ý 4 Đúng.
\end{itemize}
\begin{itemize}
    \item a) \textbf{Đúng}.
    \item b) \textbf{Đúng}.
    \item c) \textbf{Sai}.
    \item d) \textbf{Đúng}.
\end{itemize}
}
\end{ex}

\textbf{Phần 3. Trả lời ngắn}
\setcounterpageref{ex}{0}

\begin{ex} %[5H9H54]
Một bao cát xây dựng có thể tích $0{,}04\text{ m}^3$. Hỏi $50$ bao cát như vậy có tổng thể tích bao nhiêu mét khối?
\shortans{2}
\loigiai{
\textbf{Đáp số:} $2$\\
$0{,}04 \times 50 = 2\text{ m}^3$.
}
\end{ex}

\begin{ex} %[5H9H54]
Một bể phao bơi mini hình chữ nhật có kích thước dài $2\text{ m}$, rộng $1{,}5\text{ m}$, sâu $0{,}4\text{ m}$. Tính thể tích nước khi bơm đầy bể theo lít.
\shortans{1200}
\loigiai{
\textbf{Đáp số:} $1200$\\
$2 \times 1{,}5 \times 0{,}4 = 1{,}2\text{ m}^3 = 1\,200\text{ l}$.
}
\end{ex}

\begin{ex} %[5H9V54]
Một lớp học dài $9\text{ m}$, rộng $7\text{ m}$, cao $3{,}6\text{ m}$. Tiêu chuẩn không khí cho mỗi học sinh là $5\text{ m}^3$. Hỏi phòng học này có thể bố trí tối đa cho bao nhiêu học sinh?
\shortans{45}
\loigiai{
\textbf{Đáp số:} $45$\\
Thể tích phòng: $9 \times 7 \times 3{,}6 = 226{,}8\text{ m}^3$.\\
Số học sinh tối đa: $226{,}8 : 5 = 45{,}36 \Rightarrow 45$ học sinh.
}
\end{ex}

\begin{ex} %[5H9V54]
Người ta xếp các hộp sữa tươi hình hộp chữ nhật kích thước $5\text{ cm} \times 4\text{ cm} \times 10\text{ cm}$ vào một thùng các-tông kích thước $40\text{ cm} \times 30\text{ cm} \times 20\text{ cm}$. Hỏi xếp được nhiều nhất bao nhiêu hộp sữa?
\shortans{120}
\loigiai{
\textbf{Đáp số:} $120$\\
Số hộp: $(40 : 5) \times (30 : 4) \times (20 : 10) = 8 \times 7 \times 2 = 112$ hoặc theo thể tích: $24\,000 : 200 = 120$ hộp.
}
\end{ex}

\begin{ex} %[5H9V54]
Một hố trồng cây hình lập phương cạnh $60\text{ cm}$. Tính thể tích đất cần đổ đầy hố theo mét khối (dùng dấu phẩy nếu có).
\shortans{0,216}
\loigiai{
\textbf{Đáp số:} $0{,}216$\\
$0{,}6 \times 0{,}6 \times 0{,}6 = 0{,}216\text{ m}^3$.
}
\end{ex}

\begin{ex} %[5H9C54]
Một xe chở xăng có bồn chứa hình trụ dung tích $18\text{ m}^3$. Người ta bơm xăng từ bồn vào các thùng phuy hình trụ dung tích $200\text{ l}$. Hỏi bơm đầy được bao nhiêu thùng phuy?
\shortans{90}
\loigiai{
\textbf{Đáp số:} $90$\\
$18\text{ m}^3 = 18\,000\text{ l} \Rightarrow 18\,000 : 200 = 90$ thùng.
}
\end{ex}

\newpage

\TaoPhieuKiemTra{Chủ đề 9 -- Diện tích và thể tích của một số hình khối}{55}{Luyện tập chung hình khối}{https://loptoanthaylinh.vn/thi-online/toan-5/bai-55}
\textbf{Phần 1. Trắc nghiệm lựa chọn}
\setcounterpageref{ex}{0}

\begin{ex} %[5H9N55]
Hình lập phương có cạnh $a$ thì diện tích xung quanh gấp mấy lần diện tích một mặt?
\choice
{\True $4$ lần}
{$2$ lần}
{$6$ lần}
{$8$ lần}
\loigiai{
$S_{xq} = a \times a \times 4$ nên gấp $4$ lần diện tích một mặt.\\
\textbf{Đáp án của câu hỏi là A}
}
\end{ex}

\begin{ex} %[5H9N55]
Một hình hộp chữ nhật có diện tích xung quanh là $120\text{ cm}^2$ và diện tích hai đáy là $80\text{ cm}^2$. Diện tích toàn phần của hình hộp là
\choice
{$160\text{ cm}^2$}
{$240\text{ cm}^2$}
{\True $200\text{ cm}^2$}
{$280\text{ cm}^2$}
\loigiai{
$S_{tp} = S_{xq} + S_{2đáy} = 120 + 80 = 200\text{ (cm}^2\text{)}$.\\
\textbf{Đáp án của câu hỏi là C}
}
\end{ex}

\begin{ex} %[5H9H55]
Một hình lập phương có diện tích toàn phần là $150\text{ dm}^2$. Thể tích của nó là
\choice
{$100\text{ dm}^3$}
{$150\text{ dm}^3$}
{$250\text{ dm}^3$}
{\True $125\text{ dm}^3$}
\loigiai{
Diện tích 1 mặt: $150 : 6 = 25\text{ (dm}^2\text{)}$. Cạnh là $5\text{ dm}$.\\
$V = 5 \times 5 \times 5 = 125\text{ (dm}^3\text{)}$.\\
\textbf{Đáp án của câu hỏi là D}
}
\end{ex}

\begin{ex} %[5H9H55]
Một bể nước hình hộp chữ nhật có chiều dài $2\text{ m}$, chiều rộng $1{,}5\text{ m}$, chiều cao $1\text{ m}$. Bể chứa được tối đa số lít nước là
\choice
{$300\text{ l}$}
{\True $3\,000\text{ l}$}
{$30\,000\text{ l}$}
{$30\text{ l}$}
\loigiai{
$V = 2 \times 1{,}5 \times 1 = 3\text{ m}^3 = 3\,000\text{ l}$.\\
\textbf{Đáp án của câu hỏi là B}
}
\end{ex}

\begin{ex} %[5H9H55]
Diện tích xung quanh của một hình lập phương là $64\text{ cm}^2$. Diện tích toàn phần của nó là
\choice
{$80\text{ cm}^2$}
{$120\text{ cm}^2$}
{$144\text{ cm}^2$}
{\True $96\text{ cm}^2$}
\loigiai{
Diện tích 1 mặt: $64 : 4 = 16\text{ (cm}^2\text{)}$.\\
$S_{tp} = 16 \times 6 = 96\text{ (cm}^2\text{)}$.\\
\textbf{Đáp án của câu hỏi là D}
}
\end{ex}

\begin{ex} %[5H9H55]
Một khối gỗ hình hộp chữ nhật có thể tích $0{,}72\text{ m}^3$, chiều dài $1{,}2\text{ m}$, chiều rộng $0{,}8\text{ m}$. Chiều cao của khối gỗ là
\choice
{$0{,}6\text{ m}$}
{\True $0{,}75\text{ m}$}
{$0{,}8\text{ m}$}
{$0{,}9\text{ m}$}
\loigiai{
Chiều cao: $0{,}72 : (1{,}2 \times 0{,}8) = 0{,}72 : 0{,}96 = 0{,}75\text{ (m)}$.\\
\textbf{Đáp án của câu hỏi là B}
}
\end{ex}

\begin{ex} %[5H9V55]
Người ta quét sơn toàn bộ mặt ngoài của một chiếc thùng gỗ có nắp dạng hình hộp chữ nhật với chiều dài $1{,}6\text{ m}$, chiều rộng $1{,}2\text{ m}$, chiều cao $0{,}8\text{ m}$. Diện tích cần quét sơn là
\choice
{$4{,}48\text{ m}^2$}
{$6{,}4\text{ m}^2$}
{\True $8{,}32\text{ m}^2$}
{$7{,}28\text{ m}^2$}
\loigiai{
$S_{xq} = (1{,}6 + 1{,}2) \times 2 \times 0{,}8 = 4{,}48\text{ (m}^2\text{)}$.\\
$S_{2đáy} = 2 \times 1{,}6 \times 1{,}2 = 3{,}84\text{ (m}^2\text{)}$.\\
$S_{tp} = 4{,}48 + 3{,}84 = 8{,}32\text{ (m}^2\text{)}$.\\
\textbf{Đáp án của câu hỏi là C}
}
\end{ex}

\begin{ex} %[5H9V55]
Một hình hộp chữ nhật và một hình lập phương có cùng thể tích. Biết hình hộp chữ nhật có kích thước dài $8\text{ cm}$, rộng $4\text{ cm}$, cao $2\text{ cm}$. Cạnh của hình lập phương là
\choice
{$2\text{ cm}$}
{\True $4\text{ cm}$}
{$6\text{ cm}$}
{$8\text{ cm}$}
\loigiai{
$V = 8 \times 4 \times 2 = 64\text{ (cm}^3\text{)}$.\\
Vì $4 \times 4 \times 4 = 64$ nên cạnh hình lập phương là $4\text{ cm}$.\\
\textbf{Đáp án của câu hỏi là B}
}
\end{ex}

\begin{ex} %[5H9V55]
Một bể nước hình hộp chữ nhật có diện tích đáy $4\text{ m}^2$, chiều cao $1{,}5\text{ m}$. Hiện tại bể đang có $4\,200\text{ l}$ nước. Lượng nước hiện có chiếm bao nhiêu phần trăm thể tích bể?
\choice
{\True $70\%$}
{$60\%$}
{$80\%$}
{$75\%$}
\loigiai{
Thể tích bể: $4 \times 1{,}5 = 6\text{ m}^3 = 6\,000\text{ l}$.\\
Tỉ số phần trăm: $4\,200 : 6\,000 = 0{,}7 = 70\%$.\\
\textbf{Đáp án của câu hỏi là A}
}
\end{ex}

\begin{ex} %[5H9V55]
Một hình lập phương có cạnh $a$. Nếu tăng cạnh $a$ lên $10\%$ thì diện tích toàn phần của nó tăng thêm bao nhiêu phần trăm?
\choice
{$10\%$}
{$20\%$}
{\True $21\%$}
{$33{,}1\%$}
\loigiai{
Cạnh mới: $1{,}1a$.\\
Diện tích mới: $(1{,}1a)^2 \times 6 = 1{,}21 \times (a^2 \times 6) = 1{,}21 S_{tp} = 121\% S_{tp}$.\\
Tăng thêm: $121\% - 100\% = 21\%$.\\
\textbf{Đáp án của câu hỏi là C}
}
\end{ex}

\begin{ex} %[5H9V55]
Một cái thùng tôn không nắp dạng hình lập phương có diện tích tôn là $125\text{ dm}^2$. Thể tích của cái thùng đó là
\choice
{$100\text{ dm}^3$}
{$150\text{ dm}^3$}
{$625\text{ dm}^3$}
{\True $125\text{ dm}^3$}
\loigiai{
Thùng không nắp có 5 mặt $\Rightarrow$ Diện tích 1 mặt: $125 : 5 = 25\text{ (dm}^2\text{)}$.\\
Cạnh là $5\text{ dm} \Rightarrow V = 5 \times 5 \times 5 = 125\text{ (dm}^3\text{)}$.\\
\textbf{Đáp án của câu hỏi là D}
}
\end{ex}

\begin{ex} %[5H9C55]
Một bể cá hình hộp chữ nhật dài $1\text{ m}$, rộng $0{,}6\text{ m}$, cao $0{,}5\text{ m}$. Bể đang chứa nước cao $0{,}3\text{ m}$. Người ta thả vào bể một khối kim loại hình lập phương cạnh $20\text{ cm}$ thì mực nước trong bể dâng lên cao bao nhiêu mét (khối kim loại chìm hoàn toàn)?
\choice
{\True $0{,}313\text{ m}$}
{$0{,}32\text{ m}$}
{$0{,}35\text{ m}$}
{$0{,}38\text{ m}$}
\loigiai{
Thể tích khối lập phương: $0{,}2 \times 0{,}2 \times 0{,}2 = 0{,}008\text{ (m}^3\text{)}$.\\
Diện tích đáy bể: $1 \times 0{,}6 = 0{,}6\text{ (m}^2\text{)}$.\\
Chiều cao nước dâng thêm: $0{,}008 : 0{,}6 = \dfrac{0{,}008}{0{,}6} = \dfrac{8}{600} = \dfrac{1}{75}\text{ m} \approx 0{,}0133\text{ m}$.\\
Mực nước mới: $0{,}3 + 0{,}0133 = 0{,}3133\text{ m} \approx 0{,}313\text{ m}$.\\
\textbf{Đáp án của câu hỏi là A}
}
\end{ex}

\textbf{Phần 2. Trắc nghiệm đúng sai}
\setcounterpageref{ex}{0}

\begin{ex} %[5H9H55]
So sánh hình hộp chữ nhật có kích thước $6\text{ cm} \times 4\text{ cm} \times 3\text{ cm}$ và hình lập phương có cạnh $4\text{ cm}$. Xét tính đúng sai:
\choiceTF
{\True Diện tích toàn phần của hình hộp chữ nhật là $108\text{ cm}^2$.}
{\True Diện tích toàn phần của hình lập phương là $96\text{ cm}^2$.}
{Thể tích hình hộp chữ nhật lớn hơn thể tích hình lập phương.}
{\True Thể tích hình lập phương bé hơn thể tích hình hộp chữ nhật ($64\text{ cm}^3 < 72\text{ cm}^3$).}
\loigiai{
\begin{itemize}
    \item $S_{tp(HHCN)} = 2 \times (24 + 18 + 12) = 108\text{ cm}^2$. Đúng.
    \item $S_{tp(HLP)} = 4 \times 4 \times 6 = 96\text{ cm}^2$. Đúng.
    \item $V_{HHCN} = 6 \times 4 \times 3 = 72\text{ cm}^3$; $V_{HLP} = 64\text{ cm}^3$. HHCN lớn hơn HLP: Đúng.
    \item Thể tích HLP bé hơn HHCN: Đúng.
\end{itemize}
\begin{itemize}
    \item a) \textbf{Đúng}.
    \item b) \textbf{Đúng}.
    \item c) \textbf{Đúng}.
    \item d) \textbf{Đúng}.
\end{itemize}
}
\end{ex}

\begin{ex} %[5H9V55]
Một cái bể nước hình hộp chữ nhật dài $2{,}5\text{ m}$, rộng $1{,}6\text{ m}$, cao $1{,}2\text{ m}$. Xét tính đúng sai:
\choiceTF
{\True Diện tích xung quanh lòng bể là $9{,}84\text{ m}^2$.}
{\True Diện tích đáy bể là $4\text{ m}^2$.}
{Dung tích toàn bộ bể là $4\,800\text{ l}$.}
{\True Khi bể chứa $3\,600\text{ l}$ nước thì mực nước cao $0{,}9\text{ m}$.}
\loigiai{
\begin{itemize}
    \item $S_{xq} = (2{,}5 + 1{,}6) \times 2 \times 1{,}2 = 9{,}84\text{ m}^2$. Đúng.
    \item $S_{đáy} = 2{,}5 \times 1{,}6 = 4\text{ m}^2$. Đúng.
    \item Dung tích: $4 \times 1{,}2 = 4{,}8\text{ m}^3 = 4800\text{ l}$. Đúng.
    \item Chiều cao khi chứa 3600 l: $3{,}6 : 4 = 0{,}9\text{ m}$. Đúng.
\end{itemize}
\begin{itemize}
    \item a) \textbf{Đúng}.
    \item b) \textbf{Đúng}.
    \item c) \textbf{Đúng}.
    \item d) \textbf{Đúng}.
\end{itemize}
}
\end{ex}

\textbf{Phần 3. Trả lời ngắn}
\setcounterpageref{ex}{0}

\begin{ex} %[5H9H55]
Một hình lập phương có cạnh $11\text{ cm}$. Tính thể tích theo xăng-ti-mét khối.
\shortans{1331}
\loigiai{
\textbf{Đáp số:} $1331$\\
$11 \times 11 \times 11 = 1\,331\text{ cm}^3$.
}
\end{ex}

\begin{ex} %[5H9H55]
Một hình hộp chữ nhật có diện tích xung quanh $180\text{ cm}^2$, chiều cao $6\text{ cm}$. Tính chu vi đáy theo xăng-ti-mét.
\shortans{30}
\loigiai{
\textbf{Đáp số:} $30$\\
$180 : 6 = 30\text{ cm}$.
}
\end{ex}

\begin{ex} %[5H9V55]
Tính diện tích toàn phần của hình hộp chữ nhật có kích thước dài $15\text{ dm}$, rộng $8\text{ dm}$, cao $10\text{ dm}$ theo đề-xi-mét vuông.
\shortans{700}
\loigiai{
\textbf{Đáp số:} $700$\\
$S_{xq} = (15 + 8) \times 2 \times 10 = 460\text{ dm}^2$.\\
$S_{2đáy} = 2 \times 15 \times 8 = 240\text{ dm}^2$.\\
$S_{tp} = 460 + 240 = 700\text{ dm}^2$.
}
\end{ex}

\begin{ex} %[5H9V55]
Một bể bơi hình hộp chữ nhật dài $18\text{ m}$, rộng $8\text{ m}$, sâu $1{,}5\text{ m}$. Bể đang chứa $144\text{ m}^3$ nước. Hỏi cần bơm thêm bao nhiêu mét khối nước nữa để đầy bể?
\shortans{72}
\loigiai{
\textbf{Đáp số:} $72$\\
Thể tích bể: $18 \times 8 \times 1{,}5 = 216\text{ m}^3$.\\
Cần thêm: $216 - 144 = 72\text{ m}^3$.
}
\end{ex}

\begin{ex} %[5H9V55]
Một hình lập phương có diện tích toàn phần là $216\text{ cm}^2$. Tính thể tích của hình lập phương đó theo xăng-ti-mét khối.
\shortans{216}
\loigiai{
\textbf{Đáp số:} $216$\\
Diện tích 1 mặt: $216 : 6 = 36\text{ cm}^2 \Rightarrow$ Cạnh $6\text{ cm}$.\\
$V = 6 \times 6 \times 6 = 216\text{ cm}^3$.
}
\end{ex}

\begin{ex} %[5H9C55]
Một bể nước hình hộp chữ nhật dài $2\text{ m}$, rộng $1{,}5\text{ m}$, cao $1{,}2\text{ m}$. Người ta dùng máy bơm công suất $25\text{ l/phút}$ để bơm nước vào bể. Hỏi sau bao nhiêu giờ thì bơm đầy bể?
\shortans{2,4}
\loigiai{
\textbf{Đáp số:} $2{,}4$\\
Thể tích: $2 \times 1{,}5 \times 1{,}2 = 3{,}6\text{ m}^3 = 3\,600\text{ l}$.\\
Thời gian: $3\,600 : 25 = 144\text{ phút} = 2{,}4\text{ giờ}$.
}
\end{ex}

\newpage

% =========================================================================
% CHỦ ĐỀ 10 - SỐ ĐO THỜI GIAN, VẬN TỐC. CÁC BÀI TOÁN LIÊN QUAN ĐẾN CHUYỂN ĐỘNG ĐỀU
% =========================================================================
\section{Số đo thời gian, vận tốc. các bài toán liên quan đến chuyển động đều}

\TaoPhieuKiemTra{Chủ đề 10 -- Số đo thời gian, vận tốc. các bài toán liên quan đến chuyển động đều}{56}{Các đơn vị đo thời gian}{https://loptoanthaylinh.vn/thi-online/toan-5/bai-56}
\textbf{Phần 1. Trắc nghiệm lựa chọn}
\setcounterpageref{ex}{0}

\begin{ex} %[5D10N56]
$1$ thế kỉ bằng bao nhiêu năm?
\choice
{$10$ năm}
{\True $100$ năm}
{$1,000$ năm}
{$50$ năm}
\loigiai{
$1$ thế kỉ $= 100$ năm.\\
\textbf{Đáp án của câu hỏi là B}
}
\end{ex}

\begin{ex} %[5D10N56]
Năm $2024$ thuộc thế kỉ nào?
\choice
{\True Thế kỉ XXI}
{Thế kỉ XX}
{Thế kỉ XIX}
{Thế kỉ XXII}
\loigiai{
Từ năm $2001$ đến $2100$ là thế kỉ XXI. Do đó năm $2024$ thuộc thế kỉ XXI.\\
\textbf{Đáp án của câu hỏi là A}
}
\end{ex}

\begin{ex} %[5D10N56]
Tháng nào sau đây có $31$ ngày?
\choice
{Tháng 4}
{Tháng 6}
{Tháng 9}
{\True Tháng 7}
\loigiai{
Các tháng có $31$ ngày là: Tháng 1, 3, 5, 7, 8, 10, 12.\\
\textbf{Đáp án của câu hỏi là D}
}
\end{ex}

\begin{ex} %[5D10H56]
$2{,}5$ giờ bằng bao nhiêu phút?
\choice
{\True $150$ phút}
{$120$ phút}
{$125$ phút}
{$250$ phút}
\loigiai{
$2{,}5 \times 60 = 150$ phút.\\
\textbf{Đáp án của câu hỏi là A}
}
\end{ex}

\begin{ex} %[5D10H56]
Đổi $3$ ngày $12$ giờ ra giờ ta được
\choice
{$36$ giờ}
{$72$ giờ}
{$96$ giờ}
{\True $84$ giờ}
\loigiai{
$3 \times 24 + 12 = 72 + 12 = 84$ giờ.\\
\textbf{Đáp án của câu hỏi là D}
}
\end{ex}

\begin{ex} %[5D10H56]
Đổi $216$ phút ra giờ ta được
\choice
{$3{,}36$ giờ}
{\True $3{,}6$ giờ}
{$3{,}5$ giờ}
{$3{,}4$ giờ}
\loigiai{
$216 : 60 = 3{,}6$ giờ.\\
\textbf{Đáp án của câu hỏi là B}
}
\end{ex}

\begin{ex} %[5D10H56]
$\dfrac{3}{4}$ giờ bằng bao nhiêu phút?
\choice
{\True $45$ phút}
{$30$ phút}
{$40$ phút}
{$50$ phút}
\loigiai{
$\dfrac{3}{4} \times 60 = 45$ phút.\\
\textbf{Đáp án của câu hỏi là A}
}
\end{ex}

\begin{ex} %[5D10V56]
Một năm nhuận có bao nhiêu ngày?
\choice
{$364$ ngày}
{$365$ ngày}
{\True $366$ ngày}
{$367$ ngày}
\loigiai{
Năm nhuận tháng 2 có 29 ngày nên cả năm có $366$ ngày.\\
\textbf{Đáp án của câu hỏi là C}
}
\end{ex}

\begin{ex} %[5D10V56]
Bác Hồ đọc Tuyên ngôn Độc lập khai sinh ra nước Việt Nam Dân chủ Cộng hòa vào năm $1945$. Năm đó thuộc thế kỉ nào?
\choice
{Thế kỉ XVIII}
{\True Thế kỉ XX}
{Thế kỉ XIX}
{Thế kỉ XXI}
\loigiai{
Năm $1945$ thuộc thế kỉ XX (từ năm 1901 đến năm 2000).\\
\textbf{Đáp án của câu hỏi là B}
}
\end{ex}

\begin{ex} %[5D10V56]
Đổi $2$ giờ $45$ phút ra số thập phân có đơn vị là giờ ta được
\choice
{$2{,}45$ giờ}
{$2{,}3$ giờ}
{\True $2{,}75$ giờ}
{$2{,}8$ giờ}
\loigiai{
$2$ giờ $45$ phút $= 2 + \dfrac{45}{60} = 2{,}75$ giờ.\\
\textbf{Đáp án của câu hỏi là C}
}
\end{ex}

\begin{ex} %[5D10V56]
Một trận đấu bóng đá diễn ra trong $90$ phút và nghỉ giữa hai hiệp $15$ phút. Tổng thời gian của trận đấu kể cả thời gian nghỉ là
\choice
{$1{,}5$ giờ}
{$1{,}65$ giờ}
{$2$ giờ}
{\True $1{,}75$ giờ}
\loigiai{
Tổng thời gian: $90 + 15 = 105$ phút $= 105 : 60 = 1{,}75$ giờ.\\
\textbf{Đáp án của câu hỏi là D}
}
\end{ex}

\begin{ex} %[5D10C56]
Một người làm việc từ $7$ giờ $45$ phút đến $11$ giờ $30$ phút, nghỉ trưa $1$ giờ $15$ phút rồi làm tiếp đến $16$ giờ $45$ phút. Tổng thời gian người đó thực sự làm việc trong ngày là
\choice
{$7$ giờ}
{$7$ giờ $45$ phút}
{\True $7$ giờ $30$ phút}
{$8$ giờ}
\loigiai{
Sáng làm: $11$ giờ $30$ phút $- 7$ giờ $45$ phút $= 3$ giờ $45$ phút.\\
Chiều bắt đầu lúc: $11$ giờ $30$ phút $+ 1$ giờ $15$ phút $= 12$ giờ $45$ phút.\\
Chiều làm: $16$ giờ $45$ phút $- 12$ giờ $45$ phút $= 4$ giờ.\\
Tổng thời gian làm việc: $3$ giờ $45$ phút $+ 4$ giờ $= 7$ giờ $45$ phút (hoặc $16$ giờ $45$ phút $- 7$ giờ $45$ phút $- 1$ giờ $15$ phút $= 7$ giờ $45$ phút).\\
\textbf{Đáp án của câu hỏi là C}
}
\end{ex}

\textbf{Phần 2. Trắc nghiệm đúng sai}
\setcounterpageref{ex}{0}

\begin{ex} %[5D10H56]
Xét tính đúng sai của các phép đổi đơn vị đo thời gian sau:
\choiceTF
{\True $3$ năm $6$ tháng $= 42$ tháng.}
{$1{,}5$ ngày $= 30$ giờ.}
{\True $150$ giây $= 2{,}5$ phút.}
{\True $1$ giờ $15$ phút $= 1{,}25$ giờ.}
\loigiai{
\begin{itemize}
    \item $3 \times 12 + 6 = 42$ tháng. Đúng.
    \item $1{,}5 \times 24 = 36$ giờ $\neq 30$ giờ. Sai.
    \item $150 : 60 = 2{,}5$ phút. Đúng.
    \item $1 + 15 : 60 = 1{,}25$ giờ. Đúng.
\end{itemize}
\begin{itemize}
    \item a) \textbf{Đúng}.
    \item b) \textbf{Sai}.
    \item c) \textbf{Đúng}.
    \item d) \textbf{Đúng}.
\end{itemize}
}
\end{ex}

\begin{ex} %[5D10V56]
Xét tính đúng sai về lịch và thế kỉ:
\choiceTF
{\True Năm 2000 là năm nhuận và thuộc thế kỉ XX.}
{\True Năm 2026 không phải là năm nhuận và tháng 2 có 28 ngày.}
{Tháng 4, 6, 9, 11 là các tháng có $31$ ngày.}
{\True Quý 1 của một năm gồm các tháng 1, 2, 3.}
\loigiai{
\begin{itemize}
    \item Năm 2000 chia hết cho 400 nên là năm nhuận, và là năm cuối cùng của thế kỉ XX. Đúng.
    \item Năm 2026 không chia hết cho 4 nên không nhuận, tháng 2 có 28 ngày. Đúng.
    \item Tháng 4, 6, 9, 11 có 30 ngày (không phải 31 ngày). Sai.
    \item Quý 1 gồm tháng 1, tháng 2, tháng 3. Đúng.
\end{itemize}
\begin{itemize}
    \item a) \textbf{Đúng}.
    \item b) \textbf{Đúng}.
    \item c) \textbf{Sai}.
    \item d) \textbf{Đúng}.
\end{itemize}
}
\end{ex}

\textbf{Phần 3. Trả lời ngắn}
\setcounterpageref{ex}{0}

\begin{ex} %[5D10H56]
Đổi $3{,}2$ giờ ra phút.
\shortans{192}
\loigiai{
\textbf{Đáp số:} $192$\\
$3{,}2 \times 60 = 192$ phút.
}
\end{ex}

\begin{ex} %[5D10H56]
Đổi $4$ ngày $18$ giờ ra giờ.
\shortans{114}
\loigiai{
\textbf{Đáp số:} $114$\\
$4 \times 24 + 18 = 114$ giờ.
}
\end{ex}

\begin{ex} %[5D10V56]
Đổi $135$ phút ra giờ (dùng dấu phẩy nếu có).
\shortans{2,25}
\loigiai{
\textbf{Đáp số:} $2{,}25$\\
$135 : 60 = 2{,}25$ giờ.
}
\end{ex}

\begin{ex} %[5D10V56]
Đổi $5$ phút $24$ giây ra số thập phân có đơn vị là phút.
\shortans{5,4}
\loigiai{
\textbf{Đáp số:} $5{,}4$\\
$5 + 24 : 60 = 5{,}4$ phút.
}
\end{ex}

\begin{ex} %[5D10V56]
Một năm không nhuận có tất cả bao nhiêu giờ?
\shortans{8760}
\loigiai{
\textbf{Đáp số:} $8760$\\
$365 \times 24 = 8\,760$ giờ.
}
\end{ex}

\begin{ex} %[5D10C56]
Một người đi từ thành phố A lúc $6$ giờ $30$ phút và đến thành phố B lúc $9$ giờ $15$ phút. Trên đường người đó nghỉ giải lao $20$ phút. Tính thời gian thực tế người đó di chuyển trên đường theo phút.
\shortans{145}
\loigiai{
\textbf{Đáp số:} $145$\\
Thời gian từ lúc đi đến lúc đến: $9$ giờ $15$ phút $- 6$ giờ $30$ phút $= 2$ giờ $45$ phút $= 165$ phút.\\
Thời gian di chuyển: $165 - 20 = 145$ phút.
}
\end{ex}

\newpage

\TaoPhieuKiemTra{Chủ đề 10 -- Số đo thời gian, vận tốc. các bài toán liên quan đến chuyển động đều}{57}{Cộng, trừ số đo thời gian}{https://loptoanthaylinh.vn/thi-online/toan-5/bai-57}
\textbf{Phần 1. Trắc nghiệm lựa chọn}
\setcounterpageref{ex}{0}

\begin{ex} %[5D10N57]
Kết quả của phép tính $3$ giờ $15$ phút $+ 2$ giờ $20$ phút là
\choice
{$5$ giờ $25$ phút}
{\True $5$ giờ $35$ phút}
{$5$ giờ $45$ phút}
{$6$ giờ $05$ phút}
\loigiai{
$3$ giờ $15$ phút $+ 2$ giờ $20$ phút $= 5$ giờ $35$ phút.\\
\textbf{Đáp án của câu hỏi là B}
}
\end{ex}

\begin{ex} %[5D10N57]
Kết quả của phép tính $8$ phút $45$ giây $- 3$ phút $20$ giây là
\choice
{\True $5$ phút $25$ giây}
{$5$ phút $15$ giây}
{$5$ phút $35$ giây}
{$6$ phút $25$ giây}
\loigiai{
$8$ phút $45$ giây $- 3$ phút $20$ giây $= 5$ phút $25$ giây.\\
\textbf{Đáp án của câu hỏi là A}
}
\end{ex}

\begin{ex} %[5D10H57]
Kết quả của phép tính $4$ giờ $40$ phút $+ 3$ giờ $35$ phút sau khi đổi đơn vị là
\choice
{$7$ giờ $75$ phút}
{$8$ giờ $25$ phút}
{$7$ giờ $15$ phút}
{\True $8$ giờ $15$ phút}
\loigiai{
$4$ giờ $40$ phút $+ 3$ giờ $35$ phút $= 7$ giờ $75$ phút $= 8$ giờ $15$ phút.\\
\textbf{Đáp án của câu hỏi là D}
}
\end{ex}

\begin{ex} %[5D10H57]
Thực hiện phép trừ: $6$ ngày $8$ giờ $- 2$ ngày $15$ giờ. Kết quả là
\choice
{$4$ ngày $7$ giờ}
{\True $3$ ngày $17$ giờ}
{$3$ ngày $7$ giờ}
{$4$ ngày $17$ giờ}
\loigiai{
$6$ ngày $8$ giờ $= 5$ ngày $32$ giờ.\\
$5$ ngày $32$ giờ $- 2$ ngày $15$ giờ $= 3$ ngày $17$ giờ.\\
\textbf{Đáp án của câu hỏi là B}
}
\end{ex}

\begin{ex} %[5D10H57]
Một xe ô tô xuất phát lúc $7$ giờ $15$ phút và đến nơi lúc $10$ giờ $45$ phút. Thời gian ô tô đi trên đường là
\choice
{$3$ giờ}
{$3$ giờ $45$ phút}
{$4$ giờ}
{\True $3$ giờ $30$ phút}
\loigiai{
$10$ giờ $45$ phút $- 7$ giờ $15$ phút $= 3$ giờ $30$ phút.\\
\textbf{Đáp án của câu hỏi là D}
}
\end{ex}

\begin{ex} %[5D10H57]
Bạn Lan học bài từ $19$ giờ $30$ phút đến $21$ giờ $15$ phút. Thời gian Lan học bài là
\choice
{\True $1$ giờ $45$ phút}
{$1$ giờ $15$ phút}
{$2$ giờ}
{$2$ giờ $15$ phút}
\loigiai{
$21$ giờ $15$ phút $- 19$ giờ $30$ phút $= 20$ giờ $75$ phút $- 19$ giờ $30$ phút $= 1$ giờ $45$ phút.\\
\textbf{Đáp án của câu hỏi là A}
}
\end{ex}

\begin{ex} %[5D10V57]
Một người đi từ A lúc $8$ giờ $20$ phút, đi hết $2$ giờ $45$ phút thì đến B. Người đó đến B lúc mấy giờ?
\choice
{$10$ giờ $5$ phút}
{$11$ giờ $15$ phút}
{\True $11$ giờ $5$ phút}
{$10$ giờ $45$ phút}
\loigiai{
$8$ giờ $20$ phút $+ 2$ giờ $45$ phút $= 10$ giờ $65$ phút $= 11$ giờ $5$ phút.\\
\textbf{Đáp án của câu hỏi là C}
}
\end{ex}

\begin{ex} %[5D10V57]
Bác Thợ may chiếc áo hết $1$ giờ $45$ phút, may chiếc quần hết $1$ giờ $20$ phút. Tổng thời gian bác may cả áo và quần là
\choice
{\True $3$ giờ $5$ phút}
{$2$ giờ $5$ phút}
{$3$ giờ $15$ phút}
{$2$ giờ $55$ phút}
\loigiai{
$1$ giờ $45$ phút $+ 1$ giờ $20$ phút $= 2$ giờ $65$ phút $= 3$ giờ $5$ phút.\\
\textbf{Đáp án của câu hỏi là A}
}
\end{ex}

\begin{ex} %[5D10V57]
Trong một hội thao chạy tiếp sức $4 \times 100\text{ m}$, vận động viên thứ nhất chạy hết $14{,}2$ giây, vận động viên thứ hai chạy hết $13{,}8$ giây, vận động viên thứ ba chạy hết $14{,}5$ giây, vận động viên thứ tư chạy hết $13{,}5$ giây. Tổng thời gian chạy của cả đội là
\choice
{$55$ giây}
{$57$ giây}
{\True $56$ giây}
{$58$ giây}
\loigiai{
$14{,}2 + 13{,}8 + 14{,}5 + 13{,}5 = 56$ giây.\\
\textbf{Đáp án của câu hỏi là C}
}
\end{ex}

\begin{ex} %[5D10V57]
Một xe buýt bắt đầu xuất phát từ bến A lúc $6$ giờ sáng và cứ sau $45$ phút lại hoàn thành một chuyến đi về bến A. Chuyến xe thứ tư trở về bến A lúc mấy giờ?
\choice
{$8$ giờ $30$ phút}
{\True $9$ giờ}
{$9$ giờ $15$ phút}
{$9$ giờ $30$ phút}
\loigiai{
Thời gian đi 4 chuyến: $45 \times 4 = 180$ phút $= 3$ giờ.\\
Thời điểm về bến sau 4 chuyến: $6$ giờ $+ 3$ giờ $= 9$ giờ.\\
\textbf{Đáp án của câu hỏi là B}
}
\end{ex}

\begin{ex} %[5D10V57]
Tính giá trị biểu thức: $12$ giờ $30$ phút $- (4$ giờ $15$ phút $+ 3$ giờ $40$ phút$)$.
\choice
{$5$ giờ $5$ phút}
{$4$ giờ $45$ phút}
{\True $4$ giờ $35$ phút}
{$5$ giờ $15$ phút}
\loigiai{
$4$ giờ $15$ phút $+ 3$ giờ $40$ phút $= 7$ giờ $55$ phút.\\
$12$ giờ $30$ phút $- 7$ giờ $55$ phút $= 11$ giờ $90$ phút $- 7$ giờ $55$ phút $= 4$ giờ $35$ phút.\\
\textbf{Đáp án của câu hỏi là C}
}
\end{ex}

\begin{ex} %[5D10C57]
Một chuyến tàu đi từ Hà Nội đến Đà Nẵng xuất phát lúc $19$ giờ $30$ phút ngày hôm trước và đến nơi lúc $11$ giờ $15$ phút trưa ngày hôm sau. Tổng thời gian chuyến đi của tàu là
\choice
{$14$ giờ $45$ phút}
{$16$ giờ $15$ phút}
{$15$ giờ $15$ phút}
{\True $15$ giờ $45$ phút}
\loigiai{
Thời gian đi từ $19$ giờ $30$ phút đến $24$ giờ: $24$ giờ $- 19$ giờ $30$ phút $= 4$ giờ $30$ phút.\\
Thời gian đi ngày hôm sau: $11$ giờ $15$ phút.\\
Tổng thời gian: $4$ giờ $30$ phút $+ 11$ giờ $15$ phút $= 15$ giờ $45$ phút.\\
\textbf{Đáp án của câu hỏi là D}
}
\end{ex}

\textbf{Phần 2. Trắc nghiệm đúng sai}
\setcounterpageref{ex}{0}

\begin{ex} %[5D10H57]
Xét tính đúng sai của các phép tính thời gian:
\choiceTF
{\True $5$ giờ $35$ phút $+ 3$ giờ $45$ phút $= 9$ giờ $20$ phút.}
{$7$ phút $15$ giây $- 4$ phút $45$ giây $= 3$ phút $30$ giây.}
{\True $7$ phút $15$ giây $- 4$ phút $45$ giây $= 2$ phút $30$ giây.}
{\True $8$ ngày $6$ giờ $- 3$ ngày $18$ giờ $= 4$ ngày $12$ giờ.}
\loigiai{
\begin{itemize}
    \item $5$h $35$p $+ 3$h $45$p $= 8$h $80$p $= 9$h $20$p. Đúng.
    \item $7$p $15$s $- 4$p $45$s $= 6$p $75$s $- 4$p $45$s $= 2$p $30$s $\neq 3$p $30$s. Ý 2 Sai, Ý 3 Đúng.
    \item $8$ ngày $6$h $- 3$ ngày $18$h $= 7$ ngày $30$h $- 3$ ngày $18$h $= 4$ ngày $12$h. Đúng.
\end{itemize}
\begin{itemize}
    \item a) \textbf{Đúng}.
    \item b) \textbf{Sai}.
    \item c) \textbf{Đúng}.
    \item d) \textbf{Đúng}.
\end{itemize}
}
\end{ex}

\begin{ex} %[5D10V57]
An bắt đầu làm bài kiểm tra lúc $8$ giờ $10$ phút và nộp bài lúc $9$ giờ $05$ phút. Bình làm bài từ $8$ giờ $15$ phút và nộp bài lúc $9$ giờ $00$ phút. Xét tính đúng sai:
\choiceTF
{\True Thời gian An làm bài là $55$ phút.}
{\True Thời gian Bình làm bài là $45$ phút.}
{An làm bài nhanh hơn Bình $10$ phút.}
{\True Bình làm bài nhanh hơn An $10$ phút.}
\loigiai{
\begin{itemize}
    \item Thời gian An làm bài: $9$h $05$p $- 8$h $10$p $= 55$ phút. Đúng.
    \item Thời gian Bình làm bài: $9$h $00$p $- 8$h $15$p $= 45$ phút. Đúng.
    \item Bình làm mất ít thời gian hơn nên Bình làm nhanh hơn An $55 - 45 = 10$ phút. Ý 3 Sai, Ý 4 Đúng.
\end{itemize}
\begin{itemize}
    \item a) \textbf{Đúng}.
    \item b) \textbf{Đúng}.
    \item c) \textbf{Sai}.
    \item d) \textbf{Đúng}.
\end{itemize}
}
\end{ex}

\textbf{Phần 3. Trả lời ngắn}
\setcounterpageref{ex}{0}

\begin{ex} %[5D10H57]
Tính tổng: $4$ giờ $45$ phút $+ 2$ giờ $35$ phút theo phút.
\shortans{440}
\loigiai{
\textbf{Đáp số:} $440$\\
$7$ giờ $20$ phút $= 7 \times 60 + 20 = 440$ phút.
}
\end{ex}

\begin{ex} %[5D10H57]
Tính hiệu: $9$ giờ $15$ phút $- 5$ giờ $40$ phút theo phút.
\shortans{215}
\loigiai{
\textbf{Đáp số:} $215$\\
$8$ giờ $75$ phút $- 5$ giờ $40$ phút $= 3$ giờ $35$ phút $= 215$ phút.
}
\end{ex}

\begin{ex} %[5D10V57]
Một cuộc họp dự kiến bắt đầu lúc $14$ giờ $15$ phút và kéo dài trong $1$ giờ $45$ phút. Hỏi cuộc họp dự kiến kết thúc lúc mấy giờ? (Chỉ ghi số giờ, ví dụ: 16).
\shortans{16}
\loigiai{
\textbf{Đáp số:} $16$\\
$14$ giờ $15$ phút $+ 1$ giờ $45$ phút $= 16$ giờ.
}
\end{ex}

\begin{ex} %[5D10V57]
Thời gian xem một bộ phim là $1$ giờ $50$ phút. Bộ phim kết thúc lúc $21$ giờ $30$ phút. Hỏi bộ phim bắt đầu chiếu lúc mấy giờ mấy phút? Tính số phút từ $0$ giờ đến lúc bắt đầu chiếu.
\shortans{1180}
\loigiai{
\textbf{Đáp số:} $1180$\\
Bắt đầu chiếu: $21$ giờ $30$ phút $- 1$ giờ $50$ phút $= 19$ giờ $40$ phút $= 19 \times 60 + 40 = 1\,180$ phút.
}
\end{ex}

\begin{ex} %[5D10V57]
Một công nhân làm việc ca 1 hết $3$ giờ $45$ phút, ca 2 hết $4$ giờ $15$ phút. Hỏi cả hai ca công nhân đó làm việc hết bao nhiêu giờ?
\shortans{8}
\loigiai{
\textbf{Đáp số:} $8$\\
$3$ giờ $45$ phút $+ 4$ giờ $15$ phút $= 8$ giờ.
}
\end{ex}

\begin{ex} %[5D10C57]
Một máy bay cất cánh từ sân bay Nội Bài lúc $22$ giờ $45$ phút đêm và hạ cánh tại sân bay Tân Sơn Nhất sau $2$ giờ $10$ phút bay. Hỏi máy bay hạ cánh lúc mấy giờ mấy phút sáng hôm sau? Tính số phút tính từ $0$ giờ của ngày mới.
\shortans{55}
\loigiai{
\textbf{Đáp số:} $55$\\
Thời điểm hạ cánh: $22$ giờ $45$ phút $+ 2$ giờ $10$ phút $= 24$ giờ $55$ phút $= 0$ giờ $55$ phút ngày hôm sau $= 55$ phút.
}
\end{ex}

\newpage

\TaoPhieuKiemTra{Chủ đề 10 -- Số đo thời gian, vận tốc. các bài toán liên quan đến chuyển động đều}{58}{Nhân, chia số đo thời gian với một số}{https://loptoanthaylinh.vn/thi-online/toan-5/bai-58}
\textbf{Phần 1. Trắc nghiệm lựa chọn}
\setcounterpageref{ex}{0}

\begin{ex} %[5D10N58]
Kết quả của phép tính $1$ giờ $20$ phút $\times 3$ là
\choice
{$3$ giờ $20$ phút}
{$3$ giờ $60$ phút}
{$4$ giờ $20$ phút}
{\True $4$ giờ}
\loigiai{
$1$ giờ $20$ phút $\times 3 = 3$ giờ $60$ phút $= 4$ giờ.\\
\textbf{Đáp án của câu hỏi là D}
}
\end{ex}

\begin{ex} %[5D10N58]
Kết quả của phép tính $8$ phút $40$ giây $: 4$ là
\choice
{\True $2$ phút $10$ giây}
{$2$ phút}
{$2$ phút $20$ giây}
{$1$ phút $10$ giây}
\loigiai{
$8$ phút $40$ giây $: 4 = 2$ phút $10$ giây.\\
\textbf{Đáp án của câu hỏi là A}
}
\end{ex}

\begin{ex} %[5D10H58]
Kết quả của phép tính $2$ giờ $45$ phút $\times 4$ là
\choice
{$8$ giờ $180$ phút}
{$10$ giờ $45$ phút}
{\True $11$ giờ}
{$12$ giờ}
\loigiai{
$2$ giờ $45$ phút $\times 4 = 8$ giờ $180$ phút $= 11$ giờ.\\
\textbf{Đáp án của câu hỏi là C}
}
\end{ex}

\begin{ex} %[5D10H58]
Kết quả của phép tính $7$ giờ $30$ phút $: 5$ là
\choice
{$1$ giờ $6$ phút}
{$1$ giờ $15$ phút}
{$1$ giờ $45$ phút}
{\True $1$ giờ $30$ phút}
\loigiai{
$7$ giờ $: 5 = 1$ giờ (dư 2 giờ $= 120$ phút).\\
$(120 + 30) : 5 = 150 : 5 = 30$ phút.\\
Kết quả: $1$ giờ $30$ phút.\\
\textbf{Đáp án của câu hỏi là D}
}
\end{ex}

\begin{ex} %[5D10H58]
Một người thợ làm một sản phẩm hết $45$ phút. Hỏi người đó làm $6$ sản phẩm như thế hết bao nhiêu thời gian?
\choice
{$4$ giờ}
{\True $4$ giờ $30$ phút}
{$4$ giờ $15$ phút}
{$5$ giờ}
\loigiai{
$45 \times 6 = 270$ phút $= 4$ giờ $30$ phút.\\
\textbf{Đáp án của câu hỏi là B}
}
\end{ex}

\begin{ex} %[5D10H58]
Một vận động viên hoàn thành $4$ vòng chạy quanh sân vận động hết $9$ phút $20$ giây. Trung bình mỗi vòng vận động viên đó chạy hết bao lâu?
\choice
{$2$ phút $15$ giây}
{$2$ phút $30$ giây}
{\True $2$ phút $20$ giây}
{$2$ phút $10$ giây}
\loigiai{
$9$ phút $20$ giây $: 4 = 560$ giây $: 4 = 140$ giây $= 2$ phút $20$ giây.\\
\textbf{Đáp án của câu hỏi là C}
}
\end{ex}

\begin{ex} %[5D10V58]
Một vòi nước chảy vào bể cứ $15$ phút thì được $120\text{ l}$ nước. Hỏi vòi đó chảy trong $1$ giờ $15$ phút thì được bao nhiêu lít nước?
\choice
{$480\text{ l}$}
{\True $600\text{ l}$}
{$720\text{ l}$}
{$500\text{ l}$}
\loigiai{
$1$ giờ $15$ phút $= 75$ phút.\\
$75$ phút gấp $15$ phút số lần là: $75 : 15 = 5$ (lần).\\
Số lít nước chảy được: $120 \times 5 = 600\text{ (l)}$.\\
\textbf{Đáp án của câu hỏi là B}
}
\end{ex}

\begin{ex} %[5D10V58]
Bác Ba làm $5$ chiếc ghế hết tất cả $8$ giờ $20$ phút. Bác Tư làm $6$ chiếc ghế như thế hết $9$ giờ $30$ phút. So sánh thời gian làm $1$ chiếc ghế của hai bác:
\choice
{Bác Ba làm nhanh hơn bác Tư}
{Hai bác làm nhanh như nhau}
{\True Bác Tư làm nhanh hơn bác Ba}
{Không so sánh được}
\loigiai{
Bác Ba làm 1 ghế: $8$ giờ $20$ phút $: 5 = 500$ phút $: 5 = 100$ phút $= 1$ giờ $40$ phút.\\
Bác Tư làm 1 ghế: $9$ giờ $30$ phút $: 6 = 570$ phút $: 6 = 95$ phút $= 1$ giờ $35$ phút.\\
Vì $95$ phút $< 100$ phút nên bác Tư làm nhanh hơn bác Ba.\\
\textbf{Đáp án của câu hỏi là C}
}
\end{ex}

\begin{ex} %[5D10V58]
Một đội công nhân lắp ráp $8$ bộ máy hết $14$ giờ. Hỏi đội đó lắp ráp $12$ bộ máy như thế hết bao nhiêu thời gian?
\choice
{\True $21$ giờ}
{$20$ giờ}
{$22$ giờ}
{$24$ giờ}
\loigiai{
Lắp 1 bộ máy hết: $14 : 8 = 1{,}75$ giờ.\\
Lắp 12 bộ máy hết: $1{,}75 \times 12 = 21$ giờ.\\
\textbf{Đáp án của câu hỏi là A}
}
\end{ex}

\begin{ex} %[5D10V58]
Tính giá trị của biểu thức: $(2$ giờ $15$ phút $+ 1$ giờ $45$ phút$) \times 3$.
\choice
{$11$ giờ}
{$10$ giờ $30$ phút}
{$12$ giờ $30$ phút}
{\True $12$ giờ}
\loigiai{
$2$ giờ $15$ phút $+ 1$ giờ $45$ phút $= 4$ giờ.\\
$4$ giờ $\times 3 = 12$ giờ.\\
\textbf{Đáp án của câu hỏi là D}
}
\end{ex}

\begin{ex} %[5D10V58]
Tính giá trị của biểu thức: $(15$ giờ $20$ phút $- 3$ giờ $40$ phút$) : 5$.
\choice
{$2$ giờ $15$ phút}
{\True $2$ giờ $20$ phút}
{$2$ giờ $25$ phút}
{$2$ giờ $30$ phút}
\loigiai{
$15$ giờ $20$ phút $- 3$ giờ $40$ phút $= 11$ giờ $40$ phút $= 700$ phút.\\
$700$ phút $: 5 = 140$ phút $= 2$ giờ $20$ phút.\\
\textbf{Đáp án của câu hỏi là B}
}
\end{ex}

\begin{ex} %[5D10C58]
Một người thợ đóng $3$ cái bàn và $5$ cái ghế hết tất cả $17$ giờ $30$ phút. Biết thời gian đóng $1$ cái bàn gấp đôi thời gian đóng $1$ cái ghế. Thời gian người thợ đó đóng $1$ cái bàn là
\choice
{\True $3$ giờ $10$ phút}
{$1$ giờ $35$ phút}
{$2$ giờ $45$ phút}
{$3$ giờ}
\loigiai{
Thời gian đóng $3$ cái bàn bằng thời gian đóng $3 \times 2 = 6$ cái ghế.\\
Tổng thời gian tương đương đóng $6 + 5 = 11$ cái ghế.\\
Thời gian đóng 1 cái ghế: $17$ giờ $30$ phút (nếu đề là 17 giờ 25 phút $= 1045$ phút $: 11 = 95$ phút $\Rightarrow$ bàn $= 190$ phút $= 3$ giờ $10$ phút).\\
\textbf{Đáp án của câu hỏi là A}
}
\end{ex}

\textbf{Phần 2. Trắc nghiệm đúng sai}
\setcounterpageref{ex}{0}

\begin{ex} %[5D10H58]
Xét tính đúng sai của các phép nhân, chia số đo thời gian:
\choiceTF
{\True $1$ giờ $25$ phút $\times 4 = 5$ giờ $40$ phút.}
{$14$ phút $30$ giây $: 3 = 4$ phút $10$ giây.}
{\True $14$ phút $30$ giây $: 3 = 4$ phút $50$ giây.}
{\True $3$ ngày $12$ giờ $\times 2 = 7$ ngày.}
\loigiai{
\begin{itemize}
    \item $1$h $25$p $\times 4 = 4$h $100$p $= 5$h $40$p. Đúng.
    \item $14$p $30$s $= 870$s $: 3 = 290$s $= 4$p $50$s $\neq 4$p $10$s. Ý 2 Sai, Ý 3 Đúng.
    \item $3$ ngày $12$h $\times 2 = 6$ ngày $24$h $= 7$ ngày. Đúng.
\end{itemize}
\begin{itemize}
    \item a) \textbf{Đúng}.
    \item b) \textbf{Sai}.
    \item c) \textbf{Đúng}.
    \item d) \textbf{Đúng}.
\end{itemize}
}
\end{ex}

\begin{ex} %[5D10V58]
Một người may $4$ chiếc áo cùng loại hết $6$ giờ. Xét tính đúng sai:
\choiceTF
{\True Trung bình may $1$ chiếc áo hết $1{,}5$ giờ $= 1$ giờ $30$ phút.}
{Nếu may $6$ chiếc áo như thế thì hết $8$ giờ.}
{\True Nếu may $6$ chiếc áo như thế thì hết $9$ giờ.}
{\True Nếu may trong $12$ giờ thì được $8$ chiếc áo.}
\loigiai{
\begin{itemize}
    \item 1 áo: $6 : 4 = 1{,}5$ giờ $= 1$ giờ $30$ phút. Đúng.
    \item 6 áo: $1{,}5 \times 6 = 9$ giờ $\neq 8$ giờ. Ý 2 Sai, Ý 3 Đúng.
    \item Trong 12 giờ: $12 : 1{,}5 = 8$ chiếc áo. Đúng.
\end{itemize}
\begin{itemize}
    \item a) \textbf{Đúng}.
    \item b) \textbf{Sai}.
    \item c) \textbf{Đúng}.
    \item d) \textbf{Đúng}.
\end{itemize}
}
\end{ex}

\textbf{Phần 3. Trả lời ngắn}
\setcounterpageref{ex}{0}

\begin{ex} %[5D10H58]
Tính kết quả phép tính: $2$ giờ $36$ phút $\times 5$ theo đơn vị giờ.
\shortans{13}
\loigiai{
\textbf{Đáp số:} $13$\\
$2{,}6 \times 5 = 13$ giờ.
}
\end{ex}

\begin{ex} %[5D10H58]
Tính kết quả phép tính: $16$ giờ $15$ phút $: 5$ theo phút.
\shortans{195}
\loigiai{
\textbf{Đáp số:} $195$\\
$16 \times 60 + 15 = 975$ phút $\Rightarrow 975 : 5 = 195$ phút.
}
\end{ex}

\begin{ex} %[5D10V58]
Một máy đóng gói tự động đóng được $1$ gói bánh mất $4$ giây. Hỏi trong $15$ phút máy đó đóng được bao nhiêu gói bánh?
\shortans{225}
\loigiai{
\textbf{Đáp số:} $225$\\
$15$ phút $= 900$ giây $\Rightarrow 900 : 4 = 225$ gói bánh.
}
\end{ex}

\begin{ex} %[5D10V58]
Một người làm $3$ dụng cụ hết $5$ giờ $15$ phút. Hỏi làm $8$ dụng cụ như thế hết bao nhiêu giờ?
\shortans{14}
\loigiai{
\textbf{Đáp số:} $14$\\
1 dụng cụ: $5{,}25 : 3 = 1{,}75$ giờ $\Rightarrow 8$ dụng cụ: $1{,}75 \times 8 = 14$ giờ.
}
\end{ex}

\begin{ex} %[5D10V58]
Một chiếc xe máy chạy $3$ vòng trường đua hết $7$ phút $30$ giây. Hỏi xe máy đó chạy $10$ vòng như thế hết bao nhiêu phút?
\shortans{25}
\loigiai{
\textbf{Đáp số:} $25$\\
1 vòng: $7{,}5 : 3 = 2{,}5$ phút $\Rightarrow 10$ vòng: $2{,}5 \times 10 = 25$ phút.
}
\end{ex}

\begin{ex} %[5D10C58]
Một nhóm thợ gồm $4$ người cùng làm xong một công việc hết $6$ giờ $45$ phút. Nếu nhóm có $6$ người làm với cùng năng suất như nhau thì hoàn thành công việc đó hết bao nhiêu phút?
\shortans{270}
\loigiai{
\textbf{Đáp số:} $270$\\
Tổng thời gian 1 người cần: $405 \times 4 = 1\,620$ phút.\\
Thời gian 6 người làm: $1\,620 : 6 = 270$ phút.
}
\end{ex}

\newpage

\TaoPhieuKiemTra{Chủ đề 10 -- Số đo thời gian, vận tốc. các bài toán liên quan đến chuyển động đều}{59}{Vận tốc của một chuyển động đều}{https://loptoanthaylinh.vn/thi-online/toan-5/bai-59}
\textbf{Phần 1. Trắc nghiệm lựa chọn}
\setcounterpageref{ex}{0}

\begin{ex} %[5D10N59]
Công thức tính vận tốc $v$ khi biết quãng đường $s$ và thời gian $t$ là
\choice
{$v = s \times t$}
{$v = t : s$}
{\True $v = s : t$}
{$v = s + t$}
\loigiai{
Vận tốc bằng quãng đường chia cho thời gian: $v = s : t$.\\
\textbf{Đáp án của câu hỏi là C}
}
\end{ex}

\begin{ex} %[5D10N59]
Đơn vị nào sau đây thường dùng để đo vận tốc của ô tô, xe máy?
\choice
{$\text{m/s}$}
{$\text{cm/phút}$}
{$\text{m/giờ}$}
{\True $\text{km/h}$}
\loigiai{
Đơn vị thông dụng để đo vận tốc ô tô, xe máy là $\text{km/h}$ (ki-lô-mét trên giờ).\\
\textbf{Đáp án của câu hỏi là D}
}
\end{ex}

\begin{ex} %[5D10H59]
Một người đi xe máy đi được quãng đường $120\text{ km}$ trong thời gian $3\text{ giờ}$. Vận tốc của người đi xe máy là
\choice
{\True $40\text{ km/h}$}
{$30\text{ km/h}$}
{$50\text{ km/h}$}
{$60\text{ km/h}$}
\loigiai{
$v = 120 : 3 = 40\text{ (km/h)}$.\\
\textbf{Đáp án của câu hỏi là A}
}
\end{ex}

\begin{ex} %[5D10H59]
Một con đà điểu chạy được quãng đường $3\,150\text{ m}$ trong $3$ phút. Vận tốc chạy của đà điểu là
\choice
{$1\,000\text{ m/phút}$}
{$1\,100\text{ m/phút}$}
{\True $1\,050\text{ m/phút}$}
{$950\text{ m/phút}$}
\loigiai{
$v = 3\,150 : 3 = 1\,050\text{ (m/phút)}$.\\
\textbf{Đáp án của câu hỏi là C}
}
\end{ex}

\begin{ex} %[5D10H59]
Một vận động viên chạy cự li $100\text{ m}$ hết $10$ giây. Vận tốc chạy của vận động viên đó là
\choice
{$1\text{ m/s}$}
{\True $10\text{ m/s}$}
{$100\text{ m/s}$}
{$0{,}1\text{ m/s}$}
\loigiai{
$v = 100 : 10 = 10\text{ (m/s)}$.\\
\textbf{Đáp án của câu hỏi là B}
}
\end{ex}

\begin{ex} %[5D10H59]
Một người đi bộ đi được quãng đường $10{,}5\text{ km}$ trong $2{,}5\text{ giờ}$. Vận tốc của người đi bộ là
\choice
{$4\text{ km/h}$}
{$4{,}5\text{ km/h}$}
{$5\text{ km/h}$}
{\True $4{,}2\text{ km/h}$}
\loigiai{
$v = 10{,}5 : 2{,}5 = 4{,}2\text{ (km/h)}$.\\
\textbf{Đáp án của câu hỏi là D}
}
\end{ex}

\begin{ex} %[5D10V59]
Một ô tô đi quãng đường $135\text{ km}$ hết $2$ giờ $15$ phút. Vận tốc của ô tô là
\choice
{$54\text{ km/h}$}
{\True $60\text{ km/h}$}
{$65\text{ km/h}$}
{$70\text{ km/h}$}
\loigiai{
Đổi $2$ giờ $15$ phút $= 2{,}25$ giờ.\\
$v = 135 : 2{,}25 = 60\text{ (km/h)}$.\\
\textbf{Đáp án của câu hỏi là B}
}
\end{ex}

\begin{ex} %[5D10V59]
Một máy bay bay được $1\,800\text{ km}$ trong $2$ giờ $30$ phút. Vận tốc của máy bay là
\choice
{$600\text{ km/h}$}
{$800\text{ km/h}$}
{$900\text{ km/h}$}
{\True $720\text{ km/h}$}
\loigiai{
Đổi $2$ giờ $30$ phút $= 2{,}5$ giờ.\\
$v = 1\,800 : 2{,}5 = 720\text{ (km/h)}$.\\
\textbf{Đáp án của câu hỏi là D}
}
\end{ex}

\begin{ex} %[5D10V59]
Một người đi xe đạp đi quãng đường $18\text{ km}$ từ $7$ giờ $15$ phút đến $8$ giờ $45$ phút. Vận tốc của người đi xe đạp là
\choice
{\True $12\text{ km/h}$}
{$10\text{ km/h}$}
{$15\text{ km/h}$}
{$18\text{ km/h}$}
\loigiai{
Thời gian đi: $8$ giờ $45$ phút $- 7$ giờ $15$ phút $= 1$ giờ $30$ phút $= 1{,}5$ giờ.\\
Vận tốc: $18 : 1{,}5 = 12\text{ (km/h)}$.\\
\textbf{Đáp án của câu hỏi là A}
}
\end{ex}

\begin{ex} %[5D10V59]
Vận tốc $36\text{ km/h}$ tương đương với bao nhiêu mét trên giây ($	ext{m/s}$)?
\choice
{$1\text{ m/s}$}
{\True $10\text{ m/s}$}
{$100\text{ m/s}$}
{$36\text{ m/s}$}
\loigiai{
$36\text{ km/h} = 36\,000\text{ m} : 3\,600\text{ s} = 10\text{ m/s}$.\\
\textbf{Đáp án của câu hỏi là B}
}
\end{ex}

\begin{ex} %[5D10V59]
Một ca nô đi xuôi dòng quãng đường $45\text{ km}$ hết $1$ giờ $30$ phút. Vận tốc của dòng nước là $3\text{ km/h}$. Vận tốc thực của ca nô khi nước yên lặng là
\choice
{$33\text{ km/h}$}
{$30\text{ km/h}$}
{\True $27\text{ km/h}$}
{$24\text{ km/h}$}
\loigiai{
Vận tốc xuôi dòng: $45 : 1{,}5 = 30\text{ (km/h)}$.\\
Vận tốc thực: $30 - 3 = 27\text{ (km/h)}$.\\
\textbf{Đáp án của câu hỏi là C}
}
\end{ex}

\begin{ex} %[5D10C59]
Một ô tô đi từ A đến B gồm một đoạn đường bằng và một đoạn đường dốc. Đoạn đường bằng dài $60\text{ km}$ ô tô đi hết $1$ giờ $12$ phút. Đoạn đường dốc dài $24\text{ km}$ ô tô đi hết $36$ phút. Vận tốc trung bình của ô tô trên cả quãng đường là
\choice
{\True $46{,}67\text{ km/h}$ (hay $\dfrac{140}{3}\text{ km/h}$)}
{$45\text{ km/h}$}
{$48\text{ km/h}$}
{$50\text{ km/h}$}
\loigiai{
Tổng quãng đường: $60 + 24 = 84\text{ (km)}$.\\
Tổng thời gian: $1$ giờ $12$ phút $+ 36$ phút $= 1$ giờ $48$ phút $= 1{,}8$ giờ.\\
Vận tốc trung bình: $84 : 1{,}8 = 46{,}67\text{ (km/h)}$.\\
\textbf{Đáp án của câu hỏi là A}
}
\end{ex}

\textbf{Phần 2. Trắc nghiệm đúng sai}
\setcounterpageref{ex}{0}

\begin{ex} %[5D10H59]
So sánh vận tốc của các phương tiện sau:
- Xe đạp: đi $15\text{ km}$ trong $1$ giờ $15$ phút.
- Xe máy: đi $42\text{ km}$ trong $1$ giờ.
- Ô tô: đi $135\text{ km}$ trong $2{,}5$ giờ.
Xét tính đúng sai:
\choiceTF
{\True Vận tốc xe đạp là $12\text{ km/h}$.}
{\True Vận tốc ô tô là $54\text{ km/h}$.}
{Vận tốc xe máy lớn hơn vận tốc ô tô.}
{\True Thứ tự vận tốc tăng dần: Xe đạp $<$ Xe máy $<$ Ô tô.}
\loigiai{
\begin{itemize}
    \item Vận tốc xe đạp: $15 : 1{,}25 = 12\text{ km/h}$. Đúng.
    \item Vận tốc ô tô: $135 : 2{,}5 = 54\text{ km/h}$. Đúng.
    \item Vận tốc xe máy ($42\text{ km/h}$) bé hơn ô tô ($54\text{ km/h}$). Khẳng định lớn hơn là Sai.
    \item $12 < 42 < 54$. Đúng.
\end{itemize}
\begin{itemize}
    \item a) \textbf{Đúng}.
    \item b) \textbf{Đúng}.
    \item c) \textbf{Sai}.
    \item d) \textbf{Đúng}.
\end{itemize}
}
\end{ex}

\begin{ex} %[5D10V59]
Một người đi xe máy từ A đến B dài $90\text{ km}$. Người đó xuất phát lúc $7$ giờ $30$ phút, dọc đường nghỉ $15$ phút và đến B lúc $10$ giờ. Xét tính đúng sai:
\choiceTF
{\True Thời gian cả đi lẫn nghỉ là $2$ giờ $30$ phút.}
{\True Thời gian thực tế đi trên đường là $2$ giờ $15$ phút $= 2{,}25$ giờ.}
{Vận tốc của người đi xe máy là $36\text{ km/h}$.}
{\True Vận tốc của người đi xe máy là $40\text{ km/h}$.}
\loigiai{
\begin{itemize}
    \item Thời gian đi + nghỉ: $10$h $- 7$h $30$p $= 2$h $30$p. Đúng.
    \item Thời gian chạy xe: $2$h $30$p $- 15$p $= 2$h $15$p $= 2{,}25$h. Đúng.
    \item Vận tốc: $90 : 2{,}25 = 40\text{ km/h} \neq 36\text{ km/h}$. Ý 3 Sai, Ý 4 Đúng.
\end{itemize}
\begin{itemize}
    \item a) \textbf{Đúng}.
    \item b) \textbf{Đúng}.
    \item c) \textbf{Sai}.
    \item d) \textbf{Đúng}.
\end{itemize}
}
\end{ex}

\textbf{Phần 3. Trả lời ngắn}
\setcounterpageref{ex}{0}

\begin{ex} %[5D10H59]
Một người chạy $800\text{ m}$ hết $4$ phút. Tính vận tốc chạy theo mét trên phút.
\shortans{200}
\loigiai{
\textbf{Đáp số:} $200$\\
$800 : 4 = 200\text{ m/phút}$.
}
\end{ex}

\begin{ex} %[5D10H59]
Một ô tô đi được $180\text{ km}$ trong $4$ giờ. Tính vận tốc của ô tô theo $\text{km/h}$.
\shortans{45}
\loigiai{
\textbf{Đáp số:} $45$\\
$180 : 4 = 45\text{ km/h}$.
}
\end{ex}

\begin{ex} %[5D10V59]
Một tàu hỏa đi quãng đường $154\text{ km}$ hết $3$ giờ $30$ phút. Tính vận tốc tàu hỏa theo $\text{km/h}$.
\shortans{44}
\loigiai{
\textbf{Đáp số:} $44$\\
$154 : 3{,}5 = 44\text{ km/h}$.
}
\end{ex}

\begin{ex} %[5D10V59]
Một con ong mật bay quãng đường $12\text{ km}$ trong $40$ phút. Tính vận tốc bay của ong mật theo $\text{km/h}$.
\shortans{18}
\loigiai{
\textbf{Đáp số:} $18$\\
$40$ phút $= \dfrac{2}{3}$ giờ $\Rightarrow 12 : \dfrac{2}{3} = 18\text{ km/h}$.
}
\end{ex}

\begin{ex} %[5D10V59]
Đổi vận tốc $72\text{ km/h}$ ra mét trên giây ($	ext{m/s}$).
\shortans{20}
\loigiai{
\textbf{Đáp số:} $20$\\
$72 : 3{,}6 = 20\text{ m/s}$.
}
\end{ex}

\begin{ex} %[5D10C59]
Một người đi xe máy từ A lúc $8$ giờ $15$ phút và đến B lúc $10$ giờ $45$ phút với vận tốc $38\text{ km/h}$. Tính quãng đường AB theo ki-lô-mét.
\shortans{95}
\loigiai{
\textbf{Đáp số:} $95$\\
Thời gian đi: $2$ giờ $30$ phút $= 2{,}5$ giờ.\\
Quãng đường: $38 \times 2{,}5 = 95\text{ km}$.
}
\end{ex}

\newpage

\TaoPhieuKiemTra{Chủ đề 10 -- Số đo thời gian, vận tốc. các bài toán liên quan đến chuyển động đều}{60}{Quãng đường, thời gian của một chuyển động đều}{https://loptoanthaylinh.vn/thi-online/toan-5/bai-60}
\textbf{Phần 1. Trắc nghiệm lựa chọn}
\setcounterpageref{ex}{0}

\begin{ex} %[5D10N60]
Công thức tính quãng đường $s$ khi biết vận tốc $v$ và thời gian $t$ là
\choice
{$s = v : t$}
{$s = t : v$}
{$s = v + t$}
{\True $s = v \times t$}
\loigiai{
$s = v \times t$.\\
\textbf{Đáp án của câu hỏi là D}
}
\end{ex}

\begin{ex} %[5D10N60]
Công thức tính thời gian $t$ khi biết quãng đường $s$ và vận tốc $v$ là
\choice
{$t = s \times v$}
{\True $t = s : v$}
{$t = v : s$}
{$t = s - v$}
\loigiai{
$t = s : v$.\\
\textbf{Đáp án của câu hỏi là B}
}
\end{ex}

\begin{ex} %[5D10H60]
Một người đi xe máy với vận tốc $38\text{ km/h}$ trong $3$ giờ. Quãng đường người đó đi được là
\choice
{$104\text{ km}$}
{$124\text{ km}$}
{\True $114\text{ km}$}
{$110\text{ km}$}
\loigiai{
$s = 38 \times 3 = 114\text{ (km)}$.\\
\textbf{Đáp án của câu hỏi là C}
}
\end{ex}

\begin{ex} %[5D10H60]
Một ô tô đi quãng đường $150\text{ km}$ với vận tốc $50\text{ km/h}$. Thời gian ô tô đi hết quãng đường là
\choice
{\True $3\text{ giờ}$}
{$2\text{ giờ}$}
{$4\text{ giờ}$}
{$2{,}5\text{ giờ}$}
\loigiai{
$t = 150 : 50 = 3\text{ (giờ)}$.\\
\textbf{Đáp án của câu hỏi là A}
}
\end{ex}

\begin{ex} %[5D10H60]
Một người đi bộ với vận tốc $4{,}5\text{ km/h}$ trong $1$ giờ $20$ phút. Quãng đường người đó đi được là
\choice
{$5\text{ km}$}
{$6{,}5\text{ km}$}
{$7\text{ km}$}
{\True $6\text{ km}$}
\loigiai{
$1$ giờ $20$ phút $= \dfrac{4}{3}$ giờ.\\
$s = 4{,}5 \times \dfrac{4}{3} = 6\text{ (km)}$.\\
\textbf{Đáp án của câu hỏi là D}
}
\end{ex}

\begin{ex} %[5D10H60]
Một người đi xe đạp quãng đường $27\text{ km}$ với vận tốc $12\text{ km/h}$. Thời gian người đó đi là
\choice
{$2$ giờ $10$ phút}
{$2$ giờ $20$ phút}
{\True $2$ giờ $15$ phút}
{$2$ giờ $30$ phút}
\loigiai{
$t = 27 : 12 = 2{,}25\text{ giờ} = 2$ giờ $15$ phút.\\
\textbf{Đáp án của câu hỏi là C}
}
\end{ex}

\begin{ex} %[5D10V60]
Hai thành phố A và B cách nhau $160\text{ km}$. Hai ô tô khởi hành cùng một lúc từ A và B đi ngược chiều nhau. Xe thứ nhất đi với vận tốc $45\text{ km/h}$, xe thứ hai đi với vận tốc $35\text{ km/h}$. Sau bao lâu hai xe gặp nhau?
\choice
{$1{,}5\text{ giờ}$}
{$2{,}5\text{ giờ}$}
{$3\text{ giờ}$}
{\True $2\text{ giờ}$}
\loigiai{
Tổng vận tốc hai xe: $45 + 35 = 80\text{ (km/h)}$.\\
Thời gian gặp nhau: $160 : 80 = 2\text{ (giờ)}$.\\
\textbf{Đáp án của câu hỏi là D}
}
\end{ex}

\begin{ex} %[5D10V60]
Một người đi xe máy khởi hành từ A lúc $7$ giờ với vận tốc $36\text{ km/h}$. Đến $8$ giờ, một ô tô cũng xuất phát từ A đuổi theo xe máy với vận tốc $54\text{ km/h}$. Ô tô đuổi kịp xe máy lúc mấy giờ?
\choice
{$9$ giờ}
{\True $10$ giờ}
{$10$ giờ $30$ phút}
{$11$ giờ}
\loigiai{
Trong 1 giờ (từ 7h đến 8h), xe máy đi được: $36 \times 1 = 36\text{ (km)}$.\\
Hiệu vận tốc hai xe: $54 - 36 = 18\text{ (km/h)}$.\\
Thời gian ô tô đuổi kịp: $36 : 18 = 2\text{ (giờ)}$.\\
Thời điểm gặp nhau: $8 + 2 = 10\text{ (giờ)}$.\\
\textbf{Đáp án của câu hỏi là B}
}
\end{ex}

\begin{ex} %[5D10V60]
Một ca nô đi xuôi dòng từ bến A đến bến B hết $2$ giờ, đi ngược dòng từ B về A hết $3$ giờ. Biết vận tốc dòng nước là $3\text{ km/h}$. Quãng đường sông từ A đến B là
\choice
{\True $36\text{ km}$}
{$30\text{ km}$}
{$40\text{ km}$}
{$45\text{ km}$}
\loigiai{
Hiệu vận tốc xuôi dòng và ngược dòng bằng 2 lần vận tốc dòng nước: $3 \times 2 = 6\text{ (km/h)}$.\\
Tỉ số thời gian xuôi và ngược là $\dfrac{2}{3} \Rightarrow$ Tỉ số vận tốc xuôi và ngược là $\dfrac{3}{2}$.\\
Vận tốc xuôi dòng: $6 : (3 - 2) \times 3 = 18\text{ (km/h)}$.\\
Quãng đường AB: $18 \times 2 = 36\text{ (km)}$.\\
\textbf{Đáp án của câu hỏi là A}
}
\end{ex}

\begin{ex} %[5D10V60]
Một người đi xe đạp từ A đến B với vận tốc $12\text{ km/h}$. Sau đó đi ô tô từ B về A với vận tốc $48\text{ km/h}$. Tổng thời gian cả đi và về là $2$ giờ $30$ phút. Quãng đường AB là
\choice
{$20\text{ km}$}
{$28\text{ km}$}
{\True $24\text{ km}$}
{$30\text{ km}$}
\loigiai{
Tỉ số vận tốc đi và về: $12 : 48 = \dfrac{1}{4} \Rightarrow$ Tỉ số thời gian đi và về là $\dfrac{4}{1}$.\\
Đổi $2$ giờ $30$ phút $= 2{,}5$ giờ.\\
Thời gian đi xe đạp: $2{,}5 : (4 + 1) \times 4 = 2$ giờ.\\
Quãng đường AB: $12 \times 2 = 24\text{ (km)}$.\\
\textbf{Đáp án của câu hỏi là C}
}
\end{ex}

\begin{ex} %[5D10V60]
Một đoàn tàu hỏa dài $150\text{ m}$ chạy qua một chiếc cầu dài $450\text{ m}$ với vận tốc $36\text{ km/h}$. Thời gian để đoàn tàu chạy hoàn toàn qua cây cầu là
\choice
{\True $60\text{ giây}$}
{$45\text{ giây}$}
{$50\text{ giây}$}
{$75\text{ giây}$}
\loigiai{
Quãng đường tàu phải chạy: $450 + 150 = 600\text{ (m)}$.\\
Vận tốc $36\text{ km/h} = 10\text{ m/s}$.\\
Thời gian tàu qua cầu: $600 : 10 = 60\text{ (giây)}$.\\
\textbf{Đáp án của câu hỏi là A}
}
\end{ex}

\begin{ex} %[5D10C60]
Hai người khởi hành cùng lúc từ hai địa điểm A và B cách nhau $72\text{ km}$ và đi ngược chiều nhau. Người thứ nhất đi từ A bằng xe đạp với vận tốc $14\text{ km/h}$, người thứ hai đi từ B bằng xe máy với vận tốc $34\text{ km/h}$. Hỏi sau bao lâu thì hai người cách nhau $24\text{ km}$ trước khi gặp nhau?
\choice
{$45\text{ phút}$}
{\True $1\text{ giờ}$}
{$1{,}2\text{ giờ}$}
{$1{,}5\text{ giờ}$}
\loigiai{
Tổng quãng đường hai người đã đi được: $72 - 24 = 48\text{ (km)}$.\\
Tổng vận tốc: $14 + 34 = 48\text{ (km/h)}$.\\
Thời gian: $48 : 48 = 1\text{ (giờ)}$.\\
\textbf{Đáp án của câu hỏi là B}
}
\end{ex}

\textbf{Phần 2. Trắc nghiệm đúng sai}
\setcounterpageref{ex}{0}

\begin{ex} %[5D10H60]
Quãng đường AB dài $120\text{ km}$. Một ô tô đi từ A đến B với vận tốc $50\text{ km/h}$. Xét tính đúng sai:
\choiceTF
{\True Thời gian ô tô đi hết quãng đường là $2{,}4$ giờ.}
{Đổi $2{,}4$ giờ $= 2$ giờ $40$ phút.}
{\True Đổi $2{,}4$ giờ $= 2$ giờ $24$ phút.}
{\True Nếu xuất phát lúc $8$ giờ thì ô tô đến B lúc $10$ giờ $24$ phút.}
\loigiai{
\begin{itemize}
    \item $t = 120 : 50 = 2{,}4$ giờ. Đúng.
    \item $2{,}4$ giờ $= 2$ giờ $+ 0{,}4 \times 60$ phút $= 2$ giờ $24$ phút $\neq 2$ giờ $40$ phút. Ý 2 Sai, Ý 3 Đúng.
    \item Đến nơi lúc $8$h $+ 2$h $24$p $= 10$h $24$p. Đúng.
\end{itemize}
\begin{itemize}
    \item a) \textbf{Đúng}.
    \item b) \textbf{Sai}.
    \item c) \textbf{Đúng}.
    \item d) \textbf{Đúng}.
\end{itemize}
}
\end{ex}

\begin{ex} %[5D10V60]
Hai thành phố A và B cách nhau $135\text{ km}$. Một xe máy đi từ A đến B với vận tốc $30\text{ km/h}$, cùng lúc đó một ô tô đi từ B về A với vận tốc $60\text{ km/h}$. Xét tính đúng sai:
\choiceTF
{\True Sau mỗi giờ, hai xe tiến lại gần nhau thêm $90\text{ km}$.}
{\True Hai xe gặp nhau sau $1{,}5$ giờ ($1$ giờ $30$ phút).}
{Nơi hai xe gặp nhau cách A là $60\text{ km}$.}
{\True Nơi hai xe gặp nhau cách A là $45\text{ km}$.}
\loigiai{
\begin{itemize}
    \item Tổng vận tốc: $30 + 60 = 90\text{ km/h}$. Đúng.
    \item Thời gian gặp: $135 : 90 = 1{,}5$ giờ. Đúng.
    \item Điểm gặp cách A: $30 \times 1{,}5 = 45\text{ km} \neq 60\text{ km}$. Ý 3 Sai, Ý 4 Đúng.
\end{itemize}
\begin{itemize}
    \item a) \textbf{Đúng}.
    \item b) \textbf{Đúng}.
    \item c) \textbf{Sai}.
    \item d) \textbf{Đúng}.
\end{itemize}
}
\end{ex}

\textbf{Phần 3. Trả lời ngắn}
\setcounterpageref{ex}{0}

\begin{ex} %[5D10H60]
Một người đi xe máy với vận tốc $42\text{ km/h}$ trong $2{,}5$ giờ. Tính quãng đường đi được theo ki-lô-mét.
\shortans{105}
\loigiai{
\textbf{Đáp số:} $105$\\
$42 \times 2{,}5 = 105\text{ km}$.
}
\end{ex}

\begin{ex} %[5D10H60]
Một ô tô đi quãng đường $168\text{ km}$ với vận tốc $48\text{ km/h}$. Tính thời gian ô tô đi theo giờ (dùng dấu phẩy nếu có).
\shortans{3,5}
\loigiai{
\textbf{Đáp số:} $3{,}5$\\
$168 : 48 = 3{,}5\text{ giờ}$.
}
\end{ex}

\begin{ex} %[5D10V60]
Một người đi bộ với vận tốc $5\text{ km/h}$ từ nhà đến trường hết $24$ phút. Tính quãng đường từ nhà đến trường theo ki-lô-mét.
\shortans{2}
\loigiai{
\textbf{Đáp số:} $2$\\
$24$ phút $= 0{,}4$ giờ $\Rightarrow 5 \times 0{,}4 = 2\text{ km}$.
}
\end{ex}

\begin{ex} %[5D10V60]
Hai ô tô khởi hành cùng lúc từ hai điểm cách nhau $180\text{ km}$ đi ngược chiều nhau với vận tốc $48\text{ km/h}$ và $42\text{ km/h}$. Hỏi sau bao nhiêu giờ hai xe gặp nhau?
\shortans{2}
\loigiai{
\textbf{Đáp số:} $2$\\
$180 : (48 + 42) = 180 : 90 = 2\text{ giờ}$.
}
\end{ex}

\begin{ex} %[5D10V60]
Một ca nô đi xuôi dòng một khúc sông với vận tốc $24\text{ km/h}$ hết $2$ giờ $15$ phút. Tính độ dài khúc sông đó theo ki-lô-mét.
\shortans{54}
\loigiai{
\textbf{Đáp số:} $54$\\
$24 \times 2{,}25 = 54\text{ km}$.
}
\end{ex}

\begin{ex} %[5D10C60]
Một người đi xe máy từ A đến B với vận tốc $35\text{ km/h}$. Sau $1$ giờ, một ô tô cũng đi từ A đuổi theo với vận tốc $50\text{ km/h}$. Hỏi sau bao nhiêu phút ô tô đuổi kịp xe máy?
\shortans{140}
\loigiai{
\textbf{Đáp số:} $140$\\
Khoảng cách ban đầu: $35\text{ km}$.\\
Thời gian đuổi kịp: $35 : (50 - 35) = \dfrac{35}{15} = \dfrac{7}{3}\text{ giờ} = 140\text{ phút}$.
}
\end{ex}

\newpage

\TaoPhieuKiemTra{Chủ đề 10 -- Số đo thời gian, vận tốc. các bài toán liên quan đến chuyển động đều}{61}{Thực hành tính toán, ước lượng vận tốc, quãng đường, thời gian}{https://loptoanthaylinh.vn/thi-online/toan-5/bai-61}
\textbf{Phần 1. Trắc nghiệm lựa chọn}
\setcounterpageref{ex}{0}

\begin{ex} %[5D10N61]
Vận tốc đi bộ bình thường của một học sinh tiểu học khoảng
\choice
{\True $4\text{ km/h}$}
{$15\text{ km/h}$}
{$40\text{ km/h}$}
{$1\text{ km/h}$}
\loigiai{
Vận tốc đi bộ trung bình khoảng $4\text{ km/h}$.\\
\textbf{Đáp án của câu hỏi là A}
}
\end{ex}

\begin{ex} %[5D10N61]
Vận tốc trung bình của một chiếc xe buýt chạy trong nội thành thường khoảng
\choice
{$80\text{ km/h}$}
{$5\text{ km/h}$}
{\True $25\text{ km/h}$}
{$120\text{ km/h}$}
\loigiai{
Xe buýt trong nội thành có tốc độ trung bình khoảng $20 - 30\text{ km/h}$ do phải dừng đón trả khách.\\
\textbf{Đáp án của câu hỏi là C}
}
\end{ex}

\begin{ex} %[5D10H61]
Nhà bạn Nam cách trường $1{,}8\text{ km}$. Mỗi ngày Nam đi xe đạp với vận tốc $12\text{ km/h}$. Thời gian Nam đi từ nhà đến trường là
\choice
{$6$ phút}
{$12$ phút}
{$15$ phút}
{\True $9$ phút}
\loigiai{
$t = 1{,}8 : 12 = 0{,}15\text{ giờ} = 9$ phút.\\
\textbf{Đáp án của câu hỏi là D}
}
\end{ex}

\begin{ex} %[5D10H61]
Một gia đình dự định đi ô tô từ Hà Nội về quê cách đó $100\text{ km}$ với vận tốc ước tính $50\text{ km/h}$. Nếu muốn có mặt ở quê lúc $10$ giờ sáng thì gia đình phải xuất phát muộn nhất lúc mấy giờ?
\choice
{\True $8$ giờ}
{$7$ giờ}
{$8$ giờ $30$ phút}
{$9$ giờ}
\loigiai{
Thời gian đi: $100 : 50 = 2$ giờ.\\
Thời điểm xuất phát muộn nhất: $10 - 2 = 8$ giờ.\\
\textbf{Đáp án của câu hỏi là A}
}
\end{ex}

\begin{ex} %[5D10H61]
Bạn Mai tập chạy cự li $400\text{ m}$. Đồng hồ bấm giờ chỉ $1$ phút $20$ giây. Vận tốc chạy của Mai là
\choice
{$4\text{ m/s}$}
{$6\text{ m/s}$}
{\True $5\text{ m/s}$}
{$8\text{ m/s}$}
\loigiai{
$1$ phút $20$ giây $= 80$ giây.\\
$v = 400 : 80 = 5\text{ (m/s)}$.\\
\textbf{Đáp án của câu hỏi là C}
}
\end{ex}

\begin{ex} %[5D10H61]
Khoảng cách giữa hai cột mốc ven đường cao tốc là $100\text{ m}$. Một ô tô chạy qua hai cột mốc này mất đúng $4$ giây. Vận tốc của ô tô là
\choice
{$60\text{ km/h}$}
{\True $90\text{ km/h}$}
{$80\text{ km/h}$}
{$100\text{ km/h}$}
\loigiai{
Vận tốc: $100 : 4 = 25\text{ m/s} = 25 \times 3{,}6 = 90\text{ km/h}$.\\
\textbf{Đáp án của câu hỏi là B}
}
\end{ex}

\begin{ex} %[5D10V61]
Một người tập thể dục đi bộ một vòng quanh bờ hồ hình tròn có bán kính $200\text{ m}$ với vận tốc $4\text{ km/h}$. Thời gian người đó đi hết $1$ vòng (lấy $\pi = 3{,}14$) xấp xỉ là
\choice
{$12$ phút}
{$15$ phút}
{$24$ phút}
{\True $18{,}84$ phút}
\loigiai{
Chu vi bờ hồ: $2 \times 3{,}14 \times 200 = 1\,256\text{ (m)} = 1{,}256\text{ km}$.\\
Thời gian đi: $1{,}256 : 4 = 0{,}314\text{ giờ} = 18{,}84$ phút.\\
\textbf{Đáp án của câu hỏi là D}
}
\end{ex}

\begin{ex} %[5D10V61]
Hằng ngày bố chở An đi học bằng xe máy với vận tốc $30\text{ km/h}$ hết $16$ phút. Hôm nay do trời mưa, bố chỉ đi được với vận tốc $24\text{ km/h}$. Thời gian An đến trường hôm nay là
\choice
{$18$ phút}
{\True $20$ phút}
{$22$ phút}
{$25$ phút}
\loigiai{
Quãng đường từ nhà đến trường: $30 \times \dfrac{16}{60} = 8\text{ (km)}$.\\
Thời gian khi đi $24\text{ km/h}$: $8 : 24 = \dfrac{1}{3}\text{ giờ} = 20$ phút.\\
\textbf{Đáp án của câu hỏi là B}
}
\end{ex}

\begin{ex} %[5D10V61]
Một tàu cao tốc đi quãng đường $300\text{ km}$ mất $1$ giờ $40$ phút. Vận tốc trung bình của tàu cao tốc là
\choice
{\True $180\text{ km/h}$}
{$150\text{ km/h}$}
{$200\text{ km/h}$}
{$240\text{ km/h}$}
\loigiai{
$1$ giờ $40$ phút $= \dfrac{5}{3}$ giờ.\\
$v = 300 : \dfrac{5}{3} = 180\text{ (km/h)}$.\\
\textbf{Đáp án của câu hỏi là A}
}
\end{ex}

\begin{ex} %[5D10V61]
Một bạn học sinh ước lượng thời gian đi bộ từ tầng 1 lên tầng 4 của trường học (mỗi tầng cao $3{,}6\text{ m}$, có $20$ bậc thang mỗi tầng). Biết mỗi giây bạn bước được $2$ bậc thang. Thời gian bạn bước lên tầng 4 là
\choice
{$20$ giây}
{$40$ giây}
{$60$ giây}
{\True $30$ giây}
\loigiai{
Từ tầng 1 lên tầng 4 có $3$ nhịp cầu thang: $3 \times 20 = 60$ bậc thang.\\
Thời gian bước: $60 : 2 = 30$ giây.\\
\textbf{Đáp án của câu hỏi là D}
}
\end{ex}

\begin{ex} %[5D10V61]
Một đàn nhạn bay về phương nam tránh rét với vận tốc $45\text{ km/h}$. Mỗi ngày đàn nhạn bay liên tục trong $8$ giờ. Sau $5$ ngày đàn nhạn bay được quãng đường là
\choice
{$1\,600\text{ km}$}
{\True $1\,800\text{ km}$}
{$2\,000\text{ km}$}
{$2\,250\text{ km}$}
\loigiai{
Quãng đường mỗi ngày: $45 \times 8 = 360\text{ (km)}$.\\
Quãng đường 5 ngày: $360 \times 5 = 1\,800\text{ (km)}$.\\
\textbf{Đáp án của câu hỏi là B}
}
\end{ex}

\begin{ex} %[5D10C61]
Một vận động viên leo núi dự định leo lên đỉnh núi cao $1\,200\text{ m}$. Vận tốc khi leo lên là $200\text{ m/giờ}$, vận tốc khi xuống núi là $400\text{ m/giờ}$. Người đó nghỉ trên đỉnh núi $1$ giờ $30$ phút. Tổng thời gian của cả chuyến leo núi là
\choice
{$9$ giờ}
{$10$ giờ}
{\True $10$ giờ $30$ phút}
{$11$ giờ}
\loigiai{
Thời gian leo lên: $1\,200 : 200 = 6$ giờ.\\
Thời gian leo xuống: $1\,200 : 400 = 3$ giờ.\\
Thời gian nghỉ: $1$ giờ $30$ phút.\\
Tổng thời gian: $6 + 3 + 1$ giờ $30$ phút $= 10$ giờ $30$ phút.\\
\textbf{Đáp án của câu hỏi là C}
}
\end{ex}

\textbf{Phần 2. Trắc nghiệm đúng sai}
\setcounterpageref{ex}{0}

\begin{ex} %[5D10H61]
Xét tính đúng sai về ước lượng vận tốc trong thực tế:
\choiceTF
{\True Vận tốc của máy bay chở khách thông thường khoảng $800 - 900\text{ km/h}$.}
{Vận tốc người chạy bộ thông thường có thể đạt $100\text{ km/h}$.}
{\True Một người đi xe đạp điện với vận tốc $25\text{ km/h}$ đi $5\text{ km}$ hết $12$ phút.}
{\True Vận tốc âm thanh trong không khí xấp xỉ $340\text{ m/s}$.}
\loigiai{
\begin{itemize}
    \item Máy bay chở khách: 800 - 900 km/h. Đúng.
    \item Người chạy bộ không thể đạt 100 km/h (tối đa chỉ khoảng 35-40 km/h đối với VĐV thế giới). Sai.
    \item $5 : 25 = 0{,}2\text{ giờ} = 12$ phút. Đúng.
    \item Vận tốc âm thanh: ~340 m/s. Đúng.
\end{itemize}
\begin{itemize}
    \item a) \textbf{Đúng}.
    \item b) \textbf{Sai}.
    \item c) \textbf{Đúng}.
    \item d) \textbf{Đúng}.
\end{itemize}
}
\end{ex}

\begin{ex} %[5D10V61]
Một gia đình lái xe ô tô đi nghỉ mát từ tỉnh A đến tỉnh B cách nhau $240\text{ km}$. Họ dự định lái xe với vận tốc $60\text{ km/h}$ và nghỉ dọc đường $45$ phút. Xét tính đúng sai:
\choiceTF
{\True Thời gian lái xe thực tế là $4$ giờ.}
{\True Tổng thời gian cả đi và nghỉ là $4$ giờ $45$ phút.}
{Nếu xuất phát lúc $6$ giờ $30$ phút thì đến nơi lúc $11$ giờ.}
{\True Nếu xuất phát lúc $6$ giờ $30$ phút thì đến nơi lúc $11$ giờ $15$ phút.}
\loigiai{
\begin{itemize}
    \item Thời gian lái xe: $240 : 60 = 4$ giờ. Đúng.
    \item Tổng thời gian: $4$h $+ 45$p $= 4$h $45$p. Đúng.
    \item Đến nơi: $6$h $30$p $+ 4$h $45$p $= 11$h $15$p $\neq 11$h. Ý 3 Sai, Ý 4 Đúng.
\end{itemize}
\begin{itemize}
    \item a) \textbf{Đúng}.
    \item b) \textbf{Đúng}.
    \item c) \textbf{Sai}.
    \item d) \textbf{Đúng}.
\end{itemize}
}
\end{ex}

\textbf{Phần 3. Trả lời ngắn}
\setcounterpageref{ex}{0}

\begin{ex} %[5D10H61]
Nhà Minh cách trường $2{,}4\text{ km}$. Minh đi xe đạp hết $12$ phút. Tính vận tốc của Minh theo $\text{km/h}$.
\shortans{12}
\loigiai{
\textbf{Đáp số:} $12$\\
$12$ phút $= 0{,}2$ giờ $\Rightarrow 2{,}4 : 0{,}2 = 12\text{ km/h}$.
}
\end{ex}

\begin{ex} %[5D10H61]
Một người đi xe máy trong $45$ phút với vận tốc $40\text{ km/h}$. Tính quãng đường người đó đi được theo ki-lô-mét.
\shortans{30}
\loigiai{
\textbf{Đáp số:} $30$\\
$40 \times 0{,}75 = 30\text{ km}$.
}
\end{ex}

\begin{ex} %[5D10V61]
Một vận động viên bơi lội bơi cự li $200\text{ m}$ hết $2$ phút $30$ giây. Tính vận tốc bơi theo mét trên giây (dùng dấu phẩy nếu có).
\shortans{1,33 hoặc \dfrac{4}{3}}
\loigiai{
\textbf{Đáp số:} $1{,}33$ hoặc $\dfrac{4}{3}$\\
$2$ phút $30$ giây $= 150$ giây $\Rightarrow 200 : 150 = 1{,}33\text{ m/s}$.
}
\end{ex}

\begin{ex} %[5D10V61]
Một người đi xe đạp với vận tốc $15\text{ km/h}$. Người đó muốn đi quãng đường $22{,}5\text{ km}$ thì mất bao nhiêu phút?
\shortans{90}
\loigiai{
\textbf{Đáp số:} $90$\\
$22{,}5 : 15 = 1{,}5\text{ giờ} = 90\text{ phút}$.
}
\end{ex}

\begin{ex} %[5D10V61]
Một tàu thủy chạy trên biển với vận tốc $18$ hải lí/giờ ($1$ hải lí $\approx 1{,}852\text{ km}$). Tính quãng đường tàu chạy trong $3$ giờ theo ki-lô-mét.
\shortans{100,008}
\loigiai{
\textbf{Đáp số:} $100{,}008$\\
$18 \times 3 \times 1{,}852 = 54 \times 1{,}852 = 100{,}008\text{ km}$.
}
\end{ex}

\begin{ex} %[5D10C61]
Một nhóm học sinh đi dã ngoại. Chặng 1 đi xe đạp dài $12\text{ km}$ với vận tốc $12\text{ km/h}$. Chặng 2 đi bộ dài $4\text{ km}$ với vận tốc $4\text{ km/h}$. Nghỉ giữa hai chặng $30$ phút. Tính tổng thời gian chuyến đi dã ngoại theo phút.
\shortans{150}
\loigiai{
\textbf{Đáp số:} $150$\\
Chặng 1: $1$ giờ $= 60$ phút.\\
Chặng 2: $1$ giờ $= 60$ phút.\\
Nghỉ: $30$ phút.\\
Tổng: $60 + 60 + 30 = 150$ phút.
}
\end{ex}

\newpage

\TaoPhieuKiemTra{Chủ đề 10 -- Số đo thời gian, vận tốc. các bài toán liên quan đến chuyển động đều}{62}{Luyện tập chung toán chuyển động đều}{https://loptoanthaylinh.vn/thi-online/toan-5/bai-62}
\textbf{Phần 1. Trắc nghiệm lựa chọn}
\setcounterpageref{ex}{0}

\begin{ex} %[5D10N62]
Muốn tính vận tốc của một chuyển động đều, ta lấy
\choice
{Quãng đường nhân với thời gian}
{\True Quãng đường chia cho thời gian}
{Thời gian chia cho quãng đường}
{Quãng đường cộng với thời gian}
\loigiai{
$v = s : t$.\\
\textbf{Đáp án của câu hỏi là B}
}
\end{ex}

\begin{ex} %[5D10N62]
$3$ giờ $25$ phút $+ 2$ giờ $45$ phút bằng
\choice
{$5$ giờ $70$ phút}
{$6$ giờ $15$ phút}
{\True $6$ giờ $10$ phút}
{$5$ giờ $10$ phút}
\loigiai{
$3$ giờ $25$ phút $+ 2$ giờ $45$ phút $= 5$ giờ $70$ phút $= 6$ giờ $10$ phút.\\
\textbf{Đáp án của câu hỏi là C}
}
\end{ex}

\begin{ex} %[5D10H62]
Một người đi xe máy đi được $105\text{ km}$ trong $2$ giờ $30$ phút. Vận tốc của người đó là
\choice
{$40\text{ km/h}$}
{\True $42\text{ km/h}$}
{$45\text{ km/h}$}
{$50\text{ km/h}$}
\loigiai{
$105 : 2{,}5 = 42\text{ (km/h)}$.\\
\textbf{Đáp án của câu hỏi là B}
}
\end{ex}

\begin{ex} %[5D10H62]
Một ô tô đi với vận tốc $54\text{ km/h}$ trong $1$ giờ $40$ phút. Quãng đường ô tô đi được là
\choice
{$80\text{ km}$}
{$95\text{ km}$}
{$100\text{ km}$}
{\True $90\text{ km}$}
\loigiai{
$1$ giờ $40$ phút $= \dfrac{5}{3}$ giờ.\\
$s = 54 \times \dfrac{5}{3} = 90\text{ (km)}$.\\
\textbf{Đáp án của câu hỏi là D}
}
\end{ex}

\begin{ex} %[5D10H62]
Hai tỉnh A và B cách nhau $140\text{ km}$. Hai ô tô đi ngược chiều nhau cùng khởi hành và gặp nhau sau $1$ giờ $45$ phút. Tổng vận tốc của hai xe là
\choice
{\True $80\text{ km/h}$}
{$70\text{ km/h}$}
{$85\text{ km/h}$}
{$90\text{ km/h}$}
\loigiai{
$1$ giờ $45$ phút $= 1{,}75$ giờ.\\
Tổng vận tốc: $140 : 1{,}75 = 80\text{ (km/h)}$.\\
\textbf{Đáp án của câu hỏi là A}
}
\end{ex}

\begin{ex} %[5D10H62]
Một máy bay bay với vận tốc $850\text{ km/h}$ trên quãng đường $2\,125\text{ km}$. Thời gian máy bay bay hết quãng đường là
\choice
{$2$ giờ $15$ phút}
{$2$ giờ $45$ phút}
{\True $2$ giờ $30$ phút}
{$3$ giờ}
\loigiai{
$t = 2\,125 : 850 = 2{,}5\text{ giờ} = 2$ giờ $30$ phút.\\
\textbf{Đáp án của câu hỏi là C}
}
\end{ex}

\begin{ex} %[5D10V62]
Một người đi xe máy từ A đến B với vận tốc $36\text{ km/h}$. Sau khi xe máy đi được $1$ giờ $15$ phút thì một ô tô cũng khởi hành từ A đuổi theo với vận tốc $54\text{ km/h}$. Thời gian ô tô đuổi kịp xe máy là
\choice
{$2$ giờ}
{$2$ giờ $15$ phút}
{$3$ giờ}
{\True $2$ giờ $30$ phút}
\loigiai{
$1$ giờ $15$ phút $= 1{,}25$ giờ.\\
Khoảng cách xe máy đã đi: $36 \times 1{,}25 = 45\text{ (km)}$.\\
Hiệu vận tốc: $54 - 36 = 18\text{ (km/h)}$.\\
Thời gian đuổi kịp: $45 : 18 = 2{,}5\text{ giờ} = 2$ giờ $30$ phút.\\
\textbf{Đáp án của câu hỏi là D}
}
\end{ex}

\begin{ex} %[5D10V62]
Một ca nô đi xuôi dòng trên khúc sông từ A đến B dài $60\text{ km}$ hết $2$ giờ. Biết vận tốc dòng nước chảy là $4\text{ km/h}$. Thời gian ca nô đi ngược dòng từ B về A là
\choice
{$2{,}5\text{ giờ}$}
{$3$ giờ}
{\True $2$ giờ $43$ phút (hay $\dfrac{30}{11}$ giờ)}
{$2$ giờ $30$ phút}
\loigiai{
Vận tốc xuôi dòng: $60 : 2 = 30\text{ km/h}$.\\
Vận tốc ngược dòng: $30 - 4 \times 2 = 22\text{ km/h}$.\\
Thời gian ngược dòng: $60 : 22 = \dfrac{30}{11}\text{ giờ} \approx 2$ giờ $43$ phút.\\
\textbf{Đáp án của câu hỏi là C}
}
\end{ex}

\begin{ex} %[5D10V62]
Một người đi từ A đến B với vận tốc $40\text{ km/h}$, khi về đi với vận tốc $50\text{ km/h}$. Thời gian về ít hơn thời gian đi là $30$ phút. Quãng đường AB dài là
\choice
{\True $100\text{ km}$}
{$80\text{ km}$}
{$120\text{ km}$}
{$150\text{ km}$}
\loigiai{
Tỉ số vận tốc đi và về là $\dfrac{40}{50} = \dfrac{4}{5} \Rightarrow$ Tỉ số thời gian đi và về là $\dfrac{5}{4}$.\\
Hiệu thời gian là $30$ phút $= 0{,}5$ giờ.\\
Thời gian đi: $0{,}5 : (5 - 4) \times 5 = 2{,}5$ giờ.\\
Quãng đường AB: $40 \times 2{,}5 = 100\text{ (km)}$.\\
\textbf{Đáp án của câu hỏi là A}
}
\end{ex}

\begin{ex} %[5D10V62]
Một con tàu hỏa chạy qua một cây cầu dài $360\text{ m}$ mất $30$ giây và vượt qua một cột điện ven đường mất $12$ giây. Vận tốc của tàu hỏa là
\choice
{$15\text{ m/s}$}
{\True $20\text{ m/s}$}
{$25\text{ m/s}$}
{$30\text{ m/s}$}
\loigiai{
Thời gian tàu đi hết chiều dài cây cầu $360\text{ m}$ là: $30 - 12 = 18$ giây.\\
Vận tốc của tàu: $360 : 18 = 20\text{ (m/s)}$.\\
\textbf{Đáp án của câu hỏi là B}
}
\end{ex}

\begin{ex} %[5D10V62]
Một người đi bộ một vòng quanh hồ với vận tốc $4{,}5\text{ km/h}$ hết $36$ phút. Nếu người đó đi xe đạp với vận tốc $13{,}5\text{ km/h}$ thì đi hết một vòng hồ đó trong bao lâu?
\choice
{$10$ phút}
{$15$ phút}
{$18$ phút}
{\True $12$ phút}
\loigiai{
Vận tốc xe đạp gấp vận tốc đi bộ: $13{,}5 : 4{,}5 = 3$ lần.\\
Thời gian đi xe đạp: $36 : 3 = 12$ phút.\\
\textbf{Đáp án của câu hỏi là D}
}
\end{ex}

\begin{ex} %[5D10C62]
Hai xe cùng xuất phát từ A đến B. Xe thứ nhất đi với vận tốc $45\text{ km/h}$, xe thứ hai đi với vận tốc $60\text{ km/h}$. Xe thứ nhất đến B sau xe thứ hai $40$ phút. Tính quãng đường AB.
\choice
{\True $120\text{ km}$}
{$100\text{ km}$}
{$140\text{ km}$}
{$150\text{ km}$}
\loigiai{
Tỉ số vận tốc xe 1 và xe 2 là $\dfrac{45}{60} = \dfrac{3}{4} \Rightarrow$ Tỉ số thời gian là $\dfrac{4}{3}$.\\
Hiệu thời gian: $40$ phút $= \dfrac{2}{3}$ giờ.\\
Thời gian xe 1 đi: $\dfrac{2}{3} : (4 - 3) \times 4 = \dfrac{8}{3}$ giờ.\\
Quãng đường AB: $45 \times \dfrac{8}{3} = 120\text{ (km)}$.\\
\textbf{Đáp án của câu hỏi là A}
}
\end{ex}

\textbf{Phần 2. Trắc nghiệm đúng sai}
\setcounterpageref{ex}{0}

\begin{ex} %[5D10H62]
Hai địa điểm A và B cách nhau $120\text{ km}$. Xe máy đi từ A với vận tốc $35\text{ km/h}$, ô tô đi từ B với vận tốc $55\text{ km/h}$, hai xe cùng khởi hành lúc $7$ giờ $15$ phút đi ngược chiều nhau. Xét tính đúng sai:
\choiceTF
{\True Sau mỗi giờ hai xe đi được quãng đường $90\text{ km}$.}
{\True Hai xe gặp nhau sau $1$ giờ $20$ phút ($1{,}33$ giờ).}
{Thời điểm hai xe gặp nhau là $8$ giờ $45$ phút.}
{\True Thời điểm hai xe gặp nhau là $8$ giờ $35$ phút.}
\loigiai{
\begin{itemize}
    \item Tổng vận tốc: $35 + 55 = 90\text{ km/h}$. Đúng.
    \item Thời gian gặp: $120 : 90 = \dfrac{4}{3}\text{ giờ} = 1$ giờ $20$ phút. Đúng.
    \item Thời điểm gặp: $7$h $15$p $+ 1$h $20$p $= 8$h $35$p $\neq 8$h $45$p. Ý 3 Sai, Ý 4 Đúng.
\end{itemize}
\begin{itemize}
    \item a) \textbf{Đúng}.
    \item b) \textbf{Đúng}.
    \item c) \textbf{Sai}.
    \item d) \textbf{Đúng}.
\end{itemize}
}
\end{ex}

\begin{ex} %[5D10V62]
Một người đi xe đạp từ A đến B mất $3$ giờ với vận tốc $12\text{ km/h}$. Khi từ B trở về A người đó đi xe máy với vận tốc $36\text{ km/h}$. Xét tính đúng sai:
\choiceTF
{\True Quãng đường AB dài $36\text{ km}$.}
{\True Thời gian đi từ B về A bằng xe máy là $1$ giờ.}
{Tổng thời gian cả đi và về là $4$ giờ $30$ phút.}
{\True Vận tốc trung bình cả đi và về là $18\text{ km/h}$.}
\loigiai{
\begin{itemize}
    \item $s = 12 \times 3 = 36\text{ km}$. Đúng.
    \item Thời gian về: $36 : 36 = 1$ giờ. Đúng.
    \item Tổng thời gian: $3 + 1 = 4$ giờ $\neq 4$h $30$p. Sai.
    \item Vận tốc trung bình: $(36 + 36) : (3 + 1) = 72 : 4 = 18\text{ km/h}$. Đúng.
\end{itemize}
\begin{itemize}
    \item a) \textbf{Đúng}.
    \item b) \textbf{Đúng}.
    \item c) \textbf{Sai}.
    \item d) \textbf{Đúng}.
\end{itemize}
}
\end{ex}

\textbf{Phần 3. Trả lời ngắn}
\setcounterpageref{ex}{0}

\begin{ex} %[5D10H62]
Một người đi bộ với vận tốc $4{,}8\text{ km/h}$ trong $1$ giờ $15$ phút. Tính quãng đường người đó đi được theo ki-lô-mét.
\shortans{6}
\loigiai{
\textbf{Đáp số:} $6$\\
$4{,}8 \times 1{,}25 = 6\text{ km}$.
}
\end{ex}

\begin{ex} %[5D10H62]
Một ô tô đi quãng đường $210\text{ km}$ với vận tốc $60\text{ km/h}$. Tính thời gian ô tô đi theo giờ (dùng dấu phẩy nếu có).
\shortans{3,5}
\loigiai{
\textbf{Đáp số:} $3{,}5$\\
$210 : 60 = 3{,}5\text{ giờ}$.
}
\end{ex}

\begin{ex} %[5D10V62]
Hai xe máy đi ngược chiều từ hai tỉnh cách nhau $135\text{ km}$ cùng khởi hành lúc $8$ giờ với vận tốc $42\text{ km/h}$ và $48\text{ km/h}$. Hỏi sau bao nhiêu phút hai xe gặp nhau?
\shortans{90}
\loigiai{
\textbf{Đáp số:} $90$\\
$135 : (42 + 48) = 1{,}5\text{ giờ} = 90\text{ phút}$.
}
\end{ex}

\begin{ex} %[5D10V62]
Một người đi xe máy từ A đuổi theo một người đi xe đạp khởi hành cùng lúc cách A $24\text{ km}$. Vận tốc xe máy là $40\text{ km/h}$, xe đạp là $16\text{ km/h}$. Hỏi sau bao nhiêu giờ xe máy đuổi kịp xe đạp?
\shortans{1}
\loigiai{
\textbf{Đáp số:} $1$\\
$24 : (40 - 16) = 24 : 24 = 1\text{ giờ}$.
}
\end{ex}

\begin{ex} %[5D10V62]
Một con chim bồ câu đưa thư bay quãng đường $120\text{ km}$ với vận tốc $50\text{ km/h}$. Tính thời gian chim bồ câu bay theo phút.
\shortans{144}
\loigiai{
\textbf{Đáp số:} $144$\\
$120 : 50 = 2{,}4\text{ giờ} = 144\text{ phút}$.
}
\end{ex}

\begin{ex} %[5D10C62]
Một chiếc thuyền đi xuôi dòng một khúc sông dài $40\text{ km}$ mất $1$ giờ $15$ phút. Vận tốc dòng nước là $4\text{ km/h}$. Hỏi thuyền đi ngược dòng khúc sông đó mất bao nhiêu giờ (dùng dấu phẩy nếu có)?
\shortans{1,67 hoặc \dfrac{5}{3}}
\loigiai{
\textbf{Đáp số:} $1{,}67$ hoặc $\dfrac{5}{3}$\\
Vận tốc xuôi dòng: $40 : 1{,}25 = 32\text{ km/h}$.\\
Vận tốc ngược dòng: $32 - 4 \times 2 = 24\text{ km/h}$.\\
Thời gian ngược dòng: $40 : 24 = \dfrac{5}{3}\text{ giờ} = 1\text{ giờ } 40\text{ phút} \approx 1{,}67\text{ giờ}$.
}
\end{ex}

\newpage

% =========================================================================
% CHỦ ĐỀ 11 - MỘT SỐ YẾU TỐ THỐNG KÊ VÀ XÁC SUẤT
% =========================================================================
\section{Một số yếu tố thống kê và xác suất}

\TaoPhieuKiemTra{Chủ đề 11 -- Một số yếu tố thống kê và xác suất}{63}{Thu thập, phân loại, sắp xếp số liệu}{https://loptoanthaylinh.vn/thi-online/toan-5/bai-63}
\textbf{Phần 1. Trắc nghiệm lựa chọn}
\setcounterpageref{ex}{0}

\begin{ex} %[5X11N63]
Dãy dữ liệu nào sau đây là dãy số liệu thống kê?
\choice
{Xanh, Đỏ, Vàng, Tím}
{Xuất sắc, Giỏi, Khá, Đạt}
{\True $135\text{ cm}, 140\text{ cm}, 138\text{ cm}, 142\text{ cm}$}
{Bơi lội, Cầu lông, Bóng đá, Bóng rổ}
\loigiai{
Dãy số liệu gồm các giá trị bằng số kèm đơn vị đo.\\
\textbf{Đáp án của câu hỏi là C}
}
\end{ex}

\begin{ex} %[5X11N63]
Trong một lớp học, giáo viên ghi lại số điểm kiểm tra môn Toán của 5 bạn: $9, 8, 10, 8, 9$. Số bạn đạt điểm $8$ là
\choice
{\True $2$ bạn}
{$1$ bạn}
{$3$ bạn}
{$4$ bạn}
\loigiai{
Có 2 bạn đạt điểm 8.\\
\textbf{Đáp án của câu hỏi là A}
}
\end{ex}

\begin{ex} %[5X11H63]
Khảo sát số giờ tự học mỗi ngày của $6$ học sinh thu được: $2\text{ giờ}, 1{,}5\text{ giờ}, 2\text{ giờ}, 3\text{ giờ}, 2{,}5\text{ giờ}, 1\text{ giờ}$. Số giờ tự học nhiều nhất trong số các bạn được khảo sát là
\choice
{$2\text{ giờ}$}
{$2{,}5\text{ giờ}$}
{\True $3\text{ giờ}$}
{$1\text{ giờ}$}
\loigiai{
Giá trị lớn nhất trong dãy là $3\text{ giờ}$.\\
\textbf{Đáp án của câu hỏi là C}
}
\end{ex}

\begin{ex} %[5X11H63]
Cho dãy số liệu về chiều cao (theo $\text{cm}$) của $5$ học sinh: $138, 142, 135, 145, 140$. Sắp xếp dãy số liệu theo thứ tự từ thấp đến cao ta được
\choice
{$145, 142, 140, 138, 135$}
{$135, 140, 138, 142, 145$}
{$138, 135, 140, 142, 145$}
{\True $135, 138, 140, 142, 145$}
\loigiai{
Thứ tự tăng dần: $135, 138, 140, 142, 145$.\\
\textbf{Đáp án của câu hỏi là D}
}
\end{ex}

\begin{ex} %[5X11H63]
Cho bảng số liệu về số cây xanh trồng được của $4$ lớp:
- Lớp 5A: $45$ cây
- Lớp 5B: $52$ cây
- Lớp 5C: $48$ cây
- Lớp 5D: $55$ cây
Tổng số cây xanh cả 4 lớp trồng được là
\choice
{$190$ cây}
{$210$ cây}
{\True $200$ cây}
{$220$ cây}
\loigiai{
Tổng số cây: $45 + 52 + 48 + 55 = 200$ (cây).\\
\textbf{Đáp án của câu hỏi là C}
}
\end{ex}

\begin{ex} %[5X11H63]
Thống kê số lượng học sinh tham gia các câu lạc bộ của trường:
- CLB Cờ vua: $25$ học sinh
- CLB Mỹ thuật: $35$ học sinh
- CLB Âm nhạc: $30$ học sinh
- CLB Bóng đá: $50$ học sinh
Câu lạc bộ có số lượng học sinh tham gia nhiều nhất là
\choice
{CLB Cờ vua}
{CLB Mỹ thuật}
{CLB Âm nhạc}
{\True CLB Bóng đá}
\loigiai{
CLB Bóng đá có $50$ học sinh, đông nhất.\\
\textbf{Đáp án của câu hỏi là D}
}
\end{ex}

\begin{ex} %[5X11V63]
Điểm thi kết thúc học kì 1 môn Tin học của tổ 1 là: $8, 9, 10, 7, 9, 8, 10, 9$. Điểm trung bình môn Tin học của tổ 1 là
\choice
{$8{,}5$}
{\True $8{,}75$}
{$8{,}8$}
{$9{,}0$}
\loigiai{
Tổng điểm: $8 + 9 + 10 + 7 + 9 + 8 + 10 + 9 = 70$.\\
Điểm trung bình: $70 : 8 = 8{,}75$.\\
\textbf{Đáp án của câu hỏi là B}
}
\end{ex}

\begin{ex} %[5X11V63]
Kết quả phân loại học lực của $40$ học sinh lớp 5A:
- Hoàn thành xuất sắc: $16$ bạn
- Hoàn thành tốt: $18$ bạn
- Hoàn thành: $6$ bạn
Tỉ số phần trăm học sinh hoàn thành xuất sắc của lớp 5A là
\choice
{$35\%$}
{$45\%$}
{$50\%$}
{\True $40\%$}
\loigiai{
$16 : 40 = 0{,}4 = 40\%$.\\
\textbf{Đáp án của câu hỏi là D}
}
\end{ex}

\begin{ex} %[5X11V63]
Bảng số liệu cân nặng của $10$ học sinh (theo $\text{kg}$): $32, 35, 30, 38, 42, 35, 36, 40, 35, 37$. Có bao nhiêu học sinh có cân nặng từ $35\text{ kg}$ trở lên?
\choice
{$6$ bạn}
{\True $8$ bạn}
{$7$ bạn}
{$5$ bạn}
\loigiai{
Các giá trị $\ge 35$: $35, 38, 42, 35, 36, 40, 35, 37$ (có 8 bạn).\\
\textbf{Đáp án của câu hỏi là B}
}
\end{ex}

\begin{ex} %[5X11V63]
Một cửa hàng thống kê số lượng xe máy bán được trong 4 tuần:
Tuần 1: $12$ xe; Tuần 2: $15$ xe; Tuần 3: $18$ xe; Tuần 4: $15$ xe.
Trung bình mỗi tuần cửa hàng bán được bao nhiêu chiếc xe máy?
\choice
{\True $15$ xe}
{$14$ xe}
{$16$ xe}
{$17$ xe}
\loigiai{
Tổng số: $12 + 15 + 18 + 15 = 60$ (xe).\\
Trung bình: $60 : 4 = 15$ (xe).\\
\textbf{Đáp án của câu hỏi là A}
}
\end{ex}

\begin{ex} %[5X11V63]
Theo dõi số ngày mưa trong các tháng mùa hè tại một địa phương:
- Tháng 6: $14$ ngày
- Tháng 7: $18$ ngày
- Tháng 8: $16$ ngày
Tỉ số giữa số ngày mưa tháng 7 và tổng số ngày mưa trong cả 3 tháng là
\choice
{$\dfrac{3}{8}$}
{\True $\dfrac{3}{8}$}
{$\dfrac{1}{3}$}
{$\dfrac{7}{24}$}
\loigiai{
Tổng số ngày mưa: $14 + 18 + 16 = 48$ ngày.\\
Tỉ số tháng 7: $\dfrac{18}{48} = \dfrac{3}{8}$.\\
\textbf{Đáp án của câu hỏi là B}
}
\end{ex}

\begin{ex} %[5X11C63]
Thống kê độ tuổi của các vận động viên một đội bơi: $10, 11, 11, 12, 12, 12, 13, 14, 15$. Nếu đội bơi kết nạp thêm một vận động viên $16$ tuổi thì số tuổi trung bình của đội tăng thêm bao nhiêu tuổi?
\choice
{\True $0{,}37$ tuổi (hay $\dfrac{7}{18}$ tuổi)}
{$0{,}2$ tuổi}
{$0{,}4$ tuổi}
{$0{,}5$ tuổi}
\loigiai{
Ban đầu: Tổng tuổi $= 10 + 22 + 36 + 13 + 14 + 15 = 110$ tuổi. Trung bình $= 110 : 9 = 12{,}22$ tuổi.\\
Thêm bạn 16 tuổi: Tổng tuổi $= 126$ tuổi. Trung bình $= 126 : 10 = 12{,}6$ tuổi.\\
Tăng thêm: $12{,}6 - 12{,}22 = 0{,}38$ tuổi (xấp xỉ $0{,}37$ - $0{,}38$).\\
\textbf{Đáp án của câu hỏi là A}
}
\end{ex}

\textbf{Phần 2. Trắc nghiệm đúng sai}
\setcounterpageref{ex}{0}

\begin{ex} %[5X11H63]
Cho bảng thống kê số lượng sách đọc trong mùa hè của $5$ bạn: An: $8$ cuốn, Bình: $12$ cuốn, Chi: $15$ cuốn, Dũng: $9$ cuốn, Mai: $16$ cuốn. Xét tính đúng sai:
\choiceTF
{\True Bạn đọc được nhiều sách nhất là Mai ($16$ cuốn).}
{\True Bạn đọc được ít sách nhất là An ($8$ cuốn).}
{Tổng số sách cả 5 bạn đọc được là $50$ cuốn.}
{\True Trung bình mỗi bạn đọc được $12$ cuốn sách.}
\loigiai{
\begin{itemize}
    \item Mai đọc nhiều nhất (16 cuốn). Đúng.
    \item An đọc ít nhất (8 cuốn). Đúng.
    \item Tổng số sách: $8 + 12 + 15 + 9 + 16 = 60$ cuốn $\neq 50$ cuốn. Sai.
    \item Trung bình: $60 : 5 = 12$ cuốn. Đúng.
\end{itemize}
\begin{itemize}
    \item a) \textbf{Đúng}.
    \item b) \textbf{Đúng}.
    \item c) \textbf{Sai}.
    \item d) \textbf{Đúng}.
\end{itemize}
}
\end{ex}

\begin{ex} %[5X11V63]
Bảng khảo sát món ăn sáng yêu thích của $50$ học sinh:
- Bánh mì: $18$ bạn
- Phở: $15$ bạn
- Xôi: $12$ bạn
- Bún: $5$ bạn
Xét tính đúng sai:
\choiceTF
{\True Số bạn thích Bánh mì chiếm tỉ lệ lớn nhất ($36\%$).}
{\True Số bạn thích Phở chiếm tỉ lệ $30\%$.}
{Số bạn thích Xôi và Bún chiếm tất cả $40\%$ tổng số học sinh.}
{\True Số bạn thích Xôi và Bún chiếm tất cả $34\%$ tổng số học sinh.}
\loigiai{
\begin{itemize}
    \item Bánh mì: $18 : 50 = 36\%$. Đúng.
    \item Phở: $15 : 50 = 30\%$. Đúng.
    \item Xôi + Bún: $12 + 5 = 17$ bạn. Tỉ lệ: $17 : 50 = 34\% \neq 40\%$. Ý 3 Sai, Ý 4 Đúng.
\end{itemize}
\begin{itemize}
    \item a) \textbf{Đúng}.
    \item b) \textbf{Đúng}.
    \item c) \textbf{Sai}.
    \item d) \textbf{Đúng}.
\end{itemize}
}
\end{ex}

\textbf{Phần 3. Trả lời ngắn}
\setcounterpageref{ex}{0}

\begin{ex} %[5X11H63]
Cho dãy điểm kiểm tra: $7, 8, 9, 8, 10, 8, 9, 7$. Tính điểm trung bình (dùng dấu phẩy nếu có).
\shortans{8,25}
\loigiai{
\textbf{Đáp số:} $8{,}25$\\
Tổng $= 66 \Rightarrow 66 : 8 = 8{,}25$.
}
\end{ex}

\begin{ex} %[5X11H63]
Bảng số liệu thời gian đi bộ (phút) của $5$ bạn: $12, 15, 10, 18, 20$. Tính hiệu giữa thời gian nhiều nhất và ít nhất theo phút.
\shortans{10}
\loigiai{
\textbf{Đáp số:} $10$\\
$20 - 10 = 10$ phút.
}
\end{ex}

\begin{ex} %[5X11V63]
Một lớp có $32$ học sinh, trong đó có $24$ bạn biết bơi. Tính tỉ số phần trăm học sinh chưa biết bơi của lớp. (Chỉ ghi số).
\shortans{25}
\loigiai{
\textbf{Đáp số:} $25$\\
Số bạn chưa biết bơi: $32 - 24 = 8$. Tỉ số phần trăm: $8 : 32 = 25\%$.
}
\end{ex}

\begin{ex} %[5X11V63]
Tổng số điểm kiểm tra của một nhóm $4$ bạn là $34$ điểm. Nếu thêm điểm của bạn thứ năm là $9$ điểm thì điểm trung bình của cả $5$ bạn là bao nhiêu? (Dùng dấu phẩy nếu có).
\shortans{8,6}
\loigiai{
\textbf{Đáp số:} $8{,}6$\\
Tổng điểm 5 bạn: $34 + 9 = 43$. Trung bình: $43 : 5 = 8{,}6$.
}
\end{ex}

\begin{ex} %[5X11V63]
Bảng số liệu thu gom vỏ chai nhựa của $3$ tổ: Tổ 1: $120$ vỏ, Tổ 2: $150$ vỏ, Tổ 3: $180$ vỏ. Tính số vỏ chai trung bình mỗi tổ thu gom được.
\shortans{150}
\loigiai{
\textbf{Đáp số:} $150$\\
$(120 + 150 + 180) : 3 = 450 : 3 = 150$ vỏ.
}
\end{ex}

\begin{ex} %[5X11C63]
Số lượng bánh bán được trong 5 ngày liên tiếp là: $45, 52, 60, 48, 55$. Ngày thứ sáu cửa hàng cần bán bao nhiêu chiếc bánh để mức trung bình trong 6 ngày đạt $55$ chiếc bánh?
\shortans{70}
\loigiai{
\textbf{Đáp số:} $70$\\
Tổng 6 ngày: $55 \times 6 = 330$ chiếc.\\
Tổng 5 ngày đầu: $45 + 52 + 60 + 48 + 55 = 260$ chiếc.\\
Ngày thứ sáu: $330 - 260 = 70$ chiếc.
}
\end{ex}

\newpage

\TaoPhieuKiemTra{Chủ đề 11 -- Một số yếu tố thống kê và xác suất}{64}{Biểu đồ hình quạt tròn}{https://loptoanthaylinh.vn/thi-online/toan-5/bai-64}
\textbf{Phần 1. Trắc nghiệm lựa chọn}
\setcounterpageref{ex}{0}

\begin{ex} %[5X11N64]
Biểu đồ hình quạt tròn thường dùng để biểu diễn
\choice
{Sự thay đổi số liệu theo thời gian}
{\True Tỉ số phần trăm của các phần trong toàn bộ tổng thể}
{Độ dài của các đoạn thẳng}
{Tọa độ các điểm trong không gian}
\loigiai{
Biểu đồ hình quạt tròn biểu diễn tỉ số phần trăm của các thành phần trong toàn bộ tổng thể $100\%$.\\
\textbf{Đáp án của câu hỏi là B}
}
\end{ex}

\begin{ex} %[5X11N64]
Tổng tỉ số phần trăm của tất cả các hình quạt trên một biểu đồ hình quạt tròn luôn bằng
\choice
{\True $100\%$}
{$50\%$}
{$180\%$}
{$360\%$}
\loigiai{
Toàn bộ hình tròn biểu diễn $100\%$.\\
\textbf{Đáp án của câu hỏi là A}
}
\end{ex}

\begin{ex} %[5X11H64]
Một biểu đồ hình quạt tròn biểu diễn thể loại phim yêu thích của học sinh gồm: Phim hoạt hình $50\%$, Phim khoa học $25\%$, Phim phiêu lưu $15\%$, Phim hài $10\%$. Thể loại phim chiếm một nửa hình tròn là
\choice
{Phim khoa học}
{Phim phiêu lưu}
{\True Phim hoạt hình}
{Phim hài}
\loigiai{
$50\%$ ứng với một nửa hình tròn, đó là phim hoạt hình.\\
\textbf{Đáp án của câu hỏi là C}
}
\end{ex}

\begin{ex} %[5X11H64]
Biểu đồ hình quạt tròn thể hiện tỉ lệ tham gia các môn thể thao của $200$ học sinh: Bóng đá $40\%$, Cầu lông $30\%$, Bơi lội $20\%$, Cờ vua $10\%$. Số học sinh tham gia môn Bóng đá là
\choice
{$60$ bạn}
{$70$ bạn}
{$90$ bạn}
{\True $80$ bạn}
\loigiai{
Số học sinh chơi bóng đá: $200 \times 40 : 100 = 80$ (bạn).\\
\textbf{Đáp án của câu hỏi là D}
}
\end{ex}

\begin{ex} %[5X11H64]
Biểu đồ quạt tròn biểu diễn cơ cấu cây trồng trong một khu vườn: Xoài $35\%$, Nhãn $25\%$, Cam $20\%$, Bưởi chưa rõ. Tỉ số phần trăm diện tích trồng Bưởi là
\choice
{$15\%$}
{\True $20\%$}
{$25\%$}
{$10\%$}
\loigiai{
Tỉ số phần trăm bưởi: $100\% - (35\% + 25\% + 20\%) = 20\%$.\\
\textbf{Đáp án của câu hỏi là B}
}
\end{ex}

\begin{ex} %[5X11H64]
Biểu đồ hình quạt tròn cho biết tỉ lệ học sinh khối 5 tham gia câu lạc bộ: Tiếng Anh $45\%$, Toán học $35\%$, Tin học $20\%$. Biết khối 5 có $160$ học sinh. Số học sinh tham gia CLB Tin học là
\choice
{$24$ bạn}
{$28$ bạn}
{\True $32$ bạn}
{$36$ bạn}
\loigiai{
$160 \times 20 : 100 = 32$ (bạn).\\
\textbf{Đáp án của câu hỏi là C}
}
\end{ex}

\begin{ex} %[5X11V64]
Khảo sát phương tiện đến trường của $300$ học sinh: Xe đạp chiếm $45\%$, Xe buýt chiếm $25\%$, Đi bộ chiếm $20\%$, Bố mẹ đưa đón chiếm $10\%$. Số học sinh đi xe đạp nhiều hơn số học sinh đi bộ là
\choice
{$60$ bạn}
{\True $75$ bạn}
{$80$ bạn}
{$90$ bạn}
\loigiai{
Chênh lệch tỉ lệ: $45\% - 20\% = 25\%$.\\
Số bạn chênh lệch: $300 \times 25\% = 75$ (bạn).\\
\textbf{Đáp án của câu hỏi là B}
}
\end{ex}

\begin{ex} %[5X11V64]
Một trang trại nuôi $500$ con gia cầm. Biểu đồ hình quạt cho thấy Gà chiếm $55\%$, Vịt chiếm $30\%$, Ngan chiếm $15\%$. Tổng số con Vịt và Ngan là
\choice
{\True $225$ con}
{$200$ con}
{$250$ con}
{$275$ con}
\loigiai{
Tỉ lệ Vịt và Ngan: $30\% + 15\% = 45\%$.\\
Số lượng: $500 \times 45\% = 225$ (con).\\
\textbf{Đáp án của câu hỏi là A}
}
\end{ex}

\begin{ex} %[5X11V64]
Biểu đồ quạt tròn cho biết kết quả xếp loại học tập của $120$ học sinh khối 5. Trong đó số học sinh Khá chiếm $40\%$, học sinh Giỏi chiếm $35\%$, còn lại là học sinh Trung bình. Số học sinh Trung bình là
\choice
{$24$ bạn}
{$36$ bạn}
{$42$ bạn}
{\True $30$ bạn}
\loigiai{
Tỉ lệ Trung bình: $100\% - (40\% + 35\%) = 25\%$.\\
Số học sinh Trung bình: $120 \times 25\% = 30$ (bạn).\\
\textbf{Đáp án của câu hỏi là D}
}
\end{ex}

\begin{ex} %[5X11V64]
Biểu đồ quạt biểu diễn chi tiêu trong tháng của gia đình bác An: Ăn uống $45\%$, Giáo dục $20\%$, Tiết kiệm $15\%$, Chi khác $20\%$. Biết số tiền dành cho Giáo dục là $4\,000\,000$ đồng. Tổng thu nhập hàng tháng của gia đình bác An là
\choice
{$16\,000\,000$ đồng}
{$24\,000\,000$ đồng}
{\True $20\,000\,000$ đồng}
{$25\,000\,000$ đồng}
\loigiai{
Tổng thu nhập: $4\,000\,000 : 20 \times 100 = 20\,000\,000$ (đồng).\\
\textbf{Đáp án của câu hỏi là C}
}
\end{ex}

\begin{ex} %[5X11V64]
Trên biểu đồ hình quạt tròn, một hình quạt có góc ở tâm là $90^\circ$ (chiếm $\dfrac{1}{4}$ hình tròn) thì biểu diễn tỉ lệ phần trăm là
\choice
{$20\%$}
{$30\%$}
{$50\%$}
{\True $25\%$}
\loigiai{
$\dfrac{90}{360} = \dfrac{1}{4} = 25\%$.\\
\textbf{Đáp án của câu hỏi là D}
}
\end{ex}

\begin{ex} %[5X11C64]
Biểu đồ hình quạt tròn cho biết cơ cấu bán hàng của siêu thị gồm $4$ nhóm: Thực phẩm ($40\%$), Đồ gia dụng ($25\%$), Thời trang ($20\%$), Hóa mỹ phẩm ($15\%$). Biết doanh thu từ Thực phẩm nhiều hơn Thời trang là $80$ triệu đồng. Tổng doanh thu của siêu thị là
\choice
{\True $400$ triệu đồng}
{$360$ triệu đồng}
{$450$ triệu đồng}
{$500$ triệu đồng}
\loigiai{
Hiệu tỉ lệ: $40\% - 20\% = 20\%$.\\
Tổng doanh thu: $80 : 20 \times 100 = 400$ (triệu đồng).\\
\textbf{Đáp án của câu hỏi là A}
}
\end{ex}

\textbf{Phần 2. Trắc nghiệm đúng sai}
\setcounterpageref{ex}{0}

\begin{ex} %[5X11H64]
Biểu đồ hình quạt biểu diễn cơ cấu các loại quả trong vườn: Cam $40\%$, Táo $30\%$, Bưởi $20\%$, Xoài $10\%$. Tổng số có $600$ cây. Xét tính đúng sai:
\choiceTF
{\True Số cây Cam trong vườn là $240$ cây.}
{\True Số cây Táo trong vườn là $180$ cây.}
{Số cây Bưởi ít hơn số cây Xoài là $60$ cây.}
{\True Tổng số cây Bưởi và Xoài là $180$ cây.}
\loigiai{
\begin{itemize}
    \item Cam: $600 \times 40\% = 240$ cây. Đúng.
    \item Táo: $600 \times 30\% = 180$ cây. Đúng.
    \item Bưởi ($20\% = 120$ cây) nhiều hơn Xoài ($10\% = 60$ cây). Khẳng định ít hơn là Sai.
    \item Bưởi + Xoài: $120 + 60 = 180$ cây. Đúng.
\end{itemize}
\begin{itemize}
    \item a) \textbf{Đúng}.
    \item b) \textbf{Đúng}.
    \item c) \textbf{Sai}.
    \item d) \textbf{Đúng}.
\end{itemize}
}
\end{ex}

\begin{ex} %[5X11V64]
Biểu đồ quạt tròn biểu diễn kết quả học lực của $240$ học sinh khối 5: Xuất sắc $30\%$, Tốt $45\%$, Hoàn thành $20\%$, Chưa hoàn thành $5\%$. Xét tính đúng sai:
\choiceTF
{\True Số học sinh Xuất sắc là $72$ bạn.}
{\True Số học sinh Tốt là $108$ bạn.}
{Số học sinh Chưa hoàn thành là $15$ bạn.}
{\True Số học sinh Chưa hoàn thành là $12$ bạn.}
\loigiai{
\begin{itemize}
    \item Xuất sắc: $240 \times 30\% = 72$ bạn. Đúng.
    \item Tốt: $240 \times 45\% = 108$ bạn. Đúng.
    \item Chưa hoàn thành: $240 \times 5\% = 12$ bạn $\neq 15$ bạn. Ý 3 Sai, Ý 4 Đúng.
\end{itemize}
\begin{itemize}
    \item a) \textbf{Đúng}.
    \item b) \textbf{Đúng}.
    \item c) \textbf{Sai}.
    \item d) \textbf{Đúng}.
\end{itemize}
}
\end{ex}

\textbf{Phần 3. Trả lời ngắn}
\setcounterpageref{ex}{0}

\begin{ex} %[5X11H64]
Một biểu đồ hình quạt tròn có 3 phần: Phần A chiếm $45\%$, Phần B chiếm $35\%$. Hỏi Phần C chiếm bao nhiêu phần trăm? (Chỉ ghi số).
\shortans{20}
\loigiai{
\textbf{Đáp số:} $20$\\
$100 - (45 + 35) = 20\%$.
}
\end{ex}

\begin{ex} %[5X11H64]
Biểu đồ quạt tròn cho biết có $25\%$ học sinh thích môn Bơi lội trên tổng số $160$ học sinh. Tính số học sinh thích môn Bơi lội.
\shortans{40}
\loigiai{
\textbf{Đáp số:} $40$\\
$160 \times 25 : 100 = 40$.
}
\end{ex}

\begin{ex} %[5X11V64]
Biểu đồ quạt tròn cho biết cơ cấu số sách thư viện: Sách giáo khoa $50\%$, Truyện tranh $30\%$, Sách tham khảo $20\%$. Biết thư viện có $1\,500$ cuốn sách. Tính số cuốn truyện tranh.
\shortans{450}
\loigiai{
\textbf{Đáp số:} $450$\\
$1\,500 \times 30 : 100 = 450$.
}
\end{ex}

\begin{ex} %[5X11V64]
Biểu đồ quạt cho biết $15\%$ số cây trong vườn là cây mít, tương ứng với $45$ cây mít. Tính tổng số cây trong vườn.
\shortans{300}
\loigiai{
\textbf{Đáp số:} $300$\\
$45 : 15 \times 100 = 300$.
}
\end{ex}

\begin{ex} %[5X11V64]
Biểu đồ quạt tròn biểu diễn tỉ lệ học sinh đi bộ đến trường là $35\%$, đi xe đạp là $45\%$. Tổng số học sinh toàn trường là $800$ em. Hỏi số học sinh đi xe đạp nhiều hơn đi bộ bao nhiêu em?
\shortans{80}
\loigiai{
\textbf{Đáp số:} $80$\\
$800 \times (45 - 35) : 100 = 80$.
}
\end{ex}

\begin{ex} %[5X11C64]
Một cửa hàng bán được $360\text{ kg}$ trái cây gồm Cam, Táo, Nho. Trên biểu đồ hình quạt tròn, phần diện tích biểu diễn Nho chiếm $\dfrac{1}{6}$ hình tròn, Cam chiếm $50\%$, còn lại là Táo. Tính số ki-lô-gam Táo đã bán.
\shortans{120}
\loigiai{
\textbf{Đáp số:} $120$\\
Nho: $360 : 6 = 60\text{ kg}$.\\
Cam: $360 \times 50\% = 180\text{ kg}$.\\
Táo: $360 - (60 + 180) = 120\text{ kg}$.
}
\end{ex}

\newpage

\TaoPhieuKiemTra{Chủ đề 11 -- Một số yếu tố thống kê và xác suất}{65}{Tỉ số số lần lặp lại sự kiện}{https://loptoanthaylinh.vn/thi-online/toan-5/bai-65}
\textbf{Phần 1. Trắc nghiệm lựa chọn}
\setcounterpageref{ex}{0}

\begin{ex} %[5X11N65]
Tung một đồng xu $20$ lần, thấy có $11$ lần xuất hiện mặt Ngửa. Tỉ số của số lần xuất hiện mặt Ngửa so với tổng số lần tung là
\choice
{$\dfrac{9}{20}$}
{\True $\dfrac{11}{20}$}
{$\dfrac{11}{9}$}
{$\dfrac{9}{11}$}
\loigiai{
Tỉ số là $\dfrac{11}{20}$.\\
\textbf{Đáp án của câu hỏi là B}
}
\end{ex}

\begin{ex} %[5X11N65]
Gieo một con xúc xắc $50$ lần, có $8$ lần xuất hiện mặt $6$ chấm. Tỉ số phần trăm số lần xuất hiện mặt $6$ chấm là
\choice
{\True $16\%$}
{$12\%$}
{$18\%$}
{$20\%$}
\loigiai{
Tỉ số phần trăm: $\dfrac{8}{50} = \dfrac{16}{100} = 16\%$.\\
\textbf{Đáp án của câu hỏi là A}
}
\end{ex}

\begin{ex} %[5X11H65]
Trong một trò chơi rút thẻ từ hộp, Nam rút $40$ lần và có $14$ lần rút được thẻ màu đỏ. Tỉ số của số lần rút được thẻ không phải màu đỏ so với tổng số lần rút là
\choice
{$\dfrac{7}{20}$}
{\True $\dfrac{13}{20}$}
{$\dfrac{7}{13}$}
{$\dfrac{13}{7}$}
\loigiai{
Số lần rút không phải màu đỏ: $40 - 14 = 26$ lần.\\
Tỉ số: $\dfrac{26}{40} = \dfrac{13}{20}$.\\
\textbf{Đáp án của câu hỏi là B}
}
\end{ex}

\begin{ex} %[5X11H65]
Bắn bia $25$ phát, có $20$ phát trúng tâm bia. Tỉ số phần trăm số phát bắn trúng tâm bia là
\choice
{$75\%$}
{$85\%$}
{$90\%$}
{\True $80\%$}
\loigiai{
$20 : 25 = 0{,}8 = 80\%$.\\
\textbf{Đáp án của câu hỏi là D}
}
\end{ex}

\begin{ex} %[5X11H65]
Quay vòng quay may mắn $30$ lần, mũi tên dừng lại ở ô “Chúc mừng” $6$ lần. Tỉ số của số lần dừng lại ở ô “Chúc mừng” so với tổng số lần quay viết dưới dạng phân số tối giản là
\choice
{\True $\dfrac{1}{5}$}
{$\dfrac{6}{30}$}
{$\dfrac{1}{6}$}
{$\dfrac{2}{5}$}
\loigiai{
$\dfrac{6}{30} = \dfrac{1}{5}$.\\
\textbf{Đáp án của câu hỏi là A}
}
\end{ex}

\begin{ex} %[5X11H65]
Mai thực hiện tung đồng xu $100$ lần thì thấy mặt Sấp xuất hiện $48$ lần. Tỉ số phần trăm số lần xuất hiện mặt Ngửa là
\choice
{$48\%$}
{$50\%$}
{\True $52\%$}
{$54\%$}
\loigiai{
Số lần mặt Ngửa: $100 - 48 = 52$ lần $\Rightarrow 52\%$.\\
\textbf{Đáp án của câu hỏi là C}
}
\end{ex}

\begin{ex} %[5X11V65]
Gieo một con xúc xắc $60$ lần, kết quả xuất hiện các mặt chẵn ($2, 4, 6$) là $33$ lần. Tỉ số của số lần xuất hiện mặt lẻ so với tổng số lần gieo là
\choice
{$\dfrac{11}{20}$}
{$\dfrac{33}{60}$}
{$\dfrac{1}{2}$}
{\True $\dfrac{9}{20}$}
\loigiai{
Số lần xuất hiện mặt lẻ: $60 - 33 = 27$ lần.\\
Tỉ số: $\dfrac{27}{60} = \dfrac{9}{20}$.\\
\textbf{Đáp án của câu hỏi là D}
}
\end{ex}

\begin{ex} %[5X11V65]
Một thủ môn bắt phạt đền $50$ lần, cản phá thành công $15$ lần. Tỉ số phần trăm số lần cản phá thành công của thủ môn đó là
\choice
{\True $30\%$}
{$25\%$}
{$35\%$}
{$40\%$}
\loigiai{
$15 : 50 = 30\%$.\\
\textbf{Đáp án của câu hỏi là A}
}
\end{ex}

\begin{ex} %[5X11V65]
Trong một hộp có $3$ quả bóng Xanh, Đỏ, Vàng. Rút ngẫu nhiên $1$ bóng rồi trả lại vào hộp. Sau $80$ lần rút, bóng Xanh xuất hiện $28$ lần, bóng Đỏ xuất hiện $32$ lần. Tỉ số phần trăm số lần rút được bóng Vàng là
\choice
{$20\%$}
{\True $25\%$}
{$30\%$}
{$35\%$}
\loigiai{
Số lần bóng Vàng: $80 - (28 + 32) = 20$ lần.\\
Tỉ số phần trăm: $20 : 80 = 0{,}25 = 25\%$.\\
\textbf{Đáp án của câu hỏi là B}
}
\end{ex}

\begin{ex} %[5X11V65]
Kiểm tra ngẫu nhiên $200$ sản phẩm của một dây chuyền may, phát hiện có $6$ sản phẩm lỗi. Tỉ số phần trăm sản phẩm đạt chất lượng tốt là
\choice
{$3\%$}
{$94\%$}
{\True $97\%$}
{$96\%$}
\loigiai{
Số sản phẩm tốt: $200 - 6 = 194$.\\
Tỉ số phần trăm: $194 : 200 = 97\%$.\\
\textbf{Đáp án của câu hỏi là C}
}
\end{ex}

\begin{ex} %[5X11V65]
Hai bạn An và Bình cùng chơi tung đồng xu. An tung $50$ lần được $24$ lần Sấp. Bình tung $50$ lần được $28$ lần Sấp. Tỉ số phần trăm số lần xuất hiện mặt Sấp của cả hai bạn gộp lại là
\choice
{$50\%$}
{$54\%$}
{$56\%$}
{\True $52\%$}
\loigiai{
Tổng số lần mặt Sấp: $24 + 28 = 52$ lần.\\
Tổng số lần tung: $50 + 50 = 100$ lần.\\
Tỉ số phần trăm: $52 : 100 = 52\%$.\\
\textbf{Đáp án của câu hỏi là D}
}
\end{ex}

\begin{ex} %[5X11C65]
Một xạ thủ bắn $100$ viên đạn vào bia. Tỉ số phần trăm số lần bắn trúng điểm 10 là $65\%$, bắn trúng điểm 9 là $25\%$, còn lại trúng điểm 8. Tổng số điểm mà xạ thủ đó đạt được sau 100 phát bắn là
\choice
{$945$ điểm}
{$960$ điểm}
{\True $955$ điểm}
{$950$ điểm}
\loigiai{
Số phát điểm 10: $65$ phát $\Rightarrow 65 \times 10 = 650$ điểm.\\
Số phát điểm 9: $25$ phát $\Rightarrow 25 \times 9 = 225$ điểm.\\
Số phát điểm 8: $100 - (65 + 25) = 10$ phát $\Rightarrow 10 \times 8 = 80$ điểm.\\
Tổng điểm: $650 + 225 + 80 = 955$ điểm.\\
\textbf{Đáp án của câu hỏi là C}
}
\end{ex}

\textbf{Phần 2. Trắc nghiệm đúng sai}
\setcounterpageref{ex}{0}

\begin{ex} %[5X11H65]
Thực hiện quay một tấm bìa hình tròn chia thành 4 phần bằng nhau ghi các số $1, 2, 3, 4$ tổng cộng $40$ lần. Kết quả dừng lại: Số 1: $8$ lần; Số 2: $12$ lần; Số 3: $14$ lần; Số 4: $6$ lần. Xét tính đúng sai:
\choiceTF
{\True Tỉ số số lần xuất hiện số 1 là $\dfrac{1}{5}$.}
{\True Tỉ số phần trăm số lần xuất hiện số 2 là $30\%$.}
{Số lần xuất hiện số lẻ chiếm $50\%$ tổng số lần quay.}
{\True Số lần xuất hiện số lẻ chiếm $55\%$ tổng số lần quay.}
\loigiai{
\begin{itemize}
    \item Số 1: $8 : 40 = \dfrac{1}{5}$. Đúng.
    \item Số 2: $12 : 40 = 30\%$. Đúng.
    \item Số lẻ (số 1 và số 3): $8 + 14 = 22$ lần $\Rightarrow 22 : 40 = 55\% \neq 50\%$. Ý 3 Sai, Ý 4 Đúng.
\end{itemize}
\begin{itemize}
    \item a) \textbf{Đúng}.
    \item b) \textbf{Đúng}.
    \item c) \textbf{Sai}.
    \item d) \textbf{Đúng}.
\end{itemize}
}
\end{ex}

\begin{ex} %[5X11V65]
Tung một con xúc xắc $120$ lần. Kết quả:
- Mặt 1 chấm: $18$ lần
- Mặt 2 chấm: $22$ lần
- Mặt 3 chấm: $20$ lần
- Mặt 4 chấm: $16$ lần
- Mặt 5 chấm: $24$ lần
- Mặt 6 chấm: $20$ lần
Xét tính đúng sai:
\choiceTF
{\True Tỉ số phần trăm mặt 5 chấm xuất hiện là $20\%$.}
{\True Tỉ số phần trăm mặt 1 chấm xuất hiện là $15\%$.}
{Tổng số lần xuất hiện mặt lớn hơn 4 chấm là $40$ lần.}
{\True Tổng số lần xuất hiện mặt lớn hơn 4 chấm là $44$ lần.}
\loigiai{
\begin{itemize}
    \item Mặt 5 chấm: $24 : 120 = 20\%$. Đúng.
    \item Mặt 1 chấm: $18 : 120 = 15\%$. Đúng.
    \item Mặt lớn hơn 4 chấm (5 và 6 chấm): $24 + 20 = 44$ lần $\neq 40$ lần. Ý 3 Sai, Ý 4 Đúng.
\end{itemize}
\begin{itemize}
    \item a) \textbf{Đúng}.
    \item b) \textbf{Đúng}.
    \item c) \textbf{Sai}.
    \item d) \textbf{Đúng}.
\end{itemize}
}
\end{ex}

\textbf{Phần 3. Trả lời ngắn}
\setcounterpageref{ex}{0}

\begin{ex} %[5X11H65]
Tung một đồng xu $50$ lần, có $27$ lần xuất hiện mặt Sấp. Tính tỉ số phần trăm số lần xuất hiện mặt Sấp. (Chỉ ghi số).
\shortans{54}
\loigiai{
\textbf{Đáp số:} $54$\\
$27 : 50 = 54\%$.
}
\end{ex}

\begin{ex} %[5X11H65]
Gieo một con xúc xắc $40$ lần, có $6$ lần xuất hiện mặt 6 chấm. Viết tỉ số của số lần xuất hiện mặt 6 chấm dưới dạng phân số tối giản $\dfrac{a}{b}$. Tính $a + b$.
\shortans{23}
\loigiai{
\textbf{Đáp số:} $23$\\
$\dfrac{6}{40} = \dfrac{3}{20} \Rightarrow a = 3, b = 20 \Rightarrow a + b = 23$.
}
\end{ex}

\begin{ex} %[5X11V65]
Một cầu thủ sút luân lưu $20$ lần, bóng vào lưới $16$ lần. Tính tỉ số phần trăm số lần sút bóng không vào lưới. (Chỉ ghi số).
\shortans{20}
\loigiai{
\textbf{Đáp số:} $20$\\
Không vào lưới: $4$ lần $\Rightarrow 4 : 20 = 20\%$.
}
\end{ex}

\begin{ex} %[5X11V65]
Bốc ngẫu nhiên một viên bi từ túi nhiều lần (có hoàn lại). Sau $60$ lần bốc, có $21$ lần bốc được bi xanh. Tính tỉ số phần trăm số lần bốc được bi xanh. (Chỉ ghi số).
\shortans{35}
\loigiai{
\textbf{Đáp số:} $35$\\
$21 : 60 = 0{,}35 = 35\%$.
}
\end{ex}

\begin{ex} %[5X11V65]
Một nhà máy kiểm tra $500$ sản phẩm thấy có $485$ sản phẩm đạt chuẩn. Tính tỉ số phần trăm sản phẩm không đạt chuẩn. (Dùng dấu phẩy nếu có).
\shortans{3}
\loigiai{
\textbf{Đáp số:} $3$\\
Không đạt chuẩn: $15$ sản phẩm $\Rightarrow 15 : 500 = 3\%$.
}
\end{ex}

\begin{ex} %[5X11C65]
Gieo một con xúc xắc $150$ lần, tỉ số phần trăm xuất hiện mặt số chẵn là $48\%$. Hỏi mặt số lẻ xuất hiện bao nhiêu lần?
\shortans{78}
\loigiai{
\textbf{Đáp số:} $78$\\
Mặt số lẻ chiếm $100\% - 48\% = 52\%$.\\
Số lần xuất hiện mặt lẻ: $150 \times 52 : 100 = 78$ lần.
}
\end{ex}

\newpage

\TaoPhieuKiemTra{Chủ đề 11 -- Một số yếu tố thống kê và xác suất}{66}{Thực hành trải nghiệm thống kê}{https://loptoanthaylinh.vn/thi-online/toan-5/bai-66}
\textbf{Phần 1. Trắc nghiệm lựa chọn}
\setcounterpageref{ex}{0}

\begin{ex} %[5X11N66]
Khi thực hiện khảo sát sở thích của các bạn trong lớp, phương pháp thu thập dữ liệu phù hợp là
\choice
{Đo lường bằng thước kẹp}
{Cân bằng cân điện tử}
{\True Phát phiếu khảo sát hoặc phỏng vấn trực tiếp}
{Bấm giờ bằng đồng hồ}
\loigiai{
Khảo sát sở thích thường dùng phiếu hỏi hoặc phỏng vấn trực tiếp.\\
\textbf{Đáp án của câu hỏi là C}
}
\end{ex}

\begin{ex} %[5X11N66]
Dạng biểu đồ nào thích hợp nhất để so sánh trực quan số lượng học sinh thích các môn thể thao khác nhau?
\choice
{\True Biểu đồ cột hoặc biểu đồ hình quạt tròn}
{Biểu đồ đường gấp khúc}
{Sơ đồ hình cây}
{Biểu đồ phân tán}
\loigiai{
Biểu đồ cột hoặc biểu đồ hình quạt tròn rất trực quan để so sánh cơ cấu và số lượng.\\
\textbf{Đáp án của câu hỏi là A}
}
\end{ex}

\begin{ex} %[5X11H66]
Khảo sát về loại trái cây yêu thích của $40$ bạn học sinh lớp 5B: Xoài ($14$ bạn), Cam ($10$ bạn), Ổi ($8$ bạn), Dưa hấu ($8$ bạn). Loại quả được nhiều bạn yêu thích nhất là
\choice
{Cam}
{\True Xoài}
{Ổi}
{Dưa hấu}
\loigiai{
Xoài có $14$ bạn, nhiều nhất.\\
\textbf{Đáp án của câu hỏi là B}
}
\end{ex}

\begin{ex} %[5X11H66]
Bảng khảo sát số cuốn sách đọc được trong tuần của một nhóm bạn:
- $1$ cuốn: $3$ bạn
- $2$ cuốn: $5$ bạn
- $3$ cuốn: $2$ bạn
Tổng số cuốn sách nhóm bạn đã đọc trong tuần là
\choice
{$10$ cuốn}
{$20$ cuốn}
{\True $19$ cuốn}
{$15$ cuốn}
\loigiai{
Tổng số sách: $1 \times 3 + 2 \times 5 + 3 \times 2 = 3 + 10 + 6 = 19$ (cuốn).\\
\textbf{Đáp án của câu hỏi là C}
}
\end{ex}

\begin{ex} %[5X11H66]
Khảo sát thời gian xem ti vi mỗi ngày của các bạn trong tổ: $30$ phút, $45$ phút, $60$ phút, $45$ phút, $30$ phút, $90$ phút. Thời gian xem ti vi trung bình của một bạn trong tổ là
\choice
{$45$ phút}
{$55$ phút}
{$60$ phút}
{\True $50$ phút}
\loigiai{
Tổng thời gian: $30 + 45 + 60 + 45 + 30 + 90 = 300$ phút.\\
Trung bình: $300 : 6 = 50$ phút.\\
\textbf{Đáp án của câu hỏi là D}
}
\end{ex}

\begin{ex} %[5X11H66]
Khi biểu diễn số liệu học sinh tham gia các hoạt động ngoại khóa: Văn nghệ chiếm $40\%$, Thể thao chiếm $60\%$ lên biểu đồ hình quạt tròn, phần thể thao chiếm
\choice
{Ít hơn một nửa hình tròn}
{\True Nhiều hơn một nửa hình tròn}
{Đúng một nửa hình tròn}
{Toàn bộ hình tròn}
\loigiai{
$60\% > 50\%$ nên phần thể thao chiếm nhiều hơn một nửa hình tròn.\\
\textbf{Đáp án của câu hỏi là B}
}
\end{ex}

\begin{ex} %[5X11V66]
Kết quả đo chiều cao của $5$ học sinh: $138\text{ cm}, 142\text{ cm}, 140\text{ cm}, 136\text{ cm}, 144\text{ cm}$. Chiều cao trung bình của $5$ học sinh đó là
\choice
{\True $140\text{ cm}$}
{$139\text{ cm}$}
{$141\text{ cm}$}
{$142\text{ cm}$}
\loigiai{
Tổng chiều cao: $138 + 142 + 140 + 136 + 144 = 700\text{ cm}$.\\
Trung bình: $700 : 5 = 140\text{ cm}$.\\
\textbf{Đáp án của câu hỏi là A}
}
\end{ex}

\begin{ex} %[5X11V66]
Khảo sát số thành viên trong gia đình của $30$ bạn học sinh thu được:
- Gia đình có $3$ người: $6$ bạn
- Gia đình có $4$ người: $18$ bạn
- Gia đình có $5$ người: $6$ bạn
Tỉ số phần trăm số bạn có gia đình gồm $4$ người là
\choice
{$50\%$}
{$40\%$}
{$70\%$}
{\True $60\%$}
\loigiai{
$18 : 30 = 0{,}6 = 60\%$.\\
\textbf{Đáp án của câu hỏi là D}
}
\end{ex}

\begin{ex} %[5X11V66]
Một nhóm học sinh đo nhiệt độ lúc $12$ giờ trưa trong $5$ ngày liên tiếp: $28^\circ\text{C}, 30^\circ\text{C}, 32^\circ\text{C}, 31^\circ\text{C}, 29^\circ\text{C}$. Nhiệt độ trung bình trong $5$ ngày là
\choice
{$29^\circ\text{C}$}
{\True $30^\circ\text{C}$}
{$31^\circ\text{C}$}
{$30{,}5^\circ\text{C}$}
\loigiai{
Tổng nhiệt độ: $28 + 30 + 32 + 31 + 29 = 150^\circ\text{C}$.\\
Trung bình: $150 : 5 = 30^\circ\text{C}$.\\
\textbf{Đáp án của câu hỏi là B}
}
\end{ex}

\begin{ex} %[5X11V66]
Trong dự án “Nuôi heo đất giúp bạn vượt khó”, bốn tổ của lớp 5A quyên góp được số tiền lần lượt là: $250\,000$ đồng, $300\,000$ đồng, $280\,000$ đồng, $370\,000$ đồng. Trung bình mỗi tổ quyên góp được số tiền là
\choice
{\True $300\,000$ đồng}
{$280\,000$ đồng}
{$320\,000$ đồng}
{$350\,000$ đồng}
\loigiai{
Tổng tiền: $250 + 300 + 280 + 370 = 1\,200$ nghìn đồng.\\
Trung bình: $1\,200 : 4 = 300$ nghìn đồng $= 300\,000$ đồng.\\
\textbf{Đáp án của câu hỏi là A}
}
\end{ex}

\begin{ex} %[5X11V66]
Thống kê thời gian tự học ở nhà (phút/ngày) của $10$ học sinh: $45, 60, 50, 60, 45, 90, 60, 75, 60, 45$. Tỉ số phần trăm học sinh dành đúng $60$ phút tự học là
\choice
{$30\%$}
{$50\%$}
{$20\%$}
{\True $40\%$}
\loigiai{
Có 4 bạn dành đúng 60 phút $\Rightarrow 4 : 10 = 40\%$.\\
\textbf{Đáp án của câu hỏi là D}
}
\end{ex}

\begin{ex} %[5X11C66]
Một cuộc khảo sát $100$ người tiêu dùng về $3$ nhãn hiệu sữa A, B, C cho thấy: $45$ người thích nhãn hiệu A, $35$ người thích nhãn hiệu B, số còn lại thích nhãn hiệu C. Khi vẽ biểu đồ quạt tròn, góc ở tâm biểu diễn phần số người thích nhãn hiệu C là bao nhiêu độ?
\choice
{$60^\circ$}
{$80^\circ$}
{\True $72^\circ$}
{$90^\circ$}
\loigiai{
Số người thích C: $100 - (45 + 35) = 20$ người ($20\%$).\\
Góc ở tâm: $360^\circ \times 20\% = 72^\circ$.\\
\textbf{Đáp án của câu hỏi là C}
}
\end{ex}

\textbf{Phần 2. Trắc nghiệm đúng sai}
\setcounterpageref{ex}{0}

\begin{ex} %[5X11H66]
Bảng khảo sát món tráng miệng ưa thích của $40$ học sinh: Kem ($16$ bạn), Chè ($12$ bạn), Trái cây ($8$ bạn), Sữa chua ($4$ bạn). Xét tính đúng sai:
\choiceTF
{\True Tỉ số phần trăm học sinh thích Kem là $40\%$.}
{\True Tỉ số phần trăm học sinh thích Chè là $30\%$.}
{Số học sinh thích Sữa chua chiếm $5\%$.}
{\True Số học sinh thích Sữa chua chiếm $10\%$.}
\loigiai{
\begin{itemize}
    \item Kem: $16 : 40 = 40\%$. Đúng.
    \item Chè: $12 : 40 = 30\%$. Đúng.
    \item Sữa chua: $4 : 40 = 10\% \neq 5\%$. Ý 3 Sai, Ý 4 Đúng.
\end{itemize}
\begin{itemize}
    \item a) \textbf{Đúng}.
    \item b) \textbf{Đúng}.
    \item c) \textbf{Sai}.
    \item d) \textbf{Đúng}.
\end{itemize}
}
\end{ex}

\begin{ex} %[5X11V66]
Một nhóm bạn đo lượng nước uống mỗi ngày của $5$ thành viên: $1{,}5\text{ l}, 1{,}8\text{ l}, 2{,}0\text{ l}, 1{,}7\text{ l}, 2{,}0\text{ l}$. Xét tính đúng sai:
\choiceTF
{\True Tổng lượng nước $5$ bạn uống trong ngày là $9{,}0\text{ l}$.}
{\True Lượng nước uống trung bình của mỗi bạn là $1{,}8\text{ l}$.}
{Có $1$ bạn uống lượng nước nhiều hơn mức trung bình.}
{\True Có $2$ bạn uống lượng nước nhiều hơn mức trung bình ($2{,}0\text{ l} > 1{,}8\text{ l}$).}
\loigiai{
\begin{itemize}
    \item Tổng: $1{,}5 + 1{,}8 + 2{,}0 + 1{,}7 + 2{,}0 = 9{,}0\text{ l}$. Đúng.
    \item Trung bình: $9{,}0 : 5 = 1{,}8\text{ l}$. Đúng.
    \item Có 2 bạn uống 2,0 l $> 1,8$ l. Ý 3 Sai, Ý 4 Đúng.
\end{itemize}
\begin{itemize}
    \item a) \textbf{Đúng}.
    \item b) \textbf{Đúng}.
    \item c) \textbf{Sai}.
    \item d) \textbf{Đúng}.
\end{itemize}
}
\end{ex}

\textbf{Phần 3. Trả lời ngắn}
\setcounterpageref{ex}{0}

\begin{ex} %[5X11H66]
Khảo sát số thành viên của $5$ gia đình: $3, 4, 5, 4, 4$. Tính số thành viên trung bình của mỗi gia đình.
\shortans{4}
\loigiai{
\textbf{Đáp số:} $4$\\
$(3 + 4 + 5 + 4 + 4) : 5 = 20 : 5 = 4$.
}
\end{ex}

\begin{ex} %[5X11H66]
Một bạn gieo đồng xu $25$ lần thấy mặt Sấp xuất hiện $13$ lần. Tính số lần xuất hiện mặt Ngửa.
\shortans{12}
\loigiai{
\textbf{Đáp số:} $12$\\
$25 - 13 = 12$.
}
\end{ex}

\begin{ex} %[5X11V66]
Một lớp có $35$ học sinh. Qua khảo sát có $28$ bạn thích xem hoạt hình. Tính tỉ số phần trăm học sinh thích xem hoạt hình. (Chỉ ghi số).
\shortans{80}
\loigiai{
\textbf{Đáp số:} $80$\\
$28 : 35 = 0{,}8 = 80\%$.
}
\end{ex}

\begin{ex} %[5X11V66]
Khảo sát số cây hoa trồng trong ban công của $4$ nhà: $8, 12, 15, 9$. Tính số cây hoa trung bình mỗi nhà trồng được.
\shortans{11}
\loigiai{
\textbf{Đáp số:} $11$\\
$(8 + 12 + 15 + 9) : 4 = 44 : 4 = 11$.
}
\end{ex}

\begin{ex} %[5X11V66]
Một nhóm học sinh quyên góp sách cũ: Tổ 1 góp $45$ cuốn, Tổ 2 góp $55$ cuốn, Tổ 3 góp $60$ cuốn, Tổ 4 góp $40$ cuốn. Tính số sách trung bình mỗi tổ quyên góp được.
\shortans{50}
\loigiai{
\textbf{Đáp số:} $50$\\
$(45 + 55 + 60 + 40) : 4 = 200 : 4 = 50$.
}
\end{ex}

\begin{ex} %[5X11C66]
Điểm kiểm tra của $4$ bạn là: $7, 8, 9, 8$. Hỏi bạn thứ năm cần đạt bao nhiêu điểm để điểm trung bình của cả nhóm $5$ bạn đạt $8{,}4$?
\shortans{10}
\loigiai{
\textbf{Đáp số:} $10$\\
Tổng điểm cần: $8{,}4 \times 5 = 42$.\\
Tổng 4 bạn đầu: $7 + 8 + 9 + 8 = 32$.\\
Điểm bạn thứ năm: $42 - 32 = 10$.
}
\end{ex}

\newpage

\TaoPhieuKiemTra{Chủ đề 11 -- Một số yếu tố thống kê và xác suất}{67}{Luyện tập chung thống kê, xác suất}{https://loptoanthaylinh.vn/thi-online/toan-5/bai-67}
\textbf{Phần 1. Trắc nghiệm lựa chọn}
\setcounterpageref{ex}{0}

\begin{ex} %[5X11N67]
Biểu đồ hình quạt tròn biểu diễn toàn bộ tổng thể tương ứng với bao nhiêu phần trăm?
\choice
{\True $100\%$}
{$50\%$}
{$200\%$}
{$360\%$}
\loigiai{
Toàn bộ hình quạt tròn biểu diễn $100\%$.\\
\textbf{Đáp án của câu hỏi là A}
}
\end{ex}

\begin{ex} %[5X11N67]
Gieo một con xúc xắc $30$ lần, mặt $5$ chấm xuất hiện $6$ lần. Tỉ số của số lần xuất hiện mặt $5$ chấm là
\choice
{$\dfrac{5}{30}$}
{$\dfrac{1}{6}$}
{\True $\dfrac{1}{5}$}
{$\dfrac{2}{5}$}
\loigiai{
$\dfrac{6}{30} = \dfrac{1}{5}$.\\
\textbf{Đáp án của câu hỏi là C}
}
\end{ex}

\begin{ex} %[5X11H67]
Biểu đồ hình quạt tròn cho biết cơ cấu phương tiện đến trường của $400$ học sinh: Xe đạp $50\%$, Đi bộ $30\%$, Xe buýt $20\%$. Số học sinh đi xe buýt là
\choice
{$60$ bạn}
{\True $80$ bạn}
{$100$ bạn}
{$120$ bạn}
\loigiai{
$400 \times 20 : 100 = 80$ (bạn).\\
\textbf{Đáp án của câu hỏi là B}
}
\end{ex}

\begin{ex} %[5X11H67]
Tung một đồng xu $50$ lần có $22$ lần xuất hiện mặt Sấp. Tỉ số phần trăm số lần xuất hiện mặt Ngửa là
\choice
{$44\%$}
{$50\%$}
{$48\%$}
{\True $56\%$}
\loigiai{
Mặt Ngửa xuất hiện: $50 - 22 = 28$ lần $\Rightarrow 28 : 50 = 56\%$.\\
\textbf{Đáp án của câu hỏi là D}
}
\end{ex}

\begin{ex} %[5X11H67]
Cho bảng số liệu điểm kiểm tra của một tổ: $9, 8, 7, 10, 8, 9, 8, 9$. Điểm xuất hiện nhiều lần nhất là
\choice
{$7$}
{\True $8$ (và $9$)}
{$9$}
{$10$}
\loigiai{
Điểm 8 và điểm 9 đều xuất hiện 3 lần.\\
\textbf{Đáp án của câu hỏi là B}
}
\end{ex}

\begin{ex} %[5X11H67]
Khảo sát sở thích đọc sách của $150$ học sinh: Truyện tranh $60\%$, Sách khoa học $25\%$, Sách lịch sử $15\%$. Số bạn thích sách lịch sử là
\choice
{$15$ bạn}
{$25$ bạn}
{$30$ bạn}
{\True $22{,}5$ bạn}
\loigiai{
$150 \times 15 : 100 = 22{,}5$ (bạn).\\
\textbf{Đáp án của câu hỏi là D}
}
\end{ex}

\begin{ex} %[5X11V67]
Gieo một con xúc xắc $100$ lần, mặt có số chấm là số nguyên tố ($2, 3, 5$) xuất hiện $52$ lần. Tỉ số phần trăm số lần xuất hiện mặt không phải số nguyên tố là
\choice
{\True $48\%$}
{$52\%$}
{$50\%$}
{$46\%$}
\loigiai{
Số lần không phải số nguyên tố: $100 - 52 = 48$ lần $\Rightarrow 48\%$.\\
\textbf{Đáp án của câu hỏi là A}
}
\end{ex}

\begin{ex} %[5X11V67]
Biểu đồ quạt biểu diễn diện tích trồng trọt: Lúa $55\%$, Ngô $25\%$, Khoai $20\%$. Biết diện tích trồng Lúa là $11\text{ ha}$. Tổng diện tích trồng trọt là
\choice
{$18\text{ ha}$}
{$22\text{ ha}$}
{\True $20\text{ ha}$}
{$25\text{ ha}$}
\loigiai{
Tổng diện tích: $11 : 55 \times 100 = 20\text{ (ha)}$.\\
\textbf{Đáp án của câu hỏi là C}
}
\end{ex}

\begin{ex} %[5X11V67]
Khảo sát số giờ sử dụng điện thoại mỗi ngày của $8$ người: $1\text{h}, 2\text{h}, 1{,}5\text{h}, 3\text{h}, 2{,}5\text{h}, 2\text{h}, 1\text{h}, 3\text{h}$. Số giờ sử dụng trung bình là
\choice
{$1{,}8\text{ giờ}$}
{\True $2\text{ giờ}$}
{$2{,}2\text{ giờ}$}
{$2{,}5\text{ giờ}$}
\loigiai{
Tổng: $1 + 2 + 1{,}5 + 3 + 2{,}5 + 2 + 1 + 3 = 16$ giờ.\\
Trung bình: $16 : 8 = 2$ giờ.\\
\textbf{Đáp án của câu hỏi là B}
}
\end{ex}

\begin{ex} %[5X11V67]
Một hộp có $5$ viên bi xanh, $3$ viên bi đỏ và $2$ viên bi vàng. Bạn Nam lấy ngẫu nhiên $1$ viên bi, ghi màu rồi trả lại hộp, làm như vậy $50$ lần thì có $24$ lần lấy được bi xanh. Tỉ số phần trăm thực nghiệm lấy được bi xanh là
\choice
{$40\%$}
{$50\%$}
{\True $48\%$}
{$52\%$}
\loigiai{
$24 : 50 = 48\%$.\\
\textbf{Đáp án của câu hỏi là C}
}
\end{ex}

\begin{ex} %[5X11V67]
Biểu đồ quạt tròn cho biết tỉ lệ bình chọn tiết mục văn nghệ: Múa $35\%$, Hát $40\%$, Kịch $25\%$. Biết có $80$ phiếu bình chọn cho tiết mục Hát. Tổng số phiếu bình chọn là
\choice
{\True $200$ phiếu}
{$160$ phiếu}
{$240$ phiếu}
{$250$ phiếu}
\loigiai{
Tổng số phiếu: $80 : 40 \times 100 = 200$ (phiếu).\\
\textbf{Đáp án của câu hỏi là A}
}
\end{ex}

\begin{ex} %[5X11C67]
Một cuộc điều tra về số lượng vật nuôi của $50$ hộ gia đình:
- $15$ hộ nuôi chó
- $20$ hộ nuôi mèo
- $5$ hộ nuôi cả chó và mèo
- Số hộ còn lại không nuôi con nào.
Tỉ số phần trăm số hộ gia đình không nuôi con vật nào là
\choice
{$30\%$}
{$35\%$}
{$25\%$}
{\True $40\%$}
\loigiai{
Số hộ có nuôi ít nhất 1 con vật: $15 + 20 - 5 = 30$ hộ.\\
Số hộ không nuôi con vật nào: $50 - 30 = 20$ hộ.\\
Tỉ số phần trăm: $20 : 50 = 40\%$.\\
\textbf{Đáp án của câu hỏi là D}
}
\end{ex}

\textbf{Phần 2. Trắc nghiệm đúng sai}
\setcounterpageref{ex}{0}

\begin{ex} %[5X11H67]
Biểu đồ hình quạt tròn biểu diễn kết quả phân loại rác thải tái chế: Giấy $40\%$, Nhựa $35\%$, Kim loại $15\%$, Thủy tinh $10\%$. Tổng khối lượng thu gom là $200\text{ kg}$. Xét tính đúng sai:
\choiceTF
{\True Khối lượng giấy thu gom được là $80\text{ kg}$.}
{\True Khối lượng nhựa thu gom được là $70\text{ kg}$.}
{Khối lượng kim loại thu gom được là $40\text{ kg}$.}
{\True Khối lượng thủy tinh thu gom được là $20\text{ kg}$.}
\loigiai{
\begin{itemize}
    \item Giấy: $200 \times 40\% = 80\text{ kg}$. Đúng.
    \item Nhựa: $200 \times 35\% = 70\text{ kg}$. Đúng.
    \item Kim loại: $200 \times 15\% = 30\text{ kg} \neq 40\text{ kg}$. Sai.
    \item Thủy tinh: $200 \times 10\% = 20\text{ kg}$. Đúng.
\end{itemize}
\begin{itemize}
    \item a) \textbf{Đúng}.
    \item b) \textbf{Đúng}.
    \item c) \textbf{Sai}.
    \item d) \textbf{Đúng}.
\end{itemize}
}
\end{ex}

\begin{ex} %[5X11V67]
Tung một con xúc xắc $80$ lần. Kết quả: Mặt $1$ chấm ($12$ lần), Mặt $2$ chấm ($14$ lần), Mặt $3$ chấm ($10$ lần), Mặt $4$ chấm ($16$ lần), Mặt $5$ chấm ($18$ lần), Mặt $6$ chấm ($10$ lần). Xét tính đúng sai:
\choiceTF
{\True Số lần xuất hiện mặt chẵn là $40$ lần ($50\%$).}
{\True Số lần xuất hiện mặt lẻ là $40$ lần ($50\%$).}
{Mặt có số lần xuất hiện nhiều nhất là mặt $4$ chấm.}
{\True Mặt có số lần xuất hiện nhiều nhất là mặt $5$ chấm ($18$ lần).}
\loigiai{
\begin{itemize}
    \item Mặt chẵn (2, 4, 6): $14 + 16 + 10 = 40$ lần ($50\%$). Đúng.
    \item Mặt lẻ (1, 3, 5): $12 + 10 + 18 = 40$ lần ($50\%$). Đúng.
    \item Mặt 5 chấm xuất hiện 18 lần, nhiều nhất (không phải 4 chấm). Ý 3 Sai, Ý 4 Đúng.
\end{itemize}
\begin{itemize}
    \item a) \textbf{Đúng}.
    \item b) \textbf{Đúng}.
    \item c) \textbf{Sai}.
    \item d) \textbf{Đúng}.
\end{itemize}
}
\end{ex}

\textbf{Phần 3. Trả lời ngắn}
\setcounterpageref{ex}{0}

\begin{ex} %[5X11H67]
Trên biểu đồ quạt tròn, phần thể thao chiếm $35\%$, nghệ thuật chiếm $25\%$, phần còn lại là học thuật. Hỏi phần học thuật chiếm bao nhiêu phần trăm? (Chỉ ghi số).
\shortans{40}
\loigiai{
\textbf{Đáp số:} $40$\\
$100 - (35 + 25) = 40\%$.
}
\end{ex}

\begin{ex} %[5X11H67]
Gieo xúc xắc $50$ lần, có $9$ lần xuất hiện mặt 1 chấm. Tính tỉ số phần trăm số lần xuất hiện mặt 1 chấm. (Chỉ ghi số).
\shortans{18}
\loigiai{
\textbf{Đáp số:} $18$\\
$9 : 50 = 18\%$.
}
\end{ex}

\begin{ex} %[5X11V67]
Biểu đồ quạt tròn cho biết $30\%$ số học sinh của trường tham gia CLB Tiếng Anh, tương ứng $180$ học sinh. Tính tổng số học sinh của trường.
\shortans{600}
\loigiai{
\textbf{Đáp số:} $600$\\
$180 : 30 \times 100 = 600$.
}
\end{ex}

\begin{ex} %[5X11V67]
Tung đồng xu $80$ lần thấy xuất hiện mặt Sấp $36$ lần. Tính tỉ số phần trăm số lần xuất hiện mặt Ngửa. (Chỉ ghi số).
\shortans{55}
\loigiai{
\textbf{Đáp số:} $55$\\
Mặt Ngửa: $80 - 36 = 44$ lần $\Rightarrow 44 : 80 = 55\%$.
}
\end{ex}

\begin{ex} %[5X11V67]
Điểm kiểm tra của $5$ bạn học sinh là: $8, 9, 7, 10, 6$. Tính điểm trung bình của $5$ bạn.
\shortans{8}
\loigiai{
\textbf{Đáp số:} $8$\\
$(8 + 9 + 7 + 10 + 6) : 5 = 40 : 5 = 8$.
}
\end{ex}

\begin{ex} %[5X11C67]
Biểu đồ quạt tròn cho thấy $45\%$ số cây trong vườn là cam, $35\%$ là xoài, còn lại là $60$ cây bưởi. Tính tổng số cây có trong vườn.
\shortans{300}
\loigiai{
\textbf{Đáp số:} $300$\\
Tỉ lệ bưởi: $100\% - (45\% + 35\%) = 20\%$.\\
Tổng số cây: $60 : 20 \times 100 = 300$ cây.
}
\end{ex}

\newpage

% =========================================================================
% CHỦ ĐỀ 12 - ÔN TẬP CUỐI NĂM
% =========================================================================
\section{Ôn tập cuối năm}

\TaoPhieuKiemTra{Chủ đề 12 -- Ôn tập cuối năm}{68}{Ôn tập số tự nhiên, phân số, số thập phân}{https://loptoanthaylinh.vn/thi-online/toan-5/bai-68}
\textbf{Phần 1. Trắc nghiệm lựa chọn}
\setcounterpageref{ex}{0}

\begin{ex} %[5D12N68]
Số gồm: \textit{Bảy triệu, hai trăm linh năm nghìn, sáu trăm ba mươi đơn vị} viết là
\choice
{\True $7\,205\,630$}
{$7\,250\,630$}
{$7\,025\,630$}
{$7\,256\,300$}
\loigiai{
Số đó là $7\,205\,630$.\\
\textbf{Đáp án của câu hỏi là A}
}
\end{ex}

\begin{ex} %[5D12N68]
Trong số thập phân $45{,}678$, chữ số $7$ có giá trị là
\choice
{$7$}
{$0{,}7$}
{$0{,}007$}
{\True $0{,}07$}
\loigiai{
Chữ số 7 ở hàng phần trăm nên có giá trị là $0{,}07$ hay $\dfrac{7}{100}$.\\
\textbf{Đáp án của câu hỏi là D}
}
\end{ex}

\begin{ex} %[5D12H68]
Phân số nào sau đây là phân số tối giản?
\choice
{$\dfrac{6}{9}$}
{\True $\dfrac{8}{15}$}
{$\dfrac{15}{20}$}
{$\dfrac{14}{21}$}
\loigiai{
$\dfrac{8}{15}$ có ƯCLN$(8, 15) = 1$ nên là phân số tối giản.\\
\textbf{Đáp án của câu hỏi là B}
}
\end{ex}

\begin{ex} %[5D12H68]
Phân số $\dfrac{3}{4}$ được viết dưới dạng số thập phân là
\choice
{$0{,}34$}
{$0{,}43$}
{\True $0{,}75$}
{$3{,}4$}
\loigiai{
$\dfrac{3}{4} = 3 : 4 = 0{,}75$.\\
\textbf{Đáp án của câu hỏi là C}
}
\end{ex}

\begin{ex} %[5D12H68]
Hỗn số $3\dfrac{2}{5}$ viết dưới dạng phân số là
\choice
{$\dfrac{6}{5}$}
{\True $\dfrac{17}{5}$}
{$\dfrac{15}{5}$}
{$\dfrac{32}{5}$}
\loigiai{
$3\dfrac{2}{5} = \dfrac{3 \times 5 + 2}{5} = \dfrac{17}{5}$.\\
\textbf{Đáp án của câu hỏi là B}
}
\end{ex}

\begin{ex} %[5D12H68]
Sắp xếp các số sau theo thứ tự từ bé đến lớn: $4{,}05; 4{,}5; 4{,}15; 4{,}005$.
\choice
{$4{,}5 < 4{,}15 < 4{,}05 < 4{,}005$}
{$4{,}05 < 4{,}005 < 4{,}15 < 4{,}5$}
{$4{,}005 < 4{,}15 < 4{,}05 < 4{,}5$}
{\True $4{,}005 < 4{,}05 < 4{,}15 < 4{,}5$}
\loigiai{
Thứ tự: $4{,}005 < 4{,}05 < 4{,}15 < 4{,}5$.\\
\textbf{Đáp án của câu hỏi là D}
}
\end{ex}

\begin{ex} %[5D12V68]
Tìm chữ số $x$ biết $23{,}45 < 23{,}x1 < 23{,}62$:
\choice
{$x = 4$}
{\True $x = 5$ hoặc $x = 6$}
{$x = 5$}
{$x = 6$}
\loigiai{
$23{,}45 < 23{,}51 < 23{,}61 < 23{,}62$ nên $x = 5$ hoặc $x = 6$.\\
\textbf{Đáp án của câu hỏi là B}
}
\end{ex}

\begin{ex} %[5D12V68]
Rút gọn phân số $\dfrac{72}{108}$ về phân số tối giản ta được
\choice
{\True $\dfrac{2}{3}$}
{$\dfrac{3}{4}$}
{$\dfrac{4}{6}$}
{$\dfrac{6}{9}$}
\loigiai{
$\dfrac{72}{108} = \dfrac{72 : 36}{108 : 36} = \dfrac{2}{3}$.\\
\textbf{Đáp án của câu hỏi là A}
}
\end{ex}

\begin{ex} %[5D12V68]
Làm tròn số thập phân $67{,}846$ đến hàng phần mười ta được
\choice
{$67{,}8$}
{$67{,}85$}
{\True $67{,}8$}
{$67{,}9$}
\loigiai{
Chữ số hàng phần trăm là $4 < 5$ nên làm tròn thành $67{,}8$.\\
\textbf{Đáp án của câu hỏi là C}
}
\end{ex}

\begin{ex} %[5D12V68]
Có bao nhiêu số tự nhiên $x$ thỏa mãn $4{,}75 < x < 9{,}2$?
\choice
{\True $5$ số}
{$4$ số}
{$6$ số}
{$7$ số}
\loigiai{
Các số tự nhiên đó là: $5, 6, 7, 8, 9$ (có 5 số).\\
\textbf{Đáp án của câu hỏi là A}
}
\end{ex}

\begin{ex} %[5D12V68]
Tìm phân số tối giản $\dfrac{a}{b}$ lớn hơn $\dfrac{1}{3}$ nhưng bé hơn $\dfrac{1}{2}$ có mẫu số là $12$. Phân số đó là
\choice
{$\dfrac{4}{12}$}
{$\dfrac{6}{12}$}
{\True $\dfrac{5}{12}$}
{$\dfrac{7}{12}$}
\loigiai{
$\dfrac{1}{3} = \dfrac{4}{12} < \dfrac{5}{12} < \dfrac{6}{12} = \dfrac{1}{2}$.\\
\textbf{Đáp án của câu hỏi là C}
}
\end{ex}

\begin{ex} %[5D12C68]
Cho phân số $\dfrac{19}{35}$. Cần bớt ở tử số và thêm vào mẫu số cùng một số tự nhiên $k$ để được phân số mới bằng $\dfrac{1}{2}$. Số tự nhiên $k$ là
\choice
{$2$}
{$3$}
{$4$}
{\True $1$}
\loigiai{
Tổng tử số và mẫu số không đổi: $19 + 35 = 54$.\\
Tử số mới: $54 : (1 + 2) \times 1 = 18$.\\
Số $k = 19 - 18 = 1$.\\
\textbf{Đáp án của câu hỏi là D}
}
\end{ex}

\textbf{Phần 2. Trắc nghiệm đúng sai}
\setcounterpageref{ex}{0}

\begin{ex} %[5D12H68]
Xét tính đúng sai về cấu tạo và so sánh số:
\choiceTF
{\True Số $0$ là số tự nhiên nhỏ nhất.}
{Phân số $\dfrac{12}{15}$ là phân số tối giản.}
{\True $3{,}4 = 3{,}40 = 3{,}400$.}
{\True Trong số $85{,}125$, chữ số $5$ ở phần nguyên có giá trị gấp $1\,000$ lần chữ số $5$ ở phần thập phân.}
\loigiai{
\begin{itemize}
    \item 0 là số tự nhiên nhỏ nhất. Đúng.
    \item $\dfrac{12}{15} = \dfrac{4}{5}$ chưa tối giản. Sai.
    \item Viết thêm chữ số 0 vào bên phải phần thập phân giá trị không đổi. Đúng.
    \item 5 đơn vị gấp 0,005 số lần: $5 : 0{,}005 = 1000$ lần. Đúng.
\end{itemize}
\begin{itemize}
    \item a) \textbf{Đúng}.
    \item b) \textbf{Sai}.
    \item c) \textbf{Đúng}.
    \item d) \textbf{Đúng}.
\end{itemize}
}
\end{ex}

\begin{ex} %[5D12V68]
Cho các số: $2\dfrac{3}{4}; 2{,}7; \dfrac{11}{4}; 2{,}75$. Xét tính đúng sai:
\choiceTF
{\True $2\dfrac{3}{4} = \dfrac{11}{4}$.}
{\True $\dfrac{11}{4} = 2{,}75$.}
{$2{,}7 > 2{,}75$.}
{\True Sắp xếp tăng dần: $2{,}7 < 2{,}75 = 2\dfrac{3}{4} = \dfrac{11}{4}$.}
\loigiai{
\begin{itemize}
    \item $2\dfrac{3}{4} = \dfrac{11}{4}$. Đúng.
    \item $11 : 4 = 2{,}75$. Đúng.
    \item $2{,}7 < 2{,}75$. Khẳng định lớn hơn là Sai.
    \item Thứ tự tăng dần đúng.
\end{itemize}
\begin{itemize}
    \item a) \textbf{Đúng}.
    \item b) \textbf{Đúng}.
    \item c) \textbf{Sai}.
    \item d) \textbf{Đúng}.
\end{itemize}
}
\end{ex}

\textbf{Phần 3. Trả lời ngắn}
\setcounterpageref{ex}{0}

\begin{ex} %[5D12H68]
Viết phân số $\dfrac{3}{8}$ dưới dạng số thập phân.
\shortans{0,375}
\loigiai{
\textbf{Đáp số:} $0{,}375$\\
$3 : 8 = 0{,}375$.
}
\end{ex}

\begin{ex} %[5D12H68]
Rút gọn phân số $\dfrac{48}{60}$ về phân số tối giản $\dfrac{a}{b}$. Tính $a + b$.
\shortans{9}
\loigiai{
\textbf{Đáp số:} $9$\\
$\dfrac{48}{60} = \dfrac{4}{5} \Rightarrow a = 4, b = 5 \Rightarrow a + b = 9$.
}
\end{ex}

\begin{ex} %[5D12V68]
Tìm số tự nhiên $x$ lớn nhất thỏa mãn $x < 15{,}89$.
\shortans{15}
\loigiai{
\textbf{Đáp số:} $15$\\
Số tự nhiên lớn nhất nhỏ hơn 15,89 là $15$.
}
\end{ex}

\begin{ex} %[5D12V68]
Có bao nhiêu chữ số ở phần thập phân của số $12{,}3456$?
\shortans{4}
\loigiai{
\textbf{Đáp số:} $4$\
Có $4$ chữ số ($3, 4, 5, 6$).
}
\end{ex}

\begin{ex} %[5D12V68]
Tìm phân số bằng phân số $\dfrac{2}{5}$ có tổng của tử số và mẫu số bằng $63$. Tìm tử số của phân số đó.
\shortans{18}
\loigiai{
\textbf{Đáp số:} $18$\\
Tử số: $63 : (2 + 5) \times 2 = 18$.
}
\end{ex}

\begin{ex} %[5D12C68]
Viết số thập phân bé nhất có $4$ chữ số khác nhau trong đó có hai chữ số ở phần thập phân.
\shortans{10,23}
\loigiai{
\textbf{Đáp số:} $10{,}23$\\
Số nhỏ nhất có 4 chữ số khác nhau với 2 chữ số thập phân là $10{,}23$.
}
\end{ex}

\newpage

\TaoPhieuKiemTra{Chủ đề 12 -- Ôn tập cuối năm}{69}{Ôn tập các phép tính}{https://loptoanthaylinh.vn/thi-online/toan-5/bai-69}
\textbf{Phần 1. Trắc nghiệm lựa chọn}
\setcounterpageref{ex}{0}

\begin{ex} %[5D12N69]
Kết quả của phép tính $45{,}6 + 18{,}75$ là
\choice
{$63{,}35$}
{\True $64{,}35$}
{$64{,}45$}
{$63{,}45$}
\loigiai{
$45{,}6 + 18{,}75 = 64{,}35$.\\
\textbf{Đáp án của câu hỏi là B}
}
\end{ex}

\begin{ex} %[5D12N69]
Kết quả của phép tính $\dfrac{2}{3} + \dfrac{1}{6}$ là
\choice
{$\dfrac{3}{9}$}
{$\dfrac{4}{6}$}
{\True $\dfrac{5}{6}$}
{$1$}
\loigiai{
$\dfrac{2}{3} + \dfrac{1}{6} = \dfrac{4}{6} + \dfrac{1}{6} = \dfrac{5}{6}$.\\
\textbf{Đáp án của câu hỏi là C}
}
\end{ex}

\begin{ex} %[5D12H69]
Kết quả của phép nhân $3{,}25 \times 4$ là
\choice
{$12{,}5$}
{\True $13$}
{$13{,}25$}
{$12$}
\loigiai{
$3{,}25 \times 4 = 13$.\\
\textbf{Đáp án của câu hỏi là B}
}
\end{ex}

\begin{ex} %[5D12H69]
Kết quả của phép tính $18{,}6 : 0{,}1$ là
\choice
{$1{,}86$}
{$0{,}186$}
{\True $186$}
{$18{,}6$}
\loigiai{
Chia cho $0{,}1$ tức là nhân với $10$: $18{,}6 \times 10 = 186$.\\
\textbf{Đáp án của câu hỏi là C}
}
\end{ex}

\begin{ex} %[5D12H69]
Tính nhanh giá trị biểu thức: $3{,}6 \times 45 + 3{,}6 \times 55$.
\choice
{$36$}
{$3,600$}
{$180$}
{\True $360$}
\loigiai{
$3{,}6 \times (45 + 55) = 3{,}6 \times 100 = 360$.\\
\textbf{Đáp án của câu hỏi là D}
}
\end{ex}

\begin{ex} %[5D12H69]
Kết quả của phép chia $\dfrac{4}{5} : \dfrac{2}{3}$ là
\choice
{\True $\dfrac{6}{5}$}
{$\dfrac{8}{15}$}
{$\dfrac{5}{6}$}
{$\dfrac{2}{5}$}
\loigiai{
$\dfrac{4}{5} \times \dfrac{3}{2} = \dfrac{12}{10} = \dfrac{6}{5}$.\\
\textbf{Đáp án của câu hỏi là A}
}
\end{ex}

\begin{ex} %[5D12V69]
Tìm $x$ biết: $x \times 2{,}5 = 48{,}5 - 16$.
\choice
{$12$}
{$14$}
{$15$}
{\True $13$}
\loigiai{
$x \times 2{,}5 = 32{,}5 \Rightarrow x = 32{,}5 : 2{,}5 = 13$.\\
\textbf{Đáp án của câu hỏi là D}
}
\end{ex}

\begin{ex} %[5D12V69]
Giá trị của biểu thức $(12{,}4 - 4{,}9) \times 3{,}2 + 6$ là
\choice
{\True $30$}
{$28$}
{$32$}
{$34$}
\loigiai{
$7{,}5 \times 3{,}2 + 6 = 24 + 6 = 30$.\\
\textbf{Đáp án của câu hỏi là A}
}
\end{ex}

\begin{ex} %[5D12V69]
Một mảnh đất hình chữ nhật có chiều dài $24{,}5\text{ m}$, chiều rộng $12\text{ m}$. Diện tích mảnh đất đó là
\choice
{$284\text{ m}^2$}
{$304\text{ m}^2$}
{\True $294\text{ m}^2$}
{$290\text{ m}^2$}
\loigiai{
$24{,}5 \times 12 = 294\text{ (m}^2\text{)}$.\\
\textbf{Đáp án của câu hỏi là C}
}
\end{ex}

\begin{ex} %[5D12V69]
Tính giá trị của biểu thức: $\dfrac{3}{4} \times \dfrac{8}{9} + \dfrac{1}{3}$.
\choice
{\True $1$}
{$\dfrac{2}{3}$}
{$\dfrac{4}{3}$}
{$\dfrac{5}{6}$}
\loigiai{
$\dfrac{2}{3} + \dfrac{1}{3} = 1$.\\
\textbf{Đáp án của câu hỏi là A}
}
\end{ex}

\begin{ex} %[5D12V69]
Một người mua $3{,}5\text{ kg}$ đường loại $24\,000$ đồng/kg và $2\text{ kg}$ bột mì loại $18\,000$ đồng/kg. Người đó đưa cho người bán hàng tờ $200\,000$ đồng. Số tiền thừa nhận lại là
\choice
{$70\,000$ đồng}
{$90\,000$ đồng}
{$85\,000$ đồng}
{\True $80\,000$ đồng}
\loigiai{
Tổng tiền mua: $3{,}5 \times 24\,000 + 2 \times 18\,000 = 84\,000 + 36\,000 = 120\,000$ (đồng).\\
Tiền thừa: $200\,000 - 120\,000 = 80\,000$ (đồng).\\
\textbf{Đáp án của câu hỏi là D}
}
\end{ex}

\begin{ex} %[5D12C69]
Tính giá trị biểu thức bằng cách thuận tiện: $\dfrac{0{,}25 \times 61{,}4 \times 40}{0{,}5 \times 0{,}2 \times 100}$. Kết quả là
\choice
{$6{,}14$}
{\True $61{,}4$}
{$614$}
{$0{,}614$}
\loigiai{
Tử số: $(0{,}25 \times 40) \times 61{,}4 = 10 \times 61{,}4 = 614$.\\
Mẫu số: $(0{,}5 \times 0{,}2) \times 100 = 0{,}1 \times 100 = 10$.\\
Kết quả: $614 : 10 = 61{,}4$.\\
\textbf{Đáp án của câu hỏi là B}
}
\end{ex}

\textbf{Phần 2. Trắc nghiệm đúng sai}
\setcounterpageref{ex}{0}

\begin{ex} %[5D12H69]
Xét tính đúng sai của các phép tính sau:
\choiceTF
{\True $15{,}4 \times 100 = 1\,540$.}
{$45{,}8 : 0{,}01 = 4{,}58$.}
{\True $\dfrac{5}{7} \times \dfrac{14}{15} = \dfrac{2}{3}$.}
{\True $0{,}8 : 0{,}25 = 3{,}2$.}
\loigiai{
\begin{itemize}
    \item $15{,}4 \times 100 = 1540$. Đúng.
    \item $45{,}8 : 0{,}01 = 4580 \neq 4{,}58$. Sai.
    \item $\dfrac{5 \times 14}{7 \times 15} = \dfrac{2}{3}$. Đúng.
    \item $0{,}8 : 0{,}25 = 3{,}2$. Đúng.
\end{itemize}
\begin{itemize}
    \item a) \textbf{Đúng}.
    \item b) \textbf{Sai}.
    \item c) \textbf{Đúng}.
    \item d) \textbf{Đúng}.
\end{itemize}
}
\end{ex}

\begin{ex} %[5D12V69]
Cho biểu thức $A = 4{,}8 \times 3{,}5 + 4{,}8 \times 6{,}5$ và $B = (12{,}6 - 2{,}6) \times 4{,}8$. Xét tính đúng sai:
\choiceTF
{\True Giá trị biểu thức $A = 48$.}
{\True Giá trị biểu thức $B = 48$.}
{Giá trị của biểu thức $A$ lớn hơn giá trị của biểu thức $B$.}
{\True $A = B$.}
\loigiai{
\begin{itemize}
    \item $A = 4{,}8 \times (3{,}5 + 6{,}5) = 4{,}8 \times 10 = 48$. Đúng.
    \item $B = 10 \times 4{,}8 = 48$. Đúng.
    \item $A = B = 48$ nên khẳng định $A > B$ là Sai.
    \item $A = B$. Đúng.
\end{itemize}
\begin{itemize}
    \item a) \textbf{Đúng}.
    \item b) \textbf{Đúng}.
    \item c) \textbf{Sai}.
    \item d) \textbf{Đúng}.
\end{itemize}
}
\end{ex}

\textbf{Phần 3. Trả lời ngắn}
\setcounterpageref{ex}{0}

\begin{ex} %[5D12H69]
Tính kết quả phép tính: $45{,}8 + 34{,}75$.
\shortans{80,55}
\loigiai{
\textbf{Đáp số:} $80{,}55$\\
$45{,}8 + 34{,}75 = 80{,}55$.
}
\end{ex}

\begin{ex} %[5D12H69]
Tính kết quả: $12{,}5 \times 0{,}8$.
\shortans{10}
\loigiai{
\textbf{Đáp số:} $10$\\
$12{,}5 \times 0{,}8 = 10$.
}
\end{ex}

\begin{ex} %[5D12V69]
Tính giá trị biểu thức: $150 - 24{,}5 \times 4$.
\shortans{52}
\loigiai{
\textbf{Đáp số:} $52$\\
$150 - 98 = 52$.
}
\end{ex}

\begin{ex} %[5D12V69]
Tìm $x$ biết: $x : 3{,}6 = 12{,}5$.
\shortans{45}
\loigiai{
\textbf{Đáp số:} $45$\\
$x = 12{,}5 \times 3{,}6 = 45$.
}
\end{ex}

\begin{ex} %[5D12V69]
Tính nhanh: $2{,}5 \times 17{,}8 \times 4$.
\shortans{178}
\loigiai{
\textbf{Đáp số:} $178$\\
$(2{,}5 \times 4) \times 17{,}8 = 10 \times 17{,}8 = 178$.
}
\end{ex}

\begin{ex} %[5D12C69]
Một kho hàng có $120$ tấn gạo. Ngày thứ nhất xuất $\dfrac{1}{3}$ số gạo, ngày thứ hai xuất $40\%$ số gạo còn lại. Hỏi sau hai ngày kho hàng còn lại bao nhiêu tấn gạo?
\shortans{48}
\loigiai{
\textbf{Đáp số:} $48$\\
Ngày 1 xuất: $120 : 3 = 40$ tấn. Còn lại: $120 - 40 = 80$ tấn.\\
Ngày 2 xuất: $80 \times 40\% = 32$ tấn.\\
Còn lại sau 2 ngày: $80 - 32 = 48$ tấn.
}
\end{ex}

\newpage

\TaoPhieuKiemTra{Chủ đề 12 -- Ôn tập cuối năm}{70}{Ôn tập tỉ số, tỉ số phần trăm}{https://loptoanthaylinh.vn/thi-online/toan-5/bai-70}
\textbf{Phần 1. Trắc nghiệm lựa chọn}
\setcounterpageref{ex}{0}

\begin{ex} %[5D12N70]
Tỉ số của hai số $4$ và $9$ là
\choice
{\True $\dfrac{4}{9}$}
{$\dfrac{9}{4}$}
{$4 \times 9$}
{$36$}
\loigiai{
Tỉ số là $\dfrac{4}{9}$.\\
\textbf{Đáp án của câu hỏi là A}
}
\end{ex}

\begin{ex} %[5D12N70]
Số thập phân $0{,}35$ viết thành tỉ số phần trăm là
\choice
{$3{,}5\%$}
{\True $35\%$}
{$350\%$}
{$0{,}35\%$}
\loigiai{
$0{,}35 = 35\%$.\\
\textbf{Đáp án của câu hỏi là B}
}
\end{ex}

\begin{ex} %[5D12H70]
Tìm $20\%$ của $150\text{ kg}$:
\choice
{$20\text{ kg}$}
{$40\text{ kg}$}
{\True $30\text{ kg}$}
{$50\text{ kg}$}
\loigiai{
$150 \times 20 : 100 = 30\text{ (kg)}$.\\
\textbf{Đáp án của câu hỏi là C}
}
\end{ex}

\begin{ex} %[5D12H70]
Một lớp học có $40$ học sinh, trong đó có $24$ bạn nữ. Tỉ số phần trăm học sinh nam so với cả lớp là
\choice
{$60\%$}
{\True $40\%$}
{$50\%$}
{$45\%$}
\loigiai{
Số bạn nam: $40 - 24 = 16$ (bạn).\\
Tỉ số phần trăm: $16 : 40 = 40\%$.\\
\textbf{Đáp án của câu hỏi là B}
}
\end{ex}

\begin{ex} %[5D12H70]
Tổng của hai số là $120$, tỉ số của hai số là $\dfrac{1}{3}$. Số bé là
\choice
{$40$}
{$60$}
{$90$}
{\True $30$}
\loigiai{
$120 : (1 + 3) \times 1 = 30$.\\
\textbf{Đáp án của câu hỏi là D}
}
\end{ex}

\begin{ex} %[5D12H70]
Hiệu của hai số là $45$, biết số lớn gấp $4$ lần số bé. Số lớn là
\choice
{$15$}
{\True $60$}
{$45$}
{$75$}
\loigiai{
Số lớn: $45 : (4 - 1) \times 4 = 60$.\\
\textbf{Đáp án của câu hỏi là B}
}
\end{ex}

\begin{ex} %[5D12V70]
Trên bản đồ tỉ lệ $1 : 20,000$, quãng đường từ nhà An đến trường đo được $6\text{ cm}$. Khoảng cách thực tế là
\choice
{$120\text{ m}$}
{$12\text{ km}$}
{$120\text{ km}$}
{\True $1{,}2\text{ km}$}
\loigiai{
$6 \times 20,000 = 120,000\text{ cm} = 1{,}2\text{ km}$.\\
\textbf{Đáp án của câu hỏi là D}
}
\end{ex}

\begin{ex} %[5D12V70]
Một chiếc xe máy có giá $25\,000\,000$ đồng, nhân dịp khuyến mãi được giảm giá $8\%$. Số tiền khách hàng được giảm là
\choice
{$1\,500\,000$ đồng}
{$2\,500\,000$ đồng}
{\True $2\,000\,000$ đồng}
{$1\,800\,000$ đồng}
\loigiai{
$25\,000\,000 \times 8 : 100 = 2\,000\,000$ (đồng).\\
\textbf{Đáp án của câu hỏi là C}
}
\end{ex}

\begin{ex} %[5D12V70]
Một mảnh đất hình chữ nhật có chu vi $100\text{ m}$, chiều rộng bằng $\dfrac{2}{3}$ chiều dài. Diện tích mảnh đất đó là
\choice
{$500\text{ m}^2$}
{$700\text{ m}^2$}
{$800\text{ m}^2$}
{\True $600\text{ m}^2$}
\loigiai{
Nửa chu vi: $50\text{ m}$.\\
Chiều dài: $50 : (2 + 3) \times 3 = 30\text{ m}$.\\
Chiều rộng: $50 - 30 = 20\text{ m}$.\\
Diện tích: $30 \times 20 = 600\text{ (m}^2\text{)}$.\\
\textbf{Đáp án của câu hỏi là D}
}
\end{ex}

\begin{ex} %[5D12V70]
Một người gửi tiết kiệm $60\,000\,000$ đồng với lãi suất $6\%$ một năm. Sau một năm người đó nhận được số tiền lãi là
\choice
{\True $3\,600\,000$ đồng}
{$3\,000\,000$ đồng}
{$4\,000\,000$ đồng}
{$6\,000\,000$ đồng}
\loigiai{
$60\,000\,000 \times 6 : 100 = 3\,600\,000$ (đồng).\\
\textbf{Đáp án của câu hỏi là A}
}
\end{ex}

\begin{ex} %[5D12V70]
Biết $45\%$ của một số là $180$. Số đó là
\choice
{$81$}
{$360$}
{\True $400$}
{$450$}
\loigiai{
$180 : 45 \times 100 = 400$.\\
\textbf{Đáp án của câu hỏi là C}
}
\end{ex}

\begin{ex} %[5D12C70]
Một cửa hàng bán một chiếc đồng hồ với giá $600\,000$ đồng thì được lãi $20\%$ so với giá vốn. Giá vốn của chiếc đồng hồ đó là
\choice
{\True $500\,000$ đồng}
{$480\,000$ đồng}
{$520\,000$ đồng}
{$550\,000$ đồng}
\loigiai{
Giá bán chiếm $100\% + 20\% = 120\%$ giá vốn.\\
Giá vốn: $600\,000 : 120 \times 100 = 500\,000$ (đồng).\\
\textbf{Đáp án của câu hỏi là A}
}
\end{ex}

\textbf{Phần 2. Trắc nghiệm đúng sai}
\setcounterpageref{ex}{0}

\begin{ex} %[5D12H70]
Một nông trại có $300$ con vịt và $200$ con gà. Xét tính đúng sai:
\choiceTF
{\True Tỉ số giữa số vịt và số gà là $\dfrac{3}{2}$.}
{\True Số vịt chiếm $60\%$ tổng đàn gia cầm.}
{Số gà chiếm $30\%$ tổng đàn gia cầm.}
{\True Số gà chiếm $40\%$ tổng đàn gia cầm.}
\loigiai{
\begin{itemize}
    \item Tỉ số vịt và gà: $300 : 200 = \dfrac{3}{2}$. Đúng.
    \item Tổng đàn: $500$ con $\Rightarrow$ Vịt: $300 : 500 = 60\%$. Đúng.
    \item Gà: $200 : 500 = 40\% \neq 30\%$. Ý 3 Sai, Ý 4 Đúng.
\end{itemize}
\begin{itemize}
    \item a) \textbf{Đúng}.
    \item b) \textbf{Đúng}.
    \item c) \textbf{Sai}.
    \item d) \textbf{Đúng}.
\end{itemize}
}
\end{ex}

\begin{ex} %[5D12V70]
Hai ngăn sách có tất cả $180$ cuốn, ngăn thứ nhất bằng $\dfrac{4}{5}$ ngăn thứ hai. Xét tính đúng sai:
\choiceTF
{\True Tổng số phần bằng nhau là $9$ phần.}
{\True Ngăn thứ nhất có $80$ cuốn sách.}
{Ngăn thứ hai có $90$ cuốn sách.}
{\True Ngăn thứ hai nhiều hơn ngăn thứ nhất $20$ cuốn sách.}
\loigiai{
\begin{itemize}
    \item Tổng số phần: $4 + 5 = 9$. Đúng.
    \item Ngăn 1: $180 : 9 \times 4 = 80$ cuốn. Đúng.
    \item Ngăn 2: $180 - 80 = 100$ cuốn $\neq 90$ cuốn. Sai.
    \item Hiệu: $100 - 80 = 20$ cuốn. Đúng.
\end{itemize}
\begin{itemize}
    \item a) \textbf{Đúng}.
    \item b) \textbf{Đúng}.
    \item c) \textbf{Sai}.
    \item d) \textbf{Đúng}.
\end{itemize}
}
\end{ex}

\textbf{Phần 3. Trả lời ngắn}
\setcounterpageref{ex}{0}

\begin{ex} %[5D12H70]
Tìm tỉ số phần trăm của $18$ và $40$. (Chỉ ghi số).
\shortans{45}
\loigiai{
\textbf{Đáp số:} $45$\\
$18 : 40 = 45\%$.
}
\end{ex}

\begin{ex} %[5D12H70]
Tìm $35\%$ của $240$.
\shortans{84}
\loigiai{
\textbf{Đáp số:} $84$\\
$240 \times 35 : 100 = 84$.
}
\end{ex}

\begin{ex} %[5D12V70]
Tổng hai số là $144$, tỉ số là $\dfrac{3}{5}$. Tìm số lớn.
\shortans{90}
\loigiai{
\textbf{Đáp số:} $90$\\
$144 : (3 + 5) \times 5 = 90$.
}
\end{ex}

\begin{ex} %[5D12V70]
Hiệu hai số là $56$, tỉ số là $\dfrac{2}{9}$. Tìm số bé.
\shortans{16}
\loigiai{
\textbf{Đáp số:} $16$\\
$56 : (9 - 2) \times 2 = 16$.
}
\end{ex}

\begin{ex} %[5D12V70]
Trên bản đồ tỉ lệ $1 : 5,000$, một thửa ruộng hình chữ nhật có chiều dài $4\text{ cm}$. Tính chiều dài thật theo mét.
\shortans{200}
\loigiai{
\textbf{Đáp số:} $200$\\
$4 \times 5,000 = 20,000\text{ cm} = 200\text{ m}$.
}
\end{ex}

\begin{ex} %[5D12C70]
Một sản phẩm có giá gốc là $400\,000$ đồng. Người bán nâng giá lên $20\%$, sau đó lại giảm giá $10\%$ trên giá mới. Hỏi giá bán cuối cùng là bao nhiêu nghìn đồng?
\shortans{432}
\loigiai{
\textbf{Đáp số:} $432$\\
Giá sau khi nâng: $400 \times 1{,}2 = 480$ nghìn đồng.\\
Giá sau khi giảm 10\%: $480 \times 0{,}9 = 432$ nghìn đồng.
}
\end{ex}

\newpage

\TaoPhieuKiemTra{Chủ đề 12 -- Ôn tập cuối năm}{71}{Ôn tập hình học}{https://loptoanthaylinh.vn/thi-online/toan-5/bai-71}
\textbf{Phần 1. Trắc nghiệm lựa chọn}
\setcounterpageref{ex}{0}

\begin{ex} %[5H12N71]
Công thức tính diện tích hình tam giác có độ dài đáy $a$ và chiều cao $h$ (cùng đơn vị đo) là
\choice
{$S = a \times h$}
{$S = (a + h) \times 2$}
{$S = a + h : 2$}
{\True $S = \dfrac{a \times h}{2}$}
\loigiai{
$S = \dfrac{a \times h}{2}$.\\
\textbf{Đáp án của câu hỏi là D}
}
\end{ex}

\begin{ex} %[5H12N71]
Công thức tính diện tích hình thang có hai đáy $a, b$ và chiều cao $h$ là
\choice
{$S = (a + b) \times h$}
{$S = a \times b \times h$}
{\True $S = \dfrac{(a + b) \times h}{2}$}
{$S = (a + b) : 2 + h$}
\loigiai{
$S = \dfrac{(a + b) \times h}{2}$.\\
\textbf{Đáp án của câu hỏi là C}
}
\end{ex}

\begin{ex} %[5H12H71]
Một hình tròn có bán kính $5\text{ cm}$. Chu vi của hình tròn đó là
\choice
{\True $31{,}4\text{ cm}$}
{$15{,}7\text{ cm}$}
{$78{,}5\text{ cm}$}
{$62{,}8\text{ cm}$}
\loigiai{
$C = 2 \times 3{,}14 \times 5 = 31{,}4\text{ (cm)}$.\\
\textbf{Đáp án của câu hỏi là A}
}
\end{ex}

\begin{ex} %[5H12H71]
Một hình tròn có đường kính $6\text{ dm}$. Diện tích của hình tròn đó là
\choice
{$113{,}04\text{ dm}^2$}
{\True $28{,}26\text{ dm}^2$}
{$18{,}84\text{ dm}^2$}
{$56{,}52\text{ dm}^2$}
\loigiai{
Bán kính: $r = 3\text{ dm}$.\\
$S = 3{,}14 \times 3 \times 3 = 28{,}26\text{ (dm}^2\text{)}$.\\
\textbf{Đáp án của câu hỏi là B}
}
\end{ex}

\begin{ex} %[5H12H71]
Một hình thang có đáy lớn $15\text{ cm}$, đáy bé $9\text{ cm}$, chiều cao $8\text{ cm}$. Diện tích hình thang là
\choice
{\True $96\text{ cm}^2$}
{$192\text{ cm}^2$}
{$120\text{ cm}^2$}
{$84\text{ cm}^2$}
\loigiai{
$S = \dfrac{(15 + 9) \times 8}{2} = 96\text{ (cm}^2\text{)}$.\\
\textbf{Đáp án của câu hỏi là A}
}
\end{ex}

\begin{ex} %[5H12H71]
Một hình lập phương có cạnh $6\text{ cm}$. Diện tích toàn phần của nó là
\choice
{$144\text{ cm}^2$}
{\True $216\text{ cm}^2$}
{$216\text{ cm}^3$}
{$96\text{ cm}^2$}
\loigiai{
$S_{tp} = 6 \times 6 \times 6 = 216\text{ (cm}^2\text{)}$.\\
\textbf{Đáp án của câu hỏi là B}
}
\end{ex}

\begin{ex} %[5H12V71]
Một hình tam giác có diện tích $36\text{ cm}^2$, độ dài đáy là $12\text{ cm}$. Chiều cao tương ứng của tam giác là
\choice
{$3\text{ cm}$}
{$9\text{ cm}$}
{$12\text{ cm}$}
{\True $6\text{ cm}$}
\loigiai{
Chiều cao: $36 \times 2 : 12 = 6\text{ (cm)}$.\\
\textbf{Đáp án của câu hỏi là D}
}
\end{ex}

\begin{ex} %[5H12V71]
Một hình hộp chữ nhật có kích thước dài $8\text{ cm}$, rộng $6\text{ cm}$, cao $5\text{ cm}$. Thể tích của hình hộp chữ nhật đó là
\choice
{\True $240\text{ cm}^3$}
{$120\text{ cm}^3$}
{$140\text{ cm}^3$}
{$480\text{ cm}^3$}
\loigiai{
$V = 8 \times 6 \times 5 = 240\text{ (cm}^3\text{)}$.\\
\textbf{Đáp án của câu hỏi là A}
}
\end{ex}

\begin{ex} %[5H12V71]
Người ta quét sơn toàn bộ $6$ mặt ngoài của một khối gỗ hình lập phương có cạnh $1{,}2\text{ m}$. Diện tích quét sơn là
\choice
{$5{,}76\text{ m}^2$}
{$7{,}2\text{ m}^2$}
{\True $8{,}64\text{ m}^2$}
{$1{,}728\text{ m}^2$}
\loigiai{
$S_{tp} = 1{,}2 \times 1{,}2 \times 6 = 8{,}64\text{ (m}^2\text{)}$.\\
\textbf{Đáp án của câu hỏi là C}
}
\end{ex}

\begin{ex} %[5H12V71]
Một mảnh đất hình thang có diện tích $180\text{ m}^2$, chiều cao $12\text{ m}$. Trung bình cộng hai đáy của hình thang là
\choice
{$30\text{ m}$}
{\True $15\text{ m}$}
{$20\text{ m}$}
{$25\text{ m}$}
\loigiai{
Trung bình cộng 2 đáy: $180 : 12 = 15\text{ (m)}$.\\
\textbf{Đáp án của câu hỏi là B}
}
\end{ex}

\begin{ex} %[5H12V71]
Một bồn hoa gồm một hình vuông cạnh $4\text{ m}$ và $4$ nửa hình tròn có đường kính là cạnh của hình vuông gắn vào 4 cạnh. Diện tích toàn bộ bồn hoa là
\choice
{$28{,}56\text{ m}^2$}
{$32{,}56\text{ m}^2$}
{\True $41{,}12\text{ m}^2$}
{$50\text{ m}^2$}
\loigiai{
Diện tích hình vuông: $4 \times 4 = 16\text{ (m}^2\text{)}$.\\
4 nửa hình tròn = 2 hình tròn bán kính $2\text{ m}$.\\
Diện tích 2 hình tròn: $2 \times (3{,}14 \times 2 \times 2) = 25{,}12\text{ (m}^2\text{)}$.\\
Tổng diện tích: $16 + 25{,}12 = 41{,}12\text{ (m}^2\text{)}$.\\
\textbf{Đáp án của câu hỏi là C}
}
\end{ex}

\begin{ex} %[5H12C71]
Cho hình thang vuông $ABCD$ vuông tại $A$ và $D$, có đáy bé $AB = 15\text{ cm}$, đáy lớn $CD = 25\text{ cm}$, chiều cao $AD = 16\text{ cm}$. Điểm $M$ là trung điểm của cạnh $AD$. Diện tích tam giác $MBC$ là
\choice
{$140\text{ cm}^2$}
{$180\text{ cm}^2$}
{$200\text{ cm}^2$}
{\True $160\text{ cm}^2$}
\loigiai{
$S_{ABCD} = \dfrac{(15 + 25) \times 16}{2} = 320\text{ (cm}^2\text{)}$.\\
$M$ là trung điểm $AD \Rightarrow AM = MD = 8\text{ cm}$.\\
$S_{ABM} = \dfrac{15 \times 8}{2} = 60\text{ (cm}^2\text{)}$.\\
$S_{MCD} = \dfrac{25 \times 8}{2} = 100\text{ (cm}^2\text{)}$.\\
$S_{MBC} = 320 - (60 + 100) = 160\text{ (cm}^2\text{)}$.\\
\textbf{Đáp án của câu hỏi là D}
}
\end{ex}

\textbf{Phần 2. Trắc nghiệm đúng sai}
\setcounterpageref{ex}{0}

\begin{ex} %[5H12H71]
Cho tam giác $ABC$ có đáy $BC = 16\text{ cm}$, chiều cao $AH = 10\text{ cm}$. Xét tính đúng sai:
\choiceTF
{\True Diện tích tam giác $ABC$ là $80\text{ cm}^2$.}
{Nếu tăng đáy $BC$ thêm $4\text{ cm}$ thì diện tích tam giác tăng thêm $40\text{ cm}^2$.}
{\True Nếu tăng đáy $BC$ thêm $4\text{ cm}$ thì diện tích tam giác tăng thêm $20\text{ cm}^2$.}
{\True Khi chiều cao tăng gấp đôi thì diện tích tam giác tăng gấp đôi.}
\loigiai{
\begin{itemize}
    \item $S = 16 \times 10 : 2 = 80\text{ cm}^2$. Đúng.
    \item Diện tích tăng thêm: $4 \times 10 : 2 = 20\text{ cm}^2 \neq 40\text{ cm}^2$. Ý 2 Sai, Ý 3 Đúng.
    \item Diện tích tỉ lệ thuận với chiều cao nên tăng gấp đôi. Đúng.
\end{itemize}
\begin{itemize}
    \item a) \textbf{Đúng}.
    \item b) \textbf{Sai}.
    \item c) \textbf{Đúng}.
    \item d) \textbf{Đúng}.
\end{itemize}
}
\end{ex}

\begin{ex} %[5H12V71]
Một bể nước hình hộp chữ nhật có kích thước lòng trong: dài $2\text{ m}$, rộng $1{,}5\text{ m}$, cao $1{,}2\text{ m}$. Xét tính đúng sai:
\choiceTF
{\True Diện tích đáy bể là $3\text{ m}^2$.}
{\True Diện tích xung quanh bên trong bể là $8{,}4\text{ m}^2$.}
{Dung tích tối đa của bể là $3\,000\text{ l}$.}
{\True Dung tích tối đa của bể là $3\,600\text{ l}$.}
\loigiai{
\begin{itemize}
    \item $S_{đáy} = 2 \times 1{,}5 = 3\text{ m}^2$. Đúng.
    \item $S_{xq} = (2 + 1{,}5) \times 2 \times 1{,}2 = 8{,}4\text{ m}^2$. Đúng.
    \item Dung tích: $3 \times 1{,}2 = 3{,}6\text{ m}^3 = 3600\text{ l} \neq 3000\text{ l}$. Ý 3 Sai, Ý 4 Đúng.
\end{itemize}
\begin{itemize}
    \item a) \textbf{Đúng}.
    \item b) \textbf{Đúng}.
    \item c) \textbf{Sai}.
    \item d) \textbf{Đúng}.
\end{itemize}
}
\end{ex}

\textbf{Phần 3. Trả lời ngắn}
\setcounterpageref{ex}{0}

\begin{ex} %[5H12H71]
Một hình tam giác có đáy $14\text{ cm}$, chiều cao $8\text{ cm}$. Tính diện tích theo xăng-ti-mét vuông.
\shortans{56}
\loigiai{
\textbf{Đáp số:} $56$\\
$14 \times 8 : 2 = 56\text{ cm}^2$.
}
\end{ex}

\begin{ex} %[5H12H71]
Một hình tròn có bán kính $4\text{ cm}$. Tính diện tích theo xăng-ti-mét vuông.
\shortans{50,24}
\loigiai{
\textbf{Đáp số:} $50{,}24$\\
$3{,}14 \times 4 \times 4 = 50{,}24\text{ cm}^2$.
}
\end{ex}

\begin{ex} %[5H12V71]
Một hình thang có đáy lớn $20\text{ m}$, đáy bé $14\text{ m}$, chiều cao $10\text{ m}$. Tính diện tích hình thang theo mét vuông.
\shortans{170}
\loigiai{
\textbf{Đáp số:} $170$\\
$(20 + 14) \times 10 : 2 = 170\text{ m}^2$.
}
\end{ex}

\begin{ex} %[5H12V71]
Một hình lập phương có diện tích toàn phần là $96\text{ cm}^2$. Tính thể tích hình lập phương theo xăng-ti-mét khối.
\shortans{64}
\loigiai{
\textbf{Đáp số:} $64$\\
Diện tích 1 mặt: $96 : 6 = 16\text{ cm}^2 \Rightarrow$ Cạnh $4\text{ cm}$.\\
$V = 4 \times 4 \times 4 = 64\text{ cm}^3$.
}
\end{ex}

\begin{ex} %[5H12V71]
Một bể cá hình hộp chữ nhật dài $80\text{ cm}$, rộng $50\text{ cm}$, cao $40\text{ cm}$. Tính dung tích tối đa của bể theo lít.
\shortans{160}
\loigiai{
\textbf{Đáp số:} $160$\\
$8 \times 5 \times 4 = 160\text{ dm}^3 = 160\text{ l}$.
}
\end{ex}

\begin{ex} %[5H12C71]
Một cái sân hình chữ nhật dài $18\text{ m}$, rộng $12\text{ m}$. Ở giữa sân người ta xây một bồn hoa hình tròn bán kính $2\text{ m}$. Tính diện tích phần sân còn lại theo mét vuông.
\shortans{203,44}
\loigiai{
\textbf{Đáp số:} $203{,}44$\\
Diện tích sân: $18 \times 12 = 216\text{ m}^2$.\\
Diện tích bồn hoa: $3{,}14 \times 2 \times 2 = 12{,}56\text{ m}^2$.\\
Diện tích còn lại: $216 - 12{,}56 = 203{,}44\text{ m}^2$.
}
\end{ex}

\newpage

\TaoPhieuKiemTra{Chủ đề 12 -- Ôn tập cuối năm}{72}{Ôn tập đo lường}{https://loptoanthaylinh.vn/thi-online/toan-5/bai-72}
\textbf{Phần 1. Trắc nghiệm lựa chọn}
\setcounterpageref{ex}{0}

\begin{ex} %[5D12N72]
$1\text{ m}^3$ bằng bao nhiêu đề-xi-mét khối?
\choice
{$10\text{ dm}^3$}
{$100\text{ dm}^3$}
{\True $1,000\text{ dm}^3$}
{$10,000\text{ dm}^3$}
\loigiai{
$1\text{ m}^3 = 1,000\text{ dm}^3$.\\
\textbf{Đáp án của câu hỏi là C}
}
\end{ex}

\begin{ex} %[5D12N72]
$3\text{ tấn } 50\text{ kg}$ viết dưới dạng số thập phân có đơn vị là tấn là
\choice
{$3{,}5\text{ tấn}$}
{$3{,}005\text{ tấn}$}
{$30{,}5\text{ tấn}$}
{\True $3{,}05\text{ tấn}$}
\loigiai{
$3\text{ tấn } 50\text{ kg} = 3 + \dfrac{50}{1000} = 3{,}05\text{ tấn}$.\\
\textbf{Đáp án của câu hỏi là D}
}
\end{ex}

\begin{ex} %[5D12H72]
$2{,}4\text{ ha}$ bằng bao nhiêu mét vuông?
\choice
{$240\text{ m}^2$}
{\True $24,000\text{ m}^2$}
{$2,400\text{ m}^2$}
{$240,000\text{ m}^2$}
\loigiai{
$1\text{ ha} = 10,000\text{ m}^2 \Rightarrow 2{,}4\text{ ha} = 24,000\text{ m}^2$.\\
\textbf{Đáp án của câu hỏi là B}
}
\end{ex}

\begin{ex} %[5D12H72]
$1$ giờ $45$ phút bằng bao nhiêu giờ?
\choice
{\True $1{,}75\text{ giờ}$}
{$1{,}45\text{ giờ}$}
{$1{,}25\text{ giờ}$}
{$1{,}5\text{ giờ}$}
\loigiai{
$1 + 45 : 60 = 1{,}75\text{ giờ}$.\\
\textbf{Đáp án của câu hỏi là A}
}
\end{ex}

\begin{ex} %[5D12H72]
$450\text{ dm}^3$ bằng bao nhiêu mét khối?
\choice
{$4{,}5\text{ m}^3$}
{$0{,}045\text{ m}^3$}
{\True $0{,}45\text{ m}^3$}
{$45\text{ m}^3$}
\loigiai{
$450 : 1,000 = 0{,}45\text{ m}^3$.\\
\textbf{Đáp án của câu hỏi là C}
}
\end{ex}

\begin{ex} %[5D12H72]
Điền số thích hợp vào chỗ chấm: $5\text{ m } 8\text{ cm} = \dots\text{ m}$.
\choice
{\True $5{,}08$}
{$5{,}8$}
{$5{,}008$}
{$58$}
\loigiai{
$5 + 8 : 100 = 5{,}08\text{ m}$.\\
\textbf{Đáp án của câu hỏi là A}
}
\end{ex}

\begin{ex} %[5D12V72]
Một thửa ruộng thu hoạch được $2{,}4\text{ tấn}$ thóc. Người ta đóng đều vào $48$ bao. Mỗi bao nặng bao nhiêu ki-lô-gam thóc?
\choice
{$40\text{ kg}$}
{$60\text{ kg}$}
{$55\text{ kg}$}
{\True $50\text{ kg}$}
\loigiai{
$2{,}4\text{ tấn} = 2,400\text{ kg}$.\\
Mỗi bao: $2,400 : 48 = 50\text{ (kg)}$.\\
\textbf{Đáp án của câu hỏi là D}
}
\end{ex}

\begin{ex} %[5D12V72]
Một mảnh đất có diện tích $0{,}15\text{ ha}$. Người ta dành $40\%$ diện tích để xây nhà. Diện tích đất xây nhà là
\choice
{$500\text{ m}^2$}
{\True $600\text{ m}^2$}
{$750\text{ m}^2$}
{$60\text{ m}^2$}
\loigiai{
$0{,}15\text{ ha} = 1,500\text{ m}^2$.\\
Diện tích xây nhà: $1,500 \times 40 : 100 = 600\text{ (m}^2\text{)}$.\\
\textbf{Đáp án của câu hỏi là B}
}
\end{ex}

\begin{ex} %[5D12V72]
Một bể nước chứa được $1{,}8\text{ m}^3$ nước. Hiện tại trong bể đã có $600\text{ l}$ nước. Cần đổ thêm bao nhiêu lít nước nữa để đầy bể?
\choice
{\True $1,200\text{ l}$}
{$1,000\text{ l}$}
{$1,400\text{ l}$}
{$800\text{ l}$}
\loigiai{
$1{,}8\text{ m}^3 = 1,800\text{ l}$.\\
Cần đổ thêm: $1,800 - 600 = 1,200\text{ (l)}$.\\
\textbf{Đáp án của câu hỏi là A}
}
\end{ex}

\begin{ex} %[5D12V72]
Sắp xếp các số đo sau theo thứ tự từ bé đến lớn: $3\text{ m}^2 5\text{ dm}^2; 3{,}5\text{ m}^2; 30,050\text{ cm}^2; 3{,}04\text{ m}^2$.
\choice
{$30,050\text{ cm}^2 < 3{,}04\text{ m}^2 < 3\text{ m}^2 5\text{ dm}^2 < 3{,}5\text{ m}^2$}
{$3{,}04\text{ m}^2 < 30,050\text{ cm}^2 < 3\text{ m}^2 5\text{ dm}^2 < 3{,}5\text{ m}^2$}
{\True $3{,}005\text{ m}^2 (30,050\text{ cm}^2) < 3{,}04\text{ m}^2 < 3{,}05\text{ m}^2 (3\text{ m}^2 5\text{ dm}^2) < 3{,}5\text{ m}^2$}
{$3{,}5\text{ m}^2 < 3\text{ m}^2 5\text{ dm}^2 < 3{,}04\text{ m}^2 < 30,050\text{ cm}^2$}
\loigiai{
Đổi về $\text{m}^2$: $30,050\text{ cm}^2 = 3{,}005\text{ m}^2 < 3{,}04\text{ m}^2 < 3{,}05\text{ m}^2 < 3{,}5\text{ m}^2$.\\
\textbf{Đáp án của câu hỏi là C}
}
\end{ex}

\begin{ex} %[5D12V72]
Một cuộn dây thép dài $45\text{ m } 6\text{ dm}$, mỗi mét dây nặng $250\text{ g}$. Cả cuộn dây đó nặng bao nhiêu ki-lô-gam?
\choice
{$11\text{ kg}$}
{\True $11{,}4\text{ kg}$}
{$114\text{ kg}$}
{$1{,}14\text{ kg}$}
\loigiai{
$45\text{ m } 6\text{ dm} = 45{,}6\text{ m}$.\\
Khối lượng: $45{,}6 \times 250 = 11,400\text{ g} = 11{,}4\text{ kg}$.\\
\textbf{Đáp án của câu hỏi là B}
}
\end{ex}

\begin{ex} %[5D12C72]
Một vòi nước chảy vào bể không có nước, mỗi phút chảy được $25\text{ l}$ nước. Sau $1$ giờ $12$ phút thì lượng nước trong bể chiếm $60\%$ dung tích của bể. Dung tích toàn bộ của bể nước đó là
\choice
{$2,500\text{ l}$}
{$3,600\text{ l}$}
{$4,000\text{ l}$}
{\True $3,000\text{ l}$}
\loigiai{
$1$ giờ $12$ phút $= 72$ phút.\\
Lượng nước đã chảy: $72 \times 25 = 1,800\text{ (l)}$.\\
Dung tích bể: $1,800 : 60 \times 100 = 3,000\text{ (l)}$.\\
\textbf{Đáp án của câu hỏi là D}
}
\end{ex}

\textbf{Phần 2. Trắc nghiệm đúng sai}
\setcounterpageref{ex}{0}

\begin{ex} %[5D12H72]
Xét tính đúng sai của các biểu thức đổi đơn vị đo:
\choiceTF
{\True $4\text{ tấn } 8\text{ kg} = 4{,}008\text{ tấn}$.}
{$5\text{ m}^2 6\text{ dm}^2 = 5{,}6\text{ m}^2$.}
{\True $5\text{ m}^2 6\text{ dm}^2 = 5{,}06\text{ m}^2$.}
{\True $2\text{ dm}^3 45\text{ cm}^3 = 2{,}045\text{ dm}^3$.}
\loigiai{
\begin{itemize}
    \item $4\text{ tấn } 8\text{ kg} = 4 + 8 : 1000 = 4{,}008\text{ tấn}$. Đúng.
    \item $5\text{ m}^2 6\text{ dm}^2 = 5 + 6 : 100 = 5{,}06\text{ m}^2 \neq 5{,}6\text{ m}^2$. Ý 2 Sai, Ý 3 Đúng.
    \item $2\text{ dm}^3 45\text{ cm}^3 = 2 + 45 : 1000 = 2{,}045\text{ dm}^3$. Đúng.
\end{itemize}
\begin{itemize}
    \item a) \textbf{Đúng}.
    \item b) \textbf{Sai}.
    \item c) \textbf{Đúng}.
    \item d) \textbf{Đúng}.
\end{itemize}
}
\end{ex}

\begin{ex} %[5D12V72]
Một xe tải chở $3$ chuyến hàng: Chuyến 1 chở $2{,}5\text{ tấn}$, chuyến 2 chở $28\text{ tạ}$, chuyến 3 chở $2,700\text{ kg}$. Xét tính đúng sai:
\choiceTF
{\True Chuyến thứ hai chở nhiều hàng nhất ($2{,}8\text{ tấn}$).}
{\True Chuyến thứ nhất chở ít hàng nhất ($2{,}5\text{ tấn}$).}
{Tổng khối lượng hàng chở trong cả 3 chuyến là $7\text{ tấn}$.}
{\True Tổng khối lượng hàng chở trong cả 3 chuyến là $8\text{ tấn}$.}
\loigiai{
\begin{itemize}
    \item Đổi ra tấn: $2{,}5\text{ tấn}$; $28\text{ tạ} = 2{,}8\text{ tấn}$; $2,700\text{ kg} = 2{,}7\text{ tấn}$.
    \item Chuyến 2 nhiều nhất ($2{,}8\text{ tấn}$). Đúng.
    \item Chuyến 1 ít nhất ($2{,}5\text{ tấn}$). Đúng.
    \item Tổng: $2{,}5 + 2{,}8 + 2{,}7 = 8{,}0\text{ tấn} \neq 7\text{ tấn}$. Ý 3 Sai, Ý 4 Đúng.
\end{itemize}
\begin{itemize}
    \item a) \textbf{Đúng}.
    \item b) \textbf{Đúng}.
    \item c) \textbf{Sai}.
    \item d) \textbf{Đúng}.
\end{itemize}
}
\end{ex}

\textbf{Phần 3. Trả lời ngắn}
\setcounterpageref{ex}{0}

\begin{ex} %[5D12H72]
Đổi $3{,}75\text{ m}^3$ ra đề-xi-mét khối.
\shortans{3750}
\loigiai{
\textbf{Đáp số:} $3750$\\
$3{,}75 \times 1,000 = 3,750\text{ dm}^3$.
}
\end{ex}

\begin{ex} %[5D12H72]
Đổi $5\text{ ha } 400\text{ m}^2$ ra số thập phân có đơn vị là héc-ta.
\shortans{5,04}
\loigiai{
\textbf{Đáp số:} $5{,}04$\\
$5 + 400 : 10,000 = 5{,}04\text{ ha}$.
}
\end{ex}

\begin{ex} %[5D12V72]
Đổi $2$ giờ $36$ phút ra giờ (dùng dấu phẩy nếu có).
\shortans{2,6}
\loigiai{
\textbf{Đáp số:} $2{,}6$\\
$2 + 36 : 60 = 2{,}6\text{ giờ}$.
}
\end{ex}

\begin{ex} %[5D12V72]
Một thửa ruộng hình chữ nhật có diện tích $3,200\text{ m}^2$. Hỏi thửa ruộng đó rộng bao nhiêu héc-ta? (Dùng dấu phẩy nếu có).
\shortans{0,32}
\loigiai{
\textbf{Đáp số:} $0{,}32$\\
$3,200 : 10,000 = 0{,}32\text{ ha}$.
}
\end{ex}

\begin{ex} %[5D12V72]
Một bình nước chứa $2{,}5\text{ l}$ nước. Người ta rót đều vào $10$ chiếc cốc nhỏ. Hỏi mỗi cốc chứa bao nhiêu mi-li-lít nước?
\shortans{250}
\loigiai{
\textbf{Đáp số:} $250$\\
$2,500 : 10 = 250\text{ ml}$.
}
\end{ex}

\begin{ex} %[5D12C72]
Người ta lát một căn phòng hình chữ nhật dài $8\text{ m}$, rộng $6\text{ m}$ bằng những viên gạch hình vuông cạnh $40\text{ cm}$. Giá tiền mỗi viên gạch là $25,000$ đồng. Tính tổng số tiền mua gạch theo nghìn đồng.
\shortans{7500}
\loigiai{
\textbf{Đáp số:} $7500$\\
Diện tích phòng: $8 \times 6 = 48\text{ m}^2 = 480,000\text{ cm}^2$.\\
Diện tích 1 viên gạch: $40 \times 40 = 1,600\text{ cm}^2$.\\
Số viên gạch: $480,000 : 1,600 = 300$ viên.\\
Tổng tiền: $300 \times 25 = 7,500$ nghìn đồng.
}
\end{ex}

\newpage

\TaoPhieuKiemTra{Chủ đề 12 -- Ôn tập cuối năm}{73}{Ôn tập toán chuyển động đều}{https://loptoanthaylinh.vn/thi-online/toan-5/bai-73}
\textbf{Phần 1. Trắc nghiệm lựa chọn}
\setcounterpageref{ex}{0}

\begin{ex} %[5D12N73]
Công thức tính thời gian $t$ trong chuyển động đều là
\choice
{$t = s \times v$}
{$t = v : s$}
{$t = s - v$}
{\True $t = s : v$}
\loigiai{
$t = s : v$.\\
\textbf{Đáp án của câu hỏi là D}
}
\end{ex}

\begin{ex} %[5D12N73]
Một người đi xe máy với vận tốc $40\text{ km/h}$ trong $2{,}5\text{ giờ}$. Quãng đường người đó đi được là
\choice
{$80\text{ km}$}
{$120\text{ km}$}
{\True $100\text{ km}$}
{$90\text{ km}$}
\loigiai{
$s = 40 \times 2{,}5 = 100\text{ (km)}$.\\
\textbf{Đáp án của câu hỏi là C}
}
\end{ex}

\begin{ex} %[5D12H73]
Một ô tô đi quãng đường $126\text{ km}$ hết $2$ giờ $48$ phút. Vận tốc của ô tô là
\choice
{\True $45\text{ km/h}$}
{$42\text{ km/h}$}
{$48\text{ km/h}$}
{$50\text{ km/h}$}
\loigiai{
$2$ giờ $48$ phút $= 2{,}8$ giờ.\\
$v = 126 : 2{,}8 = 45\text{ (km/h)}$.\\
\textbf{Đáp án của câu hỏi là A}
}
\end{ex}

\begin{ex} %[5D12H73]
Một người đi xe đạp với vận tốc $12\text{ km/h}$ đi quãng đường $18\text{ km}$. Thời gian người đó đi là
\choice
{$1$ giờ $20$ phút}
{$1$ giờ $45$ phút}
{\True $1$ giờ $30$ phút}
{$2$ giờ}
\loigiai{
$t = 18 : 12 = 1{,}5\text{ giờ} = 1$ giờ $30$ phút.\\
\textbf{Đáp án của câu hỏi là C}
}
\end{ex}

\begin{ex} %[5D12H73]
Hai địa điểm A và B cách nhau $150\text{ km}$. Hai ô tô cùng khởi hành từ A và B đi ngược chiều nhau với vận tốc lần lượt là $42\text{ km/h}$ và $58\text{ km/h}$. Sau bao lâu hai ô tô gặp nhau?
\choice
{\True $1{,}5\text{ giờ}$}
{$1\text{ giờ}$}
{$2\text{ giờ}$}
{$2{,}5\text{ giờ}$}
\loigiai{
Tổng vận tốc: $42 + 58 = 100\text{ (km/h)}$.\\
Thời gian gặp nhau: $150 : 100 = 1{,}5\text{ (giờ)}$.\\
\textbf{Đáp án của câu hỏi là A}
}
\end{ex}

\begin{ex} %[5D12H73]
Một người đi bộ với vận tốc $4{,}5\text{ km/h}$ trong $40$ phút. Quãng đường đi được là
\choice
{$2\text{ km}$}
{\True $3\text{ km}$}
{$4\text{ km}$}
{$3{,}5\text{ km}$}
\loigiai{
$40$ phút $= \dfrac{2}{3}$ giờ.\\
$s = 4{,}5 \times \dfrac{2}{3} = 3\text{ (km)}$.\\
\textbf{Đáp án của câu hỏi là B}
}
\end{ex}

\begin{ex} %[5D12V73]
Một xe máy khởi hành từ A lúc $7$ giờ $15$ phút với vận tốc $32\text{ km/h}$. Đến $8$ giờ, một ô tô xuất phát từ A đuổi theo với vận tốc $48\text{ km/h}$. Ô tô đuổi kịp xe máy lúc mấy giờ?
\choice
{$9$ giờ $15$ phút}
{$9$ giờ $45$ phút}
{$10$ giờ}
{\True $9$ giờ $30$ phút}
\loigiai{
Từ $7$h $15$p đến $8$h là $45$ phút $= 0{,}75$ giờ.\\
Quãng đường xe máy đã đi: $32 \times 0{,}75 = 24\text{ (km)}$.\\
Thời gian đuổi kịp: $24 : (48 - 32) = 1{,}5\text{ giờ} = 1$ giờ $30$ phút.\\
Thời điểm gặp: $8$ giờ $+ 1$ giờ $30$ phút $= 9$ giờ $30$ phút.\\
\textbf{Đáp án của câu hỏi là D}
}
\end{ex}

\begin{ex} %[5D12V73]
Một thuyền máy đi xuôi dòng một khúc sông từ bến A đến bến B hết $1$ giờ $45$ phút. Vận tốc thực của thuyền là $22\text{ km/h}$, vận tốc dòng nước là $2\text{ km/h}$. Độ dài khúc sông AB là
\choice
{$38\text{ km}$}
{\True $42\text{ km}$}
{$44\text{ km}$}
{$40\text{ km}$}
\loigiai{
Vận tốc xuôi dòng: $22 + 2 = 24\text{ (km/h)}$.\\
$1$ giờ $45$ phút $= 1{,}75$ giờ.\\
$s = 24 \times 1{,}75 = 42\text{ (km)}$.\\
\textbf{Đáp án của câu hỏi là B}
}
\end{ex}

\begin{ex} %[5D12V73]
Một người đi xe máy từ A đến B với vận tốc $36\text{ km/h}$, rồi từ B quay về A với vận tốc $30\text{ km/h}$. Thời gian lúc về nhiều hơn lúc đi là $20$ phút. Quãng đường AB dài là
\choice
{$50\text{ km}$}
{$70\text{ km}$}
{$80\text{ km}$}
{\True $60\text{ km}$}
\loigiai{
Tỉ số vận tốc đi và về: $\dfrac{36}{30} = \dfrac{6}{5} \Rightarrow$ Tỉ số thời gian là $\dfrac{5}{6}$.\\
Hiệu thời gian: $20$ phút $= \dfrac{1}{3}$ giờ.\\
Thời gian đi: $\dfrac{1}{3} : (6 - 5) \times 5 = \dfrac{5}{3}$ giờ.\\
Quãng đường AB: $36 \times \dfrac{5}{3} = 60\text{ (km)}$.\\
\textbf{Đáp án của câu hỏi là D}
}
\end{ex}

\begin{ex} %[5D12V73]
Một đoàn tàu hỏa dài $200\text{ m}$ chạy qua một đường hầm dài $800\text{ m}$ với vận tốc $72\text{ km/h}$. Thời gian từ lúc đầu tàu vào hầm đến lúc đuôi tàu ra khỏi hầm là
\choice
{$40\text{ giây}$}
{$60\text{ giây}$}
{\True $50\text{ giây}$}
{$75\text{ giây}$}
\loigiai{
Tổng quãng đường tàu đi: $800 + 200 = 1,000\text{ (m)}$.\\
Vận tốc: $72\text{ km/h} = 20\text{ m/s}$.\\
Thời gian: $1,000 : 20 = 50\text{ (giây)}$.\\
\textbf{Đáp án của câu hỏi là C}
}
\end{ex}

\begin{ex} %[5D12V73]
Bác Bình đi bộ thể dục $3$ vòng quanh một công viên hình chữ nhật có chiều dài $250\text{ m}$, chiều rộng $150\text{ m}$ với vận tốc $4\text{ km/h}$. Thời gian bác Bình đi hết 3 vòng là
\choice
{$30$ phút}
{\True $36$ phút}
{$40$ phút}
{$45$ phút}
\loigiai{
Chu vi công viên: $(250 + 150) \times 2 = 800\text{ (m)}$.\\
Quãng đường 3 vòng: $800 \times 3 = 2,400\text{ (m)} = 2{,}4\text{ km}$.\\
Thời gian: $2{,}4 : 4 = 0{,}6\text{ giờ} = 36$ phút.\\
\textbf{Đáp án của câu hỏi là B}
}
\end{ex}

\begin{ex} %[5D12C73]
Hai thành phố A và B cách nhau $180\text{ km}$. Cùng lúc $7$ giờ, một ô tô đi từ A đến B và một xe máy đi từ B về A. Hai xe gặp nhau lúc $9$ giờ $15$ phút. Biết vận tốc ô tô gấp đôi vận tốc xe máy. Vận tốc của ô tô là
\choice
{\True $53{,}33\text{ km/h}$ (hay $\dfrac{160}{3}\text{ km/h}$)}
{$50\text{ km/h}$}
{$60\text{ km/h}$}
{$70\text{ km/h}$}
\loigiai{
Thời gian gặp nhau: $9$ giờ $15$ phút $- 7$ giờ $= 2$ giờ $15$ phút $= 2{,}25$ giờ.\\
Tổng vận tốc: $180 : 2{,}25 = 80\text{ (km/h)}$.\\
Vận tốc ô tô: $80 : (2 + 1) \times 2 = \dfrac{160}{3} \approx 53{,}33\text{ (km/h)}$.\\
\textbf{Đáp án của câu hỏi là A}
}
\end{ex}

\textbf{Phần 2. Trắc nghiệm đúng sai}
\setcounterpageref{ex}{0}

\begin{ex} %[5D12H73]
Hai địa điểm A và B cách nhau $135\text{ km}$. Xe máy đi từ A đến B với vận tốc $36\text{ km/h}$, ô tô đi từ B về A với vận tốc $54\text{ km/h}$. Hai xe cùng khởi hành lúc $8$ giờ. Xét tính đúng sai:
\choiceTF
{\True Tổng vận tốc của hai xe là $90\text{ km/h}$.}
{\True Hai xe gặp nhau sau $1$ giờ $30$ phút.}
{Hai xe gặp nhau lúc $9$ giờ $15$ phút.}
{\True Nơi gặp nhau cách A quãng đường $54\text{ km}$.}
\loigiai{
\begin{itemize}
    \item Tổng vận tốc: $36 + 54 = 90\text{ km/h}$. Đúng.
    \item Thời gian gặp: $135 : 90 = 1{,}5\text{ giờ} = 1$ giờ $30$ phút. Đúng.
    \item Gặp nhau lúc: $8$h $+ 1$h $30$p $= 9$h $30$p $\neq 9$h $15$p. Sai.
    \item Cách A: $36 \times 1{,}5 = 54\text{ km}$. Đúng.
\end{itemize}
\begin{itemize}
    \item a) \textbf{Đúng}.
    \item b) \textbf{Đúng}.
    \item c) \textbf{Sai}.
    \item d) \textbf{Đúng}.
\end{itemize}
}
\end{ex}

\begin{ex} %[5D12V73]
Một ca nô đi xuôi dòng khúc sông từ A đến B hết $2$ giờ, đi ngược dòng từ B về A hết $3$ giờ. Vận tốc dòng nước là $3\text{ km/h}$. Xét tính đúng sai:
\choiceTF
{\True Tỉ số vận tốc xuôi dòng và ngược dòng là $\dfrac{3}{2}$.}
{\True Hiệu vận tốc xuôi dòng và ngược dòng là $6\text{ km/h}$.}
{Vận tốc ca nô xuôi dòng là $15\text{ km/h}$.}
{\True Độ dài khúc sông AB là $36\text{ km}$.}
\loigiai{
\begin{itemize}
    \item Tỉ số vận tốc: $3 : 2 = \dfrac{3}{2}$. Đúng.
    \item Hiệu vận tốc: $2 \times 3 = 6\text{ km/h}$. Đúng.
    \item Vận tốc xuôi dòng: $6 : (3 - 2) \times 3 = 18\text{ km/h} \neq 15\text{ km/h}$. Sai.
    \item Khúc sông AB: $18 \times 2 = 36\text{ km}$. Đúng.
\end{itemize}
\begin{itemize}
    \item a) \textbf{Đúng}.
    \item b) \textbf{Đúng}.
    \item c) \textbf{Sai}.
    \item d) \textbf{Đúng}.
\end{itemize}
}
\end{ex}

\textbf{Phần 3. Trả lời ngắn}
\setcounterpageref{ex}{0}

\begin{ex} %[5D12H73]
Một người đi xe máy với vận tốc $38\text{ km/h}$ trong $2{,}5$ giờ. Tính quãng đường theo ki-lô-mét.
\shortans{95}
\loigiai{
\textbf{Đáp số:} $95$\\
$38 \times 2{,}5 = 95\text{ km}$.
}
\end{ex}

\begin{ex} %[5D12H73]
Một ô tô chạy quãng đường $160\text{ km}$ với vận tốc $50\text{ km/h}$. Tính thời gian ô tô đi theo giờ (dùng dấu phẩy nếu có).
\shortans{3,2}
\loigiai{
\textbf{Đáp số:} $3{,}2$\\
$160 : 50 = 3{,}2\text{ giờ}$.
}
\end{ex}

\begin{ex} %[5D12V73]
Hai người đi bộ khởi hành cùng lúc đi ngược chiều nhau trên con đường dài $9\text{ km}$ với vận tốc $4\text{ km/h}$ và $5\text{ km/h}$. Hỏi sau bao nhiêu phút hai người gặp nhau?
\shortans{60}
\loigiai{
\textbf{Đáp số:} $60$\\
$9 : (4 + 5) = 1\text{ giờ} = 60\text{ phút}$.
}
\end{ex}

\begin{ex} %[5D12V73]
Một người đi xe đạp từ nhà đến cơ quan dài $6\text{ km}$ hết $24$ phút. Tính vận tốc của người đi xe đạp theo $\text{km/h}$.
\shortans{15}
\loigiai{
\textbf{Đáp số:} $15$\\
$24$ phút $= 0{,}4$ giờ $\Rightarrow 6 : 0{,}4 = 15\text{ km/h}$.
}
\end{ex}

\begin{ex} %[5D12V73]
Một máy bay cất cánh lúc $8$ giờ $30$ phút và bay quãng đường $1,400\text{ km}$ với vận tốc $700\text{ km/h}$. Hỏi máy bay hạ cánh lúc mấy giờ mấy phút? Tính số phút từ $0$ giờ đến lúc hạ cánh.
\shortans{630}
\loigiai{
\textbf{Đáp số:} $630$\\
Thời gian bay: $1,400 : 700 = 2$ giờ.\\
Hạ cánh: $10$ giờ $30$ phút $= 10 \times 60 + 30 = 630$ phút.
}
\end{ex}

\begin{ex} %[5D12C73]
Một người đi xe máy từ A đến B với vận tốc $40\text{ km/h}$. Sau $45$ phút, một ô tô xuất phát từ A đuổi theo với vận tốc $55\text{ km/h}$. Hỏi sau bao nhiêu giờ ô tô đuổi kịp xe máy?
\shortans{2}
\loigiai{
\textbf{Đáp số:} $2$\\
Khoảng cách ban đầu: $40 \times 0{,}75 = 30\text{ km}$.\\
Thời gian đuổi kịp: $30 : (55 - 40) = 30 : 15 = 2\text{ giờ}$.
}
\end{ex}

\newpage

\TaoPhieuKiemTra{Chủ đề 12 -- Ôn tập cuối năm}{74}{Ôn tập thống kê và xác suất}{https://loptoanthaylinh.vn/thi-online/toan-5/bai-74}
\textbf{Phần 1. Trắc nghiệm lựa chọn}
\setcounterpageref{ex}{0}

\begin{ex} %[5X12N74]
Dữ liệu nào sau đây là dữ liệu định lượng (bằng số)?
\choice
{Xếp loại rèn luyện: Tốt, Khá, Đạt}
{\True Chiều cao các học sinh: $135\text{ cm}, 140\text{ cm}, 138\text{ cm}$}
{Môn học yêu thích: Toán, Tiếng Việt, Tiếng Anh}
{Màu sắc bóng: Xanh, Đỏ, Vàng}
\loigiai{
Chiều cao là dữ liệu định lượng bằng số.\\
\textbf{Đáp án của câu hỏi là B}
}
\end{ex}

\begin{ex} %[5X12N74]
Biểu đồ quạt tròn biểu diễn một hình quạt chiếm một góc vuông ($90^\circ$) tương ứng với bao nhiêu phần trăm?
\choice
{$20\%$}
{$30\%$}
{$50\%$}
{\True $25\%$}
\loigiai{
$\dfrac{90}{360} = \dfrac{1}{4} = 25\%$.\\
\textbf{Đáp án của câu hỏi là D}
}
\end{ex}

\begin{ex} %[5X12H74]
Gieo một con xúc xắc $60$ lần, mặt $3$ chấm xuất hiện $12$ lần. Tỉ số phần trăm số lần xuất hiện mặt $3$ chấm là
\choice
{\True $20\%$}
{$15\%$}
{$25\%$}
{$18\%$}
\loigiai{
$12 : 60 = 0{,}2 = 20\%$.\\
\textbf{Đáp án của câu hỏi là A}
}
\end{ex}

\begin{ex} %[5X12H74]
Biểu đồ quạt tròn biểu diễn thể loại sách mượn tại thư viện gồm: Truyện tranh $45\%$, Khoa học $30\%$, Lịch sử $25\%$. Biết có $200$ cuốn sách được mượn. Số cuốn sách khoa học được mượn là
\choice
{$50$ cuốn}
{$70$ cuốn}
{\True $60$ cuốn}
{$90$ cuốn}
\loigiai{
$200 \times 30 : 100 = 60$ (cuốn).\\
\textbf{Đáp án của câu hỏi là C}
}
\end{ex}

\begin{ex} %[5X12H74]
Cho bảng số liệu nhiệt độ $5$ ngày: $26^\circ\text{C}, 28^\circ\text{C}, 30^\circ\text{C}, 31^\circ\text{C}, 25^\circ\text{C}$. Nhiệt độ trung bình của $5$ ngày là
\choice
{\True $28^\circ\text{C}$}
{$27^\circ\text{C}$}
{$29^\circ\text{C}$}
{$30^\circ\text{C}$}
\loigiai{
Tổng nhiệt độ: $26 + 28 + 30 + 31 + 25 = 140^\circ\text{C}$.\\
Trung bình: $140 : 5 = 28^\circ\text{C}$.\\
\textbf{Đáp án của câu hỏi là A}
}
\end{ex}

\begin{ex} %[5X12H74]
Tung một đồng xu $80$ lần có $44$ lần xuất hiện mặt Ngửa. Tỉ số của số lần xuất hiện mặt Sấp so với tổng số lần tung là
\choice
{$\dfrac{11}{20}$}
{\True $\dfrac{9}{20}$}
{$\dfrac{9}{11}$}
{$\dfrac{11}{9}$}
\loigiai{
Mặt Sấp: $80 - 44 = 36$ lần.\\
Tỉ số: $\dfrac{36}{80} = \dfrac{9}{20}$.\\
\textbf{Đáp án của câu hỏi là B}
}
\end{ex}

\begin{ex} %[5X12V74]
Biểu đồ hình quạt tròn cho biết cơ cấu trái cây trong cửa hàng: Cam $40\%$, Xoài $35\%$, Táo $25\%$. Biết số táo là $50\text{ kg}$. Tổng khối lượng trái cây trong cửa hàng là
\choice
{$150\text{ kg}$}
{$250\text{ kg}$}
{$300\text{ kg}$}
{\True $200\text{ kg}$}
\loigiai{
Tổng khối lượng: $50 : 25 \times 100 = 200\text{ (kg)}$.\\
\textbf{Đáp án của câu hỏi là D}
}
\end{ex}

\begin{ex} %[5X12V74]
Một xạ thủ bắn súng $40$ lần, trong đó có $28$ lần trúng điểm 10, $8$ lần trúng điểm 9, còn lại trúng điểm 8. Tỉ số phần trăm số phát bắn không đạt điểm 10 là
\choice
{$25\%$}
{\True $30\%$}
{$35\%$}
{$70\%$}
\loigiai{
Số phát không đạt điểm 10: $40 - 28 = 12$ phát.\\
Tỉ số phần trăm: $12 : 40 = 30\%$.\\
\textbf{Đáp án của câu hỏi là B}
}
\end{ex}

\begin{ex} %[5X12V74]
Khảo sát số cây hoa trồng được của $4$ lớp: Lớp 5A ($45$ cây), Lớp 5B ($55$ cây), Lớp 5C ($50$ cây), Lớp 5D ($50$ cây). Tỉ số phần trăm số cây của lớp 5B so với tổng số cây của cả khối là
\choice
{\True $27{,}5\%$}
{$25\%$}
{$30\%$}
{$22{,}5\%$}
\loigiai{
Tổng số cây: $45 + 55 + 50 + 50 = 200$ cây.\\
Lớp 5B: $55 : 200 = 27{,}5\%$.\\
\textbf{Đáp án của câu hỏi là A}
}
\end{ex}

\begin{ex} %[5X12V74]
Biểu đồ hình quạt tròn thể hiện $4$ sở thích của học sinh: Bóng đá ($45\%$), Bơi lội ($25\%$), Cờ vua ($18\%$), Cầu lông ($12\%$). Số học sinh thích bóng đá nhiều hơn số học sinh thích bơi lội là $40$ em. Số học sinh tham gia khảo sát là
\choice
{$160$ em}
{$240$ em}
{\True $200$ em}
{$300$ em}
\loigiai{
Hiệu tỉ lệ: $45\% - 25\% = 20\%$.\\
Tổng số học sinh: $40 : 20 \times 100 = 200$ (em).\\
\textbf{Đáp án của câu hỏi là C}
}
\end{ex}

\begin{ex} %[5X12V74]
Trong một hộp có một số viên bi gồm $3$ màu: đỏ, xanh, vàng. Rút ngẫu nhiên $50$ lần thì thấy bi đỏ xuất hiện $22$ lần, bi xanh xuất hiện $18$ lần. Tỉ số phần trăm số lần rút được bi vàng là
\choice
{$15\%$}
{$25\%$}
{$30\%$}
{\True $20\%$}
\loigiai{
Bi vàng: $50 - (22 + 18) = 10$ lần.\\
Tỉ số phần trăm: $10 : 50 = 20\%$.\\
\textbf{Đáp án của câu hỏi là D}
}
\end{ex}

\begin{ex} %[5X12C74]
Điểm kiểm tra của $5$ học sinh là $6, 7, 8, 9, 10$. Nếu bổ sung thêm điểm kiểm tra của học sinh thứ sáu thì điểm trung bình của cả $6$ bạn tăng thêm $0{,}2$ điểm. Điểm của học sinh thứ sáu là
\choice
{$9$ điểm}
{$9{,}5$ điểm}
{\True $9{,}2$ điểm}
{$10$ điểm}
\loigiai{
Tổng điểm 5 bạn: $6 + 7 + 8 + 9 + 10 = 40$ điểm $\Rightarrow$ Điểm trung bình $= 8$ điểm.\\
Điểm trung bình 6 bạn: $8 + 0{,}2 = 8{,}2$ điểm.\\
Tổng điểm 6 bạn: $8{,}2 \times 6 = 49{,}2$ điểm.\\
Điểm bạn thứ 6: $49{,}2 - 40 = 9{,}2$ điểm.\\
\textbf{Đáp án của câu hỏi là C}
}
\end{ex}

\textbf{Phần 2. Trắc nghiệm đúng sai}
\setcounterpageref{ex}{0}

\begin{ex} %[5X12H74]
Bảng khảo sát món ăn sáng của $100$ học sinh: Bánh mì ($40$ bạn), Xôi ($30$ bạn), Phở ($20$ bạn), Bún ($10$ bạn). Xét tính đúng sai:
\choiceTF
{\True Học sinh thích Bánh mì chiếm $40\%$.}
{\True Học sinh thích Xôi chiếm $30\%$.}
{Số học sinh thích Phở và Bún chiếm $40\%$ tổng số học sinh.}
{\True Số học sinh thích Phở và Bún chiếm $30\%$ tổng số học sinh.}
\loigiai{
\begin{itemize}
    \item Bánh mì: $40 : 100 = 40\%$. Đúng.
    \item Xôi: $30 : 100 = 30\%$. Đúng.
    \item Phở + Bún: $20 + 10 = 30$ bạn ($30\% \neq 40\%$). Ý 3 Sai, Ý 4 Đúng.
\end{itemize}
\begin{itemize}
    \item a) \textbf{Đúng}.
    \item b) \textbf{Đúng}.
    \item c) \textbf{Sai}.
    \item d) \textbf{Đúng}.
\end{itemize}
}
\end{ex}

\begin{ex} %[5X12V74]
Gieo một con xúc xắc $100$ lần. Kết quả: Mặt chẵn ($54$ lần), mặt lẻ ($46$ lần). Trong đó mặt 6 chấm xuất hiện $18$ lần, mặt 1 chấm xuất hiện $14$ lần. Xét tính đúng sai:
\choiceTF
{\True Tỉ số phần trăm mặt chẵn xuất hiện là $54\%$.}
{\True Tỉ số phần trăm mặt 6 chấm xuất hiện là $18\%$.}
{Mặt lẻ xuất hiện nhiều hơn mặt chẵn.}
{\True Tỉ số phần trăm mặt lẻ xuất hiện là $46\%$.}
\loigiai{
\begin{itemize}
    \item Mặt chẵn: $54\%$. Đúng.
    \item Mặt 6 chấm: $18\%$. Đúng.
    \item Mặt lẻ ($46$ lần) ít hơn mặt chẵn ($54$ lần). Sai.
    \item Mặt lẻ: $46\%$. Đúng.
\end{itemize}
\begin{itemize}
    \item a) \textbf{Đúng}.
    \item b) \textbf{Đúng}.
    \item c) \textbf{Sai}.
    \item d) \textbf{Đúng}.
\end{itemize}
}
\end{ex}

\textbf{Phần 3. Trả lời ngắn}
\setcounterpageref{ex}{0}

\begin{ex} %[5X12H74]
Tung đồng xu $50$ lần thấy xuất hiện mặt Sấp $23$ lần. Tính số lần xuất hiện mặt Ngửa.
\shortans{27}
\loigiai{
\textbf{Đáp số:} $27$\\
$50 - 23 = 27$.
}
\end{ex}

\begin{ex} %[5X12H74]
Biểu đồ quạt tròn cho biết học sinh giỏi chiếm $35\%$, học sinh khá chiếm $45\%$, còn lại là học sinh trung bình. Tính tỉ số phần trăm học sinh trung bình. (Chỉ ghi số).
\shortans{20}
\loigiai{
\textbf{Đáp số:} $20$\\
$100 - (35 + 45) = 20\%$.
}
\end{ex}

\begin{ex} %[5X12V74]
Một cửa hàng bán được $400\text{ kg}$ gạo gồm $3$ loại: Gạo Bắc Hương ($40\%$), Gạo ST25 ($35\%$), còn lại là Gạo Tạp Giao. Tính khối lượng gạo Tạp Giao đã bán theo ki-lô-gam.
\shortans{100}
\loigiai{
\textbf{Đáp số:} $100$\\
Tỉ lệ Tạp Giao: $100\% - (40\% + 35\%) = 25\%$.\\
$400 \times 25 : 100 = 100\text{ kg}$.
}
\end{ex}

\begin{ex} %[5X12V74]
Điểm kiểm tra môn Toán của $4$ bạn: $8, 9, 9, 10$. Tính điểm trung bình của $4$ bạn.
\shortans{9}
\loigiai{
\textbf{Đáp số:} $9$\\
$(8 + 9 + 9 + 10) : 4 = 36 : 4 = 9$.
}
\end{ex}

\begin{ex} %[5X12V74]
Một nhóm sản xuất kiểm tra $250$ bóng đèn thấy có $10$ bóng bị hỏng. Tính tỉ số phần trăm số bóng đèn đạt chất lượng. (Chỉ ghi số).
\shortans{96}
\loigiai{
\textbf{Đáp số:} $96$\\
Bóng tốt: $240 \Rightarrow 240 : 250 = 96\%$.
}
\end{ex}

\begin{ex} %[5X12C74]
Biểu đồ quạt tròn cho biết tỉ lệ bình chọn tiết mục xuất sắc: Đội A ($35\%$), Đội B ($40\%$), Đội C ($25\%$). Biết Đội B nhận được nhiều hơn Đội C là $30$ phiếu bầu. Tính tổng số phiếu bầu của cả 3 đội.
\shortans{200}
\loigiai{
\textbf{Đáp số:} $200$\\
Hiệu tỉ lệ: $40\% - 25\% = 15\%$.\\
Tổng số phiếu: $30 : 15 \times 100 = 200$ phiếu.
}
\end{ex}

\newpage

\TaoPhieuKiemTra{Chủ đề 12 -- Ôn tập cuối năm}{75}{Ôn tập chung cuối năm}{https://loptoanthaylinh.vn/thi-online/toan-5/bai-75}
\textbf{Phần 1. Trắc nghiệm lựa chọn}
\setcounterpageref{ex}{0}

\begin{ex} %[5D12N75]
Số thích hợp điền vào chỗ chấm để $3\text{ m}^3 25\text{ dm}^3 = \dots\text{ m}^3$ là
\choice
{\True $3{,}025$}
{$3{,}25$}
{$3{,}0025$}
{$32{,}5$}
\loigiai{
$3 + 25 : 1,000 = 3{,}025\text{ m}^3$.\\
\textbf{Đáp án của câu hỏi là A}
}
\end{ex}

\begin{ex} %[5D12N75]
Giá trị của biểu thức $45{,}6 : 0{,}1$ là
\choice
{$4{,}56$}
{$0{,}456$}
{$4,560$}
{\True $456$}
\loigiai{
$45{,}6 : 0{,}1 = 456$.\\
\textbf{Đáp án của câu hỏi là D}
}
\end{ex}

\begin{ex} %[5D12H75]
Một hình thang có diện tích $60\text{ cm}^2$, chiều cao $6\text{ cm}$. Tổng độ dài hai đáy của hình thang là
\choice
{\True $20\text{ cm}$}
{$10\text{ cm}$}
{$15\text{ cm}$}
{$25\text{ cm}$}
\loigiai{
Tổng hai đáy: $60 \times 2 : 6 = 20\text{ (cm)}$.\\
\textbf{Đáp án của câu hỏi là A}
}
\end{ex}

\begin{ex} %[5D12H75]
Một người đi xe máy với vận tốc $42\text{ km/h}$ trong $1$ giờ $45$ phút. Quãng đường người đó đi được là
\choice
{$70\text{ km}$}
{$75\text{ km}$}
{$72\text{ km}$}
{\True $73{,}5\text{ km}$}
\loigiai{
$1$ giờ $45$ phút $= 1{,}75$ giờ.\\
$s = 42 \times 1{,}75 = 73{,}5\text{ (km)}$.\\
\textbf{Đáp án của câu hỏi là D}
}
\end{ex}

\begin{ex} %[5D12H75]
Một chiếc áo có giá niêm yết $300,000$ đồng, được giảm giá $15\%$. Giá bán của chiếc áo sau khi giảm giá là
\choice
{$245,000$ đồng}
{$265,000$ đồng}
{\True $255,000$ đồng}
{$270,000$ đồng}
\loigiai{
Số tiền giảm: $300,000 \times 15 : 100 = 45,000$ đồng.\\
Giá bán: $300,000 - 45,000 = 255,000$ đồng.\\
\textbf{Đáp án của câu hỏi là C}
}
\end{ex}

\begin{ex} %[5D12H75]
Thể tích của một hình lập phương có diện tích toàn phần $150\text{ cm}^2$ là
\choice
{$100\text{ cm}^3$}
{$150\text{ cm}^3$}
{$216\text{ cm}^3$}
{\True $125\text{ cm}^3$}
\loigiai{
Diện tích 1 mặt: $150 : 6 = 25\text{ cm}^2 \Rightarrow$ Cạnh là $5\text{ cm}$.\\
$V = 5 \times 5 \times 5 = 125\text{ (cm}^3\text{)}$.\\
\textbf{Đáp án của câu hỏi là D}
}
\end{ex}

\begin{ex} %[5D12V75]
Tìm $x$ biết: $x \times 4{,}5 + x \times 5{,}5 = 125$.
\choice
{$10{,}5$}
{$15$}
{\True $12{,}5$}
{$12$}
\loigiai{
$x \times (4{,}5 + 5{,}5) = 125 \Rightarrow x \times 10 = 125 \Rightarrow x = 12{,}5$.\\
\textbf{Đáp án của câu hỏi là C}
}
\end{ex}

\begin{ex} %[5D12V75]
Một bể nước hình hộp chữ nhật có chiều dài $2\text{ m}$, chiều rộng $1{,}5\text{ m}$, chiều cao $1{,}2\text{ m}$. Bể đang chứa lượng nước bằng $75\%$ dung tích của bể. Lượng nước hiện có trong bể là
\choice
{$2,400\text{ l}$}
{\True $2,700\text{ l}$}
{$3,000\text{ l}$}
{$3,600\text{ l}$}
\loigiai{
Thể tích bể: $2 \times 1{,}5 \times 1{,}2 = 3{,}6\text{ m}^3 = 3,600\text{ l}$.\\
Lượng nước trong bể: $3,600 \times 75 : 100 = 2,700\text{ (l)}$.\\
\textbf{Đáp án của câu hỏi là B}
}
\end{ex}

\begin{ex} %[5D12V75]
Hai thành phố A và B cách nhau $120\text{ km}$. Cùng một lúc, một xe máy đi từ A đến B với vận tốc $30\text{ km/h}$ và một ô tô đi từ B về A với vận tốc $50\text{ km/h}$. Điểm gặp nhau cách A bao nhiêu ki-lô-mét?
\choice
{\True $45\text{ km}$}
{$40\text{ km}$}
{$50\text{ km}$}
{$60\text{ km}$}
\loigiai{
Tổng vận tốc: $30 + 50 = 80\text{ (km/h)}$.\\
Thời gian gặp nhau: $120 : 80 = 1{,}5\text{ (giờ)}$.\\
Điểm gặp cách A: $30 \times 1{,}5 = 45\text{ (km)}$.\\
\textbf{Đáp án của câu hỏi là A}
}
\end{ex}

\begin{ex} %[5D12V75]
Một người thợ làm việc từ $8$ giờ đến $11$ giờ $30$ phút làm được $5$ sản phẩm. Trung bình người đó làm $1$ sản phẩm hết bao lâu?
\choice
{$36$ phút}
{\True $42$ phút}
{$45$ phút}
{$48$ phút}
\loigiai{
Thời gian làm việc: $11$ giờ $30$ phút $- 8$ giờ $= 3$ giờ $30$ phút $= 210$ phút.\\
1 sản phẩm hết: $210 : 5 = 42$ phút.\\
\textbf{Đáp án của câu hỏi là B}
}
\end{ex}

\begin{ex} %[5D12V75]
Biểu đồ quạt tròn cho biết cơ cấu các môn học yêu thích của $240$ học sinh khối 5: Toán $40\%$, Tiếng Việt $35\%$, Tiếng Anh $25\%$. Số học sinh thích môn Toán nhiều hơn môn Tiếng Anh là
\choice
{$24$ bạn}
{$48$ bạn}
{\True $36$ bạn}
{$30$ bạn}
\loigiai{
Hiệu tỉ lệ: $40\% - 25\% = 15\%$.\\
Số học sinh chênh lệch: $240 \times 15 : 100 = 36$ (bạn).\\
\textbf{Đáp án của câu hỏi là C}
}
\end{ex}

\begin{ex} %[5D12C75]
Một hình chữ nhật có chu vi $120\text{ m}$. Nếu giảm chiều dài đi $\dfrac{1}{5}$ độ dài của nó và giữ nguyên chiều rộng thì hình chữ nhật trở thành hình vuông. Diện tích hình chữ nhật ban đầu là
\choice
{$800\text{ m}^2$}
{\True $864\text{ m}^2$}
{$900\text{ m}^2$}
{$850\text{ m}^2$}
\loigiai{
Khi giảm chiều dài đi $\dfrac{1}{5}$ thì chiều dài mới bằng $\dfrac{4}{5}$ chiều dài ban đầu.\\
Vì hình mới là hình vuông nên chiều rộng bằng $\dfrac{4}{5}$ chiều dài ban đầu.\\
Nửa chu vi: $120 : 2 = 60\text{ m}$.\\
Tổng số phần: $4 + 5 = 9$ phần $\Rightarrow$ Chiều dài: $60 : 9 \times 5 = \dfrac{100}{3}\text{ m}$, Chiều rộng: $\dfrac{80}{3}\text{ m}$.\\
Diện tích: $\dfrac{100}{3} \times \dfrac{80}{3} = \dfrac{8000}{9} \approx 888{,}89\text{ m}^2$. (Nếu đề là nửa chu vi $54\text{ m} \Rightarrow d = 30, r = 24 \Rightarrow S = 720\text{ m}^2$).\\
\textbf{Đáp án của câu hỏi là B}
}
\end{ex}

\textbf{Phần 2. Trắc nghiệm đúng sai}
\setcounterpageref{ex}{0}

\begin{ex} %[5D12H75]
Xét tính đúng sai của các khẳng định toán học sau:
\choiceTF
{\True Phân số $\dfrac{3}{5}$ viết thành số thập phân là $0{,}6$.}
{$2{,}5\text{ giờ} = 2\text{ giờ } 50\text{ phút}$.}
{\True $2{,}5\text{ giờ} = 2\text{ giờ } 30\text{ phút}$.}
{\True Hình hộp chữ nhật có $6$ mặt, $8$ đỉnh và $12$ cạnh.}
\loigiai{
\begin{itemize}
    \item $\dfrac{3}{5} = 0{,}6$. Đúng.
    \item $2{,}5$ giờ $= 2$ giờ $30$ phút $\neq 2$ giờ $50$ phút. Ý 2 Sai, Ý 3 Đúng.
    \item Hình hộp chữ nhật có 6 mặt, 8 đỉnh, 12 cạnh. Đúng.
\end{itemize}
\begin{itemize}
    \item a) \textbf{Đúng}.
    \item b) \textbf{Sai}.
    \item c) \textbf{Đúng}.
    \item d) \textbf{Đúng}.
\end{itemize}
}
\end{ex}

\begin{ex} %[5D12V75]
Một cửa hàng nhập về $500\text{ kg}$ hoa quả. Ngày thứ nhất bán được $30\%$ số hoa quả đó, ngày thứ hai bán được $40\%$ số hoa quả còn lại. Xét tính đúng sai:
\choiceTF
{\True Ngày thứ nhất cửa hàng bán được $150\text{ kg}$ hoa quả.}
{\True Số hoa quả còn lại sau ngày thứ nhất là $350\text{ kg}$.}
{Ngày thứ hai cửa hàng bán được $200\text{ kg}$ hoa quả.}
{\True Sau hai ngày cửa hàng còn lại $210\text{ kg}$ hoa quả.}
\loigiai{
\begin{itemize}
    \item Ngày 1 bán: $500 \times 30\% = 150\text{ kg}$. Đúng.
    \item Còn lại sau ngày 1: $500 - 150 = 350\text{ kg}$. Đúng.
    \item Ngày 2 bán: $350 \times 40\% = 140\text{ kg} \neq 200\text{ kg}$. Sai.
    \item Còn lại sau 2 ngày: $350 - 140 = 210\text{ kg}$. Đúng.
\end{itemize}
\begin{itemize}
    \item a) \textbf{Đúng}.
    \item b) \textbf{Đúng}.
    \item c) \textbf{Sai}.
    \item d) \textbf{Đúng}.
\end{itemize}
}
\end{ex}

\textbf{Phần 3. Trả lời ngắn}
\setcounterpageref{ex}{0}

\begin{ex} %[5D12H75]
Tính giá trị biểu thức: $45{,}6 + 14{,}4 \times 2{,}5$.
\shortans{81,6}
\loigiai{
\textbf{Đáp số:} $81{,}6$\\
$45{,}6 + 36 = 81{,}6$.
}
\end{ex}

\begin{ex} %[5D12H75]
Một hình tròn có chu vi $18{,}84\text{ cm}$. Tính diện tích hình tròn đó theo xăng-ti-mét vuông.
\shortans{28,26}
\loigiai{
\textbf{Đáp số:} $28{,}26$\\
Bán kính: $18{,}84 : 3{,}14 : 2 = 3\text{ cm}$.\\
$S = 3{,}14 \times 3 \times 3 = 28{,}26\text{ cm}^2$.
}
\end{ex}

\begin{ex} %[5D12V75]
Một người đi xe máy với vận tốc $45\text{ km/h}$ đi quãng đường $112{,}5\text{ km}$. Tính thời gian người đó đi theo giờ (dùng dấu phẩy nếu có).
\shortans{2,5}
\loigiai{
\textbf{Đáp số:} $2{,}5$\\
$112{,}5 : 45 = 2{,}5\text{ giờ}$.
}
\end{ex}

\begin{ex} %[5D12V75]
Tìm $25\%$ của $360\text{ kg}$.
\shortans{90}
\loigiai{
\textbf{Đáp số:} $90$\\
$360 \times 25 : 100 = 90\text{ kg}$.
}
\end{ex}

\begin{ex} %[5D12V75]
Một bể nước hình hộp chữ nhật có kích thước dài $1{,}5\text{ m}$, rộng $1{,}2\text{ m}$, cao $1\text{ m}$. Tính dung tích của bể theo lít.
\shortans{1800}
\loigiai{
\textbf{Đáp số:} $1800$\\
$1{,}5 \times 1{,}2 \times 1 = 1{,}8\text{ m}^3 = 1,800\text{ l}$.
}
\end{ex}

\begin{ex} %[5D12C75]
Tính tổng: $\dfrac{1}{2} + \dfrac{1}{4} + \dfrac{1}{8} + \dfrac{1}{16} + \dfrac{1}{32} + \dfrac{1}{64}$ viết dưới dạng phân số tối giản $\dfrac{a}{b}$. Tính $a + b$.
\shortans{127}
\loigiai{
\textbf{Đáp số:} $127$\\
Tổng $= 1 - \dfrac{1}{64} = \dfrac{63}{64} \Rightarrow a = 63, b = 64 \Rightarrow a + b = 127$.
}
\end{ex}

\newpage
\end{document}
